% Generated; edit blueprint/src instead.
\documentclass[10pt,oneside]{book}
\usepackage[margin=1in]{geometry}
\usepackage{amsmath,amssymb,amsthm}
\newtheorem{theorem}{Theorem}[chapter]
\newtheorem{lemma}[theorem]{Lemma}
\newtheorem{definition}[theorem]{Definition}
\newtheorem{proposition}[theorem]{Proposition}
\newtheorem{corollary}[theorem]{Corollary}
\newcommand{\R}{\mathbb R}
\newcommand{\N}{\mathbb N}
\newcommand{\source}[1]{\par\smallskip\noindent\textit{Manuscript label:} \texttt{\detokenize{#1}}.\par}

\usepackage{xcolor,hyperref,longtable,booktabs,enumitem,array}
\hypersetup{colorlinks=true,linkcolor=blue!50!black,urlcolor=blue!50!black}
\newcommand{\lean}[1]{\par\smallskip\noindent{\footnotesize\textit{Lean:} \nolinkurl{#1}}\par}
\newcommand{\leanok}{\par\noindent{\footnotesize\color{green!35!black}Checked at the scope stated here.}\par}
\newcommand{\uses}[1]{\par\noindent{\footnotesize\textit{Blueprint inputs:} \nolinkurl{#1}}\par}
\newcommand{\notready}{\par\noindent{\footnotesize\color{red!70!black}OPEN: no complete Lean proof.}\par}

\usepackage{lmodern,microtype,xurl}
\setlength{\emergencystretch}{4em}
\setlength{\parskip}{.35em}
\setlength{\parindent}{0pt}
\renewcommand{\source}[1]{\par\smallskip\noindent\textit{Manuscript label:} \nolinkurl{#1}.\par}
\title{Oberwolfach formalisation\\Blueprint and partial verification}
\author{Formalisation working document}
\date{20 October 2026}
\begin{document}
\hypersetup{pageanchor=false}
\maketitle
\hypersetup{pageanchor=true}
\tableofcontents
\chapter*{Verification status}
This is a partial Lean formalisation of the working manuscript
\emph{Sharp Dirac thresholds for the generalised Oberwolfach problem}.
It is not a Lean certificate of Theorem~1.1 or of the whole paper.
The main theorem is represented by an unproved proposition.
No placeholder proof or new axiom is used in the checked modules.

The source snapshot has SHA-256
\begin{center}\ttfamily\small
ffe45b31bcfbd933828da22a9becb6aea23a64\\
2d5ffe7abf09d29f45d9d8a201.
\end{center}
The two lines above concatenate to one checksum (excluding the final period).
The Lean sources, pinned dependencies, raw verification logs, statement
inventory and reproducible verification commands accompany this document.
The complete manuscript was not formally verified in producing this blueprint.

The formalisation keeps the original objective: a proof of the main theorem
with all dependencies discharged. The completed results below are intermediate
work towards that objective. A conditional reduction is identified as such;
its hypotheses remain obligations even when its proof is accepted by Lean.
The file \texttt{CheckMain.lean} rejects a full-certificate release until an
unconditional theorem of the exact target type exists.

\chapter{The target and its mathematical representation}
\begin{definition}[Finite graphs and prescribed decompositions]
\label{bp:graphs}
\lean{Oberwolfach.EdgePartition, Oberwolfach.HasFactorDecomposition}
\leanok
All ordinary graphs are simple graphs on $\{0,\ldots,n-1\}$.
A family $(H_i)_{i<q}$ partitions the edges of $G$ if, for every $u,v$,
\[
 G(u,v)\iff \exists i<q\ H_i(u,v),\qquad
 H_i(u,v)\wedge H_j(u,v)\Longrightarrow i=j.
\]
A decomposition into $(F_i)_{i<q}$ additionally supplies an isomorphism
$H_i\cong F_i$ for every $i$. Thus the definition records multiplicities in
the prescribed list, spanning vertex sets and exact edge coverage.
\end{definition}

\begin{definition}[Ordinary cycle thresholds]
\label{bp:thresholds}
\lean{Oberwolfach.IsCycleCopy, Oberwolfach.CycleDivisible, Oberwolfach.AdmissibleCycleBound, Oberwolfach.cycleThreshold, Oberwolfach.allowedThreshold}
\uses{bp:graphs}
\leanok
An embedded $r$-cycle is specified by an injection from $\{0,\ldots,r-1\}$
whose consecutive pairs, with indices taken modulo $r$, are precisely its
edges. Other ambient vertices are isolated. A $C_r$-decomposition partitions
edges into such embedded cycles.
A graph is $C_r$-divisible if all degrees are even and $r$ divides its number
of edges. Define
\[
 \begin{gathered}
 \delta_r=\inf\{d\in[0,1]:\exists N\ \forall n\ge N\ \forall G,\\
 \ G\text{ is }C_r\text{-divisible},\ \delta(G)\ge dn
 \Longrightarrow G\text{ has a }C_r\text{-decomposition}\},\\
 \delta_S=\sup_{r\in S}\delta_r.
 \end{gathered}
\]
\end{definition}

\begin{lemma}[Elementary consequences of the infimum definition]
\label{bp:threshold-bounds}
\lean{Oberwolfach.one_admissible, Oberwolfach.cycleThreshold_nonneg, Oberwolfach.cycleThreshold_le_one, Oberwolfach.cycleThreshold_positive_margin}
\uses{bp:thresholds}
\leanok
The admissible set contains $1$, and $0\le\delta_r\le1$.
For every $\alpha>0$, all sufficiently large $C_r$-divisible graphs of
minimum degree at least $(\delta_r+\alpha)n$ have a $C_r$-decomposition.
\end{lemma}
\begin{proof}\leanok
A nonempty $n$-vertex simple graph has each degree below $n$, so the assertion
at degree coefficient $1$ is vacuous. All admissible coefficients are
nonnegative. The defining property of an infimum gives an admissible
$d<\delta_r+\alpha$; its order threshold also works for the larger degree
bound.
\end{proof}

\begin{lemma}[Finite local lemma on the original coordinate space]
\label{bp:finite-local-lemma}
\lean{Oberwolfach.finiteLaw_localLemma, Oberwolfach.finiteLaw_coordinate_localLemma, Oberwolfach.finiteLaw_disjoint_supported_events}
\uses{bp:common-event-probability}
\leanok
For finitely many bad events $A_i$, suppose $A_i$ is independent of the
avoidance of any set of events outside $N_i\cup\{i\}$. If
$0\le x_i<1$ and
$\Pr(A_i)\le x_i\prod_{j\in N_i}(1-x_j)$, then the probability of
avoiding every bad event is positive. For a finite product law, disjoint
specified coordinate supports imply the required independence.
\end{lemma}
\begin{proof}\leanok
Induction on the conditioning set bounds the probability of a bad event
intersected with that set's avoidance. Separate neighbours from
non-neighbours and use the induction hypothesis to bound the denominator
without dividing by an avoidance probability. Multiply the resulting
positive factors. Independence for coordinate supports is proved for the
given product measure.
\end{proof}

\begin{theorem}[Finite-product convex-distance inequality]
\label{bp:finite-talagrand}
\lean{Oberwolfach.finiteLaw_convexDistance_mgf, Oberwolfach.finiteLaw_convexDistance_tail, Oberwolfach.convexDistanceSq_value_certificate}
\uses{bp:categorical-concentration}
\leanok
For a configuration $x$ and a nonempty set $A$, let $D_A(x)$ be the
minimum of $\sum_i z_i^2$ over the convex hull of the vectors
$(\mathbf1_{x_i\ne y_i})_i$, $y\in A$. For any finite product law,
\[
 \mathbb E\exp(D_A/4)\le\Pr(A)^{-1}\qquad(\Pr(A)>0),
 \qquad
 \Pr(A)\Pr(D_A\ge t)\le e^{-t/4}.
\]
If coordinate certificate weights lie in $[0,c]$, have sum at most $D$,
and changing from $x$ to any $y\in A$ incurs certificate weight at least
$t\ge0$, then $t^2\le cD\,D_A(x)$.
\end{theorem}
\begin{proof}\leanok
Compactness of the finite convex hull gives the minimum. Project away one
coordinate and mix a fibre minimizer with a projected minimizer. The new
coordinate contributes at most $(1-\theta)^2$. Convexity of the exponential
gives the required interpolation of finite exponential moments. The scalar
choice of $\theta$ yields the fibre bound
$(2-\Pr(A_a)/\Pr(B))/\Pr(B)$, where $B$ is the projection.
Fibres of probability zero are treated separately. Exact product
disintegration and induction on the finite coordinate set prove the moment
bound; exponential Markov gives the tail bound. The certificate estimate
follows from a linear half-space containing the convex hull and the
Cauchy--Schwarz inequality.
\end{proof}

\begin{lemma}[Certificates and one common local event]
\label{bp:colouring-certificate-concentration}
\lean{Oberwolfach.finiteLaw_exists_median, Oberwolfach.finiteLaw_bounded_certificate_tail, Oberwolfach.finiteLaw_coordinate_localLemma_degree, Oberwolfach.local_certificate_common_event}
\uses{bp:finite-talagrand,bp:finite-local-lemma}
\leanok
Suppose $0\le R\le M$ has nonnegative coordinate certificate weights,
each at most $c>0$, with total at most $D>0$. Put
$p_t=4\exp(-t^2/(4cD))$. If $t\ge0$ and $Mp_t\le t$, then
$\Pr(|R-\mathbb ER|\ge3t)\le p_t$.
For a finite family of such variables, suppose each uses at most $r$
coordinates and each coordinate occurs in at most $s$ supports, with
$r,s\ge1$. If every tail bound is at most $1/(4rs)$, there is one
positive-mass outcome satisfying all the strict deviation bounds.
\end{lemma}
\begin{proof}\leanok
The convex-distance inequality first gives a median tail. Truncating the
absolute deviation bounds the distance between the mean and median by
$2t$. Intersecting coordinate supports give a dependency degree at most
$rs$; the finite local lemma then supplies one common outcome.
\end{proof}

\begin{lemma}[Concrete tentative-colouring counts]
\label{bp:tentative-colouring-counts}
\lean{Oberwolfach.coinRetainedColour_probability, Oberwolfach.equalizedColour_survival_probability, Oberwolfach.joint_coinRetention_probability, Oberwolfach.coinConflict_certificate, Oberwolfach.nibbleColourCount_certificate, Oberwolfach.nibbleAvailableColours_card}
\uses{bp:colouring-certificate-concentration}
\leanok
A vertex tentatively chooses a colour and a colour-dependent retention
coin. It retains the colour if the coin succeeds and no neighbour chose
that colour. The retained colours form a proper partial colouring.
All one-root and independent-set joint retention probabilities are computed
in the original product law. For lists of size $l$, let $N(v,c)$ be the
number of neighbours with $c$ in their lists. If $N(v,c)\le K$ whenever
$c$ belongs to the list at $v$, the coin probability
$(1-1/l)^{K-N(v,c)}$ makes the retention probability at every vertex
exactly $(1-1/l)^K$.

For a clique $T$, the number of rejected vertices has coordinate
certificate weights at most $2$ and total at most $3|T|$. For a family
of $k$ vertex sets and a fixed list $L$, let $M(a,b)$ count colours in
$L$ assigned in at least $a$ sets and rejected in at least $b$ sets.
This count has certificate weights at most $1$ and total at most
$3k|L|$. If the sets are cliques, the number of colours left after
removing all retained colours is exactly
$\sum_{a=0}^{k}(M(a,a)-M(a+1,a))$.
\end{lemma}
\begin{proof}\leanok
Coordinatewise probabilities give the retention identities. Certificates
use actual conflicting neighbours and the rejection coin at the assigned
vertex. Distinct colours cannot use the same witness coordinate. In a
clique, a colour assigned at several vertices is rejected at all of them.
The final identity is a finite telescoping sum over the numbers of sets
where the colour was assigned and rejected. The full iteration using these results is proved in Node~\ref{bp:near-optimal-colouring}.
\end{proof}

\begin{lemma}[Concentration of the remaining colour degrees]
\label{bp:remaining-colour-degrees}
\lean{Oberwolfach.neighbour_colour_witnesses, Oberwolfach.hypergraph_neighbourColourCount_certificate, Oberwolfach.remainingColourVertices_card, Oberwolfach.remainingColourVertices_concentration}
\uses{bp:colouring-certificate-concentration,bp:tentative-colouring-counts}
\leanok
Let $U$ be a clique in the line graph of a hypergraph of rank at most $k$
and codegree at most $B$. For each edge in $U$, consider its resource
cliques other than the clique defining $U$. For a fixed colour, count the
uncoloured edges satisfying lower bounds on the numbers of these cliques
where that colour was assigned and rejected. The coordinate certificates
have maximum weight at most $2+(k+3)kB$ and total weight at most
$(3+3k)|U|$.

The actual remaining colour degree is the alternating sum of the diagonal
and next-diagonal counts. The formal result gives a mean-centred tail
bound by applying the proved certificate inequality to these counts.
\end{lemma}
\begin{proof}\leanok
Choose a maximal matching on rejected vertices of a fixed colour. Paired
vertices witness each other's rejection; the unmatched vertices are
independent. The rank bound controls the number of unmatched neighbours
that can use one witness. Each direct witness is needed by at most $kB$
counted edges, by the codegree bound and exclusion of the defining resource.
Composing these witnesses gives the displayed certificate weights.
The exact alternating-sum identity and the finite union bound give the
concentration estimate in the original product probability space.
\end{proof}

\begin{lemma}[Joint retention and overlapping resource cliques]
\label{bp:joint-retention-calculations}
\lean{Oberwolfach.mixed_joint_coinRetention_probability, Oberwolfach.finiteLaw_avoidance_expansion, Oberwolfach.finiteLaw_clique_restriction_loss, Oberwolfach.disjoint_clique_joint_retention, Oberwolfach.hypergraph_colour_common_neighbours, Oberwolfach.restricted_hypergraph_clique_cross_degree}
\uses{bp:tentative-colouring-counts}
\leanok
For properly prescribed root colours, the exact joint retention probability
is the product of the root probabilities and the probabilities that each
remaining vertex avoids the set of colours prescribed at its root
neighbours. Different colours forbidden at one vertex are not treated as
independent events. For a common root colour, the exact overlap factor
and its two-sided bound are proved.

The joint event that every clique retains a fixed colour expands over
choices of its unique retained vertices. Removing pair intersections has
an explicitly bounded probability cost. For disjoint cliques, the
joint-retention estimate includes both the overlap error on nonconflicting
tuples and the number of conflicting tuples. Hypergraph rank and codegree
bounds imply the neighbourhood bounds required in this estimate.
\end{lemma}
\begin{proof}\leanok
Compute coordinate probabilities under the original colour-and-coin law.
Expand clique events using uniqueness of a retained vertex in a clique.
Conflicting tuples have probability zero; count them by fixing all but
one coordinate. Apply the exact same-colour joint formula to the other
tuples. Inclusion--exclusion is proved without an independence assumption.
These calculations do not yet supply the complete iterative colouring
argument or the weighted matching theorem.
\end{proof}

\begin{lemma}[Random partitions and selection of supplied colour classes]
\label{bp:matching-colour-selection}
\lean{Oberwolfach.hypergraphColorPartition_exists, Oberwolfach.weightedColor_common_upper, Oberwolfach.colorClassLottery_matching}
\uses{bp:common-event-probability,bp:categorical-concentration}
\leanok
Under their explicit finite tail-sum inequalities, the independently chosen
vertex and edge partitions simultaneously satisfy the stated degree and
fixed weighted-sum estimates. Given proper edge-colourings within the
resulting vertex-disjoint parts, independently choosing one colour class
in each part produces a matching and satisfies the supplied finite family
of weighted estimates. Existence of the required near-optimal proper
edge-colourings is not proved by this node.
\end{lemma}
\begin{proof}\leanok
Compute expectations in the original finite product law. Root-omitted
degree sums are controlled before exposing the root's colour.
For weighted edge partitions use the multiplicative exponential bound,
retaining the factor equal to the number of colours. A common-event
estimate applies to all fixed tests. Supplied proper colour classes and
disjoint vertex parts give a matching; bounded differences controls its
weighted sums.
\end{proof}

\begin{definition}[The two target propositions]
\label{bp:targets}
\lean{Oberwolfach.AllowedCycleLengths, Oberwolfach.DecompositionAt, Oberwolfach.FixedSetTheorem, Oberwolfach.MainTheorem}
\uses{bp:graphs,bp:thresholds}
\leanok
For a $2$-regular graph $F$, allowed component lengths are encoded by
\[
 \forall v,\quad |\{w:v\text{ is connected to }w\text{ in }F\}|\in S.
\]
Write $\mathcal D(S,\alpha,n)$ for the conclusion of the main theorem at
these parameters, universally quantified over the even degree $D$, the
family indexed by $\{0,\ldots,D/2-1\}$ and the $D$-regular host.
The fixed-set target is
\[
 \forall S\ne\varnothing\ \forall\alpha>0\ \exists N\ \forall n\ge N,
 \quad\mathcal D(S,\alpha,n),
\]
and the uniform target is
\[
 \forall\alpha>0\ \exists N\ \forall S\ne\varnothing\ \forall n\ge N,
 \quad\mathcal D(S,\alpha,n).
\]
In both formulas, $S\subseteq\{3,4,\ldots\}$.
These are definitions of propositions, not proofs of either proposition.
\end{definition}

\begin{theorem}[Main theorem: open target]
\label{bp:main}
\lean{Oberwolfach.MainTheorem}
\uses{bp:targets,bp:uniform,bp:fixed,bp:factor-embeddings-exhaust-host}
\notready
For every $\alpha>0$, there is an order threshold depending only on
$\alpha$ such that every sufficiently large $n$-vertex $D$-regular graph,
where $D$ is even and $D\ge(\delta_S+\alpha)n$, decomposes into every prescribed family of
$D/2$ spanning $2$-factors with component lengths in $S$.
\source{thm:new:prescribed-factor}
\end{theorem}

\chapter{A checked conditional reduction to the fixed-set theorem}
\begin{definition}[A common tail of allowed lengths]
\label{bp:tail}
\lean{Oberwolfach.SmallThresholdTail}
\leanok
For each $\alpha>0$, choose a finite $E\subseteq\{3,4,\ldots\}$ and a
nonempty $T\subseteq\{3,4,\ldots\}$ such that every $r\ge3$ outside $E$
belongs to $T$, and $\delta_r\le1/2+\alpha/2$ for $r\in T$.
\end{definition}

\begin{theorem}[Uniformisation from the fixed-set theorem alone]
\label{bp:uniform}
\lean{Oberwolfach.uniform_of_fixed_set_grid, Oberwolfach.allowedThreshold_nonneg, Oberwolfach.allowedThreshold_le_one}
\uses{bp:targets,bp:threshold-bounds}
\leanok
The fixed-set target implies the uniform target.
\end{theorem}
\begin{proof}\leanok
Fix $\alpha>0$ and put $b=\alpha/2$ and $M=\lceil1/b\rceil$.
For each integer $0\le q\le M$, let
\[
 T_q=\{r\ge3:\delta_r\le qb\}.
\]
If $T_q$ is nonempty, take its fixed-set order threshold at margin $b$;
otherwise use zero. Take the maximum before choosing $S$.
For a later nonempty $S$, put $\theta=\delta_S$ and
$q=\lceil\theta/b\rceil$. The proved bounds $0\le\theta\le1$ give
$q\le M$ and $\theta\le qb\le\theta+b$. Consequently
$S\subseteq T_q$ and $\delta_{T_q}\le qb\le\delta_S+b$.
The degree assumption for $S$ at margin $\alpha$ implies that for $T_q$
at margin $b$, without changing any prescribed factor.
This reduction does not require an estimate as the cycle length tends
to infinity. The fixed-set target remains unproved.
\end{proof}

\begin{theorem}[The manuscript's alternative common-tail reduction]
\label{bp:uniform-tail}
\lean{Oberwolfach.uniform_of_fixed_set, Oberwolfach.uniform_from_fixed_and_cofinite}
\uses{bp:targets,bp:threshold-bounds,bp:tail}
\leanok
Assume the fixed-set target, the bounds $\delta_r\ge1/2$ for $r\ge3$,
and the common-tail property. Then the uniform target holds.
It suffices for the tail property that, for every $\alpha>0$, all large
$r$ satisfy $\delta_r\le1/2+\alpha/2$.
\end{theorem}
\begin{proof}\leanok
For every $B\subseteq E$, take a fixed-set threshold for $T\cup B$ at margin
$\alpha/2$. Take the maximum of these finitely many integers before choosing
$S$. For any later $S$, put $B=S\cap E$ and $S^*=T\cup B$.
Then $S\subseteq S^*$ and
\[
 \delta_{S^*}\le\delta_S+\alpha/2.
\]
The degree hypothesis at margin $\alpha$ for $S$ therefore implies the
hypothesis at margin $\alpha/2$ for $S^*$. The prescribed factors are
unchanged. The Lean proof explicitly constructs the finite maximum and
checks the supremum inequalities and quantifier order.
\end{proof}

The fixed-set result and the lower and cofinite upper threshold estimates
are hypotheses of the common-tail reduction. The grid reduction above
removes the latter two hypotheses from the uniformisation step. Neither
reduction proves the fixed-set theorem.

\chapter{Checked finite stability arguments}
\begin{definition}[Common neighbourhoods on a fixed domain]
\label{bp:common}
\lean{Oberwolfach.common, Oberwolfach.changed}
\leanok
For finite $A,J$ and a family of finite sets $(N_j)$, put
$C(A,J,N)=\{v\in A:\forall j\in J,\ v\in N_j\}$.
The empty intersection is $A$. Put
$\operatorname{changed}(X,Y)=(X\setminus Y)\cup(Y\setminus X)$.
\end{definition}

\begin{lemma}[Neighbourhood stability and arbitrary later deletion]
\label{bp:neighbourhood-stability}
\lean{Oberwolfach.neighbourhood_stability_with_deletion}
\uses{bp:common}
\leanok
If $|N_j\triangle N'_j|\le b$ for every $j\in J$ and
$||C(A,J,N)|-\mu|\le e$, then
\[
 \big||C(A,J,N')|-\mu\big|\le e+|J|b,
 \qquad
 |C(A,J,N')\setminus D|\ge\mu-e-|J|b-|D|
\]
for every finite $D$, including a set chosen after $N,N'$.
\source{lem:pr:rt1-stability}
\end{lemma}
\begin{proof}\leanok
The symmetric difference of the two common neighbourhoods is contained
in $\bigcup_{j\in J}(N_j\triangle N'_j)$. Bound this union by the sum of
its sizes and apply the triangle inequality. Deleting $D$ removes at
most $|D|$ vertices. There is no sampling factor in this deterministic
loss and no change of the ambient normalization.
\end{proof}

\begin{definition}[Fixed weighted products]
\label{bp:patterns}
\lean{Oberwolfach.patternValue}
\leanok
For a finite coordinate set $F$, states $\rho_e$ and required states $s_e$,
put $p_F(\rho)=\prod_{e\in F}\mathbf1_{\{\rho_e=s_e\}}$.
The state type may depend on $e$. The coordinate universe itself need
not be finite, since all tested coordinate sets are finite.
Required states are supplied as total functions; in the manuscript's
formulation, partial required states can be extended outside $F$ using
the assumed nonempty state spaces. Their values outside $F$ are unused.
\end{definition}

\begin{lemma}[Changes of weighted products]
\label{bp:weighted-stability}
\lean{Oberwolfach.weighted_product_stability}
\uses{bp:patterns}
\leanok
Let $w_i\ge0$ and finite candidate coordinate sets $F_i$ be fixed.
If $\sigma_e=\tau_e$ whenever $e\in F_i\setminus D$, for every tested
candidate $i$, then
\[
 \left|\sum_iw_ip_{F_i}(\tau)-\sum_iw_ip_{F_i}(\sigma)\right|
 \le\sum_{e\in D}\sum_{i:e\in F_i}w_i.
\]
\source{lem:pr:wi-stability}
\end{lemma}
\begin{proof}\leanok
Each product belongs to $\{0,1\}$. It is unchanged if $F_i\cap D$ is
empty; otherwise its change has absolute value at most one, which is
at most $|F_i\cap D|$. Multiply by $w_i$, sum, use the triangle
inequality and interchange the two finite sums.
\end{proof}

\begin{lemma}[Omitted-root products]
\label{bp:omitted}
\lean{Oberwolfach.omitted_root_stability}
\uses{bp:weighted-stability}
\leanok
For fixed omitted coordinates $R$, replace each $F_i$ by $F_i\setminus R$.
The same estimate holds with right-hand side
$\sum_{e\in D}\sum_{i:e\in F_i\setminus R}w_i$.
The two assignments may differ arbitrarily on $R$, and neither is
conditioned on root presence. All weights remain fixed.
\end{lemma}
\begin{proof}\leanok
Apply the preceding theorem to the finite sets $F_i\setminus R$.
\end{proof}

\chapter{Checked transfers used in iteration}
\begin{lemma}[Local-scale conversion]
\label{bp:scale}
\lean{Oberwolfach.local_scale_error}
\leanok
Suppose $\epsilon>0$, $C,\sigma,n\ge0$,
$n\ge\epsilon m/2$ and $2C\sigma/\epsilon\le\delta$.
Then $\delta n+C\sigma m\le2\delta n$.
\end{lemma}
\begin{proof}\leanok
Multiply the size inequality by $C\sigma$ and the margin inequality
by $\epsilon n$. Both multipliers are nonnegative. These give
$C\sigma m\le\delta n$.
\end{proof}

\begin{lemma}[Protected restriction]
\label{bp:restriction}
\lean{Oberwolfach.protected_restriction}
\leanok
If $C\cap T=C_{\mathrm{next}}$ and $E\cap T=\varnothing$, then
$(C\setminus E)\cap T=C_{\mathrm{next}}$.
\end{lemma}
\begin{proof}\leanok
Check membership of an arbitrary element. This applies to sets of
typed occurrences as well as sets of ordinary edges; it does not
identify the two representations.
\end{proof}

\begin{lemma}[Large-cut normalization: the scalar step]
\label{bp:cuts}
\lean{Oberwolfach.large_cut_stability_arithmetic}
\leanok
Let $a_1,a_2>0$ and $x,y,b\ge0$. If $e\ge h\min(a_2x,a_1y)$ and
$e'\ge e-b\min(x,y)$, then
\[
 e'\ge\left(h-\frac b{\min(a_1,a_2)}\right)\min(a_2x,a_1y).
\]
\source{lem:pr:cut-stability}
\end{lemma}
\begin{proof}\leanok
Use
$\min(a_1,a_2)\min(x,y)\le\min(a_2x,a_1y)$.
Only this scalar implication is formalised here. The graph-theoretic
premise counting the deleted edges remains to be connected to the
maximum-degree hypothesis in Lean.
\end{proof}

\begin{lemma}[One common event and later choices]
\label{bp:common-event}
\lean{Oberwolfach.common_event_adaptive_choice}
\leanok
Let $E\subseteq\Omega$ be one event. Suppose
$\forall\omega\in E\ \forall x,\ A(\omega,x)\Rightarrow B(\omega,x)$.
If $c:\Omega\to X$ satisfies $A(\omega,c(\omega))$ for every $\omega\in E$,
then $E\subseteq\{\omega:B(\omega,c(\omega))\}$.
\end{lemma}
\begin{proof}\leanok
Substitute $x=c(\omega)$ on the same event. This is a deterministic
logical implication. It does not establish the probability of $E$,
measurability or the manuscript's sampling estimates.
\end{proof}

\chapter{Bipartite edge colourings and the last partition step}
\begin{theorem}[Perfect matchings in a regular bipartite multigraph]
\label{bp:regular-bipartite}
\lean{Oberwolfach.regularBipartite_card_eq, Oberwolfach.regularBipartite_hall, Oberwolfach.regularBipartite_perfectMatching, Oberwolfach.regularBipartite_decomposition}
\leanok
Let $X,Y$ be finite sets and $M:X\times Y\to\mathbb N$ the edge
multiplicity function of a bipartite multigraph. If every row sum and
column sum is $d$, there are bijections $\sigma_c:X\to Y$, indexed by
$c<d$, such that
\[
 M(x,y)=\sum_{c<d}\mathbf1_{\{\sigma_c(x)=y\}}.
\]
When $d=0$, the index set is empty; no equality of the sizes of $X,Y$
is required in that case.
\end{theorem}
\begin{proof}\leanok
For $d>0$, double counting gives $d|X|=d|Y|$. For every $S\subseteq X$,
counting the edges incident with $S$ gives
$d|S|\le d|N(S)|$, so Hall's theorem supplies a perfect matching.
Subtract one from the multiplicity of each of its edges. Every row sum
and column sum decreases by one. Induction on $d$ gives the asserted
identity, including parallel copies. The Lean proof uses mathlib's
proved finite Hall theorem.
\end{proof}

\begin{theorem}[K\"onig's edge-colouring theorem]
\label{bp:konig}
\lean{Oberwolfach.bipartiteMultigraph_edgeColouring, Oberwolfach.BipartiteEdges, Oberwolfach.bipartiteMultigraph_properEdgeColouring}
\uses{bp:regular-bipartite}
\leanok
If every row and column sum of $M$ is at most $d$, its individual edges
\[
 E_M=\{(x,y,i):x\in X,\ y\in Y,\ 0\le i<M(x,y)\}
\]
admit a colouring with $d$ colours such that two edges of the same
colour sharing an endpoint are equal.
\source{thm:prelim:konig}
\end{theorem}
\begin{proof}\leanok
Write $r_x=\sum_yM(x,y)$ and $s_y=\sum_xM(x,y)$. Complete $M$ to a
$d$-regular bipartite multigraph on $(X\sqcup Y,Y\sqcup X)$, with matrix
\[
 \begin{pmatrix}
 M&\operatorname{diag}(d-r_x)\\
 \operatorname{diag}(d-s_y)&M^{\mathsf T}
 \end{pmatrix}.
\]
Apply the regular decomposition and restrict each matching to the
original $X\times Y$ block. For each pair $(x,y)$, the number of retained
colours is exactly $M(x,y)$. A bijection between these colours and its
parallel-copy indices defines the edge colouring. The Lean proof
establishes injectivity on each set of incident edges.
\end{proof}

\begin{lemma}[Grouping matchings and recovering physical edges]
\label{bp:partition-grouping}
\lean{Oberwolfach.regularBipartite_groupedDecomposition, Oberwolfach.regularPartition_of_balancedOrientation}
\uses{bp:regular-bipartite,bp:graphs}
\leanok
A $kq$-regular bipartite multigraph splits into $k$ multigraphs with
degree $q$ on both sides.
Suppose moreover that its two parts are copies of the vertex set of
a simple graph $G$, and
\[
 M(x,y)+M(y,x)=\mathbf1_{\{xy\in E(G)\}}\quad\text{for all }x,y.
\]
Then grouping the perfect matchings into $k$ groups of size $q$ and
forgetting directions gives an exact edge partition of $G$ into
$2q$-regular spanning graphs.
\source{lem:family:partition}
\end{lemma}
\begin{proof}\leanok
Reindex the $kq$ perfect matchings by $[k]\times[q]$ and sum within each
group. If its resulting multiplicities are $Q_a(x,y)$, put
$W_a(x,y)=Q_a(x,y)+Q_a(y,x)$. Then
\[
 \sum_aW_a(x,y)=\mathbf1_{\{xy\in E(G)\}},\qquad
 \sum_yW_a(x,y)=2q.
\]
The first identity makes every $W_a(x,y)$ either zero or one, excludes
loops, and proves that each physical edge belongs to exactly one group.
The second gives the degree. The Lean conclusion uses the same
\texttt{EdgePartition} definition as the main target.
\end{proof}

This verifies grouping and projection. The next chapter supplies the
prescribed-degree construction and protected cores. The directed-path degree correction is checked below. The full
random balanced-orientation event and complete protected-partition assembly
are checked below.

\chapter{Prescribed degrees and the deterministic protected partition}
\begin{lemma}[Integral prescribed degrees]
\label{bp:prescribed-degrees}
\lean{Oberwolfach.quotaAssignment, Oberwolfach.bipartite_prescribedDegrees, Oberwolfach.bipartiteEdges_filter_card, Oberwolfach.bipartite_prescribedMultiplicity}
\leanok
Let $B$ be a finite bipartite multigraph with parts $X,Y$, whose parallel
edges are distinct elements of a finite set. Let $f:X\to\mathbb N$ and
$g:Y\to\mathbb N$ have equal totals. If
\[
 f(S)\le g(T)+e_B(S,Y\setminus T)\qquad(S\subseteq X,\ T\subseteq Y),
\]
then a subset of the original edges has degrees $f$ on $X$ and $g$ on $Y$.
The equivalent multiplicity-matrix conclusion also retains every capacity.
\source{lem:prelim:prescribed-degrees}
\end{lemma}
\begin{proof}\leanok
For each $x$, create $f(x)$ slots adjacent to all edges incident with $x$.
For each $y$, create $d_B(y)-g(y)$ additional slots adjacent to all edges
incident with $y$. The cut inequalities imply $g(y)\le d_B(y)$ and Hall's
condition for these slots. Their total number is $|E(B)|$, so Hall supplies
a bijection to the original edges. Select the edges assigned to $X$-slots.
The remaining edges at $y$ are precisely its $d_B(y)-g(y)$ assigned edges.
This proves the integral assertion. The real-valued assertion is proved
separately below.
\end{proof}

\begin{theorem}[Real max-flow/min-cut and fractional prescribed degrees]
\label{bp:real-flow}
\lean{Oberwolfach.real_maxFlow_minCut, Oberwolfach.networkFlow_cut_bound, Oberwolfach.bipartite_fractionalDegrees_iff, Oberwolfach.bipartite_nearUniformFractionalDegrees}
\leanok
Every finite directed network with nonnegative real capacities and distinct
source and sink has a flow whose value equals the capacity of a source--sink
cut. Consequently the bipartite cut criterion above is also sufficient and
necessary for real edge weights $0\le q_e\le c_e$. With capacities one, this
proves the fractional assertion of the prescribed-degree lemma. The real
normalized-demand application is proved with the original fixed part sizes.
\source{thm:prelim:max-flow-min-cut, lem:prelim:prescribed-degrees, lem:prelim:normalized-demands}
\end{theorem}
\begin{proof}\leanok
The feasible flows form a nonempty compact subset of the capacity box, so
net source outflow attains its maximum. For a fixed maximizing flow $f$, say
that $x$ can transfer to $y$ if some capacity-bounded $g$ changes divergence
by $\varepsilon(\mathbf1_x-\mathbf1_y)$ for $\varepsilon>0$. This relation is
reflexive and transitive: combine transfers of sizes $u,w$ by the convex
coefficients $w/(u+w),u/(u+w)$. A residual arc gives a transfer. The vertices
reachable from the source under this relation form a cut excluding the sink;
otherwise the flow value increases. Every outgoing cut arc is saturated,
and every incoming cut arc has flow zero. Summing divergence over the cut
proves equality of its capacity and the flow value. The three-layer
bipartite network then proves the prescribed-degree criterion, including
saturation of every source and sink arc. Integer integrality is supplied
by the preceding Hall argument rather than assumed for real flows.
\end{proof}

\begin{lemma}[Normalized demand cuts and integer subgraphs]
\label{bp:normalized-demands}
\lean{Oberwolfach.normalizedDemand_cut, Oberwolfach.nearUniform_cutCondition, Oberwolfach.bipartite_nearUniformDegrees, Oberwolfach.nearUniformMultiplicity_cut, Oberwolfach.intervalDemands_cut}
\uses{bp:prescribed-degrees}
\leanok
Under the normalized large-cut and demand bounds of
\texttt{lem:prelim:normalized-demands}, the required cut inequalities hold.
Consequently integer demands admit an integral subgraph. In particular,
integer demands in $[q-e,q+e]$, on two equal parts of positive order, have
relative error at most $2e/(q-e)$ provided $0\le e<q$.
Here $e$ is real and is not rounded to an integer.
\source{lem:prelim:normalized-demands}
\end{lemma}
\begin{proof}\leanok
For $s=|S|/|X|$ and $t=|T|/|Y|$, a positive demand deficit forces
$s+t>1-2\delta_0$. The deficit is at most
$(1+\delta)M\min(s,t)$, which the large-cut estimate bounds by $e_B(S,T)$.
For interval demands, their mean is at least $q-e$ and their deviation
from that mean is at most $2e$. All these inequalities, their conversion
to the integer cut condition, and the prescribed-degree application are checked.
\end{proof}

\begin{theorem}[Deterministic protected partition]
\label{bp:protected-deterministic}
\lean{Oberwolfach.regularPartition_of_orientedMultiplicities, Oberwolfach.protectedBipartiteRegularPartition, Oberwolfach.protectedRegularPartition, Oberwolfach.protectedRegularPartition_of_largeCuts}
\uses{bp:normalized-demands,bp:partition-grouping}
\leanok
Let $(M_a)$ partition the arcs of a balanced orientation of a simple graph,
with total indegree and outdegree $\sum_a m_a$. Suppose every colour has
in- and outdegrees $m_a\pm e$. Let $L_a\le M_a$ be its available arcs after
protected edges have been excluded. Suppose $L_a$ satisfies the large-cut
bound
\[
 e_{L_a}(S,T)\ge h_a n^2\min\{|S|/n,|T|/n\}
 \quad\text{if }|S|/n+|T|/n>1-2\delta_a,
\]
and
\[
 0\le e<q\le\min_a m_a,\quad
 0\le\delta_a<1/4,\quad
 \frac{2e}{q-e}\le\delta_a,\quad q+e\le h_an/2.
\]
Then the underlying graph has an exact edge partition into $2m_a$-regular
graphs. Each colour changes at most $4q+2e$ edges at any vertex. Every
edge unavailable for deletion in all colours retains its original colour.
\source{lem:family:partition}
\end{theorem}
\begin{proof}\leanok
Use deletion demands $d^+_{M_a}(v)-(m_a-q)$ and
$d^-_{M_a}(v)-(m_a-q)$ in $L_a$. Their totals agree and the normalized
lemma gives the deletion subgraphs $T_a$. The cores have in- and outdegree
$m_a-q$. The union of the $T_a$ has in- and outdegree $kq$; group its
perfect matchings into $q$ per colour and add them to the cores.
The proof checks exact multiplicities, physical edge coverage, both degree
identities, protected colours and all four addition/deletion degree bounds.
The random orientation is treated below. The complete simultaneous colour-cut
event and its assembly with all remaining parameter choices are still open.
\end{proof}

\begin{lemma}[Degree correction by directed two-edge paths]
\label{bp:orientation-routing}
\lean{Oberwolfach.degreeDiscrepancy_requests, Oberwolfach.twoPathPacking_exists, Oberwolfach.orientationRouting}
\leanok
An orientation matrix $M$ records exactly one direction of every edge of
$G$. Let prescribed integer outdegrees $t_v$ have total $e(G)$ and differ
from the current outdegrees by at most $a$. Suppose every distinct ordered
pair has at least $\beta n$ eligible directed two-edge paths, and a forbidden
graph has maximum degree at most $f$. If $L>0$ is integral and
\[
 na/L+2a+4L+2f<\beta n,
\]
then reversing eligible paths outside the forbidden graph gives precisely
the prescribed outdegrees. At most $a+2L$ physical edges change at a vertex.
\source{lem:simp:orientation-routing}
\end{lemma}
\begin{proof}\leanok
List positive and negative degree discrepancies as paired unit requests.
There are at most $na$ requests and each vertex occurs as an endpoint at
most $a$ times. Greedily choose their middle vertices, using each at most
$L$ times. At a step, saturated middle vertices exclude at most $na/L$
choices. Used outgoing and incoming arcs exclude at most $2a+2L$ choices;
the forbidden graph excludes at most $2f$. The displayed stronger margin
therefore leaves a choice. The chosen paths have disjoint directed arcs.
Since they lie in an orientation, no two use opposite directions of one
physical edge. Reversal preserves each physical-edge identity and changes
outdegree by exactly the prescribed discrepancy. Counting endpoint and
middle occurrences proves the stated physical degree bound.
\end{proof}

\chapter{Finite categorical laws and sampling implications}
\begin{lemma}[Fixed root-omitted expectations]
\label{bp:finite-expectations}
\lean{Oberwolfach.FiniteLaw, Oberwolfach.FiniteLaw.pi, Oberwolfach.finiteLaw_expectation_pattern, Oberwolfach.finiteLaw_expectation_omitted_root_sum}
\uses{bp:omitted}
\leanok
For independent categorical coordinates with probability vectors $p_e$,
a prescribed pattern on $F$ has probability $\prod_{e\in F}p_e(s_e)$.
For fixed coefficients, candidate families and omitted roots,
\[
 \mathbb E\sum_{i\in C}w_i\prod_{e\in F_i\setminus R}
 \mathbf1_{\{\sigma_e=s_{i,e}\}}
 =\sum_{i\in C}w_i\prod_{e\in F_i\setminus R}p_e(s_{i,e}).
\]
The product law is explicitly defined and proved to have total mass one.
No conditioning on root presence occurs. This does not prove the
manuscript-specific identities for physical images or their influence counts.
\end{lemma}

\begin{lemma}[Projection of one common event]
\label{bp:finite-projection}
\lean{Oberwolfach.FiniteLaw.prod, Oberwolfach.finiteLaw_common_event_projection}
\leanok
Let colours and auxiliary directions have independent finite laws.
If a joint event $E$ implies that every later admissible family admits its
required construction, then the probability of the colour event expressing
this universal conclusion is at least the probability of $E$.
\end{lemma}
\begin{proof}\leanok
Sum the product masses over the joint event. For each colour outcome its
fibre contributes at most the colour's mass; outside the universal colour
event the fibre is empty. This preserves the quantifier over all later
families on one event. It does not bound the probability of the joint event.
\end{proof}

\begin{theorem}[Concentration for the finite categorical product]
\label{bp:categorical-concentration}
\lean{Oberwolfach.boundedIndependentSum_tail, Oberwolfach.finiteLaw_pi_measure, Oberwolfach.finiteLaw_coordinate_independence, Oberwolfach.finiteLaw_bounded_coordinate_sum, Oberwolfach.categoricalDirection_cut_tail}
\uses{bp:finite-expectations}
\leanok
For independent finite coordinates $\sigma_i$, and functions
$a_i\le X_i(\sigma_i)\le b_i$,
\[
 \Pr\left(\left|\sum_i X_i-\sum_i\mathbb E X_i\right|\ge t\right)
 \le 2\exp\left(-\frac{2t^2}{\sum_i(b_i-a_i)^2}\right)
 \qquad(t\ge0).
\]
The explicit finite product is identified with the product probability
measure, proving coordinate independence. The bound follows from
mathlib's proved subgaussian Hoeffding estimate, applied to the centred
variables. In a raw directed-colour cut, one physical edge has one
categorical colour--direction state. If both directions meet the cut,
its two contributions are combined into one $[0,1]$-valued variable;
its expectation is the sum of the two corresponding $\rho_a/2$ terms.
This theorem does not yet establish the simultaneous cut event or its
asymptotic failure bound required by the protected partition.
\end{theorem}

\begin{theorem}[Bounded differences and fixed weighted patterns]
\label{bp:bounded-differences}
\lean{Oberwolfach.finiteLaw_boundedDifferences_mgf, Oberwolfach.finiteLaw_boundedDifferences, Oberwolfach.finiteLaw_zero_boundedDifferences, Oberwolfach.finiteLaw_omitted_weighted_tail}
\uses{bp:finite-expectations,bp:weighted-stability}
\leanok
Let $X$ be a function of independent finite categorical coordinates. If
changing only coordinate $e$ changes $X$ by at most $b_e\ge0$, then
\[
 \mathbb E e^{u(X-\mathbb EX)}\le
 e^{u^2\sum_e b_e^2/8},\qquad
 \Pr(|X-\mathbb EX|\ge t)\le
 2e^{-2t^2/\sum_e b_e^2}\quad(t\ge0).
\]
When all $b_e=0$, $X$ equals its expectation at every configuration.
In particular, this applies to a fixed root-omitted weighted pattern sum
with $b_e\ge\sum_{i:e\in F_i\setminus R}w_i$. Its expectation is the exact
product expression above; weights and omitted roots remain fixed.
\source{thm:prelim:bounded-differences}
\end{theorem}
\begin{proof}\leanok
Induct on the number of coordinates. Averaging over the remaining
coordinates changes the conditional mean by at most the first coordinate's
bound. The finite-variable Hoeffding exponential estimate controls this
conditional mean; the induction hypothesis controls each conditional tail.
Factoring the exponential and summing the product law gives the displayed
moment bound. Apply the proved subgaussian tail inequality. Reindexing a
finite coordinate set preserves its exact product law. The zero-influence
case follows by changing coordinates one at a time. The weighted assertion
uses the previously checked deterministic influence inequality.
\end{proof}

\begin{lemma}[One event for all fixed tests]
\label{bp:common-event-probability}
\lean{Oberwolfach.finiteLaw_union_bound, Oberwolfach.finiteLaw_common_event_lower, Oberwolfach.finiteLaw_common_event_exists, Oberwolfach.finiteLaw_simultaneous_boundedDifferences}
\uses{bp:bounded-differences}
\leanok
If fixed bad events $A_j$ have probabilities at most $p_j$, then
$\Pr(\bigcap_j A_j^c)\ge1-\sum_jp_j$. If the sum is less than one,
a positive-mass outcome satisfies every test. Applying this to bounded
differences keeps each test's own tolerance and influence bounds.
No independence between tests or uniformity over arbitrary
outcome-dependent domains is assumed.
\end{lemma}

\begin{lemma}[Fair orientations and simultaneous two-step paths]
\label{bp:random-orientation}
\lean{Oberwolfach.sampledOrientation_physical, Oberwolfach.sampledOrientation_degree_tail, Oberwolfach.sampledOrientation_twoPath_tail, Oberwolfach.sampledOrientation_common_event, Oberwolfach.sampledOrientation_balancing_event}
\uses{bp:common-event-probability,bp:orientation-routing}
\leanok
Sample one fair bit for each unordered vertex pair; bits on nonedges are
unused. This gives exactly one direction to each edge. For a fixed
eligibility graph $E\subseteq G$, the number of directed paths $x\to z\to y$
with both edges in $E$ has mean $|N_E(x)\cap N_E(y)|/4$. For $x\ne y$,
different choices of $z$ use disjoint physical coordinate pairs. The
probability of at most $|N_E(x)\cap N_E(y)|/8$ such paths is at most
$2\exp(-|N_E(x)\cap N_E(y)|/64)$.
The corresponding degree bounds and all pair bounds hold on one event
with the explicit sum of their failure bounds. On this event, the
routing conclusion holds for every later forbidden graph satisfying
its degree bound.
\source{clm:partition:balanced-orientation}
\end{lemma}

\begin{lemma}[Balanced orientation on a common event]
\label{bp:balanced-orientation}
\lean{Oberwolfach.orientationMatrix_degree, Oberwolfach.balancedOrientation_highProbability, Oberwolfach.eventually_orientation_numeric, Oberwolfach.eventually_orientation_failure_small}
\uses{bp:random-orientation,bp:orientation-routing}
\leanok
For every $\alpha>0$ there is $N(\alpha)$ such that the following holds
for every graph $G$ of order $n\ge N(\alpha)$ with even degrees and
$\delta(G)\ge(1/2+\alpha)n$. A fair independent orientation $D_0$ has,
with probability at least $1-n^{-10}$, this property: for every later
graph $F$ with $\Delta(F)\le\alpha n/100$, there is a balanced orientation
$D$ of $G$ which changes at most $4n^{4/5}$ incident edges at any vertex
and changes no edge of $F$.
The order threshold precedes the vertex set, graph and forbidden graph.
\source{clm:partition:balanced-orientation}
\end{lemma}
\begin{proof}\leanok
Minimum degree gives at least $2\alpha n$ common neighbours for every
pair. The checked orientation event has failure at most
\[
 2n\exp(-2n^{1/5})+2n^2\exp(-\alpha n/32).
\]
This is a valid weaker intermediate exponential estimate than the
manuscript's Chernoff estimate, with the same required final bound.
The checked real-power asymptotics make it at most $n^{-10}$.
Set $a=n^{3/5}$, $\beta=\alpha/4$, $f=\alpha n/100$ and
$L=\lceil n^{4/5}\rceil$. The proof includes all rounding inequalities,
shows the routing margin for large $n$, and proves $a+2L\le4n^{4/5}$.
Evenness and the degree-sum identity supply the half-degree vector and
its required total. The exact physical identity gives both the final
indegree and outdegree. The forbidden graph is chosen after the event.
\end{proof}

\begin{lemma}[The large-cut consequence of minimum degree]
\label{bp:graph-cuts}
\lean{Oberwolfach.graphCut, Oberwolfach.graphCut_eq_sum, Oberwolfach.graphCut_symm, Oberwolfach.graphCut_lower_of_degrees, Oberwolfach.graphCut_large_of_minDegree}
\leanok
Define $e_G(S,T)=\sum_{x\in S}|N_G(x)\cap T|$, counting ordered pairs.
If $\delta(G)\ge(1/2+\alpha)n$ and
$|S|+|T|>(1-2\delta)n$, then
\[
 e_G(S,T)\ge(\alpha-\delta)n\min\{|S|,|T|\}.
\]
The count is symmetric in $S,T$. If $T$ is the larger side, it has
size greater than $(1/2-\delta)n$, and each vertex in $S$ has at least
$\delta(G)+|T|-n$ neighbours in $T$. Summation proves the inequality;
the other case follows by symmetry. Overlapping sides retain both
ordered occurrences of an edge when appropriate.
\source{eq:family:raw-partition-cut}
\end{lemma}

\begin{lemma}[Simultaneous categorical directed-colour cuts]
\label{bp:colour-cuts}
\lean{Oberwolfach.sampledColourOrientation_physical, Oberwolfach.sampledColourCut_tail, Oberwolfach.sampledColourCut_all_cuts, Oberwolfach.sampledColourCut_quadratic_event, Oberwolfach.eventually_allCuts_failure_small}
\uses{bp:graph-cuts,bp:bounded-differences,bp:common-event-probability}
\leanok
Give each unordered pair one categorical state $(a,b)$ of probability
$\rho_a/2$, where $\rho_a\ge0$ and $\sum_a\rho_a=1$. On graph edges,
$a$ is the colour and $b$ the direction. The colour classes partition
the physical edges exactly. For fixed $S,T$, their directed-colour count
$X_a(S,T)$ has expectation $\rho_a e_G(S,T)/2$.
For any fixed number of colours $k$ and $q>0$, there is $N(k,q)$ such
that, for every graph of order $n\ge N(k,q)$ and every such profile,
with probability at least $1-n^{-10}$,
\[
 \left|X_a(S,T)-\rho_a e_G(S,T)/2\right|<qn^2
 \quad\text{for every }a,S,T.
\]
The threshold precedes the vertex set, host and profile. This absolute
quadratic tolerance controls large cuts; it does not replace the separate
small-cut argument using vertex degrees in the partition proof.
\source{eq:family:raw-partition-cut}
\end{lemma}
\begin{proof}\leanok
Write the ordered incidences as fixed singleton-coordinate patterns.
Their exact expectation follows from the categorical law. The fibre of
an unordered pair among ordered pairs has size at most two, proving a
coordinate influence bound of two without assuming independent opposite
arcs. Bounded differences gives a failure at most
$2\exp(-t^2/(2n^2))$. Sum over the $k4^n$ colour--cut tests and put
$t=qn^2$. The proof verifies
$2k4^n\exp(-q^2n^2/2)\le n^{-10}$ for all sufficiently large $n$.
This supplies the simultaneous raw cut estimate; assembling the degree,
protected-set, balancing and cut-loss events remains a separate obligation.
\end{proof}

\begin{lemma}[Joint sampling event and exact marginal laws]
\label{bp:protected-sampling}
\lean{Oberwolfach.protectedSample_highProbability, Oberwolfach.sampledColourDegrees_common_event, Oberwolfach.categoricalDirection_colour_probability, Oberwolfach.categoricalDirection_direction_probability, Oberwolfach.finiteLaw_pi_restrict, Oberwolfach.eventually_partition_margins, Oberwolfach.matrixCut_large_of_degree_and_quadratic, Oberwolfach.protectedPartition_from_balancing}
\uses{bp:bounded-differences,bp:common-event-probability,bp:balanced-orientation,bp:colour-cuts,bp:protected-deterministic}
\leanok
For each fixed inverse power, the raw incoming, outgoing and fixed-protected-graph
colour degrees, all directed-colour cuts, and the balancing property hold
on a single event with the asserted failure bound. The fixed protected graph
precedes the sample. The balancing property quantifies over every later
admissible forbidden graph on this same event.
The colour and direction marginals equal the original categorical and fair-bit
laws exactly. Restricting the full product to any fixed set of coordinates
also gives exactly the product law on those coordinates; unused nonedge
coordinates do not alter the physical-edge sampling law.
\end{lemma}

\chapter{Uniform fixed-size subsets and partitions}
\begin{definition}[The finite sample spaces]
\label{bp:fixed-size-laws}
\lean{Oberwolfach.FixedSubset, Oberwolfach.FixedPartition, Oberwolfach.UniformSampleSpec, Oberwolfach.UniformSampleSpec.State, Oberwolfach.UniformSampleSpec.law, Oberwolfach.UniformSampleSpec.exchange, Oberwolfach.UniformSampleSpec.exposures}
\leanok
For a finite set $V$ of size $D$, the subset space consists of all
$r$-subsets, where $r\le D$. An ordered partition with prescribed sizes
$d_1,\ldots,d_k$ is represented by a function $c:V\to[k]$ with
$|c^{-1}(a)|=d_a$. The sum condition $\sum_a d_a=D$ supplies a
partition, including when some cells are empty. Each space has its
uniform probability law. A transposition of two vertices either leaves
the object unchanged or performs one exchange in the manuscript's sense.
\end{definition}
\begin{lemma}[Exact laws under ordering and forgetting order]
\label{bp:fixed-size-pushforward}
\lean{Oberwolfach.orderedSample_cons_law, Oberwolfach.orderedSampleRange_uniform, Oberwolfach.fixedPartition_uniform, Oberwolfach.fixedSubset_nonempty, Oberwolfach.fixedPartition_nonempty}
\uses{bp:fixed-size-laws}
\leanok
A uniform injection $[r]\to V$ has a uniform first label and, conditional
on that label, a uniform continuation without replacement. Forgetting
the order gives a uniform $r$-subset. Applying a uniform permutation to
any one partition gives a uniform partition with the same cell sizes.
\end{lemma}
\begin{proof}\leanok
The first-label decomposition is an explicit equivalence, and its fibres
have equal cardinality. The fibre of a subset is equivalent to the
injections $[r]\to S$, so all subset fibres have equal cardinality.
For partitions, bijections between corresponding cells give a permutation
carrying any partition to any other. Composition by that permutation
is a bijection between the corresponding fibres of the sampling map.
All statements identify point masses of the actual finite laws.
\end{proof}
\begin{theorem}[Concentration for independent samples and partitions]
\label{bp:fixed-size-concentration}
\lean{Oberwolfach.uniformSamples_mgf, Oberwolfach.uniformSamples_absolute_tail, Oberwolfach.uniformSamples_one_sided_tails, Oberwolfach.uniformSamples_zero_variance}
\uses{bp:fixed-size-laws,bp:fixed-size-pushforward,bp:bounded-differences}
\leanok
Let independent coordinates be uniform subsets or partitions as above.
Suppose one exchange in coordinate $i$ changes $f$ by at most $b_i\ge0$.
Put $s_i=\min(r_i,D_i-r_i)$ for a subset and $s_i=D_i$ for a partition,
and put $v=\sum_i s_i b_i^2$. Then
\[
 \mathbb E\exp\bigl(t(f-\mathbb Ef)\bigr)\le\exp(t^2v/8)
 \quad(t\in\mathbb R).
\]
For $v>0$ and $a\ge0$, the two-sided tail is at most
$2\exp(-2a^2/v)$, and each one-sided tail is at most $\exp(-2a^2/v)$.
If $v=0$, $f$ equals its expectation at every outcome.
\source{lem:simp:product-sampling}
\end{theorem}
\begin{proof}\leanok
Expose an ordered sample one label at a time. A transposition couples
the continuations of any two next labels, and changes the completed
object by at most one exchange. Hence the two conditional means differ
by at most $b_i$. Conditional Hoeffding bounds and induction give the
exponential-moment estimate. Complementation of subsets gives the smaller
exposure count. Uniform permutations give the partition estimate.
A separate product induction adds the contributions of independent
coordinates. Exponential Markov bounds give the tails. In the zero case,
strictness of $1+x\le e^x$ gives equality at every outcome of positive
mass; all these uniform product laws have positive mass everywhere.
No independent-inclusion law is substituted for sampling without replacement.
\end{proof}

\begin{lemma}[Exact subset inclusion and intersection means]
\label{bp:subset-inclusion}
\lean{Oberwolfach.fixedSubset_joint_inclusion, Oberwolfach.fixedSubset_joint_inclusion_le, Oberwolfach.fixedSubset_intersection_mean}
\uses{bp:fixed-size-laws}
\leanok
For a uniform $r$-subset $S$ of an $n$-element set and a fixed set $T$,
\[
 \Pr(T\subseteq S)=\frac{\binom r{|T|}}{\binom n{|T|}}
 \le (r/n)^{|T|}.
\]
For every fixed $A$, $\mathbb E|A\cap S|=|A|r/n$.
The formulas include $r=0$ and the empty ground set, with Lean's
zero-division convention in the latter case; then $A=T=\varnothing$.
\end{lemma}
\begin{proof}\leanok
The fibre is counted by choosing the remaining elements outside $T$.
The factorial identity for binomial coefficients gives the displayed
ratio. Expressing the ratio through descending factorials bounds each
factor by $r/n$. Summing the singleton probabilities gives the mean.
\end{proof}

\begin{theorem}[Finite independent and hypergeometric Chernoff bounds]
\label{bp:chernoff}
\lean{Oberwolfach.finiteLaw_independent_chernoff, Oberwolfach.finiteLaw_independent_chernoff_large, Oberwolfach.fixedSubset_chernoff, Oberwolfach.fixedSubset_chernoff_large, Oberwolfach.finiteLaw_nonnegative_zero_mean, Oberwolfach.fixedSubset_intersection_zero_mean}
\uses{bp:subset-inclusion}
\leanok
Let $X$ be either a sum of independent finite random variables in $[0,1]$
or $|S\cap A|$ for a uniform fixed-size subset. Put $a=\mathbb EX$.
For $0\le\varepsilon\le1$,
\[
 \Pr(|X-a|\ge\varepsilon a)\le2e^{-\varepsilon^2a/3}.
\]
If $a>0$ and $t\ge a$, then $\Pr(X\ge t)\le(ea/t)^t$.
When $a=0$, $X=0$ at every outcome of positive mass.
This covers finite random variables; it does not assert a new proof
for arbitrary infinite probability spaces.
\source{lem:prelim:chernoff}
\end{theorem}
\begin{proof}\leanok
Convexity gives $\mathbb E e^{sX}\le\exp(a(e^s-1))$ for the
independent sum. For a uniform subset, expand the generating function
as a sum of joint-inclusion probabilities. The preceding lemma gives
$\mathbb E q^{|S\cap A|}\le(1+(r/n)(q-1))^{|A|}$ for $q\ge1$.
For $0<q<1$, apply this bound to the complement subset with parameter
$q^{-1}$, retaining the binomial expression until the two formulas
are identified. Setting $q=e^s$ proves the same exponential-moment
bound for every real $s$. Exponential Markov bounds and proved scalar
entropy inequalities give both tails and the large-deviation bound.
\end{proof}

\begin{lemma}[Simultaneous normalized intersection counts]
\label{bp:normalized-subset-counts}
\lean{Oberwolfach.swap_subset_intersection_change, Oberwolfach.fixedSubset_normalized_count_tail, Oberwolfach.fixedSubset_simultaneous_normalized_counts, Oberwolfach.fixedSubset_normalized_counts_exists}
\uses{bp:subset-inclusion,bp:fixed-size-concentration,bp:common-event-probability}
\leanok
For a uniform $r$-subset $S$ with $r>0$ and fixed sets $A_j\subseteq V$,
\[
 \Pr\left(\forall j\in J:\left|\frac{|A_j\cap S|}{r}
             -\frac{|A_j|}{|V|}\right|<t\right)
 \ge1-2|J|e^{-2rt^2}\qquad(t\ge0).
\]
The index type for $j$ is arbitrary; it need not be a subset of $V$.
Thus these fixed tests may describe neighbourhoods of vertices outside
the sampled domain. The graph-specific identification of the sets is
still a caller obligation.
\end{lemma}
\begin{proof}\leanok
A transposition adds at most one new element to a subset and removes
at most one old element. Hence a fixed intersection count changes by
at most one, and the normalized count changes by at most $1/r$.
Its mean is $|A_j|/|V|$. The ordered-sample exposure bound with $r$
steps gives the stated tail; a union bound gives the common event.
If its displayed lower bound is positive, a successful subset exists.
\end{proof}

\chapter{The vortex construction}
\begin{lemma}[Successive subsets and external roots]
\label{bp:vortex-step}
\lean{Oberwolfach.vortex_graph_step, Oberwolfach.vortex_set_sequence, Oberwolfach.nested_normalized_count_estimate}
\uses{bp:normalized-subset-counts}
\leanok
At each step, a set of size $\lfloor\epsilon|B|\rfloor$ can be chosen
inside $B$ so that every prescribed vertex in a set of at most
$2|B|/\epsilon$ vertices has normalized neighbourhood error less than
$|B|^{-1/3}$. The tested vertices need not belong to $B$.
Successive choices give the two-level root domains of the manuscript.
\end{lemma}
\begin{proof}\leanok
Restrict each ambient neighbourhood to the subtype corresponding to $B$.
The fixed-size concentration bound and the union bound give a common
successful outcome once $|B|$ exceeds a threshold depending only on
$\epsilon$. Construct the sequence by induction on its length.
Nesting implies that each vertex tested at a later step was tested at
every earlier step needed for its telescoping estimate.
\end{proof}

\begin{lemma}[Existence of paired vortices]
\label{bp:vortex}
\lean{Oberwolfach.Vortex, Oberwolfach.vortex_exists, Oberwolfach.vertexVortex_exists, Oberwolfach.floorOrbit_stopping, Oberwolfach.rounded_vortex_centre_error, Oberwolfach.vortex_accumulated_error}
\uses{bp:vortex-step}
\leanok
For $\beta,\rho_*>0$ and $0<\epsilon<1/4$, there is $m_0$ such that
every $D$-regular graph $P$ on $n\ge m_*\ge m_0$ vertices with $D\ge n/2$,
and every index set with $\rho=2|U|/D\in[\rho_*,1]$, admit nested sets
$(V_i,U_i)$ with independent recursions
$n_{i+1}=\lfloor\epsilon n_i\rfloor$ and
$u_{i+1}=\lfloor\epsilon u_i\rfloor$.
Their terminal vertex size lies in $[\epsilon m_*/2,m_*)$.
For $1\le i\le\ell$ and $v\in V_{\max(0,i-2)}$,
\[
 \left|2u_i/\rho-(D/n)n_i\right|\le\beta n_i,\qquad
 \left|d_P(v;V_i)-2u_i/\rho\right|\le\beta n_i.
\]
The vertex and index types are separate, and no regularity is asserted
for the induced graphs.
\source{def:completion:vortex}
\source{lem:writing:nested-pairs}
\end{lemma}
\begin{proof}\leanok
The accumulated normalized error is at most
$m_*^{-1/3}/(1-\epsilon^{1/3})$.
For $a=|U|/n$, the two floor recursions give
$|u_i-a n_i|<1/(1-\epsilon)$. Thus the error in the real degree centre
is less than $3/\rho_*$. The order threshold makes each of these errors
at most $\beta n_i/2$. Nested index sets of the required sizes are
constructed explicitly from initial segments of a finite enumeration.
\end{proof}

\begin{lemma}[A permutation preserving all vortex levels]
\label{bp:shell-permutation}
\lean{Oberwolfach.fibre_cycles, Oberwolfach.nested_shell_permutation, Oberwolfach.permutation_short_interval_injective}
\uses{bp:vortex}
\leanok
For nested finite sets starting with the whole vertex set, suppose each
nonempty difference $V_i\setminus V_{i+1}$ and the terminal set have at
least $h$ vertices. There is a permutation preserving membership in every
$V_i$ such that $\pi^k(x)\ne x$ whenever $0<k<h$.
Every orbit stays in one of these differences or the terminal set, and
$x,\pi(x),\ldots,\pi^{h-1}(x)$ are pairwise distinct.
\source{lem:writing:shell-permutation}
\end{lemma}
\begin{proof}\leanok
Assign each vertex to the largest level containing it, and choose a
cyclic permutation on each fibre. The combined permutation is explicitly
constructed with its inverse. Its return times at a vertex are exactly
the multiples of the fibre's cardinality.
\end{proof}

\begin{lemma}[Vortex estimates on difference domains]
\label{bp:vortex-differences}
\lean{Oberwolfach.Vortex.estimates_including_zero, Oberwolfach.Vortex.later_size_le_quarter, Oberwolfach.Vortex.difference_degree_estimates, Oberwolfach.nested_domains_degree_subtraction, Oberwolfach.Vortex.extended_difference_degree_estimates, Oberwolfach.Vortex.completion_difference_estimates, Oberwolfach.Vortex.local_degree_margin}
\uses{bp:vortex}
\leanok
Append $V_{\ell+1}=U_{\ell+1}=\varnothing$. For $j<k\le\ell+1$, put $A=V_j\setminus V_k$ and
$B=U_j\setminus U_k$. For $v\in V_{\max(0,k-2)}$,
\[
 \left|2|B|/\rho-(D/n)|A|\right|\le\tfrac53\beta|A|,
 \qquad
 \left|d_P(v;A)-2|B|/\rho\right|\le\tfrac53\beta|A|.
\]
The same formulas include $j=0$, using regularity of the initial graph.
The added empty level has zero degree errors; no size recursion is
asserted at this level. The caller substitutions with
$\beta\le\min\{\alpha_0/40,\delta_P/16\}$ give the stated
comparability, centre bounds and strict degree margin on each difference domain.
\source{eq:completion:vortex-differences:ordinary}
\source{eq:completion:vortex-differences:marked}
\end{lemma}
\begin{proof}\leanok
The finite-set cardinality and neighbourhood identities give exact
subtraction. The error is at most $\beta(n_j+n_k)$.
Since $n_k\le n_j/4$, this is at most
$\tfrac53\beta(n_j-n_k)$. Nesting verifies the root domains for both
estimates being subtracted.
\end{proof}

\chapter{Cyclic-partite fractional decompositions}

\begin{theorem}[Brouwer fixed point]
\label{bp:brouwer}
\lean{Oberwolfach.compact_convex_fixed_point}
\leanok
A continuous self-map of a nonempty compact convex subset of a
finite-dimensional real normed space has a fixed point.
\source{thm:prelim:brouwer}
\end{theorem}
\begin{proof}\leanok
The vendored proof from \texttt{harfe/fixed-point-theorems-lean4}
is ported to the pinned library. Its cubical Sperner argument and
finite-dimensional homeomorphism reduction are included in the kernel
check. No fixed-point axiom is introduced. The MIT license and source
retrieval and port records accompany the source.
\end{proof}

\begin{lemma}[Transportation domains and residual endpoint maps]
\label{bp:transport-residual}
\lean{Oberwolfach.transportMatrices_isCompact,Oberwolfach.intervalResidualMap_mem,Oberwolfach.intervalResidualMap_continuous}
\leanok
The fixed-marginal transportation domains are nonempty compact convex
sets. For actual consecutive interfaces, the reserve respects every
edge capacity and the scaled residual endpoint matrix belongs to the
same domain. This map is continuous, including at zero residual degree.
\source{doLpart:eq:reserve-weight}
\source{doLpart:eq:H-map}
\end{lemma}
\begin{proof}\leanok
Prove the reserve-load bound using the common-neighbourhood lower
bound. Exact residual row and column sums give the prescribed marginals.
At zero residual degree, bound the path weight between zero and that
degree. No division by a possibly zero number is cancelled.
\end{proof}

\begin{lemma}[Complementary chains and conditional path realization]
\label{bp:cyclic-path-assembly}
\uses{bp:brouwer,bp:transport-residual}
\lean{Oberwolfach.complementary_transport_fixed_point,Oberwolfach.weighted_conditional_realization,Oberwolfach.fractional_cycle_path_assembly}
\leanok
There is one fixed point of the complementary-chain map. The explicit
chain weighting has the required endpoint and consecutive-pair
marginals. Conditional realization preserves each interval's path
weighting and forces its endpoints to agree with the boundary variables.
Summing over the distinguished intervals gives marginal
$\rho_i+(k-1)\widehat\Lambda_i$ on each interval.
\source{doLpart:eq:fixed-point}
\source{doLpart:lem:conditional-realization}
\end{lemma}
\begin{proof}\leanok
Use Brouwer on the product transportation domain. Prove all chain
marginals by finite summation, including the cyclic reindexing. Construct
the conditional product law, then sum out its boundary coordinates.
Zero-mass fibres contribute zero; they do not require a positive
conditional denominator.
\end{proof}

\begin{theorem}[Fractional cycles in the interface representation]
\label{bp:cyclic-fractional-even}
\uses{bp:cyclic-path-assembly}
\lean{Oberwolfach.cyclicInterfaceGraph_has_fractional_cycles}
\leanok
For $k\ge2$, specify the interfaces $X_iY_i$ and $Y_iX_{i+1}$ by
adjacency relations. Assume local balance at both types of vertex,
nonempty finite classes, and endpoint degrees at least $\delta$, where
$k|Y_i|\le(2k-1)\delta$ and all middle classes have the same order.
Then the resulting graph has a fractional decomposition into actual
part-respecting copies of $C_{2k}$.
The formal representation uses copies of fixed finite types $X,Y$ for
the classes. Transport from an arbitrary named partition of a host
vertex set is a separate semantic bridge.
\source{doLpart:thm:main}
\end{theorem}
\begin{proof}\leanok
Local balance proves the common interface mass. Apply the fixed-point
path construction with marginals calculated from the given graph.
Reserve plus residual weights recover both interface capacities exactly.
Tag the classes and embed the cycle injectively, checking every
consecutive edge. Positive-weight tuples lie in the host graph.
Aggregate weights over all tuples representing each actual subgraph;
the edge equations are preserved, including zero loads on nonedges.
\end{proof}

\begin{corollary}[Positive weights and pair loads in the cyclic graph]
\label{bp:cyclic-spread}
\uses{bp:cyclic-fractional-even}
\lean{Oberwolfach.cyclicInterfaceGraph_fractional_cycles_with_pair_bounds}
\leanok
With equal class sizes $n$ and minimum cyclic degree at least
$(k/(2k-1)+\alpha)n$, the same construction gives weights on actual
part-respecting cycles between $a_-n^{2-2k}$ and $a_+n^{2-2k}$,
where $a_-,a_+>0$ depend only on $k,\alpha$.
The total weight through two distinct adjacent edges is at most
$a_+/n$, and through two vertex-disjoint edges at most $a_+/n^2$.
The formal statement uses tagged copies of a common finite type.
\source{doLpart:cor:spread}
\end{corollary}
\begin{proof}\leanok
Preserve the explicit conditional-product formula through the existence
construction. Positive slack bounds all residual path weights on both
sides and makes every residual endpoint entry positive. Cancel the
endpoint weights in each sector, including the empty internal product
for four-cycles. Compatible tuples have unique cycle-graph descriptions,
so aggregation introduces no multiplicity factor. Count selections with
prescribed vertices by their fixed coordinates; conflicting prescriptions
in one class give zero selections. Apply this count to the three or four
endpoints of the specified edges.
\end{proof}

\begin{lemma}[Regularization with controlled edge-deletion probabilities]
\label{bp:random-regularization}
\lean{Oberwolfach.bipartite_matrix_random_regularization}
\lean{Oberwolfach.retainedBipartiteGraph_regular}
\lean{Oberwolfach.retainedBipartiteGraph_le}
\lean{Oberwolfach.retainedBipartiteGraph_deleted_edge}
\leanok
For fixed $\beta>0$, choose $C,\eta_0>0$, then for each
$0<\eta<\eta_0$ choose $m_0$. A bipartite graph with equal parts
of size $m\ge m_0$, minimum degree $(1/2+\beta)m$ and degrees
in $[d-\eta m,d+\eta m]$ has a random spanning
$\lfloor d-\sqrt\eta m\rfloor$-regular subgraph in which each
original edge is deleted with probability at most $C\sqrt\eta$.
The formal graph uses two tagged copies of the label type.
\source{lem:four:random-regularization}
\end{lemma}
\begin{proof}\leanok
Prove the required large-cut estimate from the minimum degree.
Explicit floor and scale inequalities give the scaled demands used by
the real-flow theorem. For rounding, construct a doubly stochastic
auxiliary matching matrix with degree slots and a private complement
for each original edge. Birkhoff's theorem gives a finite law with the
prescribed edge marginals. Supported matchings satisfy the exact
integer demands. Their complementary original edge sets give the
regular graph. No independence among deleted edges is asserted.
\end{proof}

\begin{lemma}[Conditional partition degree estimates]
\label{bp:conditional-partition-degrees}
\lean{Oberwolfach.rooted_graph_equal_cells_good}
\lean{Oberwolfach.fixedPartition_joint_placement}
\lean{Oberwolfach.fixedPartition_two_placement}
\leanok
Condition on prescribed classes for a fixed vertex set $X$. For selected
classes of size $m$, a fixed family $J$ of vertex permutations, and every
vertex $v$, test the numbers of labels $z$ in a class for which
$v\sigma(z)$ is an edge. If $\varepsilon\le\eta/4$, $\eta\le1$ and
$2|X|\le\eta m/4$, all these counts differ from $Dm/n$ by less than
$\eta m$ on one event of probability at least
\[
1-2|\mathcal C||J|n\exp(-\eta^2m/2).
\]
The two-label placement probability in two distinct classes is exactly
$m^2/(n(n-1))$.
\source{thm:four:successor-kernel}
\end{lemma}
\begin{proof}\leanok
The remaining labels have the exact uniform fixed-size partition law.
A cell therefore has the uniform fixed-size subset law. Bound the
conditional mean shift by $2|X|$ and apply the proved exchange
concentration inequality. Union-bound the fixed graph-neighbourhood
and inverse-image tests. The vertices in these tests range over the
entire original domain. Count labelled injections to obtain the exact
placement probabilities.
\end{proof}

\begin{lemma}[Positive weights after cyclic interface regularization]
\label{bp:near-regular-cyclic-weights}
\lean{Oberwolfach.near_regular_cyclic_matrix_weights}
\lean{Oberwolfach.stochastic_cyclic_weights}
\leanok
Near-regular cyclic zero-one interface matrices admit nonnegative
weights on their transversal tuples. If their common regularized degree
$q=\lfloor d-\sqrt\eta m\rfloor$ exceeds the positive-margin cyclic
threshold, every positional load is exactly $q$. Every original
interface edge has load between $1-C\sqrt\eta$ and $1$; nonedges have
zero load. All tuple weights have a common upper bound
$B m^{2-c}$, and every valid tuple has weight at least
$(1-cC\sqrt\eta)A m^{2-c}$, with the same positive cyclic constants.
\source{thm:four:successor-kernel}
\end{lemma}
\begin{proof}\leanok
Sample the independently regularized interfaces and apply the proved
positive cyclic weighting to each positive-mass outcome. Removing
zero-mass outcomes changes no probabilities or expectations. Average
these weights under the same finite law. Exact incidence equations give
the positional and edge loads. A union bound over the specified tuple's
$c$ interface edges gives its survival probability; no independence
among edges of an interface is used.
\end{proof}

\begin{lemma}[Weight lost by excluding close pairs]
\label{bp:close-tuple-loss}
\lean{Oberwolfach.permutation_close_weight_loss}
\lean{Oberwolfach.separated_successor_edges_injective}
\leanok
For tuples with $c$ coordinates, fix $a$ pairwise separated labels,
where $1\le a<c$. If each tuple has weight at most $C n^{2-c}$,
the total weight of tuples satisfying these prescriptions and having
a pair at permutation distance at most $b$ is at most
\[
c^2(2b+1)C/n^{a-1}.
\]
A separated successor tuple has distinct physical edges, including
its permuted closing edge.
\source{thm:four:successor-kernel}
\end{lemma}
\begin{proof}\leanok
A violating pair must include a coordinate outside the prescribed set.
Expose all other coordinates; the remaining label has at most $2b+1$
choices. Sum over ordered pairs and apply the pointwise weight bound.
For edge distinctness, use injectivity of the tuple labels and the
separation of the successor of the first label from all other labels.
\end{proof}

\chapter{The remaining main-proof dependencies}
The next nodes record principal proof obligations, not completed Lean
interfaces. Their original statements and source locations are preserved
in \texttt{manuscript\_inventory.json}. The extracted references supplement
this map but are not a verified exhaustive dependency graph.

\begin{proposition}[Fixed-set decomposition: open]
\label{bp:fixed}
\lean{Oberwolfach.FixedSetTheorem}
\uses{bp:ordinary-preparation,bp:marked-preparation,bp:partition,bp:single}
\notready
Formalise the classification of factors by abundant bounded cycle lengths
or a positive proportion of long cycles, the sequential removal of small
classes, and the common-event application of preparation and partition.
The order threshold may depend on $S$ at this stage.
\source{thm:family:fixed-set}
\end{proposition}

\begin{proposition}[Ordinary preparation: open]
\label{bp:ordinary-preparation}
\uses{bp:ordinary-completion,bp:packing,bp:reconstruction}
\notready
The small exceptional vertex set is chosen before the independent edge
sample. With the asserted probability, the sample supplies a bounded
protected edge set. On that same outcome, every later regular graph
satisfying the stated symmetric-difference and protection conditions has
the prescribed decomposition. Formalise this order of quantifiers,
the sample law, all coefficient functions, the matching selection and
its potential-load bounds.
\source{prop:family:ordinary-preparation}
\end{proposition}

\begin{proposition}[Long-cycle preparation: open]
\label{bp:marked-preparation}
\uses{bp:marked-completion,bp:packing,bp:reconstruction}
\notready
Formalise typed directed occurrences, successor twists, the categorical
states of each physical edge, raw interval assignments and subsequent
retention. Preserve the fixed successor coefficients, omitted roots and
full external root domains. Prove the exact projection and multiplicity
identities used to recover the prescribed components.
\source{prop:family:long-preparation}
\end{proposition}

\begin{lemma}[Protected regular edge partition]
\label{bp:partition}
\lean{Oberwolfach.protectedRegularPartition_highProbability, Oberwolfach.protectedRegularPartition_family_highProbability, Oberwolfach.regularDegree_colourDensities}
\uses{bp:protected-deterministic,bp:balanced-orientation,bp:colour-cuts,bp:protected-sampling}
\leanok
Fix $\alpha>0$, $k\ge1$, positive $r_a,\xi_a$, and an integer $B\ge0$.
There is $N$ depending only on these data such that the following holds
for every $n\ge N$. Let $G$ be an even $D$-regular graph of order $n$ with
$D\ge(1/2+\alpha)n$. Let $\sum_a m_a=D/2$ and
$\rho_a=2m_a/D\ge r_a$. Fix sets $I_b$ with
$\sum_b|I_b|\le\alpha n/100000$ before colouring the edges independently
with probabilities $(\rho_a)_a$, obtaining $S_a^0$.
With probability at least $1-n^{-10}$, for every later family of
subgraphs $A_a\subseteq S_a^0$ with $\sum_a e(A_a)\le B$, there is an
exact edge partition $(S_a)_a$ of $G$ such that
\[
 d_{S_a}(v)=2m_a,\qquad A_a\subseteq S_a,\qquad
 S_a[I_b]=S_a^0[I_b],\qquad
 \Delta(S_a\triangle S_a^0)\le\xi_a n.
\]
The more general checked form protects any fixed graph of maximum degree
at most $\alpha n/100000$, and every later graph of maximum degree at most $B$.
\source{lem:family:partition}
\end{lemma}
\begin{proof}\leanok
Use the common degree, cut and balancing event. Small cut sides are
handled by subtracting omitted labels from the row or column degree;
large sides use the quadratic cut estimate. Set
$\delta_a=\min\{\alpha/32,\rho_a/512\}$ and $h_a=\alpha\rho_a/128$.
The raw cut is at least $2h_an\min\{|S|,|T|\}$.
Excluding protected edges and reversing directions loses at most
\[
 (\rho_a\alpha n/100000+n^{3/5}+B+4n^{4/5})\min\{|S|,|T|\}.
\]
The checked numerical choices make this at most $h_an\min\{|S|,|T|\}$.
Choose $q=\lfloor q_*n\rfloor$ with the source bounds on $q_*$.
All rounding and normalized-demand margins are proved uniformly before
$n,G,(m_a),(\rho_a)$ are specified. Apply the deterministic partition
with real error $e=5n^{4/5}$ and obtain the bound $4q+2e$.
The probability estimates are first proved with failure $n^{-12}$;
their union has failure at most $n^{-10}$. Project out the auxiliary
directions only after the universal quantifier over later protected graphs.
\end{proof}


\begin{lemma}[Complete fixed vortex geometry]
\label{bp:vortex-geometry}
\lean{Oberwolfach.Vortex.index_size_upper, Oberwolfach.Vortex.index_size_lower, Oberwolfach.Vortex.consecutive_size_bounds, Oberwolfach.Vortex.two_level_size_bound, Oberwolfach.Vortex.difference_size_lower, Oberwolfach.vortex_with_permutation_exists}
\uses{bp:vortex,bp:shell-permutation,bp:vortex-differences}
\leanok
The independent vertex and index floor recursions give $2u_i\le n_i$
and $u_i\ge\rho_*n_i/8$ under the stated degree and error bounds.
For consecutive nonempty levels, $\epsilon n_i/2\le n_{i+1}\le\epsilon n_i$.
The two-level size bound includes the empty pair. Difference domains
have order at least $3\epsilon m_*/8$. Enlarging $m_*$ before choosing
the input graph ensures a preserving permutation with no cycle shorter
than the prescribed fixed length. This constructs the geometry only;
the fractional coefficients and their later sampling are separate obligations.
\source{clm:completion:nested:ordinary}
\source{clm:completion:nested:marked}
\end{lemma}

\begin{lemma}[Ordinary threshold lower bound of one half]
\label{bp:threshold-half}
\lean{Oberwolfach.cycleThreshold_half_le, Oberwolfach.allowedThreshold_half_le, Oberwolfach.cycle_divisible_dense_counterexample, Oberwolfach.HasCycleDecomposition.closed_set_edge_divisible}
\uses{bp:thresholds}
\leanok
For every $r\ge3$, $\delta_r\ge1/2$. Thus $\delta_S\ge1/2$
for every nonempty allowed set $S\subseteq\{3,4,\ldots\}$.
\source{prop:four:lower}
\end{lemma}
\begin{proof}\leanok
For $t=2rk$, use $K_{t-1}$ and $K_{t+1}$ minus a Hamilton cycle
on disjoint vertex sets. The resulting graph is $(t-2)$-regular
on $2t$ vertices and is $C_r$-divisible. Its first component has
$\binom{t-1}{2}$ edges, not divisible by $r$. A cycle decomposition
would force divisibility of the edge count in every closed vertex set.
The degree ratio tends to $1/2$. The formal construction chooses $k$
explicitly for each requested order and each density below $1/2$.
This checks only the half-density consequence of the cited proposition.
\end{proof}

\begin{lemma}[Odd-cycle threshold lower bounds]
\label{bp:threshold-odd}
\lean{Oberwolfach.cycleThreshold_odd_sharp_lower, Oberwolfach.cycleThreshold_triangle_lower, Oberwolfach.odd_threshold_counterexample, Oberwolfach.HasCycleDecomposition.edge_count_le_obstruction}
\uses{bp:thresholds,bp:graphs}
\leanok
For odd $r\ge3$, $\delta_r\ge r/(2r-2)$; in particular,
$\delta_3\ge3/4$. Take two parts of size $(r-1)s$,
where $s=r(2k+1)$, join all pairs between parts, and within each part
use $r-1$ disjoint cliques of order $s$. The degree is $rs-1$,
the graph is divisible, and each odd cycle must consume an internal edge.
The internal-edge count is too small for a decomposition.
This is a proved alternative explicit construction for the ordinary bound.
The resolvable-threshold bridge is not asserted by this node.
\source{prop:four:lower}
\end{lemma}

\begin{lemma}[Components of a finite two-factor]
\label{bp:factor-components}
\lean{Oberwolfach.AllowedCycleLengths_iff_component_cycles, Oberwolfach.regular_two_component_cycle, Oberwolfach.component_regular_two_cycle_iso, Oberwolfach.connected_regular_two_cycle_iso, Oberwolfach.two_regular_cycle_partition, Oberwolfach.AllowedCycleLengths_iff_component_orders}
\uses{bp:graphs}
\leanok
Each component of a finite 2-regular graph is isomorphic to the cycle
graph on its order. The main target's condition on component cardinalities
is equivalent to the stated condition on cycle lengths. The cycle covers
precisely the reachable vertices, and its edges exhaust that component.
\end{lemma}

\begin{lemma}[Four-cycle lower bound and local threshold comparisons]
\label{bp:threshold-four}
\lean{Oberwolfach.cycleThreshold_four_lower, Oberwolfach.four_threshold_counterexample, Oberwolfach.cycleThreshold_ge_transformerDensity, Oberwolfach.cycleThreshold_ge_cycleFactorDensity, Oberwolfach.cycleThreshold_ge_ordinaryCycleLowerDensity}
\uses{bp:threshold-half,bp:threshold-odd}
\leanok
We have $\delta_4\ge2/3$. For $t=4k+7$, take cliques $A,C$ of size $t$
and an independent set $B$ of size $t-1$, add all $AB$ and $BC$ edges,
and delete a six-cycle inside $C$. All degrees are even, the edge count
is divisible by four, and the minimum degree is at least $2t-4$.
Every four-cycle uses an even number of edges inside $A$, whereas
$\binom t2$ is odd. The degree ratio tends to $2/3$.
Together with the half and odd bounds, this proves all lower comparisons
with the cycle-factor and transformer constants. This example proves an
ordinary threshold bound and does not assert regularity.
\end{lemma}

\begin{lemma}[Remaining ordinary threshold estimates: open]
\label{bp:threshold-inputs}
\uses{bp:threshold-half}
\notready
The ordinary upper thresholds remain to be formalised. The lower bounds
and the cycle-factor/transformer comparisons are checked. The common-tail reduction requires a cofinite upper bound;
the finite-grid reduction does not. The fixed-set proof also calls
local-factor results whose thresholds need the stronger lower bounds.
\end{lemma}

\begin{lemma}[One spanning factor: open]
\label{bp:single}
\uses{bp:ordinary-preparation,bp:marked-preparation}
\notready
Formalise the density-one application of ordinary or long-cycle
preparation, prescribing the same factor at every index.
The edge sample is then the entire input graph deterministically.
A positive-probability conclusion of preparation gives a decomposition
of that graph, hence one spanning copy. The cutoff is fixed before taking
the maximum of the finitely many ordinary order thresholds.
\source{lem:family:single-copy}
\end{lemma}

\begin{proposition}[Ordinary completion: open]
\label{bp:ordinary-completion}
\uses{bp:ordinary-cover,bp:absorbers,bp:sampling,bp:scale,bp:restriction,bp:vortex}
\notready
The vortex is proved. Construct the simultaneous reference estimates and
protected terminal absorber. Prove the deterministic iterative-absorption theorem
and verify all substitutions into its local estimates. The checked scalar
and set identities above do not establish these graph hypotheses.
\source{prop:completion:ordinary-construction}
\source{thm:family:completion:ordinary}
\end{proposition}

\begin{proposition}[Marked completion: open]
\label{bp:marked-completion}
\uses{bp:marked-cover,bp:absorbers,bp:sampling,bp:scale,bp:restriction,bp:omitted,bp:vortex,bp:vortex-fixed-absorber-completion,bp:nested-vortex-active-weights,bp:actual-active-local-load,bp:actual-reference-active-load,bp:actual-reference-active-parent-load,bp:actual-vortex-balance-domains,bp:joint-actual-balance,bp:joint-actual-inner-cuts,bp:actual-inner-cut-fields,bp:joint-actual-all-levels,bp:actual-balance-record-inputs,bp:parent-load-record-inputs,bp:joint-actual-parent-configurations,bp:ordered-common-configurations,bp:retained-load-fixed-weights,bp:active-nondeeper-record-inputs,bp:joint-active-record-events,bp:joint-active-record-inputs,bp:vortex-level-geometry-inputs,bp:joint-full-parent-neighbourhood,bp:joint-parent-spoke-agreement,bp:parent-neighbourhood-record-inputs,bp:two-step-vortex-and-endpoints,bp:reference-parent-normalization,bp:level-input-bridges,bp:reference-vortex-record-inputs,bp:nested-weight-target,bp:actual-reference-cut-density,bp:joint-record-error-margins,bp:finite-reference-chain,bp:sampled-fixed-absorber-completion,bp:uniform-inner-margins,bp:uniform-global-reference,bp:uniform-parent-active-cover,bp:uniform-protected-cover,bp:uniform-reference-level-density,bp:uniform-finite-completion-density,bp:ordered-level-configuration,bp:joint-parent-record-fields,bp:joint-analytic-record-fields,bp:ordered-record-error-choices,bp:sample-reference-density-hierarchy}
\notready
Keep the four marked degree balances, physical injectivity, successor
disjointness and the fixed-domain potential loads through the vortex.
Formalise the separate typed and physical deletion estimates and the
common marked absorber construction.
\source{prop:completion:marked-construction}
\source{thm:family:completion:marked}
\end{proposition}

\begin{lemma}[Ordinary cover-down: open]
\label{bp:ordinary-cover}
\uses{bp:packing,bp:local-factors,bp:neighbourhood-stability,bp:cuts}
\notready
Formalise correction of divisibility, inner correction, boundary covering,
restoration on the protected pair and each resulting degree bound.
\source{lem:completion:ordinary-cover-down}
\end{lemma}

\begin{lemma}[Marked cover-down: open]
\label{bp:marked-cover}
\uses{bp:packing,bp:local-factors,bp:neighbourhood-stability,bp:cuts,bp:omitted}
\notready
Formalise the corresponding marked operations without merging types,
directions, occurrences or physical images. Prove the four balance
identities for each operation before iterating it.
\source{lem:completion:marked-cover-down}
\end{lemma}

\begin{theorem}[Terminal absorbers: open]
\label{bp:absorbers}
\uses{bp:local-factors}
\notready
Formalise ordinary and marked transformers and the two alternative
wheel decompositions of each resulting absorber. Prove the embedding,
avoidance, divisibility and bounded-size conclusions; cite no literature
result as a Lean axiom.
\source{thm:new:ordinary-relative-absorber}
\source{thm:new:marked-relative-absorber}
\end{theorem}

\begin{theorem}[Packing applications: partly open]
\label{bp:packing}
\uses{bp:ordinary-near-balance-packing,bp:marked-near-balance-packing,bp:fractional,bp:weighted-matching,bp:weighted-stability}
\notready
Both ordinary and marked near-balance packing theorems are proved,
including sparse balancing deletion, fixed-weight restriction, positive
correction and the matching application. Complete both models' sampled
initial weights and remaining potential-load estimates. The fixed candidate
families and coefficients must be preserved. A bound for present roots
does not replace a bound for potential roots.
\source{thm:pr:ordinary-wheel-weights}
\source{thm:weights:marked-numerical}
\source{lem:new:relative-ordinary-packing}
\source{lem:new:relative-marked-packing}
\end{theorem}

\begin{theorem}[Positive successor weights for even lengths]
\label{bp:successor-complete}
\lean{Oberwolfach.successor_weights_theorem}
\leanok
For every even length at least four and positive degree margin, positive lower
and upper weight constants are fixed first. For every positive accuracy there
is a degree tolerance, then for each separation distance there is an order
threshold. The resulting weights on ordered tuples are zero unless every
ordinary edge and the closing edge to the permuted first vertex exist and
all labels are separated. Each admissible tuple has the asserted two-sided
weight bound. Every admissible position-edge has load between one minus the
accuracy and one; every positional vertex load has the corresponding bounds
relative to the supplied degree. The original partition law, one common good
event, fixed weights, close-pair deletion loss and the tuple representation
are all included in the proof.
\source{thm:four:successor-kernel}
\end{theorem}

\begin{lemma}[Exact bounded coefficients]
\label{bp:bounded-coefficient-exactness}
\lean{Oberwolfach.subspace_bounded_coefficients_exact}
\leanok
If every target in the prescribed neighbourhood within the specified subspace
has an approximation with coefficients in the same bounded box, with error
strictly smaller than the neighbourhood radius, the centre has an exact
representation with coefficients in that box. Compactness, convexity and the
separating-functional argument are proved; no approximation outside the
subspace is used.
\source{lem:four:zonotope}
\end{lemma}

\begin{lemma}[Long-interval mixing and path reserves]
\label{bp:long-interval-reserves}
\lean{Oberwolfach.long_interval_mixing}
\lean{Oberwolfach.long_reserved_path_edge_bound}
\lean{Oberwolfach.reservedPathWeight_endpoints}
\leanok
Finite transition products satisfy the relative-error contraction and the
explicit logarithmic mixing cutoff. One exceptional interface is allowed.
A clean subinterval suffices even when arbitrary interfaces precede or follow
it. For every edge of a long interval, a clean prefix or suffix gives the
reserve bound eleven twenty-sevenths of the original capacity. The reserve
is a weighting of actual paths with exactly the prescribed endpoint matrix.
These are inputs to the fixed-point and positive-weight assembly below.
\source{thm:four:long-capacities}
\end{lemma}

\begin{lemma}[Exact assembly of three capacity intervals]
\label{bp:long-capacity-assembly}
\uses{bp:long-interval-reserves}
\lean{Oberwolfach.long_capacity_interval_assembly}
\lean{Oberwolfach.halfResidualPathWeight_properties}
\lean{Oberwolfach.halfResidualPathWeight_continuous}
\lean{Oberwolfach.threeIntervalSector_formula}
\leanok
For locally balanced capacities, normalized transitions give the required
stationary profiles. The actual path reserve leaves positive residual degrees;
the half-residual endpoint map is continuous and retains the specified
boundary marginals. The proved fixed-point construction synchronizes three
intervals, and one reserve plus two half-residual contributions recovers every
capacity exactly. The pointwise sector formula cancels the endpoint matrices.
The uniform scale estimates and the representation by directed
transversal cycles are completed in the next node.
\source{thm:four:long-capacities}
\end{lemma}

\begin{theorem}[Long cyclic capacity weights on labelled classes]
\label{bp:long-capacity-complete}
\uses{bp:long-capacity-assembly}
\lean{Oberwolfach.long_cyclic_capacity_theorem}
\lean{Oberwolfach.cyclic_capacity_three_interval_weights}
\lean{Oberwolfach.threeBlock_next_position}
\leanok
The positive lower and upper coefficients are chosen before the class size
and the common finite label set. Every directed transversal cycle supported
by positive capacities receives a positive weight with the required power
of the class size, and every directed edge has exactly its prescribed load.
The original local balance and one-exception mixing hypotheses imply all
intermediate reserve and fixed-point hypotheses. Three compatible paths and
one cyclic sequence are in bijection; every edge, including the closing edge,
is preserved. Zero-capacity sequences carry zero weight. This covers copies
of an arbitrary common finite label set, as required by the marked-capacity
caller. The marked position splitting is proved below; its subsequent exact
correction and final caller remain open.
\source{thm:four:long-capacities}
\end{theorem}

\begin{lemma}[Ordinary capacities split among marked-cycle positions]
\label{bp:marked-position-splitting}
\lean{Oberwolfach.marked_capacity_position_splitting}
\lean{Oberwolfach.capacity_position_splitting}
\lean{Oberwolfach.bounded_signed_degree_correction}
\leanok
The graph degree and large-cut hypotheses give bounded signed corrections
with both endpoint demands. The last capacity is defined by subtraction;
its accumulated error proves the same positive bounds. The accuracy constants
are chosen before the graph, degree target and weights. Typed edge occurrences
are retained throughout.
\source{clm:capacity:position-splitting}
\end{lemma}

\begin{lemma}[Cycle weights in the copied marked graph]
\label{bp:marked-copied-weights}
\uses{bp:long-capacity-complete,bp:marked-position-splitting}
\lean{Oberwolfach.long_marked_copied_weights}
\lean{Oberwolfach.copied_marked_transition_hypotheses}
\lean{Oberwolfach.ordinary_three_step_minorization}
\leanok
The constructed capacities satisfy the original long-cycle hypotheses.
Two-step common-neighbourhood bounds only apply to distinct labels; the
three-step estimate explicitly subtracts the omitted diagonal contribution.
Uniform cycle-weight coefficients are chosen before the actual density,
graph, or capacities. The weights extend by zero to the fixed sequence space.
The subsequent exclusion and exact-correction caller is not yet complete.
\source{thm:capacity:marked-numerical}
\end{lemma}

\begin{lemma}[Typed marked-cycle representation]
\label{bp:marked-cycle-representation}
\lean{Oberwolfach.markedCycleIncidence_balanced}
\lean{Oberwolfach.markedCycleIncidence_zero_one}
\lean{Oberwolfach.markedCycleTypedEdge_marked_iff}
\lean{Oberwolfach.markedCycle_physical_images_injective}
\leanok
Canonical sequences start immediately after the unique marked arc. Their
incidence vectors have entries zero or one and satisfy both required balance
identities. The typed edge map is injective; composition with the supplied
injective physical map retains all multiplicities.
\end{lemma}

\begin{lemma}[Positive exact coefficients]
\label{bp:positive-exact-coefficients}
\uses{bp:bounded-coefficient-exactness}
\lean{Oberwolfach.subspace_positive_bounded_coefficients}
\leanok
A uniform positive preliminary weighting is retained when the remaining
balanced target is made exact. The shifted neighbourhood is contained in
the original approximation neighbourhood by explicit norm inequalities.
The same coefficient box is used for every nearby balanced target.
\end{lemma}

\begin{theorem}[Fractional cycle and successor weights: open]
\label{bp:fractional}
\notready
The even cyclic-partite result in the explicit interface representation
is proved in Node~\ref{bp:cyclic-fractional-even}, and its positive bounded
weights and pair loads in Node~\ref{bp:cyclic-spread}. A host with a supplied
class-labelling equivalence is also covered by the graph representation bridge.
The regularization input is proved in Node~\ref{bp:random-regularization};
Nodes~\ref{bp:conditional-partition-degrees}--\ref{bp:close-tuple-loss}
prove its conditional sampling, averaged cyclic weights and exclusion estimates.
The complete successor-weight theorem, including its ordered-tuple representation,
is proved in Node~\ref{bp:successor-complete}. The long cyclic capacity result
on copied classes is proved in Node~\ref{bp:long-capacity-complete}. Still formalise
the marked capacity theorem and the remaining partition/representation callers. The separate odd-length extension is not
used by the successor-weight construction. Imported literature inputs need either existing matching
mathlib statements or complete formal proofs.
\source{doLpart:thm:main}
\source{thm:four:long-capacities}
\source{thm:four:successor-kernel}
\source{thm:four:dense-spread}
\end{theorem}

\begin{theorem}[Weighted hypergraph matching input: open]
\label{bp:weighted-matching}
\notready
A fixed-error weighted theorem with polynomial size and codegree hypotheses is
proved in Node~\ref{bp:uniform-weighted-matching}. The stronger power-rate error
in the cited Ehard--Glock--Joos theorem is not claimed here. The fractional-resource application is proved independently in
Node~\ref{bp:fractional-resource-matching}, using the fixed-error theorem.
A reference is not a formal proof.
\source{thm:prelim:weighted-matching}
\end{theorem}

\begin{lemma}[Local factors and extensions: open]
\label{bp:local-factors}
\notready
The marked one-hub factor, including the actual physical wheel conclusions, is proved in Node~\ref{bp:marked-hub-factor-physical}. The ordinary local factor, including its fixed coefficients and root-omitted sums, is proved in Node~\ref{bp:ordinary-fixed-weight-factor}. The endpoint-triangle factor from the supplied parent matching and root-omitted potential estimates is proved in Node~\ref{bp:endpoint-fixed-weight-factor}. The ordinary local-extension conclusion is proved in Node~\ref{bp:ordinary-local-extensions}. The marked extension bounds on the same common event are proved in Node~\ref{bp:marked-local-extensions}. The ordinary successive-hub step and its exact remainder are proved in Nodes~\ref{bp:ordinary-hub-cover} and~\ref{bp:ordinary-hub-cover-remainder}. The triangle root conversion and the capped one-row-root construction for longer cycles are proved in Nodes~\ref{bp:endpoint-physical-root} and~\ref{bp:capped-root-wheels}. The common-event and comparison-host substitutions for successive roots, the boundary stages, and the full graph-specific cover-down assembly remain open. Separation restrictions and external roots must be preserved.
\source{lem:completion:ordinary-local-factor}
\source{lem:completion:marked-hub-factor}
\source{lem:completion:local-extensions}
\source{lem:completion:marked-local-extensions}
\end{lemma}

\begin{lemma}[Sampling and simultaneous estimates: open]
\label{bp:sampling}
\uses{bp:weighted-stability,bp:omitted,bp:common-event,bp:bounded-differences,bp:common-event-probability,bp:fixed-size-concentration,bp:chernoff,bp:raw-interval-law,bp:ordinary-raw-sampling,bp:ordinary-seed-means,bp:raw-neighbourhood-counts,bp:rooted-dense-configurations,bp:rooted-physical-survival,bp:isolator-count-conversion,bp:raw-isolator-event,bp:greedy-transfer-count,bp:transfer-physical-event,bp:ordinary-fullhost-reference,bp:ordinary-reference-cuts,bp:marked-first-exposure,bp:marked-reference-cuts,bp:marked-base-cycles,bp:marked-sample-physical-map,bp:marked-configurations,bp:marked-seed-loads,bp:raw-marked-restriction,bp:marked-reference-statistics,bp:vortex-common-events,bp:marked-reference-common,bp:marked-reference-vortex,bp:interval-retention,bp:retained-interval-identities,bp:marked-local-configurations,bp:marked-local-vortex,bp:marked-restricted-balance,bp:marked-restricted-balance-vortex,bp:marked-row-patterns,bp:marked-rows-vortex,bp:marked-indices-vortex,bp:marked-local-load-vortex,bp:marked-packing-load-vortex,bp:marked-retention-law,bp:marked-spoke-change-estimates,bp:marked-retained-pattern-vortex,bp:marked-retained-balance-load-stability,bp:marked-retained-configurations-vortex,bp:marked-retained-cut-vortex,bp:marked-retained-balance-vortex,bp:marked-retained-load-realization,bp:geometric-restriction-family-size,bp:marked-retained-common-configurations,bp:thinned-common-weight-witness,bp:common-joint-witness,bp:retained-common-joint-realization}
\notready
The finite independent categorical laws and the uniform fixed-size subset
and partition concentration theorem are checked above. Formalise the
remaining fixed-root substitutions, retention and trimming estimates,
configuration counts and all ordinary/marked realisation substitutions.
The original raw interval law, its exact marginals, separated-label and
hub-column independence, and ordinary fixed-coefficient potential-load
concentration, seed normalization, conflict-exclusion errors and fixed-domain
common-neighbourhood counts are proved in the nodes below. Dense conflict-free
rooted-set lower bounds and physical-edge survival are also proved, including
the external-root case. The complete potential isolator event and the
transfer lower counts and two-exposure potential event are now proved below.
The combined ordinary full-host reference event, including its parent cycle weights and uniform order threshold, is now proved below. The ordinary and marked reference cuts, marked configuration events, and fixed successor-weight potential loads are now proved below. The single-colour uniform-retention restriction theorem is now proved below, including the original two-stage law and all fixed restrictions. The one-hub local cycle/path events, including all ambient roots and their common realization on fixed vortex domains, are also proved below. Restricted balance and every bounded row/index pattern are now proved on fixed domains and fixed vortex levels below. Domain-specific local and packing weights, their positive normalized coefficients and potential-load events on fixed vortex domains are now proved below. The actual preparation consumers retain only interval colour 1. The actual marked retained law, row/index events across the vortex, and deterministic balance/load transfer are now proved below. The remaining common-witness transfer and complete joint preparation substitutions remain open. Assemble the remaining common events
on the original fixed domains. A bound for independent categorical
coordinates does not supply the law of sampling without replacement.
\source{lem:pr:rt1-realisation}
\source{lem:pr:marked-realisation}
\source{lem:pr:wi-interactions}
\source{lem:completion:trimming-restrictions}
\end{lemma}

\begin{lemma}[Supplied-vortex common outcome: checked conditional form]
\label{bp:supplied-vortex-weighted-common-outcome}
\lean{Oberwolfach.long_assigned_adaptive_supplied_vortex_weighted_fixed_ordered}
\leanok
For a supplied compatible vortex and the displayed regular-host, density,
weight and matching hypotheses, the theorem constructs the level weights
and one retained outcome satisfying the agreement-inclusive adaptive
property together with the stated matching and ordinary partial-map
conclusions. Its spacing is chosen before the weight bounds, and the
absorber tolerance before the later vortex tolerances. The hypotheses and
quantifier order in the Lean declaration are part of this statement.
\end{lemma}

\begin{lemma}[Common regular-host vortex and weight inputs: checked]
\label{bp:weighted-vortex-host-inputs}
\lean{Oberwolfach.long_host_matching_weighted_vortex}
\leanok
At the displayed original allowed-cycle regular-host threshold, the
construction supplies a vortex and a preserving short-cycle-free
permutation, together with matching weights for that same permutation.
The vortex size is selected before the allowed set and the host. The
complete invocation combining this result with the preceding common
outcome remains open.
\end{lemma}

\begin{lemma}[Reserved paths: checked finite-sequence and walk forms]
\label{bp:reserved-path-walks}
\lean{Oberwolfach.long_reserved_prescribed_subwalks, Oberwolfach.long_reserved_prescribed_subwalks_of_manuscript_bounds}
\leanok
For specified two-factor components with supplied cyclic orderings, the
summed block-capacity hypothesis yields the required injective interior
sequences and disjoint extended sequences. Consecutive vertices are
adjacent, and the sequences are represented by simple graph walks with
the displayed lengths, and each interior walk is a subwalk of its extension.
The numerical floor condition in the manuscript implies that capacity bound.
\source{lem:reconstruction:reserved-paths}
\end{lemma}

\begin{lemma}[One marked interval: checked conditional reconstruction]
\label{bp:marked-one-interval-path}
\lean{Oberwolfach.marked_hub_factor_interval_path_with_edges}
\leanok
Given the displayed vertex-disjoint marked hub factor, disjoint spoke
types, injective successor interval, terminal avoidance and physical-host
edge hypotheses, the theorem selects rooted wheels and gives a simple
path of the prescribed length with the selected wheel vertices and
terminal vertex as its support. Its physical edge set is exactly the union
of the selected wheels' physical rim edge sets. It does not yet establish
the simultaneous edge partition or all of the reconstruction lemma.
\source{lem:reconstruction:paths}
\end{lemma}

\begin{lemma}[Marked intervals at one hub: checked factor-wide assembly]
\label{bp:marked-hub-interval-family}
\lean{Oberwolfach.marked_hub_factor_interval_family_with_vertices}
\leanok
Under the displayed hub-factor, spoke-type, interval-cover, disjointness,
terminal-avoidance and physical-edge injectivity hypotheses, all supplied
intervals at one hub have simple physical paths of the stated lengths.
Their vertex sets are pairwise disjoint and cover the hub factor vertices
and interval endpoints. Their physical edge sets are pairwise disjoint and
cover precisely the hub factor's physical rim edges. The Lean declaration
records all hypotheses and the indexed support and edge-set identities.
The assembly across distinct hubs and the full reconstruction remain open.
\source{lem:reconstruction:paths}
\end{lemma}

\begin{lemma}[Marked permutation-interval reconstruction: checked form]
\label{bp:marked-global-permutation-reconstruction}
\lean{Oberwolfach.marked_decomposition_permutation_reconstruction_directed, Oberwolfach.marked_interval_partition_count}
\leanok
For the exact physical-resource decomposition and the displayed loopless,
injective physical rim map, suppose each prescribed permutation interval
is injectively indexed, the extended intervals at a hub are disjoint, their
terminal vertices avoid the covered set, and the marked spokes are exactly
the interval vertices. The theorem constructs all paths in the physical
graph, with the prescribed endpoints and lengths, hubwise vertex
disjointness and vertex union, and global disjoint coverage of both
unordered physical edges and ordered physical arcs. It also proves the
hubwise marked-vertex count; the complementary count is a separate
checked lemma. The Lean declaration records the complete hypotheses and
all indexed conclusions. Its application to the completion and iteration
interfaces remains open.
\source{lem:reconstruction:paths}
\end{lemma}

\begin{lemma}[Reconstruction and prescribed embeddings: open]
\label{bp:reconstruction}
\uses{bp:extending-list-path-pieces,bp:pinned-cycle-correction,bp:prescribed-image-invariants,bp:source-family-path-updates,bp:five-component-steps,bp:family-component-invariant,bp:actual-forbidden-source-graph,bp:family-state-resources,bp:path-segment-schedule,bp:family-original-advances,bp:actual-family-induction,bp:ordered-prescribed-embedding,bp:dependent-prescribed-embedding}
\notready
Formalise segment selection, gluing maps, inverse images, endpoint
prescriptions, list embeddings and edge assignment. Establish injectivity
and exact multiplicity before interpreting auxiliary wheels as physical
factor edges.
\source{lem:reconstruction:reserved-paths}
\source{lem:images:residual-structure}
\source{lem:new:prescribed-embedding}
\source{lem:new:prescribed-edge-assignment}
\end{lemma}

\chapter{Release criteria and scope}
A full certificate must contain a theorem whose type is definitionally
identical to \texttt{Oberwolfach.MainTheorem}, with no additional theorem
hypotheses and no unapproved axioms. The permitted foundational axioms
are \texttt{propext}, \texttt{Classical.choice} and \texttt{Quot.sound}.
Their presence is disclosed; ``no additional axioms'' does not mean an
axiom-free foundation.

The verification procedure builds the modules, enumerates the exported
project theorems and their transitive axioms, and replays declarations with
Lean's kernel. Kernel replay is not an independent implementation of the
kernel. Faithfulness of the definitions to the manuscript remains a separate
mathematical review. In particular, no successful build or green blueprint
node certifies an unimplemented sampling, absorption or literature theorem.

The present project has no unconditional Lean proof of the main theorem.
The ready-to-share package must therefore be described as a blueprint and
partial formal verification, not as a certificate of the paper.

\appendix
\chapter{Whole-manuscript statement inventory}
This table indexes the 137 public statement and definition environments in the
snapshot. The JSON inventory additionally records 33 proof-local claims,
one background problem and all 133 proof environments. Inclusion is not a claim
that a statement has been audited or formalised. Source line numbers refer to
the archived snapshot.
\small
\begin{longtable}{p{.69\textwidth}r >{\raggedright\arraybackslash}p{.18\textwidth}}
\toprule Source label & Line & Lean status \\ \midrule\endhead
\nolinkurl{conj:gen_oberwolfach} & 228 & Open \\
\nolinkurl{thm:new:prescribed-factor} & 342 & Open \\
\nolinkurl{cor:new:factor-consequences} & 357 & Open \\
\nolinkurl{thm:main:resolution} & 384 & Open \\
\nolinkurl{doLpart:thm:main} & 425 & Open \\
\nolinkurl{conj:odd_cyclic_threshold} & 444 & Open \\
\nolinkurl{def:new:permutation-intervals} & 1934 & Open \\
\nolinkurl{def:typed-graph} & 2324 & Open \\
\nolinkurl{thm:prelim:bounded-differences} & 2612 & Finite case checked \\
\nolinkurl{lem:prelim:chernoff} & 2631 & Open \\
\nolinkurl{lem:simp:product-sampling} & 2672 & Open \\
\nolinkurl{thm:prelim:max-flow-min-cut} & 2731 & Real case checked \\
\nolinkurl{lem:prelim:prescribed-degrees} & 2746 & Bipartite forms checked \\
\nolinkurl{thm:prelim:hall} & 2771 & Open \\
\nolinkurl{thm:prelim:konig} & 2780 & Multigraph theorem checked \\
\nolinkurl{def:pr:cut} & 2790 & Open \\
\nolinkurl{lem:pr:cut-stability} & 2811 & Open \\
\nolinkurl{lem:prelim:normalized-demands} & 2828 & Bipartite forms checked \\
\nolinkurl{thm:prelim:separation} & 2863 & Open \\
\nolinkurl{thm:prelim:caratheodory} & 2877 & Open \\
\nolinkurl{thm:prelim:brouwer} & 2886 & Open \\
\nolinkurl{thm:prelim:weighted-matching} & 2901 & Open \\
\nolinkurl{lem:new:weighted-resource-matching} & 2935 & Open \\
\nolinkurl{prop:four:lower} & 3048 & Open \\
\nolinkurl{thm:common:cycle-factor} & 3127 & Open \\
\nolinkurl{lem:common:fixed-length-path} & 3154 & Open \\
\nolinkurl{prop:family:ordinary-preparation} & 3187 & Open \\
\nolinkurl{prop:family:long-preparation} & 3217 & Open \\
\nolinkurl{lem:family:single-copy} & 3247 & Open \\
\nolinkurl{lem:family:partition} & 3257 & Partition theorem checked \\
\nolinkurl{thm:family:fixed-set} & 3306 & Open \\
\nolinkurl{thm:four:two-regular-fractional-threshold} & 3568 & Open \\
\nolinkurl{thm:four:DHLP-balanced} & 3750 & Open \\
\nolinkurl{thm:four:dense-spread} & 3765 & Open \\
\nolinkurl{lem:four:wheel-rounding} & 3820 & Open \\
\nolinkurl{doLpart:lem:conditional-realization} & 3973 & Open \\
\nolinkurl{doLpart:cor:spread} & 4361 & Open \\
\nolinkurl{doLpart:thm:odd} & 4473 & Open \\
\nolinkurl{thm:four:long-capacities} & 4636 & Open \\
\nolinkurl{lem:four:zonotope} & 4863 & Open \\
\nolinkurl{lem:four:random-regularization} & 4908 & Open \\
\nolinkurl{thm:four:successor-kernel} & 4971 & Open \\
\nolinkurl{thm:capacity:marked-numerical} & 5143 & Open \\
\nolinkurl{def:pr:rt1} & 5497 & Open \\
\nolinkurl{def:pr:marked-typicality} & 5640 & Open \\
\nolinkurl{lem:pr:rt1-stability} & 5677 & Generic lemma checked \\
\nolinkurl{cor:cn:relative-typicality-large-cuts} & 5901 & Open \\
\nolinkurl{def:four:local-interval-law} & 6015 & Open \\
\nolinkurl{lem:interval:raw-law} & 6062 & Open \\
\nolinkurl{def:interval:retention} & 6090 & Open \\
\nolinkurl{lem:pr:rt1-realisation} & 6136 & Open \\
\nolinkurl{lem:pr:marked-realisation} & 6187 & Open \\
\nolinkurl{lem:pr:wcodeg-realisation} & 6296 & Open \\
\nolinkurl{lem:four:sparse-alternating} & 6345 & Open \\
\nolinkurl{def:completion:ordinary-host} & 6551 & Open \\
\nolinkurl{def:ordinary-local-reference} & 6580 & Open \\
\nolinkurl{def:completion:marked-host} & 6611 & Open \\
\nolinkurl{def:completion:marked-sample} & 6681 & Open \\
\nolinkurl{lem:completion:trimming-restrictions} & 6800 & Open \\
\nolinkurl{lem:pr:wi-stability} & 7185 & Generic lemma checked \\
\nolinkurl{lem:pr:wi-interactions} & 7288 & Open \\
\nolinkurl{lem:four:edge-isolator} & 7404 & Open \\
\nolinkurl{thm:pr:ordinary-wheel-weights} & 7497 & Open \\
\nolinkurl{def:marked:common-paths} & 7760 & Open \\
\nolinkurl{def:marked:wheel-estimates} & 7828 & Open \\
\nolinkurl{thm:weights:marked-numerical} & 7883 & Open \\
\nolinkurl{def:marked:potential-loads} & 8096 & Open \\
\nolinkurl{lem:four:typed-correction} & 8179 & Open \\
\nolinkurl{lem:new:relative-ordinary-correction} & 8335 & Open \\
\nolinkurl{lem:packing:ordinary-deletion-stability} & 8441 & Open \\
\nolinkurl{lem:new:relative-ordinary-packing} & 8546 & Open \\
\nolinkurl{lem:packing:marked-deletion-stability} & 8646 & Open \\
\nolinkurl{lem:packing:marked-potential-stability} & 8718 & Open \\
\nolinkurl{lem:new:relative-marked-packing} & 8774 & Open \\
\nolinkurl{lem:new:bounded-degree-selection} & 9020 & Open \\
\nolinkurl{lem:four:transformer} & 9129 & Open \\
\nolinkurl{thm:four:absorber} & 9344 & Open \\
\nolinkurl{lem:new:triangle-transformer-explicit} & 9500 & Open \\
\nolinkurl{lem:new:triangle-absorber} & 9767 & Open \\
\nolinkurl{def:class-profile:data} & 9928 & Open \\
\nolinkurl{def:cycle:weight-absorber-estimates} & 9999 & Open \\
\nolinkurl{lem:simp:profile-marginals} & 10031 & Finite-resource proof checked \\
\nolinkurl{lem:simp:profile-factor} & 10111 & Finite-resource proof checked \\
\nolinkurl{thm:simp:profile-finite-cleanup} & 10309 & Open \\
\nolinkurl{lem:new:bernoulli-factor} & 10490 & Open \\
\nolinkurl{def:common:ordinary-weights} & 10691 & Open \\
\nolinkurl{thm:new:ordinary-relative-absorber} & 10725 & Open \\
\nolinkurl{lem:local:marked-hub-absorbers} & 11004 & Open \\
\nolinkurl{def:endpoints:assignment-graphs} & 11063 & Open \\
\nolinkurl{lem:pairing:degree-cut-hall} & 11146 & Open \\
\nolinkurl{lem:four:typed-pairing} & 11172 & Open \\
\nolinkurl{lem:new:cuts-from-codegrees} & 11285 & Open \\
\nolinkurl{thm:new:fresh-candidacy} & 11319 & Open \\
\nolinkurl{lem:new:marked-endpoint-counts} & 11513 & Open \\
\nolinkurl{lem:pr:ordinary-endpoints} & 11711 & Open \\
\nolinkurl{lem:completion:ordinary-local-factor} & 11802 & Open \\
\nolinkurl{lem:completion:marked-hub-factor} & 11986 & Open \\
\nolinkurl{lem:completion:triangle-endpoint-factor} & 12111 & Open \\
\nolinkurl{lem:completion:local-extensions} & 12273 & Open \\
\nolinkurl{lem:completion:marked-local-extensions} & 12356 & Open \\
\nolinkurl{lem:completion:marked-estimate-changes} & 12498 & Open \\
\nolinkurl{def:common:marked-estimates} & 12694 & Open \\
\nolinkurl{thm:new:marked-relative-absorber} & 12755 & Open \\
\nolinkurl{lem:four:common-neighbourhood-balancing} & 12818 & Open \\
\nolinkurl{prop:completion:ordinary-construction} & 13695 & Open \\
\nolinkurl{prop:completion:marked-construction} & 13756 & Open \\
\nolinkurl{def:completion:vortex} & 13826 & Open \\
\nolinkurl{lem:writing:nested-pairs} & 13858 & Open \\
\nolinkurl{lem:writing:shell-permutation} & 13929 & Open \\
\nolinkurl{def:completion:ordinary-local-estimates} & 13979 & Open \\
\nolinkurl{def:completion:ordinary-estimates} & 14029 & Open \\
\nolinkurl{lem:completion:ordinary-estimate-changes} & 14143 & Open \\
\nolinkurl{lem:completion:ordinary-cover-down} & 14249 & Open \\
\nolinkurl{thm:family:completion:ordinary} & 15168 & Open \\
\nolinkurl{def:completion:marked-local-estimates} & 15379 & Open \\
\nolinkurl{def:completion:marked-estimates} & 15445 & Open \\
\nolinkurl{lem:completion:marked-cover-down} & 15615 & Open \\
\nolinkurl{thm:family:completion:marked} & 16819 & Open \\
\nolinkurl{lem:completion:ordinary-reference} & 17066 & Open \\
\nolinkurl{lem:completion:ordinary-reference-cuts} & 17421 & Open \\
\nolinkurl{lem:completion:marked-reference} & 17519 & Open \\
\nolinkurl{lem:completion:marked-reference-configurations} & 17698 & Open \\
\nolinkurl{lem:completion:marked-reference-loads} & 17856 & Open \\
\nolinkurl{lem:reconstruction:reserved-paths} & 20386 & Open \\
\nolinkurl{def:reconstruction:interval-assignment} & 20431 & Open \\
\nolinkurl{lem:reconstruction:paths} & 20456 & Open \\
\nolinkurl{lem:images:residual-structure} & 20554 & Open \\
\nolinkurl{lem:images:residual-prescriptions} & 20641 & Open \\
\nolinkurl{lem:simp:orientation-routing} & 20731 & Incidence form checked \\
\nolinkurl{lem:new:marked-indegree-correction} & 20785 & Open \\
\nolinkurl{def:lists:extending} & 20857 & Open \\
\nolinkurl{lem:new:prescribed-embedding} & 20870 & Open \\
\nolinkurl{lem:new:prescribed-positions} & 20983 & Open \\
\nolinkurl{lem:new:prescribed-edge-assignment} & 21068 & Open \\
\nolinkurl{lem:list:raw} & 21145 & Open \\
\nolinkurl{lem:new:ordinary-list-extensions} & 21219 & Open \\
\nolinkurl{doLpart:prop:sharp} & 25255 & Open \\
\bottomrule\end{longtable}



\chapter{Marked capacities and the ordinary fractional-threshold bridge}
\begin{theorem}[Marked-capacity weights on canonical cycles]
\label{bp:marked-capacity-complete}
\lean{Oberwolfach.marked_capacity_cycle_weights_uniform_length, Oberwolfach.markedCycleIncidence_injective, Oberwolfach.markedCycle_distinct_physical_edges}
\uses{bp:marked-position-splitting,bp:marked-copied-weights,bp:marked-cycle-representation}
\leanok
For fixed positive lower density and expansion constants, there is a length
threshold. For each larger length, positive error and weight constants can
be chosen before the separation parameter. The actual degree, two-path,
large-cut and separation hypotheses give positive weights on all canonical
separated marked cycles with exactly one marked arc. The weights have order
$n^{2-c}$ and exactly realise every supplied balanced vector sufficiently
close to the all-one vector. The proof establishes approximation uniformly
for every nearby balanced vector, then proves the positive exact correction.
It also checks that distinct canonical cycles have distinct incidence
vectors and that injectivity of the physical edge map gives distinct
physical edges on each cycle.
\source{thm:capacity:marked-numerical}
\end{theorem}
\begin{proof}\leanok
The position splitting and copied-cycle construction supply exact weights
before separation. Deleting non-separated cycles loses at most the explicit
fixed coefficient times $(2b+1)/n$ in each edge equation. A positive baseline
on every admissible cycle and the proved subspace coefficient theorem give
exact positive coefficients. The size threshold is chosen after $b$ to make
the displayed error smaller than one quarter of the correction radius.
\end{proof}

\begin{lemma}[Averaging and undirected path packing]
\label{bp:fractional-averaging-packing}
\lean{Oberwolfach.HasCycleDecomposition.fractional, Oberwolfach.fractionalDecomposition_of_uniform_cover, Oberwolfach.fractionalDecomposition_of_divisible_complements, Oberwolfach.undirectedTwoPathPacking_exists}
\leanok
Integral cycle decompositions give fractional decompositions. Fractional
weights on a family of subgraphs covering every edge with the same positive
multiplicity can be averaged to give weights on the original graph.
For complements of an edge partition into $R$ classes, that multiplicity
is exactly $R-1$. Prescribed endpoint pairs can also be connected by paths
of length two with pairwise disjoint undirected edges and bounded middle
usage under the explicit candidate-size inequality. The complete divisibility repair and its fractional-threshold application
are proved in Node~\ref{bp:ordinary-fractional-threshold}.
\end{lemma}


\begin{theorem}[Ordinary cycle fractional threshold]
\label{bp:ordinary-fractional-threshold}
\lean{Oberwolfach.cycle_fractional_positive_margin, Oberwolfach.dense_divisible_complement_colouring, Oberwolfach.divisible_colour_repair}
\uses{bp:fractional-averaging-packing}
\leanok
For every fixed $r\ge3$ and $\alpha>0$, every sufficiently large graph of
minimum degree at least $(\delta_{C_r}+\alpha)n$ has a fractional
$C_r$-decomposition. No divisibility assumptions are needed here.
\source{eq:threshold:cycle-fractional}
\end{theorem}
\begin{proof}\leanok
Choose an even multiple $q$ of $r$ before the order threshold. One common
colouring event controls all class degrees and all ordered colour--reserve
paths. Recolour distinct reserve edges to correct edge-count residues;
then pair the odd-degree vertices and switch along physically edge-disjoint
paths to correct all degree parities. The reserve complement is even because
$q$ is even. Quantitative degree losses are absorbed at the fixed positive
margin. Each of the $q+1$ complements is $C_r$-divisible and meets the
ordinary decomposition threshold. Averaging their decompositions covers
each original edge with its exact multiplicity $q$.
\end{proof}

\begin{lemma}[Signed isolation of one edge]
\label{bp:signed-cycle-edge-isolator}
\lean{Oberwolfach.signed_cycle_edge_isolator, Oberwolfach.signed_cycle_vertex_load, Oberwolfach.fractional_cycle_coefficient_le_one}
\uses{bp:ordinary-fractional-threshold}
\leanok
At a fixed positive margin above the ordinary cycle threshold, every edge
can be isolated by a signed weighting of actual cycle subgraphs. Each
coefficient has absolute value at most one. The vertex load is one half at
each endpoint of the chosen edge and zero elsewhere.
\source{lem:four:edge-isolator}
\end{lemma}
\begin{proof}\leanok
Subtract fractional decompositions of the original graph and the graph
with the specified edge deleted. The one-edge degree loss is proved
explicitly. Every positive fractional coefficient is at most one by its
edge-load equation. Double-counting the incident edges of each actual cycle
gives the stated vertex-load identity.
\end{proof}


\begin{theorem}[Positive ordinary cycle weights and pair loads]
\label{bp:positive-ordinary-cycle-weights}
\lean{Oberwolfach.positive_ordinary_cycle_weights, Oberwolfach.positive_ordinary_fractional_cycle_weights, Oberwolfach.ordinary_cycle_adjacent_pair_load, Oberwolfach.ordinary_cycle_disjoint_pair_load}
\uses{bp:ordinary-fractional-threshold}
\leanok
For each fixed cycle length and positive degree margin, there are positive
constants and an order threshold, chosen before the host graph. Every
supplied edge vector sufficiently close to one has exact positive cycle
coefficients bounded above and below by constants times $n^{2-r}$.
In particular, the all-one vector gives a fractional cycle decomposition.
The total weight through two incident edges is $O(1/n)$, and through two
disjoint edges it is $O(1/n^2)$, for these same coefficients.
\source{thm:four:dense-spread}
\end{theorem}
\begin{proof}\leanok
Sample a fixed number of vertices and independently retain physical edges
with probabilities determined by the supplied vector. The successful
induced graph has a fractional decomposition by the proved cycle threshold.
Conditional-root concentration bounds the failed samples containing each
specified edge. Average using the exact two-vertex inclusion probability;
the $r$-vertex inclusion probability bounds each coefficient. The sample
size is fixed before the host graph so that the approximation error is
smaller than the positive correction radius. The proved convex coefficient
theorem, with a positive baseline on every actual cycle, makes the edge
loads exact. Injective encodings of physical cycles give the two pair-load
bounds. This formal proof establishes the needed ordinary cycle result
directly; it does not assume the general balanced-decomposition theorem
quoted in the manuscript.
\end{proof}

\begin{theorem}[Positive ordinary wheel weights]
\label{bp:positive-ordinary-wheel-weights}
\lean{Oberwolfach.positive_ordinary_wheel_weights, Oberwolfach.ordinary_wheel_correction, Oberwolfach.ordinary_spoke_correction}
\uses{bp:signed-cycle-edge-isolator,bp:ordinary-transfer-identities}
\leanok
The admissible induced-subgraph and transfer-configuration counts in the
manuscript imply an exact positive ordinary wheel weighting. The tolerance
is chosen before the host graphs and initial weights. Every final wheel
weight lies between $b_-n^{1-r}/2$ and $(b_++b_-/2)n^{1-r}$.
\source{thm:pr:ordinary-wheel-weights}
\end{theorem}
\begin{proof}\leanok
Lift the proved signed edge isolator to wheels at each admissible hub and
average over the actual configurations. Count configurations containing the
fixed rim together with the isolated edge to bound every coefficient.
Sum over ordered endpoint pairs with a factor of one half, retaining the
undirected-edge multiplicity explicitly. Exact rim loads and the degree
balance imply zero spoke-residual sum at every rim vertex. The labelled
base-cycle and rooted-row construction gives an actual single-vertex
transfer. Average over its configurations and over pairs of neighbouring
hubs. Only pairs containing the fixed wheel's hub contribute, cancelling
the normalization by the spoke degree. Sum the resulting bounds and choose
the tolerance to preserve the positive baseline on every wheel.
\end{proof}

\begin{lemma}[Ordinary wheel resources and potential-load stability]
\label{bp:ordinary-wheel-potential-stability}
\lean{Oberwolfach.ordinaryWheelResources_card, Oberwolfach.ordinaryWheelResources_injective, Oberwolfach.ordinary_wheel_physical_pair_load, Oberwolfach.ordinary_active_exact_loads, Oberwolfach.ordinary_active_pair_loads, Oberwolfach.ordinaryPotentialLoad_rim_degree_stability, Oberwolfach.ordinaryPotentialLoad_spoke_degree_stability}
\uses{bp:positive-ordinary-wheel-weights,bp:weighted-stability}
\leanok
An ordinary $r$-wheel has exactly $2r$ physical edges, and different wheels
have different edge sets. Restricting the hypergraph vertices to actual host
edges gives exact fractional degree one and pair load at most $r^4B/n$
whenever the wheel coefficients are at most $Bn^{1-r}$.
On a fixed parent wheel family with these coefficients, changing at most
$d$ incident edges per vertex changes a root-omitted potential load by at
most $8r^4Bd/n$ for a rim root and $4r^4Bd/n$ for a spoke root.
The root need not be present in either host. This node does not assert the
isolator-configuration or transfer-configuration potential estimates.
\source{eq:loads:ordinary-potential-seed}
\end{lemma}
\begin{proof}\leanok
Injective encodings bound cycles through prescribed vertices; the fixed
hub and free hub cases give the actual wheel pair loads. The potential
load is a product over every physical edge except the root. Its difference
is bounded by the total weight through the root and each changed edge.
For changed edges meeting root vertices use the $n^{-1}$ pair bound; for
other changed edges use $n^{-2}$. Sum these bounds using the actual row and
column degrees of the two difference relations. Candidate families,
coefficients and root omissions remain fixed throughout.
\end{proof}

\begin{theorem}[Alternating corrections from intersection estimates]
\label{bp:sparse-alternating-correction}
\lean{Oberwolfach.sparse_alternating_correction}
\leanok
For every positive lower density bound there are positive constants
$\delta_0,C$ such that the degree, pair-intersection, triple-intersection
and exceptional weighted-codegree hypotheses of the manuscript imply:
every edge vector with zero row and column sums is a linear combination
of the alternating vectors of actual ordered four-cycles. The sum of the
absolute coefficients incident with each edge is at most $C\|z\|_\infty$.
The constants precede the graph, supplied density profile, error and
exceptional sets.
\source{lem:four:sparse-alternating}
\end{theorem}
\begin{proof}\leanok
The actual four-corner averaging operator is written as $T=I+K+L$.
The two negative contributions cancel by the zero column and row sums.
The weighted-codegree, degree and pair-intersection estimates bound $L$.
Expanding $K^2$ gives the actual triple-intersection kernel. Repeated
column indices and exceptional row pairs are bounded separately, giving
an error proportional to $\delta+D/|A|+1/|B|$. The support factor of $K$
is retained on missing edges.

The operator $E=I-T$ preserves the zero-marginal subspace and satisfies
$\|E^2z\|\le\|z\|/2$. In the formal proof, finite-dimensional linear
algebra replaces the manuscript's infinite series: this bound makes $T$
injective and therefore surjective on that subspace. If $Tw=z$, the
identity $w=z+Ez+E^2w$ bounds $\|w\|$ by $2(1+\|E\|)\|z\|$.
The four-cycle coefficient bounds for $Tw$ give the claimed incident load.
No extra intersection assumption is introduced.
\end{proof}


\begin{theorem}[Positive marked wheel weights]
\label{bp:positive-marked-wheel-weights}
\lean{Oberwolfach.positive_marked_wheel_weights_numerical, Oberwolfach.marked_wheel_correction, Oberwolfach.marked_spoke_correction, Oberwolfach.markedWheelRimLoad_physical, Oberwolfach.markedWheelSpokeLoad_physical_exact}
\uses{bp:sparse-alternating-correction}
\leanok
The marked numerical host estimates yield positive exact wheel weights.
The length threshold precedes the cycle length; all coefficient and error
constants precede the separation parameter and the host order. The spoke
profiles are the supplied constant densities. The common-path family uses
root-first coordinates and omits the four root spokes. These spokes are
supplied only by the four-cycle used in the correction.
\source{thm:weights:marked-numerical}
\end{theorem}
\begin{proof}\leanok
The proved marked capacity theorem corrects rim loads. Averaging cycles
compatible with two hubs corrects spoke column sums. The remaining spoke
residuals have zero row and column sums. The proved sparse alternating
correction expresses each residual by actual four-cycles. A single common
path tuple, used with both roots and both hubs, supplies its four-wheel
transfer. The exact signed spoke equations and zero rim equations are
proved on that same tuple family. Counting only four-cycles incident with
the fixed wheel's designated position bounds each resulting coefficient.
The tolerance is then chosen to retain half of the initial lower bound.

The root-first orders are explicitly permuted into canonical cycle
coordinates. The ordinary-root order places the marked arc from the second
to the third free label. The conversion is injective and does not require
separation of the two roots. Separation and physical-map injectivity give
loopless, distinct physical rim edges and the cross-collision condition.
The physical load identities, with disjoint spoke types, identify the typed
exact equations with the manuscript's physical edge equations. Constant
spoke intersection estimates imply the inverse-density estimates by an
explicit normalization, with the original exceptional sets preserved.
\end{proof}


\chapter{Further checked matching ingredients}
\begin{lemma}[One common colouring outcome]
\label{bp:colour-round-realization}
\lean{Oberwolfach.colour_round_realization, Oberwolfach.colour_round_common_concentration, Oberwolfach.hypergraph_remaining_colour_expectation_equal_degrees, Oberwolfach.paddedColourClique_card_original, Oberwolfach.paddedResources_codegree}
\uses{bp:tentative-colouring-counts,bp:colouring-certificate-concentration}
\leanok
Let indexed edges have at most $k$ resources and lists of size $l$, and
let the designated real edges have exactly $k$ resources. Suppose each
resource--colour degree is at most $s\le l$, with equality for every colour
in a real edge's list at each of its resources. Write
\[
 \rho=(1-1/l)^{ks},\qquad \kappa=1-s\rho/l.
\]
Under the explicit codegree and exponential-tail inequalities in the Lean
statement, one outcome of the original product experiment simultaneously
approximates every real available-list size by $l\kappa^k$ and every
root-omitted remaining colour degree by
$|U|(1-\rho)\kappa^{k-1}$. The error is the stated expectation error plus
$6(k+1)t$. The retained colours belong to their supplied lists.
Finite padding equalizes the original resource--colour degrees without
increasing codegrees between distinct resources.
\end{lemma}
\begin{proof}\leanok
The exact product law supplies all expectations and tails. Each test uses
only two successive closed neighbourhoods of a specified base set. Bounding
base-set sizes and their coordinate incidence gives one application of the
finite local lemma. The same positive-mass outcome works for all tests.
The proof retains the omitted resource in every remaining-degree test.
\end{proof}
The finite iteration and its assembly into near-optimal hypergraph colouring
are now proved below. The uniform choices for the fixed-error weighted
theorem and its fractional-resource rounding application are also proved below.

\begin{lemma}[Cycle weights for vertex-disjoint packing]
\label{bp:cycle-vertex-matching-weights}
\lean{Oberwolfach.ordinaryCycleVertices_card, Oberwolfach.ordinary_cycle_vertex_weight_identity, Oberwolfach.ordinary_cycle_vertex_matching_weights}
\uses{bp:positive-ordinary-cycle-weights}
\leanok
Suppose $a_C\ge0$ is an edge-fractional $C_r$-decomposition of $Q$ with
$a_C\le B n^{2-r}$. The hyperedge indexed by $C$ has resource set $V(C)$,
which has size $r$. If $d\ge n/2$ and
$(1-\eta)d\le d_Q(v)\le d$, then $b_C=2a_C/d$ satisfies
\[
 1-\eta\le\sum_{C\ni v}b_C\le1,\qquad
 \sum_{C\ni x,y}b_C\le\frac{4r^2B}{n}\quad(x\ne y).
\]
\end{lemma}
\begin{proof}\leanok
Summing the edge equations at $v$ counts each cycle containing $v$ twice,
so $2\sum_{C\ni v}a_C=d_Q(v)$. Counting actual cycle subgraphs through
two fixed vertices bounds their total $a$-weight by $r^2B$. Rescaling
and $d\ge n/2$ give the assertions.
\end{proof}
The almost-regular hypothesis is not supplied by every current application
of the manuscript's cycle-factor theorem. This lemma therefore does not
by itself replace that theorem or establish the required absorber construction.

\begin{lemma}[Bounded perfect fractional matchings]
\label{bp:bounded-perfect-fractional-matching}
\lean{Oberwolfach.bounded_bipartite_fractional_matching,Oberwolfach.bounded_perfect_fractional_matching}
\leanok
If $0<\beta<1/2$, $m\ge4/\beta$, and $Q$ is an $m$-vertex graph
with minimum degree at least $(1/2+\beta)m$, then $Q$ has a perfect
fractional matching with every edge coefficient in
$[\beta/(5m),5/(\beta m)]$.
\source{lem:selection:bounded-fractional-matching (live-source amendment)}
\end{lemma}
\begin{proof}\leanok
The Lean proof uses the previously proved real-flow cut criterion on the
bipartite double cover. A baseline $\beta/(5m)$ is placed on every edge;
the remaining row and column demands lie in $[1-\beta/5,1]$.
The dense cut inequalities give an additional coefficient at most
$4/(\beta m)$ on each edge. Symmetrisation gives vertex sums one in $Q$.
The live manuscript instead contains the authors' shorter Hall proof.
Both proofs establish the same statement, without the cited literature lemma.
\end{proof}

\begin{theorem}[Fractional almost-perfect matching]
\label{bp:fractional-almost-perfect-matching}
\lean{Oberwolfach.almost_perfect_matching_of_fractional_weights}
\leanok
Let $r\ge2$, $0<\xi\le1$, and
$\mu=(\xi/(100r))^{20r}$. For a finite indexed $r$-uniform hypergraph,
suppose nonnegative edge weights have vertex loads in $[1-\mu/4,1]$
and pair loads at most $\mu/4$. Then a matching covers at least
$(1-\xi)|V|$ vertices.
\end{theorem}
\begin{proof}\leanok
Replace edge $e$ by $\lfloor Tw_e\rfloor$ separately indexed copies,
where $\mu T>4(|E|+1)$. The degree rounding loss is at most $|E|$.
With $D_0=T$ and $B=\lceil\mu T/2\rceil$, the checked matching theorem
of Norin and Steiner applies. Its six-file dependency closure is included,
with licence, pinned revision, and the compatibility patch recorded.
Projection back to original indices preserves disjointness and coverage.
The displayed constant is the one in the formal development, with an
additional margin for rounding; it is not the larger paper constant.
This theorem alone does not control several prescribed weight sums.
\end{proof}

\begin{theorem}[Almost-spanning cycle packing]
\label{bp:almost-spanning-cycle-packing}
\lean{Oberwolfach.almost_spanning_cycle_packing}
\uses{bp:positive-ordinary-cycle-weights,bp:cycle-vertex-matching-weights,bp:fractional-almost-perfect-matching}
\leanok
For $r\ge3$, $\alpha>0$ and $0<\xi\le1$, there are $\eta>0$ and $N$
such that every graph $Q$ of order $n\ge N$ with
$\delta(Q)\ge(\delta_{C_r}+\alpha)n$ and
$(1-\eta)d\le d_Q(v)\le d$ for all $v$ has vertex-disjoint
$r$-cycles covering at least $(1-\xi)n$ vertices.
\end{theorem}
\begin{proof}\leanok
Use positive bounded fractional cycle weights, convert edge loads to
vertex loads, and apply the preceding theorem. The order threshold bounds
$4r^2B/n$ by $\mu/4$. All constants precede the graph and the degree scale.
The conclusion is an approximate factor, not an exact factor or an absorber.
\end{proof}

\begin{lemma}[Uniform numerical conditions for a colouring round]
\label{bp:colour-round-numerics}
\lean{Oberwolfach.colouring_round_eventually_valid,Oberwolfach.colouring_round_ratio_progress,Oberwolfach.colouring_round_list_lower}
\uses{bp:colour-round-realization}
\leanok
For fixed $k,K,\delta,\tau>0$ with $K\ge1$ and $\delta\le1$, one order
threshold ensures all the explicit concentration conditions for every
$l$ above that threshold, every $0<s\le l$, every codegree bound
$B\le K l^{1-\delta}$, and every dependency-degree bound $d\le K l^2$.
The robust ratio-progress and list-size inequalities are proved separately.
\end{lemma}
\begin{proof}\leanok
The concentration denominator is at most $A_k l(B+1)$ and the dependency
factor is at most $(k+1)(K+1)^4l^{10}$. A stretched exponential in
$l^\delta$ dominates this polynomial uniformly over all supplied choices.
\end{proof}
The finite iteration and its fixed-error prescribed-weight application are proved below.


\chapter{Completed colouring iteration and finite weighted matching}
\begin{theorem}[Near-optimal colouring with a fixed error]
\label{bp:near-optimal-colouring}
\lean{Oberwolfach.uniform_improving_colour_round, Oberwolfach.near_optimal_hypergraph_list_colouring, Oberwolfach.near_optimal_hypergraph_edge_colouring}
\uses{bp:colour-round-realization,bp:colour-round-numerics}
\leanok
For fixed $k\ge1$, $0<\epsilon\le1/2$ and $0<\delta\le1$, there is
$D_0$ such that every finite indexed $k$-uniform hypergraph of maximum degree
at most $D\ge D_0$ and codegree at most $D^{1-\delta}$ has a proper edge
colouring with $\lceil(1+\epsilon)D\rceil$ colours. Parallel edge indices
remain distinct. A list-colouring form is proved as well.
\end{theorem}
\begin{proof}\leanok
The one-round estimates give integer list sizes and degree bounds while
increasing the list-to-degree ratio by a fixed positive amount. A positive
fraction of each list survives each round. Both constants are chosen before
the auxiliary degree bounds. A fixed finite induction maintains these
bounds on the actual uncoloured edge indices, and the final local-lemma
colouring extends the retained partial colouring. All floor and ceiling
errors and all exponential-tail requirements are proved. No near-optimal
colouring theorem is assumed.
\end{proof}
The error $\epsilon$ is fixed before $D_0$. This theorem alone does not
establish a power-rate error $D^{-a}$.

\begin{theorem}[Finite simultaneous weighted-matching construction]
\label{bp:finite-weighted-matching}
\lean{Oberwolfach.weighted_edge_partition_exists, Oberwolfach.partition_colouring_weighted_matching, Oberwolfach.weighted_matching_on_disjoint_parts, Oberwolfach.finite_weighted_hypergraph_matching}
\uses{bp:near-optimal-colouring,bp:common-event-probability,bp:categorical-concentration}
\leanok
Fix nonnegative hyperedge weights $w_a$ and put $W_a=\sum_e w_a(e)$.
Under the three explicit numerical common-event bounds in the Lean statement,
there is a matching satisfying, simultaneously for every supplied test $a$,
\[
 \left|w_a(M)-\frac{W_a}{p^{k-1}q\ell}\right|
 < t_a+\frac{p s_a}{q\ell},
 \qquad \ell=\lceil(1+\epsilon)D\rceil.
\]
Here $p$ is the number of vertex parts, $q$ the number of edge parts,
$s_a$ the error for each first-stage part, and $t_a$ the final selection
error. The precise degree, codegree, weight and tail bounds are stated in
the linked Lean theorem.
\end{theorem}
\begin{proof}\leanok
One vertex-colouring outcome controls every omitted-root degree and every
fixed weight sum. One subsequent edge-colouring outcome controls all
part degrees, codegrees and weight totals. Proper colourings of these parts
are constructed by the preceding theorem. Independent colour-class choices
then give one matching satisfying every test. Projection to the original
hyperedge indices is injective and preserves the selected weight sums.
Summing the first-stage errors and dividing by the actual colour count gives
the displayed error. All three existence steps use proved concentration
inequalities and a common-event argument.
\end{proof}
The numerical hypotheses are discharged in the fixed-error uniform theorem
below. The support reduction, integer-copy and projection steps are proved
in Node~\ref{bp:fractional-resource-matching}.


\chapter{Uniform weighted matching and support reduction}
\begin{theorem}[Uniform weighted matching with fixed relative error]
\label{bp:uniform-weighted-matching}
\lean{Oberwolfach.matching_integer_root_scale, Oberwolfach.matching_edge_split_scale, Oberwolfach.matching_power_scales, Oberwolfach.matching_normalization_error, Oberwolfach.weighted_matching_large_scales, Oberwolfach.uniform_weighted_hypergraph_matching}
\uses{bp:finite-weighted-matching,bp:near-optimal-colouring}
\leanok
Fix integers $k\ge2$, $B\ge0$ and real numbers $K,a,\gamma>0$.
There is $n_0$ such that the following holds whenever $n\ge n_0$
and $n\le D\le n^B$. Let $\mathcal H$ be a finite indexed
$k$-uniform hypergraph with no isolated vertices, at most $n^B$ vertices,
maximum degree at most $D$ and maximum codegree at most $KD/n$.
For a family of at most $n^B$ functions $w_i:E(\mathcal H)\to[0,b_i]$,
where $b_i>0$, put $W_i=\sum_e w_i(e)$ and suppose $W_i\ge a b_iDn$.
There is a matching $M$ such that
\[
 |w_i(M)-W_i/D|<\gamma W_i/D
 \qquad\text{for every prescribed }i.
\]
The threshold is chosen before the hypergraph and the weight functions.
\end{theorem}
\begin{proof}\leanok
Choose $\eta=\min\{\gamma/20,1/100\}$ and
$h=\lfloor n^{1/(64k)}\rfloor$ using the integer root operation.
Take $p=h^{16}$ vertex parts, $q=\lfloor D/h^{16k-12}\rfloor$ edge
parts, cell degree bound $d=\lceil(1+4\eta)h^4\rceil$ and
$\ell=\lceil(1+\eta)d\rceil$ colours.
The six exponent bounds for the three common events are at most $-ch$
for a fixed $c>0$. Their union bounds are dominated by a fixed polynomial
in $h$ times $e^{-ch}$, and hence are less than one for a threshold chosen
before the host and tests. The integer inequalities give
$D\le p^{k-1}q\ell\le(1+8\eta)D$. The accumulated weight error and
this normalization error are together less than $10\eta W_i/D$.
All integer rounding, positivity and asymptotic estimates are proved in Lean.
\end{proof}

\begin{lemma}[Preserving finitely many measurements with sparse support]
\label{bp:measurement-sparsification}
\lean{Oberwolfach.finite_measurement_sparsification}
\leanok
For finite sets $E,J$, nonnegative weights $w:E\to\mathbb R$ with positive
total mass and vectors $z_e\in\mathbb R^J$, there are nonnegative weights
$w'$ on at most $|J|+1$ original indices such that
\[
 \sum_e w'(e)=\sum_e w(e),\qquad
 \sum_e w'(e)z_e=\sum_e w(e)z_e.
\]
Moreover, $w(e)=0$ implies $w'(e)=0$.
\end{lemma}
\begin{proof}\leanok
Normalize by the positive total mass and apply the proved Caratheodory
convex-hull theorem to the vectors indexed by the original positive support.
Choose original indices for its affine-independent representation and
aggregate their coefficients. The affine dimension bounds the new support.
\end{proof}
The following theorem completes the fractional-resource rounding argument.
Preparation, cover-down, terminal absorption and the main decomposition
theorem still require their graph-specific proofs.


\begin{theorem}[Rounding fractional weights with prescribed local estimates]
\label{bp:fractional-resource-matching}
\lean{Oberwolfach.fractional_resource_matching}
\uses{bp:uniform-weighted-matching,bp:measurement-sparsification}
\leanok
Fix integers $r\ge2$, $B\ge1$ and $K,a,\gamma>0$.
There are $\eta_0>0$ and $n_0$ such that the following holds for $n\ge n_0$.
Let an indexed hypergraph have at most $n^B$ vertices and edge sizes
between $2$ and $r$. Suppose nonnegative edge weights $w$ have vertex
loads in $[1-\eta_0,1+\eta_0]$ and pair loads at most $K/n$.
Let $\mathcal S$ be at most $n^B$ vertex subsets of size at least $an$,
and let $\mathcal F$ be at most $n^B$ fixed functions $g:E\to[0,1]$
with $W_g=\sum_e g(e)w(e)\ge an$. There is a matching $M$ such that
\[
 |S\setminus V(M)|\le\gamma |S|\quad(S\in\mathcal S),\qquad
 (1-\gamma)W_g\le g(M)\le(1+\gamma)W_g\quad(g\in\mathcal F).
\]
The thresholds are independent of the number of original edge identifiers.
\source{lem:new:weighted-resource-matching}
\end{theorem}
\begin{proof}\leanok
Use $\eta_0=\min\{\gamma/20,1/100\}$. Preserve all vertex, ordered-pair
and prescribed-weight sums on at most $4n^{2B}$ indices. Replace each
weight by $\lfloor Tw'(e)\rfloor$ copies, where $T=n^{2B+2}$, and pad
each copy to rank $r$ with private vertices. Original degrees lie between
$(1-2\eta_0)T$ and $(1+\eta_0)T$; private degrees are one.
Take $D=\lceil(1+\eta_0)T\rceil$. The enlarged vertex set and test family
have size at most $n^{3B+4}$, and the pair loads are at most $(K+1)D/n$.
The functions $e\mapsto |e\cap S|$ and the lifts of $g$ satisfy the
uniform weighted theorem's mass bound with coefficient $a/(4r)$.
Apply that theorem with relative error $\eta_0$ and project its matching.
Nonempty original edges make this projection injective, so selected test
sums are preserved exactly. The proved floor and normalization bounds give
the displayed conclusions. The empty vertex set is handled separately.
\end{proof}


\begin{theorem}[Ordinary wheel packing with at most $n$ hubs]
\label{bp:ordinary-wheel-packing-bounded-hubs}
\lean{Oberwolfach.ordinaryResourceIncident, Oberwolfach.ordinaryIncidentEdges_card_bound, Oberwolfach.ordinaryActiveResources_matching_physical, Oberwolfach.ordinary_wheel_packing_bounded_hubs}
\uses{bp:fractional-resource-matching,bp:ordinary-wheel-potential-stability}
\leanok
Fix $r\ge3$ and $C,\gamma>0$. For sufficiently large $n$, let an
ordinary extended host have $n$ rim vertices and at most $n$ hubs.
Suppose its allowed $r$-wheels have nonnegative weights at most
$Cn^{1-r}$ with exact load one on every present rim and spoke edge.
There is a family of physically edge-disjoint allowed wheels leaving at
most $\gamma n$ uncovered incident edges at every rim vertex and hub.
This is the bounded-hub version; the manuscript also states a more general
version with an arbitrary fixed bound $C_Vn$ on the total number of vertices.
\end{theorem}
\begin{proof}\leanok
Use actual present rim and spoke edges as resources. Every wheel has
$2r$ resources, the fractional degrees are one, and the pair loads are at
most $r^4C/n$. There are at most $2n^2$ resources and $2n$ physical vertices.
Apply fractional-resource matching with exponent $3$ to the incident sets
of size at least $\gamma n/2$, using relative error
$\min\{1/2,\gamma/4\}$. Every incident set has size at most $2n$.
The omitted smaller sets already have fewer than $\gamma n/2$ edges.
The formal projection proves disjointness of the actual wheel edge sets.
\end{proof}

\begin{lemma}[Integer demands for sparse corrections]
\label{bp:sparse-integer-demands}
\lean{Oberwolfach.balanced_integer_half_degrees, Oberwolfach.near_uniform_integer_demands, Oberwolfach.residue_in_integer_interval, Oberwolfach.signed_balanced_degree_demands, Oberwolfach.residue_degree_demands}
\leanok
A nonnegative integer degree sequence with even total has two integer
halves with equal totals and pointwise difference at most one. A prescribed
total can be distributed over a nonempty finite set within one of its mean.
Every interval of $r$ consecutive integers contains the required residue
modulo $r$. These choices give the spoke demands and the two oriented rim
demands used in sparse ordinary correction. The signed imbalance is
preserved exactly.
\end{lemma}
\begin{proof}\leanok
Choose half of the vertices with odd demands for upward rounding. For a
prescribed total $T$ on $n$ vertices use $\lfloor T/n\rfloor$, adding one
at $T\bmod n$ vertices. For the hub demands choose the representative of
the prescribed residue in the given interval. The formal mean estimates
retain the rounding errors and the signed total imbalance.
\end{proof}

\begin{theorem}[Sparse ordinary degree correction]
\label{bp:sparse-ordinary-correction}
\lean{Oberwolfach.sparse_integer_correction_matrices, Oberwolfach.sparse_correction_numeric, Oberwolfach.graph_of_suborientation, Oberwolfach.sparse_ordinary_degree_correction}
\uses{bp:sparse-integer-demands,bp:normalized-demands}
\leanok
The manuscript's sparse ordinary degree-correction lemma is proved for
actual simple rim graphs, spoke relations and a supplied orientation of
the allowed rim edges. The deletions avoid the forbidden edges, make
$d_{Q-J_Q}(v)=2d_{H-J_H}(v)$ at every rim vertex, leave every hub degree
divisible by $r$, and remove at most $10\xi n$ incident edges at any vertex.
The order threshold precedes the graphs, orientation and forbidden sets.
\source{lem:new:relative-ordinary-correction}
\end{theorem}
\begin{proof}\leanok
For $0<\xi\le h/100$, use $\eta_0=\delta_0\omega\xi/1000$ and choose
$n_0$ so that $100(r+2)\le\delta_0\omega\xi n$ for $n\ge n_0$.
Choose hub deletion degrees in
$[\lfloor\xi n\rfloor,\lfloor\xi n\rfloor+r)$ with the required
residues; distribute their total almost uniformly over the rim vertices.
Split $2h_v+d_Q(v)-2d_H(v)$ into equal-total oriented demands.
Apply the proved integral normalized-cut theorem to both allowed incidence
matrices. The supplied orientation has only one arc per physical edge,
so the selected matrix gives a simple deletion graph. The formal proof
checks every capacity, relative-error, parity and degree bound.
\end{proof}

\begin{theorem}[Ordinary configuration counts after sparse deletions]
\label{bp:ordinary-configuration-deletion}
\lean{Oberwolfach.finite_configuration_loss, Oberwolfach.ordinary_transfer_deletion_stability, Oberwolfach.ordinary_isolator_deletion_stability}
\uses{bp:positive-ordinary-wheel-weights}
\leanok
Suppose the deleted rim graph and spoke relation have maximum degree at
most $d$, and there are $n\ge1$ rim vertices and at most $n$ hubs.
At a surviving rim root $xy$, the number of ordinary $k$-set isolator
configurations decreases by at most $8dn^{k-2}$, for $k\ge4$.
At a rim vertex and two surviving root spokes, the number of labelled
$r$-wheel transfer configurations decreases by at most
$5r^4dn^{r^2-1}$, for $r\ge3$.
The parent graph, forbidden graph, minimum-degree conditions and label
representations are unchanged. This statement concerns surviving roots;
it does not yet assert the manuscript's absent-root potential identities.
\end{theorem}
\begin{proof}\leanok
Map every original configuration avoiding the deleted incidences into
the new configuration family, retaining its labels and hub. This map is
injective. Count failures by a particular deleted edge. An edge incident
with a fixed root fixes one free label; an edge with two free endpoints
fixes two. Sum the resulting bounds using the actual row and column
degrees of the deleted relation. The transfer's two fixed root spokes
are explicitly required to survive.
\end{proof}

\begin{theorem}[Restriction of fixed ordinary wheel weights]
\label{bp:ordinary-wheel-restriction}
\lean{Oberwolfach.ordinaryWheelRestrictionEquiv, Oberwolfach.ordinary_wheel_restricted_resource_load, Oberwolfach.ordinary_restricted_rim_load_stability, Oberwolfach.ordinary_restricted_spoke_load_stability, Oberwolfach.ordinary_packing_transport}
\uses{bp:ordinary-wheel-potential-stability}
\leanok
Restriction to a smaller host retains each surviving wheel and its supplied
coefficient. If $0\le w(W)\le Bn^{1-r}$ and the deleted rim and spoke
maximum degrees are at most $d$, the load at a surviving rim edge changes
by at most $8r^4Bd/n$, and the load at a surviving spoke changes by at most
$4r^4Bd/n$. A packing in the smaller host is a packing in the original
host. At each vertex, adding back deleted edges increases the uncovered
incident-edge count by at most the degree of the deletion.
\end{theorem}
\begin{proof}\leanok
Use the bijection between surviving parent wheels and wheels in the
smaller host. The physical edge sets and coefficients are unchanged.
Apply the proved root-omitted load bounds and identify physical incident
sets with the active-edge representation. The final comparison is the
finite-set inequality
$|S\setminus C|\le |(S\setminus D)\setminus C|+|D|$.
\end{proof}

\begin{theorem}[Ordinary near-balance wheel packing]
\label{bp:ordinary-near-balance-packing}
\lean{Oberwolfach.ordinary_near_balance_wheel_packing}
\uses{bp:sparse-ordinary-correction,bp:ordinary-configuration-deletion,bp:ordinary-wheel-restriction,bp:positive-ordinary-wheel-weights,bp:ordinary-wheel-packing-bounded-hubs}
\leanok
The ordinary packing lemma is proved with the manuscript's near-balance,
large-cut, configuration-count and positive fixed-weight hypotheses.
The initial load tolerance $\varepsilon_*>0$ is chosen before the requested
remainder bound $\gamma>0$. The resulting packing consists of actual
pairwise physically edge-disjoint wheels and leaves at most $\gamma n$
uncovered incident edges at every rim or hub vertex.
\source{lem:new:relative-ordinary-packing}
\end{theorem}
\begin{proof}\leanok
First choose the exact-weight tolerance $\varepsilon_W$ with configuration
supply $\beta/2$, and put $\varepsilon_*=\varepsilon_W/4$.
After $\gamma$, choose $\xi$ so that
$10\xi\le\gamma/2$, $80\xi\le\beta/2$,
$50r^4\xi\le\beta/2$ and
$80r^4b_+\xi\le\varepsilon_W/4$.
The sparse correction has maximum degree at most $10\xi n$.
The proved deletion estimates preserve the required supplies and load
errors. Correct the restricted weighting exactly, apply the proved wheel
matching theorem with remainder $\gamma n/2$, and transport its wheels
back to the original host. All thresholds precede the host and weighting.
\end{proof}


\begin{lemma}[Marked physical resources and pair loads]
\label{bp:marked-wheel-resources}
\lean{Oberwolfach.markedCyclePhysicalEdges_injective, Oberwolfach.markedWheelPhysicalResources_card, Oberwolfach.markedActivePhysical_injective, Oberwolfach.markedActiveResources_image, Oberwolfach.marked_active_exact_loads, Oberwolfach.marked_cycle_arc_pair_root_card, Oberwolfach.marked_wheel_active_pair_load}
\uses{bp:marked-cycle-representation}
\leanok
Let $c\ge4$ and let the physical map on the typed rim arcs be injective.
Use the present rim arcs and spoke pairs as resources. A marked $c$-wheel
uses exactly $2c$ such resources and its image is its physical edge set.
For weights at most $Cn^{1-c}$ and at most $n$ hubs, any two distinct
resources have total wheel weight at most $c^4C/n$. If the spoke types are
disjoint, exact typed rim and spoke loads give load one at each resource.
\end{lemma}
\begin{proof}\leanok
Two distinct arcs in a common auxiliary cycle fix at least three vertices:
coinciding unordered endpoint pairs would force equal positions or a
$2$-cycle. For $s$ fixed auxiliary vertices and a hub set $T$, the event
weight is at most $|T|c^s n^{c-s}Cn^{1-c}$.
Apply this bound to rim--rim, rim--spoke and spoke--spoke pairs.
The map from active resources to physical edges is injective. No marked
head is identified with its physical image before applying the supplied
permutation.
\end{proof}

\begin{theorem}[Marked wheel packing with both leftover bounds]
\label{bp:marked-wheel-packing-bounded-hubs}
\lean{Oberwolfach.markedResourceIncident, Oberwolfach.markedIncidentEdges_card_bound, Oberwolfach.markedIncidentEdges_physical_image, Oberwolfach.marked_leftover_card_physical, Oberwolfach.marked_wheel_packing_bounded_hubs}
\uses{bp:marked-wheel-resources,bp:fractional-resource-matching}
\leanok
Fix $c\ge4$ and $C,\gamma>0$. For sufficiently large $n$, consider a marked
host with $n$ rim vertices, at most $n$ hubs, disjoint spoke types and an
injective physical rim map. Nonnegative weights at most $Cn^{1-c}$ with
exact typed loads one yield a family of physically edge-disjoint marked
wheels. Both auxiliary and physical uncovered incident-edge counts are at
most $\gamma n$ at every vertex. The length threshold and size threshold
precede the host, permutation, separation parameter and weights.
\end{theorem}
\begin{proof}\leanok
Apply fractional-resource matching once to both families of incident sets.
There are at most $3n^2$ resources, at most $4n$ tests, and each incident
set has size at most $5n$. Include sets of size at least $\gamma n/2$ and
use relative error $\min\{1/2,\gamma/10\}$. Smaller sets satisfy the bound
without a matching estimate. The physical resource map preserves both
set differences and their cardinalities, proving the physical conclusion
for the same selected matching. This proves the exact-load rounding step used in the near-balance
assembly below.
\end{proof}


\begin{lemma}[Four sparse typed deletion matrices]
\label{bp:sparse-typed-correction-matrices}
\lean{Oberwolfach.typed_degree_demands, Oberwolfach.typed_integer_correction_matrices, Oberwolfach.typed_correction_numeric, Oberwolfach.sparse_typed_correction_matrices, Oberwolfach.TypedMatrixBalanced, Oberwolfach.typed_matrix_remainder_balanced, Oberwolfach.sparse_typed_balanced_remainder, Oberwolfach.marked_permuted_incidence_bound}
\uses{bp:sparse-integer-demands,bp:normalized-demands}
\leanok
The four typed correction matrices can be constructed inside the original
allowed matrices. Their removed outdegree, indegree, marked-indegree and
hub-ratio differences equal the four supplied integer imbalances.
Subtracting these matrices from the original host makes all four degree
equations exact. With $q=c-1$, each individual row or column sum is at most
$2c\xi n$. The parameters are chosen before the host and its forbidden
submatrices. Physical incidence counts are bounded after any injective
embedding into an ambient set and any permutation of that set; no
invariance of the original vertex subset is required.
\source{lem:four:typed-correction}
\end{lemma}
\begin{proof}\leanok
Put $k=\lfloor\xi n\rfloor$, $s=|W|$ and $\mu=sk/n$.
Choose integer spoke rows with totals $sk$ and
$H=qsk+\sum_w a^W(w)$. The marked rim indegrees are $a_v+a^K(v)$;
choose its outdegrees $k_v$ with the same total and error at most one.
The ordinary rim demands are
$h_v+a_v+a^+(v)-k_v$ and $h_v+a^-(v)-a^K(v)$.
Their totals agree since $\sum a^+=\sum a^-$.
The marked demand errors are at most $1+2e$ and the ordinary errors at
most $3+2e$ and $1+4e$, where $e=\eta n$.
The explicit inequalities
$3+4e\le\delta(\mu-3e)$ and $qk+3e\le hn/2$
allow four applications of the proved normalized integral cut theorem.
All degree identities use integer subtraction only after proving the
subgraph inequalities. Individual degree bounds follow from the same
mean estimates. Injectivity bounds each physical endpoint fiber by one,
including when the permutation maps a vertex outside the original subset.
The incidence-matrix conclusion is realized on the original host and
used in the near-balance packing assembly below.
\end{proof}


\begin{lemma}[Restriction of marked wheels and fixed weights]
\label{bp:marked-wheel-restriction}
\lean{Oberwolfach.markedWheelRestrictionEquiv, Oberwolfach.marked_wheel_restricted_resource_load, Oberwolfach.markedWheelToLarger_matching, Oberwolfach.markedPotentialLoad_restricted, Oberwolfach.markedWheelSurvives_iff_typed_occurrences}
\uses{bp:marked-wheel-resources}
\leanok
Restricting either rim type and either spoke type gives a bijection between
surviving parent wheels and the canonical wheels of the smaller host. The
vertex tuple, hub, permutation and separation parameter are unchanged.
The inclusion preserves physical resource sets, fixed-weight resource sums,
and physically edge-disjoint matchings. The typed survival criterion needs
only endpoint injectivity separately in each rim type, not global physical
injectivity or disjoint parent spoke relations.
\end{lemma}
\begin{proof}\leanok
Use the same tuple and hub in both directions. Endpoint injectivity
identifies the chosen typed rim witnesses. At position zero retain the
marked spoke; at each other position retain the ordinary spoke. Transport
sums through this equivalence without recomputing weights.
\end{proof}

\begin{lemma}[Root-omitted marked loads under sparse deletions]
\label{bp:typed-marked-potential-stability}
\lean{Oberwolfach.resourceRootOmittedLoad_change, Oberwolfach.resourceRootOmittedLoad_mono, Oberwolfach.resourceOmittedPairWeight_image_le, Oberwolfach.typedMarkedOccurrenceResources_forget_inj, Oberwolfach.typedMarkedPotentialLoad_deletion, Oberwolfach.potentialMarkedHubs_deletion}
\uses{bp:marked-wheel-resources}
\leanok
Let $c\ge4$, $|V|=n>0$, $|U|\le n$, and fix canonical marked wheels with
weights in $[0,Bn^{1-c}]$. The rim occurrence universe has separate ordinary
and marked types; each spoke also retains its type. The fixed candidate
universe need not have globally injective physical rim images or disjoint
spoke relations. If the two states are nested and the removed rim
indegrees and outdegrees and the removed degrees in each spoke type are at
most $tn$, then for every typed potential root $e$,
\[
0\le L_e-L'_e\le 11c^4Bt.
\]
Here the root coordinate itself is omitted in each load. For every fixed
$c$-tuple, the compatible-hub loss is at most $ctn$. Neither the root nor
the tuple's rim edges need belong to either state.
\source{lem:packing:marked-potential-stability}
\end{lemma}
\begin{proof}\leanok
A changed non-root occurrence contributes its fixed pair weight. Two
distinct co-occurring rim arcs fix at least three auxiliary vertices;
disjoint auxiliary endpoint pairs fix four. A spoke fixes the hub.
These give costs at most $c^4B/n$ for edges incident with a fixed root
endpoint or hub, and $c^4B/n^2$ otherwise. Summing actual deletion degrees
gives $11c^4Bd/n$ at a rim root and $5c^4Bd/n$ at a spoke root.
Forgetting the spoke label is injective within each individual wheel, so
the pair estimate never identifies two co-occurring typed resources.
The bound is deterministic and simultaneous in the choice of root and
states; it does not by itself prove a probabilistic common event or the
marked near-divisibility packing theorem.
\end{proof}

\begin{lemma}[Common marked configurations after deletion]
\label{bp:marked-common-count-deletion}
\lean{Oberwolfach.commonMarkedCycleConstraintsEquiv, Oberwolfach.commonMarkedPathConstraintsEquiv, Oberwolfach.commonMarkedCycle_deletion_count, Oberwolfach.commonMarkedPath_deletion_count}
\uses{bp:marked-cycle-representation}
\leanok
Let $c=m+1\ge2$ and $|V|=n$. Delete at most $d$ arcs entering
or leaving each vertex in either rim type, and at most $d$ edges at each
hub in either spoke type. On the unchanged vertex and hub domains, the
number of common marked cycles decreases by at most $3cd n^{c-1}$.
For any two specified alternative root vertices, root position and two
hubs, the number of canonical common marked paths decreases by at most
$4cd n^{c-2}$. These path conditions omit all four root spokes.
No global physical-image injectivity is required for either estimate.
\source{lem:packing:marked-deletion-stability}
\end{lemma}
\begin{proof}\leanok
Use exact equivalences with tuples satisfying the unchanged separation
conditions and explicit rim and spoke constraints. A changed binary
constraint costs at most $d n^{s-1}$ on $s$ free coordinates. With one fixed
endpoint, use the corresponding incoming or outgoing degree bound.
A changed spoke constraint has the same bound. There are $3c$ constraints
for common cycles, and $2c+2(c-1)\le4c$ for common paths. This counts
actual canonical configurations in the restricted arc subtypes.
\end{proof}

\begin{lemma}[Realization of sparse typed corrections]
\label{bp:sparse-typed-host-correction}
\lean{Oberwolfach.occurrenceMultiplicity_matrix_restriction, Oberwolfach.occurrence_matrix_deleted_degree, Oberwolfach.occurrence_matrix_deleted_indegree, Oberwolfach.typed_matrix_remainder_realized, Oberwolfach.typed_matrix_balanced_occurrence_degrees, Oberwolfach.sparse_typed_host_correction}
\uses{bp:sparse-typed-correction-matrices}
\leanok
For an actual typed host satisfying the near-balance and allowed-cut
hypotheses of the sparse matrix correction, the four correction matrices
are realized by deletions of its original arcs and spokes. Ordered
endpoints are injective within each rim type. The remaining host satisfies
all four exact degree identities required by marked-wheel correction.
The removed degrees equal the selected matrix row and column sums, and
retain their bounds $2c\xi n$. Allowed subgraphs, endpoint maps and the
order of the parameter choices are unchanged.
\end{lemma}
\begin{proof}\leanok
Each original occurrence matrix entry is zero or one. An arc is retained
exactly when its correction entry is zero; a spoke is retained exactly
when it originally existed and its correction entry is zero. Pointwise
subtraction agrees with the actual remaining multiplicities. Sum these
identities to obtain the degree and balance conclusions. No new edge
identifiers or representation maps are chosen.
\end{proof}

\begin{lemma}[Fixed-profile estimates after deletion]
\label{bp:marked-host-estimate-deletion}
\lean{Oberwolfach.occurrenceDegree_restriction_target, Oberwolfach.occurrence_matrix_cut_deletion, Oberwolfach.represented_two_step_deletion, Oberwolfach.SparseIntersectionData.deletion}
\uses{bp:sparse-typed-host-correction}
\leanok
Deleting degree at most $d$ changes each directed degree by at most $d$,
each directed two-step count by at most $2d$, and an ordered cut between
$S,T$ by at most $d\min\{|S|,|T|\}$. For a bipartite spoke relation,
one-, two- and three-column intersections change by at most $d,2d,3d$.
If the fixed inverse-density profile is bounded by $L$, its weighted
row codegrees change by at most $2L^2d$. Consequently sparse intersection
data with error $\delta$ remain valid with error $\delta+\rho$, provided
$3d\le\rho|V|$ and $2L^2d\le\rho|U|$. The supplied profiles and exceptional
sets are unchanged, as are the two original normalising domains.
\end{lemma}
\begin{proof}\leanok
Sum the actual deleted incidences. For products of adjacency indicators,
charge a lost product to one of its changed factors. Apply the same
argument with the original coefficient $q(u)^2$ for weighted codegrees.
The cut estimate sums the deleted row or column bounds, whichever is
smaller. Triangle inequalities retain the original comparison targets.
These deterministic results do not assert the probabilistic common event
or the complete near-divisibility packing theorem.
\end{proof}


\begin{lemma}[Uniform marked correction and balanced packing]
\label{bp:marked-uniform-balanced-packing}
\lean{Oberwolfach.marked_rim_correction_uniform_parameters, Oberwolfach.positive_marked_wheel_weights_uniform_parameters, Oberwolfach.positive_marked_wheel_weights_numerical_uniform, Oberwolfach.marked_balanced_wheel_packing}
\uses{bp:positive-marked-wheel-weights, bp:marked-wheel-packing-bounded-hubs}
\leanok
The length threshold depends only on $\alpha,p_0$. After the length and
the remaining density, extension and weight parameters are given, choose
the host and weight error tolerances. For every separation parameter and
positive leftover tolerance, a sufficiently large exactly balanced marked
host satisfying the stated extension and sparse-intersection bounds has
one physical edge-disjoint wheel packing with the required auxiliary and
physical degree bounds on the leftover.
\end{lemma}
\begin{proof}\leanok
The capacity theorem chooses the length threshold before the later
parameters. Expose that order in the correction theorem, retaining the
previous declarations as corollaries. Correct the positive approximate
weighting, then apply the proved marked exact-load packing theorem to all
auxiliary and physical incident tests simultaneously.
\end{proof}

\begin{lemma}[Restriction of the fixed marked weighting]
\label{bp:marked-actual-load-restriction}
\lean{Oberwolfach.markedPotentialLoad_parent_rim, Oberwolfach.markedPotentialLoad_parent_spoke, Oberwolfach.markedPotentialLoad_restricted_rim, Oberwolfach.markedPotentialLoad_restricted_spoke, Oberwolfach.marked_restricted_rim_load_stability, Oberwolfach.marked_restricted_spoke_load_stability, Oberwolfach.marked_approximate_weights_restriction}
\uses{bp:marked-host-estimate-deletion}
\leanok
For a physically injective marked host with disjoint spoke types, restricting
a fixed weighting to the retained wheels preserves its individual lower
and upper bounds. If all aggregate deleted in- and outdegrees and spoke
row and column degrees are at most $d$, the actual rim loads change by at
most $8c^4Bd/n$ and spoke loads by at most $4c^4Bd/n$. Thus an initial
residual bound $\eta$ becomes at most $\eta+8c^4Bd/n$ at every retained
edge; absent spoke residuals are zero.
\end{lemma}
\begin{proof}\leanok
Identify each root-omitted potential with its actual load when that root
survives. Use the parent spoke union, rather than all possible spokes,
and the proved potential degree estimates. Apply the triangle inequality
with the original unit target. No coefficient is recomputed.
\end{proof}

\begin{lemma}[Retained marked extension counts]
\label{bp:marked-extension-lower-restriction}
\lean{Oberwolfach.marked_hub_count_restriction, Oberwolfach.marked_common_cycle_count_restriction, Oberwolfach.marked_replacement_pair_count_restriction}
\uses{bp:marked-common-count-deletion}
\leanok
The doubled extension bounds of the parent host imply the single-margin
bounds in the retained host whenever $4cd\le\beta n$, where each deleted
type has degree at most $d$. The root position, permutation and separation
condition are unchanged. The four root spokes in the smaller graph are
supplied explicitly when passing from common paths to replacement pairs.
\end{lemma}
\begin{proof}\leanok
Subtract the compatible-hub loss $cd$, the common-cycle loss
$3cd n^{c-1}$, and the common-path loss $4cd n^{c-2}$.
The path-to-pair map is injective and restores precisely the four root
spokes omitted from the path count.
\end{proof}

\begin{lemma}[Transport to the original marked host]
\label{bp:marked-packing-parent-transport}
\lean{Oberwolfach.markedActiveToLarger_injective, Oberwolfach.markedActiveToLarger_physical, Oberwolfach.markedActiveToLarger_incident, Oberwolfach.markedActiveResources_toLarger, Oberwolfach.marked_leftover_transport, Oberwolfach.marked_deleted_incidence_bound, Oberwolfach.marked_packing_transport_bounds}
\uses{bp:marked-wheel-packing-bounded-hubs, bp:sparse-typed-host-correction}
\leanok
A physical edge-disjoint packing in the retained host is the same physical
packing in the parent host. If every deleted degree in each of the four
types is at most $d$, its auxiliary and physical uncovered degrees increase
by at most $6d$. All original arc identifiers and the permutation are
preserved.
\end{lemma}
\begin{proof}\leanok
Use the injective inclusion of active resources to identify covered sets
and the retained incident sets in both representations. The uncovered
parent incident set is contained in the union of its uncovered retained
part and its deleted part. At a vertex there are four rim and two spoke
incidences; at a hub only the two spoke types occur. In the physical marked
indegree, use the auxiliary indegree at $\pi^{-1}(v)$.
\end{proof}

\begin{lemma}[Root-first and canonical common paths]
\label{bp:marked-common-path-coordinates}
\lean{Oberwolfach.rootedCycleOrder_reinsert, Oberwolfach.orderedCommonMarkedPath_to_canonical, Oberwolfach.orderedCommonMarkedPath_canonical_card_le, Oberwolfach.rootFirstCommonMarkedPath_canonical_card_le}
\uses{bp:marked-cycle-representation}
\leanok
Deleting the first coordinate in a root-first cyclic order and reindexing
its free coordinates gives an injective map into canonical common paths.
Reinserting either root commutes with this reindexing. No root spoke is
required by the map or its cardinality inequality.
\end{lemma}
\begin{proof}\leanok
Use the rooted cyclic order equivalence to recover each free coordinate.
After either root is reinserted, the rim conditions, the remaining spoke
conditions and separation are exactly the rotated original conditions.
The inverse coordinate map proves injectivity.
\end{proof}

\begin{theorem}[Marked near-balance packing]
\label{bp:marked-near-balance-packing}
\lean{Oberwolfach.marked_small_deletion_packing, Oberwolfach.marked_near_balance_wheel_packing, Oberwolfach.marked_near_balance_wheel_packing_numerical}
\uses{bp:sparse-typed-host-correction,bp:marked-uniform-balanced-packing,bp:marked-actual-load-restriction,bp:marked-extension-lower-restriction,bp:marked-packing-parent-transport,bp:marked-common-path-coordinates,bp:marked-host-estimate-deletion}
\leanok
For the supplied marked host, construct a sparse correction of all four
balance equations, retaining the original arc identifiers, permutation
and cycle weights. The doubled parent cut, two-step and extension bounds
supply the corresponding single-margin bounds in the retained host.
There is one physical edge-disjoint wheel packing with both auxiliary
and physical uncovered degree at most $\gamma n$. The capacity threshold
is fixed before the cycle length and all later weight and spoke parameters;
host and weight tolerances precede $\gamma$; the balancing tolerance
precedes the separation parameter, and the order threshold follows it.
The numerical form uses constant spoke densities, root-first common-path
counts and the original physical edge-load tests. Its sparse-intersection
smallness input is the combined sum of the error and the two finite-order
terms; the separate manuscript bounds imply it after halving their common
tolerance. The host permutation is on its full rim vertex set.
\source{lem:new:relative-marked-packing}
\end{theorem}
\begin{proof}\leanok
Choose a deletion fraction $\sigma$ satisfying all degree, cut, extension,
weight and intersection error bounds before constructing the correction.
For $c=k+2$, take $\xi=\min\{\sigma/(2c),h/(200c)\}$ in the sparse
four-type correction theorem. Every deleted row and column is at most
$2c\xi n\le\sigma n$. Apply balanced packing with error $\gamma/2$
after restricting the original weights, then add at most $6\sigma n$
when transferring the same wheels to the parent host. Take
$\sigma\le\gamma/12$. The physical and auxiliary conclusions use
this same matching. Root-first reindexing and constant-density
normalization give the numerical form without changing the weights.
\end{proof}


\begin{lemma}[Local dependence for fixed root-omitted patterns]
\label{bp:resource-pattern-locality}
\lean{Oberwolfach.resourceOmittedPairWeight_square_sum_le, Oberwolfach.locality_square_sum, Oberwolfach.resourcePatternLoad_supported_tail, Oberwolfach.resourcePatternLoad_supported_expectation, Oberwolfach.finiteLaw_pi_injective_projection_event}
\uses{bp:omitted,bp:bounded-differences,bp:common-event-probability}
\leanok
Let fixed nonnegative coefficients weight finite resource patterns, and
omit the tested root resource from each product. If the total coefficient
at a root is at most $L$, each remaining pair weight is at most $a$, and
each pattern has at most $q$ remaining resources, the sum of squared
resource influences is at most $aqL$. If each resource reads at most $B$
independent primitive coordinates and each primitive coordinate belongs
to at most $B$ resource supports, the squared primitive influence sum is
at most $B^2aqL$. Consequently the two-sided deviation probability at
$t>0$ is at most $2\exp(-2t^2/(B^2aqL))$ when this supplied variance
bound is positive. Exact expectations factor when the non-root resource
supports in each pattern are disjoint. Restriction to any fixed injective
subfamily of primitive coordinates preserves its product law.
\end{lemma}
\begin{proof}\leanok
Sum the pair weights by patterns, multiply by their maximum, and use
Cauchy--Schwarz on the fixed resource-coordinate incidence relation.
Bounded differences gives the tail, including a separate proof for zero
actual influence. Disjoint coordinate supports give the expectation
identity. These conclusions use a fixed support relation, not merely an
outcome-dependent bound on the number of changed states.
\end{proof}

\begin{lemma}[Permutation windows and raw interval geometry]
\label{bp:raw-interval-geometry}
\lean{Oberwolfach.permutationWindow_card, Oberwolfach.permutationWindow_comp_subset, Oberwolfach.rawInterval_extended_disjoint, Oberwolfach.rawIntervalState_supported, Oberwolfach.rawIntervalFamilyCoordinates_reverse_card}
\leanok
The ambient permutation window $[-a,b]$ has at most $a+b+1$ labels,
with equality when the permutation has no sufficiently short cycle.
Distinct isolated starts have disjoint successor-extended intervals.
The raw state at $x$ reads activation coordinates in $[-2d,d]$, start
colours in $[-d,0]$, and the original colour at $x$. Both the number read
and the number of label states that can read any coordinate are at most
$4d+3$. Labels at permutation distance greater than $3d$ have disjoint
coordinate supports.
\source{def:four:local-interval-law}
\end{lemma}
\begin{proof}\leanok
Represent windows by powers of the permutation and prove their duality,
composition and disjointness. Isolated starts cannot occur within distance
$d$. The assignment is determined by its possible starts, their isolation
tests, and the original label colour. Count these exact coordinate sets
and their reverse incidence sets.
\end{proof}

\begin{lemma}[The original raw interval law]
\label{bp:raw-interval-law}
\lean{Oberwolfach.isolatedStartProbability_scaled_lower, Oberwolfach.rawIntervalCovered_probability, Oberwolfach.rawIntervalUnblocked_probability, Oberwolfach.rawIntervalState_label_probability, Oberwolfach.rawIntervalState_separated_probability, Oberwolfach.rawIntervalFamily_column_independence, Oberwolfach.rawIntervalFamily_labels_probability}
\uses{bp:raw-interval-geometry,bp:resource-pattern-locality}
\leanok
Use the activation law $a_d=1/(2d+1)$, isolated-start probability
$r_d=a_d(1-a_d)^{2d}$, start-colour probabilities $q_g/(dr_d)$, and
original label-colour probabilities $s_j/\overline q$, where
$\overline q=1-(1+1/d)\sum_gq_g$. For $d\ge2$ and
$\sum_gq_g\le1/8$ and $\sum_js_j\le\overline q$, with nonnegative
$q_g,s_j$, all these laws are valid. Suppose the permutation has no
cycle of length at most $2d$. A label receives interval
colour $g$ with probability $q_g$, is outside all extended intervals with
probability $\overline q$, and receives ordinary colour $j$ with
probability $s_j$. States at separated labels are independent; the entire
vectors at distinct hubs are independent. These are the raw states before
uniform retention, whose law is not asserted to be a product law.
\source{lem:interval:raw-law}
\end{lemma}
\begin{proof}\leanok
Prove $dr_d>2/15$ and the positivity of $\overline q$. Exactly $d$
possible starts cover a label, and $d+1$ extended intervals contain it.
Their occurrence events are mutually exclusive and have probabilities
$q_g/d$. The uncovered event is independent of the original label colour.
Fixed coordinate supports prove the two independence statements.
\end{proof}

\begin{lemma}[Ordinary potential loads under the actual raw experiment]
\label{bp:ordinary-raw-sampling}
\lean{Oberwolfach.ordinary_wheel_resource_load_bound, Oberwolfach.ordinary_wheel_omitted_pair_bound, Oberwolfach.rawWheelCoordinates_candidate_disjoint, Oberwolfach.ordinaryRawPotentialLoad_expectation, Oberwolfach.ordinaryRawPotentialLoad_common_event, Oberwolfach.ordinaryRawPotentialLoad_adaptive_common_event}
\uses{bp:ordinary-wheel-potential-stability,bp:resource-pattern-locality,bp:raw-interval-law}
\leanok
Independently sample physical rim edges with fixed supplied probabilities,
and sample spokes using the original raw interval law. For ordinary
$r$-wheels avoiding the permutation conflict graph and fixed coefficients
$0\le w(W)\le Cn^{1-r}$, each root-omitted expected load is exactly the
sum of $w(W)$ times the product of all non-root marginal probabilities.
With $|U|\le n$, the total root coefficient is at most $r^2C$ and
each root/non-root pair coefficient at most $r^4C/n$. The variance bound
is therefore $(4d+3)^2(r^4C/n)(2r)(r^2C)$.
One union-bound event controls all fixed physical rim and spoke roots.
On that same event, every later modification of maximum rim and both
spoke difference degrees at most $a$ incurs additional load error at most
$8r^4Ca/n$ for a rim root and $4r^4Ca/n$ for a spoke root.
The seed-specific normalization and distance from the target one are
proved in the following node.
\source{lem:completion:ordinary-reference}
\end{lemma}
\begin{proof}\leanok
Physical rim coordinates and the three sorts of interval coordinates
remain distinct. Candidate separation proves disjoint primitive supports
for all remaining resources of a wheel. Apply the exact raw marginal
identities and the fixed-support concentration result. The deterministic
potential-load stability inequalities hold for every later modification,
so no new event or sampling-density multiplier is introduced.
\end{proof}


\begin{lemma}[Admissible cycle weights and full-host seed means]
\label{bp:ordinary-seed-means}
\lean{Oberwolfach.ordinary_cycle_rim_conflict_loss, Oberwolfach.ordinary_cycle_vertex_conflict_loss, Oberwolfach.ordinary_admissible_cycle_edge_load, Oberwolfach.ordinary_admissible_cycle_vertex_load, Oberwolfach.permutationConflictGraph_degree_le, Oberwolfach.ordinaryFullHostSeed_rim_mean, Oberwolfach.ordinaryFullHostSeed_spoke_mean, Oberwolfach.ordinaryFullHostSeed_bounds, Oberwolfach.ordinaryFullHostSeed_common_event, Oberwolfach.ordinary_load_failure_eventually_small, Oberwolfach.ordinaryPotentialLoad_delete_conflicts}
\uses{bp:positive-ordinary-cycle-weights,bp:ordinary-raw-sampling,bp:raw-interval-geometry}
\leanok
Let $\lambda$ be a nonnegative fractional $r$-cycle decomposition of $P$,
with $\lambda(C)\le C_Pn^{2-r}$, and let $\Delta(F)\le L$.
At an edge $xy\in E(P)\setminus E(F)$ the weight of cycles containing an
$F$-edge among their vertices is at most $5r^4C_PL/n$; at a fixed vertex
it is at most $3r^4C_PL$. In particular $F_{\pi,b}$ has maximum degree
at most $2b$.
For each admissible cycle and each of the $m$ hubs use the unchanged
coefficient $\lambda(C)/(mp_R^{r-1}p_H^r)$. Its expected potential rim
load equals the admissible cycle weight through the edge. Its expected
potential spoke load is $p_R/(mp_H)$ times the admissible cycle weight
through the vertex. If $p_RD_0=2mp_H$, $D_0\ge n/2$ and
$|d_P(v)-D_0|\le\varepsilon n$, the two errors from one are at most
$5r^4C_PL/n$ and $2\varepsilon+12r^4C_PL/n$.
The lower and upper seed bounds are $c_Pn^{1-r}$ and
$C_P/(\omega_*p_*^{2r-1})\,n^{1-r}$, under the manuscript's supplied
size and density bounds. These constants precede $d$.
A single event controls all potential roots with these mean errors plus
any prescribed positive deviation $t$. Its explicit exponential failure
bound is eventually at most $n^{-q}$ for each fixed $q$.
Deleting conflict edges from the physical sample leaves every candidate
potential load exactly unchanged.
\source{lem:completion:ordinary-reference}
\end{lemma}
\begin{proof}\leanok
Bound excluded cycles by their specified vertices and sum over the
conflict edges, distinguishing the number of distinct prescribed vertices.
Reindex full-hub wheels by admissible cycles and hubs. The omitted resource
removes exactly one rim or spoke marginal, which cancels the corresponding
seed normalization. Apply the previously proved simultaneous concentration
bound. Every candidate rim avoids $F$, proving the exact deletion identity.
\end{proof}

\begin{lemma}[Common-neighbourhood counts in the raw experiment]
\label{bp:raw-neighbourhood-counts}
\lean{Oberwolfach.supportedPatternSum_tail, Oberwolfach.rawIntervalNeighbourCount_expectation, Oberwolfach.rawIntervalNeighbourCount_tail, Oberwolfach.rawIntervalNeighbourCount_common_event, Oberwolfach.rawIntervalHubNeighbourhood_common_event}
\uses{bp:resource-pattern-locality,bp:raw-interval-law}
\leanok
For a fixed label set $A$ and a nonempty fixed hub set $J$, count labels
in $A$ receiving the selected ordinary colour at every hub in $J$.
Its mean is $|A|\prod_{u\in J}s_u$. Its deviation at $t>0$ has probability
at most $2\exp(-2t^2/((4d+3)^2|J|n))$.
One event gives the corresponding estimates for any specified finite
family of fixed domains and hub sets. Applying it to $A=V$ and
$A=N_P(v)$ for every ambient vertex $v$ gives the sizes of the actual
random common neighbourhoods and every degree into them. The vertex $v$
is not required to belong to the random neighbourhood.
\source{eq:family:OF.A7}
\end{lemma}
\begin{proof}\leanok
Patterns at a fixed label involve distinct hub columns, so their
probabilities multiply. Each pair resource occurs in at most one pattern;
the sum of its squared coefficients is at most $|J|n$. Use the actual
fixed coordinate supports and the locality bound $4d+3$, then take a
finite union bound on the original fixed domains.
\end{proof}


\begin{lemma}[Dense conflict-free sets with prescribed roots]
\label{bp:rooted-dense-configurations}
\lean{Oberwolfach.fixedSubset_rooted_dense_probability, Oberwolfach.rootedSubset_additional_inclusion, Oberwolfach.rootedSubset_conflict_probability, Oberwolfach.rootedSubset_dense_conflict_free_card, Oberwolfach.rooted_dense_set_scales, Oberwolfach.rootedSubset_dense_conflict_free_many, Oberwolfach.enlarged_induced_degree, Oberwolfach.dense_rooted_sets_within, Oberwolfach.external_rooted_dense_sets, Oberwolfach.ordinary_isolator_candidate_many}
\uses{bp:fixed-size-concentration,bp:subset-inclusion,bp:raw-neighbourhood-counts}
\leanok
Fix $s$ prescribed roots, a positive density margin $\alpha$, and a positive
neighbourhood-size proportion $\beta$. There is $K$, chosen before the
conflict degree bound $L$, such that for each free-set size $k\ge K$ there
is $c>0$, also chosen before $L$, with the following property for sufficiently
large $n$. Suppose $|A|\ge\beta n$, the root set $Z$ is independent in $F$,
$\Delta(F)\le L$, and every vertex of $A\cup Z$ has at least
$(\theta+\alpha/2)|A|$ neighbours in $A$. For $0\le\theta\le1$ and
$0<\alpha\le1$, there are at least $cn^k$ sets $S$ of size $s+k$ with
$Z\subseteq S\subseteq A\cup Z$, with $F[S]$ empty, and with
$\delta(P[S])\ge(\theta+\alpha/4)(s+k)$.
In particular the prescribed roots need not belong to $A$.
\source{eq:family:OF.A17}
\end{lemma}
\begin{proof}\leanok
For a graph of minimum degree at least $(\delta+\varepsilon)n$, the
uniform $(s+k)$-set containing the roots has degree failure at most
$2(s+k)e^{-\varepsilon^2(s+k)/8}$, provided
$s+1\le\varepsilon(s+k)/4$; the target degree is
$(\delta+\varepsilon/2)(s+k)$. Additional
vertices have inclusion probability at most $(k/(n-s))^j$ for $j$ distinct
additional vertices. Summing over the conflict edges bounds their failure
probability by $L(2sq+nq^2)$, where $q=k/(n-s)$. Choose $k$ before $L$,
then choose $n$ so each failure is at most $1/4$. The exact count of rooted
sets is $\binom{n-s}{k}$, with lower bound $n^k/(2^k k!)$ when
$n\ge2(s+k)$. Enlarge $A$ by its prescribed roots and retain the explicit
bounded degree loss. The subtype-to-ambient map preserves cardinality,
all induced degree counts and conflict independence. This gives the stated
family of potential isolator candidates with only the free spokes required.
\end{proof}

\begin{lemma}[Simultaneous root-omitted physical survival]
\label{bp:rooted-physical-survival}
\lean{Oberwolfach.fixedSetFamily_distinct_edge_count, Oberwolfach.fixedSetFamily_normalized_edge_load, Oberwolfach.finite_pattern_incidence_square_sum, Oberwolfach.rootedInducedSurvivalSum_lower_tail, Oberwolfach.rootedInducedSurvivalSum_common_event, Oberwolfach.rootedInducedSurvivalSum_two_exposures, Oberwolfach.finiteLaw_pi_sum_probability, Oberwolfach.rawWheelRootedSurvival_common_event}
\uses{bp:rooted-dense-configurations,bp:bounded-differences,bp:common-event-probability,bp:raw-interval-law}
\leanok
Let an injectively represented family of $k$-sets contain fixed distinct
vertices $x,y$, where $k\ge4$, and give every set weight $n^{-(k-2)}$.
Test exactly the physical edges of $P[S]$ other than $xy$. The total weight
is at most one, each tested edge has incident weight at most $1/n$, and
the squared-influence sum is at most $k^2/n$.
If the family has at least $cn^{k-2}$ members and every supplied physical
edge probability is at least $p\in[0,1]$, the probability that the normalized
survival sum is at most $p^{k^2}c-t$ is bounded by
$2\exp(-2t^2/(k^2/n))$ for $t>0$.
One event controls every member of a fixed finite test family.
The candidate families may depend on the first, spoke exposure; the physical
edge law remains its original independent law. The joint good-event failure
is at most the first-exposure failure plus the physical union bound.
\source{lem:completion:ordinary-reference}
\end{lemma}
\begin{proof}\leanok
A physical edge distinct from $xy$ fixes at least one additional vertex,
so it belongs to at most $n^{k-3}$ candidate sets. Each candidate tests at
most $k^2$ edges. Sum the nonnegative edge incidences, multiply by their
maximum, and apply bounded differences. The exact product of the supplied
edge probabilities gives the expectation lower bound. Finite probability
slices combine both exposures, and an exact product-law identity transfers
the result to the original raw wheel experiment.
\end{proof}

\begin{lemma}[Potential counts and actual isolator configurations]
\label{bp:isolator-count-conversion}
\lean{Oberwolfach.rootOmittedInducedEdges_survive, Oberwolfach.rootedInducedSurvivalSum_card, Oberwolfach.ordinary_isolator_potential_count_eq}
\uses{bp:rooted-physical-survival}
\leanok
Once the root edge and its two endpoint spokes are present, the normalized
potential count equals $n^{-(k-2)}$ times the number of actual isolator
configurations at the specified hub. Deleting every conflict edge leaves
this equality unchanged.
\source{eq:reference:ordinary-isolator-sum}
\end{lemma}
\begin{proof}\leanok
Construct inverse maps between surviving potential candidates and the
physical configurations. All non-root spokes are already in the candidate
definition; the two remaining spokes and the root edge supply exactly the
omitted conditions. Each candidate is independent in the conflict graph,
so every edge of its induced physical graph survives conflict-edge deletion.
\end{proof}


\begin{lemma}[The ordinary potential isolator event]
\label{bp:raw-isolator-event}
\lean{Oberwolfach.rawIntervalCommonHubs_common_event, Oberwolfach.rawWheelCommonHubs_event, Oberwolfach.raw_ordinary_isolator_potential_event, Oberwolfach.ordinary_isolator_potential_lower, Oberwolfach.ordinary_isolator_total_count_lower, Oberwolfach.raw_ordinary_isolator_count_event}
\uses{bp:raw-neighbourhood-counts,bp:rooted-dense-configurations,bp:rooted-physical-survival,bp:isolator-count-conversion}
\leanok
The original interval/physical sampling law has one event giving positive
potential isolator counts for every admissible pair of distinct vertices.
The bounded set size and positive candidate constant are chosen before the
conflict-degree bound. The graph-order threshold is chosen afterwards.
The root edge is omitted in each per-hub survival test. The two root spokes
are retained when summing over hubs. Thus the normalized potential sum is
bounded below even if its physical root edge is absent. If that edge is
present, the event also gives the corresponding actual configuration count
in the sampled graph with the conflict edges deleted.
\source{lem:completion:ordinary-reference}
\end{lemma}
\begin{proof}\leanok
First obtain the sizes and every ambient-root degree into each sampled hub
neighbourhood. These imply the hypotheses of the proved rooted dense-set
count. Apply the physical-edge survival theorem conditionally, for all
potential pairs and hubs together. A separate common event counts the hubs
adjacent to both endpoints: its mean is $p_H^2|U|$, and the spoke dependencies
are treated through their original coordinate supports. Intersect the two
events without asserting independence. Sum the per-hub lower bounds, and
use the exact candidate/configuration correspondence only after imposing
root-edge presence.
\end{proof}

\begin{lemma}[Counted construction of the actual transfer configurations]
\label{bp:greedy-transfer-count}
\lean{Oberwolfach.sequential_choices_count, Oberwolfach.dense_common_neighbour_avoiding_used, Oberwolfach.rooted_pattern_embedding_count, Oberwolfach.transferVertexPosition_injective, Oberwolfach.transferPredecessors_card, Oberwolfach.transferConstructionGraph_base_cycle, Oberwolfach.transferConstructionGraph_row_cycle, Oberwolfach.transferLabelsOfPattern_rim, Oberwolfach.ordinary_transfer_greedy_count, Oberwolfach.ordinary_transfer_configurations_many, Oberwolfach.ordinary_transfer_candidates_many}
\uses{bp:ordinary-transfer-identities}
\leanok
Put $b=3$ for $r=3$ and $b=2$ for $r\ge4$. Suppose every vertex of
$A\cup\{v\}$ has at least $\delta|A|$ neighbours in $A$, the conflict graph
has maximum degree at most $L$, and the required free-label spokes are
present on $A$. Whenever
\[
0\le a\le [1-b(1-\delta)]|A|-(r^2+1)(L+1),
\]
there are at least $a^{r^2}$ labelled transfer configurations, with the
root spokes omitted. In particular, if $|A|\ge\beta n$ and
$\gamma=1-b(1-\delta)>0$, then for sufficiently large $n$ the lower bound
is $cn^{r^2}$, where $c=(\gamma\beta/2)^{r^2}>0$ is chosen before $L$.
\source{eq:loads:ordinary-potential-transfers}
\end{lemma}
\begin{proof}\leanok
Order the base-cycle vertices first. In each subsequent row, order the two
path endpoints first and then the internal vertices along the path.
The explicit predecessor sets have cardinality at most two when $r\ge4$;
in the triangle case the second endpoint has three predecessors.
The ordered positions are bijective with the $r^2$ free labels. Each base,
row and root-replaced row cycle is proved to lie in the construction graph.
At every step, common-neighbour inclusion--exclusion gives at least
$[1-b(1-\delta)]|A|$ candidates. Remove the used vertices and all their
conflict neighbours, losing at most $(r^2+1)(L+1)$. An induction counts the
extensions. The resulting map into the manuscript's labelled transfer
configurations is injective.
\end{proof}

\begin{lemma}[The ordinary potential transfer event]
\label{bp:transfer-physical-event}
\lean{Oberwolfach.transferRequiredEdges_card, Oberwolfach.transfer_required_edge_incidence, Oberwolfach.transferSurvivalSum_lower_tail, Oberwolfach.transferSurvivalSum_two_exposures, Oberwolfach.transferSurvivalSum_card, Oberwolfach.transfer_candidate_survival_deleted_count, Oberwolfach.raw_ordinary_transfer_potential_event}
\uses{bp:greedy-transfer-count,bp:raw-neighbourhood-counts,bp:bounded-differences,bp:raw-interval-law,bp:common-event-probability}
\leanok
For a fixed potential transfer family, give each labelled configuration
weight $n^{-r^2}$. Its required physical-edge set is the union of its
constituent cycle edge sets, of size at most $q=(1+2r)r$.
The total weight is at most one and each edge has incident weight at most
$r^2/n$. If the family has at least $cn^{r^2}$ members and all supplied edge
probabilities are at least $p\in[0,1]$, the lower-tail failure at
$p^q c-t$ is at most $2\exp(-2t^2/(r^2q/n))$.
In the original interval/physical sampling law, one event gives these
positive potential transfer counts for every vertex and every pair of
distinct hubs, under the stated minimum-degree and numerical margins.
The two root spokes are omitted. If both are present, the normalized sum
is exactly $n^{-r^2}$ times the number of actual transfer configurations,
also after conflict-edge deletion.
\source{lem:completion:ordinary-reference}
\end{lemma}
\begin{proof}\leanok
A required edge has at least one non-root endpoint, fixing a free label.
The existing function-count bound gives normalized incidence at most
$r^2/n$. Bound the squared-influence sum by maximum incidence times total
incidence, and use the actual product of the edge probabilities for the
expectation. Shared edges are tested once. The candidate count follows
from the preceding greedy construction applied to the sampled common hub
neighbourhood; all its ambient-root degree hypotheses follow from the
first common event. Integrate the conditional physical event over that
first exposure and identify the product with the original raw wheel law.
Exact inverse maps establish the final physical configuration identity.
\end{proof}


\begin{lemma}[Uniform ordinary full-host sampling]
\label{bp:ordinary-fullhost-reference}
\lean{Oberwolfach.raw_ordinary_isolator_uniform_potential_event, Oberwolfach.raw_ordinary_transfer_uniform_potential_event, Oberwolfach.raw_ordinary_isolator_uniform_probability, Oberwolfach.raw_ordinary_transfer_uniform_probability, Oberwolfach.rawWheelSpokeDegrees_uniform_probability, Oberwolfach.ordinaryFullHostSeed_uniform_probability, Oberwolfach.ordinary_full_host_reference_event, Oberwolfach.ordinary_full_host_sampling}
\uses{bp:positive-ordinary-cycle-weights,bp:ordinary-seed-means,bp:raw-isolator-event,bp:transfer-physical-event,bp:common-event-probability}
\leanok
For $r\ge3$, $\theta\ge\delta_{C_r}$, $\alpha,\omega,p_*>0$ and
$p_*\le1$, there are $k\ge r+2$ and positive constants $\beta,c,B$.
For each error tolerance $\eta>0$ there is a degree-error tolerance;
for each later conflict-degree bound, interval length and failure power,
there is a uniform order threshold. In every sufficiently large host
with minimum degree at least $(\theta+\alpha)n$, positive bounded
fractional cycle weights can be chosen before the hub set, conflict graph,
permutation and sample arrays. For all feasible supplied marginals and
the exact rim/spoke degree-balance identity, the resulting seed coefficients
have the proved two-sided bounds. One event gives spoke minimum degree
$\beta n$, all potential loads within $\eta$ of one, and potential isolator
and transfer counts at least $\beta$. Its failure probability is at most
$n^{-q}$ for any prescribed fixed $q$. Roots retain their original domains;
the physical isolator root and the two transfer root spokes are omitted.
\source{lem:completion:ordinary-reference}
\end{lemma}
\begin{proof}\leanok
Apply the positive fractional cycle-weight theorem to the parent graph.
In the configuration estimates use $p_*/2$ and $p_*^2/2$ as the lower
neighbourhood-size parameters, choosing the configuration constants before
the actual spoke density. The proved exponential bounds absorb the
polynomial numbers of root tests with thresholds uniform over all supplied
arrays. The fixed conflict-deletion errors are bounded by increasing the
order threshold, without changing any coefficient. Intersect the degree,
load, isolator and transfer events and add their failure probabilities;
no independence of these four events is assumed. The case
$\theta+\alpha\ge1$ is excluded by the strict simple-graph degree bound,
so no artificial upper restriction on $\alpha$ remains.
\end{proof}

\begin{lemma}[The labelled transfer identity and upper counts]
\label{bp:ordinary-transfer-identities}
\lean{Oberwolfach.transfer_index_vertex_load, Oberwolfach.OrdinaryTransferConfiguration.loads, Oberwolfach.ordinary_transfer_configuration_root_count}
\leanok
The signed base, row and root-replaced row cycles give vertex load one
at the specified root and zero elsewhere. Placing the same cycle weights
at two hubs with opposite signs gives zero rim load and transfers one unit
between the two root spokes. A prescribed set of $s$ distinct free labels
occurs in at most $(r^2)^s n^{r^2-s}$ labelled configurations.
\end{lemma}
\begin{proof}\leanok
The two versions of each row cancel at all its internal vertices. Their
remaining difference is the root indicator minus the base-vertex indicator.
Sum over the base cycle, retaining the coefficient $1/r$, and lift the
cycle identity to the two hubs. For the count, inject each configuration
into its function on the $r^2$ free positions, assign the prescribed labels
to positions, and count the remaining functions.
\end{proof}


\begin{lemma}[Simultaneous ordinary completion cut estimates]
\label{bp:ordinary-reference-cuts}
\lean{Oberwolfach.rawIntervalCutCount_quadratic_event, Oberwolfach.bipartiteMatrixCut_large_of_degree_and_quadratic, Oberwolfach.rawIntervalSpokeMatrix_degree_event, Oberwolfach.rawWheelSpokeCuts_uniform_probability, Oberwolfach.finiteLaw_pi_prod_probability, Oberwolfach.retention_orientation_probability, Oberwolfach.rawWheel_orientation_probability, Oberwolfach.sampledColourOrientation_large_cuts, Oberwolfach.ordinary_reference_large_cuts, Oberwolfach.ordinary_completion_large_cuts}
\leanok
For fixed positive $\alpha,p_*,\omega$ with $\alpha\le1/2$, there
are $0<\delta<1/4$ and $h>0$ before the interval length is chosen.
The ordinary spoke graph and the independently oriented admissible
physical sample simultaneously satisfy
\[
 e(X,Y)\ge2h\min\{|A_2||X|,|A_1||Y|\}
 \quad\hbox{if } |X|/|A_1|+|Y|/|A_2|>1-2\delta.
\]
The probability is at least $1-n^{-q}$ for every fixed exponent $q$,
with a uniform order threshold over hosts and supplied marginal arrays.
The general form allows any fixed conflict graph of maximum degree $L$;
the completion form discharges this input for $F_{\pi,3d}$ using
$\Delta(F_{\pi,3d})\le6d$. No independence of the two success events
is required.
\source{lem:completion:ordinary-reference-cuts}
\end{lemma}
\begin{proof}\leanok
For each fixed spoke cut the bounded coordinate supports give a
quadratic-deviation tail $2\exp(-c_d n^2)$. Union over at most $4^n$
subset pairs, and combine with minimum-degree estimates for cuts with a
small side. For physical cuts use independent colour and direction
coordinates. Exact pushforward identities identify this law with the
original retention coins and independent fair orientations, including
orientations of absent edges. Remove the conflict graph using its
explicit degree loss before applying the dense-host estimate. Intersect
the spoke and physical events with error exponent $q+1$.
This is an alternative proof of the cut conclusion; it does not use the
manuscript's relative-codegree corollary.
\end{proof}

\begin{lemma}[The first marked exposure]
\label{bp:marked-first-exposure}
\lean{Oberwolfach.rawIntervalFamily_selected_patterns_probability, Oberwolfach.rawIntervalSelectedNeighbourCount_common_event, Oberwolfach.rawIntervalSelectedLabelCount_common_event, Oberwolfach.rawMarkedHubs_uniform_probability_explicit, Oberwolfach.rawMarkedNeighbourhood_graph_event, Oberwolfach.rawMarkedNeighbourhood_dense_event, Oberwolfach.rawMarkedSpokeClasses_disjoint, Oberwolfach.rawMarked_first_event}
\uses{bp:raw-interval-law,bp:bounded-differences,bp:common-event-probability}
\leanok
Fix $m$, $0<\alpha\le1/2$, and positive lower densities $s_*,t_*,\omega$.
In a host of minimum degree at least $(1/2+\alpha)n$, the original
interval experiment has one event controlling every separated
$(m+1)$-tuple and every pair of distinct hubs. Each tuple has at least
$\omega t_*s_*^m n/2$ compatible hubs. Both spoke-intersection classes
have size at least $\min\{s_*,t_*\}^2n/2$. Every ambient vertex has at
least $(1/2+\alpha/2)|T|$ neighbours in either such class $T$ and in
$\pi(T)$. The ordinary class is disjoint from the marked class and its
successor image. The probability is at least $1-n^{-q}$, uniformly over
all feasible marginal arrays, for each fixed interval length and $q$.
\source{lem:completion:marked-reference-configurations, first exposure only}
\end{lemma}
\begin{proof}\leanok
Compute exact means using the original supported coordinate patterns.
Apply bounded differences to the fixed domains, including neighbourhoods
of every ambient root and their permutation preimages. A polynomial union
controls all tuples and hub pairs together. Convert additive errors into
the displayed relative degree bounds. The two disjointness assertions
hold for every raw outcome: an interval-covered label and its successor
are blocked from receiving an ordinary label at the same hub.
This does not yet establish the subsequent greedy configuration counts.
\end{proof}

\begin{lemma}[The four initial marked cut bounds]
\label{bp:marked-reference-cuts}
\lean{Oberwolfach.markedPhysicalCollapse_law, Oberwolfach.markedPhysical_orientation_probability, Oberwolfach.matrixCut_large_of_degree_and_quadratic_strong, Oberwolfach.permutationSeparatedMatrix_cut_loss, Oberwolfach.markedPhysical_large_cuts, Oberwolfach.rawSelectedSpokeMatrix_large_cuts, Oberwolfach.rawMarked_spoke_cuts, Oberwolfach.marked_reference_large_cuts}
\uses{bp:raw-interval-law,bp:ordinary-reference-cuts,bp:common-event-probability}
\leanok
Fix positive $\alpha,k_*,t_*,\omega$, with $\alpha\le1/2$ and
$t_*\le1$. Put
\[
 \delta_D=\min\{\alpha,k_*,t_*\}/256,\qquad
 h_D=\min\{\alpha k_*,t_*\}/128.
\]
For actual directed physical probabilities $p,k\ge k_*$ and ordinary
and marked spoke probabilities at least $t_*$, all four incidence graphs
satisfy the $(\delta_D,16h_D)$ large-cut inequality on one sample.
The host has minimum degree at least $(1/2+\alpha)n$, and short-distance
auxiliary arcs are removed at radius $3d$. The probability is at least
$1-n^{-q}$ with a threshold uniform over the actual probabilities and
feasible arrays. The physical result also retains the coefficients
$\alpha p/4$ and $\alpha k/4$ for the respective types and every
admissible cut margin.
\source{lem:completion:marked-reference, cut conclusions}
\end{lemma}
\begin{proof}\leanok
Prove equality between the five-state physical law and the projection
of a three-colour direction law, merging both absent directions.
The marked auxiliary cut is exactly a physical cut after permuting its
second side. Removing close pairs loses at most $(6d+1)\min\{|S|,|T|\}$,
without a sampling-density factor. Stronger direct quadratic cut bounds
leave the required coefficient after this deletion. For each spoke type,
combine degree estimates and a simultaneous quadratic error over every
cut. Intersect the two spoke events, then combine the independent physical
exposure, retaining the fixed values of $\delta_D$ and $h_D$.
This node records the cut conclusion. The full reference event with the
degree and intersection estimates is assembled below.
\end{proof}


\begin{lemma}[Marked base cycles in the original sample]
\label{bp:marked-base-cycles}
\lean{Oberwolfach.permutedNeighbourChoices_card_lower, Oberwolfach.mixed_greedy_sequence_count, Oberwolfach.dense_marked_base_cycle_count, Oberwolfach.denseMarkedBase_physical_injective, Oberwolfach.markedBaseSurvivalSum_lower_tail, Oberwolfach.markedBaseSurvivalSum_count, Oberwolfach.markedBase_uniform_event, Oberwolfach.rawMarked_base_cycles_event, Oberwolfach.survivingBase_card_le_common_cycle}
\uses{bp:marked-first-exposure,bp:bounded-differences,bp:common-event-probability}
\leanok
Write $c=m+1\ge3$ and $a=\alpha\min\{s_*,t_*\}^2/8$.
On one event of probability at least $1-n^{-q}$, every separated
$c$-tuple has the compatible-hub bound of the first exposure, and every
pair of distinct hubs has at least $k_*^c a^c n^c/2$ common sampled base
cycles. Their canonical tuples start at the marked head. The physical
closing edge joins the last ordinary label to the successor of that head.
The coefficients precede the interval length and actual feasible arrays.
\source{lem:completion:marked-reference-configurations, hubs and base cycles}
\end{lemma}
\begin{proof}\leanok
In the two spoke classes, use the ambient-root neighbourhood bounds to
choose a separated tuple. The last choice satisfies the two required
neighbour constraints. Prove injectivity of its physical-edge map before
assigning states. Every physical edge belongs to at most $2c n^{c-1}$
candidate tuples. After normalization this gives variance bound
$2c^2/n$ and a lower tail exponential in $n$. This deliberately uses a
weaker intermediate bound than the manuscript's quadratic exponent;
it suffices for a polynomial union over hub pairs. The pattern sum is
proved equal to the actual sampled-arc count. Integrate the common
physical event over the good first exposure and embed the counted tuples
into the canonical common-cycle family. The common-path argument is recorded
next, followed by its full exponential interface.
\end{proof}

\begin{lemma}[The sampled physical map is injective]
\label{bp:marked-sample-physical-map}
\lean{Oberwolfach.markedPhysicalRim_eq_one_iff, Oberwolfach.markedPhysicalRim_physical_unique, Oberwolfach.markedSampleArc_physical_injective}
\uses{bp:marked-reference-cuts}
\leanok
For every outcome of the five-state physical experiment, the physical
map on the actual ordinary and marked arc types is injective.
This is not asserted for all potential arcs before sampling.
\end{lemma}
\begin{proof}\leanok
Equal unordered physical edges have the same sampled type and orientation.
Reversing the endpoints changes the orientation bit, so their ordered
physical endpoints agree. The successor permutation is injective, giving
equality of the auxiliary heads as well.
\end{proof}


\begin{lemma}[Common marked paths without root spokes]
\label{bp:marked-common-paths}
\lean{Oberwolfach.marked_path_choice_count, Oberwolfach.markedPathChoice_to_parent, Oberwolfach.denseRootMarkedPath_uniform_density, Oberwolfach.pathPhysicalEdge_injective, Oberwolfach.markedPathSurvivalSum_lower_tail, Oberwolfach.markedPathSurvivalSum_count, Oberwolfach.markedPath_uniform_event, Oberwolfach.rawMarked_paths_event, Oberwolfach.survivingPath_card_le_common_path}
\uses{bp:marked-first-exposure,bp:marked-base-cycles,bp:marked-sample-physical-map,bp:bounded-differences}
\leanok
For cycle length $c=r+5$ with $r\ge1$, put
$a=\alpha\min\{s_*,t_*\}^2/8$.
One raw-sampling event of probability at least $1-n^{-q}$ controls all
distinct root pairs, distinct hub pairs and both root types. Every such
choice has at least $k_*^{c+2}a^{c-1}n^{c-1}/2$ common sampled paths.
The four root spokes are omitted. The roots need not be separated from
one another under the permutation; all free labels are separated from
both roots and from one another.
\source{lem:completion:marked-reference-configurations, common paths}
\end{lemma}
\begin{proof}\leanok
Choose the two path endpoints first, so that each of the remaining choices
has at most two neighbour constraints. Reorder these choices into the
actual common path and verify the ordinary and marked physical endpoints
separately. The $c+2$ required physical edges are distinct even when the
roots are close: the internal edges use free labels, while the four root
edges use distinct roots and endpoints. Apply the original physical law,
then a polynomial union over all fixed roots, hubs and types. An explicit
injection takes every surviving tuple to the root-first common-path family
used by the correction argument. No root-spoke presence is assumed.
\end{proof}

\begin{lemma}[One event for all marked configurations]
\label{bp:marked-configurations}
\lean{Oberwolfach.rawMarked_configurations_event, Oberwolfach.rawMarkedHubs_exponential_probability, Oberwolfach.rawMarkedNeighbourhood_exponential_dense_event, Oberwolfach.markedBase_exponential_event, Oberwolfach.markedPath_exponential_event, Oberwolfach.rawMarked_first_exponential_event, Oberwolfach.rawMarked_base_cycles_exponential_event, Oberwolfach.rawMarked_paths_exponential_event, Oberwolfach.markedConfigurationExponent_pos, Oberwolfach.rawMarked_configurations_exponential_event, Oberwolfach.MarkedReferenceConfigurations, Oberwolfach.marked_configurations_raw_event}
\uses{bp:marked-first-exposure,bp:marked-base-cycles,bp:marked-common-paths,bp:common-event-probability}
\leanok
For each fixed $c=r+5\ge6$ and positive input density bounds there is
$\beta>0$, chosen before the interval length and actual marginal arrays,
such that one original raw sample has all the following properties with
failure probability at most $34n^{c+3}e^{-a_dn}$, where $a_d>0$: every separated $c$-tuple has at least
$32\beta n$ compatible hubs; every distinct hub pair has at least
$16\beta n^c$ common cycles; every distinct root pair, distinct hub pair
and either root type has at least $16\beta n^{c-1}$ common paths with
root spokes omitted. The counted families are the actual sampled-arc
families, not abstract independent indicators.
\source{lem:completion:marked-reference-configurations}
\end{lemma}
\begin{proof}\leanok
Intersect the base and path events on the same interval--physical product
space. Choose $\beta$ below all three previously established coefficients.
Keep the explicit exponential bounds in both exposures. The first exposure
has failure at most $14n^{c+3}e^{-a n}$; the physical cycle and path
estimates contribute at most $2n^2e^{-b n}$ and $4n^4e^{-b'n}$.
Combine the cycle and path events by a union bound. An exact product-law
identity expresses the result on the dependent raw-coordinate space used
for the reference degrees and loads. The inverse-power interface is retained
as a weaker earlier consequence.
\end{proof}


\begin{lemma}[Exact marked seed means]
\label{bp:marked-seed-means}
\lean{Oberwolfach.typedMarkedResource_product_omitted, Oberwolfach.rawMarkedPotentialLoad_expectation, Oberwolfach.rawMarkedSeed_expectation, Oberwolfach.markedParent_spoke_position_load, Oberwolfach.markedParent_marked_rim_position_load, Oberwolfach.markedParent_ordinary_rim_position_load, Oberwolfach.rawMarkedSeed_spoke_mean, Oberwolfach.rawMarkedSeed_rim_mean}
\uses{bp:successor-complete,bp:marked-wheel-resources,bp:typed-marked-potential-stability,bp:raw-interval-law}
\leanok
For fixed successor tuple weights with positional error $\tau$, use
$w(C,u)=\lambda(C)/(Dp^{c-1}ks^{c-1})$ and omit precisely the tested
root indicator. On the original raw law the mean of each eligible
potential rim load and each potential spoke load lies between $1-\tau$
and $1$. A spoke root may use any hub in a fixed external domain.
The denominator $D$ is unchanged by that extension.
\source{lem:completion:marked-reference-loads, expectation calculation}
\end{lemma}
\begin{proof}\leanok
Within each parent wheel, verify disjoint primitive supports for its
physical rims and interval spokes. The product over the remaining
resources equals the full probability product divided by the omitted
root marginal. A rim root sums over the inner hub set; a spoke root fixes
its hub. Convert canonical wheel sums exactly to positional tuple sums.
The successor theorem bounds each of these sums, with all $c-1$ ordinary
positions counted. No conditioning on root presence and no change of
parent coefficients occurs.
\end{proof}

\begin{lemma}[Marked seed loads on nested hub domains]
\label{bp:marked-seed-loads}
\lean{Oberwolfach.rawMarkedResourceCoordinates_card, Oberwolfach.rawMarkedResourceCoordinates_reverse_card, Oberwolfach.typedMarkedPotentialLoad_supported_common_event, Oberwolfach.typedMarkedSpoke_supported_common_event, Oberwolfach.successor_seed_nonneg_upper, Oberwolfach.successor_seed_positive_lower, Oberwolfach.rawMarkedSeed_load_event, Oberwolfach.rawMarkedSeed_external_spoke_event, Oberwolfach.rawMarkedSeed_nested_hubs_event, Oberwolfach.marked_reference_loads}
\uses{bp:marked-seed-means,bp:resource-pattern-locality,bp:bounded-differences,bp:common-event-probability,bp:raw-marked-restriction}
\leanok
Put $m=c-1$, $p=\varrho s/2$, $k=p/m$, $t=s/m$, and
$D=2|B|/\varrho\in[n/2,2n]$. Fix $B\subseteq Y$ with
$|B|\le n$ and $|Y|\le K_{\rm out}n$, and fix the successor weighting
before exposure. For $p_* =\varrho_*s_* /(2m)$, put
$B_+=2C_\lambda/p_*^{2m+1}$ and
\[
 c_* = \frac{2\eta^2}{(8d+6)^2\,2(m+1)^7B_+^2}>0.
\]
One original raw sample has all potential loads on $(V,B)$ and all
spoke loads on $(V,Y)$ within $\tau+\eta$ of $1$, except with probability
at most $(8+4K_{\rm out})n^2e^{-c_*n}$. The normalized coefficients
are between $(a_\lambda/2)n^{-m}$ and $B_+n^{-m}$ on their positive
parent support. These are fixed-domain assertions; there is no union over
outcome-dependent weightings or domains.
\source{lem:completion:marked-reference-loads}
\end{lemma}
\begin{proof}\leanok
A primitive coordinate affects at most $8d+6$ typed resources.
Root-omitted pair loads are bounded by $(m+1)^4B_+/n$.
For external spokes, the tested root fixes the hub, so that bound needs
no upper bound on the number of external hubs. Bounded differences and
a union over all fixed roots give the displayed exponential estimate.
Project the original hub coordinates injectively onto $B$ and intersect
its load event with the external-spoke event on the same ambient law.
The actual density identities give the normalization factors and the
uniform two-sided coefficient bounds.
\end{proof}

\begin{lemma}[Exact restrictions of raw marked samples]
\label{bp:raw-marked-restriction}
\lean{Oberwolfach.rawMarked_hub_restriction_probability, Oberwolfach.rawMarked_vertex_restriction_probability, Oberwolfach.invariantSetPermutation, Oberwolfach.invariantSetPermutation_noshort, Oberwolfach.permutationWindow_image, Oberwolfach.rawIntervalState_restrict, Oberwolfach.rawMarkedRestrictVertices_rims, Oberwolfach.rawMarkedRestrictVertices_spokes}
\uses{bp:raw-interval-law,bp:common-event-probability}
\leanok
Restricting fixed hub and vertex coordinate sets preserves their exact
product laws. If the vertex embedding commutes with the successor
permutations, its interval states agree with the original interval states.
An order-preserving embedding also preserves the ordinary and marked
physical rim states. The induced order on a subset has this property.
A genuinely permutation-invariant subset carries the restricted
permutation, with the original lower bound on cycle lengths.
\end{lemma}
\begin{proof}\leanok
Use injective projection of the independent primitive coordinates.
Every permutation window on the embedded domain is exactly the image
of the restricted window. This identifies competing starts, covered
vertices, blocked vertices and both spoke types. The physical orientation
bit is preserved by the order embedding. Merely restricting to an arbitrary
non-invariant subset does not justify these state identities.
\end{proof}


\begin{lemma}[Marked reference degrees and intersections]
\label{bp:marked-reference-statistics}
\lean{Oberwolfach.permutationSeparatedMatrix_row_loss, Oberwolfach.permutationSeparatedMatrix_column_loss, Oberwolfach.markedPhysicalRim_degree_event, Oberwolfach.sampledColourOrientation_twoPath_tail, Oberwolfach.sampledColourOrientation_twoPaths_common_event, Oberwolfach.permutationSeparatedMatrix_two_step_loss, Oberwolfach.markedPhysicalRim_reference_paths, Oberwolfach.markedPhysicalRim_uniform_capacity_event, Oberwolfach.rawMarkedNeighbourhood_exponential_event, Oberwolfach.rawMarkedCodegree_normalized_event, Oberwolfach.rawMarked_spoke_reference_event, Oberwolfach.MarkedReferenceStatistics, Oberwolfach.markedReferenceExponent_pos, Oberwolfach.marked_reference_statistics}
\uses{bp:marked-reference-cuts,bp:raw-neighbourhood-counts,bp:raw-interval-law,bp:common-event-probability,bp:categorical-concentration}
\leanok
For fixed positive errors $\varepsilon,\delta$, minimum-degree margin
$\gamma$, lower densities $p_*,s_*$ and hub ratio $\kappa$, the order
threshold precedes the graph and actual densities. On the original raw
marked law, one event controls both rim in/outdegrees around $pD,kD$,
all ordinary directed two-step paths between distinct roots, intersections
of one, two or three distinct hub neighbourhoods, and the normalized
spoke codegrees outside the $3d$ permutation neighbourhoods.
Its failure probability is at most $36n^3e^{-cn}$, where
\[
 c=\min\left\{\frac{\varepsilon^2}{2},
 \frac{\gamma^2p_*^4}{16},
 \frac{2\delta^2}{3(4d+3)^2},
 \frac{\delta^2s_*^4\kappa}{(4d+3)^2}\right\}>0.
\]
The parent degrees must be within $\delta_Pn$ of $D$, with
$0\le\delta_P\le\varepsilon/4$, and $\kappa n\le |U|\le n$.
The spoke marginals are fixed and each at least $s_*$; the ordinary
rim marginal is at least $p_*$. This node does not include the cut or
later retention conclusions.
\source{lem:completion:marked-reference, degrees and intersections}
\end{lemma}
\begin{proof}\leanok
Incoming marked arcs at $x$ use parent degree $d_P(\pi(x))$.
Deleting permutation-close arcs costs at most $6d+1$ in each degree
and at most $2(6d+1)$ intermediate vertices in a realised two-step count.
No sampling-density factor is attached to this loss. For distinct roots,
different two-step paths use disjoint physical edges. The interval
neighbourhood and codegree estimates retain their original local supports.
Normalize the codegrees by the fixed squared spoke marginal and the
original hub-domain size. Exact product-law splitting combines all these
estimates on the same raw outcome.
\end{proof}

\begin{lemma}[One event at all fixed vortex levels]
\label{bp:vortex-common-events}
\lean{Oberwolfach.vortex_inverse_sum_bound, Oberwolfach.vortex_fixed_domain_common_event, Oberwolfach.vortex_exponential_common_event}
\uses{bp:vortex,bp:common-event-probability}
\leanok
Fix $0<\varepsilon<1$, $c,\chi>0$, $C\in\mathbb R$ and
$K\in\mathbb N$. There is a terminal-size threshold $M$ such that,
for every fixed finite sequence with $n_{i+1}\le\varepsilon n_i$
and $n_i\ge M$, events on one probability space with failure bounds
$Cn_i^Ke^{-cn_i}$ hold simultaneously with probability at least
$1-\chi$. The threshold is independent of the initial order and
number of levels. Independence between levels is not assumed.
\end{lemma}
\begin{proof}\leanok
Choose $M$ so the individual bounds are at most $1/n_i$.
The geometric size decrease implies
$\sum_i1/n_i\le((1-\varepsilon)M)^{-1}$. Apply the finite union
bound on the common sample space and enlarge $M$ to make this at most
$\chi$. This proves a fixed-sequence event, not uniformity over arbitrary
outcome-dependent domain choices.
\end{proof}


\begin{lemma}[The complete marked reference event]
\label{bp:marked-reference-common}
\lean{Oberwolfach.marked_reference_cut_event, Oberwolfach.markedReferenceCutFamilies_fixed_cuts, Oberwolfach.markedReferenceCutFamilies_ordinary_cut, Oberwolfach.marked_reference_event, Oberwolfach.marked_reference_actual_densities, Oberwolfach.markedReferenceCommonExponent_pos, Oberwolfach.marked_reference_common_event}
\uses{bp:marked-reference-statistics,bp:marked-reference-cuts,bp:marked-configurations,bp:common-event-probability}
\leanok
On one original marked sample, the degree, intersection and cut estimates
hold with failure at most $72n^3e^{-an}$. The fixed cut parameters are
$\delta_D=\min\{\gamma,k_*,t_*\}/256$ and
$h_D=\min\{\gamma k_*,t_*\}/128$, with coefficient $16h_D$.
The ordinary rim also has cut coefficient $\gamma p/4$ at cutoff
$\min\{\gamma/64,p/128\}$. The actual density substitution
$p=\varrho s/2$, $k=p/(c-1)$, $t=s/(c-1)$ and $D=2|U|/\varrho$
is checked, with the order threshold preceding the actual densities.
Intersecting with the configuration event gives failure at most
$106n^{c+3}e^{-a'n}$ for a fixed positive $a'$.
The configuration coefficient precedes the errors and the interval length.
These assertions concern the raw sample, before retention.
\source{lem:completion:marked-reference, lem:completion:marked-reference-configurations}
\end{lemma}
\begin{proof}\leanok
Retain exponential bounds throughout the cut proofs, including the union
across all cuts. Convert the stronger cut families to the source constants
by explicit numerical inequalities. Use a union bound on the same raw
coordinate law for statistics, cuts and configurations. The exponent is
the minimum of their positive fixed exponents. No independence of these
test events is assumed, and none of the marginal arrays is recomputed.
\end{proof}

\begin{lemma}[Marked reference estimates on fixed vortex domains]
\label{bp:marked-reference-vortex}
\lean{Oberwolfach.rawMarked_pair_restriction_probability, Oberwolfach.rawMarkedRestrictPair_rims, Oberwolfach.rawMarkedRestrictPair_spokes, Oberwolfach.marked_reference_vortex_event}
\uses{bp:marked-reference-common,bp:raw-marked-restriction,bp:vortex-common-events}
\leanok
For a fixed sequence of vertex and hub domains with geometrically decreasing
vertex orders, the reference degrees, intersections, cuts and configurations
hold simultaneously on one ambient sample with probability at least $1-\chi$.
Each induced host must satisfy the stated near-regularity, minimum-degree
and hub-size hypotheses. Each vertex embedding commutes with the supplied
permutations. The terminal-order threshold is independent of the initial
order and the number of domains. It precedes the hosts, domain embeddings
and actual marginal arrays. This does not include external-root local-factor
estimates or the retention step.
\end{lemma}
\begin{proof}\leanok
Restrict hubs and vertices by exact coordinate projections. The induced
product law is unchanged; commutation with the permutations identifies
all interval states and physical orientations. Apply the uniform reference
event on each fixed domain and sum its exponential failure bounds along
the size sequence. Overlaps between domains require no independence.
\end{proof}


\begin{lemma}[Conditional retention and state changes]
\label{bp:retention-conditional}
\lean{Oberwolfach.startRetentionLaw, Oberwolfach.startRemovedHubs_joint_probability_le, Oberwolfach.startRemovedHubs_row_tail, Oberwolfach.retainedRawFamily_row_changes, Oberwolfach.retainedRawFamily_index_changes, Oberwolfach.retainedRawFamily_row_tail, Oberwolfach.retainedRawFamily_index_tail, Oberwolfach.startRetentionMinimum_polynomial_rate}
\uses{bp:raw-interval-law,bp:common-event-probability,bp:fixed-size-concentration}
\leanok
For a fixed raw outcome, put $A_a=\{u:a\text{ is a retained raw start of the selected colour at }u\}$
and $t=\min_a|A_a|$. Independently over $a$, choose a uniform $t$-subset
$J_a\subseteq A_a$. For $W\subseteq A_a$,
\[
 \Pr(W\cap J_a=\varnothing)\le
 \left(\frac{|A_a|-t}{|A_a|}\right)^{|W|}.
\]
A deleted start changes at most $d+1$ states, in its interval and successor.
For a permutation-invariant set $S$, only starts in $S$ can change states in
$S$. At a fixed vertex, only the $d+1$ possible preceding starts can change
its state. If every raw start count differs from its common mean
$\mu\ge pn/d$ by at most $n^{3/5}$, the deletion probability is at most
$(4d/p)n^{-2/5}$ for sufficiently large $n$.
\end{lemma}
\begin{proof}\leanok
Keep the raw activation and ordinary-colour variables fixed. Independence
across starts gives the row tail. Within a start, the exact uniform-subset
joint exclusion bound gives the index tail. No independence of hub deletions
within that subset is assumed. The minimum-count identity bounds every
removed set by $2n^{3/5}$.
\end{proof}

\begin{lemma}[Uniform retention on all fixed restrictions]
\label{bp:interval-retention}
\lean{Oberwolfach.rawStartHubSet_common_event, Oberwolfach.uniform_binomial_restriction_tail, Oberwolfach.rawUniformRetentionLaw, Oberwolfach.raw_uniform_retention_restrictions, Oberwolfach.raw_uniform_retention_restrictions_real, Oberwolfach.rawUniformRetentionLaw_projection, Oberwolfach.rawUniformRetentionLaw_common_event}
\uses{bp:retention-conditional,bp:raw-interval-law,bp:common-event-probability,bp:categorical-concentration}
\leanok
For fixed $d\ge2$, $p,\zeta>0$ and real exponents $C,Q$, there are
$m,N$ such that the following holds. In the original raw interval law on
$|V|=n\ge N$, $|U|\le n$, the selected start-colour marginal is constant
$t_0$ and $t_0|U|\ge pn$. Retain only that colour as above.
Fix at most $n^C$ pairs $(S,T)$ with $S$ permutation-invariant and
$|S|,|T|\ge m$. With probability at least $1-n^{-Q}$, all raw start counts
are within $n^{3/5}$ of $t_0|U|/d$, and every row changes at most
$\zeta|S|$ states on $S$, while every index count changes at most
$\zeta|T|$ states on $T$, simultaneously for all these pairs.
The thresholds precede the host, domains, actual marginal arrays and outcome.
The retained law projects exactly to the original raw law, so any other
common raw event can be intersected with this event by an additive union
bound. Domains must be fixed before exposure.
\source{lem:completion:trimming-restrictions}
\end{lemma}
\begin{proof}\leanok
Use Hoeffding over independent hub columns for the raw counts. Conditional
on that event, the row and index tails are bounded by
$\binom{M}{k}p_n^k$, with the additional factor $d+1$ for index counts.
For $k=\lceil\zeta M/(d+1)\rceil$, these bounds decay polynomially in $n$
with an exponent proportional to $M$. Choose $m$ before $n$, take one union
bound over the supplied family, and integrate the conditional law. The proof
does not need comparability between $|S|$ and $|T|$; it proves the stated
stronger restriction interface directly.
\end{proof}

\begin{lemma}[Exact retained interval degrees]
\label{bp:retained-interval-identities}
\lean{Oberwolfach.retainedIntervalState_interval_iff, Oberwolfach.retainedRawStart_selected_iff, Oberwolfach.retainedRawStart_card, Oberwolfach.retainedRawFamily_covered_hubs, Oberwolfach.retainedRawFamily_covered_card}
\uses{bp:retention-conditional,bp:raw-interval-law}
\leanok
Every start of the trimmed colour occurs at exactly $t$ hubs. Every vertex
is covered by that retained colour at exactly $dt$ hubs, provided the
permutation has no cycle of length less than $d$. These equalities hold
for every retention outcome, not only in expectation.
\source{eq:interval:retained-identities}
\end{lemma}
\begin{proof}\leanok
The possible starts covering a vertex form a reverse permutation interval
of size $d$. Their retained hub sets each have size $t$. They are pairwise
disjoint: a hub in two such sets would give two conflicting raw kept starts.
Thus their union has size $dt$ and is exactly the set of hubs whose retained
state at that vertex is the selected interval colour.
\end{proof}

\begin{lemma}[Local spoke classes with ambient roots]
\label{bp:marked-local-classes}
\lean{Oberwolfach.rawMarkedLocal_domain_event, Oberwolfach.rawMarkedLocal_graph_event, Oberwolfach.rawMarkedLocal_classes_disjoint, Oberwolfach.rawMarkedLocal_dense_event}
\uses{bp:raw-interval-law,bp:categorical-concentration}
\leanok
Let $A$ be a fixed invariant inner set, and suppose every ambient vertex
has at least $(1/2+\alpha)|A|$ neighbours in $A$. Assume the ambient
vertex and hub orders are at most $K|A|$, and both selected spoke
probabilities are at least $t_*>0$. On one event for every hub $u$, the
classes $T_0(u),T_K(u)\subseteq A$ each have size at least $t_*|A|/2$,
and every ambient root has at least $(1/2+\alpha/2)|T_j(u)|$ neighbours
in $T_j(u)$ and in $\pi(T_j(u))$. The classes are disjoint, and
$T_0(u)\cap\pi(T_K(u))=\varnothing$. The failure is at most
$12K^2|A|^2\exp(-c|A|)$, with explicit $c>0$.
\end{lemma}
\begin{proof}\leanok
Apply the original raw-law count theorem on $A$, $N_P(x)\cap A$ and
$A\cap\pi^{-1}(N_P(x))$, for every ambient $x$. Invariance of $A$
identifies the last parent count. The parent hypotheses apply to external
roots explicitly; they are not inferred from $P[A]$. Raw state exclusion
gives both disjointness statements deterministically.
\end{proof}

\begin{lemma}[Local greedy counts and physical survival]
\label{bp:marked-local-physical}
\lean{Oberwolfach.denseMarkedBase_local_density, Oberwolfach.denseRootMarkedPath_local_density, Oberwolfach.localScale_exponential_tail, Oberwolfach.markedBase_local_exponential_event, Oberwolfach.markedPath_local_exponential_event, Oberwolfach.markedLocal_physical_event}
\uses{bp:marked-local-classes,bp:marked-base-cycles,bp:marked-common-paths}
\leanok
The marked base-cycle and common-path lower counts hold at an arbitrary
supplied inner order $a$, provided the two classes have linear size in $a$
and the required neighbourhood bounds hold at every ambient root. Their
physical survival counts hold simultaneously for every hub and every pair
of distinct ambient roots in both spoke types. No root spokes are required.
The leading coefficients depend on the degree and marginal lower bounds
and cycle length, but not on the ambient-to-inner order ratio $K$.
The failure bound is $(2K+4K^3)a^3\exp(-c a)$, where $c>0$ may depend on $K$.
\end{lemma}
\begin{proof}\leanok
Reuse the actual greedy choices at scale $a$. Convert the physical-survival
sums normalized at ambient order $n$ by the exact factor $(a/n)^k$.
The inequality $a\le n\le Ka$ changes only the exponential constant.
All requirements use compatible states of actual physical edges, with the
original five-state law. A union bound uses one physical outcome for all tests.
\end{proof}

\begin{lemma}[Local configurations in the original marked sample]
\label{bp:marked-local-configurations}
\lean{Oberwolfach.RawMarkedLocalConfigurations, Oberwolfach.markedLocalConfigurationExponent_pos, Oberwolfach.rawMarked_local_configurations_event}
\uses{bp:marked-local-classes,bp:marked-local-physical,bp:common-event-probability}
\leanok
For each fixed invariant inner set $A$, the local cycle and root-omitted
path counts hold in the original raw marked sample with failure at most
$(12K^2+2K+4K^3)|A|^3\exp(-c|A|)$. All hubs and all distinct ambient
vertex roots are included on this event. The order threshold precedes
the host, domain and actual feasible marginal arrays.
\end{lemma}
\begin{proof}\leanok
Condition on the one-hub spoke-class event, apply the simultaneous physical
survival theorem, and integrate over the original first exposure. Exact
product-law regrouping identifies the result with the raw coordinate law;
no exposure is repeated and no independent hub columns are identified.
\end{proof}

\begin{lemma}[Local configurations on every fixed vortex level]
\label{bp:marked-local-vortex}
\lean{Oberwolfach.marked_local_configurations_vortex_event}
\uses{bp:marked-local-configurations,bp:raw-marked-restriction,bp:vortex-common-events}
\leanok
For a fixed geometrically decreasing sequence of inner sets, each contained
in a supplied outer domain of comparable order, all local configuration
counts hold on one ambient sample with probability at least $1-\chi$.
All roots in each outer domain are allowed. Each inner set is invariant
under its supplied permutation; the embeddings of outer domains commute
with the ambient permutation. The terminal-order threshold is independent
of the initial order and level count. This statement concerns fixed domains,
not arbitrary domains chosen from the sample.
\end{lemma}
\begin{proof}\leanok
Use exact raw-law restriction to each outer domain. Apply the local event
there and sum the exponential bounds along the inner-order sequence.
No independence between overlapping levels is used.
\end{proof}

\begin{lemma}[Restricted marked degrees and balance]
\label{bp:marked-restricted-balance}
\lean{Oberwolfach.markedRestrictedDegree_tail, Oberwolfach.markedRestrictedDegree_common_event, Oberwolfach.rawIntervalSelectedLabelCountOn_common_event, Oberwolfach.rawMarkedRestrictedSpoke_common_event, Oberwolfach.rawMarked_restricted_balance_event, Oberwolfach.marked_restricted_balance_uniform_event}
\uses{bp:raw-interval-law,bp:raw-marked-restriction,bp:categorical-concentration}
\leanok
For fixed nonempty vertex and hub domains $A,B$, with $A$ invariant,
restricted rim and spoke degrees concentrate simultaneously at every
ambient root. The incoming marked parent degree is $d_P(\pi(v),A)$.
If $|d_P(v,A)-D|\le\delta M$ for all ambient $v$, the identities
$pD=s_0|B|$, $kD=t_0|B|$, $s_0=mt_0$ give all restricted balance
errors at most $\varepsilon M$, on an event of failure at most
$16KM\exp(-cM)$. Here $|A|,|B|\le M$, the ambient orders are at most
$KM$, and $\delta\le\varepsilon/4$. The positive constant and threshold
precede the host, domains and feasible marginal arrays.
\end{lemma}
\begin{proof}\leanok
Each fixed-direction rim count uses injective physical coordinates.
Deleting labels at permutation distance at most $3d$ changes the parent
centre by at most $6d+1$. Exact restriction to the hub subtype gives
spoke row counts; column counts use the original interval law. Union
bounds and the displayed mean identities give the balance inequalities.
The fixed deletion term is absorbed by the uniform order threshold.
\end{proof}

\begin{lemma}[Restricted balance on fixed vortex domains]
\label{bp:marked-restricted-balance-vortex}
\lean{Oberwolfach.rawMarkedRestrictedBalance_empty, Oberwolfach.marked_restricted_balance_vortex_event}
\uses{bp:marked-restricted-balance,bp:raw-marked-restriction,bp:vortex-common-events}
\leanok
All restricted balance inequalities hold on one ambient sample at every
fixed geometrically decreasing local scale, with failure at most $\chi$.
Every root in the supplied outer domains is included. The threshold is
independent of the initial order and number of levels. The empty terminal
pair has identically zero degree differences.
\end{lemma}
\begin{proof}\leanok
Restrict the original law exactly to each outer domain and sum its
exponential failure bounds. The parent degree assumptions are imposed
at all outer roots into the inner set; no external-root estimate is
inferred from an induced subgraph.
\end{proof}

\begin{lemma}[Actual mixed marked row patterns]
\label{bp:marked-row-patterns}
\lean{Oberwolfach.rawMarkedResourceCoordinates_disjoint, Oberwolfach.rawMarkedPatternSum_expectation, Oberwolfach.rawMarkedPatternSum_tail, Oberwolfach.markedRowResources_disjoint, Oberwolfach.markedRowResources_load, Oberwolfach.rawMarkedRowCount_patterns, Oberwolfach.markedRowResources_target, Oberwolfach.rawMarkedRowCount_tail}
\uses{bp:raw-interval-law,bp:raw-interval-geometry,bp:marked-sample-physical-map,bp:categorical-concentration}
\leanok
The actual mixed row count has the product target of the manuscript:
ordinary and marked directional marginals for the sampled arcs, the
supplied spoke marginals for the requested hubs, and the cardinality
of the fixed admissible row domain. Required physical edges must be
pairwise distinct. Fixed parent-neighbour conditions add no probability
factor. If the original outer set is $A$ and the number of sampled rim
and spoke requests is at most $h$, the deviation tail at $t>0$ is at most
\[
 2\exp\left(-\frac{2t^2}{(8d+6)^2\,2(h+1)|A|}\right).
\]
The admissible domain may be empty and the roots may lie outside $A$.
\end{lemma}
\begin{proof}\leanok
Write each summand as a pattern of actual typed resources. Their primitive
supports are disjoint within one admissible candidate: physical coordinates
are distinct, rim and spoke coordinates are disjoint, and distinct hubs
use disjoint interval coordinates. A resource occurs in at most two row
summands. The proved squared-influence bound gives the tail at the original
outer scale, without dividing by the admissible-domain size.
\end{proof}

\begin{lemma}[All bounded row patterns, including fixed vortex levels]
\label{bp:marked-rows-vortex}
\lean{Oberwolfach.boundedFinset_card, Oberwolfach.rawMarkedRows_common_event, Oberwolfach.rawMarkedRows_exponential_event, Oberwolfach.marked_rows_vortex_event}
\uses{bp:marked-row-patterns,bp:raw-marked-restriction,bp:vortex-common-events}
\leanok
One event controls every row pattern with at most $h$ sampled rim
requests, $h$ fixed parent requests, and $h$ spoke requests. On a fixed
outer host of orders at most $K|A|$, the failure is at most
$2(4K+1)^{2h}(2K+1)^h|A|^{3h}\exp(-c|A|)$, for explicit $c>0$.
The analogous event holds simultaneously at all fixed geometric vortex
levels, with every outer root and a terminal threshold independent of
initial size and level count. All supplied marginal arrays are preserved.
\end{lemma}
\begin{proof}\leanok
Bounded subsets are encoded by tuples padded with absent entries. This
proves the polynomial number of all tests before exposure. Apply the row
tail, exact original-law restriction, and the geometric exponential sum.
Later choices from these fixed root domains use this same event.
\end{proof}

\begin{lemma}[All marked index patterns on fixed domains]
\label{bp:marked-indices-vortex}
\lean{Oberwolfach.markedIndexResources_disjoint, Oberwolfach.markedIndexResources_load, Oberwolfach.rawMarkedIndexCount_patterns, Oberwolfach.markedIndexResources_target, Oberwolfach.rawMarkedIndexCount_tail, Oberwolfach.rawMarkedIndices_common_event, Oberwolfach.rawMarkedIndices_exponential_event, Oberwolfach.marked_indices_vortex_event}
\uses{bp:marked-row-patterns,bp:marked-rows-vortex,bp:raw-marked-restriction,bp:vortex-common-events}
\leanok
For a fixed hub set $B$, all separated typed vertex-label sets of size at
most $h$ have their actual index counts within $\varepsilon|B|$ of the
sum over $u\in B$ of the corresponding supplied marginal products.
Empty requests are included. If the ambient vertex order is at most
$K|B|$, the common failure is at most
$2(2K+1)^h|B|^h\exp(-c|B|)$, with explicit $c>0$.
The same conclusions hold on every fixed geometric vortex hub domain
in one ambient sample, including all labels in each outer vertex domain.
\end{lemma}
\begin{proof}\leanok
Each typed spoke resource occurs in at most one hub summand. Separated
labels have disjoint interval supports within a summand. Enumerate every
bounded label set before exposure, then use exact restriction and the
geometric exponential sum. No uniformity over outcome-dependent domains
is asserted.
\end{proof}

\begin{lemma}[Constructed local weights and omitted-root loads]
\label{bp:marked-local-load-vortex}
\lean{Oberwolfach.successor_sequence_weight_families, Oberwolfach.rawMarkedLocalSpokeLoads_exponential_event, Oberwolfach.marked_local_spoke_loads_vortex_event, Oberwolfach.marked_local_spoke_realization}
\uses{bp:marked-seed-loads,bp:raw-marked-restriction,bp:vortex-common-events}
\leanok
For fixed vortex vertex domains whose induced parent graphs satisfy the
proved successor-weight degree hypotheses, choose a separate parent tuple
weighting on each domain. These weightings are chosen before the feasible
sampling marginals. Their positive coefficient bounds precede the positional
error, and the near-regularity tolerance precedes the separation parameter.
After normalization by the actual densities, the coefficients retain the
proved two-sided bounds. One ambient sample gives every local spoke load
within $\tau+\varepsilon$ of one, at all inner vertex roots, both spoke
types, and all hubs in the supplied larger hub domains. Each tested root
spoke is omitted regardless of its presence in the sample.
\end{lemma}
\begin{proof}\leanok
Construct the domain-specific successor weights by the proved existence
theorem. Normalize using the supplied degree centre and actual marginals.
Use the original-law potential expectation and concentration, restrict
exactly to each fixed domain, and sum the exponential failures. The terminal
threshold is uniform in the initial order and level count. Weights on
different domains are not identified or recomputed after sampling.
\end{proof}

\begin{lemma}[Constructed packing weights and nested hub loads]
\label{bp:marked-packing-load-vortex}
\lean{Oberwolfach.marked_packing_loads_vortex_event, Oberwolfach.marked_packing_load_realization}
\uses{bp:marked-local-load-vortex,bp:marked-seed-loads,bp:raw-marked-restriction,bp:vortex-common-events}
\leanok
The same construction gives one event for every eligible inner rim and
spoke root in the packing domains, together with all potential spoke roots
in the larger hub domains. The exact inner-hub mean identities are required
and used. The seed normalization remains that of the inner packing domain,
including for outer hub roots. Each fixed domain has its own coefficient
function, with two-sided bounds independent of the later sampling outcome.
\end{lemma}
\begin{proof}\leanok
Combine the already proved inner-root and external-spoke event before
summing across the fixed vortex domains. The embedding of inner hubs into
the larger hub set is explicit. The original product law restricts exactly,
and the potential load omits only its tested root indicator.
\end{proof}

\begin{lemma}[The global marked retention law]
\label{bp:marked-retention-law}
\lean{Oberwolfach.finiteLaw_sigma_projection_probability, Oberwolfach.finiteLaw_pi_sigma_injective_projection, Oberwolfach.rawMarkedRetentionLaw_interval_projection, Oberwolfach.rawMarkedRetentionLaw_raw_projection, Oberwolfach.rawMarkedRetentionLaw_common_event, Oberwolfach.raw_marked_retention_common_event}
\uses{bp:interval-retention,bp:marked-first-exposure,bp:common-event-probability}
\leanok
First expose all raw physical and interval coordinates, then retain the
same number of starts of one selected interval colour at every vertex,
using the conditional uniform subsets. The projection onto raw coordinates
has exactly the original marked law. The projection onto interval and
retention coordinates has exactly the proved two-stage interval law.
Consequently an arbitrary common raw event of failure at most $a$ and
the fixed-family retention event of failure at most $n^{-Q}$ hold together
with failure at most $a+n^{-Q}$.
\end{lemma}
\begin{proof}\leanok
Use the exact pushforward identity for a dependent second-stage law.
The product-coordinate projection preserves the interval experiment.
Intersect events using the union bound; no independence of the retention
choices from their raw outcome is assumed.
\end{proof}

\begin{lemma}[Neighbourhood estimates under spoke changes]
\label{bp:marked-spoke-change-estimates}
\lean{Oberwolfach.finiteConjunction_sum_change_at_outer, Oberwolfach.retainedMarkedSpokes_restriction_changes, Oberwolfach.markedRowCountWithSpokes_change, Oberwolfach.markedIndexCountWithSpokes_change, Oberwolfach.markedRows_spoke_changes, Oberwolfach.markedIndices_spoke_changes, Oberwolfach.retainedMarked_restricted_col_changes, Oberwolfach.retainedMarked_restricted_row_changes}
\uses{bp:marked-row-patterns,bp:marked-indices-vortex,bp:interval-retention,bp:raw-marked-restriction}
\leanok
Keep the physical rims, admissible domains and marginal-product targets
fixed. If each requested spoke relation changes at at most $\delta$ points
of the original domain, a test with at most $h$ requested spokes changes
by at most $h\delta$. This holds for every bounded row or index test.
For a retained assignment the changes of either spoke type are bounded
by the number of changed pair labels. These bounds restrict to smaller
fixed domains under the commuting vertex and injective hub embeddings.
\end{lemma}
\begin{proof}\leanok
A changed conjunction has at least one changed predicate. Sum the resulting
indicator inequality, retaining the original outer domain even if the
admissible row set is smaller. Removing intervals may expose preassigned
ordinary labels, so both insertions and deletions are counted. Restrictions
transport this deterministic inequality and do not assert a new local
uniform-retention law.
\end{proof}

\begin{lemma}[Retained neighbourhood estimates across the vortex]
\label{bp:marked-retained-pattern-vortex}
\lean{Oberwolfach.permutation_invariant_image, Oberwolfach.marked_retained_rows_vortex_event, Oberwolfach.marked_retained_indices_vortex_event}
\uses{bp:marked-retention-law,bp:marked-spoke-change-estimates,bp:marked-rows-vortex,bp:marked-indices-vortex}
\leanok
On the single globally retained marked experiment, the bounded row
estimates hold at every fixed vortex level with failure at most
$\chi+n^{-Q}$ and error $(\varepsilon+h\zeta)|A_i|$.
The corresponding index event has failure at most $\chi+n^{-Q}$ and
error $(\varepsilon+h\zeta)|B_i|$. Their targets use the original supplied
marginals. The finite level count is bounded by the prescribed polynomial,
the relevant domains exceed the uniform retention threshold, and the
vertex restrictions are invariant. Both events concern the same law and
may be intersected with additive failure bounds.
\end{lemma}
\begin{proof}\leanok
Combine the proved raw vortex event with the global retention event on
images of the fixed domains. Commuting embeddings preserve their
invariance and cardinalities. Apply the deterministic spoke-change bounds
at each level. Restricted retained states are taken from that global
outcome, without resampling or replacing their distribution.
\end{proof}

\begin{lemma}[Balance and fixed potential loads after retention]
\label{bp:marked-retained-balance-load-stability}
\lean{Oberwolfach.markedRestrictedBalance_spoke_changes, Oberwolfach.typedMarkedPotentialLoad_spoke_degree_stability, Oberwolfach.retainedMarkedPotential_spoke_stability, Oberwolfach.retainedMarkedPotential_all_roots_stability, Oberwolfach.retainedMarkedDomain_spoke_stability, Oberwolfach.retainedMarkedDomain_all_roots_stability}
\uses{bp:marked-restricted-balance,bp:marked-spoke-change-estimates,bp:weighted-stability,bp:omitted}
\leanok
With physical rims fixed, spoke changes of degree at most $\delta$ increase
restricted balance errors by at most $(m+2)\delta$. For fixed parent wheel
coefficients bounded by $B/n^m$, the root-omitted potential load changes
by at most $5(m+1)^4B\delta/n$ at any spoke root, including external hubs.
This bound requires no bound on the total number of hubs. For all rim and
spoke roots in a domain with at most $n$ hubs, the bound is
$11(m+1)^4B\delta/n$. These conclusions also apply to restrictions of a
global retained assignment, with its actual change bounds on the images
of the domains.
\end{lemma}
\begin{proof}\leanok
Control individual degree changes by changed-predicate indicators and
apply the exact balance identities. For potential loads, use the proved
fixed-coefficient pair-incidence inequalities. The spoke-root proof does
not use the total hub-count assumption, which was removed and its caller
updated. The tested root remains omitted, even if absent; no sampling
factor multiplies a realised change.
\end{proof}


\begin{lemma}[Local configurations after global retention]
\label{bp:marked-retained-configurations-vortex}
\lean{Oberwolfach.fixedDomainTuple_predicate_loss, Oberwolfach.markedLocalBase_class_loss, Oberwolfach.markedLocalPath_class_loss, Oberwolfach.markedLocalConfigurations_spoke_changes, Oberwolfach.marked_retained_configurations_vortex_event}
\uses{bp:marked-local-vortex,bp:marked-retention-law,bp:marked-spoke-change-estimates}
\leanok
If the two spoke classes change by at most $\zeta|A_i|$ vertices per hub,
the local cycle and path coefficients decrease by at most $(r+5)\zeta$
and $(r+4)\zeta$, respectively. All free coordinates lie in the original
inner set $A_i$; the fixed path roots can lie in the larger ambient set.
The corresponding estimates hold simultaneously at every fixed vortex
level of the globally retained sample, with failure at most $\chi+n^{-Q}$.
\end{lemma}
\begin{proof}\leanok
Every lost tuple has a coordinate in the old spoke class but not the new
one. Count the remaining free coordinates in $A_i$. The physical-edge,
separation and distinctness conditions are unchanged, and no root-spoke
condition is inserted. Apply this deterministic bound on the common raw
and retention event, using the actual restrictions of the global sample.
\end{proof}

\begin{lemma}[Cut families after global retention]
\label{bp:marked-retained-cut-vortex}
\lean{Oberwolfach.bipartite_cut_change_degrees, Oberwolfach.bipartite_cut_lower_after_changes, Oberwolfach.markedCutFamilies_spoke_changes, Oberwolfach.marked_cut_vortex_event, Oberwolfach.marked_retained_cut_vortex_event}
\uses{bp:marked-reference-cuts,bp:vortex-common-events,bp:marked-retention-law,bp:marked-spoke-change-estimates}
\leanok
If each row and column changes at a fraction at most $\zeta$ of its full
domain, every spoke cut changes by at most
$\zeta\min\{|U||X|,|V||Y|\}$. Thus the reference spoke coefficient
$t_*/8$ becomes $t_*/8-\zeta$ at the original cutoff. The physical rim
cut inequalities are unchanged. All four cut families hold on every
fixed vortex domain after global retention, with failure at most
$\chi+n^{-Q}$ under the stated size and degree hypotheses.
\end{lemma}
\begin{proof}\leanok
Sum changed-edge indicators first over rows and then over columns and
take the smaller bound. Prove the raw cut estimates at each domain's own
order and sum their exponential failure bounds across the vortex. Transfer
these estimates on the global retention event. No local cut is inferred
by restricting an ambient cut inequality.
\end{proof}

\begin{lemma}[Restricted balance after global retention]
\label{bp:marked-retained-balance-vortex}
\lean{Oberwolfach.marked_retained_balance_vortex_event}
\uses{bp:marked-restricted-balance-vortex,bp:marked-retained-balance-load-stability,bp:marked-retention-law}
\leanok
For fixed nonempty restriction pairs $A_i,B_i$ with both sizes at most
$M_i$, the four balance inequalities after global retention have error
$(\varepsilon+(m+2)\zeta)M_i$, simultaneously at all fixed levels and
all outer roots. The failure bound is $\chi+n^{-Q}$. The original
parent-degree centres and the exact identities between those centres
and the supplied spoke densities remain hypotheses.
\end{lemma}
\begin{proof}\leanok
Intersect the raw balance event with retention on the images of $A_i,B_i$.
Their change bounds are at most $\zeta M_i$. Apply the deterministic
balance-transfer theorem; the error is not rescaled by the smaller
admissible domains. The separate empty-pair identity covers empty pairs.
\end{proof}


\begin{lemma}[Constructed weights and retained potential loads]
\label{bp:marked-retained-load-realization}
\lean{Oberwolfach.markedLocalSpokeLoads_spoke_changes, Oberwolfach.markedPackingLoads_spoke_changes, Oberwolfach.marked_retained_local_load_vortex_event, Oberwolfach.marked_retained_packing_load_vortex_event, Oberwolfach.marked_retained_local_spoke_realization, Oberwolfach.marked_retained_packing_realization}
\uses{bp:marked-local-load-vortex,bp:marked-packing-load-vortex,bp:marked-retention-law,bp:marked-retained-balance-load-stability}
\leanok
Construct a separate successor weighting on each fixed parent domain
before supplying the actual marginals. Keep its normalization unchanged
through raw sampling and global retention. If the outer hub order is
at most $K$ times the local vertex order and the normalized coefficients
are bounded above by $B_+/n^m$, the additional potential-load error is
at most $5(m+1)^4B_+K\zeta$ for external spoke roots and
$11(m+1)^4B_+K\zeta$ for the combined packing and external-root tests.
The packing hub set is explicitly embedded in the larger root domain
and has at most $n$ hubs. The retained estimates hold on all fixed vortex
levels with failure at most $\chi+n^{-Q}$; the weights retain their proved
two-sided bounds. The parent degree tolerance precedes the separation
parameter, and all order thresholds precede the hosts and marginal arrays.
\end{lemma}
\begin{proof}\leanok
Use the proved successor-weight existence theorem and normalized bounds.
Apply the raw load event and the global retention event on the fixed
domain images. Transfer loads by their actual spoke symmetric differences.
Restriction to the inner hub set is injective and cannot increase a change
sum. Keep the tested root indicator omitted and the original eligibility
condition for potential rim roots. No weight is recomputed after exposure.
\end{proof}


\begin{lemma}[Number of fixed vortex restrictions]
\label{bp:geometric-restriction-family-size}
\lean{Oberwolfach.geometric_level_count_le, Oberwolfach.geometric_restriction_family_quadratic}
\uses{bp:vortex-common-events}
\leanok
For positive integer scales $M_0,\ldots,M_L$ with
$M_{i+1}\le\eta M_i$ and $\eta<1$, one has $L+1\le M_0$.
Consequently, if $M_0\le n$ and $k\le n$, a family with $k$ fixed
restriction types at each level has at most $n^2$ members.
\end{lemma}
\begin{proof}\leanok
The integer scales decrease by at least one at each step. Induction gives
$M_i+i\le M_0$, and the terminal scale is at least one. Multiply the
level bound by the number of fixed restriction types.
\end{proof}


\begin{lemma}[Two-hub configurations after global retention]
\label{bp:marked-retained-common-configurations}
\lean{Oberwolfach.twoHubTuple_spoke_loss,Oberwolfach.rootFirstCommonMarkedPath_spoke_changes,Oberwolfach.commonMarkedCycle_spoke_changes,Oberwolfach.markedCommonConfigurations_spoke_changes,Oberwolfach.marked_common_configurations_vortex_event,Oberwolfach.marked_retained_common_configurations_vortex_event}
\uses{bp:marked-configurations,bp:marked-retention-law,bp:raw-marked-restriction,bp:vortex-common-events}
\leanok
For the fixed vortex domains, compatible-hub counts, two-hub cycle counts
and root-omitted common-path counts hold after the single global retention,
with probability at least $1-\chi-n^{-Q}$. The raw coefficient $\beta>0$
is chosen before $\eta,\chi,d$ and the hosts. If each spoke row and column
changes at most $\zeta n_i$ entries, the three retained lower coefficients
are respectively
\[
32\beta-(r+5)\zeta,\qquad
16\beta-2(r+5)\zeta,\qquad
16\beta-2(r+4)\zeta.
\]
The common-path family omits all four root spokes. Its roots need not be
spoke neighbours, and no additional separation between those roots is
imposed. The actual sampled physical arcs are unchanged.
\end{lemma}
\begin{proof}
\leanok
A lost tuple has a free coordinate incident to a changed spoke at one of
the two fixed hubs. Summing over the free coordinate gives the stated
losses. Compatible hubs are bounded by the row symmetric differences.
The original product law restricts to each fixed domain; the retained
sample is restricted from the one global experiment, not resampled locally.
A geometric sum controls raw failures, and the existing conditional
retention event controls all fixed restrictions simultaneously.
\end{proof}


\begin{lemma}[Common-witness potential loads with fixed selectors]
\label{bp:thinned-common-weight-witness}
\lean{Oberwolfach.thinnedMarked_spoke_probability,Oberwolfach.thinnedMarkedCoordinates_candidate_disjoint,Oberwolfach.thinnedMarkedPotentialLoad_expectation,Oberwolfach.thinnedMarkedPotentialLoad_mean_event,Oberwolfach.thinnedMarkedSeed_spoke_mean,Oberwolfach.thinnedMarkedSeed_spoke_event,Oberwolfach.commonThinnedSpokes_potential_load_change,Oberwolfach.rawMarked_common_weight_witness_event,Oberwolfach.rawMarked_common_weight_exponential_event,Oberwolfach.retainedMarked_common_weight_event,Oberwolfach.retainedMarked_common_weight_realization}
\uses{bp:marked-seed-loads,bp:marked-retention-law,bp:weighted-stability,bp:common-event,bp:bounded-differences}
\leanok
For each distinct hub pair, give each vertex an independent thinning coin
of mean $1/m$, where the cycle length is $m+1$. The ordinary common class
has marginal $a_0=s^2/m$ and the marked common class has marginal
$q=(s/m)^2=a_0/m$. With the parent weighting fixed before exposure,
the root-omitted common potential sums have means within $\tau$ of one
and exponentially small deviations.

After the original global retention, one event supplies selectors for
all distinct hub pairs with potential-load error
\[
 \tau+\varepsilon+10(m+1)^4B_+\zeta,
 \qquad B_+=2B/p_*^{2m+1}.
\]
For a fixed inverse-power exponent $Q$, the order threshold precedes
the host and the marginal arrays, and the failure probability is at most
$2n^{-Q}$. The shared parent weighting is constructed from the proved
successor-weight theorem, before the sampling marginals are supplied.
\end{lemma}
\begin{proof}
\leanok
The common spoke is represented by its two original hub-column indicators
and, for an ordinary spoke, its independent coin. Actual coordinate supports
are disjoint within each separated parent cycle. Omitting a root removes
both of its spokes and its coin. Positional normalization retains the
parent error $\tau$. The coordinate read and influence multiplicities are
at most $16d+14$, giving an exponential concentration bound. A finite union
and projection of the auxiliary coins produce one raw event for every
hub pair. Each selector and the parent weighting are then kept fixed;
column symmetric differences give the stated deterministic transfer loss.
The final order threshold also ensures the required minimum size of the
hub domain.

This proves the weighted component only. The additional class and
cycle/path count conclusions must use the same selectors and are not
asserted by this lemma. Their joint raw realization and retained realization are proved below.
\end{proof}


\begin{lemma}[One selector for common sizes, counts and loads]
\label{bp:common-joint-witness}
\lean{Oberwolfach.supportedIndicatorCount_tail,Oberwolfach.thinnedCommonClassCount_common_event,Oberwolfach.thinnedCommonClass_disjoint,Oberwolfach.thinnedCommonClass_dense_event,Oberwolfach.thinnedCommonClass_size_event,Oberwolfach.thinnedMarked_two_exposures,Oberwolfach.thinnedCommon_configurations_event,Oberwolfach.thinnedCommon_joint_event,Oberwolfach.rawMarked_common_joint_witness_event}
\uses{bp:thinned-common-weight-witness,bp:marked-configurations,bp:bounded-differences,bp:common-event}
\leanok
On one event in the original raw sample, of probability at least
$1-n^{-Q}$, every distinct hub pair has a selector that simultaneously
satisfies the common-class size estimates, the dense ambient-root
neighbourhood bounds, the ordered common-cycle and root-omitted path
counts, and the common potential-spoke load estimates. The shared parent
weighting is fixed across all pairs. The order threshold precedes the
host and the original marginal arrays. In particular, the selector in
the count inequalities is the selector in the load inequalities.
\end{lemma}
\begin{proof}\leanok
Bounded differences controls common-class indicators on fixed domains,
including neighbourhoods of all ambient roots and their inverse images
under the permutation. The actual experiment equals first exposing
interval coordinates and vertex coins, then exposing the original
physical states. The proved greedy and physical-survival arguments apply
to these classes in that second exposure. Accurate class sizes are also
retained. Intersect these events with the weighted-load event before
projecting to existence of a selector. A union bound over hub pairs,
without an independence assumption, gives the stated probability.

The configuration counts here concern explicit ordered candidate
families. The retained realization below converts these to the canonical
common cycles and root-first paths. The full preparation and absorber
assembly remain separate obligations.
\end{proof}


\begin{lemma}[Constructed common witnesses after global retention]
\label{bp:retained-common-joint-realization}
\lean{Oberwolfach.selectedBase_card_le_common_cycle,Oberwolfach.selectedPath_card_le_common_path,Oberwolfach.commonSelectedClass_card_change,Oberwolfach.commonSelectedClass_overlap_change,Oberwolfach.rawCommonWitness_transfer,Oberwolfach.commonWitness_error_margins,Oberwolfach.retainedMarked_common_joint_witness_event,Oberwolfach.retainedMarked_common_joint_realization,Oberwolfach.successorCycleRatio_eventually_margin,Oberwolfach.marked_common_density_bounds,Oberwolfach.marked_common_completion_realization,Oberwolfach.retainedMarkedSpokes_disjoint,Oberwolfach.retainedMarkedSpokes_successor_disjoint}
\uses{bp:common-joint-witness,bp:marked-retention-law,bp:marked-spoke-change-estimates,bp:marked-sample-physical-map}
\leanok
The original globally retained experiment has one event, of probability
at least $1-2n^{-Q}$, on which every distinct hub pair admits a single
selector satisfying the common-class size, overlap, potential-load,
canonical cycle and root-first path estimates. The parent weight is
constructed before the sampling marginals and has positive two-sided
bounds on every valid separated parent cycle. The same weight is used
for every pair.

The positive cycle and path constants precede the requested error
$\eta$. The near-regularity tolerance precedes the separation length;
the order threshold precedes the host and all marginal arrays.
The actual completion densities $p=\rho s/2$, $k=p/m$, $t=s/m$
meet the required lower bounds uniformly over the prescribed density
ranges. The condition on the successor-cycle ratio follows for all
sufficiently large even lengths from the minimum-degree margin.
\end{lemma}
\begin{proof}\leanok
Inject the surviving ordered configurations into the actual common-cycle
and root-first path types. A changed spoke column of size at most
$\zeta n$ changes either common class by at most $2\zeta n$;
the general comparison bounds successor overlap by $4\zeta n$.
The common potential-load loss is $10(m+1)^4B_+\zeta$, while the
cycle and path coefficients lose at most $2(m+1)\zeta$ and
$2m\zeta$, respectively. Choose $\zeta$ so that the errors fit
$\eta$ and the two count coefficients retain half their positive
raw values. Intersect the raw event with the proved global retention
event, keeping the same selector and coefficient function.

The retained spoke types are deterministically disjoint, as are the
ordinary class and the successor image of the marked class. This uses
the retained intervals and the original ordinary-label variables.
The existing categorical physical-edge representation is unchanged.
The full absorber and preparation constructions remain unformalised.
\end{proof}


\begin{lemma}[Simultaneous selection with a fixed degree bound]
\label{bp:bounded-degree-selection}
\lean{Oberwolfach.bounded_degree_selection,Oberwolfach.sequential_select_bounded_hits,Oberwolfach.graph_candidates_hit_bound,Oberwolfach.graph_iSup_local_degree_bound,Oberwolfach.bounded_degree_selection_order,Oberwolfach.greedy_graph_candidates,Oberwolfach.bounded_degree_selection_from_single_copy,Oberwolfach.cover_down_candidate_budget,Oberwolfach.mul_sqrt_eq_rpow_three_halves}
\leanok
Fix $p>0$. An order threshold is chosen before the host graph, the
number $N$ of supports, and all candidate families. Suppose
$|S_i\cap S_j|\le p^2n$ for $i\ne j$, and every vertex belongs to at
most $pn$ supports. If every deletion of maximum degree at most
$Dp^{3/2}n$ from the $i$th support permits
$\lceil4p^{3/2}n\rceil$ edge-disjoint candidates, each of maximum
degree at most $D$, then one can select a candidate from every family
so that they are pairwise edge-disjoint and their union has maximum
degree at most $Dpn$.
\end{lemma}
\begin{proof}\leanok
The fixed tests are precisely $(i,v)$ with $v\in S_i$; there are at
most $pn^2$ of them, irrespective of $N$. Each test records the number
of earlier chosen graphs using an edge from $v$ into $S_i$.
Edge-disjoint candidates and the support intersection bound show
that at most a fraction $\sqrt p/4$ of the next candidates hit any
active test. Conditional averaging of
\[
 \Phi=\sum_x 2^{Z_x}(1+\sqrt p/4)^{R_x}
\]
allows a next choice without increasing $\Phi$. Here $Z_x$ counts
actual previous hits and $R_x$ counts remaining indices that can hit
that test. The initial potential is less than
$2^{\lfloor p^{3/2}n\rfloor+1}$ for all sufficiently large $n$,
uniformly in the supports. Thus every needed local deletion has
degree at most $D\lfloor p^{3/2}n\rfloor$ and the candidate
hypothesis applies. No independence between successive choices is
assumed. The final degree bound uses support memberships.
Repeated single-copy existence supplies the candidate family with
all prior edges included in the deletion budget. For the two
cover-down callers, $D=4$ and $n_S\le2a$ give the checked bound
$4p^{3/2}n_S+4\lceil4p^{3/2}n_S\rceil\le40p^{3/2}a+4$.
The exact factor and endpoint-matching hypotheses at those callers
remain separate obligations; this result does not establish them.
\end{proof}


\begin{lemma}[Prescribed paths and rooted cycle embeddings]
\label{bp:prescribed-paths}
\lean{Oberwolfach.greedy_simple_path_avoiding,Oberwolfach.fixed_length_path,Oberwolfach.embedded_cycle_from_simple_path,Oberwolfach.fixed_cycle_through_edge,Oberwolfach.fixed_cycle_through_vertex}
\uses{bp:graphs}
\leanok
For fixed $\beta>0$, $k\ge2$ and forbidden-set bound $b$, the order
threshold is chosen before the graph, endpoints and forbidden set.
Every graph on at most $n$ vertices with minimum degree at least
$(1/2+\beta)n$ contains a simple path of length $k$ between prescribed
distinct endpoints, avoiding a prescribed set of at most $b$ other
vertices. Closing a simple path along an endpoint edge gives an
actual embedded cycle graph. Corresponding cycle embeddings through
an edge or a vertex are also proved, with the same avoidance scope.
\end{lemma}
\begin{proof}\leanok
Greedy extension avoids the entire previous path support and the
forbidden set. The final two edges use a common neighbour outside
that support. The common-neighbour lower bound is $2\beta n$.
The cycle conversion proves an injective map from the actual cycle
graph, including its closing edge; a walk with repeated vertices is
not used as a cycle.
\end{proof}

\begin{lemma}[Wheel attachments at prescribed rim roots]
\label{bp:ordinary-wheel-attachments}
\lean{Oberwolfach.common_spoke_hub_avoiding,Oberwolfach.ordinary_rooted_wheel_embedding,Oberwolfach.ordinaryWheelResources_disjoint_iff,Oberwolfach.ordinaryWheels_disjoint_at_one_root,Oberwolfach.ordinary_wheel_attachments}
\uses{bp:prescribed-paths}
\leanok
A spoke degree of at least $(1-1/(2r))|U|$ at every rim vertex
suffices to choose a common hub for a prescribed $r$-cycle outside
any hub set of size less than $|U|/2$.
Consequently one can attach $N$ wheels at specified rim roots,
allowing a root to repeat. The order threshold depends on
$\beta,r,N$ and the bounded root-set size; the spoke-density
constant depends only on $r$. The wheels have distinct hubs,
are edge-disjoint, meet the original rim-root set only at their own
specified roots, and have no common new rim vertices.
\end{lemma}
\begin{proof}\leanok
At each step the forbidden rim set is the union of the root set
and all previous cycle supports, with the current root removed.
It has size at most $q+Nr$. The forbidden hub set includes every
previous hub. Two wheels sharing at most one rim vertex and having
different hubs are proved edge-disjoint in the physical edge
representation. Thus repeated prescribed roots do not create
repeated edges.
\end{proof}

\begin{lemma}[Exact finite transformer and absorber assembly]
\label{bp:resource-absorber-assembly}
\lean{Oberwolfach.HasResourceDecomposition,Oberwolfach.IsResourceTransformer,Oberwolfach.IsResourceAbsorber,Oberwolfach.resource_transformer_of_indicator_identity,Oberwolfach.resource_transformer_compose,Oberwolfach.resource_transformer_chain,Oberwolfach.resource_absorber_chain_with_attachment,Oberwolfach.resource_absorber_chain_avoids}
\leanok
An exact decomposition is a finite family of pairwise disjoint
physical edge sets whose union equals the specified edge set.
The two transformer decompositions compose along any finite chain,
including the manuscript's chain of four transformers. Each
intermediate remainder occurs exactly once. A decomposable final
remainder and the decomposable attachment edges then give the two
absorber decompositions. The corresponding forbidden-edge avoidance
statement is preserved.
An equality of edge indicator differences also yields the exact
common transformer edge set, under the proved containment and
edge-disjointness hypotheses stated in that interface.
\end{lemma}
\begin{proof}\leanok
Each composition partitions its edges into two disjoint pairs:
$(T\cup L)\cup(T'\cup M)$ and
$(T\cup M)\cup(T'\cup N)$. Induction retains every intermediate
edge set exactly once. Removing the decomposable attachment part
from the initial remainder leaves an absorber for the original
remainder. These are finite assembly results: the host embeddings
and existence of the individual transformers are still open.
\end{proof}

\begin{lemma}[Ordinary physical decomposition representation]
\label{bp:ordinary-exact-decompositions}
\lean{Oberwolfach.ordinaryHostResources_union,Oberwolfach.ordinaryHostResources_sdiff,Oberwolfach.ordinaryHostResources_injective,Oberwolfach.HasOrdinaryWheelDecomposition,Oberwolfach.ordinary_packing_exact_decomposition,Oberwolfach.ordinary_wheel_decomposition_union,Oberwolfach.ordinary_wheel_decomposition_exact_loads,Oberwolfach.ordinary_wheel_decomposition_balance}
\uses{bp:resource-absorber-assembly}
\leanok
Ordinary rim edges are unordered pairs in $V$; spokes are pairs in
$V\times U$. Their tagged physical sets preserve unions, differences
and graph equality. The exact wheel decomposition uses actual
embedded cycles and hubs, and its indicator sums give unit load on
each host edge. It follows that rim degree is twice spoke degree
at each rim vertex and every hub degree is divisible by $r$.
These identities are derived from an exact partition, not inferred
from a fractional decomposition.
\end{lemma}

\begin{lemma}[A small reservoir for every later union of root sets]
\label{bp:universal-absorber-reservoir}
\lean{Oberwolfach.injective_assignment_of_many_choices,Oberwolfach.assigned_absorber_reservoir,Oberwolfach.disjoint_configuration_selection_finite,Oberwolfach.configuration_reservoir_order,Oberwolfach.universal_absorber_reservoir,Oberwolfach.small_universal_absorber,Oberwolfach.profile_root_partition,Oberwolfach.paired_absorber_decompositions,Oberwolfach.paired_absorber_free_card}
\uses{bp:resource-absorber-assembly}
\leanok
Suppose a fixed family of configurations has bounded free-set size,
polynomial total and intersecting-pair counts, and a polynomial-size
family of root sets each admits a positive proportion of compatible
configurations. For any prescribed positive size allowance, there are
positive constants and a uniform order threshold giving a disjoint
configuration family whose free vertices fit that allowance and which
retains linearly many configurations for every root set.
If every configuration has an internal factor, the resulting free set
absorbs every later disjoint union of sufficiently few admissible root
sets lying outside it. Exact prescribed class sizes allow a leftover set
to be partitioned into such root sets. The two factors and the free-set
size in the paired absorber construction are also verified.
\end{lemma}
\begin{proof}\leanok
Select configurations independently with the same probability, before
any leftover is chosen. Chernoff's inequality holds simultaneously for
all fixed compatibility sets; Markov's inequality bounds the total
selection and selected intersecting pairs. Remove every configuration
incident with such a pair, charging the loss to its two endpoints.
A single remaining family therefore retains the asserted count for
every root. Hall's theorem gives distinct compatible configurations
for all roots in a later leftover. The two exact vertex partitions are
assembled, retaining the internal factors of unused configurations.
This proves the reservoir selection and finite assembly. The required
graph-specific configuration counts, weight correction and remaining
factor-extraction steps are separate obligations.
\end{proof}

\begin{lemma}[Exact correction from signed absorber trades]
\label{bp:profile-signed-correction}
\lean{Oberwolfach.fixedSubset_centred_moment,Oberwolfach.profileDemandCoefficient_vertex_identity,Oberwolfach.profileDemandCoefficient_bound,Oberwolfach.profile_defect_centred_data,Oberwolfach.profileRootSet_hit_probability,Oberwolfach.averaged_absorber_trade_incidence,Oberwolfach.averaged_absorber_trade_bound,Oberwolfach.profile_trade_correction,Oberwolfach.correctedTradeWeight_two_sided}
\uses{bp:resource-absorber-assembly}
\leanok
Let class $i$ have $a_i k$ vertices, and let every indexed cycle
contain $a_i$ vertices of that class. Choose an $a_i$-subset uniformly
in each class, independently between classes. The explicit signed
root coefficients correct any initial vertex-load error at most
$\eta$, provided the supplied trade has incidence equal to the root
indicator. If the trade coefficient at a cycle is at most
$K/k^c+\mathbf1_{\{C\cap X\ne\varnothing\}}K/k^{c-1}$ in absolute
value, then the change in each cycle weight is at most
$(1+4c)(c+1)K\eta/k^{c-1}$.
The two actual absorber partitions give the required incidence identity.
The following lemma supplies the combinatorial occurrence estimate
from the manuscript's labelled absorber and multiplicity hypotheses.
\end{lemma}
\begin{proof}\leanok
The fixed-size one-vertex and two-vertex sampling probabilities give
the exact centred moment. Independence between classes eliminates
cross terms. The probability that a root meets a fixed cycle is at
most $c/k$, by the union bound. Averaging the signed trades therefore
gives both the exact vertex sums and the displayed coefficient bound.
A lower bound on the original weights and a sufficiently small error
then preserve two-sided positive coefficient bounds.
\end{proof}

\begin{lemma}[Vertex-sum correction from labelled absorber abundance]
\label{bp:profile-marginal-correction}
\lean{Oberwolfach.rooted_label_maps_constrained_count,Oberwolfach.labelled_cycle_records_occurrence_card,Oberwolfach.labelled_absorber_factors,Oberwolfach.labelled_absorber_occurrence_disjoint,Oberwolfach.labelled_absorber_occurrence_general,Oberwolfach.labelled_family_trade_spread,Oberwolfach.profile_marginal_correction}
\uses{bp:profile-signed-correction}
\leanok
The vertex-sum correction lemma follows from the stated lower bound on
labelled absorber counts, the finite template family and the multiplicity
bound on indexed cycles with a given vertex set. The error threshold is
chosen before the vertex classes and indexed cycle family. Both positive
coefficient bounds in the conclusion are retained. No trade-spread
assumption is added.
\end{lemma}
\begin{proof}\leanok
Keep the template, full injective label map and every cycle identifier
in each record. A specified cycle disjoint from the roots fixes all its
$c$ free positions. A cycle meeting the roots fixes at least $c-1$
free positions, since each template occurrence has at most one root.
The remaining cycle identifiers each have at most $J$ choices on their
specified vertex set. These counts give the signed-trade bound after
dividing by the number of labelled absorbers. The exact fixed-size root
coefficient calculation gives unit vertex sums. A sufficiently small
error, depending only on the fixed parameters, keeps every corrected
cycle weight between the required two bounds.
\end{proof}


\begin{lemma}[Exact profile-factor extraction]
\label{bp:profile-factor-extraction}
\lean{Oberwolfach.profile_factor_extraction,Oberwolfach.profile_absorber_reservoir,Oberwolfach.labelled_absorber_configuration_reservoir,Oberwolfach.restricted_absorber_abundance,Oberwolfach.profile_fractional_factor,Oberwolfach.resource_restriction_absorption_completion}
\uses{bp:profile-marginal-correction,bp:universal-absorber-reservoir}
\leanok
For a fixed positive integer class profile, fixed template family and
bounded multiplicity of cycle identifiers per vertex set, sufficiently
accurate positive vertex weights and sufficiently many labelled absorbers
imply an exact vertex factor. All tolerance and order constants precede
the host classes and the indexed cycle family.
\end{lemma}
\begin{proof}
\leanok
Count root-free configurations and intersecting pairs while retaining
their internal cycle identifiers. The map forgetting a root has bounded
fibres. Select a small disjoint reservoir with exact class proportions;
it works simultaneously for every later sufficiently small leftover.
Deleting the reservoir loses at most its size times the pair-weight
bound at each surviving vertex. A separate labelled count bounds the
lost absorbers. Restriction preserves the actual cycle identifiers,
template maps and supplied weights. Correct the surviving vertex sums
exactly, apply the proved fractional-resource matching theorem with
vertices as resources, and absorb the uncovered class-profile set.
This proves the finite-resource form of the manuscript's profile-factor
lemma. The graph-specific construction and counting of the absorber
families needed by later callers remain separate obligations.
\end{proof}


\begin{lemma}[Paired-absorber representation and fixed-size sampling]
\label{bp:profile-factor-sampling}
\lean{Oberwolfach.pairedCycleRecordToLabelled_injective,Oberwolfach.paired_profile_factor_extraction,Oberwolfach.profile_induced_degree_tail,Oberwolfach.absorber_profile_count_tail,Oberwolfach.embedded_profile_weighted_degree_expectation,Oberwolfach.profile_root_conditional_probability,Oberwolfach.profile_roots_common_event}
\uses{bp:profile-factor-extraction}
\leanok
Paired cycle records inject into the labelled template records used by
factor extraction. Independent uniform class samples have exact
hypergeometric inclusion probabilities. Rooted weighted cycle sums and
labelled absorber counts satisfy the stated exchange concentration bounds,
with the sampled ground size and original multiplicities. Conditioning
on fixed roots gives uniform samples of the remaining populations.
\end{lemma}
\begin{proof}\leanok
Retain the ordered free labels and every cycle identifier in the paired
record. For weighted degrees, bound the changed cycles by prescribing the
root and an exchanged label. For absorber counts, prescribe an exchanged
free position; count all remaining labels in the selected union. The
proved fixed-size exposure inequality gives concentration. Explicit
bijections identify the conditioned law. Weighted union bounds use the
number of roots contained in each outcome, and therefore introduce no
factor for all ambient root labels. These statements are ingredients for
the full finite-resource bounded-extension theorem in Node~\ref{bp:profile-bounded-extension}. Graph-specific applications still require their actual absorber and weight estimates.
\end{proof}

\begin{lemma}[Root-omitted means and conditional class sampling]
\label{bp:rooted-profile-means}
\lean{Oberwolfach.profile_rooted_weighted_degree_tail,Oberwolfach.profile_rooted_absorber_count_tail,Oberwolfach.normalized_profile_inclusion_error_uniform,Oberwolfach.rooted_deletion_profile_expectation,Oberwolfach.rooted_deletion_rescaled_mean_error,Oberwolfach.absorber_deletion_profile_mean_lower}
\uses{bp:profile-factor-sampling}
\leanok
Conditioning on the prescribed root labels preserves the uniform
fixed-size law on the remaining populations. A cycle through a root
loses one sampled label only in that root's class. Its expected weighted
contribution is the exact product of the corresponding hypergeometric
ratios, times its original retained weight. The error separates the
sample-size correction from the ambient-population correction.
\end{lemma}
\begin{proof}\leanok
Remove cycles meeting another fixed forbidden vertex and retain the
weight on each surviving identifier. Count their omitted supports in
each class and apply the exact inclusion formula. The vertex-deletion
loss is subtracted before sampling. For absorbers, count the original
records whose free labels avoid the forbidden set; only their free
labels are sampled. Bound the product-ratio error and the number of
deleted records separately. The conditioned concentration inequalities
are transported back to the original sample space with the probability
of including the roots retained explicitly.
\end{proof}

\begin{lemma}[Factors avoiding sparse displayed pairs]
\label{bp:sparse-pair-factors}
\lean{Oberwolfach.labelled_absorber_local_free_pair_count,Oberwolfach.labelled_absorber_occurrence_root_unique,Oberwolfach.labelled_absorber_bounded_pair_loss,Oberwolfach.resource_bounded_pairs_degree_loss,Oberwolfach.candidate_absorber_card_loss,Oberwolfach.profile_factor_avoiding_pairs,Oberwolfach.profile_two_disjoint_factors}
\uses{bp:profile-factor-extraction}
\leanok
Under the profile-factor hypotheses, a factor exists that avoids every
pair in a fixed set of displayed endpoint pairs having bounded incidence
at each vertex and at most as many pairs as vertices. The weight tolerance
is chosen before the incidence bound. The order threshold may depend on
that bound. In particular, if each indexed cycle has at most two displayed
pairs at each vertex and at most $c$ displayed pairs in total, there are
two vertex factors with disjoint displayed pair sets.
\end{lemma}
\begin{proof}\leanok
A forbidden pair meeting a root set fixes a free absorber label, because
each occurrence contains at most one root. A pair disjoint from the roots
fixes two distinct free labels. Explicit counts retaining all cycle
identifiers give an absorber loss of order $k^{c^2-1}$. The local
weighted-degree loss is bounded by two- and three-vertex cycle counts,
and is of order $1/k$. Restrict the indexed cycle family, transport the
surviving records injectively, and apply the proved factor-extraction
theorem. For the second conclusion, construct the first factor and prove
the required bounds on the union of its displayed pairs. Apply the
avoidance theorem to obtain the second factor. The avoidance predicate
forbids both endpoints in a retained cycle, which is stronger than
forbidding that pair as one of its edges; the same explicit deletion
bounds cover this stronger condition. This is a same-ground-set result;
the full bounded-extension theorem with a prescribed additional set
remains a separate obligation.
\end{proof}

\begin{lemma}[One class sample and deterministic extension factors]
\label{bp:bounded-extension-stages}
\lean{Oberwolfach.profile_fixed_size_conditional_probability,Oberwolfach.profile_all_sampling_tests_exists,Oberwolfach.profile_sample_order_union_margin,Oberwolfach.profile_remaining_forbidden_pairs_probability,Oberwolfach.profile_external_root_conditioned_degree_tail,Oberwolfach.uniform_profile_linear_rate_tail,Oberwolfach.absorber_extension_ground_count,Oberwolfach.extension_resource_root_degree,Oberwolfach.extension_resource_degree_change,Oberwolfach.profile_factor_on_subset,Oberwolfach.profile_extension_factor}
\uses{bp:rooted-profile-means,bp:profile-factor-extraction}
\leanok
The class sampling law admits simultaneous vertex and eligible-root
estimates with explicit total failure bounds. Forbidden endpoint pairs
contribute at most $4Dt^2/k$ when their number is at most $Dk$.
For disjoint sets $A,Z$, explicit root-omitted weight estimates and
absorbers with free vertices in $A$ imply a factor on $A\cup Z$ in
which each cycle meets $Z$ at most once. The tolerance is chosen before
the bound on $|Z|$.
\end{lemma}
\begin{proof}\leanok
Use exact without-replacement inclusion probabilities on the remaining
class populations. The eligible-root union bound counts roots present
in each outcome. Bound the sample order before the ambient order.
For deterministic extraction, map the prescribed vertex subset to its
class subtypes and preserve every original cycle identifier. An
absorber occurrence can meet $Z$ only at its roots, and contains at
most one root. Its actual recorded cycles therefore survive the
restriction. The vertex sums at $Z$ are the exact root-omitted sums;
at $A$ their change is bounded by local pair counts. Apply the proved
factor theorem and lift the resulting factor to the original family.
The complete finite-resource bounded-extension assembly is now proved
in the later bounded-extension node. Graph-specific constructions of
the supplied cycle families and absorber abundance remain separate.
\end{proof}

\begin{lemma}[Conditional means on the original sample space]
\label{bp:nested-profile-means}
\lean{Oberwolfach.nested_profile_embedding_range,Oberwolfach.nested_profile_remaining_card,Oberwolfach.embedded_rooted_original_weight_expectation,Oberwolfach.conditional_root_original_mean,Oberwolfach.conditional_absorber_original_mean,Oberwolfach.profile_fixed_size_weighted_linear_tail,Oberwolfach.profile_fixed_size_absorber_linear_tail}
\uses{bp:rooted-profile-means,bp:bounded-extension-stages}
\leanok
After fixed deletion and root conditioning, the remaining populations
are exactly the complements of the union of the two fixed sets.
The conditional weighted-cycle and absorber means are exact products
of hypergeometric ratios times the corresponding original retained
weight sum or labelled-record count. Both functions have conditional
exponential tails whose constants do not depend on the ambient order.
\end{lemma}
\begin{proof}\leanok
The composed embedding preserves each class index. Count preimages
in each class without changing labels. A root-omitted cycle loses one
position only in its root's class. Restricting its support to the sample
and that root already excludes the other deleted vertices, so its
supplied weight is unchanged. Absorber counts test only free labels;
fixed roots in the sample do not alter that test. Bound the effects of
one exchanged sample label and apply the fixed-size concentration
inequality. The original root-inclusion probability remains a factor
in each joint failure bound.
\end{proof}

\begin{lemma}[Two extension factors from explicit local estimates]
\label{bp:two-extension-factors}
\lean{Oberwolfach.candidate_absorber_ground_loss,Oberwolfach.candidate_absorber_ground_pair_loss_scale,Oberwolfach.candidate_resource_local_pair_degree_loss_scale,Oberwolfach.profile_extension_avoiding_pairs,Oberwolfach.profile_two_extension_factors_from_counts}
\uses{bp:bounded-extension-stages,bp:sparse-pair-factors}
\leanok
For disjoint sets with class sizes $a_i t$ and $a_i z$, where $z$ is
bounded and $t$ is sufficiently large, the explicit root-omitted weight
estimates and absorber counts imply factors on $A$ and $A\cup Z$.
Their displayed edge sets are disjoint, and every cycle of the latter
factor meets $Z$ at most once. The error tolerance precedes the bound
on $z$ in the quantifier order.
\end{lemma}
\begin{proof}\leanok
Extract the first factor on $A$. Its displayed pairs have cardinality
at most $|A|$ and incidence at most two at each vertex. Restrict the
original indexed family to cycles avoiding these endpoint pairs.
Both the weight loss and the absorber loss are measured on the actual
local vertex set; the counted absorbers still have all free vertices
in $A$. Apply the enlarged-set factor theorem to the retained family
and lift its cycle identifiers back. The later bounded-extension node
constructs one sample supplying all these estimates simultaneously.
\end{proof}


\begin{theorem}[Two bounded factors from the original host estimates]
\label{bp:profile-bounded-extension}
\lean{Oberwolfach.conditional_root_mean_bound,Oberwolfach.conditional_absorber_mean_lower_from_host,Oberwolfach.profileRequiredRoots_eligibility,Oberwolfach.profile_host_vertex_tail,Oberwolfach.profile_host_absorber_tail,Oberwolfach.profile_host_good_sample,Oberwolfach.profile_sampling_numerical_margins,Oberwolfach.profile_two_bounded_factors}
\uses{bp:nested-profile-means,bp:two-extension-factors}
\leanok
The finite indexed-resource version of the bounded-extension theorem
holds with the original host weight and labelled-absorber estimates.
The tolerance precedes the additional-set bound; the bounded sample
order precedes the forbidden-set and forbidden-pair bounds. Both
factors, on the sample and on its union with the prescribed additional
set, are constructed. Their displayed edge sets are disjoint, the
sample avoids all forbidden pairs, and each cycle of the second
factor meets the additional set at most once.
\end{theorem}
\begin{proof}\leanok
Condition on exactly those fixed root labels which must belong to the
sample. Incompatible roots cannot occur. Derive the two conditional
means from the original weights and records, with the actual deleted
sets and class populations. Apply the fixed-size exponential tails.
Weight each joint failure probability by its root-inclusion probability,
then sum over roots present in an outcome. The forbidden-pair event is
included in this same common event. Explicit numerical thresholds
preserve the stated quantifier order. Rescale the unchanged cycle
weights, transfer the sampled counts from scale $t$ to $t+z$, and
apply the deterministic two-factor theorem. Graph-specific verification
of the original host hypotheses and the final main-theorem assembly
are not supplied by this result.
\end{proof}


\begin{lemma}[Counting actual paired hub absorbers]
\label{bp:marked-hub-paired-absorbers}
\lean{Oberwolfach.sequential_disjoint_tuple_count,Oberwolfach.paired_absorber_abundance,Oberwolfach.rootFirstCommonMarkedPath_hub_cycles,Oberwolfach.marked_hub_paired_absorber_abundance}
\uses{bp:profile-factor-sampling,bp:marked-common-path-coordinates}
\leanok
Base-cycle and common-path counts in the actual marked hub graph imply a uniform positive lower bound for labelled paired absorbers. Common paths are required only for distinct roots of the same type. The constants precede the host, permutation, separation parameter and root labels.
\source{lem:local:marked-hub-absorbers}
\end{lemma}
\begin{proof}\leanok
Choose an ordered base avoiding all prescribed roots. For each position, choose a common replacement path avoiding all earlier labels. Bound the excluded tuples by prescribing a forbidden label and its position. Keep each free label and each cycle identifier; these choices inject into the labelled absorber records. Both root-replacement cycles use the supplied spoke conditions and the original ambient permutation.
\end{proof}

\begin{lemma}[Profile factors on embedded classes]
\label{bp:embedded-profile-factor}
\lean{Oberwolfach.all_supported_absorber_family_card,Oberwolfach.supported_profile_factor_extraction,Oberwolfach.marked_hub_factor}
\uses{bp:profile-factor-extraction,bp:marked-hub-paired-absorbers}
\leanok
Profile factor extraction applies to classes embedded in an ambient vertex set. All resource identifiers and coefficients are retained. For a marked hub this gives a factor on the union of its two spoke neighbourhoods, without assuming that this union is invariant under the ambient permutation.
\end{lemma}
\begin{proof}\leanok
Pull back only the vertex supports along the class embedding. Every free absorber label lies on a root-free resource, so its support also pulls back. The original coefficients have the same incident sums on the corresponding vertices. Recover ordered root labels and bound the multiplicity of cycle identifiers with a given vertex set before applying factor extraction.
\end{proof}

\begin{lemma}[Marked one-hub factor and physical wheels]
\label{bp:marked-hub-factor-physical}
\lean{Oberwolfach.marked_hub_factor_and_wheels,Oberwolfach.marked_factor_physical_edges_disjoint,Oberwolfach.marked_factor_physical_spoke_once,Oberwolfach.marked_factor_physical_degree}
\uses{bp:embedded-profile-factor}
\leanok
Under the manuscript's hub balance, ordinary--successor separation, physical-map injectivity, actual cycle/path counts and positive almost-unit vertex weights, there is a marked hub factor. Its wheels cover each incident spoke exactly once, their physical rim edges are distinct, and the graph formed by these rim edges has maximum degree two.
\source{lem:completion:marked-hub-factor}
\end{lemma}
\begin{proof}\leanok
Let the marked spoke neighbourhood have size $k$ and the ordinary one size $mk$. Convert the supplied bounds $b/n^m\leq w\leq B/n^m$ to $b\beta^m/k^m\leq w\leq B/k^m$ using $\beta n\leq k\leq n$. Apply the embedded-class factor theorem. The jointly injective labels of its cycles give unique spoke coverage. Injectivity of the physical map prevents a repeated rim edge. The twisted head labels are injective because the ordinary neighbourhood avoids the permutation image of the marked neighbourhood; injectivity of tails and heads gives degree at most two.
\end{proof}


\begin{lemma}[Paired absorbers in ordinary local cycle families]
\label{bp:ordinary-local-absorbers}
\lean{Oberwolfach.sparse_greedy_sequence_count,Oberwolfach.ordinary_common_path_cycle,Oberwolfach.ordinary_local_paired_absorber_count}
\uses{bp:marked-hub-paired-absorbers}
\leanok
The supplied positive common-neighbourhood estimates in the sparse graph construct paired labelled absorbers for its conflict-free ordinary cycles. For triangles, three common neighbours are required; for longer cycles, two suffice. Prescribed roots need not be mutually nonconflicting.
\end{lemma}
\begin{proof}\leanok
Count actual injective embeddings, excluding used vertices and their conflict neighbours at each step. Choose the two path endpoints before its interior. This order gives the required common paths without demanding dense minimum degree in the sparse graph. Close each path separately with its corresponding root and retain every cycle identifier in the paired template.
\end{proof}

\begin{lemma}[Ordinary local factors from relative neighbourhood estimates]
\label{bp:ordinary-relative-factor}
\lean{Oberwolfach.single_class_factor_extraction,Oberwolfach.ordinary_sparse_local_factor,Oberwolfach.ordinary_relative_local_factor}
\uses{bp:ordinary-local-absorbers,bp:embedded-profile-factor}
\leanok
Actual relative common-neighbourhood bounds, positive cycle weights with almost-unit incident sums and divisibility yield an ordinary cycle factor. The constants precede the graphs, densities, spoke count and error parameters. Forbidden-neighbourhood deletion is deducted before applying the supplied density product.
\end{lemma}
\begin{proof}\leanok
Derive linear common-neighbourhood bounds from the dense parent and the supplied relative estimates. Construct the paired absorbers, then apply one-class factor extraction. Convert ambient-scale coefficients to the scale of the actual local vertex set using its proved linear size bound.
\end{proof}

\begin{lemma}[The ordinary factor with the original fixed coefficients]
\label{bp:ordinary-fixed-weight-factor}
\lean{Oberwolfach.ordinaryLocalPotential_present,Oberwolfach.ordinaryLocalTarget_error,Oberwolfach.ordinary_local_coefficient_bounds,Oberwolfach.ordinary_local_factor_from_fixed_weights}
\uses{bp:ordinary-relative-factor,bp:omitted}
\leanok
The factor conclusion of the ordinary local-factor lemma follows from the original fractional cycle decomposition, its prescribed density normalization, near regularity, and the root-omitted potential-load estimates. No absorber-count or factor-existence premise is added.
\source{lem:completion:ordinary-local-factor}
\end{lemma}
\begin{proof}\leanok
Reindex the surviving cycle family without changing its weights or counting an undirected cycle more than once. If the root is present, its root-omitted sum is exactly the local incident sum. The fixed target differs from one by at most $2\delta+12r^4c_+L/n$. Bound each normalized weight above and below, add the supplied potential error, and apply the relative local-factor theorem. This local result does not prove the separate minimum-degree-only cycle-factor theorem used by the terminal transformer construction.
\end{proof}


\begin{lemma}[Endpoint triangles and their paired absorbers]
\label{bp:endpoint-paired-absorbers}
\lean{Oberwolfach.endpoint_triangle_from_vertices,Oberwolfach.endpoint_base_count,Oberwolfach.endpoint_path_count,Oberwolfach.endpoint_path_triangle,Oberwolfach.endpoint_paired_absorber_count,Oberwolfach.endpoint_triangle_factor}
\uses{bp:ordinary-local-absorbers,bp:embedded-profile-factor}
\leanok
Endpoint triangles are unordered sets with two row vertices and one index. Their actual row and index common-neighbourhood estimates construct paired absorbers and, with positive almost-unit incident weights, an exact factor when the row class is twice the index class.
\end{lemma}
\begin{proof}\leanok
Choose the base triangle outside the conflict neighbourhoods of all prescribed roots. For an index replacement choose two rows; for a row replacement choose one index and one row. Only the latter choice uses a common-index estimate for two rows, and that pair consists of one root and its nonconflicting base vertex. The two original row roots are never put together in this test. Inject the resulting nine-free-vertex constructions into the paired template, retaining actual triangle identifiers, and apply the two-class factor theorem.
\end{proof}

\begin{lemma}[Endpoint factors from the source neighbourhood inequalities]
\label{bp:endpoint-source-factor}
\lean{Oberwolfach.endpointHostGraph_row_common,Oberwolfach.endpointHostGraph_index_common,Oberwolfach.common_neighbour_deletion_lower,Oberwolfach.endpoint_relative_row_lower,Oberwolfach.endpoint_source_counts,Oberwolfach.endpoint_source_factor}
\uses{bp:endpoint-paired-absorbers}
\leanok
The source row and index inequalities imply the candidate counts required above. Row edges and spokes are represented in an actual auxiliary graph. External roots, original row-domain normalization and the singleton-or-three-row scope of the index estimates are retained. At most three times as many ambient labels as row vertices suffice for the scale conversion.
\end{lemma}
\begin{proof}\leanok
Identify the auxiliary common neighbourhoods with the original row and spoke intersections. Intersect at most three parent neighbourhoods and subtract the actual conflict-degree loss before multiplying by the density product. Use the index estimate only on the root singleton and on the root together with one nonconflicting row pair. Convert positive cycle coefficients to the ambient scale and invoke the proved endpoint factor theorem. This intermediate result assumes almost-unit incident triangle sums. Node~\ref{bp:endpoint-fixed-weight-factor} constructs these weights and proves their bounds from the supplied parent matching and potential-load estimates. The random-event application remains a separate caller obligation.
\end{proof}

\begin{theorem}[Minimum-degree triangle factors above three quarters]
\label{bp:minimum-degree-triangle-factor}
\lean{Oberwolfach.dense_bipartite_perfect_matching,Oberwolfach.tripartite_triangle_bijections,Oberwolfach.minimum_degree_three_partition,Oberwolfach.triangle_factor_from_bijections,Oberwolfach.minimum_degree_triangle_factor,Oberwolfach.bounded_augmentation_triangle_factor}
\uses{bp:fixed-size-concentration}
\leanok
For every $\alpha>0$, every sufficiently large graph of order divisible by three and minimum degree at least $(3/4+\alpha)n$ has a triangle factor. No regularity or fractional-decomposition assumption is needed.
\end{theorem}
\begin{proof}\leanok
A uniform balanced three-part partition retains degree at least three quarters of each part at every vertex. Hall gives a perfect matching between the first two classes. The bipartite graph between these matched pairs and the third class, requiring adjacency to both endpoints, has minimum degree at least half on both sides. A second application of Hall gives the triangles; their physical subgraphs are disjoint on vertices and cover the original graph. The bounded-augmentation corollary checks the stronger sample-size margin needed to use this result in the triangular transformer. The sharper $2/3$ triangle-factor threshold and the minimum-degree-only inputs for other cycle lengths are not proved here.
\end{proof}


\begin{lemma}[The sampled pair of factors in the triangle transformer]
\label{bp:triangle-transformer-sample-factors}
\lean{Oberwolfach.ordinaryFactorGraph_degree_le_two,Oberwolfach.ordinary_factor_lift,Oberwolfach.ordinaryCycleLift_disjoint_edges,Oberwolfach.two_triangle_factors_at_hub,Oberwolfach.triangle_transformer_sample,Oberwolfach.triangle_transformer_sample_factors}
\uses{bp:minimum-degree-triangle-factor}
\leanok
For a prescribed bounded set $Z$ of order divisible by three, a sufficiently large bounded set $S$ can be chosen in the available common spoke neighbourhood so that $G[S]$ and $G[S\cup Z]$ have edge-disjoint triangle factors. The second factor has no edge within $Z$. The hypotheses are the actual degrees into the available set, including degrees at external roots.
\end{lemma}
\begin{proof}\leanok
The rooted uniform-subset argument supplies a fresh set and simultaneous degree bounds at its vertices and at all prescribed roots. The graph of the first factor has maximum degree two. Delete its edges and all edges inside $Z$, and apply the proved minimum-degree triangle-factor theorem with the checked bounded-augmentation margin. Lift both families to physical subgraphs of the original graph, preserving exact vertex partitions and edge-disjointness. This proves the two-factor step, but not the remaining wheel identities or the full triangular transformer.
\end{proof}

\begin{lemma}[Endpoint matching-weight targets on the original domain]
\label{bp:endpoint-matching-targets}
\lean{Oberwolfach.endpoint_matching_row_target,Oberwolfach.endpoint_matching_total_target,Oberwolfach.weighted_graph_handshake,Oberwolfach.endpoint_eligible_ordered_mass,Oberwolfach.endpoint_unordered_matching_mass_error}
\uses{bp:endpoint-source-factor}
\leanok
With unchanged positive perfect fractional matching weights bounded above by $c_+/a$, conflict degree at most $L$ changes an eligible endpoint sum by at most $2c_+L/a$. The total unordered eligible-edge weight differs from $m_x/2$ by at most $(1+2c_+)L/2$.
\end{lemma}
\begin{proof}\leanok
For a fixed endpoint, charge only the conflict neighbours of it and the external root. Sum the row losses, and separately charge the removed root-neighbour rows using their original sum one. Weighted incidence double counting identifies half the ordered mass with the sum over genuine unordered eligible edges. The density normalization and potential-sum identification are proved in Node~\ref{bp:endpoint-fixed-weight-factor}. The common random-event application remains a separate obligation.
\end{proof}


\begin{lemma}[Endpoint factors from the unchanged parent matching]
\label{bp:endpoint-fixed-weight-factor}
\lean{Oberwolfach.endpoint_triangle_coefficient_bounds,Oberwolfach.endpoint_index_target_error,Oberwolfach.endpoint_triangle_row_load,Oberwolfach.endpoint_triangle_index_load,Oberwolfach.endpoint_row_potential_present,Oberwolfach.endpoint_index_potential_present,Oberwolfach.endpoint_fixed_weight_estimates,Oberwolfach.graphEdgeWeightExtension_incident,Oberwolfach.endpoint_parent_matching_extension,Oberwolfach.endpoint_row_target_as_edges,Oberwolfach.endpoint_source_domain_card,Oberwolfach.endpoint_conflict_no_index,Oberwolfach.endpoint_factor_from_parent_matching}
\uses{bp:endpoint-matching-targets,bp:endpoint-source-factor}
\leanok
The supplied perfect fractional matching on the original parent induced graph, the row/index estimates and the root-omitted potential bounds imply the endpoint-triangle factor conclusion. The weights are divided by the fixed denominator $bq^2s^3$, without recomputing the matching after sampling.
\end{lemma}
\begin{proof}
\leanok
The row and index coordinate maps are bijections with the actual incident triangle families. They preserve the omitted root edge and count each unordered row edge once. Zero extension of the parent edge weights preserves all incident sums. The coefficient bounds hold at scale $a^{-2}$, and the incident error is at most $\delta+(2c_++(1+2c_+)/(2\kappa p_*))L/a+\xi$. The source domain gives the ambient-size bound, and the row-supported conflict graph has no index incidences. Apply the proved endpoint factor theorem. This closes the deterministic local-factor implication; it does not assert the potential hypotheses hold in the later common sampling event.
\end{proof}

\begin{lemma}[Ordinary local extensions with external roots]
\label{bp:ordinary-local-extensions}
\lean{Oberwolfach.ordinaryExtensionRootBound_threshold,Oberwolfach.extension_parent_common_lower,Oberwolfach.extension_relative_row_lower,Oberwolfach.extension_relative_index_lower,Oberwolfach.ordinary_local_extensions}
\uses{bp:threshold-inputs}
\leanok
The original row and index neighbourhood hypotheses imply positive linear lower bounds for every ordinary extension set. The coefficients and error tolerance precede the permutation-conflict radius $d$; only the minimum order depends on $d$.
\end{lemma}
\begin{proof}
\leanok
For $k_3=4$, $k_4=3$ and $k_r=2$ otherwise, the proved threshold lower bounds give $1-1/k_r\le\delta_{C_r}$. Intersect the actual parent neighbourhoods, subtract at most $6d|S|$ conflict neighbours, and only then multiply by the density product. The index conclusion follows directly from its fixed separated-root hypothesis. All roots range over the ambient set, including roots outside the sampled domain.
\end{proof}


\begin{lemma}[Marked local extensions on the original common event]
\label{bp:marked-local-extensions}
\lean{Oberwolfach.marked_two_extension_physical_injective,Oberwolfach.marked_local_extensions,Oberwolfach.marked_rows_supply_extension_patterns,Oberwolfach.marked_row_event_extension_sets,Oberwolfach.marked_index_event_extension_lower,Oberwolfach.marked_extension_marginal_lower,Oberwolfach.marked_common_event_local_extensions}
\uses{bp:marked-row-patterns}
\leanok
The original marked row and index estimates imply the one-edge, two-edge and mixed-index extension bounds on the same event. All endpoint roots stay in the fixed ambient domain. The coefficients and error tolerance precede the conflict radius; only the order threshold depends on that radius.
\end{lemma}
\begin{proof}\leanok
Represent each required ordinary or marked edge by its actual directed request. The physical head of a marked request is the prescribed permutation of its index head. The hypotheses exclude the physical collisions in two-edge requests. Use the original row-domain size in the error, and subtract the realised conflict-neighbourhood loss before applying the density product. Bound the supplied hub-dependent marginals from below; do not replace them with a new constant-density law. The row and index conclusions are obtained from the same realization, uniformly for the later root choices. This proves the extension step, not the later path construction or full cover-down theorem.
\end{proof}


\begin{lemma}[Ordinary local factors after realised rim deletions]
\label{bp:ordinary-robust-local-factor}
\lean{Oberwolfach.ordinaryLocalPotential_sdiff_loss,Oberwolfach.ordinaryLocalPotential_sdiff_bound,Oberwolfach.ordinary_common_sdiff_bound,Oberwolfach.ordinary_robust_local_factor}
\uses{bp:ordinary-relative-factor,bp:ordinary-fixed-weight-factor}
\leanok
The local factor remains available after every sufficiently small maximum-degree rim deletion. The constant controlling the deletion is chosen before the graph and the deletion. Parent weights, root domains and original density products remain fixed.
\end{lemma}
\begin{proof}\leanok
If the deleted graph has maximum degree $L$, a common-neighbourhood count with root set $S$ loses at most $|S|L$. Fixed cycle coefficients bounded by $B/n^{r-1}$ give potential-load loss at most $3r^4BL/n$ at every ambient root. These are realised losses with no additional sampling-density multiplier. Choose the deletion tolerance using the already proved local-factor tolerances and apply the relative factor theorem to the remaining graph.
\end{proof}

\begin{theorem}[Successive ordinary hub factors and physical wheels]
\label{bp:ordinary-hub-cover}
\lean{Oberwolfach.ordinary_factor_wheels,Oberwolfach.ordinary_factor_wheels_avoiding,Oberwolfach.ordinary_wheel_packings_disjoint,Oberwolfach.sequential_ordinary_hub_wheels,Oberwolfach.ordinary_hub_cover_from_fixed_weights}
\uses{bp:ordinary-robust-local-factor}
\leanok
For the prescribed distinct hubs, the fixed-weight local hypotheses give one physical wheel packing covering each prescribed spoke exactly once. Its rim graph has maximum degree at most twice the number of hubs, provided that this quantity lies within the proved deletion allowance. Forbidden pairs remain excluded from every rim cycle.
\end{theorem}
\begin{proof}\leanok
Construct each factor in the graph obtained by removing all previously selected rims. The graph degree bound proves the next local-factor input. Lift the resulting wheels back to the original host without altering any physical resources. Disjoint rims and distinct hubs imply disjoint physical wheel edge sets. Their exact union gives the spoke partition. The final theorem supplies the factors from the original fractional cycle weights and potential targets; the intermediate sequential selection lemma's factor premise is discharged.
\source{clm:protected:deep:ordinary}
This is the hub portion of the deep-pair cover-down claim. The row-root construction, boundary cover-down and protected output are not proved by this theorem.
\end{proof}

\begin{lemma}[Exact ordinary hub remainder]
\label{bp:ordinary-hub-cover-remainder}
\lean{Oberwolfach.spoke_row_sdiff_card,Oberwolfach.spoke_col_sdiff_card,Oberwolfach.ordinary_decomposition_remainder_balance,Oberwolfach.ordinary_hub_packing_decomposition,Oberwolfach.ordinary_hub_cover_remainder}
\uses{bp:ordinary-hub-cover}
\leanok
When the local factors cover the entire original neighbourhood of each processed hub, their union removes all and only those spoke columns. Its deletion from an ordinary wheel-divisible host leaves an ordinary wheel-divisible extended graph.
\end{lemma}
\begin{proof}\leanok
Identify the union with the actual rim graph and spoke predicate. Use the physical wheel decomposition to derive the exact row balance and hub congruences. Subtract the contained rim graph degrees and spoke sets. The resulting edge-set difference is exactly the packing remainder, not merely an equality of total edge counts.
\end{proof}


\begin{lemma}[Physical triangle wheels at a row root]
\label{bp:endpoint-physical-root}
\lean{Oberwolfach.endpoint_triangle_root_wheel,Oberwolfach.endpoint_factor_root_wheel_packing,Oberwolfach.endpoint_factor_completed_wheels,Oberwolfach.endpoint_root_wheels_from_parent_matching}
\uses{bp:endpoint-fixed-weight-factor}
\leanok
The endpoint triangle factor gives actual ordinary $3$-wheels, even though all their rims contain the row root. Every prescribed root rim edge and root spoke is covered exactly once. Rim degree outside the root is at most two, spoke row degree outside the root is at most one, and each hub has spoke degree at most three.
\end{lemma}
\begin{proof}\leanok
Identify the exact rim support as the root together with the two row vertices of the endpoint triangle. A shared rim edge would force a shared nonroot row vertex; a shared spoke would force a shared index. Both contradict the endpoint factor's vertex partition. Retain the exact support and hub identities for the later comparison-host argument. The parent matching and fixed potential estimates supply the factor through the proved local theorem. The later common-event and restriction substitutions are still separate obligations.
\end{proof}

\begin{lemma}[Assigning forbidden-pair-free neighbour pairs]
\label{bp:ordinary-neighbour-pairs}
\lean{Oberwolfach.degree_cut_perfect_matching,Oberwolfach.degree_cut_matching_after_deletion,Oberwolfach.bipartite_balanced_split,Oberwolfach.ordinary_neighbour_pair_assignment}
\uses{bp:fixed-size-concentration}
\leanok
For a bipartite graph with parts of size $2m,m$, positive linear minimum degrees and the source lower bounds on all sufficiently large cuts, assign two distinct neighbours to each column. The pairs partition the row part and avoid every edge of a fixed bounded-degree forbidden graph. The cut threshold precedes the forbidden-degree bound; only the order threshold depends on that bound.
\end{lemma}
\begin{proof}\leanok
A uniform balanced split retains all column degree bounds simultaneously. Hall's condition follows from minimum degrees for the small sets and from the large-cut condition for a failed intermediate cut. After the first perfect matching, remove the forbidden partners of its endpoints. Injectivity bounds the deletion on both sides. The checked cut and degree margins give the second Hall matching.
\source{lem:pr:ordinary-endpoints}
\end{proof}

\begin{lemma}[Fixed paths and physical root wheels in a sparse host]
\label{bp:sparse-root-wheels}
\lean{Oberwolfach.sparse_common_neighbour_avoiding,Oberwolfach.sparse_fixed_length_forbidden_path_within,Oberwolfach.embedded_cycle_from_path_with_support,Oberwolfach.ordinary_root_wheel_from_common_counts_within,Oberwolfach.two_common_neighbours_after_deletion,Oberwolfach.ordinary_root_wheel_after_local_deletion,Oberwolfach.ordinary_root_wheels_from_local_extensions}
\uses{bp:ordinary-local-extensions}
\leanok
Actual common-neighbour counts give a simple path of the required length, avoiding the prescribed vertices and every forbidden pair. Closing the path through a row root produces a physical wheel with its prescribed two root rim edges and root spoke. The root can be excluded from the queried domain, so no small deletion degree at that root is assumed.
\end{lemma}
\begin{proof}\leanok
Subtract the realised excluded set and all forbidden neighbours of used vertices. Close the simple path using a final common neighbour. The cycle map retains exact support and all consecutive path edges. Deleting a graph $D$ costs at most $d_D(a)+d_D(b)$ in a two-root common neighbourhood, using only the two queried roots. The local-extension wrapper preserves the parameter order and supplies the sparse count bounds from the original parent hypotheses.
\end{proof}

\begin{theorem}[Capped ordinary wheels at one row root]
\label{bp:capped-root-wheels}
\lean{Oberwolfach.rootWheelInternalLoad_total,Oberwolfach.rootWheelInternal_saturated,Oberwolfach.rootWheelFamily_degree,Oberwolfach.root_wheel_family_exact_edges,Oberwolfach.capped_ordinary_root_wheels,Oberwolfach.root_wheel_cap_parameters,Oberwolfach.ordinary_capped_root_wheels_from_local_extensions}
\uses{bp:sparse-root-wheels,bp:ordinary-neighbour-pairs}
\leanok
For distinct assigned hubs and disjoint prescribed endpoint pairs, the sparse common-neighbour hypotheses construct one physical wheel packing through the row root. With internal-use cap $K$, its rim degree outside the root is at most $2(K+1)$ and spoke row degree at most $K+1$; each hub has degree at most $r$. The root rim and spoke sets are exactly the assigned pairs and hubs.
\end{theorem}
\begin{proof}\leanok
At most $Nr/K$ vertices have reached the cap during $N$ wheel choices. A vertex other than the root occurs as a prescribed endpoint at most once, giving deleted rim degree at most $2(K+1)$. The next path avoids saturated vertices and all previous rims. The explicit inequality $Nr/K+(r+3)(L+1)+4(K+1)<b$ supplies a choice. The exact two neighbours of the root in every previous rim ensure that future prescribed root edges survive. Distinct hubs separate the spokes. Induction proves the cap and all physical edge-disjointness assertions. The cap and common-neighbour tolerance are chosen before the conflict radius. The local-extension application proves the numerical inequality from the original parent hypotheses. The deterministic successive-root construction is proved in the later row-cover theorem. Its common-event substitutions and the complete protected cover-down conclusion remain open.
\source{clm:protected:deep:ordinary}
\end{proof}

\begin{lemma}[Endpoint factors inside an arbitrary ambient graph]
\label{bp:endpoint-domain-transport}
\lean{Oberwolfach.embedding_comap_common,Oberwolfach.endpoint_counts_embedding,Oberwolfach.endpointTriangleLift_bijective,Oberwolfach.endpointTriangleLift_weighted_degree,Oberwolfach.endpointTriangleLift_decomposition,Oberwolfach.endpoint_triangle_factor_on_domain}
\uses{bp:endpoint-fixed-weight-factor}
\leanok
An endpoint triangle factor may be constructed on the union of its row and index classes inside a larger graph. Restriction preserves all internal common-neighbourhood counts, each triangle coefficient, and every weighted degree exactly; lifting gives the original physical vertex partition.
\end{lemma}
\begin{proof}\leanok
Pull the two classes back through the subtype embedding of their union. The image map is a bijection on endpoint triangles. Reindex the supplied weighted sums through that bijection, apply the proved factor theorem on the finite domain, and lift its disjoint vertex sets. The degree of the restricted forbidden graph is at most its ambient degree. The source factor and all three direct wrappers now use this construction and no longer require the entire ambient vertex set to equal the local domain. The original potential sums, external row root, fixed matching and normalization are unchanged.
\end{proof}

\begin{theorem}[Endpoint factors after realised bounded-degree deletions]
\label{bp:endpoint-factor-deletion}
\lean{Oberwolfach.endpoint_row_removed_count,Oberwolfach.endpoint_index_removed_count,Oberwolfach.endpoint_potential_deletion_bounds,Oberwolfach.endpoint_potential_errors_after_deletion,Oberwolfach.mixed_common_sdiff_bound,Oberwolfach.endpoint_factor_after_bounded_deletion}
\uses{bp:endpoint-domain-transport,bp:endpoint-fixed-weight-factor}
\leanok
Write $a=|A|$ and suppose $|B|\le a$. If every normalised fixed coefficient is at most $C/a^2$, deleting rim and spoke graphs of maximum degree at most $L$ changes a row potential by at most $5CL/a$ and an index potential by at most $3CL/a$. These estimates hold at every tested root, including an absent root edge or spoke. The exact endpoint-factor theorem remains applicable with the explicitly increased neighbourhood and potential errors and with its exact final class-size relation.
\end{theorem}
\begin{proof}\leanok
A row candidate is lost only through one of its two tested rim edges or three tested spokes. Count these realised losses in $B\times A$. For an index candidate, charge its internal rim edge or an endpoint in the deleted neighbourhood of the row root or hub. The parent graph, forbidden graph, eligible edge family, weights and denominator stay fixed. Mixed neighbourhoods lose at most $(|X|+|J|)L$ vertices; index neighbourhoods lose at most $|X|L$. The caller uses the original parent matching and error targets and constructs an actual endpoint triangle factor after deletion. No separate event for the later chosen deletion is introduced. Supplying the original estimates simultaneously from the manuscript's sampling experiment remains a separate caller obligation.
\end{proof}

\begin{theorem}[Hub covers with an incidence degree bound]
\label{bp:hub-cover-incidence}
\lean{Oberwolfach.ordinary_hub_packing_rim_support,Oberwolfach.graph_family_degree_by_support,Oberwolfach.sequential_ordinary_hub_wheels_by_incidence,Oberwolfach.ordinary_hub_cover_from_fixed_weights_by_incidence}
\uses{bp:ordinary-hub-cover,bp:ordinary-robust-local-factor}
\leanok
For distinct hubs with prescribed neighbourhoods $T_i$, the successive factor construction has rim degree at a vertex $v$ at most $2|\{i:v\in T_i\}|$. It suffices that this pointwise bound lies within the local deletion tolerance. The original parent-cycle coefficients, targets and fractional decomposition supply every factor.
\end{theorem}
\begin{proof}\leanok
The exact spoke union at a single hub forces the corresponding rim support to lie in $T_i$. Summing its degree bound two only over those $i$ with $v\in T_i$ gives the stated estimate at every partial stage. The proved robust local factor theorem supplies the next factor after that actual rim deletion. This extends the earlier few-hubs conclusion; it does not identify a gap in its original small-hub-set caller.
\end{proof}

\begin{theorem}[Successive ordinary row-root covers]
\label{bp:ordinary-successive-row-covers}
\lean{Oberwolfach.ordinary_neighbour_pairs_on_domains,Oberwolfach.ordinary_row_root_cover,Oberwolfach.bipartite_cut_deleted_column,Oberwolfach.two_rim_one_spoke_deleted,Oberwolfach.ordinary_row_root_after_deletion,Oberwolfach.ordinary_wheel_comparison_disjoint,Oberwolfach.ordinary_row_covers_from_fixed_counts,Oberwolfach.successive_root_parameters,Oberwolfach.ordinary_small_set_row_covers}
\uses{bp:capped-root-wheels,bp:ordinary-neighbour-pairs}
\leanok
For $r\ge4$, let the distinct row roots lie outside a fixed set $A$, with all their rim neighbours in $A$. Exact row balance and the stated positive degree, large-cut and two-rim-one-spoke common-neighbour bounds give edge-disjoint physical wheel packings covering every rim edge and spoke at every root. For $h$ roots the rim degree on $A$ is at most $2h(K+1)$, the spoke row degree on $A$ at most $h(K+1)$, and the spoke column degree at most $hr$. The cap $K$ and permitted relative number of roots are chosen before the forbidden-degree bound and the final order threshold.
\end{theorem}
\begin{proof}\leanok
The Hall construction assigns the actual endpoint pairs at each root. Deleting spoke edges with column degree at most $L$ loses at most $L|J|$ edges in a cut with column set $J$. A two-rim-one-spoke neighbourhood loses at most $3L$ vertices. These bounds supply every hypothesis of the capped single-root construction after the previously selected edges are deleted on $A$. Edges at earlier roots are restored only in this comparison: the new wheel's support is contained in $A$ together with its own root, and the formal physical-edge argument proves disjointness from every previous wheel. Summing the fixed per-root cap gives the bounds throughout the iteration. The final theorem proves the numerical margins from a sufficiently small relative root-set size, and discharges the intermediate selection lemma's existence premise.
\source{clm:protected:deep:ordinary}
This establishes the deterministic ordinary row-root construction for $r\ge4$ under explicit fixed count and cut bounds. Deriving those inputs from the manuscript's actual common event, combining the edge-prescription and hub stages, the triangle and marked counterparts, and the protected cover-down conclusion remain separate obligations.
\end{proof}

\begin{lemma}[Exact union and remainder of ordinary row covers]
\label{bp:ordinary-row-cover-remainder}
\lean{Oberwolfach.ordinaryPacking_resources,Oberwolfach.ordinaryPacking_decomposition,Oberwolfach.row_cover_union_packing,Oberwolfach.row_cover_union_degree_bounds,Oberwolfach.row_cover_remainder}
\uses{bp:ordinary-successive-row-covers,bp:ordinary-hub-cover-remainder}
\leanok
The union of the successive row-root packings is an actual physical wheel decomposition of its rim graph and spoke relation. It covers every original edge at each processed root. Removing it from an ordinary wheel-divisible host preserves the exact row degree equalities and hub congruences and leaves no edge at any processed root.
\end{lemma}
\begin{proof}\leanok
Within-root and between-root physical disjointness give one packing. Identify its two resource types exactly with the rim union and spoke predicate. The generic wheel balance theorem and exact degree/set subtraction give the remainder identities. No multiplicity is discarded and no approximate degree identity replaces the exact divisibility statement.
\end{proof}

\begin{lemma}[Endpoint cuts from degrees and separated codegrees]
\label{bp:ordinary-codegree-cuts}
\lean{Oberwolfach.degree_codegree_variance,Oberwolfach.degree_codegree_cut_square,Oberwolfach.codegree_cut_numeric,Oberwolfach.bipartite_cuts_from_degrees_and_codegrees,Oberwolfach.degree_codegree_cuts_on_domains,Oberwolfach.ordinary_neighbour_pairs_from_codegrees,Oberwolfach.ordinary_small_set_row_covers_from_codegrees}
\uses{bp:ordinary-neighbour-pairs,bp:ordinary-successive-row-covers}
\leanok
Minimum row degrees and upper bounds on nonexceptional row codegrees imply the required cut bounds. If every row has at most $d$ exceptional partners, the variance calculation includes $d+1$ to account for the diagonal pairs. The resulting endpoint assignments and successive ordinary row covers no longer assume a cut bound separately.
\end{lemma}
\begin{proof}\leanok
Expand the square of the deviation of the column degrees from $p|S|$. Bound the contribution of the exceptional pairs explicitly. The variance is at most $3\varepsilon|S|^2|B|+(d+1)|S||B|$. Cauchy--Schwarz converts this to a cut bound. The error tolerance depends only on a fixed positive density lower bound; the order threshold is chosen after $d$. Pullbacks to the actual finite domains preserve the sums and cardinalities exactly.
\end{proof}

\begin{theorem}[Successive row covers from relative neighbourhood estimates]
\label{bp:ordinary-relative-row-covers}
\lean{Oberwolfach.ordinary_endpoint_counts_from_common,Oberwolfach.endpoint_spoke_normalisation,Oberwolfach.endpoint_column_normalisation,Oberwolfach.ordinary_endpoint_counts_normalised,Oberwolfach.ordinary_small_set_row_covers_from_relative_counts}
\uses{bp:ordinary-codegree-cuts,bp:ordinary-local-extensions}
\leanok
For $r\ge4$, the original relative common-neighbourhood estimates, the eligible parent degree bound, exact row balance and a sufficiently small set of row roots give the successive physical row covers. The cap, root-set tolerance and count-error tolerance precede the forbidden degree.
\end{theorem}
\begin{proof}\leanok
Apply the index estimates to the independent root sets $\{x\}$, $\{x,v\}$ and $\{x,v,w\}$ and the row estimates to one rim and one spoke condition. Retain the original error scales $|A|$ and $|B|$ when passing to the actual endpoint parts. The positive parent degree margin gives the two-rim-one-spoke extension bounds. Invoke the checked degree/codegree and successive-row constructions, thereby supplying all their existence and count premises.
\end{proof}

\begin{theorem}[One sampling event for later ordinary row covers]
\label{bp:ordinary-adaptive-row-event}
\lean{Oberwolfach.rawOrdinaryPatternSum_expectation,Oberwolfach.rawOrdinaryPatternSum_tail,Oberwolfach.ordinaryRowResources_incidence,Oberwolfach.ordinaryRowResources_disjoint,Oberwolfach.rawOrdinaryRowCount_tail,Oberwolfach.rawOrdinaryRows_common_event,Oberwolfach.rawOrdinaryIndices_common_event,Oberwolfach.ordinary_neighbourhood_common_event,Oberwolfach.mixed_neighbourhood_stability,Oberwolfach.ordinary_relative_counts_under_changes,Oberwolfach.ordinary_adaptive_row_cover_event}
\uses{bp:ordinary-raw-sampling,bp:raw-interval-law,bp:ordinary-relative-row-covers,bp:neighbourhood-stability}
\leanok
Under the original independent rim and interval-spoke law, one event of probability at least $1-|A|^{-q}$ supplies every bounded row and separated index count on the fixed outside domain $A$. On that same event, every subsequent bounded modification with balanced row roots and the prescribed support admits the successive ordinary row covers. The modified graphs and the root list may be chosen after the sample; the outside domain and the density targets are fixed before it.
\end{theorem}
\begin{proof}\leanok
Represent each row test by its actual physical rim and spoke resources. Distinct rim resources have distinct coordinates; spokes at a common row and distinct hubs have disjoint coordinate supports. A physical rim occurs in at most two candidate positions. These statements give the exact mean and the concentration bound under the unchanged law. Enumerate every bounded test before sampling, including empty tests, and combine the row and independent-index events. Changing each tested neighbourhood by at most $L$ changes a row count by at most $(|S|+|J|)L$ and an index count by at most $|S|L$. The universally quantified deterministic implication applies to all later modifications on the same event and supplies the row covers. There is no density factor on the realised change cost.

This completes this ordinary $r\ge4$ row stage on fixed domains. Preparation must still supply its actual eligible parent, domain, degree and balance hypotheses, and combine it with prescribed-edge and hub covers. It does not prove the triangle or marked successive-row stages, protected cover-down, terminal absorption, or Theorem~1.1.
\end{proof}

\begin{theorem}[Successive triangle row covers with the original weights]
\label{bp:triangle-successive-row-covers}
\lean{Oberwolfach.bounded_fractional_matching_on_subset,Oberwolfach.endpoint_parent_matching_exists,Oberwolfach.endpoint_factor_row_cover,Oberwolfach.triangle_comparison_cover_of_factor,Oberwolfach.triangle_row_cover_parameters,Oberwolfach.triangle_row_covers_from_fixed_weights}
\uses{bp:endpoint-factor-deletion,bp:ordinary-row-cover-remainder}
\leanok
Let $A$ be the original domain for the fixed parent matching weights and let $X\subseteq A$ be the free rim set. For sufficiently few distinct row roots outside $X$, the original count and potential estimates, exact root balance, and the stated support conditions give pairwise edge-disjoint triangle-wheel covers. For $h$ roots their total rim degree on $X$ is at most $2h$, their spoke row degree on $X$ at most $h$, and their spoke column degree at most $3h$. The roots may belong to $A$. All errors remain normalised by $|A|$ and $|B|$.
\end{theorem}
\begin{proof}\leanok
The parent degree assumption supplies a bounded fractional matching on $A\cap N_P(x)$ before sampling: its edge weights lie between $1/(30|A|)$ and $40/|A|$. During the successive construction delete only the previously used rim edges inside $X$ and the previously used spokes between $X$ and $B$. Their maximum degrees are at most $3h$. The proved bounded-deletion endpoint-factor theorem constructs the next factor with the same weights. Its actual triangles become physical wheels. The support argument excludes every old root and proves resource disjointness even though comparison restores edges at those roots. The explicit choice of the root-set tolerance absorbs the count and potential changes. The fixed potential estimates still have to be supplied simultaneously by the preparation sampling caller; they have not been replaced with new random weights.
\end{proof}

\begin{theorem}[A fixed-domain event permits later free-set restrictions]
\label{bp:adaptive-free-domain-row-covers}
\lean{Oberwolfach.ordinary_bounded_neighbourhood_common_event,Oberwolfach.ordinary_scaled_set_row_covers,Oberwolfach.ordinary_scaled_set_row_covers_from_codegrees,Oberwolfach.ordinary_fixed_domain_row_covers,Oberwolfach.ordinary_adaptive_free_domain_row_cover_event}
\uses{bp:ordinary-adaptive-row-event,bp:ordinary-codegree-cuts}
\leanok
The ordinary row-cover common event is uniform over every subsequent free set $X\subseteq A$ with the stated bound on $|A\setminus X|$, as well as the later bounded graph changes and balanced root choices. Here $A$ remains the original fixed sampling domain. The order and error scales are $|A|$, not $|X|$.
\end{theorem}
\begin{proof}\leanok
The raw law gives one simultaneous event for any prescribed bounded neighbourhood-test order. The deterministic row-cover proof now separates the reference order from its free vertex set. Restricting a two-rim-one-spoke candidate set from $A$ to $X$ loses at most $|A\setminus X|$ vertices. This loss is deducted explicitly from its positive lower bound. The index and endpoint counts continue to use the fixed domains, and the endpoint sets are contained in $X$ by the genuine root-support hypothesis. Consequently no new concentration statement on the possibly outcome-dependent set $X$ is used. This repairs an insufficient interface in the earlier Lean caller; it is not a finding that the manuscript's fixed-domain argument was false.
\end{proof}

\begin{theorem}[Covering prescribed inner edges by actual wheels]
\label{bp:ordinary-prescribed-edge-packing}
\lean{Oberwolfach.ordinary_rim_edge_wheel_from_common_counts,Oberwolfach.ordinary_spoke_wheel_from_common_counts,Oberwolfach.wheelRimLoad_saturated,Oberwolfach.wheelHubLoad_saturated,Oberwolfach.ordinary_rim_edge_wheel_after_deletion,Oberwolfach.ordinary_spoke_wheel_after_deletion,Oberwolfach.prescribed_rim_wheel_packing,Oberwolfach.prescribed_spoke_wheel_packing,Oberwolfach.prescribed_rim_wheel_inner_intersection,Oberwolfach.prescribed_spoke_wheel_inner_intersection,Oberwolfach.prescribed_inner_edge_wheel_packing,Oberwolfach.prescribed_wheel_family_exact_union,Oberwolfach.prescribed_edge_cap_parameters,Oberwolfach.prescribed_edge_extension_counts}
\uses{bp:sparse-root-wheels,bp:ordinary-local-extensions}
\leanok
Explicit positive common-neighbour counts supply a physical ordinary wheel packing through prescribed rim edges and spokes in an inner vertex pair $(I,J)$. Every other rim vertex lies in a disjoint free set $X$. A rim-edge wheel has its hub in a free hub set disjoint from $J$; a spoke wheel keeps its prescribed hub in $J$. Every wheel meets the inner host in exactly its designated edge. The union is an actual wheel decomposition of its used edges and intersects the inner host in exactly the prescribed edge set.
\end{theorem}
\begin{proof}\leanok
Count the number of selected wheels containing each rim vertex and the number using each hub. After $N$ selections, at most $Nr/K$ rim vertices and $N/K$ hubs have reached cap $K$. The prescribed endpoints contribute at most their supplied incidence bound $D$. The rim-edge construction therefore has deleted rim degree at most $2(K+D)$, spoke row degree at most $K+D$, and spoke column degree at most $Kr$. The spoke construction permits repeated prescribed hubs, bounded by $D$, and has column degree at most $Dr$. Avoid the actual previously used edges and spokes at each choice. The second construction runs after the first rim union has been deleted, with the additional common-neighbour loss explicitly deducted.

The combined packing has rim degree at most $4(K+D)$, spoke row degree at most $2(K+D)$, and spoke column degree at most $(K+D)r$. The identities for the inner intersections are proved on the physical resource sets, not inferred from informal support descriptions. With $K=\lceil\alpha m\rceil$, $N\le\alpha^2m^2$ and $D\le\alpha m$, the numerical theorem chooses $\alpha$ before the forbidden degree and the final order threshold and verifies both stages' budgets. Restricting the original common-neighbour counts to the later free sets loses at most the sizes of the deleted vertex sets. No sampling density multiplies this realised loss.
\source{clm:protected:deep:ordinary}
These results prove the finite prescribed-edge stage under explicit extension counts. The actual preparation substitutions, its combination with the hub and row stages, the marked counterpart, and the full protected cover-down result remain to be formalised.
\end{proof}


\begin{theorem}[Ordinary row covers on two fixed counting domains]
\label{bp:two-domain-row-cover-event}
\lean{Oberwolfach.common_inter_domain,Oberwolfach.fixed_domain_mixed_neighbourhood_stability,Oberwolfach.fixed_domain_mixed_target_stability,Oberwolfach.ordinary_relative_counts_under_domain_changes,Oberwolfach.ordinary_domain_scoped_row_cover_event,Oberwolfach.rawIntervalCommonHubsOnDomain_common_event,Oberwolfach.rawOrdinaryDomainIndices_common_event,Oberwolfach.ordinary_two_domain_neighbourhood_common_event,Oberwolfach.ordinaryWheelOfHubDomain_injective,Oberwolfach.ordinaryHubDomainPacking_rim,Oberwolfach.ordinaryHubDomainPacking_spokes,Oberwolfach.ordinaryHubDomainPacking_row_card,Oberwolfach.ordinaryHubDomainPacking_column_card,Oberwolfach.ordinaryWheelOfHubDomain_disjoint,Oberwolfach.ordinaryHubDomainPacking_disjoint,Oberwolfach.ordinaryRowCover_of_hub_domain,Oberwolfach.common_hub_domain_card,Oberwolfach.ordinary_two_domain_counts_under_changes,Oberwolfach.ordinary_two_fixed_domain_row_covers,Oberwolfach.ordinary_two_domain_row_cover_event}
\uses{bp:adaptive-free-domain-row-covers,bp:ordinary-adaptive-row-event}
\leanok
Fix the original vertex domain $A$ and hub domain $B$ before the raw interval experiment. The ordinary $r\ge4$ row-cover theorem holds on one event for all subsequent permitted graph changes, free vertex sets $X\subseteq A$, and balanced lists of row roots. Only changes inside $A$ or $B$ enter the corresponding neighbourhood error. The roots may be outside $A$. The resulting wheels have their original ambient vertex and hub endpoints.
\end{theorem}
\begin{proof}\leanok
Intersect each changing neighbourhood with its counting domain before applying the symmetric-difference inequality. This charges at most $|S|L+|J|L$ for the mixed row test and $|S|L$ for the index test. Apply the proved injective-coordinate projection to the fixed hub subset without conditioning on any graph event. The original product law gives simultaneous row and index estimates, and the deterministic stability bound then permits all later modifications on that same event. Apply the row-cover construction over the hub subtype and embed its wheels back into the ambient host. Exact rim and spoke identities preserve the degree bounds, physical disjointness and root coverage. Preparation must still provide the actual domain, minimum-degree and balanced-root inputs.
\end{proof}

\begin{theorem}[The prescribed-edge stage on the original common event]
\label{bp:prescribed-packing-common-event}
\lean{Oberwolfach.prescribed_relative_extension_counts,Oberwolfach.prescribed_inner_packing_from_relative_counts,Oberwolfach.prescribed_inner_packing_common_event}
\uses{bp:ordinary-prescribed-edge-packing,bp:two-domain-row-cover-event}
\leanok
For fixed counting domains $A,U_0$ and positive density and degree margins, the original-law neighbourhood event supplies every subsequent admissible prescribed-inner-edge packing. The forbidden degree is fixed after the permitted prescribed-edge density, and the final order threshold follows these choices. The prescribed edges, protected sets and excluded free sets may depend on the outcome, subject to the explicit support, cardinality and incidence bounds.
\end{theorem}
\begin{proof}\leanok
The parent degree bound yields the two-rim-one-spoke extension lower bound. The original independent-root index estimate yields the common hub bound. Removing candidates outside the later free sets loses at most $|A\setminus X|$ and $|U_0\setminus B|$, without a sampling-density factor. Substitute these bounds into the proved cap construction, with $K=\lceil\alpha|A|\rceil$. The assembled packing has rim degree at most $12\alpha|A|$, spoke row degree at most $6\alpha|A|$ and spoke column degree at most $3r\alpha|A|$. Every wheel meets the protected inner host in exactly its prescribed edge. Both relative-count stability and this construction are deterministic consequences of the one raw event. This supplies the prescribed-edge common-event caller, not the remaining hub/row assembly, marked case, boundary cover-down or terminal absorption.
\end{proof}

\begin{theorem}[The triangle potential event with constructed parent weights]
\label{bp:triangle-parent-potential-event}
\lean{Oberwolfach.endpoint_parent_weight_family,Oberwolfach.raw_endpoint_row_potential_tail,Oberwolfach.raw_endpoint_index_potential_tail,Oberwolfach.endpoint_potential_common_event,Oberwolfach.triangle_sampling_with_parent_weights}
\uses{bp:endpoint-matching-targets,bp:endpoint-fixed-weight-factor,bp:two-domain-row-cover-event}
\leanok
The three-quarter parent degree hypothesis supplies one fixed family of fractional matching weights. For every fixed positive error tolerance, the original raw law has one event of probability at least $1-|A|^{-q}$ on which both endpoint potential sums and all triangle neighbourhood tests have the required errors. The weights precede the density arrays and random outcome; the order threshold is uniform over the supplied arrays.
\end{theorem}
\begin{proof}\leanok
A row potential omits its root rim and tests two other rims and three spokes. An index potential omits its root spoke and tests three rims and two other spokes. Their exact pattern sums have means given by the original eligible row mass and $p_R/(|B|p_H)$ times the original unordered eligible-edge mass, respectively. The fixed normalization is $|B|p_R^2p_H^3$. Each pattern uses at most five resources. Candidate/resource incidence bounds give variance of order $1/|A|$ when $|B|\ge\omega|A|$ and the densities have a fixed positive lower bound. Apply the proved concentration inequality and a union bound over the fixed root pairs. Intersect with the three-root neighbourhood event; the weaker two-root event is insufficient for the triangle caller. Parent weights are supplied by the proved bounded fractional matching theorem, with exact vertex sums and positive lower bounds on every parent edge.
\end{proof}

\begin{theorem}[Conflict restriction and later deletions in the triangle caller]
\label{bp:triangle-conflict-deletion-transport}
\lean{Oberwolfach.endpoint_potentials_conflict_restriction,Oberwolfach.restricted_conflict_has_no_hub_edges,Oberwolfach.endpoint_potentials_conflict_deletion,Oberwolfach.endpoint_potential_estimates_after_deletion,Oberwolfach.triangle_relative_counts_after_deletion,Oberwolfach.ordinary_row_covers_full_conflict}
\uses{bp:triangle-parent-potential-event,bp:endpoint-factor-deletion}
\leanok
Let the fixed rim/root domain $Y$ be disjoint from the hub domain $B$. Restrict the sampling conflict graph to $Y$ when applying the endpoint factor theorem. Eligible edges, root domains, weight targets and potential sums remain unchanged. The same event controls every later degree-bounded rim and spoke deletion, with the original denominators and explicit additive losses. Any resulting row covers supported in $Y$ transfer back to the full rim conflict graph with identical physical resources and degrees.
\end{theorem}
\begin{proof}\leanok
All vertices of each eligible edge and its rim root lie in $Y$, so restricting the conflict graph changes none of the candidate tests. It makes every hub isolated, as the endpoint factor requires. Deleting conflict edges from the sampled rim leaves both potential sums unchanged because their candidates were already conflict-free. Later deletions have potential losses at most $5CL/|A|$ and $3CL/|A|$, for the normalized coefficient constant $C$. The triangle row-count loss is at most $3L+4L$, and the index-count loss is at most $3L$. These deterministic deductions permit all later choices on the same event, without a sampling-density factor. On each constructed rim support the restricted and full conflict relations coincide, so the final wheel transport preserves every edge and spoke. Complete preparation and the final cover-down assembly remain separate obligations.
\end{proof}

\begin{theorem}[Physical row covers on the triangle sampling event]
\label{bp:triangle-row-cover-common-event}
\lean{Oberwolfach.triangle_row_covers_on_sampling_event,Oberwolfach.triangle_row_cover_common_event,Oberwolfach.OrdinaryRowCover.filter_root,Oberwolfach.rooted_ordinary_row_cover_family}
\uses{bp:triangle-conflict-deletion-transport,bp:triangle-parent-potential-event,bp:two-domain-row-cover-event}
\leanok
One event under the original raw law supplies actual successive triangle row covers for every later permitted bounded deletion, free set and balanced root list. Parent matching weights are constructed inside the proof. Every retained wheel contains exactly one designated row root. The same exact-root conclusion now holds for the longer-cycle row-cover event.
\end{theorem}
\begin{proof}\leanok
Choose the count tolerance, potential tolerance and deletion margin before the graph and supplied density arrays. Apply the joint triangle sampling theorem, then its deterministic deletion bounds and the proved successive endpoint-factor construction. Restrict the conflict graph only for that construction and restore its full condition on the resulting rim supports. Finally remove wheels that miss their designated root. Such a wheel covers no edge or spoke at that root, so its removal preserves coverage; every degree and physical resource set only decreases. Since all other rim vertices lie in the free set, which excludes every designated root, each retained wheel contains exactly its own root from the list. This completes the event-to-row-cover implication. It does not yet provide the preparation geometry and balance, the hub/row combination or full cover-down.
\end{proof}

\begin{theorem}[Hub covers on the original sampling event]
\label{bp:hub-cover-common-event}
\lean{Oberwolfach.ordinary_spoke_omitted_resources,Oberwolfach.ordinary_spoke_potential_eq_local,Oberwolfach.ordinary_hub_counts_under_changes,Oberwolfach.ordinary_hub_cover_from_spoke_potentials,Oberwolfach.ordinary_hub_cover_after_changes,Oberwolfach.ordinary_hub_cover_common_event,Oberwolfach.ordinary_hub_packing_wheel_support,Oberwolfach.ordinary_hub_packing_row_degree}
\uses{bp:two-domain-row-cover-event,bp:ordinary-robust-local-factor,bp:positive-ordinary-cycle-weights}
\leanok
For an almost-regular parent above the ordinary cycle-decomposition threshold, with the supplied rim and spoke probabilities satisfying the degree-balance equation, one original-law event supplies actual successive hub-wheel covers. It permits every subsequent bounded host change and every hub list satisfying the explicit incidence and divisibility bounds. Parent cycle weights are constructed before sampling; the factor needed at each hub is proved, not assumed.
\end{theorem}
\begin{proof}\leanok
Reindex the root-spoke-omitted wheel load by its rim cycle and hub. The resulting sum is exactly the local cycle potential on that hub's neighbourhood; the root spoke is not tested. The simultaneous neighbourhood and potential estimates therefore provide the local-factor inputs. Later changes have additive count and potential losses with the original vertex-order denominator. During successive hub choices, previously used rim edges have degree at most twice the actual incidence with the chosen hub neighbourhoods. The proved robust local-factor theorem supplies each next factor. Exact resource identities show that every wheel uses a selected hub and that its rim lies in that hub's specified neighbourhood. The full preparation substitutions and protected cover-down remain open.
\end{proof}

\begin{theorem}[The row event after an earlier hub cover]
\label{bp:row-event-after-hubs}
\lean{Oberwolfach.ordinary_hub_cover_avoids_other_rows,Oberwolfach.ordinary_hub_deletion_preserves_row_balance,Oberwolfach.ordinary_hub_deletion_fixed_domain,Oberwolfach.ordinary_wheel_avoids_removed_packing,Oberwolfach.ordinary_row_cover_event_after_hubs}
\uses{bp:hub-cover-common-event,bp:two-domain-row-cover-event}
\leanok
The longer-cycle row event supplies actual row covers after every earlier hub cover whose rim vertices lie in the free vertex set and whose hubs lie outside the fixed row-counting hub domain. The resulting wheels are physically disjoint from that hub cover. The original neighbourhood targets, counting domains and law are unchanged.
\end{theorem}
\begin{proof}\leanok
The hub resource identity implies that later row-root stars are untouched, so their exact degree balance is preserved. Spoke rows restricted to the fixed free hub domain, and every column in that domain, are also unchanged. Only the actual bounded rim deletion contributes to the new error budget. The symmetric-difference triangle inequality accounts for it. Apply the original row event to the remaining host, then use its edge and spoke exclusions to prove resource disjointness. The row-count lemmas have also been strengthened: forbidden-pair exclusion is required only on the retained root edges, not on all parent edges. All callers of that change are checked. This does not yet combine the prescribed-edge stage or prove the triangle, marked or boundary versions of the complete cover-down lemma.
\end{proof}

\begin{theorem}[Exact cycle factors in almost-regular graphs]
\label{bp:almost-regular-cycle-factor}
\lean{Oberwolfach.ordinaryLocalPotential_full_domain,Oberwolfach.almost_regular_cycle_factor,Oberwolfach.ordinary_factor_graph_regular,Oberwolfach.regular_cycle_factor}
\uses{bp:ordinary-fixed-weight-factor,bp:positive-ordinary-cycle-weights}
\leanok
For each fixed cycle length $r\ge3$ and positive margin $\alpha$, sufficiently large graphs of minimum degree at least $(\delta_{C_r}+\alpha)n$ and sufficiently small degree variation have a spanning $C_r$-factor, provided $r\mid n$. In particular this applies to regular graphs at that threshold. The factor is an actual vertex-disjoint family of cycles whose union is a spanning two-regular subgraph.
\end{theorem}
\begin{proof}\leanok
Construct positive bounded fractional cycle weights from the ordinary threshold. On the full vertex domain the local potential with coefficients $2w(C)/D$ equals its original target exactly. Near-regularity controls the remaining error from that target to one. All common-neighbourhood inputs of the proved local-factor theorem hold identically with density one. Its absorption construction then gives an exact vertex factor. Mapping the local cycle family to ordinary cycles preserves the vertex partition; the physical union is two-regular. Near-regularity is a genuine hypothesis of this argument. The manuscript's minimum-degree-only transformer applications are not discharged by this result.
\end{proof}


\begin{theorem}[The prescribed-edge, hub and row stages on one event]
\label{bp:ordinary-three-stage-event}
\lean{Oberwolfach.ordinary_hub_and_row_common_event,Oberwolfach.ordinary_prescribed_hub_row_common_event,Oberwolfach.prescribed_inner_edge_enumeration,Oberwolfach.ordinary_deep_wheel_unique_root,Oberwolfach.prescribed_deep_packing_union}
\uses{bp:prescribed-packing-common-event,bp:hub-cover-common-event,bp:row-event-after-hubs}
\leanok
For $r\ge4$, one event under the original raw law constructs a prescribed-inner-edge packing, then the hub and row covers in its actual remainder. The fixed excluded hub set precedes sampling. The small inner rim set, prescribed edges, bounded comparison-host changes and balanced root lists may be chosen afterwards. The recombined packing covers every designated root star and meets the protected inner host in exactly the prescribed edge set.
\end{theorem}
\begin{proof}\leanok
Intersect the prescribed-edge and hub/row events using error exponent $q+1$. Exact degree bounds for the prescribed family control the realised changes to rims, spoke rows and spoke columns. These losses have no sampling-density multiplier. Removing wheels preserves balance and divisibility at the designated roots. The prescribed-edge list exhausts their inner incidences, so all remaining neighbours are in the free sets. The hub/row theorem constructs the next packing there. Its wheel origins give the support and protected-intersection statements; physical resource identities give disjointness and the final union. A separate exact enumeration counts every unordered rim edge once, and a physical vertex-set calculation gives one designated root per deep wheel. Translating the manuscript's separate rim and spoke degree bounds to the enumeration's combined row-incidence bound requires a factor of two. Full preparation substitutions, the triangle and marked compositions, the boundary cover-down and terminal absorption remain open.
\end{proof}

\begin{theorem}[Finite prescribed sets and physical boundary supports]
\label{bp:finite-deep-cover-boundary-supports}
\lean{Oberwolfach.ordinary_deep_cover_finite_sets_common_event,Oberwolfach.prescribed_physical_incidence_bound,Oberwolfach.ordinary_physical_retention_law,Oberwolfach.fresh_boundary_support_common_event,Oberwolfach.boundary_support_estimates_after_changes}
\uses{bp:ordinary-three-stage-event}
\leanok
The ordinary three-stage theorem is expressed using actual finite prescribed-edge and root sets. Its packing meets the protected inner host exactly in the prescribed edges, and every subsequent wheel has exactly one designated root. The separate rim and spoke degree bounds imply a combined physical incidence bound with a factor of two. Independently reserving boundary edges has the stated law on actual physical rim and spoke edges. One event controls all support sizes, pair intersections and vertex memberships on the original fixed domains.
\end{theorem}
\begin{proof}\leanok
Enumerate the finite prescribed sets without duplicating unordered rim edges and lift the indexed construction. Use the disjoint union of rim vertices and hub vertices to represent the physical graph. Its resource map is injective, so projecting independent edge coins gives the original rim/spoke retention law. Common-neighbourhood concentration with one or two roots yields the support estimates. A later neighbourhood change of size at most $L$ changes the one-root counts by at most $L$ and the two-root counts by at most $2L$, without multiplying either loss by the retention density. Full boundary factor and endpoint-selection callers remain separate obligations.
\end{proof}

\begin{theorem}[Bernoulli vertex loads and labelled absorber counts]
\label{bp:bernoulli-factor-counts}
\lean{Oberwolfach.bernoulli_root_load_mean,Oberwolfach.bernoulli_absorber_mean,Oberwolfach.bernoulli_factor_counts_common_event,Oberwolfach.root_omitted_load_local_change,Oberwolfach.absorber_ground_local_change}
\leanok
Fix the indexed resource family, its weights, absorber templates and class domains before independent vertex retention. One event controls the weighted load at every original root, every original absorber-root set of the prescribed size, and every supplied class size. Roots need not lie in the retained set. Later changes are measured by the size of the union of the retained and altered sets.
\end{theorem}
\begin{proof}\leanok
Omit the root trial in each load. Its expectation is exactly the original weighted degree times $p^{c-1}$. Each labelled absorber has exactly $c^2$ distinct free vertices, so its expected retained count is its original count times $p^{c^2}$. Bounded differences and the actual multiplicity bounds yield exponential tails. A union bound over the fixed roots and class domains gives one event; later choices add no events. For deterministic alterations, insert an omitted root into the local ground set when expressing its load as an induced weighted degree. This establishes the bound even when that root was absent from the sample.
\end{proof}

\begin{theorem}[Factor extraction after local alterations and sparse pair deletions]
\label{bp:factor-after-local-changes}
\lean{Oberwolfach.candidate_absorber_ground_general_pair_loss,Oberwolfach.resource_sparse_pairs_degree_loss_scale,Oberwolfach.absorber_ground_sparse_pairs_scale,Oberwolfach.profile_factor_after_local_changes}
\uses{bp:bernoulli-factor-counts}
\leanok
Given bounded positive indexed weights, near-unit omitted-root loads on a selected set, and sufficient labelled absorbers there, sufficiently small vertex alterations and sparse forbidden-pair sets preserve an exact factor whenever the final class sizes have the required ratio. The forbidden pairs may have quadratic total size and linear maximum incidence at the local order. The conclusion constructs the factor with the original surviving resource identifiers.
\end{theorem}
\begin{proof}\leanok
Separate forbidden pairs incident to a root from pairs fixing two free labels. Their absorber losses have powers $c^2-1$ and $c^2-2$, respectively. The analogous weighted-degree losses fix one or two additional cycle labels. Bound all unfixed labels using the local ground-set order, so the coefficients do not involve a sampling density. Choose the alteration tolerance from these explicit coefficients and the factor-extraction tolerance. Restrict the candidate family to identifiers avoiding the forbidden pairs, preserve every surviving weight, and apply the proved profile factor-extraction theorem. No factor-existence premise is introduced.
\end{proof}

\begin{theorem}[Factors on one Bernoulli vertex event]
\label{bp:bernoulli-profile-factor-event}
\lean{Oberwolfach.bernoulli_profile_factor_common_event}
\uses{bp:bernoulli-factor-counts,bp:factor-after-local-changes}
\leanok
For a fixed positive class profile, indexed resource multiplicity and absorber family, positive bounded weights with sufficiently small degree error and sufficient absorber counts admit the following conclusion. Choose the degree-error and alteration tolerances before the vertex-retention probability. For each fixed retention probability and requested inverse-power failure bound, all sufficiently large hosts have one event on which every admissible altered set with the exact final class profile has a factor avoiding every admissible sparse forbidden-pair set.
\end{theorem}
\begin{proof}\leanok
Use the common event for loads, absorber counts and total retained size. If $n$ is the original total order and $s=pn$, the final profile scale $k$ satisfies $s\le2ck$ and $k\le2s$. Rescale each original weight by exactly $p^{1-c}$. Both individual bounds, the omitted-root error and the absorber lower bound convert to the final scale using these inequalities. Shrink the alteration allowance before $p$ to account for its linear and quadratic conversions. The deterministic factor theorem then constructs every required factor. Only the order needed for concentration and the final factor threshold depend on $p$. The physical graph and marked-successor caller substitutions remain distinct obligations.
\end{proof}

\begin{theorem}[Physical forbidden edges and fixed domains]
\label{bp:fresh-physical-domains}
\lean{Oberwolfach.bernoulli_graph_avoiding_factor,Oberwolfach.supported_bernoulli_graph_avoiding_factor,Oberwolfach.supported_bernoulli_physical_factor,Oberwolfach.representedForbiddenGraph_degree_bound}
\uses{bp:bernoulli-profile-factor-event}
\leanok
The fresh factor event applies to a forbidden simple graph and to original resources supported on fixed embedded domains. Two injective physical representations of each auxiliary label, including identity and successor representations, are permitted. The resulting factor avoids every forbidden physical edge between distinct labels represented in either way.
\end{theorem}
\begin{proof}\leanok
The ordered forbidden pairs have total size at most $|Y|\Delta(F)$ and incidence at most $2\Delta(F)$. The sample-size estimate bounds $|Y|$ at the sampled scale. Restriction retains original resource identifiers and coefficients; an injective coordinate projection preserves the independent law exactly. Lift the constructed factor while preserving its avoidance property. For two representations the pulled-back forbidden graph has degree at most $4\Delta(F)$; reduce the allowance by four before sampling. No restricted permutation is substituted for the original map.
\end{proof}

\begin{theorem}[Fresh ordinary factors from neighbourhood counts]
\label{bp:fresh-ordinary-factors}
\lean{Oberwolfach.single_class_bernoulli_factor,Oberwolfach.ordinary_fresh_local_factor,Oberwolfach.ordinary_fresh_boundary_factor,Oberwolfach.incident_edge_retention_event}
\uses{bp:fresh-physical-domains}
\leanok
On a fixed domain of positive linear size, the sparse ordinary common-neighbourhood bounds and fixed positive near-unit vertex weights imply a fresh factor event. It includes every later small vertex alteration satisfying cycle-length divisibility and every later bounded-degree forbidden graph. The original-domain order determines the failure probability. The same conclusion holds for a neighbourhood sampled by its actual incident physical edges.
\end{theorem}
\begin{proof}\leanok
Construct paired absorbers using the proved ordinary common-path counts. Convert individual coefficient bounds from ambient order to domain size using its stated positive lower density. Apply the one-class fresh factor theorem and transport the altered set exactly. Distinct neighbours of a fixed root correspond to distinct undirected-edge coordinates, so the incident-edge law is exactly the required Bernoulli vertex law.
\end{proof}

\begin{theorem}[Simultaneous ordinary factor choices]
\label{bp:fresh-ordinary-simultaneous}
\lean{Oberwolfach.ordinary_fresh_domains_common_event,Oberwolfach.ordinary_fresh_simultaneous_factors,Oberwolfach.select_local_ordinary_factors,Oberwolfach.fixed_domain_common_event}
\uses{bp:fresh-ordinary-factors}
\leanok
One physical-coordinate event supplies robust factors for all supplied fixed domains. On that event, later altered domains satisfying the stated overlap, membership and numerical deletion budgets admit simultaneous edge-disjoint ordinary factors, with total graph degree at most $2p_{\rm sel}n$.
\end{theorem}
\begin{proof}\leanok
Increase the individual failure exponent by two and use the positive domain-size lower bound in a union bound; independence between domain events is unnecessary. Each constructed factor gives its actual graph with maximum degree two, exact support and forbidden-edge avoidance. Repeated bounded-deletion choices supply the candidate families of the proved selection theorem. The output retains the original cycle identifiers. The actual cover-down hypotheses must still supply the final domains, overlap bounds and numerical budgets.
\end{proof}

\begin{theorem}[Fresh endpoint triangle factors]
\label{bp:fresh-endpoint-factors}
\lean{Oberwolfach.endpoint_triangle_fresh_factor,Oberwolfach.endpointFactorGraph_choice}
\uses{bp:fresh-physical-domains}
\leanok
The endpoint neighbourhood estimates and original positive triangle weights imply a fresh two-class factor event. An exact two-to-one ratio is required only on the final altered classes. Every later forbidden graph within the degree allowance is avoided, and the returned endpoint records determine their actual three-edge graphs with maximum degree two.
\end{theorem}
\begin{proof}\leanok
Construct the paired absorbers from the endpoint row and index counts. Embed the two disjoint classes into the parent domain, preserving all triangle identifiers and weights. Their total size is bounded above by ambient order and below by the index-class estimate, giving the required coefficient and absorber conversions. Apply the fixed-profile fresh theorem with profile $(2,1)$, then lift the factor and its avoidance property. Realize each endpoint record as the triangle on its own three vertices.
\end{proof}

\begin{theorem}[Simultaneous endpoint triangle choices]
\label{bp:fresh-endpoint-simultaneous}
\lean{Oberwolfach.endpoint_fresh_domains_common_event,Oberwolfach.endpoint_fresh_simultaneous_factors,Oberwolfach.select_endpoint_factors}
\uses{bp:fresh-endpoint-factors,bp:fresh-ordinary-simultaneous}
\leanok
On one event in the actual independent physical-coordinate experiment, every eligible family of later altered endpoint domains satisfying the stated overlap and deletion budgets admits edge-disjoint endpoint triangle factors. Their union has degree at most $2p_{\rm sel}n$.
\end{theorem}
\begin{proof}\leanok
Project the original coordinate sample injectively for each fixed root and apply the domain-size union bound. Apply the proved simultaneous-selection argument to the actual endpoint triangle graphs, retaining the corresponding endpoint records. This establishes the factor selection component, not the full cover-down theorem, marked completion, terminal absorption or the main decomposition theorem.
\end{proof}


\begin{theorem}[Boundary supports after packing and covering]
\label{bp:boundary-after-packing}
\lean{Oberwolfach.boundary_after_packing_change_degree,Oberwolfach.boundary_after_packing_estimates,Oberwolfach.boundary_selection_supports,Oberwolfach.boundary_indexed_selection_supports,Oberwolfach.boundary_remainder_selection_common_event}
\leanok
\uses{bp:finite-deep-cover-boundary-supports}
Let $R\subseteq G$ be the reserved boundary graph, let $C$ avoid $R$, and let $D$ be the subsequent covering. The neighbourhood difference between $R$ and $(G-C)-D$ at $v$ is at most $d_{(G-R)-C}(v)+d_D(v)$. One original reservation event gives the support sizes, pair intersections and membership bounds for every later such $C,D$. The bounds are transported to the actual inner vertex subtype with each outside root counted once.
\end{theorem}
\begin{proof}\leanok
Every changed edge is either an uncovered edge outside the reservation or an edge of $D$. Apply the deterministic one- and two-neighbourhood loss estimates. If $|A|\ge\kappa n$ and the total count error is at most $\tau^2|A|$, the support sizes are at most $2\tau|A|$, intersections at most $(3\tau/\kappa)^2|A|$, and memberships at most $(3\tau/\kappa)|A|$. The losses are not multiplied by $\tau$. For the manuscript application $G$ here is the boundary graph, rather than the whole host with its reserved inner edges.
\end{proof}

\begin{theorem}[Joint hub and endpoint factors on the physical inner domain]
\label{bp:mixed-fresh-boundary}
\lean{Oberwolfach.fixed_domain_quadratic_common_event,Oberwolfach.outer_roots_fit_quadratic_local_order,Oberwolfach.ordinary_fresh_domains_quadratic_event,Oberwolfach.endpoint_fresh_domains_quadratic_event,Oberwolfach.select_mixed_boundary_factors,Oberwolfach.mixed_fresh_boundary_factors}
\leanok
\uses{bp:fresh-ordinary-factors,bp:fresh-endpoint-factors,bp:boundary-after-packing}
The same independent physical-coordinate event constructs ordinary triangle factors at hub roots and endpoint triangle factors at row roots, simultaneously for every permitted later change of their classes. Their physical graphs are mutually edge-disjoint across both families and have bounded union degree. The event permits quadratically many fixed root domains in the local order, with the tolerances chosen before the sampling probability.
\end{theorem}
\begin{proof}\leanok
Map ordinary rim factors into the physical inner domain; degrees are preserved and a forbidden physical graph pulls back with no larger degree. Intersect the two fresh-factor events without asserting independence and apply the proved selection theorem to their union of families. Local failure exponent $q+3$ allows quadratically many domains. This corrects an insufficient linear-domain interface: the number of outside roots can exceed the inner order, but is at most its square once the order exceeds the inverse size ratio. Original weights and cycle identifiers are unchanged. The actual preparation weights and numerical deletion budgets remain caller obligations.
\end{proof}

\begin{theorem}[Physical completion of the boundary triangles]
\label{bp:triangle-boundary-completion}
\lean{Oberwolfach.ordinaryWheelPhysicalGraph_edges,Oberwolfach.physical_triangle_wheel,Oberwolfach.root_completed_cliques_disjoint,Oberwolfach.root_completed_cliques_inner,Oberwolfach.root_completed_cliques_at_root,Oberwolfach.triangle_boundary_completion}
\leanok
\uses{bp:mixed-fresh-boundary}
An inner triangle completed by its outside root gives an actual ordinary three-wheel whenever its four physical vertices comprise three rim vertices and one hub, with the required host adjacencies and rim conflict avoidance. Vertex-disjoint triangles at one root and edge-disjoint triangles at different roots give edge-disjoint wheel resources. Their union consumes exactly the selected triangle edges inside the inner domain and has precisely the prescribed incidences at outside roots.
\end{theorem}
\begin{proof}\leanok
Construct the rim triangle and its spokes explicitly. The six physical edges are exactly the injective images of the rim and spoke resources. A common edge of two completed cliques either lies inside the inner domain or has an outside endpoint; the respective edge-disjointness or same-root vertex-disjointness excludes it. Restriction to the inner domain removes exactly the added root incidences. The remaining preparation and representation substitutions are not asserted by this conversion theorem.
\end{proof}


\begin{theorem}[Endpoint assignments after random restriction]
\label{bp:fresh-endpoint-pairings}
\lean{Oberwolfach.bernoulli_fixed_sets_common_event,Oberwolfach.bernoulli_graph_domain_common_event,Oberwolfach.sampled_domain_degree_after_changes,Oberwolfach.sampled_domain_codegree_upper,Oberwolfach.fresh_bipartite_common_event,Oberwolfach.fresh_neighbour_pair_common_event,Oberwolfach.fresh_neighbour_pair_physical_event}
\leanok
\uses{bp:boundary-after-packing}
Fix positive density and part-size lower bounds. There are positive parent and deletion tolerances, chosen before the exceptional-pair bound and sampling probability, such that one Bernoulli event supplies degree and cut bounds for every allowed subsequent change of the two vertex classes and every degree-bounded edge deletion. When the final classes have sizes $2m,m$, it supplies two distinct nonconflicting row neighbours for each hub, covering every row once. The same conclusion holds for quadratically many fixed roots in the original physical-coordinate experiment.
\end{theorem}
\begin{proof}\leanok
Concentrate the sizes, all ambient-root degrees and all two-root counts simultaneously. Realised vertex and edge changes cost their full cardinalities. Derive cut bounds from the degree and codegree estimates, then apply the proved two-neighbour assignment theorem. Projection onto each root's injective physical-coordinate map preserves the law. The union bound does not require independent root events. This proves the degree/cut and assignment assertions; the manuscript's displayed exponential envelopes and fixed-size coupling are not asserted by these direct Bernoulli lemmas.
\end{proof}

\begin{theorem}[Joint hub factors and endpoint assignments]
\label{bp:mixed-fresh-boundary-pairings}
\lean{Oberwolfach.EndpointPairing.graph_degree,Oberwolfach.EndpointPairing.graph_support,Oberwolfach.EndpointPairing.deleted_graph_choice,Oberwolfach.select_mixed_boundary_pairings,Oberwolfach.mixed_fresh_boundary_pairings}
\leanok
\uses{bp:fresh-endpoint-pairings,bp:fresh-ordinary-factors}
For ordinary wheel lengths at least four, one event constructs hub cycle factors and boundary endpoint assignments simultaneously on their physical inner domains. Subject to the explicit overlap, membership and numerical deletion budgets, their graphs are mutually edge-disjoint and the union has bounded maximum degree. All original cycle and pair witnesses are retained.
\end{theorem}
\begin{proof}\leanok
Intersect the fresh-factor and endpoint-assignment events with an increased failure exponent. Apply simultaneous selection to their union of families. Every single-copy premise is supplied by the corresponding common-event conclusion, including avoidance of the graphs arising during selection. The graph-specific preparation, changes after approximate packing, subsequent boundary paths and full cover-down theorem still require their remaining caller substitutions.
\end{proof}

\begin{theorem}[Prescribed boundary paths with bounded incidence]
\label{bp:boundary-path-completion}
\lean{Oberwolfach.boundary_wheel_after_deletion,Oberwolfach.boundary_path_wheel_packing,Oberwolfach.indexed_boundary_path_wheel_packing,Oberwolfach.boundary_path_exact_edges,Oberwolfach.boundary_finish_parameters}
\leanok
\uses{bp:mixed-fresh-boundary-pairings}
For $r\ge4$, simultaneously complete all prescribed root--endpoint pairs to ordinary $r$-wheels. Reserve the designated spokes for their own triples. Cap the number of uses of each free vertex. The resulting physical resources are disjoint, their degrees satisfy the displayed bounds, and the outside incidences are exactly the prescribed pairs and root spokes. The reservation density is chosen before the conflict-degree bound; only the subsequent order threshold depends on that bound.
\end{theorem}
\begin{proof}\leanok
Induct on the prescribed triples. Exclude vertices attaining the free-use cap and subtract all previously used edges at their realised degrees. Initial rim-deletion degrees are needed only at the two queried inner vertices. Reserve all assigned spokes, except those belonging to the current triple. The endpoint and hub incidence bounds control the remaining degree losses. An explicit square-root cap and order threshold satisfy the full extension inequality. No deletion loss is multiplied by a sampling probability.
\end{proof}

\begin{theorem}[Ordinary boundary completion from selected factors]
\label{bp:selected-ordinary-boundary-finish}
\lean{Oberwolfach.typed_endpoint_pairing,Oberwolfach.boundary_pair_incidences,Oberwolfach.endpoint_boundary_completion,Oberwolfach.selected_hub_factor_wheels,Oberwolfach.ordinary_selected_boundary_finish}
\leanok
\uses{bp:boundary-path-completion,bp:mixed-fresh-boundary-pairings}
For $r\ge4$, actual selected hub cycle factors and endpoint assignments determine an ordinary wheel decomposition covering every exterior incidence. The remainder is supported on the prescribed inner vertex and hub sets. The selected cycle identifiers are unchanged, and spokes are required only on the sets actually covered by those factors. The conclusion includes the rim-degree, inner-spoke-row and inner-hub-column bounds.
\end{theorem}
\begin{proof}\leanok
Interpret each endpoint pairing in the two physical summands. Its two row embeddings have disjoint images covering the prescribed row set. Mutual physical disjointness supplies the incidence bounds needed for path completion. Turn the selected hub factors into wheels and use the boundary path theorem in their actual remainder. The root and hub coverage identities give the support conclusion. This completes the deterministic final boundary stage, not the full preparation or cover-down theorem.
\end{proof}

\begin{theorem}[One boundary event for all later forbidden graphs]
\label{bp:boundary-late-forbidden-graph}
\lean{Oberwolfach.select_mixed_boundary_pairings_avoiding,Oberwolfach.mixed_fresh_boundary_pairings_avoiding,Oberwolfach.selected_factor_change_host,Oberwolfach.EndpointPairing.changeHost_after_deletion}
\leanok
\uses{bp:mixed-fresh-boundary-pairings}
The joint original-law factor and endpoint event permits every later forbidden physical graph satisfying the explicit degree budget. Selected cycles and assignments transfer to their actual remaining host with identical graphs, identifiers and physical resources.
\end{theorem}
\begin{proof}\leanok
At each selection step, avoid the union of the preceding selected graph and the later supplied forbidden graph. Its degree costs the sum of their actual degrees. The single-copy conclusions are universally quantified on the same previously proved event. Host transfer retains the chosen records and does not resample or recompute coefficients.
\end{proof}

\begin{theorem}[Bounded actual cycle factors for the ordinary transformer]
\label{bp:ordinary-bounded-two-factors}
\lean{Oberwolfach.supported_bounded_factors,Oberwolfach.single_class_bounded_factors,Oberwolfach.canonicalGraphPairs_disjoint_iff,Oberwolfach.ordinary_bounded_two_factors}
\leanok
\uses{bp:profile-bounded-extension,bp:embedded-profile-factor}
Given the common positive cycle weights and labelled-absorber counts, construct edge-disjoint ordinary cycle factors on $A$ and $A\cup Z$, where $A$ has bounded size, avoids the supplied forbidden set and contains no conflict pair. Every cycle of the enlarged factor contains at most one vertex of $Z$. The tolerance precedes the root bound, and the size of $A$ precedes the bounded forbidden-set and conflict bounds.
\end{theorem}
\begin{proof}\leanok
Apply the proved profile theorem in one-class coordinates on the actual fixed domain. Transport the same resource identifiers and weights through the embedding. Represent every undirected edge by its fixed quotient representative, shared by all cycles. Consequently disjoint displayed pair sets are disjoint physical rim-edge sets, even when two cycle representations orient an edge differently. This argument uses the supplied profile hypotheses and does not assume a minimum-degree-only cycle-factor theorem.
\end{proof}

\begin{theorem}[Constructed bounded hub switch]
\label{bp:ordinary-bounded-hub-trade}
\lean{Oberwolfach.selected_factor_hub_decomposition,Oberwolfach.ordinary_hub_factor_trade,Oberwolfach.ordinary_bounded_hub_trade}
\leanok
\uses{bp:ordinary-bounded-two-factors}
At two distinct hubs with the prescribed common vertex domain, the bounded factors give two actual wheel decompositions using the same rim graph. Their signed physical edge difference is exactly the difference between the two sets of spokes on $Z$. Every rim edge meets $A$, and the union rim graph has maximum degree at most four.
\end{theorem}
\begin{proof}\leanok
Use each unchanged factor at each hub. Swap the hubs of the two factors; their rim indicators cancel, as do the spokes on $A$. Disjointness of $A,Z$ and of the factors proves the remaining identity over integer edge indicators. The at-most-one-root condition excludes all rim edges inside $Z$. This supplies the two-factor switch in the relative transformer; the other path gadgets and their full transformer assembly remain separate obligations.
\end{proof}


\begin{theorem}[Ordinary triangle boundary with its selected factors]
\label{bp:ordinary-triangle-selected-boundary}
\lean{Oberwolfach.typed_endpoint_triangle,Oberwolfach.endpoint_triangle_complete_at_row,Oberwolfach.ordinary_triangle_row_boundary,Oberwolfach.ordinary_triangle_boundary_finish,Oberwolfach.ordinary_boundary_degree}
\leanok
\uses{bp:triangle-boundary-completion,bp:selected-ordinary-boundary-finish}
The actual selected endpoint triangles and hub triangle factors give a
three-wheel decomposition covering every exterior incidence. Its induced
inner physical graph is exactly the union of the selected factor graphs.
If that union has degree at most $d$ and a physical vertex belongs to at
most $k$ of the joint hub and row domains, the completed graph has degree
at most $d+k$ at every inner vertex.
\end{theorem}
\begin{proof}\leanok
An endpoint triangle has exactly two rim vertices and one hub. Adjoining
its prescribed row root gives an actual wheel. For a fixed root, the
selected triangles are vertex-disjoint; for distinct roots, their inner
edge sets are disjoint. Complete the unchanged hub factors at their
prescribed hubs. Mixed physical disjointness excludes shared inner rim
edges, and the two families use disjoint sets of hubs for their spokes.
The exact exterior neighbourhoods give the remainder support. Each inner
vertex gains one edge for each exterior root whose selected domain
contains it, giving the stated degree bound. This deterministic theorem
does not assert the remaining full cover-down substitutions.
\end{proof}


\begin{theorem}[Paired ordinary wheels with exact root differences]
\label{bp:ordinary-paired-root-wheels}
\lean{Oberwolfach.graph_swap_root_indicator,Oberwolfach.ordinary_wheel_root_replacement,Oberwolfach.ordinary_paired_root_wheels}
\leanok
\uses{bp:boundary-path-completion}
Common-neighbour bounds construct two ordinary wheels with the same
internal path, prescribed alternative roots and prescribed hubs. Their
rim indicators differ exactly on the two root-to-path edges at each
root. Their spoke indicators differ exactly on the root spokes and the
spokes of the common path at the two hubs. Both wheels avoid the original
forbidden vertices and the original conflict graph.
\end{theorem}
\begin{proof}\leanok
Exclude the second root and its forbidden-graph neighbourhood while
constructing the first wheel. This costs at most $L+1$ extra vertices.
The cycle's degree two identifies its two neighbours at the old root.
Relabel that root to the new root in the same cycle; the given endpoint
adjacencies and spoke conditions put the result in the original host.
Compute the changed rim edges and spokes using integer indicators.
Neither a conflict condition between the two roots nor spokes at the
wrong root--hub pairs are imposed. The subsequent actual vertex-stage
construction assembles these wheels with the rim-edge and spoke-edge families.
\end{proof}

\begin{theorem}[Actual two-root rim-edge wheel pairs]
\label{bp:ordinary-two-root-wheels}
\lean{Oberwolfach.graph_two_root_indicator,Oberwolfach.graph_two_root_subgraph,Oberwolfach.ordinary_wheel_two_root_replacement,Oberwolfach.ordinary_edge_path_wheels,Oberwolfach.ordinary_rim_edge_pair}
\leanok
\uses{bp:ordinary-paired-root-wheels}
The two ends of a prescribed rim edge can be replaced along one common
path, with an exact signed identity for the root edge, the two attachment
edges and the root spokes. The four-cycle construction uses a
three-neighbour intersection; longer cycles require two-neighbour
intersections. A common hub is chosen adjacent to all prescribed labels,
whose total number is at most $r+6$.
\end{theorem}
\begin{proof}\leanok
Construct the actual path and close it with each prescribed root edge.
Relabel both roots in the same cycle. The intermediate cross-edge cancels
in the signed calculation; its presence in the host is not required.
Choose the common hub using the specified common-neighbour lower bound
and the actual forbidden-hub budget.
\end{proof}

\begin{theorem}[Simultaneous ordinary transformer paths]
\label{bp:ordinary-transformer-edge-paths}
\lean{Oberwolfach.finite_disjoint_choices,Oberwolfach.finite_conflict_separated_choices,Oberwolfach.separated_transformer_paths,Oberwolfach.ordinary_transformer_edge_paths,Oberwolfach.fresh_transformer_hubs}
\leanok
\uses{bp:ordinary-two-root-wheels}
The rim and spoke paths are selected in one family of mutually disjoint
paths. No conflict edge joins vertices on different paths. All path
vertices avoid the old roots and their conflict neighbours. A spoke path
lies in the common neighbourhood of its two specified hubs. Fresh hubs
for the rim paths are selected after their additional attachment labels
are fixed.
\end{theorem}
\begin{proof}\leanok
Enlarge the forbidden set by the actual previously chosen supports and
their conflict neighbours. Charge at most $L+1$ vertices for each used
label, and apply the fixed-length path result with that explicit budget.
Use lengths $r-3$ and $r-2$ for the rim and spoke paths, respectively.
A second finite selection chooses distinct hubs from the common
neighbourhoods containing the future attachment labels.
\end{proof}

\begin{theorem}[Disjoint rim-edge and spoke-edge wheel families]
\label{bp:ordinary-transformer-edge-wheel-families}
\lean{Oberwolfach.ordinary_one_root_path_wheels,Oberwolfach.transformer_wheel_pair_disjoint,Oberwolfach.ordinary_rim_wheel_families,Oberwolfach.ordinary_spoke_wheel_families,Oberwolfach.IsEmbeddedCycle.no_triangle}
\leanok
\uses{bp:ordinary-transformer-edge-paths,bp:ordinary-two-root-wheels}
The selected paths give actual paired rim-edge and spoke-edge wheel
families, with exact root and path supports and signed physical-edge
identities. Each family is edge-disjoint on both sides. Spoke hubs may
repeat; distinct original spoke edges exclude the remaining collisions.
\end{theorem}
\begin{proof}\leanok
Construct the cycles on the same paths and use the prescribed hubs.
Private path supports reduce a possible shared edge to an edge or spoke
among the old roots. Distinct root edges or root spokes exclude that
case. Cycles of length at least four have no triangle, which also excludes
an edge between the two fixed neighbours of a vertex-wheel root. This
statement does not yet assemble the vertex-wheel families or the complete
transformer.
\end{proof}

\begin{theorem}[Transformer endpoint assignments and attachment cancellation]
\label{bp:ordinary-transformer-attachment-identity}
\lean{Oberwolfach.transformer_endpoint_assignments,Oberwolfach.canonical_edge_endpoint_sum_subtype,Oberwolfach.transformer_attachment_identity,Oberwolfach.shifted_hub_root_injective,Oberwolfach.shifted_hub_avoids_own_path}
\leanok
\uses{bp:ordinary-transformer-edge-wheel-families}
The equality $d_G(v)=2d_H(v)$ gives a bijection from incident rim edges
to the two endpoints of each incident spoke path. A permutation without
fixed points shifts the assigned rim-edge hubs. The resulting root-spoke
map is injective even after a root map that identifies nonadjacent
vertices. The corresponding attachment terms cancel exactly over the
integers.
\end{theorem}
\begin{proof}\leanok
Rotate a finite enumeration of each nonempty even-degree incidence set.
Sum endpoint contributions using one canonical representative of every
undirected edge. Reindex the spoke endpoints through the bijection and
the hub terms through the permutation. If two shifted root spokes
coincide, hub injectivity identifies the underlying rim edge; distinct
images of its two ends exclude the reversed case. This proves the
algebraic cancellation needed for identity (A). The later physical
identity applies it to the actual vertex-wheel assembly.
\end{proof}

\begin{theorem}[Actual spoke-path sets for the hub switch]
\label{bp:ordinary-transformer-spoke-sets}
\lean{Oberwolfach.spoke_index_fiber_card,Oberwolfach.ordinary_transformer_spoke_sets,Oberwolfach.balanced_spoke_root_separation}
\leanok
\uses{bp:ordinary-transformer-edge-paths,bp:ordinary-bounded-hub-trade}
Let $Z_u$ be the union of the actual paths corresponding to original
spokes at $u$. Then $|Z_u|=(r-1)d_H(u)$, the sets $Z_u$ are disjoint and
conflict-free, and every vertex of $Z_u$ is adjacent to both $u$ and its
image. Wheel divisibility implies $r\mid |Z_u|$.
\end{theorem}
\begin{proof}\leanok
Count the original spoke indices in each hub fibre, and use disjointness
and the exact length of each path. The same path sets supply the
common-hub adjacencies. Original hubs remain separate indices even when
their images coincide.
\end{proof}

\begin{theorem}[Actual endpoint pairs and four additional hub neighbours]
\label{bp:actual-transformer-endpoints}
\lean{Oberwolfach.canonicalEdgeDart_bijective,Oberwolfach.transformer_endpoint_realization,Oberwolfach.actual_rim_endpoint_adjacency,Oberwolfach.actual_spoke_endpoint_adjacency,Oberwolfach.transformerHubAttachments_card,Oberwolfach.actual_shifted_root_spokes_injective,Oberwolfach.ordinary_actual_rim_wheels}
\leanok
\uses{bp:ordinary-transformer-edge-paths,bp:ordinary-transformer-attachment-identity}
The actual disjoint paths determine two injective endpoint maps with
disjoint images. Pair these endpoints using the incidence bijections.
For each rim edge, its shifted assignments prescribe at most four extra
hub neighbours. The resulting fresh hubs belong to the common
neighbourhood of these labels and both root cycles.
\end{theorem}
\begin{proof}\leanok
Use one canonical orientation per undirected edge to identify its two
ends with directed incidences. The selected paths are simple, have
positive length, and have disjoint vertex sets. Their endpoint
occurrences are therefore injective. Apply the inverse shifted incidence
equivalence to the two ends of each rim edge to obtain its four labels.
The existing common-neighbour construction then chooses its actual hub.
\end{proof}

\begin{theorem}[Actual simultaneous vertex-wheel construction]
\label{bp:actual-transformer-vertex-wheels}
\lean{Oberwolfach.ordinary_vertex_wheel_pair,Oberwolfach.simultaneous_ordinary_vertex_wheel_pairs,Oberwolfach.ordinary_vertex_wheel_packing,Oberwolfach.ordinary_actual_vertex_wheels}
\leanok
\uses{bp:actual-transformer-endpoints,bp:ordinary-paired-root-wheels}
The prescribed endpoint pairs and shifted hubs give paired vertex wheels
with mutually disjoint internal vertex sets. These internal vertices
avoid all earlier roots and path vertices. Both families are physically
edge-disjoint, including when the root map identifies nonadjacent
vertices. The required count budget is
$|W|+|I|r+(r+4)(L+1)<t$.
\end{theorem}
\begin{proof}\leanok
Construct the actual alternative-root wheels, excluding the previously
used internal vertices at each step. Their endpoint pairs are disjoint.
The shifted root-spoke assignments are injective on both sides, which
excludes the only possible spoke collisions at old roots. The remaining
resources lie on private internal vertices or disjoint endpoint pairs.
\end{proof}

\begin{theorem}[Cross-family physical disjointness]
\label{bp:actual-transformer-cross-packing}
\lean{Oberwolfach.edge_vertex_wheels_disjoint,Oberwolfach.actual_attachment_spokes_unused,Oberwolfach.ordinary_actual_rim_vertex_packing,Oberwolfach.ordinary_actual_spoke_cross_packing}
\leanok
\uses{bp:actual-transformer-vertex-wheels,bp:ordinary-transformer-edge-wheel-families}
First edge wheels are disjoint from second vertex wheels, and conversely.
The rim-edge and spoke-edge wheel families are also mutually disjoint.
No injectivity of the root map on all vertices or of the map on all hubs
is imposed.
\end{theorem}
\begin{proof}\leanok
At a reused endpoint, the shifted hub differs from the hub of that
endpoint's rim path. Spoke-path endpoints lie on different paths from
all rim paths. The new vertex-wheel interiors avoid the earlier wheels.
A shared rim edge would then join the two neighbours of the vertex-wheel
root, producing a triangle in an $r$-cycle with $r\ge4$. Fresh hubs avoid
the original and image spoke hubs. These checks treat both decompositions
separately and retain physical edge multiplicity.
\end{proof}

\begin{theorem}[Physical attachment identity and its exact remaining term]
\label{bp:actual-transformer-physical-identity}
\lean{Oberwolfach.actual_transformer_attachment_cancellation,Oberwolfach.ordinary_actual_attachment_identity,Oberwolfach.transformer_original_edge_identity,Oberwolfach.transformer_spoke_defect_identity}
\leanok
\uses{bp:actual-transformer-vertex-wheels,bp:ordinary-transformer-spoke-sets}
For the actual constructed wheel pairs, identity (A) is the old physical
edge-set difference plus precisely the spoke difference on the sets
$Z_u$. Each physical old edge has multiplicity one. Different original
hubs remain separate summation indices even when their images coincide.
\end{theorem}
\begin{proof}\leanok
Apply the proved signed edge differences of each actual pair. Reindex
through the canonical incidence equivalence, endpoint bijection and
hub permutation. Injectivity on physical rim edges and spokes converts
the old-edge sums into indicator vectors. Disjoint spoke paths convert
the remaining sum into the indicator difference on each $Z_u$.
\end{proof}

\begin{theorem}[Realised ordinary edge and vertex stage]
\label{bp:ordinary-transformer-vertex-stage-realization}
\lean{Oberwolfach.ordinary_transformer_edge_stage,Oberwolfach.ordinary_transformer_vertex_stage_realization}
\leanok
\uses{bp:actual-transformer-cross-packing,bp:actual-transformer-physical-identity}
Given the selected edge paths and their constructed rim/spoke wheel
pairs, the stated common-neighbour estimates construct all vertex wheels
and give two actual wheel decompositions. Their physical difference is
$\chi_L-\chi_{L'}$ plus the explicitly identified hub-switch term.
\end{theorem}
\begin{proof}\leanok
Construct the vertex wheels with the actual shifted hubs and prove every
within-family and cross-family disjointness condition. The indicator of
each union equals the sum of its wheel indicators. Apply the actual
physical attachment identity. This completes the edge/vertex stage;
the hub trades and their size/root checks are completed in the subsequent
ordinary profile-host transformer theorem. Full terminal absorption and
the marked counterparts remain separate obligations.
\end{proof}

\begin{theorem}[Simultaneous host-preserving hub trades]
\label{bp:simultaneous-ordinary-hub-trades}
\lean{Oberwolfach.wheelFamily_host_decomposition,Oberwolfach.selected_factor_host_decomposition,Oberwolfach.ordinary_hub_factor_host_trade,Oberwolfach.ordinary_bounded_host_hub_trade,Oberwolfach.ordinary_simultaneous_hub_trades}
\leanok
\uses{bp:ordinary-bounded-hub-trade,bp:ordinary-bounded-two-factors}
The supplied common cycle weights and absorber counts construct all hub
trades simultaneously. Their new row sets are mutually disjoint, avoid
the earlier labels, and satisfy the required forbidden-pair conditions.
The size parameter is chosen before the forbidden-set and conflict
bounds. The existing sets $Z_i$ remain in their factor domains.
\end{theorem}
\begin{proof}\leanok
Retain the actual host wheels of the selected factors when forming both
decompositions. Apply bounded two-factor selection successively, excluding
previous new labels and their conflict neighbours. Its size parameter
does not depend on this bounded exclusion allowance. Disjoint $Z_i$ and
avoidance of the fixed old set give disjoint total row supports.
\end{proof}

\begin{theorem}[Exact spoke intersections and trade assembly]
\label{bp:ordinary-hub-trade-assembly}
\lean{Oberwolfach.ordinary_transformer_old_spokes,Oberwolfach.transformerSpokeRows_active_sum,Oberwolfach.ordinary_transformer_hub_trade_spoke_avoidance,Oberwolfach.ordinary_hub_trade_avoids_previous,Oberwolfach.resource_transformer_adjoin_trades,Oberwolfach.ordinary_profile_transformer_completion}
\leanok
\uses{bp:simultaneous-ordinary-hub-trades,bp:ordinary-transformer-vertex-stage-realization}
The two actual edge/vertex wheel families have, on each spoke path,
exactly its original and image old-hub spokes. The opposite hub trades
are physically disjoint from these families. Adjoining the constructed
trades cancels the full $Z_u$ term and gives a transformer in the
original host, retaining every wheel's forbidden-pair condition.
\end{theorem}
\begin{proof}\leanok
Fresh rim-edge and vertex-wheel hubs avoid the old hub set. Among
spoke-edge wheels, disjoint paths identify the unique contributing
index. Inactive hub indices contribute zero. Each trade rim edge meets
a newly selected row vertex, and its $Z_u$ spokes use the opposite hub.
The exact integer identity then yields the two required decompositions.
No injectivity of the hub map is imposed.
\end{proof}

\begin{theorem}[Ordinary transformer root restrictions and bounded order]
\label{bp:ordinary-transformer-root-order}
\lean{Oberwolfach.transformer_root_resources_disjoint,Oberwolfach.ordinary_transformer_row_buffer,Oberwolfach.ordinary_transformer_positive_root_restriction,Oberwolfach.ordinary_transformer_positive_row_avoidance,Oberwolfach.balanced_root_edge_counts,Oberwolfach.ordinary_resource_subgraph,Oberwolfach.ordinary_transformer_graph_realization,Oberwolfach.ordinary_resource_subgraph_order}
\leanok
\uses{bp:ordinary-hub-trade-assembly}
The transformer has no edge inside the old root sets. Every other used
row vertex avoids the earlier forbidden labels; every used hub in the
forbidden hub set is an original or image root. If at most $p$ old hubs
are active and each has degree at most $q$, the constructed row and hub
sets have total size at most
\[
 p\bigl(r(t+5q)+q+2\bigr).
\]
This bound does not count unused forbidden vertices.
\end{theorem}
\begin{proof}\leanok
On old roots the positive wheel family contains only original edges.
Each trade edge has a new endpoint or a spoke-path endpoint. Subtracting
the original edges therefore gives the stated restriction. Degree
balance gives equal rim/spoke edge counts and twice that many rim
incidences. Count the actual wheel supports, trade sets and used hubs.
Convert the physical resource set to its genuine simple rim graph and
bipartite spoke relation without changing any edge.
\end{proof}

\begin{theorem}[Complete ordinary transformer from fixed cycle profiles]
\label{bp:ordinary-profile-host-transformer}
\lean{Oberwolfach.ordinary_transformer_edge_realization,Oberwolfach.ordinary_transformer_from_paths_and_profiles,Oberwolfach.ordinary_profile_host_transformer}
\leanok
\uses{bp:ordinary-transformer-edge-paths,bp:actual-transformer-endpoints,bp:ordinary-transformer-vertex-stage-realization,bp:ordinary-hub-trade-assembly,bp:ordinary-transformer-root-order}
For $r\ge4$, the explicit two-row and two-hub common-neighbour bounds,
the additional three-row bound for $r=4$, the fresh-hub common-neighbour
bound, and the fixed positive cycle-weight and absorber-count profiles
give an actual ordinary transformer. Original and image root graphs are
row-disjoint, the map is injective on physical edges, and every active
hub differs from its own image. No injectivity on all vertices or hubs
is required. All displayed selection budgets are explicit hypotheses.
The tolerance precedes the root bounds; the transformer size precedes
the forbidden-set allowance.
\end{theorem}
\begin{proof}\leanok
Construct the incidence bijections and permutations, then all disjoint
rim and spoke paths. Choose the fresh hubs after their shifted endpoint
neighbours are known. Construct their edge wheels and the vertex wheels.
Use the unchanged profiles on the actual two-hub domains to construct
the simultaneous trades. Apply their exact disjointness and cancellation,
and count only vertices used. This proves the single-transformer result.
The full relative-absorber substitutions, finite four-transformer
construction, triangular/marked counterparts and main theorem remain open.
\end{proof}

\begin{theorem}[Explicit finite expansion graphs]
\label{bp:finite-expansion-degrees}
\lean{Oberwolfach.ordinaryExpansionRim_degree_old,Oberwolfach.ordinaryExpansionRim_degree_private,Oberwolfach.ordinaryExpansionSpokes_degree_private_hub,Oberwolfach.actual_edge_spoke_expansion_balance,Oberwolfach.actual_edge_spoke_expansion_divisibility,Oberwolfach.ordinaryExpansion_central_cycle_divisible,Oberwolfach.ordinaryExpansion_central_order}
\leanok
Each rim or spoke identifier has its own private wheel vertices and hub.
Replacing its distinguished edge by crossed edges gives a simple graph,
even when the original formal rim pairs are loops or repeat.
Actual degree balance and hub divisibility are preserved. The central
expansion of equal numbers of formal loops and spokes is wheel-divisible
when their number is divisible by the cycle length.
\end{theorem}
\begin{proof}\leanok
Compute both old endpoint counts, every private rim degree, every spoke
row degree and every hub degree. Identify the old endpoint counts with
actual graph degrees. Private vertices retain their wheel degrees.
\end{proof}

\begin{theorem}[Collapse maps on finite expansions]
\label{bp:finite-expansion-collapse}
\lean{Oberwolfach.ordinaryExpansion_rim_edge_injective,Oberwolfach.ordinaryExpansion_spoke_injective,Oberwolfach.ordinaryExpansion_resources_image,Oberwolfach.ordinaryExpansion_resources_injective,Oberwolfach.ordinaryExpansion_collapse_resources,Oberwolfach.ordinaryExpansion_reindex_rim,Oberwolfach.ordinaryExpansion_reindex_spoke}
\leanok
\uses{bp:finite-expansion-degrees}
Arbitrary maps on the old rows and hubs, retaining each private wheel
index, induce a bijection on physical expansion edges. Bijections of
private rim and spoke identifiers give the corresponding reindexings.
\end{theorem}
\begin{proof}\leanok
Every crossed edge has a private endpoint identifying its original edge
and its endpoint role. Distinct distinguished rim positions prevent a
collision. Private edges retain both their index and endpoints. Construct
preimages for each target rim edge and spoke.
\end{proof}

\begin{theorem}[Expansion-to-attachment edge bijection]
\label{bp:finite-attachment-bijection}
\lean{Oberwolfach.ordinaryAttachmentRim_eq,Oberwolfach.ordinaryExpansion_attachment_rim_hom,Oberwolfach.ordinaryExpansion_attachment_spoke_hom,Oberwolfach.ordinaryExpansion_attachment_rim_preimage,Oberwolfach.ordinaryExpansion_attachment_spoke_preimage,Oberwolfach.ordinaryExpansion_attachment_rim_injective,Oberwolfach.ordinaryExpansion_attachment_spoke_injective,Oberwolfach.actual_expansion_attachment_bijection}
\leanok
\uses{bp:finite-expansion-degrees}
For the actual edges and spokes of a graph pair, identifying each
private distinguished rim vertex with its selected old root maps the
expansion bijectively onto the attachment graph, preserving both parts.
The latter rim graph is the original graph together with its attached
cycle graphs.
\end{theorem}
\begin{proof}\leanok
Prove adjacency preservation and explicit preimages. An edge whose image
has a private endpoint is identified by that endpoint and its role.
Images with two old endpoints are identified by the original physical
edge. The deleted private edge prevents a crossed-edge collision. The
canonical actual graph encodings supply the required injectivity and
coverage, rather than leaving them as assumptions of the final result.
The host embeddings are proved below. The full four-transformer application remains open.
\end{proof}


\begin{theorem}[Attachment decompositions and common central resources]
\label{bp:attachment-decomposition-central-resources}
\lean{Oberwolfach.ordinaryAttachment_extra_wheel_decomposition,Oberwolfach.ordinaryAttachment_divisible,Oberwolfach.ordinaryExpansion_reindex_resources,Oberwolfach.ordinaryExpansion_common_central_resources}
\leanok
\uses{bp:finite-attachment-bijection,bp:finite-expansion-collapse}
The extra attachment edges decompose into their private wheels, even
when old roots repeat. Adding them preserves row balance and hub
divisibility. Collapsing old roots and reindexing equal-size edge families
gives exact edge bijections to a common central expansion.
\end{theorem}
\begin{proof}\leanok
Construct the actual wheel family, prove disjointness of its physical
edges and identify its exact rim and spoke unions. Compute degrees in
the disjoint union. Compose the verified collapse and reindex maps,
retaining both injectivity and the full image of the resource set.
\end{proof}

\begin{theorem}[Complete embeddings in a dense ordinary host]
\label{bp:dense-expansion-attachment-embeddings}
\lean{Oberwolfach.ordinary_expansion_host_embedding,Oberwolfach.ordinary_attachment_host_embedding}
\leanok
\uses{bp:finite-expansion-degrees,bp:finite-attachment-bijection}
Under the explicit dense-host degree and spoke bounds in the Lean
statements, finite expansions and attachments embed with the prescribed
old vertices fixed. All new rows and hubs are injective and avoid the
specified sets. Formal rim loops and repeated pairs are allowed.
\end{theorem}
\begin{proof}\leanok
Construct labelled cycles with their distinguished positions retained.
Choose the replacement cross-edges in the host, then choose private hubs.
Successive choices exclude all earlier private labels. For attachments,
place the unused private zero label at an additional isolated new vertex
so that the embedding of the full finite vertex type remains injective.
\end{proof}

\begin{theorem}[Complete embeddings in a relative ordinary host]
\label{bp:relative-expansion-attachment-embeddings}
\lean{Oberwolfach.sparse_rooted_cycle_embedding,Oberwolfach.sparse_rim_expansion_embedding,Oberwolfach.sparse_spoke_expansion_embedding,Oberwolfach.ordinary_relative_expansion_embedding,Oberwolfach.ordinary_relative_attachment_embedding}
\leanok
\uses{bp:finite-expansion-degrees,bp:finite-attachment-bijection}
The explicit common-neighbour bounds and exclusion budgets suffice for
complete expansion and attachment embeddings. All private rows are
separated from one another and from the old row set. The old row set
itself need not be separated. Dense minimum-degree or dense spoke-degree
conditions are not assumed here.
\end{theorem}
\begin{proof}\leanok
Choose paths using the actual two-neighbour intersections, subtracting
all forbidden labels and their separation neighbourhoods. Close the paths
into labelled cycles. Choose hubs from the common neighbourhood of each
separated cycle. For a spoke replacement, choose its new hub first and
construct the cycle inside that hub neighbourhood, retaining the old-hub
cross-spoke at position zero. Successive choices preserve all cross-family
separation conditions. The full four-transformer absorber and subsequent
preparation and cover-down applications remain open.
\end{proof}

\begin{theorem}[Ambient edge maps for the finite absorber graphs]
\label{bp:ambient-absorber-edge-maps}
\lean{Oberwolfach.ordinaryAmbient_resource_bijection,Oberwolfach.actual_embedded_expansion_attachment_bijection,Oberwolfach.actual_embedded_expansion_central_bijection,Oberwolfach.ordinaryAmbientMap_edge_roots_separated}
\leanok
\uses{bp:finite-attachment-bijection,bp:attachment-decomposition-central-resources}
The finite expansion maps induce the actual edge bijections between embedded
attachment, expansion and common central graph pairs. They preserve the
canonical undirected-edge representation, and their rim and spoke
injectivity and adjacency conclusions are proved. The corresponding root
separation follows from separation between the two embedded ranges and
separation of each displayed rim edge. The vertex maps need not be injective.
\end{theorem}
\begin{proof}\leanok
Extend the finite row and hub maps to the ambient types and prove that the
physical edge maps commute with the embeddings. Transport the finite edge
bijections, then recover the precise edge-subtype statements used by the
transformer. Check the four old row labels directly, including repeated
images and all cross-pairs.
\end{proof}

\begin{theorem}[Finite comparison graph and bounded active root sets]
\label{bp:finite-comparison-active-roots}
\lean{Oberwolfach.ordinary_comparison_from_divisibility,Oberwolfach.ordinary_comparison_identifier_equivalences,Oberwolfach.original_central_cycle_divisible,Oberwolfach.ordinarySupport_balance_divisibility,Oberwolfach.ordinary_active_root_order,Oberwolfach.actual_expansion_order_bound,Oberwolfach.ordinary_embedded_spoke_column_bound,Oberwolfach.ordinary_embedded_active_hub_bound}
\leanok
\uses{bp:finite-expansion-degrees,bp:attachment-decomposition-central-resources}
Local balance and hub divisibility give a decomposable disjoint-wheel
comparison graph with the required equal rim and spoke counts, and hence
the actual identifier bijections to the common central graph. Restricting
the original graph pair to its active vertices recovers it exactly after
embedding and bounds the finite source order by its physical support.
Expansion orders, root spoke degrees and active hub counts have bounds on
these finite types, independent of the ambient host order.
\end{theorem}
\begin{proof}\leanok
Sum the local degree identities. Construct the comparison wheels and their
edge partitions explicitly. Use equal finite cardinalities for the two
identifier equivalences. Restrict to the rim support and the incident hubs;
balance shows that every spoke row belongs to the rim support. Compute the
embedded degrees and bound the finite edge sets by pairs of source labels.
\end{proof}

\begin{theorem}[Fixed cycle profiles after row and rim deletions]
\label{bp:ordinary-relative-profile-restriction}
\lean{Oberwolfach.ordinary_restricted_profile_loads,Oberwolfach.ordinary_restricted_profile_absorbers,Oberwolfach.ordinary_relative_profile_restriction}
\leanok
The original common-cycle coefficients and labelled absorber families
supply the profiles on the actual surviving cycle family after a linear
row deletion and a bounded rim deletion. With a vertex deletion set $Z$
and deleted rim graph $D$, the weighted-degree loss is at most
$(|Z|+e(D))r^rB/n$. The absorber-count loss is at most
$(|Z|+e(D))C n^{r^2-1}$ for the explicit template-dependent constant in Lean.
The tolerance is chosen before the bounded edge allowance. The unchanged
weights are rescaled to $k=|T\setminus Z|/r$ with both lower and upper bounds.
\end{theorem}
\begin{proof}\leanok
Identify the surviving local cycle family by a support and edge predicate.
Bound each lost weight by a fixed pair weight and sum over actual deleted
vertices and canonical edge representatives. For absorber records, retained
roots force a deleted label to be free; a forbidden pair fixes at least one
free label in an occurrence. Transport all surviving template records
injectively. A cycle containing both endpoints of a deleted edge is used
only for an upper bound on loss; it is not identified with a cycle using
that edge. Choose the deletion tolerance first and the order threshold
after the edge allowance. The full transformer applications must still
supply their actual domains, edge sets and neighbourhood budgets.
\end{proof}

\begin{theorem}[Attachment decomposition in the actual separated host]
\label{bp:attachment-host-decomposition}
\lean{Oberwolfach.ordinary_attachment_extra_host_decomposition,Oberwolfach.ordinary_decomposition_embedding,Oberwolfach.ordinary_decomposition_recover_host}
\leanok
\uses{bp:relative-expansion-attachment-embeddings,bp:attachment-decomposition-central-resources}
The private attachment wheels decompose the actual embedded extra edges
in the given rim and spoke host, retaining its forbidden-pair condition.
The old row roots need not be mutually separated.
\end{theorem}
\begin{proof}\leanok
Each attachment wheel contains at most one old root. Its other row labels
are separated from all old roots and from one another. Embed the explicit
private wheel family, retain its edge disjointness and exact union, and
recover its membership in the actual host. No arbitrary universal-wheel
decomposition is presumed to satisfy the forbidden-pair condition.
\end{proof}


\section{Relative ordinary transformer from the original profiles}
\begin{lemma}[Endpoint exclusions and exact row location]
\label{bp:ordinary-transformer-row-location}
\lean{Oberwolfach.ordinary_transformer_forbidden_edges,Oberwolfach.ordinary_transformer_positive_row_location}
\leanok
The positive wheel family uses only old roots and the available row set.
The completed transformer's new factor vertices stay in the supplied cycle
domains. If all row and hub endpoints of a forbidden edge set are included
in the forbidden sets, the transformer's proved avoidance conclusions imply
disjointness from that edge set.
\end{lemma}
\begin{proof}
\uses{bp:ordinary-profile-host-transformer}
\leanok
Trace the supports of the edge wheels, vertex wheels and hub trades.
A shared rim edge would have both endpoints in the old row set; a shared
spoke would join an old row to a forbidden hub. Both are excluded by the
proved transformer conclusions.
\end{proof}

\begin{lemma}[Relative transformer with original weights]
\label{bp:ordinary-relative-host-transformer}
\lean{Oberwolfach.ordinary_relative_host_transformer,Oberwolfach.ordinary_finite_pair_transformer}
\leanok
For cycle length at least four, the original ambient-scale common-cycle
profiles and relative common-neighbourhood bounds construct an actual
transformer. A sufficiently small linear row deletion is allowed. The
linear hub exclusion is included in the forbidden hub set, with the old
spokes retained. The tolerance precedes the root size and conflict degree;
the size bound precedes the bounded forbidden-set allowance. New rows lie
in the available set or the specified old row set. Finite physical-edge
bijections supply the ambient map inputs without requiring their vertex
maps to be injective.
\end{lemma}
\begin{proof}
\uses{bp:ordinary-transformer-row-location,bp:ordinary-relative-profile-restriction,bp:ambient-absorber-edge-maps}
\leanok
Restrict the unchanged coefficients and convert the ambient scale to the
actual common-cycle domain. Subtract the deleted row count from each
neighbourhood bound. Local balance bounds the rim, spoke and incidence
counts, and the explicit numerical inequalities discharge every path,
vertex-wheel and hub margin. Apply the checked transformer construction.
\end{proof}

\begin{lemma}[Fresh expansions and the decomposable comparison graph]
\label{bp:ordinary-fresh-finite-host-models}
\lean{Oberwolfach.ordinary_fresh_expansion_embedding,Oberwolfach.ordinary_fresh_comparison_embedding,Oberwolfach.ordinary_relative_attachment_realization,Oberwolfach.ordinary_embedded_attachment_decomposition}
\leanok
The relative host contains fresh copies of the finite expansion and
comparison graphs, with all new rows separated from the prescribed old
rows and from one another. Attached wheels preserve the original pair,
and their actual physical edges form a disjoint, decomposable extra set.
Attaching to a decomposable comparison pair preserves its decomposition
in the supplied host.
\end{lemma}
\begin{proof}
\uses{bp:finite-comparison-active-roots,bp:attachment-host-decomposition}
\leanok
Choose the old-label copies greedily using the conflict-degree bound, then
apply the simultaneous private embeddings. Transport the exact attachment
edge partition. Every attached wheel contains at most one old row.
\end{proof}


\begin{lemma}[A disjoint family of actual relative transformers]
\label{bp:ordinary-relative-transformer-family}
\lean{Oberwolfach.ordinary_sequential_resource_selection,Oberwolfach.ordinary_relative_transformer_family}
\leanok
For a finite family of bounded root pairs satisfying the actual relative
transformer hypotheses in one host, there are mutually edge-disjoint
transformers. The order threshold accounts for all previously selected
row and hub labels. Every cycle coefficient is the originally supplied
coefficient. The initial tolerances precede the root bound and the
forbidden-set bound.
\end{lemma}
\begin{proof}
\uses{bp:all-length-relative-transformer}
\leanok
Before each choice add the endpoints of every earlier physical edge to
the forbidden row and hub sets. Their respective cardinalities are at
most the number of earlier choices times the proved row and hub bounds.
Apply the single-transformer theorem with this enlarged allowance.
\end{proof}

\begin{lemma}[Four-transformer assembly for an embedded diagram]
\label{bp:ordinary-four-transformer-assembly}
\lean{Oberwolfach.finite_four_transformer_absorber,Oberwolfach.ordinary_relative_four_transformer_absorber}
\leanok
Given the five embedded, mutually edge-disjoint graph pairs, their four
physical-edge maps, a decomposable extra attachment at the first pair,
and a host-wheel decomposition of the last pair, the original relative
profiles construct an absorber. The transformers are constructed rather
than assumed. Each intermediate graph enters the resulting absorber once.
\end{lemma}
\begin{proof}
\uses{bp:ordinary-relative-transformer-family,bp:attachment-host-decomposition}
\leanok
Include all five graph pairs in the common forbidden edge set. Construct
the four disjoint transformers, reverse the two required orientations,
and apply the exact transformer-chain identity and attachment cancellation.
\end{proof}

\begin{lemma}[The five finite graph pairs in one host]
\label{bp:ordinary-five-host-embeddings}
\lean{Oberwolfach.ordinary_fresh_comparison_attachment,Oberwolfach.ordinary_joint_expansion_embeddings,Oberwolfach.ordinary_diagram_model_size,Oberwolfach.ordinary_five_host_embeddings}
\leanok
The actual finite models of the ordinary absorber diagram can be embedded
jointly. The first attachment preserves the given original pair and has
a decomposable extra edge set. The last attachment is decomposable in the
actual host. The three expansions have disjoint row and hub label sets,
separated from one another and from both attachments. A common finite
model bound depends only on the original finite order and the cycle length.
\end{lemma}
\begin{proof}
\uses{bp:ordinary-fresh-finite-host-models,bp:finite-comparison-active-roots}
\leanok
Construct the first attachment, then a fresh decomposable comparison
attachment. Construct all three expansions successively while excluding
every label of those attachments and all earlier expansions. The finite
model bound discharges the entire accumulated allowance.
\end{proof}

\begin{lemma}[Concrete transformer inputs of the finite diagram]
\label{bp:ordinary-diagram-transformer-inputs}
\lean{Oberwolfach.ordinary_actual_diagram_divisibility,Oberwolfach.actual_finite_expansion_central_bijection,Oberwolfach.ordinary_embedded_transformer_data}
\leanok
All three concrete expansions satisfy row balance and hub divisibility.
Their maps to the common central graph are physical-edge bijections.
A finite physical-edge bijection between actual host embeddings supplies
the ambient transformer hypotheses, allowing noninjective vertex maps.
\end{lemma}
\begin{proof}
\uses{bp:ordinary-five-host-embeddings,bp:ambient-absorber-edge-maps}
\leanok
Use the original local divisibility to construct the decomposable
comparison graph and identify its edge and spoke index sets. Transfer the
finite maps through the injective host embeddings. Count active roots
using the finite source types, not the ambient vertex types.
\end{proof}

\begin{theorem}[Constructed ordinary host diagram]
\label{bp:ordinary-constructed-host-diagram}
\lean{Oberwolfach.ordinary_four_diagram_inputs,Oberwolfach.ordinary_host_diagram,Oberwolfach.ordinary_attachment_forbidden_resources}
\leanok
The original balanced divisible remainder and the host neighbourhood
hypotheses construct all five embedded graph pairs, their four ambient
transformer inputs, their physical-edge partitions and their exclusions.
The total numbers of row and hub roots are bounded before the host order
is chosen. Added attachment edges have a new endpoint.
\end{theorem}
\begin{proof}
\uses{bp:ordinary-five-host-embeddings,bp:ordinary-diagram-transformer-inputs}
\leanok
Apply each concrete expansion map to its actual embeddings. Form the
unions of their row and hub images, check every map orientation, and
bound the unions by the finite model orders. The attachment partition
excludes resources supported entirely on the forbidden old sets.
\end{proof}

\begin{theorem}[Constructed triangular relative transformer]
\label{bp:triangular-relative-transformer}
\lean{Oberwolfach.triangular_transformer_edge_realization,Oberwolfach.triangular_transformer_vertex_stage_realization,Oberwolfach.triangular_profile_host_transformer,Oberwolfach.triangular_parent_profile_host_transformer,Oberwolfach.triangular_relative_host_transformer}
\leanok
Four row-neighbour conditions, separated common-hub counts and the
original common-pair cycle profiles construct a triangular transformer.
The two incidences of each root edge share one new vertex. No injectivity
of these shared endpoints is assumed. The forbidden-set allowance is
chosen after the transformer size bound.
\end{theorem}
\begin{proof}
\uses{bp:ordinary-hub-trade-assembly,bp:ordinary-relative-profile-restriction}
\leanok
Choose the shared vertices and then the ordered spoke endpoint pairs.
Construct the actual edge and vertex wheel pairs. Prove disjointness
using the injectivity of endpoint--root and endpoint--hub pairs and the
absence of the shared attachment edge from the earlier wheel. Sum the
signed identities and apply the fixed-profile two-factor completion.
Restrict the original coefficients without renormalizing their scale.
\end{proof}

\begin{theorem}[Ordinary relative transformers for every $r\ge3$]
\label{bp:all-length-relative-transformer}
\lean{Oberwolfach.all_lengths_relative_host_transformer}
\leanok
The triangular construction and the longer-cycle construction have the
same resource, order and exclusion conclusions. For $r=3$ the row inputs
include four prescribed neighbours; for $r=4$ they include three; the
longer-cycle construction keeps its two-neighbour inputs.
\end{theorem}
\begin{proof}
\uses{bp:triangular-relative-transformer,bp:ordinary-relative-host-transformer}
\leanok
Split according to whether $r=3$ before choosing the tolerances and size
parameters. Apply the corresponding proved construction with its actual
row hypotheses and the original cycle profiles.
\end{proof}

\begin{theorem}[Relative ordinary host absorber for $r\ge3$]
\label{bp:ordinary-relative-host-absorber}
\lean{Oberwolfach.ordinary_diagram_absorber,Oberwolfach.ordinary_relative_host_absorber,Oberwolfach.ordinary_diagram_assembly_bounds}
\leanok
For $r\ge3$, the original common-pair cycle profiles, absorber counts,
row common-neighbourhood counts and separated hub counts construct a
bounded ordinary wheel absorber. Its order bound precedes the bounded
forbidden-set allowance. New rows avoid the linear row deletion set;
rows or hubs in the other forbidden sets must be old roots. Its physical
resources avoid every resource supported on the forbidden row and hub sets.
\end{theorem}
\begin{proof}
\uses{bp:ordinary-constructed-host-diagram,bp:ordinary-four-transformer-assembly,bp:all-length-relative-transformer}
\leanok
Construct the diagram, then all four transformers using the same original
cycle coefficients. Exclude the complete resource set on the diagram's
row roots and forbidden hub set. Apply the exact alternating chain
identity, and count the endpoint unions of that same absorber. All
numerical embedding and transformer margins are derived by choosing the
host order last. Profiles are required only for distinct hub pairs.
\end{proof}

\begin{corollary}[An actual ambient remainder and absorber graph]
\label{bp:ordinary-ambient-absorber}
\lean{Oberwolfach.ordinary_relative_ambient_absorber,Oberwolfach.ordinary_absorber_graph_realization}
\leanok
The relative construction applies to any bounded physical remainder by
restricting it to its active rows and hubs. The resulting resources recover
a genuine extended subgraph with the two required wheel decompositions,
the physical order bound, and the old-root edge exclusions.
\end{corollary}
\begin{proof}
\uses{bp:ordinary-relative-host-absorber,bp:ambient-absorber-edge-maps}
\leanok
The active row and hub embeddings recover the original remainder exactly.
Convert the absorber resource set to its rim graph and spoke relation,
retaining the actual host and forbidden-pair condition in both decompositions.
\end{proof}

\begin{theorem}[Marked paths and actual wheel pairs]
\label{bp:marked-transformer-wheel-stages}
\lean{Oberwolfach.marked_rim_cycle_families,Oberwolfach.marked_spoke_wheel_families,Oberwolfach.marked_rim_wheel_hubs,Oberwolfach.marked_fixed_attachment_wheels,Oberwolfach.marked_actual_rim_wheel_packing,Oberwolfach.marked_actual_spoke_wheel_packing}
\leanok
The row and hub common-neighbourhood bounds construct directed marked
paths, paired rim wheels and paired spoke wheels with the prescribed types.
Fixed attachment vertices can be retained when constructing a vertex-wheel
pair. The rim and spoke families separately have disjoint physical edges.
\end{theorem}
\begin{proof}
\uses{bp:marked-row-patterns,bp:marked-local-physical}
\leanok
Choose paths with specified types at each position and avoid the permutation
conflict neighbourhoods. The final connector is used only at distinct
endpoints. Align the shared path positions to prove that the hub's type
prescriptions are disjoint. Physical injectivity and the absence of other
old-root arcs prove the packing assertions, including repeated spoke hubs.
\end{proof}

\begin{theorem}[Marked assignments and exact attachment cancellation]
\label{bp:marked-attachment-cancellation}
\lean{Oberwolfach.marked_attachment_pairings,Oberwolfach.marked_rim_pair_integer_identity,Oberwolfach.marked_spoke_pair_integer_identity,Oberwolfach.marked_attachment_integer_identity}
\leanok
Outside the exceptional unpadded row of degree one with one marked
incidence, Hall's theorem assigns the hub occurrences away from their own
paths with the required incoming types. The signed wheel identities and
incidence bijections cancel the attachment edges exactly.
\end{theorem}
\begin{proof}
\uses{bp:marked-transformer-wheel-stages}
\leanok
Apply Hall twice to complements of partial matchings. Reindex each cycle
by its canonical marked rotation. Its physical marked edge retains the
permutation on the head. Sum incoming and outgoing incidences separately;
the hub assignment is a permutation of all incidences at the fixed row.
\end{proof}

\begin{theorem}[Marked two-factor hub switch]
\label{bp:marked-two-factor-hub-switch}
\lean{Oberwolfach.marked_factor_at_hub_decomposition,Oberwolfach.marked_hub_factor_trade}
\leanok
Two physical-edge-disjoint marked factors on $A$ and
$A\cup Z$, where $A\cap Z=\varnothing$, whose cycles are compatible with both distinct hubs $u,u'$,
give two exact wheel decompositions. Their signed difference is precisely
the spokes on $Z$ at $u$ minus those on $Z$ at $u'$.
\end{theorem}
\begin{proof}
\uses{bp:marked-transformer-wheel-stages}
\leanok
Attach each hub to the actual cycle family. Vertex disjointness and
physical injectivity prove disjointness within each decomposition. The
same rim factors occur on both sides and cancel; the spokes on $A$ also
cancel. This result does not assume a wheel decomposition, but its two
factor inputs must still be supplied by the common-profile construction.
\end{proof}

The complete marked transformer and five-graph absorber remain open.
The joint packing and the completion of a supplied packed diagram are
proved below; padding and the full diagram construction are still needed. The common-profile factor caller is
constructed below from its supplied original cycle estimates. The
remaining preparation and cover-down substitutions, iteration and
prescribed-family reconstruction also remain open. These checked stages
are not a certificate of Theorem~1.1.


\begin{theorem}[Balancing common marked classes with fixed weights]
\label{bp:marked-common-class-balancing}
\lean{Oberwolfach.common_class_balanced_profiles,Oberwolfach.class_marked_absorber_counts,Oberwolfach.typed_class_load_change_normalized,Oberwolfach.class_marked_path_count_change}
\leanok
The original common-cycle estimates construct two balanced classes containing
all prescribed roots. The successor exclusion, original cycle coefficients,
ambient error scale and omitted root in each weighted sum are preserved.
The tolerances precede the forbidden-set bound, host, permutation and
separation distance. Paired absorber counts are derived from the actual
cycle and common-path counts for these classes.
\end{theorem}
\begin{proof}
\uses{bp:marked-transformer-wheel-stages}
\leanok
Remove the forbidden vertices and successor conflicts, then trim the two
classes while retaining their prescribed roots. Charge each lost cycle or
path to a changed free label. The weighted comparison holds at every
ambient typed root; convert it to a degree sum only once the root belongs
to the new class. Order each root set by its type and construct the paired
absorber records using the proved common-path construction.
\end{proof}

\begin{theorem}[Constructed marked common-profile hub switch]
\label{bp:constructed-marked-common-hub-switch}
\lean{Oberwolfach.marked_common_hub_switch,Oberwolfach.marked_bounded_two_factors,Oberwolfach.marked_fixed_attachment_families}
\leanok
The supplied original common-cycle and omitted-root estimates construct
both bounded marked factors and their two exact physical wheel
decompositions. No factor or absorber-existence premise is added. The
size bound precedes the forbidden-set and conflict-degree bounds. Vertex
wheel families can also be constructed through their previously assigned
attachments by choosing only the internal path vertices.
\end{theorem}
\begin{proof}
\uses{bp:marked-common-class-balancing,bp:marked-two-factor-hub-switch}
\leanok
Apply the balanced-class construction. Rescale the coefficient bounds
from the ambient order to the class size without changing a coefficient.
Apply the checked bounded-factor construction, convert auxiliary pair
disjointness to physical edge disjointness using the original injective
map, and attach the two hubs. The resulting integer identity transfers
exactly the prescribed root spokes between those hubs.
\end{proof}


\begin{theorem}[Physical packing and exact marked diagram identity]
\label{bp:marked-packed-diagram}
\lean{Oberwolfach.marked_three_family_packing,Oberwolfach.marked_packed_attachment_identity,Oberwolfach.marked_initial_family_old_spokes,Oberwolfach.marked_active_hub_card}
\leanok
The actual rim, spoke and vertex wheel families have disjoint physical
edges on each side under the stated path, root and hub exclusions. Their
edge sets have signed difference equal to the two original edge sets plus
the prescribed spoke loads on the private paths. Only hubs incident with
original spokes contribute to the remaining sum.
\end{theorem}
\begin{proof}
\uses{bp:marked-transformer-wheel-stages,bp:marked-attachment-cancellation}
\leanok
Pull each common physical rim edge back to its original directed auxiliary
endpoints. Fixed attachment vertices identify their own rim path, and the
assigned hub avoids that path. For each side, convert the disjoint wheel
family into an edge set before using the integer cancellation identity.
At a private spoke-path label, fresh hubs cannot be old hubs, and disjoint
paths determine the unique old spoke-path index.
\end{proof}

\begin{theorem}[Attachment assignments in the caller's direction]
\label{bp:marked-assigned-attachment-types}
\lean{Oberwolfach.marked_row_assignment_types,Oberwolfach.marked_attachment_endpoints_injective,Oberwolfach.marked_assigned_attachment_rows_avoid_rim_path,Oberwolfach.marked_assigned_incidence_spoke_types}
\leanok
For rows satisfying the directed degree and marked-degree equalities,
outside the unpadded degree pair $(1,1)$, attachment slots can be assigned
to hub occurrences with exactly the required spoke types. Their prescribed
endpoint pairs are disjoint, and the four extra ordinary requests at each
rim hub avoid that hub's own rim path.
\end{theorem}
\begin{proof}
\uses{bp:marked-attachment-cancellation,bp:marked-transformer-wheel-stages}
\leanok
Invert the Hall assignment, which maps hub occurrences to slots. Handle
empty rows explicitly. Read the root spoke types from the unrotated rim
cycles and preserve the incoming/outgoing distinction. The two attachment
bijections and the disjoint private paths prove endpoint injectivity.
\end{proof}

\begin{theorem}[Simultaneously constructed marked hub switches]
\label{bp:marked-simultaneous-constructed-switches}
\lean{Oberwolfach.marked_constructed_switch_family,Oberwolfach.marked_switch_avoids_forbidden_resources,Oberwolfach.marked_forbidden_resource_rows_card}
\leanok
The supplied original common-cycle profiles construct all bounded hub
switches with disjoint physical resource sets. The size bound is fixed
before the forbidden-set bound, conflict degree and number of switches.
New rim edges avoid prescribed physical edges when the new auxiliary
labels avoid their endpoints and inverse-permutation endpoints; retained
root spokes are checked separately at their actual hubs.
\end{theorem}
\begin{proof}
\uses{bp:constructed-marked-common-hub-switch}
\leanok
Select the switches sequentially, excluding previous new labels and
conflict neighbours while retaining the fixed roots. Every factor rim
edge has a new auxiliary endpoint because the retained roots are
independent in the larger factor. A forbidden physical rim edge excludes
at most four auxiliary labels. Disjoint auxiliary supports imply physical
edge disjointness through the original injective representation.
\end{proof}

\begin{theorem}[Completion of a packed marked diagram from profiles]
\label{bp:marked-profile-transformer-completion}
\lean{Oberwolfach.marked_profile_transformer_completion,Oberwolfach.marked_transformer_switch_assembly}
\leanok
Given the actual initial packed wheel diagram, its exact original-edge
identity and the common-cycle estimates, the marked hub switches construct
a bounded transformer with two exact physical wheel decompositions.
Factor and absorber existence are not input hypotheses of this completion
result. Constructing the complete initial diagram, including padding,
remains a caller obligation.
\end{theorem}
\begin{proof}
\uses{bp:marked-packed-diagram,bp:marked-simultaneous-constructed-switches}
\leanok
The switch identity determines exactly which spokes occur on each
retained root. Auxiliary support and the opposite-hub exclusions prove
disjointness from the initial diagram. Adjoin the switches on opposite
sides and cancel their signed loads. The final indicator identity gives
the original-edge inclusions and exclusions and hence the transformer.
\end{proof}

\begin{theorem}[Padding wheels and their exact degree effects]
\label{bp:marked-padding-construction}
\lean{Oberwolfach.marked_padding_wheel_families,Oberwolfach.padding_row_balance,Oberwolfach.padding_hub_balance,Oberwolfach.padding_constructed_row_assignments,Oberwolfach.marked_padding_family_packing,Oberwolfach.marked_padding_avoids_old_resources}
\leanok
Host neighbourhood counts construct padding wheels at prescribed marked
heads, with private free vertices and fresh distinct hubs. Their occurrence
union preserves directed row balance and the marked row equality. Every
new hub has one marked spoke and exactly $m$ ordinary spokes. Padding all
rows of degree pair $(1,1)$ eliminates the Hall-assignment obstruction.
The padding wheels are physically disjoint and avoid the prescribed old
resources under the explicit row and hub exclusions.
\end{theorem}
\begin{proof}
\uses{bp:marked-transformer-wheel-stages,bp:marked-assigned-attachment-types}
\leanok
Construct the directed cycles first and choose their actual typed hubs
from the common-neighbourhood bounds. Count incoming occurrences by the
cycle rotation. A marked incoming arc and its marked spoke have the same
head. The old and new occurrence types are disjoint sums. If a padded row
had degree pair $(1,1)$, either its added marked-head count is zero, contrary
to coverage of every old exceptional row, or it is positive and its total
degree is at least two. Apply the checked assignment to these derived degrees.
\end{proof}

\begin{theorem}[Actual path and attachment constructions]
\label{bp:marked-path-attachment-construction}
\lean{Oberwolfach.marked_initial_paths,Oberwolfach.marked_attachment_hub_construction,Oberwolfach.marked_actual_attachment_arcs,Oberwolfach.marked_vertex_wheel_construction,Oberwolfach.markedPathUnion_card}
\leanok
The actual host counts construct the rim paths, spoke paths, their rim hubs
and all vertex-wheel paths. The four additional ordinary requests at each
rim hub are calculated from the chosen attachment bijections. The fixed
endpoint arcs and root spoke types used by the vertex wheels follow from
the earlier cycles and the incidence assignment. Only internal vertices
of a fixed-endpoint path are newly chosen. Separation from every earlier
path vertex, including unused attachment vertices, is retained in the result.
\end{theorem}
\begin{proof}
\uses{bp:marked-transformer-wheel-stages,bp:marked-assigned-attachment-types}
\leanok
Include all rim-path vertices among the forbidden roots when choosing
spoke paths. Their union has the explicit cardinality bound used by the
next choice. Endpoint ownership and the derangement prove that the four
hub requests avoid the hub's own rim path. Read the directed attachment
arcs in the unrotated cycle coordinates and the root spokes from the
canonical marked-head coordinates. The total number of incidence slots
is twice the number of original arcs.
\end{proof}

\begin{theorem}[Packing of the constructed families]
\label{bp:marked-constructed-family-packing}
\lean{Oberwolfach.marked_constructed_family_packing}
\leanok
The constructed rim, spoke and opposite-root vertex wheels form one
physical packing. The proof permits noninjective old row maps, requiring
only the original physical-edge injectivity and distinct images of the
two ends of each original arc.
\end{theorem}
\begin{proof}
\uses{bp:marked-path-attachment-construction,bp:marked-packed-diagram}
\leanok
Private paths and original edge injectivity give the rim and spoke
packings. The assigned root--hub pairs are injective even for a noninjective
row map. Shared attachment vertices identify their original rim incidence;
the derangement gives a different assigned hub. This discharges all
within-family and cross-family inputs of the joint packing theorem.
\end{proof}


\begin{theorem}[Constructed marked diagram and its completion]
\label{bp:marked-initial-diagram-completion}
\lean{Oberwolfach.marked_initial_diagram_construction,Oberwolfach.marked_diagram_exact_identity,Oberwolfach.marked_diagram_supported,Oberwolfach.marked_diagram_rows_card,Oberwolfach.marked_diagram_hub_profiles,Oberwolfach.marked_diagram_completion}
\leanok
The padded root graph and actual host counts construct the complete initial
diagram. Its six wheel families give both physical packings, the exact
signed edge identity, and explicit auxiliary row support. The old-hub
root classes are derived from these wheels; their type ratio, separation
and disjointness are proved. The supplied original common-cycle estimates
then complete the diagram to a transformer.
\end{theorem}
\begin{proof}
\uses{bp:marked-path-attachment-construction,bp:marked-constructed-family-packing,bp:marked-packed-diagram}
\leanok
Apply the proved path, incidence assignment, attachment-hub and vertex-wheel
constructions. Read the free-spoke types from the actual wheel coordinates,
restrict the exact residual sum to active old hubs and construct every hub
switch. Physical rim support is pulled back through the original injective
representation, so a marked physical head need not belong to the auxiliary
row support. The forbidden set may have size $\epsilon n+F_0$: balancing
removes it from the classes before the factors are chosen. The order
threshold absorbs $F_0$, with the tolerance and factor size chosen first.
\end{proof}

\begin{theorem}[Paired padding wheels and map extensions]
\label{bp:marked-paired-padding}
\lean{Oberwolfach.marked_padding_pair_construction,Oberwolfach.padding_cycle_map_extension,Oberwolfach.padding_hub_map_extension,Oberwolfach.marked_wheel_family_rim_injective,Oberwolfach.marked_wheel_family_spoke_injective,Oberwolfach.resource_transformer_remove_padding}
\leanok
Both families of padding wheels are constructed together. Their private
labels and hubs extend the old maps, preserve old values, and form a
physical packing on both sides. Old vertices may be identified by the
map. Physical rim and spoke occurrence maps are injective. Decomposable
padding can be removed from both transformer endpoints by retaining its
two wheel decompositions in the transformer.
\end{theorem}
\begin{proof}
\uses{bp:marked-padding-construction}
\leanok
Use the doubled index set for the simultaneous wheel construction. Extend
the map on new private labels, which are disjoint from all old labels.
Packing separates distinct wheels and the cycle coordinates separate
occurrences in one wheel. Combine each padded endpoint decomposition
with the padding decomposition on the opposite side.
\end{proof}


\begin{theorem}[Padded marked host transformer with geometric support]
\label{bp:marked-padded-host-transformer}
\lean{Oberwolfach.marked_padded_host_transformer,Oberwolfach.marked_diagram_plus_old_resources,Oberwolfach.marked_diagram_free_rows_separated,Oberwolfach.marked_diagram_hubs_card}
\leanok
The explicit host intersection counts and original common-cycle profiles
construct a transformer for the padded root graphs. It has bounded row
and hub supports, avoids all edges wholly among old roots, and its free
path rows are permutation-separated from one another and the old rows.
The switch rows are independent in the supplied forbidden graph and avoid
its neighbours at every initial row. The tolerance is chosen before the
root size, and the switch size before the bounded exclusion allowance.
\end{theorem}
\begin{proof}
\uses{bp:marked-initial-diagram-completion}
\leanok
Construct the initial diagram and forbid its rows and their neighbours
before selecting the common-hub switches. Bound this enlarged forbidden
set by its actual degree sum. Every old-root edge in the positive initial
packing belongs to the original graph and is absent from the transformer.
All new switch edges have a new endpoint. Switches introduce no new hubs.
\end{proof}

\begin{theorem}[Actual augmented root-graph identities and geometry]
\label{bp:marked-augmented-root-inputs}
\lean{Oberwolfach.padded_graph_resources,Oberwolfach.padded_graph_occurrences_injective,Oberwolfach.padded_maps_compatible,Oberwolfach.actual_padding_hub_balance,Oberwolfach.padding_all_exceptional_rows,Oberwolfach.marked_paired_padding_root_geometry,Oberwolfach.padded_typed_presence,Oberwolfach.marked_transformer_unpad}
\leanok
The actual paired padding wheels preserve row and hub balances, give
injective physical occurrence maps and compatible extended maps, remove
every exceptional marked row, and satisfy the paired-root separation and
typed edge hypotheses. Their two wheel decompositions remove the padding
from the endpoints of a completed transformer.
\end{theorem}
\begin{proof}
\uses{bp:marked-paired-padding}
\leanok
Compute the degrees on the disjoint occurrence union. Push the exact hub
ratio through the actual hub map, allowing identified old hubs. Private
padding labels are disjoint from all old labels and from the other side.
The canonical wheel coordinates give the types and physical edges. The
unpadding conclusion retains the two padding families in the transformer.
This node does not yet discharge the full augmented-host numerical and
common-profile application or the marked five-graph absorber.
\end{proof}


\begin{theorem}[Constructed marked host transformer]
\label{bp:marked-host-transformer}
\lean{Oberwolfach.marked_host_transformer,Oberwolfach.marked_root_padding_construction,Oberwolfach.MarkedRootPair.pad,Oberwolfach.MarkedRootPadding.remove,Oberwolfach.MarkedRootPadding.extended_rows_separated,Oberwolfach.marked_transformer_avoids_forbidden_resources}
\leanok
The original typed root graphs have a marked wheel transformer under the
explicit host intersection counts and original common-cycle profiles.
The construction supplies row and hub bounds, forbidden-row avoidance,
old-root edge exclusion and separation. Every hub in the forbidden old
hub set is already an endpoint hub. The tolerance precedes root size;
the switch size precedes the bounded exclusion allowance.
\end{theorem}
\begin{proof}
\uses{bp:marked-padded-host-transformer,bp:marked-augmented-root-inputs}
\leanok
Construct padding at precisely the exceptional source rows. Verify every
augmented occurrence, map, type and balance identity and substitute its
actual cardinalities in the padded host transformer. Retain the two
padding decompositions to remove padding from the endpoints. Combining
the separation and support conclusions proves avoidance of prescribed
physical resources by excluding their physical endpoints and inverse
permutation images. All original weights and omitted-root domains remain
unchanged. The five-graph marked absorber and the preparation applications
are still separate obligations.
\end{proof}


\begin{theorem}[Uniform marked transformer]
\label{bp:marked-uniform-transformer}
\lean{Oberwolfach.marked_uniform_host_transformer}
\leanok
The marked transformer has a size bound chosen before the bounded
forbidden-set allowance, permutation separation radius and forbidden-graph
degree bound. Only the host-order threshold depends on those later inputs.
\end{theorem}
\begin{proof}
\uses{bp:marked-host-transformer}
\leanok
Bound the actual padding sizes by the original occurrence and row bounds,
then absorb every fixed exclusion into the positive row and hub count
margins. The original common-cycle classes and coefficient functions are
unchanged when their allowed error is increased.
\end{proof}

\begin{theorem}[Finite marked expansion and attachment diagram]
\label{bp:marked-finite-expansion}
\lean{Oberwolfach.indexedMarkedExpansion_balanced,Oberwolfach.indexedMarkedExpansion_no_loops,Oberwolfach.indexedMarkedExpansion_rim_injective,Oberwolfach.indexedMarkedExpansion_spoke_injective,Oberwolfach.indexedMarkedExpansion_to_attachment,Oberwolfach.indexedMarkedComparison_balanced,Oberwolfach.indexedMarkedComparisonExpansion_to_central,Oberwolfach.indexedMarkedRootPair}
\leanok
Indexed marked edge expansions preserve all four balance identities,
remove formal loops and repeated directed endpoint pairs, and map to the
attached-wheel graph by an explicit edge permutation. A disjoint union of
marked wheels with the original four type counts is constructed. Both
expansions map to the same central expansion. Actual realizations supply
the transformer input with a possibly noninjective image vertex map.
\end{theorem}
\begin{proof}
\leanok
Exchange the heads or hubs of each old edge and its private same-type
wheel edge. The exchanges are typed involutions. Sum the local balance
identities to construct the typed occurrence bijections to the disjoint
wheel union. The central collapse retains every private copy index.
The host embeddings and the four-transformer absorber assembly remain
separate obligations; this theorem does not assume their existence.
\end{proof}


\begin{theorem}[Actual marked host realisations]
\label{bp:marked-five-host-realisations}
\lean{Oberwolfach.marked_five_host_realizations,Oberwolfach.marked_mapped_realizations_root_pair}
\leanok
The two attachments and all three intermediate expansions of the finite
marked absorber diagram have simultaneous realisations in the actual host
under its extension counts and the displayed finite exclusion budgets.
The first is the original graph together with a disjoint union of wheels;
the last has a wheel decomposition. The three intermediate row and hub
maps are injective. The attachment maps may identify vertices, while all
physical rim and spoke occurrence maps remain injective.
\end{theorem}
\begin{proof}
\uses{bp:marked-finite-expansion}
\leanok
Construct the prescribed marked cycles and hubs with the fixed physical
representation. Private expansion vertices are separated from old rows,
without assuming that all old rows are mutually separated. The precise
attachment permutations identify the resource sets. A joint finite choice
constructs the three expansions in disjoint row and hub sets. Pullback
along the central vertex maps preserves the actual physical resources.
The four-transformer assembly and preparation callers are separate work.
\end{proof}


\begin{theorem}[Constructed marked relative host absorber]
\label{bp:marked-relative-host-absorber}
\lean{Oberwolfach.marked_attachment_realization_separated,Oberwolfach.marked_comparison_attachment_separated,Oberwolfach.marked_host_diagram,Oberwolfach.marked_uniform_transformer_family,Oberwolfach.marked_diagram_absorber,Oberwolfach.marked_relative_host_absorber}
\leanok
From the actual marked host extension counts, the fixed original common-pair
cycle profiles and a bounded balanced indexed remainder realised in the
host, construct its marked wheel absorber. All five diagram graphs and
all four transformers are constructed. The absorber and its union with
the original remainder both have exact wheel decompositions. Its physical
resources avoid the specified old-resource set. New rows are mutually
separated and separated from the original roots; the original roots need
not be mutually separated. Row and hub bounds precede the bounded forbidden
allowance, permutation radius and forbidden-graph degree bound.
\end{theorem}
\begin{proof}
\uses{bp:marked-five-host-realisations,bp:marked-uniform-transformer}
\leanok
Retain separation in the attachment and comparison interfaces, construct
the four physical edge-preserving maps, and choose the four transformers
jointly. Enlarge permutation exclusions only around bounded sets of added
rows, preserving the order of the small error and later separation radius.
The alternating four-transformer chain gives both exact partitions. The
same finite supports give the order and exclusion bounds.
Preparation must still supply the actual host counts and original profiles;
this theorem is not a certificate for Theorem 1.1.
\end{proof}


\begin{theorem}[Absorbing an actual finite marked remainder]
\label{bp:marked-finite-remainder-absorber}
\lean{Oberwolfach.marked_finite_remainder_absorber}
\leanok
A finite set of actual arcs and spokes satisfying the four marked balance
identities and a bound on its active vertices has a relative absorber in
a host satisfying the extension counts and fixed common-pair profiles.
Both wheel decompositions, physical exclusions, vertex bounds and the
separation of the new rows are constructed. The absorber size precedes the
bounded forbidden allowance and the permutation separation radius.
\end{theorem}
\begin{proof}
\uses{bp:marked-relative-host-absorber}
\leanok
Restrict the indexed graph to the endpoints and hubs occurring in the given
remainder. Physical injectivity bounds its number of occurrences. Include
the rows, inverse-permutation heads and hubs of the finitely many forbidden
physical resources among the exclusions, and apply the constructed relative
absorber to this realised graph.
\end{proof}

\begin{theorem}[Fixed common profiles on one retained sample]
\label{bp:marked-retained-absorber-inputs}
\lean{Oberwolfach.marked_common_root_load_transfer,Oberwolfach.marked_common_completion_profile_realization,Oberwolfach.marked_completion_absorber_input_realization}
\leanok
The original retained sampling law supplies the marked host extension counts
and common-pair cycle profiles on a single event of probability at least
$1-2/n^Q$. The original coefficients are chosen before the actual marginal
arrays. Root-omitted estimates include every ambient root; the row errors
retain the original outer-domain scale.
\end{theorem}
\begin{proof}
\uses{bp:retained-common-joint-realization,bp:raw-neighbourhood-counts}
\leanok
Use exact equivalences between the sampled cycle classes and the surviving
parent cycles. Transport fixed weights through these equivalences, omitting
precisely the tested root spoke. Retained row and index tests supply all
extension counts. Intersect these with the common-profile event under the
same sampling law, using failure exponents chosen one larger than required.
\end{proof}

\begin{theorem}[Simultaneous marked absorbers after sampling]
\label{bp:marked-sampled-finite-absorbers}
\lean{Oberwolfach.marked_sampled_finite_absorbers}
\leanok
Under the displayed near-regular parent, density, permutation and sampling
hypotheses, one retained sample admits the constructed relative absorber
for every later bounded divisible marked remainder and every permitted
bounded physical forbidden set. The error constants precede the remainder
size; the absorber size precedes the forbidden allowance and radius. The
probability is at least $1-2/n^Q$ under the original law.
\end{theorem}
\begin{proof}
\uses{bp:marked-finite-remainder-absorber,bp:marked-retained-absorber-inputs}
\leanok
Choose the absorber constants from the fixed lower density bounds, then the
sampling error and the order threshold. On the common event, the concrete
sample has injective physical arcs, separated arc endpoints and disjoint
spoke types. These facts and the realised counts and profiles allow the
finite-remainder absorber construction for every subsequent permitted
choice. No new sampling or independent event is introduced for that choice.
The complete preparation, cover-down and main-theorem callers remain open.
\end{proof}


\begin{theorem}[Marked hub factors after physical deletions]
\label{bp:marked-deletion-robust-hub-factor}
\lean{Oberwolfach.marked_deletion_robust_hub_factor}
\leanok
A marked hub with balanced spoke classes, positive cycle and same-type
common-path counts, and a fixed two-sided approximate fractional weighting
admits a factor after every sufficiently small maximum-degree deletion of
physical rim edges. The tolerance precedes the permutation separation radius.
\end{theorem}
\begin{proof}
\leanok
Injectivity of physical arcs bounds all four auxiliary deletion degrees;
for marked incoming arcs the physical degree is evaluated at the successor
of the auxiliary head. Coordinate equivalences bound the lost cycle and
root-omitted path tuples. The unchanged weighting loses at most the proved
spoke-potential error. Apply the constructed marked hub factor to the
restricted host with half the original count margin.
\end{proof}

\begin{theorem}[Successive marked hub cover]
\label{bp:marked-successive-hub-cover}
\lean{Oberwolfach.marked_hub_cover_from_fixed_weights,Oberwolfach.marked_sample_hub_cover}
\leanok
For a sufficiently small prescribed set of distinct hubs, construct
physically disjoint wheels covering every spoke at those hubs exactly once.
The physical rim graph has degree at most twice the number of hubs.
On an actual physical sample, the fixed successor weighting and the
realised local configuration and spoke-load bounds supply these factors.
\end{theorem}
\begin{proof}
\uses{bp:marked-deletion-robust-hub-factor}
\leanok
At each hub delete the rim graph of the preceding factors, apply the
robust factor construction, and lift its wheels to the original host.
The distinct hubs give spoke disjointness. Exact sum identities identify
the sampled hub degree with the original root-omitted spoke potential;
injections transport the realised cycle and path counts. The complete
marked row and boundary cover-down callers remain open.
\end{proof}


\begin{theorem}[Marked prescribed inner-edge packing]
\label{bp:marked-prescribed-inner-edge-packing}
\lean{Oberwolfach.marked_prescribed_inner_edge_packing}
\leanok
Given the displayed row and hub intersection counts, an invariant inner
row set and inner hub set, construct pairwise physically disjoint marked
wheels for prescribed inner rim edges and spokes. Each wheel has exactly
its prescribed edge inside the inner pair. The numbers of prescribed
edges, their endpoint multiplicities and the free-use cap occur explicitly
in the three selection inequalities. Physical rim and spoke degrees satisfy
the stated bounds at all ambient vertices.
\end{theorem}
\begin{proof}
\leanok
First complete prescribed rim arcs using separated free paths and compatible
hubs outside the inner hub set. Saturated rows and hubs are excluded using
their total usage. Injectivity of the physical arc map and the exact old-arc
identity ensure that no later prescribed rim edge has been used. Then delete
these actual physical rims and spokes; the extension losses are bounded by
their realized degrees, without sampling factors. Complete the prescribed
spokes in this restricted host and lift their wheels. Marked physical heads
are charged at their successor vertices; inner membership uses invariance
under the ambient permutation. The numerical preparation substitutions and
full marked row and boundary cover-down remain separate obligations.
\end{proof}


\begin{theorem}[Sparse prescribed inner-edge packing]
\label{bp:marked-sparse-inner-edge-packing}
\lean{Oberwolfach.marked_sparse_inner_edge_packing}
\leanok
The explicit inner-edge selection inequalities hold when the numbers of
prescribed edges are at most $\eta^2n^2$ and their endpoint multiplicities
are at most $\eta^2n$. The admissible upper bound on $\eta$ depends only on
the extension density and cycle length; the order threshold is then chosen
after the permutation radius. The protected sets have their separately
specified size bounds. Both physical rim and spoke degrees are controlled.
\end{theorem}
\begin{proof}
\leanok
Use the free-use cap $\lceil\eta n\rceil$, derive each of the three displayed
selection budgets, and apply the prescribed inner-edge packing theorem.
\end{proof}

\begin{theorem}[Compatible typed endpoint pairings]
\label{bp:marked-typed-endpoint-pairing}
\lean{Oberwolfach.marked_endpoint_pairing_threshold}
\leanok
The five specified bipartite blocks admit the two compatible endpoint
bijections under the source degree, cut, class-size and label-fibre bounds.
The two labels paired at any hub are separated under the original permutation.
\end{theorem}
\begin{proof}
\leanok
Match the two incoming classes separately. Delete outgoing choices whose
labels are too close to their assigned incoming label. The fibre and
permutation-ball counts bound each row and column deletion. Quantitative
cut bounds survive this deletion. Hall's condition on the three remaining
blocks follows from the surplus of the ordinary hub class over the marked
outgoing class. This also forbids two marked arcs at one hub.
\end{proof}

\begin{theorem}[A marked wheel through assigned row attachments]
\label{bp:marked-assigned-row-wheel}
\lean{Oberwolfach.marked_assigned_row_wheel,Oberwolfach.marked_assigned_row_root_resources}
\leanok
Given the assigned root arcs, spoke and separated attachments, actual
extension counts produce a marked wheel containing all three prescribed
physical resources. Every added path vertex avoids the displayed forbidden
sets. A marked outgoing arc requires the marked spoke at its first attachment;
the other two type cases retain their original root or internal marked spoke.
\end{theorem}
\begin{proof}
\leanok
Choose a separated directed path with the required position-specific spoke
types, close it with the supplied arcs and rotate to the canonical marked
order. The exact physical root edges follow from this coordinate map.
Packing all assigned wheels, checking the actual cover-down input domains,
and assembling the final iteration remain open beyond this construction.
\end{proof}


\begin{theorem}[Marked row cover from endpoint bounds]
\label{bp:marked-row-cover-from-endpoint-bounds}
\lean{Oberwolfach.marked_row_cover_from_endpoint_bounds}
\leanok
Given actual incoming and outgoing arc indices at a row, its distinct spoke
hubs, the five quantitative endpoint blocks and the displayed host extension
counts, construct a physically disjoint marked wheel family covering every
prescribed incoming arc, outgoing arc and root spoke. The physical rim degree
inside the free row set, the spoke degrees at free rows and the spoke degrees
at each hub satisfy the stated bounds.
\end{theorem}
\begin{proof}
\leanok
Construct compatible endpoint matchings, retaining the three possible arc
type combinations. Distinct actual arcs and the injective physical map
prove disjointness of the prescribed root rims. Pack the assigned wheels
sequentially. Delete only the previously used physical rim edges inside the
invariant free set; each remaining rim edge is one of the two assigned root
edges. The degree bound separates forced attachment occurrences from free
choices. Distinct assigned hubs ensure spoke disjointness within this row.
The free-use cap is chosen before the permutation radius and fibre bound;
the order threshold pays their additive costs. Actual preparation and
multi-row cover-down substitutions, boundary completion and iteration are
not concluded by this theorem.
\end{proof}


\begin{theorem}[Actual marked endpoint blocks]
\label{bp:marked-actual-endpoint-blocks}
\lean{Oberwolfach.marked_endpoint_block_bounds, Oberwolfach.marked_row_cover_on_labelled_parts, Oberwolfach.marked_row_comparison_resources_disjoint, Oberwolfach.bipartite_block_after_deletion}
\leanok
The marked row and index intersection estimates give the five quantitative
endpoint blocks on their actual label sets. Under the explicit degree,
class-balance and packing inequalities, these blocks give a physically
edge-disjoint row cover with the asserted free-set degree bounds. Its rows
lie in the free set together with the prescribed root. Previously chosen
wheels supported on other roots and that free set can intersect it only in
resources belonging to the free set. Bounded later spoke deletions preserve
the displayed degree and large-cut bounds with the stated additive losses.
\end{theorem}
\begin{proof}
\leanok
Use singleton, pair and triple index estimates, retaining the original
error scales. The exceptional close pairs are bounded by the degree of the
permutation conflict graph. Transport the five blocks to the two tagged
endpoint classes; each label occurs at most twice. Apply the compatible
endpoint matchings and assigned-wheel packing. Physical-map injectivity
pulls a common physical rim back to the same auxiliary arc. Invariance of
the free set under the permutation then puts both physical endpoints in
that set. The deletion calculation subtracts actual row and column degrees
and actual cut losses. This does not yet supply the successive-row,
boundary, terminal or iteration callers of the main theorem.
\end{proof}


\begin{theorem}[Successive marked row covers]
\label{bp:marked-successive-row-covers}
\lean{Oberwolfach.marked_row_cover_after_free_deletion, Oberwolfach.marked_row_covers_from_endpoint_bounds, Oberwolfach.marked_endpoint_tagged_arc_embedding_exhaustive}
\leanok
The actual five endpoint blocks, root balances, separation conditions and
explicit numerical margins give physically disjoint covers for all
prescribed distinct roots outside an invariant free set. Every supplied
incoming arc, outgoing arc and root spoke is covered. Each row cover has
the displayed physical free-rim, free-row and hub-column degree bounds.
The sampled endpoint maps enumerate all actual root arcs with their other
auxiliary endpoint in the specified domain.
\end{theorem}
\begin{proof}
\leanok
Delete only the previously used physical rims inside the free set and
previously used spokes whose row belongs to that set. Root arcs survive:
for an incoming marked arc, its physical head is the successor of the root,
which is still outside the invariant free set. Transfer extension counts
with the realized rim and spoke losses, and transfer each endpoint block
by its row, column and cut inequalities. Apply the proved one-row
construction and lift the resulting wheels without changing any physical
resource. Physical-map injectivity and the row-support statements show
that any intersection with an earlier wheel would belong to the deleted
free resources. The induction sums the actual free degrees and covers all
roots. No row-cover existence premise remains in the concluding theorem.
Preparation substitutions, combination with the hub stage, boundary
completion, terminal reservation and the main theorem remain open.
\end{proof}


\begin{theorem}[Combined marked hub and row covers]
\label{bp:marked-hub-row-stages}
\lean{Oberwolfach.marked_hub_and_row_covers_from_weights, Oberwolfach.marked_successive_row_numeric_threshold, Oberwolfach.marked_row_cover_unique_auxiliary_root}
\leanok
The fixed positive cycle weights, common-path counts, actual endpoint
blocks and explicit deletion margins construct a single physical wheel
packing covering the prescribed hub spokes and all prescribed root arcs
and root spokes. Each row wheel contains its designated auxiliary root,
all its other auxiliary rows belong to the invariant free set, and its
hub belongs to the assigned hub set. The row-use cap and small-set
tolerances precede the permutation radius; only the order threshold
is chosen afterwards.
\end{theorem}
\begin{proof}
\leanok
Construct the hub factors using the checked fixed-weight theorem. Their
exact spoke union implies that every hub wheel has its rows in the free
set, and hence its physical rim also lies there. Delete this actual rim
graph, then construct the successive row covers with the already proved
endpoint-block theorem. Root arcs survive the deletion because the free
set is invariant under the permutation. Lift all wheels back to the
original host. The rim-avoidance conclusion and disjoint assigned hub
sets prove physical disjointness of the two stages. Each row wheel's
required root spoke proves root membership, and its row support proves
that no other designated auxiliary root occurs. This is not a statement
about the physical endpoints of twisted arcs. Preparation substitutions,
boundary completion, protected iteration, terminal reservation and the
unconditional main theorem remain open.
\end{proof}


\begin{theorem}[Marked boundary factors and endpoint sampling]
\label{bp:marked-boundary-sampling}
\lean{Oberwolfach.marked_fresh_hub_common_event, Oberwolfach.marked_fresh_endpoint_blocks, Oberwolfach.select_marked_endpoint_pairings, Oberwolfach.marked_factor_physical_supported}
\leanok
The original physical Bernoulli law gives a common event for all specified
hub factors and, separately, a common event for all specified endpoint
blocks. Each event holds for every subsequent allowed change of the
classes and every allowed physical edge deletion. Positive fixed cycle
weights and the actual cycle and root-omitted path counts construct the
hub factors. The five degree and cut blocks construct separated typed
endpoint assignments, with physical degree at most four.
\end{theorem}
\begin{proof}
\leanok
Injective maps associate each endpoint coordinate with its actual physical
edge. All endpoint blocks at one outside row use the same spoke coins.
Projection preserves the original law, and union bounds require no
independence between the events. The final class ratio is imposed after
sampling. Ambient-normalized weights are converted to the fixed parent
domain scale by its positive lower density. Physical forbidden edges pull
back to a graph on auxiliary labels with at most three times the degree.
A marked factor covers auxiliary set $Y$ and has physical support in
$Y\cup\pi(Y)$. Endpoint assignments retain both typed copies of a label;
their fibre bound gives physical maximum degree four. The simultaneous
assignment theorem discharges its single-choice premise from the five
actual blocks. Mixed factor/assignment selection, boundary path completion,
preparation substitutions and the unconditional main theorem remain open.
\end{proof}

\begin{theorem}[Joint marked boundary selection]
\label{bp:marked-mixed-boundary-selection}
\lean{Oberwolfach.marked_fresh_mixed_boundary_selection, Oberwolfach.marked_selected_hub_packing}
\leanok
On the original common physical sample, select hub factors and typed
endpoint assignments with pairwise disjoint physical graphs and an
explicit maximum-degree bound. The factors give actual wheel packings at
the distinct prescribed hubs, covering precisely the selected spokes.
\end{theorem}
\begin{proof}
\leanok
\uses{bp:marked-boundary-sampling}
The common hub event supplies a factor after every allowed physical rim
deletion. The five endpoint blocks supply each typed assignment after its
actual physical deletion. Their supports lie in the prescribed physical
row/hub set; use the proved bounded-degree selection theorem jointly.
Lift each selected factor to wheels. Distinct hubs separate their spoke
sets, and the selected physical rim graphs separate their rim edges.
\end{proof}

\begin{theorem}[Marked boundary paths from selected assignments]
\label{bp:marked-assigned-boundary-packing}
\lean{Oberwolfach.marked_boundary_packing_from_assignments, Oberwolfach.marked_boundary_counts_from_local_extensions, Oberwolfach.marked_boundary_finish_parameters}
\leanok
The actual typed endpoint bijections, together with the inner-set local
extension counts and the displayed degree and numerical bounds, give
physically disjoint boundary wheels. All prescribed incoming and outgoing
arcs and outside-root spokes are covered. The packing avoids the selected
initial rim graph and satisfies the stated inner rim, spoke-row and
spoke-column bounds. Roots and hubs may repeat.
\end{theorem}
\begin{proof}
\leanok
\uses{bp:marked-mixed-boundary-selection}
The selected physical degree bounds control endpoint and hub occurrences.
Arc injectivity follows from the original incidence maps; spoke injectivity
follows from the endpoint bijections and the distinct outside roots.
Reserve every assigned endpoint spoke before constructing any path.
At each step exclude saturated free rows, reserved spokes belonging to
other paths, and the actual earlier physical rim edges. Equal hubs have
disjoint assigned endpoint sets, so their reserved spokes cannot collide.
The local extension counts are taken directly on the inner set, without
charging its entire ambient complement. A square-root free-row cap and an
order threshold chosen after the separation radius give the numerical
margins. The original local-extension sets imply these two required counts.
\end{proof}

\begin{theorem}[Exact divisible marked boundary remainder]
\label{bp:marked-boundary-remainder}
\lean{Oberwolfach.marked_packing_remainder_divisible, Oberwolfach.marked_boundary_coverage_from_incidence_lists, Oberwolfach.marked_full_host_boundary_remainder}
\leanok
Subtracting a physical marked wheel packing preserves all four integer
balance identities on the original typed arc and spoke sets. If its
incidence lists cover the boundary, its remaining arcs have both auxiliary
endpoints in the inner row set and its remaining spokes lie in the inner
row/hub pair. Its physical resource set is exactly the original resource
set minus the packing.
\end{theorem}
\begin{proof}
\leanok
Each individual wheel has one marked arc and one marked spoke and obeys
the four local identities. Physical edge disjointness implies disjointness
of the actual arc and spoke sets. Sum their finite cardinalities and
subtract the balanced subgraph. Physical injectivity translates rim
coverage back to the unchanged typed arc identities. The incoming,
outgoing and hub lists exhaust the exterior incidences.
These results do not yet discharge the full preparation, protected
cover-down, terminal reservation, iteration or main-theorem callers.
\end{proof}

\begin{theorem}[Constructed marked boundary packing and remainder]
\label{bp:marked-selected-boundary-completion}
\lean{Oberwolfach.marked_selected_boundary_completion, Oberwolfach.marked_selected_boundary_divisible_remainder}
\leanok
The selected outside-hub factors and typed boundary assignments give one
physical wheel packing. It covers all boundary resources in the supplied
complete incidence lists. Its remaining typed arcs and spokes lie in the
inner row/hub pair and obey all four divisibility identities. The physical
remaining resource set is exactly the original set minus the packing.
The inner rim and both spoke-degree bounds are retained.
\end{theorem}
\begin{proof}
\leanok
\uses{bp:marked-mixed-boundary-selection,bp:marked-assigned-boundary-packing,bp:marked-boundary-remainder}
Construct the factors' wheels and then the assigned paths, forbidding the
actual factor rim. The two families use disjoint sets of hubs, so their
spoke resources are disjoint. The local extension hypotheses count new
rows directly in the inner set. Combine their exact coverage and degree
bounds, then subtract the contained divisible packing. The theorem still
requires the actual selection and local-count inputs; it does not certify
full protected cover-down or the main theorem.
\end{proof}

\begin{theorem}[Physical support event for the marked boundary]
\label{bp:marked-boundary-support-event}
\lean{Oberwolfach.marked_boundary_support_common_event, Oberwolfach.marked_boundary_support_at_row, Oberwolfach.marked_boundary_support_at_hub}
\leanok
Sample each physical edge once with probability $p$. One event, of
probability at least $1-n^{-q}$, bounds every boundary support by $13pn$,
every intersection at distinct outside indices by $145p^2n$, and every
label membership count by $13pn$, where $n=|V|+|U|$.
The supports are the actual incoming/outgoing boundary neighbourhoods
and their permutation images, together with the actual spoke neighbours.
\end{theorem}
\begin{proof}
\leanok
Every crossing typed edge has exactly one outside auxiliary endpoint and
one inside label. Its physical coordinate is injective. Thus different
outside indices have disjoint sets of coins. A coordinate affects at most
two support labels. The three directed edge types give an owner/label
fibre bound of six, hence at most twelve coordinates per support label.
The proved bounded-difference inequality and a single union bound give
all three estimates, without independence between the resulting events.
The ordinary and marked arcs keep their supplied directions and physical
maps. The successor acts only on rows, and fixes the hub summand.
\end{proof}

\begin{theorem}[Boundary support changes after the sample]
\label{bp:marked-boundary-support-changes}
\lean{Oberwolfach.boundary_label_changed_witness, Oberwolfach.boundary_label_row_changes, Oberwolfach.boundary_label_membership_changes, Oberwolfach.support_selection_margins}
\leanok
Each later changed membership has an actual changed physical coordinate.
A row support changes by at most twice the number of changed edges owned
by that row. A label membership count changes by at most the numbers of
changed edges at that label and its predecessor. If these row and label
losses are at most $p^2n$ and $pn$, the support family satisfies the
selection bounds with parameter $14p$.
\end{theorem}
\begin{proof}
\leanok
\uses{bp:marked-boundary-support-event}
Choose a coordinate witnessing the membership in one support but not the
other. Its retention indicators differ. Take images of these actual
changed edges and count their two possible labels. Intersections lose at
most the sum of the two row changes. These inequalities are deterministic
and hold for every permitted later choice on the original event.
\end{proof}

\begin{theorem}[Marked supports under later boundary changes]
\label{bp:marked-robust-boundary-support}
\lean{Oberwolfach.directed_boundary_support_changes, Oberwolfach.marked_boundary_robust_support_event, Oberwolfach.marked_boundary_selection_supports}
\leanok
One physical sampling event controls every later boundary mask with at
most $p^2n/2$ changed incidences at each auxiliary vertex, where
$n=|V|+|U|$. The actual mixed factor and endpoint supports then have
pair intersections at most $(14p)^2n$ and label membership at most $14pn$.
Only crossing boundary edges are charged; changes on other coordinates
are irrelevant to these supports.
\end{theorem}
\begin{proof}
\leanok
\uses{bp:marked-boundary-support-event,bp:marked-boundary-support-changes}
Each changed support label is charged to a changed crossing edge.
The factor and endpoint supports are subsets of the actual neighbourhood
supports at distinct outside hubs and rows. The embedding of these roots
preserves the pair and membership bounds.
\end{proof}

\begin{theorem}[Joint marked selection avoiding earlier edges]
\label{bp:marked-supported-boundary-selection}
\lean{Oberwolfach.select_mixed_marked_boundary_avoiding, Oberwolfach.marked_fresh_mixed_boundary_selection_avoiding, Oberwolfach.marked_fresh_supported_boundary_selection}
\leanok
The fixed-weight hub event and actual support event together supply
physically disjoint hub factors and endpoint assignments, avoiding a
supplied earlier used graph. Its full maximum degree is added to the
selection deletion budget. The five endpoint block hypotheses and the
class-balance and numerical conditions remain explicit inputs.
The event probability is at least
$1-|V|^{-q}-(|V|+|U|)^{-q}$, with every later mask and class choice
quantified inside this event.
\end{theorem}
\begin{proof}
\leanok
\uses{bp:marked-mixed-boundary-selection,bp:marked-robust-boundary-support}
Delete the union of the earlier graph and the choices made so far.
Use the fixed-weight factor event and the five Hall blocks on the
remaining physical resources. Intersect the factor and support events
by their individual failure bounds, without asserting independence.
\end{proof}

\begin{theorem}[Fresh endpoint blocks from marked profiles]
\label{bp:marked-fresh-endpoint-profiles}
\lean{Oberwolfach.marked_fresh_endpoint_profiles}
\leanok
The marked row and index profile estimates imply the fresh endpoint
block event on the original physical coordinates. The absolute errors
are compared with the lower bounds for the original endpoint classes.
The exceptional-pair graph has maximum degree at most $6d$; the
accuracy and later-change allowance precede the radius $d$ and the
sampling probability. All subsequent class changes and physical
edge deletions are quantified on the same event.
\end{theorem}
\begin{proof}
\leanok
\uses{bp:marked-actual-endpoint-blocks}
Two- and three-spoke index estimates give the row degrees and codegrees;
one-arc, one-spoke row estimates give column degrees. Embed the
permutation conflict graph in the physical row summand. Apply the
proved physical-coordinate endpoint sampling theorem to these actual
counts, without changing their domains or density targets.
\end{proof}

\begin{theorem}[The five endpoint substitutions]
\label{bp:marked-five-endpoint-profiles}
\lean{Oberwolfach.marked_fresh_five_endpoint_profiles, Oberwolfach.marked_endpoint_pairing_from_five_deleted_blocks}
\leanok
The five endpoint graphs use the explicit incoming/outgoing and ordinary/marked
tags from the manuscript. The repeated endpoint and hub classes use the same
physical coins. Their common event supplies endpoint assignments after each
permitted physical deletion, under the stated balance and separation margins.
\end{theorem}
\begin{proof}
\leanok
\uses{bp:marked-fresh-endpoint-profiles,bp:marked-typed-endpoint-pairing}
Index the five blocks at every outside root, apply the common profile event,
and substitute their actual row, hub and spoke types into the proved Hall
construction. Lift the assignments from the deleted host; their graphs remain
disjoint from the deleted physical graph.
\end{proof}

\begin{theorem}[Finish a selected marked boundary]
\label{bp:marked-mixed-boundary-finish}
\lean{Oberwolfach.marked_endpoint_assignment_total, Oberwolfach.marked_mixed_selection_bounds, Oberwolfach.marked_boundary_finish_parameters_upto, Oberwolfach.marked_mixed_boundary_finish, Oberwolfach.indexed_selection_small_density_budget}
\leanok
For every prescribed positive local extension bound and inner degree allowance,
all sufficiently small reservation probabilities permit the selected boundary
to be completed with that allowance. The parameter bound precedes the
permutation radius; the subsequent order threshold absorbs integer rounding.
The actual selected graphs supply the rim and endpoint degree bounds and the
total number of paths. The conclusion includes all four divisibility identities
and the exact physical inner remainder. Local extension counts, complete
incidence lists and the initial divisibility identities remain explicit inputs.
\end{theorem}
\begin{proof}
\leanok
\uses{bp:marked-supported-boundary-selection}
Separate the mixed physical union into its factor rims and endpoint graphs.
Count assignments by their hubs using the selected degree bound. Round
$56pn$ up to an integer, bounded by $57pn$ once $pn\ge1$. Apply the
proved boundary construction with free-row cap $\lceil\sqrt p\,n\rceil$.
The parameter proof holds for every $p$ below a fixed upper bound, so the
same $p$ may satisfy the earlier selection deletion budget.
\end{proof}

\begin{theorem}[One absorber for every terminal remainder]
\label{bp:marked-universal-terminal-absorber}
\lean{Oberwolfach.alternative_absorber_union, Oberwolfach.marked_supported_edges_card, Oberwolfach.marked_terminal_resource_count, Oberwolfach.marked_finite_terminal_absorber, Oberwolfach.marked_universal_terminal_absorber}
\leanok
In a marked host satisfying the proved finite-remainder absorber property,
reserve one fixed absorber for every later divisible subgraph of a specified
terminal host on at most $s$ auxiliary vertices. There are at most
$2^{2s^2}$ alternatives. If individual absorbers have row and hub bounds
$a_0,b_0$, their union has at most $2^{2s^2}(2a_0^2+a_0b_0)$ physical
resources, $2^{2s^2}a_0$ rows and $2^{2s^2}b_0$ hubs. It avoids the supplied
forbidden resource set and every resource contained in the terminal pair.
\end{theorem}
\begin{proof}
\leanok
Ordinary and twisted rims on a row set $X$ give at most $2|X|^2$ possible
physical resources; spokes give at most $|X||Z|$. Physical injectivity
bounds the number of typed terminal resources by $2s^2$. For each divisible
alternative, choose its proved finite absorber while excluding the resources
already used. The total exclusion budget is fixed before these choices.
Every absorber avoids every terminal resource, so their union absorbs each
alternative, even though the alternative remainders themselves may overlap.
Both decompositions are assembled as exact finite resource partitions.
\end{proof}

\begin{theorem}[Terminal absorbers on the original retained event]
\label{bp:marked-sampled-terminal-absorber}
\lean{Oberwolfach.marked_finite_to_terminal_absorber, Oberwolfach.marked_sampled_terminal_absorbers}
\leanok
The original marked retention experiment supplies one common event on which,
for every permitted terminal pair and bounded forbidden resource set, one
fixed absorber handles all later divisible subgraphs of that terminal host.
The degree tolerance precedes the terminal size; the absorber bounds precede
the forbidden allowance and permutation radius. No additional sample is used.
\end{theorem}
\begin{proof}
\leanok
\uses{bp:marked-universal-terminal-absorber}
Enumerate the terminal alternatives, increase the finite-absorber forbidden
allowance by the total size of the preceding absorbers, and enlarge the order
threshold to ensure that the terminal pair fits the original root allowance.
Apply event inclusion to the already proved finite-absorber common event.
\end{proof}

\begin{theorem}[Exact terminal completion]
\label{bp:marked-terminal-completion}
\lean{Oberwolfach.marked_terminal_completion}
\leanok
A fixed terminal absorber completes every later physical wheel packing that
avoids it and covers the remaining host outside the terminal pair. The host
is initially divisible; no divisibility assumption is imposed on the final
remainder. The conclusion is an exact physical resource decomposition.
\end{theorem}
\begin{proof}
\leanok
\uses{bp:marked-sampled-terminal-absorber}
Subtract a decomposition of the reserved absorber together with the later
packing. The four integer identities for wheel packings imply divisibility
of the exact typed remainder. Boundary coverage places it in the terminal
host. Absorb it, then unite the two disjoint physical decompositions.
\end{proof}

\begin{theorem}[Boundary completion after physical deletions]
\label{bp:marked-boundary-physical-deletion}
\lean{Oberwolfach.marked_boundary_counts_after_spoke_deletion, Oberwolfach.marked_boundary_counts_after_physical_deletion, Oberwolfach.marked_hub_factor_physical_restriction, Oberwolfach.marked_endpoint_physical_restriction, Oberwolfach.marked_restricted_packing_physical_avoidance, Oberwolfach.marked_selected_boundary_after_physical_deletion}
\leanok
A physical deletion of maximum degree $\Delta$ reduces the local boundary
extension lower bounds by at most $3\Delta$. Selected hub factors and endpoint
bijections that avoid the deleted edges remain valid in the remaining typed
host. The boundary construction there lifts to a packing in the original
host with the same resources and quantitative bounds, avoiding the deletion
on both rims and spokes. All prescribed root arcs and spokes are explicitly
required to remain present.
\end{theorem}
\begin{proof}
\leanok
\uses{bp:marked-mixed-boundary-finish}
Charge two rim constraints and one spoke constraint at their actual physical
endpoints. Restrict the typed occurrence sets and spoke relations. Preserve
factor coordinates and both endpoint bijections, construct the boundary
packing in this restricted host, and lift it through the original occurrence
maps. The loss is deterministic and has no sampling-density factor.
\end{proof}

\begin{theorem}[Boundary remainder after the preceding packing]
\label{bp:marked-boundary-after-packing}
\lean{Oberwolfach.marked_packings_disjoint_of_physical_graphs, Oberwolfach.marked_boundary_remainder_after_packing}
\leanok
Physically disjoint preceding and boundary wheel packings have a physically
disjoint union. If the latter covers every remaining outside incidence, the
exact typed remainder is supported on the inner pair and satisfies all four
divisibility equations of the original host.
\end{theorem}
\begin{proof}
\leanok
\uses{bp:marked-boundary-physical-deletion}
Translate physical graph disjointness into resource disjointness, using that
rim resources have two distinct endpoints. Each outside incidence is either
already covered or covered by the boundary packing. Subtract their actual
union and use the finite wheel divisibility identities.
\end{proof}

\begin{theorem}[Marked decompositions in another typed host]
\label{bp:marked-local-host-transfer}
\lean{Oberwolfach.marked_wheel_transfer_on_resources, Oberwolfach.marked_decomposition_transfer_on_resources, Oberwolfach.marked_absorber_transfer_on_resources, Oberwolfach.marked_decomposition_transfer_subhosts}
\leanok
A marked decomposition transfers to another host if every ordinary and
marked rim edge it uses has the same directed endpoints and type there,
and each used spoke retains its type. The cycle coordinates, hub and
physical resource sets are unchanged. No agreement is required outside
those resources, and no common physically injective superhost is assumed.
\end{theorem}
\begin{proof}
\leanok
Rebuild each cycle's arc witnesses with the supplied endpoint equalities.
Keep its vertex tuple and hub, and identify its physical rim and spoke sets.
Take the image of the finite decomposition. For restrictions of one host,
physical injectivity identifies the surviving arc identifiers.
\end{proof}

\begin{theorem}[Finite marked remainders between hosts]
\label{bp:marked-finite-remainder-host-transfer}
\lean{Oberwolfach.marked_remainder_divisible_subhost_iff, Oberwolfach.marked_finite_remainder_image_divisible, Oberwolfach.marked_finite_remainder_map_injective, Oberwolfach.marked_finite_remainder_image_resources, Oberwolfach.marked_terminal_remainder_absorption_between_hosts}
\leanok
An endpoint- and type-preserving map of the actual finite rim occurrences
preserves their physical resources and all four divisibility equations.
Physical injectivity in the original host proves injectivity of this map.
A terminal remainder in a second host can therefore be absorbed using the
first host's fixed reservoir when their terminal arcs and spoke types agree.
\end{theorem}
\begin{proof}
\leanok
\uses{bp:marked-local-host-transfer}
Compare each filtered finite set under its injective image. Construct the
image remainder, apply the first host's absorber, and transfer the two
resulting decompositions back on precisely the reservoir and remainder.
\end{proof}

\begin{theorem}[Fixed terminal completion in a modified host]
\label{bp:marked-terminal-completion-between-hosts}
\lean{Oberwolfach.marked_fixed_terminal_completion, Oberwolfach.marked_terminal_absorber_in_ambient_host, Oberwolfach.marked_terminal_completion_in_modified_subhost, Oberwolfach.marked_terminal_completion_between_hosts}
\leanok
The terminal absorber is chosen before a later modified host and its
covering packing. The modified host retains the absorber, its terminal
arcs are available in the sampled host with the same endpoints and types,
and the terminal spoke types agree. Every later physical wheel packing
that avoids the absorber and covers the outside resources extends to an
exact decomposition of the modified host. The final remainder's four
balances are derived; they are not extra assumptions.
\end{theorem}
\begin{proof}
\leanok
\uses{bp:marked-sampled-terminal-absorber,bp:marked-finite-remainder-host-transfer,bp:marked-terminal-completion}
Enumerate the sampled terminal host before reserving the absorber. Subtract
a decomposition of that absorber together with the later packing from the
modified host. Transfer the resulting divisible terminal remainder, absorb
it, and combine disjoint decompositions. This proves the terminal adapter;
construction of the covering packing by protected iteration remains a
separate obligation.
\end{proof}

\begin{theorem}[Terminal absorption outside a larger pair]
\label{bp:marked-terminal-larger-exclusion}
\lean{Oberwolfach.marked_finite_terminal_absorber_avoiding, Oberwolfach.marked_universal_terminal_absorber_avoiding, Oberwolfach.MarkedReservedTerminalAbsorberProperty, Oberwolfach.marked_finite_to_reserved_terminal_absorber, Oberwolfach.marked_sampled_reserved_terminal_absorbers, Oberwolfach.marked_reserved_terminal_completion_between_hosts, Oberwolfach.marked_reserved_absorber_avoids_inner}
\leanok
A bounded terminal pair $(X,Z)$ may lie inside a larger excluded pair
$(W,B)$. The number of possible terminal remainders depends only on
$|X|+|Z|\le s$, while $|W|,|B|\le\varepsilon n+f$ is permitted.
One reserved absorber avoids every resource internal to $(W,B)$ and
absorbs all later divisible subgraphs of the terminal host. The original
retained sampling law supplies this property simultaneously for all such
pairs and later forbidden sets. The constants $\varepsilon,\delta$ precede
$s$; the absorber bounds precede $f,d$ and the order threshold.
\end{theorem}
\begin{proof}
\leanok
\uses{bp:marked-sampled-terminal-absorber,bp:marked-terminal-completion-between-hosts}
Bound the support of each finite alternative, not the larger excluded sets.
Apply the finite-remainder construction with those larger sets, and take
the disjoint union of the individual absorbers. Count the alternatives
using only the terminal pair. Under invariance of $W$ under $\pi$, the
auxiliary endpoint exclusion is precisely the physical induced-edge
exclusion on supported resources. Transfer the final decompositions only
on resources used in the two hosts. This strengthens the reservation
interface without changing the random experiment or the main statement.
\end{proof}

\begin{theorem}[Protected hub and row covers and exact restoration]
\label{bp:marked-protected-stage-restoration}
\lean{Oberwolfach.marked_hub_row_avoids_inner, Oberwolfach.marked_hub_row_free_degree_bounds, Oberwolfach.marked_hub_and_row_covers_from_weights, Oberwolfach.marked_protected_resource_restoration, Oberwolfach.marked_protected_restoration_margin, Oberwolfach.marked_three_stage_restoration}
\leanok
The constructed hub and row covers have the displayed cumulative rim,
spoke-row and spoke-column bounds. If their free rows and allowed hubs
avoid the protected pair, their physical packing uses no edge internal
to that pair. For actual disjoint prescribed-edge, deep-vertex and active
packings, their union leaves the exact divisible remainder. The removed
inner resources are the disjoint union of the prescribed inner set and
the active packing's inner resources. Their deeper-star intersection is
exactly the prescribed set, and their deeper induced intersection is empty.
\end{theorem}
\begin{proof}
\leanok
Check every wheel's rows and hub. Bound the union degrees and count spokes
on the same actual packing. Subtract only whole wheels and identify the
physical complement. Filter the three unions by the inner, star and deeper
predicates; use the constructed exclusions and then the explicit numerical
margin. These assembly results still require the full cover-down caller
to supply its prescribed-edge, deep and active constructions together.
\end{proof}

\begin{theorem}[A packing for prescribed inner occurrences]
\label{bp:marked-prescribed-finite-subgraph}
\lean{Oberwolfach.marked_two_families_packing, Oberwolfach.marked_two_families_inner_intersection, Oberwolfach.marked_two_families_degree_bounds, Oberwolfach.marked_sparse_inner_packing, Oberwolfach.marked_sparse_prescribed_subgraph_packing}
\leanok
Under the proved extension-count hypotheses and sparse incidence bounds,
a supplied finite inner set of typed rim occurrences and spokes is covered
by one physical marked-wheel packing. Its inner intersection is exactly
that supplied resource set. The physical rim, spoke-row and spoke-column
degrees are at most $24\eta n$, $6\eta n$ and $3(l+5)\eta n$ respectively.
The tolerance precedes the permutation separation radius.
\end{theorem}
\begin{proof}
\leanok
Enumerate the actual finite occurrences by equivalences with finite ordinals.
Preserve every endpoint, type, physical image and filtered incidence count.
Apply the constructed sparse rim and spoke families, then unite their images.
The family disjointness proves physical packing, and exact filtered unions
and degree inequalities give all conclusions for the same packing.
\end{proof}

\begin{theorem}[Finite divisibility after sparse typed correction]
\label{bp:marked-finite-correction}
\lean{Oberwolfach.marked_matrix_balanced_divisible, Oberwolfach.marked_matrix_correction_finite_divisible, Oberwolfach.marked_matrix_deletion_avoids_protected_pair, Oberwolfach.sparse_typed_finite_host_correction}
\leanok
The actual sparse correction matrices yield a finite retained host satisfying
the four marked divisibility identities. Original occurrence identifiers
and spoke types are preserved. If the allowed matrices vanish on a protected
pair, the actual correction deletes no physical edge internal to that pair,
provided its row set is invariant under $\pi$.
\end{theorem}
\begin{proof}
\leanok
Convert the matrix row and column sums into filtered counts. Disjointness of
the two spoke types prevents double counting. Transport the restricted arc
occurrences by their forgetful injection. The correction matrices are bounded
by the allowed matrices, so their entries vanish on the forbidden pair.
\end{proof}

\begin{theorem}[Balances at vertices of a temporarily reserved star]
\label{bp:marked-reserved-star-balances}
\lean{Oberwolfach.marked_star_restriction_row_counts, Oberwolfach.marked_star_restriction_hub_counts, Oberwolfach.marked_remove_reserved_star_deep_balances}
\leanok
Restricting a finite marked host to edges incident with specified rows or
hubs preserves all relevant incidences at those rows and hubs. Subtracting
the star of a divisible inner subgraph from a divisible host therefore
preserves the three row balances on the specified rows and the hub balance
on the specified hubs. Divisibility of the entire star is not needed or claimed.
\end{theorem}
\begin{proof}
\leanok
Every incidence counted at a specified row or hub already belongs to the
star. Identify the filtered finite sets, then subtract their cardinalities
using the actual subgraph containments and the two original balances.
\end{proof}

\begin{theorem}[Actual endpoint lists in a modified marked host]
\label{bp:marked-actual-deep-endpoints}
\lean{Oberwolfach.occurrenceNeighbourEquiv, Oberwolfach.occurrenceNeighbourEmbedding_exhaustive, Oberwolfach.occurrenceRootNeighbours_card, Oberwolfach.markedNeighbourArcEmbedding_exhaustive, Oberwolfach.marked_neighbour_row_balances, Oberwolfach.marked_actual_row_required_iff, Oberwolfach.marked_actual_star_covered, Oberwolfach.marked_finite_host_local_row_balances, Oberwolfach.marked_finite_host_local_hub_balance}
\leanok
The actual incoming and outgoing neighbourhoods enumerate every retained
occurrence at their root, preserving its type and endpoints. The three
local row balances give the required endpoint cardinalities. Covering these
lists and the actual hub spokes covers the entire physical star of an
invariant row set. Viewing retained finite sets as a typed subhost preserves
these local balances; no global balance is required at other roots.
\end{theorem}
\begin{proof}
\leanok
Ordered-endpoint injectivity identifies each neighbour with its unique
incident occurrence. The resulting equivalences prove both exhaustion and
cardinality. Split the ordinary and marked arc sums, and count the disjoint
spoke types. For a marked incoming edge, permutation invariance identifies
physical star membership with auxiliary star membership.
\end{proof}

\begin{theorem}[Deep support after temporary reservation and prescribed covering]
\label{bp:marked-deep-reservation-clearing}
\lean{Oberwolfach.marked_reserved_star_clearing, Oberwolfach.marked_cleared_star_neighbours, Oberwolfach.marked_local_balances_after_packing, Oberwolfach.marked_reserved_star_and_packing_deep_balances}
\leanok
Reserve the star of the corrected inner graph and cover the correction
edges in that star by the prescribed packing. The remaining inner graph
has no edge incident with a deeper row or hub. Every remaining neighbour
at such a row or hub therefore lies outside the inner pair. The local
balance identities survive both the reservation and the actual wheel packing.
\end{theorem}
\begin{proof}
\leanok
The exact inner-star intersection before the prescribed covering is the
prescribed correction set. Remove it and use the physical/auxiliary
identities on invariant row sets. Subtract the exact finite incidence
counts of the already proved wheel packing from the local star balances.
\end{proof}

\begin{theorem}[Constructed hub and actual row covers with full support bounds]
\label{bp:marked-actual-deep-constructor}
\lean{Oberwolfach.marked_hub_and_actual_row_covers, Oberwolfach.marked_hub_row_supported_avoids_star, Oberwolfach.marked_hub_row_avoids_remaining_inner_star, Oberwolfach.marked_hub_row_full_free_rim_degree}
\leanok
Under the displayed fixed-weight, extension, endpoint-block and numerical
hypotheses, the hub-and-row construction covers the actual incident edges.
It retains its physical disjointness and degree bounds and records the
support of every wheel. This proves avoidance of the entire star of the
other inner vertices. At a free vertex, the full physical rim degree is
at most $2h+n(4(K+4))+2n$, where $h$ hubs and $n$ rows are processed.
The additional $2n$ counts edges to the row roots and their successor images.
\end{theorem}
\begin{proof}
\leanok
Apply the constructed hub factors and row covers to the exhaustive actual
neighbour lists. Retain each hub wheel's free row support and each row
wheel's single designated root. The physical graph is supported on the
free set, the row roots and their successor images. Partition each free
vertex's neighbours between the free set and those at most $2n$ labels.
The manuscript-specific profile substitutions and complete cover-down
assembly are still separate obligations.
\end{proof}

\begin{theorem}[Physical deletion and fixed-domain marked profiles]
\label{bp:marked-physical-deletion-profiles}
\lean{Oberwolfach.rawMarkedDeleteRims_rims, Oberwolfach.marked_row_physical_coordinate_degree, Oberwolfach.marked_neighbourhoods_after_physical_deletion}
\leanok
Deleting a physical rim graph of maximum degree $\Delta_R$ and a physical
spoke graph whose row and column degrees are at most $\Delta_V,\Delta_U$
changes the marked row error from $E_R$ to $E_R+h(\Delta_R+\Delta_U)$,
and the index error from $E_I$ to $E_I+h\Delta_V$. The original domains,
coefficients and comparison products are unchanged. Roots range over the
whole ambient sets. The result holds for every later permitted deletion.
\end{theorem}
\begin{proof}
\leanok
Compare the finite conjunctions defining the actual counts. An incoming
marked request at $v$ loses at most $d_D(\pi(v))$ rim choices; the other
requests lose at most the degree at their physical root. Charge deleted
spokes to their actual row or column. Sum the losses without any sampling
factor. No probability law or event is changed.
\end{proof}

\begin{theorem}[Deep endpoint classes and assignments from fixed profiles]
\label{bp:marked-deep-endpoint-profiles}
\lean{Oberwolfach.marked_endpoint_class_sizes, Oberwolfach.marked_deep_endpoint_profiles, Oberwolfach.marked_deep_row_assignment, Oberwolfach.marked_sample_row_cardinal_balances}
\leanok
Assume the displayed parent-degree, original-profile and physical-deletion
bounds, with $c\ge4$, marked spoke probability $t$ and ordinary spoke
probability $(c-1)t$. The admissible error is fixed before the separation
radius $d$. For sufficiently large ambient order, the actual modified
endpoint classes have the stated positive sizes, all five assignment
graphs have the proved degree and cut bounds, and the marked-outgoing
class satisfies the pairing gap. Local marked divisibility supplies the
three cardinal balances. The two compatible endpoint bijections then
exist in the actual modified graph.
\end{theorem}
\begin{proof}
\leanok
Single-arc estimates, with the shifted physical centre for an incoming
marked arc, and single-spoke estimates give the class sizes and gap.
Two- and three-spoke estimates give the degree and codegree bounds.
Apply the proved cut estimate and typed Hall construction. The endpoint
sets equal the actual occurrence neighbourhoods, so no new occurrence
identifiers or assumed endpoint matchings are used. This closes the deep
endpoint-assignment input; the complete cover-down and iteration callers
remain separate obligations.
\end{proof}

\begin{theorem}[Part restriction and deep-row numerical choices]
\label{bp:marked-deep-restriction-budget}
\lean{Oberwolfach.bipartite_block_after_part_and_edge_deletion, Oberwolfach.marked_neighbour_physical_deletion_loss, Oberwolfach.marked_deep_row_numeric_budget}
\leanok
Restricting the two classes of a bipartite block loses at most the number
of removed opposite vertices from each degree. Its original large-cut
bound survives. Subsequent bounded-degree edge deletion has the stated
additional loss. The deep-row free-use cap and permitted processed-root
fraction can be chosen before the separation radius; the final order
threshold absorbs its additive cost in the path-extension budget.
\end{theorem}
\begin{proof}
\leanok
Partition each old neighbourhood into retained and removed vertices.
Apply the actual edge-deletion cut bound. Choose the free-use cap first,
then the permitted fraction, and finally the order threshold. These are
explicit finite inequalities; their availability does not by itself
prove the full manuscript cover-down caller.
\end{proof}

\begin{lemma}[Successive marked row covers from the fixed profiles]
\label{bp:marked-sampled-row-covers}
\lean{Oberwolfach.marked_sample_row_covers_from_profiles,Oberwolfach.marked_sample_row_covers_after_hub_rims,Oberwolfach.marked_sample_row_covers_cover_star}
\leanok
For the current typed sampled host, the parent degree bounds and the fixed
row and index profiles imply the endpoint class sizes, the five bipartite
degree and cut estimates, and the path extension counts needed for successive
row covers. The tolerance, per-row bound and permitted root fraction are
chosen before the separation radius. The row covers have disjoint physical
resources and cover every physical edge incident with the specified
permutation-invariant row set. An initial hub rim graph supported on the
free set is avoided, with its degree loss included in the path inequality.
\end{lemma}
\begin{proof}
\uses{bp:marked-deep-endpoint-profiles,bp:marked-deep-restriction-budget}
\leanok
The buffered endpoint estimates absorb the cumulative degree losses.
The actual extension counts come from the fixed row and index profiles.
After deleting the initial hub rims, the proved successive-row construction
is applied and its wheels are lifted to the original occurrence types.
Exhaustion of the endpoint embeddings gives physical star coverage.
\end{proof}

\begin{lemma}[Physical deletion preserves retained sampled occurrences]
\label{bp:marked-sampled-deletion-occurrences}
\lean{Oberwolfach.markedSampleDeletedArcEquiv,Oberwolfach.markedSampleDeletedWheelToOriginal_resources}
\leanok
Deleting a set of physical rim coordinates identifies each remaining
sampled arc with its original directed typed occurrence. The resulting map
on marked wheels preserves its full set of physical resources.
\end{lemma}
\begin{proof}
\leanok
The two arc subtypes have equivalent defining membership conditions. The
wheel map keeps every row and hub and uses this equivalence on its arcs.
\end{proof}

\begin{theorem}[Sampled marked hub and row covers after physical deletions]
\label{bp:marked-sampled-deep-cover}
\lean{Oberwolfach.marked_sample_hub_and_row_covers,Oberwolfach.marked_sample_deep_cover_after_physical_deletion}
\leanok
The fixed parent degree estimates, row and index profiles, local
cycle/path counts and root-omitted spoke weights construct a packing
covering all edges incident with the specified deep rows and hubs.
The local balance and support conditions refer to the actual current
graph. Endpoint matchings, path counts, hub factors and hub class lower
bounds are derived. The resulting physical packing has the displayed
free rim and spoke degree bounds and avoids the protected inner pair.
One positive physical deletion allowance is chosen before the separation
radius and works for every later deletion satisfying its degree bounds.
\end{theorem}
\begin{proof}
\uses{bp:marked-sampled-row-covers,bp:marked-deep-endpoint-profiles}
\leanok
Construct the hub cover first. Its rim graph is supported on the free
row set and has degree at most twice the number of covered hubs. Apply
the successive-row construction after this rim deletion. The permitted
row hubs exclude the covered hubs, so the two packings are physically
disjoint. Root support gives the full free degree bound, including edges
to designated roots and their successor images. Fixed-domain deletion
estimates and the original root-omitted weights give one common numerical
allowance. This result does not yet discharge the preparation, complete
protected cover-down or iteration callers.
\end{proof}

\begin{lemma}[Fixed weighted sums and configuration counts after deletion]
\label{bp:marked-fixed-profiles-physical-deletion}
\lean{Oberwolfach.marked_local_spoke_loads_after_physical_deletion,Oberwolfach.marked_base_physical_rim_deletion_loss,Oberwolfach.marked_path_physical_rim_deletion_loss,Oberwolfach.marked_local_configurations_after_physical_deletion,Oberwolfach.marked_sample_hub_class_lower}
\leanok
Physical rim and spoke deletions have explicit deterministic losses for
the fixed omitted-root spoke sums and the actual ambient cycle and common
path counts. The path roots may be external to the free tuple and their
spokes are omitted. Hub class sizes follow from the one-spoke row profile.
\end{lemma}
\begin{proof}
\leanok
An incoming marked request is charged at its shifted physical endpoint.
For each lost configuration, charge a deleted physical edge with a free
endpoint or a deleted spoke class vertex. Sum over the unchanged tuple
family. The weighted estimate uses the unchanged parent coefficients
and normalization, with no sampling factor multiplying a realised loss.
\end{proof}

\begin{lemma}[Exact remaining occurrences and deep balances]
\label{bp:marked-exact-remaining-deep-balances}
\lean{Oberwolfach.markedSampleArcSetGraph_coordinate,Oberwolfach.marked_sample_reserved_star_packing_balances,Oberwolfach.marked_sample_cleared_star_neighbours}
\leanok
Physical deletion of the graphs of the reserved occurrence sets and the
prescribed-edge wheel packing removes exactly those original arcs and
spokes. The remaining sampled graph has all three row balance identities
at the specified deep rows and the required ordinary-to-marked spoke ratio
at the specified deep hubs. An empty inner star gives the neighbour
exclusions used by the deep construction.
\end{lemma}
\begin{proof}
\uses{bp:marked-sampled-deep-cover}
\leanok
The sampled arc map is physically injective and the canonical deleted
arc embedding preserves sources, auxiliary heads and types. Transfer the
finite incidence identities to these retained occurrences. At a deep
vertex the reserved star has the full incidences of the corrected inner
graph; subtracting the actual wheel packing preserves the balances.
The exact prescribed inner intersection clears the inner star, giving
both rim and spoke neighbour exclusions. Global divisibility of the
temporary star is not required.
\end{proof}

\begin{lemma}[Exact physical transport and fixed-domain changed edges]
\label{bp:marked-exact-physical-transport-domain-loss}
\lean{Oberwolfach.markedPhysicalDeletedSamplePacking_lift,Oberwolfach.markedPackingResourceEquality_graphs,Oberwolfach.marked_finite_prescribed_star_clearing,Oberwolfach.fixedDomainTuple_edge_loss,Oberwolfach.physical_deletion_invariant_domain_sum}
\leanok
Returning a wheel packing from the deleted sample to its original host
preserves its exact physical resource union, rim graph and spoke relation.
It avoids the deleted rim and spoke graphs. Changed-edge tuple counts
on a fixed free-coordinate domain use that domain's cardinality; fixed
roots may be external. The required degree bound counts neighbours in
the domain for every ambient root.
\end{lemma}
\begin{proof}
\leanok
Keep every wheel tuple and hub unchanged. The actual arc equivalence
preserves the physical endpoint of each marked closing arc. Exact finite
resource subtraction gives the cleared-star identity. For a changed-edge
count fix all but one free coordinate and sum its actual deleted
neighbours in the domain. An induced-subgraph bound alone is not used
for an external root.
\end{proof}

\begin{lemma}[Configurations and fixed weights on smaller domains]
\label{bp:marked-fixed-domain-deletion}
\lean{Oberwolfach.marked_local_configurations_after_physical_deletion_on_domain,Oberwolfach.marked_domain_spoke_loads_after_global_physical_deletion}
\leanok
The original fixed domain controls losses of marked cycle and two-root path
configurations under physical deletion. The same deletion transfers the
fixed potential-spoke weighted loads on each invariant domain, including
all tested roots for which the spoke is absent.
\end{lemma}
\begin{proof}
\uses{bp:marked-exact-physical-transport-domain-loss}
\leanok
For each changed edge, fix every other free coordinate and sum its deleted
neighbours in the original domain. Count all four root-incident edges of
the two-root path. Restriction commutes with actual physical deletion;
restricted graph degrees equal neighbour counts into the domain image.
Apply the fixed-weight deletion estimate to these degrees, retaining the
original omitted roots and normalization.
\end{proof}

\begin{lemma}[The two marked deep-stage packings from the actual correction]
\label{bp:marked-corrected-sample-deep-cover}
\lean{Oberwolfach.marked_corrected_sample_deep_cover}
\leanok
The fixed sampled neighbourhood profiles, original fixed weights and a
sparse correction of the inner graph construct both the prescribed-edge
packing and the packing covering the remaining deep rows and hubs.
The corrected inner graph, and the outer graph, satisfy the four marked
divisibility identities. The auxiliary constants precede the permutation
separation radius and the order threshold.
\end{lemma}
\begin{proof}
\uses{bp:marked-sampled-deep-cover,bp:marked-exact-remaining-deep-balances,bp:marked-exact-physical-transport-domain-loss}
\leanok
Derive extension counts after physical deletion from the original profiles
and the hub-size ratio. Reserve the star of the corrected inner graph.
Its physical degrees are at most the corresponding inner-part sizes, and
it is disjoint from the correction star. Pack the correction star, preserving
its exact original typed occurrences and inner intersection. The proved
rim, spoke-row and spoke-column degrees bound the next physical deletion.
Exact remaining incidences give the deep balances and neighbour exclusions.
Apply the sampled hub-and-row constructor. Its support proves avoidance of
every edge incident with the remaining inner rows and hubs, not only of
edges induced by the inner pair. The full active/boundary cover-down and
preparation callers are not asserted here.
\end{proof}

\begin{lemma}[Actual active and inner deletion costs]
\label{bp:marked-active-inner-deletion}
\lean{Oberwolfach.marked_active_total_deletion_bounds,Oberwolfach.marked_reserved_star_active_graphs,Oberwolfach.marked_inner_configurations_after_both_packings,Oberwolfach.marked_inner_spoke_loads_after_both_packings}
\leanok
After the prescribed-edge and deep packings, the remaining active host has
bounded physical rim and spoke deletion degrees. The temporary reserved
star is absent on the active pair. On the fixed inner row domain, the deep
packing changes no rim or spoke used by the configuration and potential-load
estimates, so only the prescribed-edge packing contributes to those losses.
\end{lemma}
\begin{proof}
\uses{bp:marked-corrected-sample-deep-cover,bp:marked-fixed-domain-deletion}
\leanok
Use the constructed deep packing's star avoidance to identify its neighbour
intersections with the remaining inner row and hub sets as empty, for every
ambient root. Outside those sets, apply its free-degree bounds. Restriction
removes every edge of the reserved star. The actual union degrees are
bounded by the sum of the two packing degrees. For inner counts and fixed
weighted sums, deletion of the deep packing has exactly zero effect on the
fixed inner domain; the previously checked physical-deletion bounds apply
with the original coefficients, omissions and normalization.
\end{proof}

\begin{lemma}[Construction of the protected inner correction and deep packings]
\label{bp:marked-constructed-inner-deep}
\lean{Oberwolfach.marked_protected_inner_correction,Oberwolfach.marked_induced_inner_correction,Oberwolfach.marked_constructed_inner_deep_cover}
\leanok
The actual four near-balance equations and cut bounds of the induced inner
host construct a sparse correction that avoids the deeper protected pair.
The corrected remainder satisfies all four marked divisibility equations.
With the fixed sampled profiles and weights, this constructs the prescribed
packing and the subsequent deep-row and deep-hub packing.
\end{lemma}
\begin{proof}
\uses{bp:marked-corrected-sample-deep-cover}
\leanok
Delete the protected rectangles from the allowed correction matrices;
row and column losses reduce the cut coefficient from $2h$ to $h$.
Apply the proved matrix correction and select its actual typed occurrences.
Lift the induced correction to the original identifiers, preserving each
incidence, spoke type and divisibility identity. The auxiliary incidence
bound is $8(q+1)\xi |I|$ and the spoke row and column bounds are
$4(q+1)\xi |I|$, including all ambient roots. Choose $\xi$ before the
separation radius to meet both the prescribed inner allowance and the
subsequent packing allowance. The correction itself need not be divisible.
\end{proof}

\begin{lemma}[Inner factor host and external boundary roots]
\label{bp:marked-inner-boundary-selection}
\lean{Oberwolfach.marked_fresh_inner_boundary_selection}
\leanok
The fresh mixed factor and endpoint selection remains valid when the
host supplying its cycles and endpoint relations differs from the host
supplying boundary supports. Both use the same row and hub types and the
same physical coin law; the inner host may omit the external boundary roots.
The stated count, weight, five-block and later-change hypotheses are retained.
\end{lemma}
\begin{proof}
\leanok
Intersect the common factor event for the inner host with the support event
for the boundary host on the same physical-coordinate probability space.
All later masks and class choices stay inside this one event. The actual
boundary class inclusions give the selector's overlap and membership bounds.
The input verification for the full cover-down caller remains separate.
\end{proof}

\begin{lemma}[Inner factors in the ambient physical host]
\label{bp:marked-inner-factor-embedding}
\lean{Oberwolfach.permutationClose_embedding_iff,Oberwolfach.hubCycleEmbeddingEquiv,Oberwolfach.hubCycleEmbedding_weighted_degree,Oberwolfach.hubFactorEmbedding_rim,Oberwolfach.marked_hub_factor_host_lift}
\leanok
An injective row map commuting with the successor permutations preserves
separation. It transports hub cycles bijectively to their image host,
preserves weighted degrees at image vertices, and sends the factor's
physical rim graph to its exact graph image. The lift to a larger typed
host preserves the physical graph.
\end{lemma}
\begin{proof}
\leanok
Commute each permutation power with the row embedding. Every image-host
hub cycle lies in the row image, giving an inverse on its coordinates.
Transport the same weights through this bijection. Physical rim edges are
pairs of a coordinate and its ordinary or successor-shifted head; hence
the finite edge sets and their graphs commute with the embedding. Lift
actual occurrence witnesses to the larger host, keeping cycle coordinates.
\end{proof}

\begin{lemma}[Boundary selection with the original inner tolerance]
\label{bp:marked-local-boundary-domain}
\lean{Oberwolfach.marked_embedded_fresh_hub_common_event,Oberwolfach.marked_local_domain_boundary_selection,Oberwolfach.marked_local_boundary_selection_in_host}
\leanok
The fresh hub-factor tolerance and allowed relative class change are chosen
on the original inner row domain before its embedding into the ambient
host. One event in the original ambient physical coin law supplies the
factors and five-block endpoint assignments for every permitted later mask,
class choice and forbidden graph. The selected objects belong to the actual
ambient typed host and retain their physical disjointness and degree bound.
\end{lemma}
\begin{proof}
\uses{bp:marked-inner-factor-embedding,bp:marked-inner-boundary-selection}
\leanok
The injection on unordered physical edges pulls the ambient product law
back to exactly the inner product law. Restrict an ambient forbidden graph
to the inner rows, construct the factor there and embed it back. Intersect
this event with the original boundary-support event; use actual boundary
class inclusions and all five degree/cut blocks in the mixed selector.
Lift the chosen factors and endpoint assignments to the actual host without
changing their graphs. The error tolerance has no dependence on the later
inner/ambient size ratio. The actual five-block and post-packing caller
hypotheses must still be verified in the complete cover-down construction.
\end{proof}

\begin{lemma}[Boundary selection from the original sample]
\label{bp:marked-sample-boundary-inputs}
\lean{Oberwolfach.marked_sample_local_boundary_selection}
\leanok
The actual inner row profile, configuration estimates and fixed potential-spoke loads supply the fresh factor inputs. Sampled inner arc occurrences embed into the actual ambient host, preserving endpoints and physical edges. One common ambient boundary event supplies factors and endpoint pairings for all later classes satisfying the displayed degree, cut and change bounds.
\uses{bp:marked-local-boundary-domain}
\end{lemma}
\begin{proof}
\leanok
Derive the hub-class lower bound and root-omitted degree estimate, apply the original-inner-domain factor event, intersect it with the ambient support event, and lift the selected cycles along the actual sample arc embedding. Apply mixed selection with all five endpoint blocks.
\end{proof}

\begin{lemma}[Five endpoint profiles and factors on one event]
\label{bp:marked-boundary-profile-selection}
\lean{Oberwolfach.marked_sample_profile_boundary_selection}
\leanok
The original row and index profiles imply all five endpoint degree and cut estimates on the same boundary sampling event as the inner factors. The estimates hold for every later permitted change of the five shared classes. The factor weights and inner normalization are unchanged.
\uses{bp:marked-sample-boundary-inputs}
\end{lemma}
\begin{proof}
\leanok
Intersect the two common events and apply the five endpoint estimates to the actual shared classes. Their cut density is twice the density supplied to the mixed selector. Sum the failure probabilities; independence between these events is not required.
\end{proof}
\begin{lemma}[Actual marked boundary remainder changes]
\label{bp:marked-boundary-remainder-changes}
\lean{Oberwolfach.marked_sample_boundary_classes_after_packing, Oberwolfach.marked_sample_boundary_remainder_common_event}
\leanok
Let $G$ be the actual physical boundary graph, $R$ its independent reserved sample, $C$ a later packing disjoint from $R$, and $D$ a subsequent covering. If $\Delta((G-R)-C)\le\ell$ and $\Delta(D)\le\delta$, each fixed endpoint or spoke class changes by at most $\ell+\delta$ when replaced by its class in $(G-C)-D$. Each auxiliary-root mask changes at most $2(\ell+\delta)$ times. One original sampling event controls the support, overlap and membership bounds for every later such pair $C,D$ satisfying the numerical allowance.
\uses{bp:marked-boundary-profile-selection}
\end{lemma}
\begin{proof}
\leanok
A changed physical edge belongs to the approximate leftover or to $D$. Injectivity of the actual physical edge map bounds each incidence count. Incoming marked edges are charged at the image of the auxiliary root under the permutation. Permutation invariance of the inner row set proves that the exact physical boundary graph contains every auxiliary crossing edge.
\end{proof}

\begin{lemma}[Marked boundary completion after the earlier packings]
\label{bp:marked-boundary-prior-packings}
Let the earlier approximate packing avoid all resources inside the inner pair, and let the subsequent covering have the displayed small physical maximum degree. The selected boundary factors and endpoint assignments produce a packing disjoint from both. Their union has the exact divisible remainder supported on the inner pair. The prescribed arcs and spokes must be the actual remaining incidence lists. No small global degree assumption is imposed on the approximate packing.
\lean{Oberwolfach.marked_selected_boundary_after_packings}
\leanok
\end{lemma}
\begin{proof}
The strengthened endpoint constructor classifies every rim edge as inner or one of its prescribed boundary arcs. Together with the row and hub supports, this proves avoidance of the approximate packing. Physical restriction excludes the subsequent covering while retaining all occurrence identities. The finite divisibility and exact physical remainder follow by subtracting the constructed wheel packing. The hypotheses supplying the earlier packings and selected assignments remain separate obligations in the full cover-down caller.
\leanok
\end{proof}

\begin{lemma}[Balances and total loss of the remaining marked boundary]
\label{bp:marked-remaining-boundary-balances}
For the actual sample after physical deletion, finite remainder divisibility implies the three endpoint-class balances and the hub ratio. The endpoint lists embed into the original sample, preserve both endpoints and arc types, and enumerate exactly the surviving incident arcs. The selected boundary construction after the earlier two packings gives an inner divisible remainder whose combined rim degrees, spoke rows, and inner spoke columns obey the stated loss bound. Its reservation parameter is chosen before the separation radius.
\lean{Oberwolfach.marked_remaining_boundary_row_balances, Oberwolfach.marked_remaining_boundary_hub_ratio, Oberwolfach.marked_deleted_endpoint_embedding_exhaustive, Oberwolfach.marked_boundary_total_inner_loss, Oberwolfach.marked_mixed_boundary_finish_after_packings}
\leanok
\end{lemma}
\begin{proof}
The deleted-sample embeddings transfer finite incidence balances to the actual endpoint classes. Restriction to the inner row domain gives the hub equation. For the combined loss, the first packing contributes no inner resources; the second packing is charged by its physical maximum degree, and the constructed boundary packing by its proved local bounds. The numerical completion proof derives its degree and path-count budgets before invoking this calculation. Construction of the earlier packings and the common event supplying the selection assumptions remain obligations of the full caller.
\leanok
\end{proof}

\begin{lemma}[Physical boundary retention and inner-domain spoke covering]
\label{bp:marked-inner-boundary-sampling}
On one event for independent physical-edge coins, all fixed marked one-arc and two-arc extension counts on the inner set retain the displayed lower bound, proportional to $p^3$. The event includes all ambient root tuples. These bounds transfer to the host containing all inner resources and the retained boundary. The local extension construction then packs prescribed spokes with all free row vertices in the inner set, with the displayed row, rim-degree and hub-incidence bounds.
\lean{Oberwolfach.retained_pattern_lower_common_event, Oberwolfach.marked_boundary_extension_retention_event, Oberwolfach.marked_retained_boundary_free_counts, Oberwolfach.marked_inner_extension_counts_after_rim_deletion, Oberwolfach.marked_inner_prescribed_spoke_wheel_packing}
\leanok
\end{lemma}
\begin{proof}
Each candidate reads at most three physical coins and each coin affects at most three candidates. The expected count is at least $p^3$ times its fixed parent count; bounded differences and a polynomial union bound give one common event. Typed occurrence restrictions preserve all endpoints and marks. A path is constructed using spoke conditions restricted to the inner set, so no loss for the complement of that set is incurred. Physical rim deletion costs at most twice its maximum degree in the local intersections. Greedy packing then excludes saturated inner rows and the already used spokes at the actual prescribed hub. The earlier approximate packing, rim-edge cover, parameter substitutions and final cover-down caller remain separate obligations.
\leanok
\end{proof}

\begin{lemma}[Prescribed outside edges through the inner pair]
\label{bp:marked-outside-cover-construction}
\lean{Oberwolfach.marked_wheel_from_two_root_path,Oberwolfach.marked_rim_wheel_on_inner_domain,Oberwolfach.marked_inner_rim_wheel_after_deletion,Oberwolfach.marked_inner_prescribed_rim_wheel_packing,Oberwolfach.marked_inner_prescribed_edge_packing,Oberwolfach.marked_inner_domain_packing_numeric_budgets,Oberwolfach.marked_inner_domain_order_threshold,Oberwolfach.marked_sparse_inner_domain_edge_packing,Oberwolfach.marked_sparse_inner_domain_packing,Oberwolfach.marked_sparse_outside_subgraph_packing}
\leanok
Inner-domain extension counts and compatible inner hubs construct a physical wheel packing covering the actual prescribed outside arcs and spokes. The rim stage uses inner hubs, and the spoke stage uses outside hubs after deletion of the actual first-stage rim graph. Their outside intersection is exactly the prescribed subgraph. Both stages and their union satisfy explicit rim, spoke-row and spoke-column bounds. The required tolerance inequality is explicit and precedes the separation radius.
\end{lemma}
\begin{proof}
\leanok
Construct the rooted cycles, keep the physical deletion losses at their actual roots, cap repeated inner-row and hub use, and combine the two families. Enumerating the supplied finite edge sets preserves their actual incidences and occurrence identifiers.
\end{proof}

\begin{lemma}[One sampled event for later outside covers]
\label{bp:marked-sampled-outside-cover}
\lean{Oberwolfach.marked_boundary_hub_coins_size,Oberwolfach.marked_boundary_hub_coins_influence,Oberwolfach.marked_boundary_hub_retention_event,Oberwolfach.marked_retained_cover_free_counts,Oberwolfach.marked_retained_cover_hub_counts,Oberwolfach.marked_boundary_cover_common_event,Oberwolfach.marked_boundary_cover_scale_choices,Oberwolfach.marked_sampled_outside_cover,Oberwolfach.marked_cover_resource_outside_if_forbidden,Oberwolfach.marked_cover_packing_avoids_earlier}
\leanok
The original physical-edge sample simultaneously supplies all path extensions and compatible hubs. On that single event, every later sparse prescribed outside subgraph has the preceding packing in the retained cover host. The all-root extension bound has coefficient $p^3$; choosing the free-use scale $p^4$ permits a $p^9$ remainder. Avoidance of an earlier approximate packing uses its exclusion from the inner pair and retained boundary and exact outside coverage, without a small global degree hypothesis on that earlier packing.
\end{lemma}
\begin{proof}
\leanok
Each candidate reads at most three physical coins, with reverse incidence at most three. Apply the proved bounded-difference estimate and join the path and hub events under the same original law. Transfer their counts to the actual restricted typed host, then apply the deterministic construction. The earlier packing is excluded by classifying each physical resource, rather than charging its global degrees.
\end{proof}

These statements do not yet construct the earlier approximate packing, discharge the complete cover-down caller, or prove Theorem~1.1. Its release gate remains closed.

\begin{lemma}[The original outside remainder and actual physical degrees]
\label{bp:marked-original-outside-cover}
\lean{Oberwolfach.marked_outside_arc_restriction,Oberwolfach.marked_sampled_outside_cover_original,Oberwolfach.ordinary_physical_degree_row,Oberwolfach.ordinary_physical_degree_column,Oberwolfach.ordinary_physical_degree_bound,Oberwolfach.marked_sampled_small_outside_cover,Oberwolfach.marked_packing_rim_pair_adj,Oberwolfach.marked_approximate_packing_reservation_exclusions}
\leanok
The sampled outside cover lifts to the original typed host, preserving occurrence identifiers, directions, types, physical resources and degree bounds. At the explicit $p^4$ usage scale, a later prescribed outside subgraph with degrees at most $p^9n$ has a covering packing of physical maximum degree at most both $p^3n$ and $\nu p^3n$, for sufficiently small $p$ chosen before the separation radius. The original approximate packing's actual exclusion from the inner pair and retained boundary supplies the required spatial avoidance.
\end{lemma}
\begin{proof}
\leanok
Restrict and lift the actual finite arc set, preserving every endpoint and type. Convert the separate rim, spoke-row and spoke-column estimates to exact physical graph degrees. Classify the earlier packing's resources using the original boundary graph and coin values.
\end{proof}

\begin{lemma}[Balances after inner and boundary reservation]
\label{bp:boundary-reservation-balances}
\lean{Oberwolfach.retained_coin_count_tail,Oberwolfach.retained_coin_counts_common_event,Oberwolfach.reservation_remaining_outer_count,Oberwolfach.reservation_remaining_inner_count,Oberwolfach.weighted_pointwise_error,Oberwolfach.reservation_remaining_balance,Oberwolfach.boundary_reservation_balance_common_event}
\leanok
For fixed full and inner incidence sets and fixed signed coefficients, the balance after reservation is concentrated about $B_0-pB_1$ at an outside root and $(1-p)(B_0-B_1)$ at an inner root. One event controls every supplied root and test under the original physical-coordinate law. Its order threshold is uniform over the retention probability.
\end{lemma}
\begin{proof}
\leanok
Count the actual set differences, apply the singleton-coordinate concentration theorem, and sum the individual errors with the absolute values of their fixed coefficients. This identity still has to be instantiated for the complete approximate-packing caller.
\end{proof}

\begin{lemma}[Outside completion from the restricted-host remainder]
\label{bp:marked-restricted-outside-completion}
\lean{Oberwolfach.marked_incident_remainder_filter,Oberwolfach.marked_unused_incidence_bounds,Oberwolfach.marked_restricted_unused_card,Oberwolfach.marked_restricted_unused_bounds,Oberwolfach.marked_boundary_rim_allowed_physical_congr,Oberwolfach.marked_reservation_complement_excludes,Oberwolfach.marked_sampled_outside_completion}
\leanok
On the original sampling event, every later packing in the host complementary to the reserved inner and boundary resources, with auxiliary leftover degrees at most $p^9n$, admits a physically disjoint outside completion. The completion covers every remaining outside rim and spoke resource and has physical maximum degree at most both $p^3n$ and $\nu p^3n$. No bound on the earlier packing's total degree is required.
\end{lemma}
\begin{proof}
\leanok
Identify the restricted unused resources with retained original occurrences. Form the actual outside remainder and derive its sizes and incidence bounds. Apply the sampled outside constructor and prove physical avoidance using the reserved-resource exclusions and permutation invariance of the inner row set.
\end{proof}

\begin{lemma}[Actual sampled typed-host balances]
\label{bp:marked-actual-reserved-balances}
\lean{Oberwolfach.incident_reservation_image,Oberwolfach.incident_reservation_balance_common_event,Oberwolfach.incident_reservation_count_common_event,Oberwolfach.markedBalanceTermSet_incident,Oberwolfach.markedBalanceTermSet_card,Oberwolfach.marked_active_reservation_complement,Oberwolfach.marked_balance_reservation_common_event,Oberwolfach.markedBalanceTermSet_row_degrees,Oberwolfach.markedBalanceTermSet_hub_degrees,Oberwolfach.markedBalanceTerm_toLarger,Oberwolfach.markedBalanceTermSet_restriction_card,Oberwolfach.marked_reservation_balance_errors,Oberwolfach.markedSixDegreeBalance_row,Oberwolfach.markedSixDegreeBalance_hub,Oberwolfach.marked_reserved_host_balance_event,Oberwolfach.marked_balance_inner_term_restriction}
\leanok
The six actual degree terms of the complementary typed host satisfy their four near-balance inequalities on one original-coordinate event, uniformly over the retention probability. If the original and inner-neighbour balances are bounded by $\eta$ and $\eta'$, and each degree count has error at most $\varepsilon|V\cup U|$, the resulting error is at most $\eta+\eta'+(m+4)\varepsilon|V\cup U|$. The hub coefficient is the original $m$. The inner-neighbour term retains all ambient roots and restricts only the opposite endpoint.
\end{lemma}
\begin{proof}
\leanok
Use the exact original physical coin map and injectivity on each incidence set. Identify the six counts with outgoing, incoming and spoke degrees in the actual restricted host. Sum their errors with the absolute values of the original four coefficient vectors. The formula at an outside root is $B_0-pB_1$, and at an inner root it is $(1-p)(B_0-B_1)$.
\end{proof}

\begin{lemma}[Canonical common paths after deletion]
\label{bp:marked-canonical-path-restriction}
\lean{Oberwolfach.marked_canonical_path_count_restriction}
\leanok
A realized deletion of at most $d$ incidences at each relevant root removes at most $4(m+1)dn^{m-1}$ canonical common paths. The four root spokes are omitted from the path count. If the original count is at least $2\beta n^m$ and $4(m+1)d\le\beta n$, at least $\beta n^m$ paths remain.
\end{lemma}
\begin{proof}
\leanok
Apply the established common-path deletion estimate on the unchanged original vertex domain and multiply its degree budget by $n^{m-1}$.
\end{proof}

\begin{lemma}[Packing after a reservation]
\label{bp:marked-restricted-near-packing}
\lean{Oberwolfach.bipartiteLargeCut_coefficient_mono,Oberwolfach.bipartiteLargeCut_subrelation,Oberwolfach.occurrence_relation_restriction_change,Oberwolfach.occurrence_bipartiteLargeCut_restriction,Oberwolfach.occurrence_cut_restriction_loss,Oberwolfach.marked_restricted_near_balance_packing}
\leanok
The marked near-balance packing theorem applies to the actual restricted host after a sufficiently small reservation. The admissible reservation degree is fixed before the desired uncovered-degree tolerance. The four original cut bounds, the ordinary two-step counts, common cycle and root-omitted path counts, spoke profiles and fixed original weights supply the corresponding hypotheses in the restricted host. Only its actual four near-balance equations remain to be supplied by the reservation event.
\end{lemma}
\begin{proof}
\leanok
Bound every cut loss by the realized row and column deletions. Restrict the original weight function and preserve its two-sided bounds; the load change is charged on the original vertex scale. Retain the supplied spoke profiles and exceptional-pair bound. Apply the proved near-balance theorem in the restricted host. Reserved edges are not added to its uncovered-edge count.
\end{proof}

\begin{lemma}[The actual sampled complementary packing]
\label{bp:marked-sampled-complement-approximation}
\lean{Oberwolfach.incident_opposite_endpoint_bound,Oberwolfach.marked_balance_inner_term_card,Oberwolfach.reservation_count_loss_upper,Oberwolfach.marked_reservation_deleted_count_event,Oberwolfach.relation_row_restriction_loss,Oberwolfach.marked_deleted_term_degree_bounds,Oberwolfach.marked_sampled_complement_approximation}
\leanok
Under the displayed parent-host and inner-neighbour balance hypotheses, the original boundary sample admits an approximate wheel packing in its complementary typed host, with both auxiliary and physical leftover degrees at most $\tau^9 n$. The inner-pair size allowance precedes $\tau$; the balance accuracy follows $\tau$ and precedes the separation radius and order threshold. One event has probability at least $1-2/|V\cup U|^q$.
\end{lemma}
\begin{proof}
\leanok
An incident occurrence with opposite endpoint in the inner pair is encoded by its type, orientation and that endpoint. This bounds its deterministic inner loss. The original six-term sampling event supplies the remaining degree loss and the four actual balance errors. Intersect these events without assuming independence, then apply the restricted-host packing theorem with uncovered-degree tolerance $\tau^9$. No approximate-packing existence hypothesis is inserted.
\end{proof}

\begin{lemma}[Constructed approximate packing and outside cover]
\label{bp:marked-actual-outside-boundary}
\lean{Oberwolfach.marked_complement_packing_avoids_inner,Oberwolfach.markedPairResource,Oberwolfach.graph_degree_le_marked_incident,Oberwolfach.graph_degree_le_unused_marked_incident,Oberwolfach.marked_complement_boundary_remainder,Oberwolfach.marked_sampled_approximate_outside_cover}
\leanok
The original physical boundary sample constructs an approximate packing in its complement and a disjoint outside cover. Both use wheel parameter $l+4$. The first packing avoids the inner pair and the retained boundary. Its unretained boundary remainder has degree at most $\tau^9 n$. The outside cover covers every original outside resource and has physical degree at most $\tau^3 n$ and $\nu\tau^3 n$, with the stated probability and parameter order.
\end{lemma}
\begin{proof}
\uses{bp:marked-sampled-complement-approximation}
\leanok
Intersect the two original-law events. Lift the actual complementary packing by retaining all occurrence identifiers. Its resource set is unchanged. Each neighbour in the unretained boundary remainder determines a distinct unused physical resource of the restricted host. The permutation-invariant inner set equates physical and auxiliary rim crossings. Apply the proved outside-completion construction and its exclusion theorem.
\end{proof}

\begin{lemma}[Actual remaining boundary classes]
\label{bp:marked-actual-boundary-classes}
\lean{Oberwolfach.marked_outside_packing_remainder,Oberwolfach.marked_outside_packing_boundary_balances,Oberwolfach.marked_remaining_endpoint_class,Oberwolfach.marked_remaining_spoke_class,Oberwolfach.marked_remaining_spoke_column,Oberwolfach.marked_actual_boundary_class_changes}
\leanok
Subtracting an actual outside wheel packing preserves all four divisibility identities. Its remaining endpoint classes have the three exact row balances, and each outside hub has the required ordinary-to-marked ratio. The classes obtained by deleting the actual physical graphs equal the classes in the remaining boundary graph. Their changes from the original reserved classes are bounded by the sum of the actual complementary leftover and subsequent covering degrees.
\end{lemma}
\begin{proof}
\leanok
Translate physical membership back to the original typed arcs and spoke pairs before subtracting their integer incidence counts. Outside coverage forces remaining neighbours of every outside vertex into the inner pair. The physical endpoint coin retains the marked-head permutation. Exact class identities and the fixed-domain change theorem then give the stated bounds, without multiplying the realized deletion loss by a sampling probability.
\end{proof}

\begin{lemma}[Numerical and sampled inputs to the marked boundary selector]
\label{bp:marked-robust-boundary-selector}
\lean{Oberwolfach.physical_boundary_class_size_event,Oberwolfach.marked_boundary_class_size_margins,Oberwolfach.marked_boundary_class_margins_event,Oberwolfach.marked_boundary_selector_numerics,Oberwolfach.marked_sample_quantitative_boundary_selection,Oberwolfach.marked_sample_robust_boundary_selection}
\leanok
The fixed parent endpoint and spoke classes supply the actual class-size and pairing-gap estimates on one original physical sampling event, uniformly over subsequent changes bounded by $(p^9+p^3)n$. The numerical degree, separation, rounding and deletion inequalities required by the mixed selector follow for sufficiently small $p$ and sufficiently large order. The selector accuracy is fixed before its deletion allowance; no reverse dependence is used.
\end{lemma}
\begin{proof}
\leanok
Apply concentration to the fixed physical-coordinate images. For every later class, its cardinality differs from that of the retained parent class by at most the cardinality of their symmetric difference. Choose the accuracy from the parent lower bounds and gap, then choose the reservation density after the selector's deletion allowance. Explicit powers of $p$ and the order threshold supply the remaining numerical inequalities. Intersect the events under the unchanged physical-coin law.
\end{proof}

\begin{lemma}[Exact support and balance transport for the two outside packings]
\label{bp:marked-union-boundary-balances}
\lean{Oberwolfach.marked_endpoint_class_boundary_support,Oberwolfach.marked_deleted_endpoint_subset,Oberwolfach.marked_spoke_class_boundary_support,Oberwolfach.embedded_filtered_class_image,Oberwolfach.embedded_filtered_class_card,Oberwolfach.embedded_filtered_class_change_card,Oberwolfach.embedded_boundary_hub_class_image,Oberwolfach.marked_packing_union_is_matching,Oberwolfach.markedPackingPhysicalGraph_union,Oberwolfach.marked_union_packing_boundary_balances}
\leanok
The union of the two physically disjoint outside packings has exactly the union of their physical graphs. Its actual remaining classes satisfy the three row balances and the hub ratio. Passage to a fixed inner vertex set preserves class cardinalities and symmetric differences. Endpoint support uses each original directed occurrence and its physical marked-head coordinate.
\end{lemma}
\begin{proof}
\uses{bp:marked-actual-boundary-classes}
\leanok
Apply the remaining-boundary balance identities to the actual union packing, and prove the rim and spoke graph equalities coordinate by coordinate. Injectivity of the fixed inner embedding preserves the two class cardinalities and their symmetric difference. Retain the original arc witnesses when checking boundary support.
\end{proof}

\begin{lemma}[Selection and completion on the actual remaining boundary]
\label{bp:marked-actual-boundary-completion}
\lean{Oberwolfach.marked_actual_boundary_support,Oberwolfach.marked_sample_actual_boundary_selection,Oberwolfach.marked_packing_arc_coordinate_iff,Oberwolfach.marked_deleted_endpoint_avoids_packing,Oberwolfach.marked_deleted_endpoint_exhausts_packing,Oberwolfach.marked_packing_spoke_adj_iff,Oberwolfach.marked_remaining_spokes_exact,Oberwolfach.marked_actual_boundary_finish,Oberwolfach.MarkedHubFactorOn.changeSet,Oberwolfach.MarkedHubFactorOn.changeSet_graph,Oberwolfach.MarkedEndpointPairing.changeSets,Oberwolfach.MarkedEndpointPairing.changeSets_graph,Oberwolfach.marked_sample_actual_boundary_completion}
\leanok
The two actual outside packings determine the endpoint classes, hub classes, their three row balances, the hub ratio and the remaining boundary mask. The original sampling event supplies a mixed selection on these classes. Its boundary completion avoids both earlier packings and leaves an exactly divisible remainder supported on the inner pair. The sum of all three packings has the prescribed inner rim, spoke-row and spoke-column degree bounds.
\end{lemma}
\begin{proof}
\uses{bp:marked-robust-boundary-selector,bp:marked-union-boundary-balances,bp:marked-actual-outside-boundary}
\leanok
Apply the class-change estimates to the actual physical graphs. Restrict hub classes through the fixed embedding, preserving cardinalities and symmetric differences. The original typed surviving arcs are enumerated by explicit incidence embeddings. Their physical coordinates and every remaining root spoke avoid the union packing. Outside coverage forces their other endpoints into the inner pair. Equal class representations transport the selected cycles and bijections without changing their physical graphs. The proved boundary finishing theorem then applies, with both its deletion and total-loss inequalities checked explicitly.
\end{proof}

\begin{lemma}[Constructed active marked cover-down]
\label{bp:marked-active-cover-down}
\lean{Oberwolfach.marked_sampled_approximate_outside_cover_ordered,Oberwolfach.marked_active_cover_down}
\leanok
Under the displayed parent degree, cut, weight, configuration, inner-domain profile and divisibility hypotheses, an actual physical marked-wheel packing covers every resource outside the fixed inner pair and leaves an exactly divisible inner remainder. Its used inner rim, spoke-row and spoke-column degrees satisfy the stated loss bound. The allowed inner-set size is fixed before the extension density on that set; the reservation density precedes the near-balance accuracy and the graph-order threshold.
\end{lemma}
\begin{proof}
\uses{bp:marked-actual-boundary-completion,bp:marked-actual-outside-boundary}
\leanok
The earlier approximate packing and the outside cover are constructed, rather than assumed. Intersect their event with the boundary-completion event under the same physical-coin law. Convert vertex-domain degree bounds to ambient-domain bounds using the actual cardinality inequality. For exponent one and sufficiently large inner order the common event has positive probability, so it supplies all three packings simultaneously. Their union has the exact asserted remainder. Preparation must still instantiate the displayed hypotheses, and protected restoration and iteration remain separate callers.
\end{proof}

\begin{lemma}[Outside balance after an inner reservation]
\label{bp:marked-outside-balance-reservation}
\lean{Oberwolfach.MarkedOutsideBalance,Oberwolfach.MarkedRemainderDivisible.outside,Oberwolfach.MarkedOutsideBalance.sdiff,Oberwolfach.MarkedOutsideBalance.remove_inner,Oberwolfach.marked_reserved_outside_balance,Oberwolfach.marked_outside_balance_vertex_pullback}
\leanok
The three row identities and the hub identity outside a fixed inner pair survive deletion of any edges wholly inside that pair and removal of a contained physical wheel packing. An injective change of row and hub domains preserves these identities when the retained spoke set is the exact image and its types agree.
\end{lemma}
\begin{proof}
\leanok
An inner edge contributes zero to every tested outside incidence. Subtract the exact finite incidence counts of the wheel packing; each wheel has all four balance identities. Injectivity identifies every filter cardinality under the change of domains. No balance at an inner vertex, and no divisibility of the reserved star, is inferred.
\end{proof}

\begin{lemma}[Marked active cover-down with only outside balance]
\label{bp:marked-active-cover-down-local}
\lean{Oberwolfach.marked_boundary_remainder_support,Oberwolfach.marked_boundary_remainder_after_packing_support,Oberwolfach.marked_selected_boundary_after_packings_support,Oberwolfach.marked_mixed_boundary_finish_after_packings_support,Oberwolfach.marked_actual_boundary_finish_support,Oberwolfach.marked_remaining_boundary_row_balances_local,Oberwolfach.marked_outside_packing_remainder_local,Oberwolfach.marked_outside_packing_boundary_balances_local,Oberwolfach.marked_union_packing_boundary_balances_local,Oberwolfach.marked_sample_actual_boundary_selection_local,Oberwolfach.marked_sample_actual_boundary_completion_local,Oberwolfach.marked_active_cover_down_local}
\leanok
The active cover-down constructor requires exact balance only outside its inner pair, in addition to its unchanged parent profile, weight, configuration, near-balance and numerical inputs. It constructs all three physical wheel packings, covers every edge outside the inner pair, and proves the total inner rim, spoke-row and spoke-column loss bounds. The remainder is exactly the original physical resource set minus these whole wheels. Global divisibility of this temporary active host is not asserted.
\end{lemma}
\begin{proof}
\uses{bp:marked-outside-balance-reservation,bp:marked-active-cover-down,bp:marked-robust-boundary-selector}
\leanok
The actual preceding packings preserve the outside identities. These yield precisely the three endpoint balances and the hub ratio used by the selector. Neither the actual path construction nor the support and complement conclusions require global divisibility. Repeat the original common-event construction with those local inputs, keeping all coefficient functions, sampling laws, root domains and probability bounds unchanged. Global divisibility in the protected caller is instead proved after restoring its reservation and subtracting only whole wheels.
\end{proof}

\begin{lemma}[Exact embedding of active sample packings]
\label{bp:marked-active-sample-packing-lift}
\lean{Oberwolfach.markedResourceEmbedding,Oberwolfach.markedWheelVertexEmbedding,Oberwolfach.markedWheelVertexEmbedding_resources,Oberwolfach.markedPackingVertexEmbedding_resources,Oberwolfach.markedPackingVertexEmbedding_matching,Oberwolfach.markedResourceEmbedding_inner_iff,Oberwolfach.markedResourceEmbedding_inner_filter,Oberwolfach.markedResourceEmbedding_avoids_star,Oberwolfach.markedResourceEmbedding_active_inner_filter,Oberwolfach.marked_active_deleted_packing_lift,Oberwolfach.markedRestrictedSampleArcEmbedding,Oberwolfach.markedRestrictedSampleArcEmbedding_range}
\leanok
An active packing in the actual physically deleted restricted sample embeds into the original sample with exactly the image of its physical resources. It avoids both prescribed deletion graphs. Its intersection with the original inner pair equals its intersection with the active inner pair, and it avoids the protected star. The restricted sample contains exactly the original typed arcs whose two auxiliary endpoints belong to the active row domain.
\end{lemma}
\begin{proof}
\leanok
The row embedding commutes with the successor permutation, preserving separation and the marked physical closing edge. The two arc maps preserve their endpoints and types. The resource map is injective, so matching is preserved. The physical-deletion lift preserves the tuple, hub and resources and proves actual exclusion from both deletion graphs. The inner and star identities follow by taking preimages of the specified finite vertex sets.
\end{proof}

\begin{lemma}[Outside balance in the actual cleared active sample]
\label{bp:marked-cleared-active-balance}
\lean{Oberwolfach.marked_outside_balance_arc_pullback,Oberwolfach.MarkedOutsideBalance.congr_spoke_types,Oberwolfach.marked_deleted_sample_outside_balance,Oberwolfach.marked_restricted_sample_outside_balance,Oberwolfach.marked_sample_packing_deletion_outside_balance,Oberwolfach.marked_reserved_sample_outside_balance,Oberwolfach.marked_star_cover_deleted_support,Oberwolfach.marked_star_cover_active_balance}
\leanok
Physical deletion transports the outside balance identities through the exact surviving typed occurrence sets. Covering the entire protected star leaves every surviving auxiliary arc and spoke endpoint in the active row and hub domains. Restricting to those exact domains preserves the outside balance identities, the original arc directions and marks, and the original outcome.
\end{lemma}
\begin{proof}
\uses{bp:marked-outside-balance-reservation,bp:marked-active-sample-packing-lift}
\leanok
The deleted-arc embedding has precisely the undeleted physical coordinates as its range. On surviving spokes the original and deleted marked predicates agree. Whole wheels have the four exact finite balance identities. The protected-star coverage and permutation invariance force both auxiliary endpoints of every surviving rim arc outside the protected rows. The remaining spokes avoid the protected rows and hubs. Exact range equalities for the active embeddings identify all relevant finite filters and their cardinalities.
\end{proof}

\begin{lemma}[Constructed correction and deep cover supply active balance]
\label{bp:marked-constructed-deep-active-balance}
\lean{Oberwolfach.marked_constructed_deep_active_balance}
\leanok
Under the existing parent profiles, fixed weights, inner near-balance equations, inner cuts and numerical hypotheses, construct the inner correction, prescribed-edge packing and protected-star packing. In addition to their original support and degree conclusions, their actual remaining active sample satisfies the exact outside balance equations required by marked active cover-down. This holds for every invariant embedding onto the specified active complements.
\end{lemma}
\begin{proof}
\uses{bp:marked-cleared-active-balance}
\leanok
Apply the already constructed correction and deep-cover theorem. The reserved arcs and spokes lie wholly in the inner pair, so removing them preserves all outside identities. Remove the actual prescribed packing, then apply the protected-star coverage and active restriction result to the actual deep packing. All returned constants and their quantifier order are unchanged. This verifies the exact balance input; the remaining approximate balance, profile and numerical inputs to active cover-down are separate obligations.
\end{proof}

\begin{lemma}[Exact balance after the protected deletions]
\label{bp:marked-finite-active-balance-values}
\lean{Oberwolfach.markedFiniteBalance_sdiff,Oberwolfach.markedFiniteBalance_eq_zero_iff,Oberwolfach.markedFiniteBalance_subtract_divisible,Oberwolfach.markedFiniteDomainBalance_domain_sdiff,Oberwolfach.markedFiniteDomainBalance_sdiff,Oberwolfach.markedFiniteBalance_bound,Oberwolfach.markedFiniteDomainBalance_bound,Oberwolfach.markedFiniteDomainBalance_zero_of_avoids,Oberwolfach.markedFiniteBalance_star_outside,Oberwolfach.markedFiniteBalance_star_inside,Oberwolfach.markedFiniteDomainBalance_star_active_zero,Oberwolfach.marked_finite_active_remainder_balances,Oberwolfach.marked_finite_active_inner_error}
\leanok
Write the three row balance differences and the hub difference as real-valued functions of the finite typed arc and spoke sets. Subtracting a contained divisible packing leaves the full balance unchanged. At an active root, the reserved star's full balance equals the corrected inner graph's opposite-endpoint balance on the deeper sets. Its balance restricted to the active inner pair is zero. A packing avoiding that pair's entire star contributes zero to its restricted balance at every root. Consequently the restricted balance after both packings is the original inner balance minus the deeper-set balance minus the prescribed packing's restricted balance. Actual incidence bounds $\delta$ bound the last contribution by $(m+2)\delta$.
\end{lemma}
\begin{proof}
\uses{bp:marked-constructed-deep-active-balance}
\leanok
Use exact finite-set subtraction with containment, then subtract nested opposite-endpoint domains. Evaluate the reserved-star filters separately at roots in and outside the deeper sets. Preserve the hub coefficient $m$ and bound the realised deleted incidences directly, without a sampling factor. The caller must still supply its original restricted-balance estimates and numerical substitutions.
\end{proof}

\begin{lemma}[Balance values in the actual representations]
\label{bp:marked-finite-balance-representations}
\lean{Oberwolfach.markedFiniteBalance_full_row,Oberwolfach.markedFiniteBalance_full_hub,Oberwolfach.markedBalanceTermSet_domain_card,Oberwolfach.markedFiniteBalance_subtype_restriction,Oberwolfach.markedFiniteDomainBalance_row,Oberwolfach.markedFiniteDomainBalance_hub,Oberwolfach.markedFiniteBalance_embedding,Oberwolfach.markedFiniteDomainBalance_embedding,Oberwolfach.marked_restricted_sample_balance_values,Oberwolfach.marked_deleted_sample_balance_values}
\leanok
The finite balance values equal the six-term balance expressions used by approximate wheel packing. For the restricted expressions, only the opposite endpoint is restricted; all ambient roots remain. Injective maps of typed occurrences, rows and hubs preserve these values at mapped roots, when the spoke types agree on the supplied spoke set and the inner domains are exact preimages. These identities are instantiated for the actual physically deleted sample and its cleared active-domain restriction.
\end{lemma}
\begin{proof}
\uses{bp:marked-finite-active-balance-values}
\leanok
Count the two arc types and the disjoint spoke types separately. Compare each actual subtype's finite incidence sets through its injective occurrence map. The same root-specific restriction works for all six degree terms. No event is enlarged to outcome-dependent domains.
\end{proof}

\begin{lemma}[Exact active resources and inner degrees]
\label{bp:marked-active-resource-degree-transfer}
\lean{Oberwolfach.markedFiniteResources_embedding,Oberwolfach.marked_restricted_sample_resources,Oberwolfach.marked_deleted_sample_resources,Oberwolfach.markedPackingResources_embedding_graphs,Oberwolfach.markedPackingResources_embedding_rim_degree,Oberwolfach.spanningRestriction_map_preimage,Oberwolfach.markedPackingResources_embedding_inner_degrees}
\leanok
Deletion of the actual finite rim and spoke sets gives their exact physical resource complement. Restriction to active domains represents all surviving resources when their endpoints lie in those domains. Embedding a packing back into the original vertices preserves the mapped rim graph, spoke relation, inner rim degrees and inner spoke row and column counts exactly.
\end{lemma}
\begin{proof}
\uses{bp:marked-active-sample-packing-lift}
\leanok
Identify the exact surviving occurrence ranges and the permutation-compatible physical rim map. Physical resource equality determines rim graphs and spoke relations. Map their neighbour sets and restricted spoke sets injectively; the resulting cardinalities are unchanged.
\end{proof}

\begin{lemma}[Quantitative balance in the constructed active sample]
\label{bp:marked-protected-active-estimates}
\lean{Oberwolfach.marked_finite_physical_incidence_bound,Oberwolfach.markedFiniteDomainBalance_induced_inside,Oberwolfach.marked_finite_corrected_inner_domain_bound,Oberwolfach.marked_sample_packing_domain_balance_bound,Oberwolfach.markedFiniteBalance_zero_of_support,Oberwolfach.markedFiniteDomainBalance_zero_of_support,Oberwolfach.marked_protected_finite_balance_bounds,Oberwolfach.marked_packing_finite_star_avoidance,Oberwolfach.marked_protected_sample_balance_bounds,Oberwolfach.marked_canonical_deleted_packing,Oberwolfach.marked_canonical_deleted_balance_values,Oberwolfach.marked_canonical_active_balance_bounds,Oberwolfach.marked_two_stage_active_balance_bounds,Oberwolfach.marked_constructed_deep_active_estimates}
\leanok
The actual inner correction and the two constructed deep-stage packings supply both full and inner-restricted balance bounds at every active root. The full error is at most $\varepsilon_T+(m+2)\delta_0$, where $\delta_0$ bounds the inner correction's incidences. The inner-restricted error is at most $\varepsilon_I+\varepsilon_T+(m+2)(4\delta_R+\delta_S)$, where $\delta_R,\delta_S$ bound the prescribed packing's physical rim and spoke degrees. The deep packing contributes zero to the inner-restricted balance by its proved avoidance of the entire active inner star.
\end{lemma}
\begin{proof}
\uses{bp:marked-finite-active-balance-values,bp:marked-finite-balance-representations,bp:marked-constructed-deep-active-balance}
\leanok
Compare the actual original, deleted and active occurrence sets through their explicit injective maps. A physical rim degree bounds the relevant auxiliary incidences by four times that degree, including the shifted root for incoming marked arcs. Apply the finite subtraction identities, use induced-domain equality at inner roots and zero support at outside roots, and transfer through the exact active embeddings. The canonical lift of the actual second packing supplies the required containment and resource equalities. Original opposite-endpoint domains and ambient roots are retained.
\end{proof}

\begin{lemma}[Control of the deep-stage error parameter]
\label{bp:marked-deep-error-controlled}
\lean{Oberwolfach.marked_corrected_sample_deep_cover_controlled,Oberwolfach.marked_constructed_inner_deep_cover_controlled,Oberwolfach.marked_constructed_deep_active_balance_controlled,Oberwolfach.marked_constructed_deep_active_estimates_controlled}
\leanok
The correction and deep-cover constructors admit every positive error parameter below a fixed positive upper bound. Their parent tolerances, packing constants and inner-size allowance precede this choice; the subsequent correction tolerance and order thresholds follow it. The previous public statements are consequences with the error parameter equal to that upper bound and have unchanged hypotheses and conclusions.
\end{lemma}
\begin{proof}
\uses{bp:marked-protected-active-estimates}
\leanok
In the original construction choose an upper bound satisfying both required deep-cover inequalities, then run the same proof for an arbitrary smaller positive error. All subsequent thresholds are chosen in the same order. This exposes the available quantifier needed to make the loss small relative to an inner-domain ratio fixed later; it does not infer that a single previously fixed error is arbitrarily small. The original three exported statement texts were checked unchanged before integration.
\end{proof}

\begin{lemma}[Physical deletions on the active domains]
\label{bp:marked-protected-active-deletion}
\lean{Oberwolfach.rawMarkedDeleteRims_union,Oberwolfach.rawMarkedDeleteRims_sup,Oberwolfach.markedDeletePhysicalSpokes_sup,Oberwolfach.marked_protected_active_deletion}
\leanok
Successive physical deletions agree with deletion of their union, retaining all directions and types. On the active domains the reserved inner star is absent. The actual rim, spoke-row and spoke-column degrees deleted there are bounded by the sum of the preceding packing's bounds and the deep packing's free-domain bounds; no global degree bound on the deep packing is imposed.
\end{lemma}
\begin{proof}
\uses{bp:marked-active-resource-degree-transfer}
\leanok
Check each surviving physical coordinate and the exact active-domain preimages. The deep packing avoids every edge incident with the protected inner pair, so its contribution is zero at those roots. At every other active root apply its proved free-domain bounds and the union degree inequalities. These bounds are inputs for profile and fixed-weight transfer; the complete protected cover-down caller remains to be assembled.
\end{proof}

\begin{lemma}[Typed potential sums and actual sampled weights]
\label{bp:marked-typed-fixed-seed-loads}
\lean{Oberwolfach.markedCycleIncidence_canonical_indicator,Oberwolfach.markedWheel_spoke_typed_indicator,Oberwolfach.marked_wheel_sum_over_survivors,Oberwolfach.typedMarkedPotentialLoad_surviving_sum,Oberwolfach.typedMarkedPotentialLoad_restricted_rim,Oberwolfach.typedMarkedPotentialLoad_restricted_spoke,Oberwolfach.marked_sample_seed_rim_load,Oberwolfach.marked_sample_seed_spoke_load,Oberwolfach.marked_sample_seed_witness,Oberwolfach.marked_sample_seed_present_root_bounds}
\leanok
The original root-omitted typed potential sums equal the actual sampled wheel loads at present roots. Original successor coefficients yield a positive approximate weighting of the actual sampled wheels, with the original normalization and two-sided individual bounds. The complete pair universe need not have an injective physical-edge map or disjoint spoke types.
\end{lemma}
\begin{proof}
\leanok
Identify directed incidence with canonical typed membership, sum over surviving cycles and hubs, and restore the omitted root only when that root is present. At absent spokes the actual residual is zero. The individual lower bound applies to each actual sampled cycle through its parent successor validity.
\end{proof}

\begin{lemma}[Fixed-weight losses under physical and vertex deletions]
\label{bp:marked-fixed-weight-deletion-losses}
\lean{Oberwolfach.marked_packing_potential_after_physical_deletion,Oberwolfach.marked_deleted_sample_seed_witness,Oberwolfach.canonical_marked_rim_extra_row_weight,Oberwolfach.canonical_marked_rim_fixed_hub_weight,Oberwolfach.typedMarked_spoke_extra_row_weight,Oberwolfach.typedMarked_spoke_other_hub_weight_zero,Oberwolfach.finite_weighted_exception_union,Oberwolfach.resourceRootOmittedLoad_exceptional_candidates,Oberwolfach.typedMarked_forbidden_set_weight,Oberwolfach.typedMarkedPotentialLoad_vertex_deletion,Oberwolfach.marked_induced_sample_seed_witness}
\leanok
Physical deletions with rim degree at most $\delta_R$ and spoke row and column degree at most $\delta_S$ change each original potential load by at most $11(m+1)^4B_*(2\delta_R+\delta_S)/n$, where $B_*=2B/p_*^{2m+1}$. Removing an invariant row set $T$ and a hub set $Z$ changes loads at retained roots by at most $(|T|+|Z|)(m+1)^3B_*/n$. These are losses of the fixed weights, without a sampling-density multiplier. Their sum with the original load error bounds the actual sampled weights on the induced domains; surviving individual weights retain their original bounds.
\end{lemma}
\begin{proof}
\uses{bp:marked-typed-fixed-seed-loads}
\leanok
For edge deletions count changed auxiliary incidences, retaining the shifted incoming marked root. For vertex deletions bound the weight of candidates using each forbidden row or hub; at a retained root an extra forbidden row is distinct from the prescribed root rows. Apply the finite weighted union bound to candidates whose survival indicators can change. This does not impose bounded deletion degree at the removed vertices.
\end{proof}

\begin{lemma}[Exact transfer to the active vertex domains]
\label{bp:marked-active-fixed-weight-transfer}
\lean{Oberwolfach.markedWheelVertexEmbedding_surjective,Oberwolfach.markedCycleIncidence_vertex_embedding,Oberwolfach.markedWheelVertexEmbedding_loads,Oberwolfach.marked_restricted_sample_seed_loads,Oberwolfach.markedPhysicalRim_spanningRestriction_iff,Oberwolfach.rawMarkedPairRims_spanningRestriction,Oberwolfach.markedSampleArc_spanningRestriction_support,Oberwolfach.marked_restricted_sample_seed_bounds}
\leanok
When all surviving rim endpoints and spoke endpoints belong to the active domains, restriction induces a bijection on actual wheels. Composing the supplied coefficients with this embedding preserves every rim and spoke load and both individual coefficient bounds. Inducing the physical graph on a permutation-invariant set restricts both auxiliary endpoints, including the shifted physical endpoint of a marked arc.
\end{lemma}
\begin{proof}
\uses{bp:marked-typed-fixed-seed-loads,bp:marked-fixed-weight-deletion-losses}
\leanok
Construct the inverse tuple and hub through the injective row and hub embeddings and prove their arc and spoke predicates. Transport incidence and sums through the resulting bijection. No coefficient is recomputed and no bound normalized by the original order is silently replaced by a bound normalized by the active order. The full protected cover-down caller remains a separate obligation.
\end{proof}

\begin{lemma}[Actual active weights after protected deletions]
\label{bp:marked-active-weight-caller}
\lean{Oberwolfach.marked_complement_seed_transfer,Oberwolfach.marked_fixed_weight_potential_physical_deletion,Oberwolfach.marked_sample_fixed_weight_deletion,Oberwolfach.marked_active_sample_seed_witness,Oberwolfach.marked_protected_active_seed_witness,Oberwolfach.fixed_weight_order_bounds}
\leanok
The supplied parent successor coefficients give an actual approximate wheel weighting on the complementary row and hub domains, after any subsequent physical deletion whose degrees on those active domains have the displayed bounds. If $n_A$ is the active row order, the load error is at most
\[
\eta+\frac{(|T|+|Z|)(m+1)^3B_*}{n}
 +\frac{11(m+1)^4B_*(2\delta_R+\delta_S)}{n_A},
\qquad B_*=\frac{2B}{p_*^{2m+1}}.
\]
Every surviving coefficient has its original two-sided bounds. When $n_A\le n\le2n_A$, their conversion to the active order loses a factor at most $2^m$ in the lower coefficient and no factor in the upper coefficient.
\end{lemma}
\begin{proof}
\uses{bp:marked-active-fixed-weight-transfer,bp:marked-fixed-weight-deletion-losses,bp:marked-protected-active-deletion}
\leanok
First remove the forbidden vertices, then use the exact bijection onto the complementary sample. Apply the deterministic physical-deletion bound there and commute restriction with deletion of the original physical coordinates. The upper coefficient bound on the full potential family follows from the original successor coefficients; no successor-weight condition is assumed for the smaller graph. The remaining full cover-down caller must supply its already constructed deletion graphs and degree bounds.
\end{proof}

\begin{lemma}[Neighbourhood estimates on the actual active sample]
\label{bp:marked-active-neighbourhood-caller}
\lean{Oberwolfach.markedRowRim_embedding,Oberwolfach.markedRow_parent_embedding,Oberwolfach.markedRowEligible_embedding,Oberwolfach.markedRowPhysical_embedding,Oberwolfach.markedRow_physical_injective_embedding,Oberwolfach.markedRowDomain_image,Oberwolfach.markedRowCountWithSpokes_embedding,Oberwolfach.markedIndexCountWithSpokes_embedding,Oberwolfach.markedRowsWithSpokes_embedding,Oberwolfach.markedIndicesWithSpokes_embedding,Oberwolfach.finite_bounded_sum_restriction,Oberwolfach.finite_centered_count_restriction,Oberwolfach.markedRowsWithSpokes_domain_restriction,Oberwolfach.markedIndicesWithSpokes_domain_restriction,Oberwolfach.marked_protected_active_neighbourhoods}
\leanok
Exact row and hub embeddings preserve bounded request sets, permutation separation, physical-edge distinctness, eligible-domain sizes and density products. Restricting an outer summation domain adds at most its number of removed vertices to the absolute count error. Subsequent active physical deletions add at most $h(\delta_R+\delta_U)$ to row errors and $h\delta_V$ to index errors. These conclusions hold for all active roots using the original all-root estimates and retain the original density profiles.
\end{lemma}
\begin{proof}
\leanok
Map each typed directed request explicitly and use injectivity of the induced map on unordered physical pairs. Reindex the exact finite sums. For a removed coordinate the centered summand, an indicator minus a number in $[0,1]$, has absolute value at most one. Apply the checked physical neighbourhood-deletion bounds after transferring to the actual active types. There is no division by the eligible-domain size and no inference about external roots from an induced graph alone.
\end{proof}

\begin{lemma}[Active spoke intersections]
\label{bp:active-spoke-data}
\lean{Oberwolfach.markedRows_spoke_intersections,Oberwolfach.markedIndices_spoke_codegree,Oberwolfach.marked_neighbourhoods_sparse_intersection_data,Oberwolfach.marked_protected_active_spoke_data}
\leanok
The actual active spoke relations satisfy the required intersection estimates, with constant density profiles inherited from the original neighbourhood event. The exceptional set for a root is exactly its permutation-distance-$3d$ neighbourhood, of size at most $6d+1$. No codegree conclusion is imposed inside that set.
\end{lemma}
\begin{proof}
\leanok
Use the exact representation identities and deterministic finite-count inequalities in the named Lean declarations. These conclusions retain the original physical sampling law and permit every later deletion satisfying the displayed bounds on the same event.
\end{proof}

\begin{lemma}[Actual active cuts and degrees]
\label{bp:active-cuts-degrees}
\lean{Oberwolfach.bipartiteLargeCut_sum_restrict,Oberwolfach.bipartiteMatrixCut_embedding,Oberwolfach.bipartiteLargeCut_embedding,Oberwolfach.markedSampleArc_indicator_iff,Oberwolfach.markedPhysicalRim_indicator,Oberwolfach.markedSampleArc_multiplicity,Oberwolfach.markedSampleArc_degrees,Oberwolfach.markedPhysicalRim_restrictPair_eq,Oberwolfach.marked_rim_largeCut_after_physical_deletion,Oberwolfach.marked_protected_active_cuts,Oberwolfach.bounded_matrix_vertex_restriction,Oberwolfach.finset_indicator_sum_change,Oberwolfach.markedPhysicalRim_real_le_one,Oberwolfach.marked_rim_degree_physical_deletion,Oberwolfach.marked_protected_active_degrees}
\leanok
The sampled typed arc identifiers have exactly the zero-one physical adjacency matrices. Restriction to fixed row and hub embeddings and later physical deletion give all four cut bounds. The cutoff change is explicit: $1-2\delta\le\kappa(1-2\delta_A)$, where both retained part fractions are at least $\kappa$. Realized deletion degrees cost $\zeta$ in the cut coefficient. A rim degree loses at most the number of removed vertices and its relevant physical deletion degree. The incoming marked degree is charged at the shifted physical root. Any change in degree target is charged separately by its absolute difference.
\end{lemma}
\begin{proof}
\leanok
Use the exact representation identities and deterministic finite-count inequalities in the named Lean declarations. These conclusions retain the original physical sampling law and permit every later deletion satisfying the displayed bounds on the same event.
\end{proof}

\begin{lemma}[Two-step counts at all distinct active roots]
\label{bp:active-two-step}
\lean{Oberwolfach.relation_two_step_count_eq,Oberwolfach.relation_two_step_vertex_restriction,Oberwolfach.relation_two_step_changes,Oberwolfach.marked_rim_two_step_physical_deletion,Oberwolfach.marked_protected_active_two_step}
\leanok
The actual two-step count differs from its original count by at most the number of removed intermediate vertices plus twice the physical degree deletion bound. This holds for all distinct endpoints, with no separation hypothesis.
\end{lemma}
\begin{proof}
\leanok
Use the exact representation identities and deterministic finite-count inequalities in the named Lean declarations. These conclusions retain the original physical sampling law and permit every later deletion satisfying the displayed bounds on the same event.
\end{proof}

\begin{lemma}[Actual active common-cycle and path counts]
\label{bp:active-common-configurations}
\lean{Oberwolfach.embeddedTuplePredicateEquiv,Oberwolfach.embedded_tuple_predicate_card,Oberwolfach.markedCycle_endpoint_transport_iff,Oberwolfach.tuple_embedding_insertNth,Oberwolfach.commonMarkedCycle_vertex_card,Oberwolfach.canonicalCommonMarkedPath_vertex_card,Oberwolfach.canonicalCommonMarkedPath_spoke_changes,Oberwolfach.marked_rim_relation_deletion_bounds,Oberwolfach.marked_spoke_relation_deletion_bound,Oberwolfach.marked_sample_common_cycle_physical_deletion,Oberwolfach.marked_sample_common_path_physical_deletion,Oberwolfach.potentialMarkedHubs_embedding_loss,Oberwolfach.marked_protected_active_compatible_hubs,Oberwolfach.markedSampleArc_restrictPair_exists_iff,Oberwolfach.spoke_domain_change_bound,Oberwolfach.marked_sample_common_cycle_vertex_loss,Oberwolfach.marked_sample_common_path_vertex_loss,Oberwolfach.marked_protected_active_common_configurations,Oberwolfach.marked_active_cycle_margin,Oberwolfach.marked_active_path_margin,Oberwolfach.marked_active_cut_parameters}
\leanok
Fixed vertex embeddings give exact common-cycle and common-path transport once their free coordinates are restricted to the image. If $t$ vertices are removed, the cycle loss is at most $2(m+1)t n^m$ and the canonical path loss at most $2mt n^{m-1}$. A further physical deletion of active degree at most $\Delta$ costs at most $3(m+1)\Delta n_A^m$ and $4(m+1)\Delta n_A^{m-1}$, respectively. The two root spokes are never imposed on a common path. Compatible-hub loss is at most the number of removed hubs plus $(m+1)$ times the active spoke-row deletion bound. The original and active orders remain distinct in these formulas. Explicit scalar margins turn these comparisons into lower bounds when their displayed numerical budgets hold; the complete cover-down caller must still verify those budgets.
\end{lemma}
\begin{proof}
\leanok
Use the exact representation identities and deterministic finite-count inequalities in the named Lean declarations. These conclusions retain the original physical sampling law and permit every later deletion satisfying the displayed bounds on the same event.
\end{proof}

\begin{lemma}[Reference lower bounds on the active graph]
\label{bp:active-reference-lower-bounds}
\lean{Oberwolfach.marked_protected_active_configuration_lower_bounds,Oberwolfach.markedPhysicalRim_pos_iff_one,Oberwolfach.marked_protected_active_reference_statistics,Oberwolfach.marked_degree_target_removed_hubs,Oberwolfach.square_bipartiteLargeCut_iff,Oberwolfach.marked_protected_active_ordinary_cut,Oberwolfach.marked_reference_four_bipartite_cuts,Oberwolfach.marked_protected_active_four_cuts}
\leanok
The parent reference degree, two-step, cut and common-configuration inequalities imply the corresponding active cover-down inputs. Every inequality retains its stated numerical budget. The path correspondence preserves the omitted root spokes. The change in degree target is charged to the removed hubs.
\end{lemma}
\begin{proof}
\leanok
The named Lean declarations derive these conclusions from the original sampled profiles and the explicit realized deletion bounds. The complete protected cover-down and preparation callers remain open.
\end{proof}

\begin{lemma}[Inner profiles after the prescribed and deep packings]
\label{bp:inner-after-both-packings}
\lean{Oberwolfach.marked_neighbourhoods_star_avoiding_deletion,Oberwolfach.marked_inner_neighbourhoods_after_both_packings,Oberwolfach.marked_inner_profiles_after_both_packings}
\leanok
Deep wheels avoiding the entire inner star change neither the inner neighbourhood tests nor the fixed-domain configurations or weighted spoke loads. The prescribed-edge deletion contributes the explicit error. The weight function and its original normalization are unchanged, and external roots are included.
\end{lemma}
\begin{proof}
\leanok
The named Lean declarations derive these conclusions from the original sampled profiles and the explicit realized deletion bounds. The complete protected cover-down and preparation callers remain open.
\end{proof}

\begin{lemma}[Actual sampled extension and endpoint inputs]
\label{bp:actual-cover-down-extensions}
\lean{Oberwolfach.marked_sample_arc_present_pair,Oberwolfach.marked_sample_arc_separated,Oberwolfach.marked_common_two_vertex_hubs,Oberwolfach.marked_sample_extension_counts,Oberwolfach.marked_cover_down_extension_inputs,Oberwolfach.marked_inner_endpoint_sizes_after_packings}
\leanok
Actual sampled arcs satisfy the separation used for common-hub counts. The fixed-domain row and index estimates imply the exact step, join and hub-extension hypotheses of active cover-down, with constants chosen before the separation radius. Endpoint class sizes and the marked-outgoing pairing gap are derived in the host after both earlier packings; the deep packing has zero cost on this domain.
\end{lemma}
\begin{proof}
\leanok
The named Lean declarations derive these conclusions from the original sampled profiles and the explicit realized deletion bounds. The complete protected cover-down and preparation callers remain open.
\end{proof}


\begin{lemma}[Actual active packing lift and coverage]
\label{bp:active-packing-lift-coverage}
\lean{Oberwolfach.marked_active_deleted_rim_preimage,Oberwolfach.marked_active_deleted_spoke_preimage,Oberwolfach.marked_packing_covers_outside_inner,Oberwolfach.marked_active_cover_packing_lift,Oberwolfach.marked_packing_disjoint_from_deletion,Oberwolfach.markedPackingResources_embedding_inner_bounds,Oberwolfach.marked_deleted_star_cover_transport}
\leanok
The active packing lifts to the original typed host with its exact physical resource image. Its surviving outside resources are covered, physical deletions are avoided, and all ambient inner-degree bounds are retained.
\end{lemma}
\begin{proof}
\leanok
The named declarations prove these conclusions from their explicit inputs. The complete protected cover-down, preparation and iteration callers remain open.
\end{proof}

\begin{lemma}[Three-stage protected completion]
\label{bp:three-stage-protected-completion}
\lean{Oberwolfach.marked_protected_three_stage_coverage,Oberwolfach.marked_protected_three_stage_completion}
\leanok
The prescribed, deep and active packings form one physical packing. Coverage and cross-disjointness follow from the actual deleted hosts. Subtraction gives the exact divisible inner remainder and the protected-star identities; temporary reservations are not counted as removed resources.
\end{lemma}
\begin{proof}
\leanok
The named declarations prove these conclusions from their explicit inputs. The complete protected cover-down, preparation and iteration callers remain open.
\end{proof}

\begin{lemma}[Uniform active reference substitutions]
\label{bp:uniform-active-reference-substitutions}
\lean{Oberwolfach.marked_active_domain_size_bounds,Oberwolfach.marked_active_reference_budget_from_scale,Oberwolfach.marked_active_reference_scale_exists,Oberwolfach.marked_active_reference_geometry_from_scale,Oberwolfach.marked_active_exception_order,Oberwolfach.marked_active_seed_from_scale,Oberwolfach.marked_two_stage_deletion_scale,Oberwolfach.marked_active_reference_degrees_from_scale,Oberwolfach.marked_active_reference_cuts_from_scale,Oberwolfach.marked_active_configurations_from_scale,Oberwolfach.marked_active_spoke_data_from_scale}
\leanok
One positive scale simultaneously supplies the degree, cut, common-cycle, root-omitted path, spoke-intersection and fixed-weight error budgets on the active domain. It is chosen before the graph order and separation radius. The original weight function and its normalization remain fixed.
\end{lemma}
\begin{proof}
\leanok
The named declarations prove these conclusions from their explicit inputs. The complete protected cover-down, preparation and iteration callers remain open.
\end{proof}

\begin{lemma}[Correction parameters after the deep-cover constants]
\label{bp:deep-constants-before-correction}
\lean{Oberwolfach.marked_constructed_inner_deep_cover_ordered,Oberwolfach.marked_constructed_deep_active_balance_ordered,Oberwolfach.marked_constructed_deep_active_estimates_ordered}
\leanok
The inner-correction parameters, including its deletion tolerance, are universally quantified after the deep-cover constants. The construction proves this stronger order of choices. The earlier public statements follow by specialization.
\end{lemma}
\begin{proof}
\leanok
The named declarations prove these conclusions from their explicit inputs. The complete protected cover-down, preparation and iteration callers remain open.
\end{proof}


\begin{lemma}[Combined global reference inputs]
\label{bp:global-reference-active-inputs}
\lean{Oberwolfach.marked_active_global_reference_conditions,Oberwolfach.marked_active_cover_down_global_conditions,Oberwolfach.marked_active_cover_down_canonical_roots}
\leanok
The fixed parent reference estimates supply every global input of active cover-down, uniformly over the later permitted physical deletions. The cover constructor uses actual complement subtypes for its outside roots and hubs; their cardinality bounds follow from its size hypothesis.
\end{lemma}
\begin{proof}
\leanok
The named declarations prove these conclusions from their explicit hypotheses. The full protected cover-down caller and main theorem remain open.
\end{proof}

\begin{lemma}[Inner factors and neighbourhoods on active types]
\label{bp:inner-active-factor-neighbourhoods}
\lean{Oberwolfach.marked_inner_factor_deletion_tolerance,Oberwolfach.marked_inner_factor_inputs_after_both_packings,Oberwolfach.rawMarkedRestrictVertices_pair_comp,Oberwolfach.marked_inner_factor_active_restriction,Oberwolfach.marked_inner_active_neighbourhoods,Oberwolfach.marked_active_local_extension_inputs}
\leanok
The original inner coefficients and outcome coordinates give the factor and neighbourhood inputs after both earlier packings. Roots outside the inner domains remain included. Nested restrictions agree pointwise.
\end{lemma}
\begin{proof}
\leanok
The named declarations prove these conclusions from their explicit hypotheses. The full protected cover-down caller and main theorem remain open.
\end{proof}

\begin{lemma}[Actual balance and separate deletion allowances]
\label{bp:protected-active-balance-and-deletion}
\lean{Oberwolfach.marked_active_balance_parameter_allowances,Oberwolfach.marked_protected_active_balance_inputs,Oberwolfach.marked_active_total_deletion_bounds_separate,Oberwolfach.marked_protected_active_deletion_separate,Oberwolfach.marked_constructed_active_deletion_bounds}
\leanok
The constructed finite balance estimates imply the four active balance inputs. Distinct row and hub deletion allowances are retained, and the actual two-packing bounds meet the chosen active tolerance.
\end{lemma}
\begin{proof}
\leanok
The named declarations prove these conclusions from their explicit hypotheses. The full protected cover-down caller and main theorem remain open.
\end{proof}

\begin{lemma}[Active endpoint counts and pairing gap]
\label{bp:active-endpoint-counts}
\lean{Oberwolfach.markedEndpointLabels_embedding,Oberwolfach.markedEndpointLabels_embedding_card,Oberwolfach.markedSpokeNeighbours_embedding_card,Oberwolfach.marked_active_endpoint_sizes_after_packings}
\leanok
Exact embedding identities preserve endpoint directions and marked types, spoke-neighbour counts and the marked-outgoing pairing gap. The scale conversion uses the proved ambient cardinality bound.
\end{lemma}
\begin{proof}
\leanok
The named declarations prove these conclusions from their explicit hypotheses. The full protected cover-down caller and main theorem remain open.
\end{proof}

\begin{lemma}[Ordered active cover-down constructor]
\label{bp:ordered-active-cover-constructor}
\lean{Oberwolfach.marked_active_cover_down_reference,Oberwolfach.marked_active_cover_down_constructed_balance,Oberwolfach.marked_active_cover_down_uniform_densities,Oberwolfach.marked_active_cover_down_local_extensions,Oberwolfach.marked_active_cover_down_endpoint_inputs}
\leanok
The active packing is constructed from original reference estimates, finite balance estimates, fixed density profiles and actual inner inputs. Extension counts and endpoint class estimates are derived. The cycle-length threshold depends on an ordinary-rim density bound fixed before the length; the lower bound for all densities is chosen afterwards.
\end{lemma}
\begin{proof}
\leanok
The named declarations prove the stated conclusions from their explicit hypotheses. The full protected constructor and main theorem remain open.
\end{proof}

\begin{lemma}[Original inner inputs after constructed deletions]
\label{bp:original-inner-input-deletions}
\lean{Oberwolfach.marked_active_inner_inputs_after_packings,Oberwolfach.marked_active_inner_parent_degrees,Oberwolfach.graph_neighbourhood_embedding_card,Oberwolfach.marked_prescribed_inner_deletion}
\leanok
The actual earlier packing degrees imply inner deletion allowances. The original fixed coefficients supply factor rows, configurations and loads after both packings; inner neighbourhood estimates include every ambient root, with the lost hub target charged explicitly.
\end{lemma}
\begin{proof}
\leanok
The named declarations prove the stated conclusions from their explicit hypotheses. The full protected constructor and main theorem remain open.
\end{proof}

\begin{lemma}[Local parameter scale]
\label{bp:protected-local-parameter-scale}
\lean{Oberwolfach.marked_protected_local_scale_exists}
\leanok
One positive scale chosen after the active tolerances simultaneously satisfies the listed local deletion, extension, endpoint and balance inequalities.
\end{lemma}
\begin{proof}
\leanok
The named declarations prove the stated conclusions from their explicit hypotheses. The full protected constructor and main theorem remain open.
\end{proof}


\begin{lemma}[Marked Protected Active Geometry]
\label{bp:markedprotectedactivegeometry}
\lean{Oberwolfach.marked_protected_active_geometry}
\leanok
The actual complement embeddings give lower bounds for both remaining inner parts and an upper bound for their combined size.
\end{lemma}
\begin{proof}
\leanok
The named declarations prove these conclusions under their explicit hypotheses. This does not complete the full protected construction or Theorem 1.1.
\end{proof}

\begin{lemma}[Marked Protected Inner Budgets]
\label{bp:markedprotectedinnerbudgets}
\lean{Oberwolfach.marked_protected_inner_budget_scale,Oberwolfach.marked_protected_inner_error_budget,Oberwolfach.marked_protected_inner_exception_order}
\leanok
One positive error scale, chosen before the separation and order parameters, gives the required global deletion, inner factor, profile, endpoint and finite balance margins. The changed hub target contributes twice the deleted hub count divided by the original density.
\end{lemma}
\begin{proof}
\leanok
The named declarations prove these conclusions under their explicit hypotheses. This does not complete the full protected construction or Theorem 1.1.
\end{proof}

\begin{lemma}[Marked Protected Completion From Active Packing]
\label{bp:markedprotectedcompletionfromactivepacking}
\lean{Oberwolfach.marked_protected_completion_from_active_packing}
\leanok
The actual active packing lifts to the original sample and combines with the two earlier packings to leave a divisible inner remainder. The protected star contains exactly the prescribed correction, and the deepest induced subgraph is unchanged.
\end{lemma}
\begin{proof}
\leanok
The named declarations prove these conclusions under their explicit hypotheses. This does not complete the full protected construction or Theorem 1.1.
\end{proof}

\begin{lemma}[Marked Protected Canonical Domains]
\label{bp:markedprotectedcanonicaldomains}
\lean{Oberwolfach.marked_protected_canonical_domains}
\leanok
Complement and inner domains, their embeddings and restricted permutations are constructed as actual subtypes, with exact images and cardinalities.
\end{lemma}
\begin{proof}
\leanok
The named declarations prove these conclusions under their explicit hypotheses. This does not complete the full protected construction or Theorem 1.1.
\end{proof}

\begin{lemma}[Marked Active Cover Down After Packings]
\label{bp:markedactivecoverdownafterpackings}
\lean{Oberwolfach.marked_active_cover_down_after_packings}
\leanok
The active packing is constructed after the two specified earlier packings from original fixed-domain estimates. Their actual physical degrees give the deletion costs. The deep packing adds no inner-star cost.
\end{lemma}
\begin{proof}
\leanok
The named declarations prove these conclusions under their explicit hypotheses. This does not complete the full protected construction or Theorem 1.1.
\end{proof}

\begin{lemma}[Marked Physical Auxiliary Incidence]
\label{bp:markedphysicalauxiliaryincidence}
\lean{Oberwolfach.marked_auxiliary_incident_physical_union,Oberwolfach.marked_incidence_bound_from_physical_resources}
\leanok
An auxiliary incidence at a row vertex is a physical incidence at that vertex or its image under the permutation. Uniform physical incidence bounds therefore give twice that bound for either representation.
\end{lemma}
\begin{proof}
\leanok
The named declarations prove these conclusions under their explicit hypotheses. This does not complete the full protected construction or Theorem 1.1.
\end{proof}

\begin{lemma}[Marked Resource Degree Upper]
\label{bp:markedresourcedegreeupper}
\lean{Oberwolfach.marked_incident_card_le_graph_degree}
\leanok
Physical resources with two-element rim edges are counted by their actual neighbours in a containing simple graph.
\end{lemma}
\begin{proof}
\leanok
The named declarations prove these conclusions under their explicit hypotheses. This does not complete the full protected construction or Theorem 1.1.
\end{proof}

\begin{lemma}[Marked Packing Inner Incidence]
\label{bp:markedpackinginnerincidence}
\lean{Oberwolfach.marked_packing_inner_physical_incidence_bound,Oberwolfach.marked_packing_inner_both_incidence_bound}
\leanok
The three actual inner rim, spoke-row and spoke-column degree bounds imply bounds for all physical and auxiliary incidences at every ambient root.
\end{lemma}
\begin{proof}
\leanok
The named declarations prove these conclusions under their explicit hypotheses. This does not complete the full protected construction or Theorem 1.1.
\end{proof}

\begin{lemma}[Marked Protected Completion Incidences]
\label{bp:markedprotectedcompletionincidences}
\lean{Oberwolfach.marked_protected_completion_with_incidence_bounds}
\leanok
The three-stage restoration preserves the stated incidence budgets in both representations, in addition to divisibility, exact inner remainder and protected-star identities.
\end{lemma}
\begin{proof}
\leanok
The named declarations prove these conclusions under their explicit hypotheses. This does not complete the full protected construction or Theorem 1.1.
\end{proof}

\begin{lemma}[Marked Active Cover Down From Deep Bounds]
\label{bp:markedactivecoverdownfromdeepbounds}
\lean{Oberwolfach.marked_active_cover_down_from_deep_bounds}
\leanok
The actual free-domain degrees of the deep packing and global degrees of the correction-edge packing give all three active deletion bounds. The remaining parent target includes both removed vertex and removed hub terms.
\end{lemma}
\begin{proof}
\leanok
The named declarations establish this result under their displayed hypotheses. Preparation, iteration and the main theorem are not certified by these declarations.
\end{proof}

\begin{lemma}[Marked Sample Balance Congruence]
\label{bp:markedsamplebalancecongruence}
\lean{Oberwolfach.marked_sample_finite_balance_congr,Oberwolfach.marked_sample_outside_balance_congr}
\leanok
Equality of the actual restricted outcomes and spoke predicates transfers full and restricted finite balances and exact outside balance without changing the probability law.
\end{lemma}
\begin{proof}
\leanok
The named declarations establish this result under their displayed hypotheses. Preparation, iteration and the main theorem are not certified by these declarations.
\end{proof}

\begin{lemma}[Marked Constructed Deep Canonical Estimates]
\label{bp:markedconstructeddeepcanonicalestimates}
\lean{Oberwolfach.marked_constructed_deep_canonical_estimates}
\leanok
The sparse correction and both preceding packings are constructed in original occurrence types. Their active balance estimates are transferred to deletion of precisely those packings.
\end{lemma}
\begin{proof}
\leanok
The named declarations establish this result under their displayed hypotheses. Preparation, iteration and the main theorem are not certified by these declarations.
\end{proof}

\begin{lemma}[Marked Finite Correction Incidence]
\label{bp:markedfinitecorrectionincidence}
\lean{Oberwolfach.marked_finite_correction_physical_incidence,Oberwolfach.marked_finite_correction_restricted_incidence}
\leanok
The correction has physical incidence at most three times its auxiliary arc and spoke bounds, at every ambient root. Marked heads are charged at the inverse permutation root.
\end{lemma}
\begin{proof}
\leanok
The named declarations establish this result under their displayed hypotheses. Preparation, iteration and the main theorem are not certified by these declarations.
\end{proof}

\begin{lemma}[Marked Protected Final Margins]
\label{bp:markedprotectedfinalmargins}
\lean{Oberwolfach.marked_protected_final_incidence_budget,Oberwolfach.marked_protected_final_positive_scales}
\leanok
The actual correction and active packing bounds imply the requested total inner and protected-star incidence bounds, with an explicit ratio between inner and parent orders.
\end{lemma}
\begin{proof}
\leanok
The named declarations establish this result under their displayed hypotheses. Preparation, iteration and the main theorem are not certified by these declarations.
\end{proof}

\begin{lemma}[Marked Finite Inner Graph Support]
\label{bp:markedfiniteinnergraphsupport}
\lean{Oberwolfach.marked_finite_inner_graph_support}
\leanok
An inner auxiliary arc set gives a physical graph supported on the same permutation-invariant inner set. The analogous spoke support is exact.
\end{lemma}
\begin{proof}
\leanok
The named declarations establish this result under their displayed hypotheses. Preparation, iteration and the main theorem are not certified by these declarations.
\end{proof}

\begin{lemma}[Marked Protected Separate Geometry Budgets]
\label{bp:markedprotectedseparategeometrybudgets}
\lean{Oberwolfach.marked_protected_inner_error_budget_separate_geometry}
\leanok
The deeper-set size allowance is kept separate from the smaller correction and weighted-profile error. Both the removed-row and removed-hub target changes are retained.
\end{lemma}
\begin{proof}
\leanok
The named declarations establish this result under their displayed hypotheses. Preparation, iteration and the main theorem are not certified by these declarations.
\end{proof}

\begin{lemma}[Marked Protected Vortex Geometry Scale]
\label{bp:markedprotectedvortexgeometryscale}
\lean{Oberwolfach.marked_protected_vortex_geometry_scale}
\leanok
A vortex ratio chosen before the correction scale satisfies the independent deep-star size budgets. The deeper fraction is the square of the ratio, and the active inner size lower bound is linear in it.
\end{lemma}
\begin{proof}
\leanok
The named declarations establish this result under their displayed hypotheses. Preparation, iteration and the main theorem are not certified by these declarations.
\end{proof}

\begin{lemma}[Marked Active Cover Down Separate Geometry]
\label{bp:markedactivecoverdownseparategeometry}
\lean{Oberwolfach.marked_active_cover_down_separate_geometry}
\leanok
The active packing is constructed with separate geometric and correction allowances, preserving all fixed weights, original domains and outside-balance conditions.
\end{lemma}
\begin{proof}
\leanok
The named declarations establish this result under their displayed hypotheses. Preparation, iteration and the main theorem are not certified by these declarations.
\end{proof}

\begin{lemma}[Marked Protected Cover Down Separate Geometry From Packings]
\label{bp:markedprotectedcoverdownseparategeometryfrompackings}
\lean{Oberwolfach.marked_protected_cover_down_separate_geometry_from_packings}
\leanok
The active packing and original-host restoration follow from the two earlier packings under the independent geometric allowances.
\end{lemma}
\begin{proof}
\leanok
The named declarations establish this result under their displayed hypotheses. Preparation, iteration and the main theorem are not certified by these declarations.
\end{proof}

\begin{lemma}[Marked Protected Cover Down Constructed Separate Geometry]
\label{bp:markedprotectedcoverdownconstructedseparategeometry}
\lean{Oberwolfach.marked_protected_cover_down_constructed_separate_geometry}
\leanok
The original host and its fixed profile inputs construct the sparse correction and all three packings. Their exact divisible remainder, deepest exclusion and physical and auxiliary incidence bounds are proved.
\end{lemma}
\begin{proof}
\leanok
The named declarations establish this result under their displayed hypotheses. Preparation, iteration and the main theorem are not certified by these declarations.
\end{proof}

\begin{lemma}[Marked Protected Cover Down]
\label{bp:markedprotectedcoverdown}
\lean{Oberwolfach.marked_protected_cover_down}
\leanok
The constructed packing satisfies the requested inner and protected-star degree bounds in both representations. Its inputs still have to be supplied by preparation and iteration; this is not the main theorem.
\end{lemma}
\begin{proof}
\leanok
The named declarations establish this result under their displayed hypotheses. Preparation, iteration and the main theorem are not certified by these declarations.
\end{proof}

\begin{lemma}[Vortex Protected Geometry]
\label{bp:vortexprotectedgeometry}
\lean{Oberwolfach.Vortex.extended_successor_size_bounds,Oberwolfach.Vortex.protected_size_bounds}
\leanok
The actual paired vortex satisfies every displayed inner, deeper and active-domain size bound, including the last empty deeper pair. Separate vertex and index floor recurrences are preserved.
\end{lemma}
\begin{proof}
\leanok
The named declarations prove the displayed conditional results. They do not establish the complete preparation, iteration caller or main theorem.
\end{proof}

\begin{lemma}[Marked Reference After Physical Deletion]
\label{bp:markedreferenceafterphysicaldeletion}
\lean{Oberwolfach.marked_reference_statistics_after_physical_deletion,Oberwolfach.marked_rim_cut_physical_deletion_bound,Oberwolfach.marked_reference_cuts_after_physical_deletion}
\leanok
Reference degrees, two-step intersections and directed cut bounds survive a bounded physical rim deletion with explicit deterministic losses.
\end{lemma}
\begin{proof}
\leanok
The named declarations prove the displayed conditional results. They do not establish the complete preparation, iteration caller or main theorem.
\end{proof}

\begin{lemma}[Marked Iteration Balance Preservation]
\label{bp:markediterationbalancepreservation}
\lean{Oberwolfach.markedFiniteDomainBalance_star,Oberwolfach.markedFiniteDomainBalance_zero_inside_of_empty_induced,Oberwolfach.marked_iteration_restricted_balance_preservation,Oberwolfach.marked_iteration_balance_normalized}
\leanok
Exact finite balance differences charge only deleted edges meeting the protected inner pair. If no inner edge is removed, balances at inner roots and the induced graph are unchanged.
\end{lemma}
\begin{proof}
\leanok
The named declarations prove the displayed conditional results. They do not establish the complete preparation, iteration caller or main theorem.
\end{proof}

\begin{lemma}[Marked Reference Adaptive Vortex Event]
\label{bp:markedreferenceadaptivevortexevent}
\lean{Oberwolfach.marked_reference_adaptive_vortex_event}
\leanok
One original retained experiment supplies the reference estimates at all fixed levels and for every later permitted physical rim deletion. No additional randomness is assigned to the deletion.
\end{lemma}
\begin{proof}
\leanok
The named declarations prove the displayed conditional results. They do not establish the complete preparation, iteration caller or main theorem.
\end{proof}

\begin{lemma}[Marked Finite Packing Iteration]
\label{bp:markedfinitepackingiteration}
\lean{Oberwolfach.markedFiniteResources_mono,Oberwolfach.marked_finite_packing_iteration,Oberwolfach.marked_finite_iteration_with_reserved_absorber}
\leanok
Given a proved packing step preserving a specified invariant, finite induction gives the exact removed-resource decomposition and terminal divisibility. This is conditional induction; the analytic step hypotheses remain to be supplied.
\end{lemma}
\begin{proof}
\leanok
The named declarations prove the displayed conditional results. They do not establish the complete preparation, iteration caller or main theorem.
\end{proof}

\begin{lemma}[Vortex Marked Parent Degrees]
\label{bp:vortexmarkedparentdegrees}
\lean{Oberwolfach.vortex_degree_margins,Oberwolfach.Vortex.marked_parent_degrees,Oberwolfach.Vortex.marked_difference_parent_degrees,Oberwolfach.Vortex.cover_down_parent_orders}
\leanok
The actual full and difference vortex domains satisfy the parent degree and target bounds, including the order thresholds at all cover-down steps.
\end{lemma}
\begin{proof}
\leanok
The named declarations prove the displayed conditional results. They do not establish the complete preparation, iteration caller or main theorem.
\end{proof}

\begin{lemma}[Marked Finite Counts From Incidence]
\label{bp:markedfinitecountsfromincidence}
\lean{Oberwolfach.marked_finite_counts_from_incidence,Oberwolfach.marked_finite_star_counts_from_incidence}
\leanok
The actual physical and auxiliary incidence conclusions imply finite arc, row and hub deletion bounds, both globally and on the protected star.
\end{lemma}
\begin{proof}
\leanok
The named declarations prove the displayed conditional results. They do not establish the complete preparation, iteration caller or main theorem.
\end{proof}

\begin{lemma}[Marked Iteration Balance From Resources]
\label{bp:markediterationbalancefromresources}
\lean{Oberwolfach.marked_iteration_balance_from_resource_bounds}
\leanok
The cover-down resource bounds imply the next restricted balance bounds, with exact inheritance at inner roots and no change to the inner graph.
\end{lemma}
\begin{proof}
\leanok
The named declarations prove the displayed conditional results. They do not establish the complete preparation, iteration caller or main theorem.
\end{proof}

\begin{lemma}[Marked Finite Deletion Graph Bounds]
\label{bp:markedfinitedeletiongraphbounds}
\lean{Oberwolfach.finite_physical_arc_graph_degree,Oberwolfach.finite_spoke_graph_degree_counts,Oberwolfach.finite_physical_arc_graph_into_set,Oberwolfach.finite_spoke_graph_into_sets,Oberwolfach.marked_finite_deletion_graph_bounds}
\leanok
The finite deleted arc and spoke sets determine deletion graphs with the required global degrees and neighbourhood bounds into the protected sets, including external roots.
\end{lemma}
\begin{proof}
\leanok
The named declarations prove the displayed conditional results. They do not establish the complete preparation, iteration caller or main theorem.
\end{proof}

\begin{lemma}[Marked Neighbourhood Deletion On Domain]
\label{bp:markedneighbourhooddeletionondomain}
\lean{Oberwolfach.marked_row_rim_deletion_into_domain,Oberwolfach.marked_neighbourhoods_after_deletion_on_domain}
\leanok
Deletion neighbourhood bounds into fixed invariant domains preserve the original row and index estimates with explicit losses. Original targets and ambient root sets are retained.
\end{lemma}
\begin{proof}
\leanok
The named declarations prove the displayed conditional results. They do not establish the complete preparation, iteration caller or main theorem.
\end{proof}

\begin{lemma}[Marked Iteration Neighbourhood From Resources]
\label{bp:markediterationneighbourhoodfromresources}
\lean{Oberwolfach.marked_iteration_neighbourhoods_from_resource_bounds}
\leanok
The actual cover-down incidence output supplies both global and smaller-domain neighbourhood estimates after deletion of the constructed finite arc and spoke sets.
\end{lemma}
\begin{proof}
\leanok
The named declarations prove the displayed conditional results. They do not establish the complete preparation, iteration caller or main theorem.
\end{proof}

\begin{lemma}[Marked Neighbourhood Adaptive Vortex Event]
\label{bp:markedneighbourhoodadaptivevortexevent}
\lean{Oberwolfach.marked_neighbourhood_adaptive_vortex_event}
\leanok
One retained experiment supplies both neighbourhood families on all fixed levels and after every later admissible deletion. The deletion quantifier follows the sampled outcome and no union bound over choices is used.
\end{lemma}
\begin{proof}
\leanok
The named declarations prove the displayed conditional results. They do not establish the complete preparation, iteration caller or main theorem.
\end{proof}

\begin{lemma}[Marked Iteration Local Loads From Resources]
\label{bp:markediterationlocalloadsfromresources}
\lean{Oberwolfach.finite_spoke_graph_comap_star,Oberwolfach.marked_iteration_local_loads_from_resource_bounds}
\leanok
The protected-star resource bound preserves fixed local potential-spoke weights after actual deletion. The normalization remains the original local order and the tested root need not survive.
\end{lemma}
\begin{proof}
\leanok
The named declarations prove the displayed conditional results. They do not establish the complete preparation, iteration caller or main theorem.
\end{proof}

\begin{lemma}[Vortex Sampling Domain Sizes]
\label{bp:vortexsamplingdomainsizes}
\lean{Oberwolfach.nested_contracting_difference_sizes,Oberwolfach.Vortex.number_of_levels_le_order,Oberwolfach.Vortex.sampling_difference_sizes}
\leanok
The actual full and difference vortex domains have the required minimum sizes, separate vertex and index contractions, and number-of-levels bound.
\end{lemma}
\begin{proof}
\leanok
The named declarations prove the stated results under their displayed hypotheses. Preparation, the full analytic iteration and the main theorem remain separate obligations.
\end{proof}

\begin{lemma}[Vortex Successor Weight Families]
\label{bp:vortexsuccessorweightfamilies}
\lean{Oberwolfach.vortex_difference_successor_weight_families}
\leanok
Successor weights are constructed on each actual difference domain before the sampling arrays are supplied. The lengths one and two correspond to the manuscript domains and preserve the external-root degree hypothesis.
\end{lemma}
\begin{proof}
\leanok
The named declarations prove the stated results under their displayed hypotheses. Preparation, the full analytic iteration and the main theorem remain separate obligations.
\end{proof}

\begin{lemma}[Vortex Full Successor Weight Families]
\label{bp:vortexfullsuccessorweightfamilies}
\lean{Oberwolfach.vortex_full_successor_weight_families}
\leanok
Successor weights are constructed separately on every full vortex domain before sampling, with the original degree target and uniform constants.
\end{lemma}
\begin{proof}
\leanok
The named declarations prove the stated results under their displayed hypotheses. Preparation, the full analytic iteration and the main theorem remain separate obligations.
\end{proof}

\begin{lemma}[Marked Iteration Configurations From Resources]
\label{bp:markediterationconfigurationsfromresources}
\lean{Oberwolfach.marked_iteration_configurations_from_resource_bounds}
\leanok
Actual protected-star resource bounds give local cycle and root-omitted path estimates at every ambient root.
\end{lemma}
\begin{proof}
\leanok
The named declarations prove the stated results under their displayed hypotheses. Preparation, the full analytic iteration and the main theorem remain separate obligations.
\end{proof}

\begin{lemma}[Marked Common Configurations After Physical Deletion]
\label{bp:markedcommonconfigurationsafterphysicaldeletion}
\lean{Oberwolfach.marked_common_configurations_after_physical_deletion}
\leanok
Common hub, cycle and root-omitted path lower bounds survive physical rim and spoke deletion, with error increment $4\delta/n$.
\end{lemma}
\begin{proof}
\leanok
The named declarations prove the stated results under their displayed hypotheses. Preparation, the full analytic iteration and the main theorem remain separate obligations.
\end{proof}

\begin{lemma}[Marked Iteration Global Estimates From Resources]
\label{bp:markediterationglobalestimatesfromresources}
\lean{Oberwolfach.marked_finite_global_deletion_graph_bounds,Oberwolfach.marked_iteration_global_loads_from_resource_bounds,Oberwolfach.marked_iteration_common_configurations_from_resource_bounds}
\leanok
Actual finite resource incidence bounds imply the deletion-graph degrees, all potential-root load bounds with fixed coefficients, and common-configuration bounds.
\end{lemma}
\begin{proof}
\leanok
The named declarations prove the stated results under their displayed hypotheses. Preparation, the full analytic iteration and the main theorem remain separate obligations.
\end{proof}

\begin{lemma}[Marked Common Adaptive Vortex Event]
\label{bp:markedcommonadaptivevortexevent}
\lean{Oberwolfach.marked_common_adaptive_vortex_event}
\leanok
One original retained event controls common configurations on every fixed domain for all later permitted deletion graphs, including outcome-dependent deletions.
\end{lemma}
\begin{proof}
\leanok
The named declarations prove the stated results under their displayed hypotheses. Preparation, the full analytic iteration and the main theorem remain separate obligations.
\end{proof}

\begin{lemma}[Marked Protected Cover Down Separated Tolerances]
\label{bp:markedprotectedcoverdownseparatedtolerances}
\lean{Oberwolfach.marked_protected_cover_down_separated_tolerances}
\leanok
The complete protected cover construction permits an analytic tolerance fixed before a separate, arbitrarily smaller correction allowance. Its actual star loss depends only on the correction allowance.
\end{lemma}
\begin{proof}
\leanok
The named declarations prove the stated results under their displayed hypotheses. Preparation, the full analytic iteration and the main theorem remain separate obligations.
\end{proof}

\begin{lemma}[Marked Protected Cover Down Iterable]
\label{bp:markedprotectedcoverdowniterable}
\lean{Oberwolfach.marked_protected_cover_down_iterable}
\leanok
The complete protected cover construction allows the requested star-loss bound to be chosen after the analytic tolerance. This repairs a circular dependency in the earlier Lean iteration interface; it does not assert completion of the iteration.
\end{lemma}
\begin{proof}
\leanok
The named declarations prove the stated results under their displayed hypotheses. Preparation, the full analytic iteration and the main theorem remain separate obligations.
\end{proof}

\begin{lemma}[Marked Reference Input Stability]
\label{bp:markedreferenceinputstability}
\lean{Oberwolfach.marked_spoke_reference_change_after_deletion,Oberwolfach.marked_reference_input_stability}
\leanok
The full-domain degree, cut, common-configuration, row, index and reference-spoke bounds transfer together under explicit numerical margins and actual physical deletion.
\end{lemma}
\begin{proof}
\leanok
The named declarations prove the stated results under their displayed hypotheses. Preparation, the full analytic iteration and the main theorem remain separate obligations.
\end{proof}

\begin{lemma}[Marked Iteration Error Scales]
\label{bp:markediterationerrorscales}
\lean{Oberwolfach.marked_iteration_reference_scale,Oberwolfach.marked_iteration_star_scale,Oberwolfach.marked_iteration_local_normalization,Oberwolfach.marked_iteration_all_root_balance_budget}
\leanok
The reference deletion allowance is chosen before the graph and the star allowance is chosen after the analytic tolerance. The stated local normalization and all-root balance inequalities follow with explicit constants.
\end{lemma}
\begin{proof}
\uses{bp:markedreferenceinputstability}
\leanok
The Lean proofs establish these assertions with the hypotheses stated in the source. The complete analytic iteration and the main theorem are separate, unfinished obligations.
\end{proof}

\begin{lemma}[Marked Packing Loads After Physical Deletion]
\label{bp:markedpackingloadsafterphysicaldeletion}
\lean{Oberwolfach.rawMarkedRestrictHubs_deleteRims,Oberwolfach.marked_packing_loads_after_physical_deletion}
\leanok
Physical rim and spoke deletions change the full potential-load predicate by at most $33(m+1)^4(2B/p_*^{2m+1})\delta/n$, with the same coefficients, hub embedding and potential roots.
\end{lemma}
\begin{proof}
\uses{bp:markediterationglobalestimatesfromresources}
\leanok
The Lean proofs establish these assertions with the hypotheses stated in the source. The complete analytic iteration and the main theorem are separate, unfinished obligations.
\end{proof}

\begin{lemma}[Marked Packing Adaptive Vortex Event]
\label{bp:markedpackingadaptivevortexevent}
\lean{Oberwolfach.marked_packing_adaptive_vortex_event}
\leanok
One event of the original retained experiment controls the fixed weights on every supplied domain and all later permitted physical deletions, even when the deletions depend on the outcome.
\end{lemma}
\begin{proof}
\uses{bp:markedpackingloadsafterphysicaldeletion,bp:marked-retained-load-realization}
\leanok
The Lean proofs establish these assertions with the hypotheses stated in the source. The complete analytic iteration and the main theorem are separate, unfinished obligations.
\end{proof}

\begin{lemma}[Marked Active Subhost Incidence]
\label{bp:markedactivesubhostincidence}
\lean{Oberwolfach.markedActiveSubhostEmbedding_physical,Oberwolfach.markedActiveSubhostEmbedding_incident,Oberwolfach.markedActiveSubhostEmbedding_incidence_card}
\leanok
An endpoint-preserving embedding of a remaining host into its original host preserves both auxiliary and physical incidences. The filtered incidence cardinalities agree for every resource set contained in the remaining host.
\end{lemma}
\begin{proof}
\leanok
The Lean proofs establish these assertions with the hypotheses stated in the source. The complete analytic iteration and the main theorem are separate, unfinished obligations.
\end{proof}

\begin{lemma}[Marked Deleted Host Incidence]
\label{bp:markeddeletedhostincidence}
\lean{Oberwolfach.marked_full_resources_active_image,Oberwolfach.marked_deleted_host_incidence_card}
\leanok
For the actual physically deleted sample, the original and remaining hosts give exactly the same incidence counts on any further removed resources.
\end{lemma}
\begin{proof}
\uses{bp:markedactivesubhostincidence,bp:marked-active-resource-degree-transfer}
\leanok
The Lean proofs establish these assertions with the hypotheses stated in the source. The complete analytic iteration and the main theorem are separate, unfinished obligations.
\end{proof}

\begin{lemma}[Marked Fixed Reference Cover Remainder]
\label{bp:markedfixedreferencecoverremainder}
\lean{Oberwolfach.fixed_reference_after_cover,Oberwolfach.marked_cover_remainder_fixed_reference}
\leanok
The next remainder is the next fixed reference graph minus exactly the current covering consumption. Earlier deletions that avoid the inner pair contribute no additional loss, and both incidence bounds transfer without an extra factor.
\end{lemma}
\begin{proof}
\uses{bp:markeddeletedhostincidence}
\leanok
The Lean proofs establish these assertions with the hypotheses stated in the source. The complete analytic iteration and the main theorem are separate, unfinished obligations.
\end{proof}

\begin{lemma}[Marked Unchanged Inner Sample]
\label{bp:markedunchangedinnersample}
\lean{Oberwolfach.marked_removed_resources_empty_inner,Oberwolfach.marked_raw_pair_unchanged_on_inner}
\leanok
An empty physical deletion on an inner pair leaves both the restricted raw sample and its restricted spoke relations exactly unchanged.
\end{lemma}
\begin{proof}
\uses{bp:marked-active-resource-degree-transfer}
\leanok
The Lean proofs establish these assertions with the hypotheses stated in the source. The complete analytic iteration and the main theorem are separate, unfinished obligations.
\end{proof}

\begin{lemma}[Marked Induced Occurrence Matrices]
\label{bp:markedinducedoccurrencematrices}
\lean{Oberwolfach.occurrenceMultiplicity_induced,Oberwolfach.markedSampleArc_induced_multiplicity,Oberwolfach.marked_induced_rim_matrix_unchanged,Oberwolfach.marked_induced_occurrence_cuts_unchanged}
\leanok
The induced occurrence matrices are the restrictions of the original matrices. On a permutation-invariant inner set, an empty physical deletion preserves these matrices and their cut bounds.
\end{lemma}
\begin{proof}
\uses{bp:markedunchangedinnersample}
\leanok
The Lean proofs establish these assertions with the hypotheses stated in the source. The complete analytic iteration and the main theorem are separate, unfinished obligations.
\end{proof}

\begin{lemma}[Marked Actual Iteration Balances]
\label{bp:markedactualiterationbalances}
\lean{Oberwolfach.marked_actual_iteration_balances}
\leanok
The exact canonical deleted sample satisfies the next all-root restricted balance bound. At inner roots, its balance values equal those of the original reference sample.
\end{lemma}
\begin{proof}
\uses{bp:markediterationbalancefromresources,bp:marked-protected-active-estimates}
\leanok
The Lean proofs establish these assertions with the hypotheses stated in the source. The complete analytic iteration and the main theorem are separate, unfinished obligations.
\end{proof}

\begin{lemma}[Marked Induced Balance Inputs]
\label{bp:markedinducedbalanceinputs}
\lean{Oberwolfach.markedFiniteDomainBalance_univ,Oberwolfach.marked_induced_full_balance_eq,Oberwolfach.marked_induced_balance_inputs}
\leanok
The four numerical near-balance hypotheses for inner correction follow from the original fixed-domain balances at inner roots, with the marked coefficient and occurrence identifiers preserved.
\end{lemma}
\begin{proof}
\uses{bp:marked-protected-active-estimates}
\leanok
The Lean proofs establish these assertions with the hypotheses stated in the source. The complete analytic iteration and the main theorem are separate, unfinished obligations.
\end{proof}

\section{Actual marked balance events and local iteration profiles}

\begin{lemma}[Actual inner correction inputs]\label{bp:markedactualinnercorrectioninputs}
\lean{Oberwolfach.marked_actual_inner_correction_inputs}
\leanok
The four near-balance inequalities and all four cuts in the actual remaining sample follow from the fixed reference when the earlier deletion avoids the inner pair.
\end{lemma}
\begin{proof}\leanok
The linked Lean declaration verifies the stated finite objects and their exact substitutions. This result does not by itself prove the main theorem.
\end{proof}

\begin{lemma}[Restricted profiles with external hubs]\label{bp:markediterationrestrictedprofiles}
\lean{Oberwolfach.marked_iteration_restricted_profiles}
\leanok
The original fixed inner row and cycle/path profiles survive every later protected-star deletion. The hub embedding is unrestricted, and the error retains its original local scale.
\end{lemma}
\begin{proof}\leanok
The linked Lean declaration verifies the stated finite objects and their exact substitutions. This result does not by itself prove the main theorem.
\end{proof}

\begin{lemma}[Sampled directed degrees as occurrence degrees]\label{bp:markedrestrictedsampledegrees}
\lean{Oberwolfach.markedSampleArc_opposite_domain_degrees}
\leanok
Degrees into a fixed opposite-endpoint domain equal the degrees of the explicitly restricted sampled occurrence types.
\end{lemma}
\begin{proof}\leanok
The linked Lean declaration verifies the stated finite objects and their exact substitutions. This result does not by itself prove the main theorem.
\end{proof}

\begin{lemma}[Finite balance values from sampled profiles]\label{bp:markedrestrictedsamplebalance}
\lean{Oberwolfach.marked_restricted_sample_balance}
\leanok
The four sampled restricted-balance inequalities give the four finite balance differences at every ambient root, with exactly the same error.
\end{lemma}
\begin{proof}\leanok
The linked Lean declaration verifies the stated finite objects and their exact substitutions. This result does not by itself prove the main theorem.
\end{proof}

\begin{lemma}[Actual deletion balance errors]\label{bp:markedactualiterationbalanceerror}
\lean{Oberwolfach.marked_actual_iteration_balance_error}
\leanok
A later deletion of protected-star incidence at most delta changes the balance by at most (m+2)delta. Inner-root values are unchanged when no inner resource is deleted.
\end{lemma}
\begin{proof}\leanok
The linked Lean declaration verifies the stated finite objects and their exact substitutions. This result does not by itself prove the main theorem.
\end{proof}

\begin{lemma}[One common finite-balance event]\label{bp:markedfinitebalanceadaptivevortexevent}
\lean{Oberwolfach.marked_finite_balance_adaptive_vortex_event}
\leanok
The original retained sampling law supplies finite balance bounds on every fixed vortex domain and after every later eligible deletion. No resampling or outcome-dependent domain enlargement is used.
\end{lemma}
\begin{proof}\leanok
The linked Lean declaration verifies the stated finite objects and their exact substitutions. This result does not by itself prove the main theorem.
\end{proof}

\section{Reindexing the marked remainder and preserving iteration margins}

\begin{lemma}[Active Vertex Incidence]\label{bp:markedactivevertexincidence}
\lean{Oberwolfach.markedActiveVertexEmbedding_physical,Oberwolfach.markedActiveVertexEmbedding_incident,Oberwolfach.markedActiveVertexEmbedding_incidence_card_le}
\leanok
Reindexing a typed host on smaller invariant vertex sets preserves both physical and auxiliary incidences. In particular, restriction does not enlarge a bound on the incidence of a specified resource set.
\end{lemma}
\begin{proof}\leanok The stated transport or numerical bound is proved for the actual finite objects in the corresponding Lean module. It does not assert completion of the main theorem. \end{proof}

\begin{lemma}[Resource Star Embedding]\label{bp:markedresourcestarembedding}
\lean{Oberwolfach.markedResourceEmbedding_star_iff,Oberwolfach.markedResourceEmbedding_star_filter}
\leanok
Protected stars commute with the actual embeddings of the nested row and hub domains.
\end{lemma}
\begin{proof}\leanok The stated transport or numerical bound is proved for the actual finite objects in the corresponding Lean module. It does not assert completion of the main theorem. \end{proof}

\begin{lemma}[Reindexed Fixed Remainder]\label{bp:markedreindexedfixedremainder}
\lean{Oberwolfach.marked_reindexed_fixed_remainder}
\leanok
Transport the next fixed-reference error to its actual smaller vertex types. Both incidence bounds retain their constants, and the deletion is empty on the next protected inner pair.
\end{lemma}
\begin{proof}\leanok The stated transport or numerical bound is proved for the actual finite objects in the corresponding Lean module. It does not assert completion of the main theorem. \end{proof}

\begin{lemma}[Restricted Finite Images]\label{bp:markedrestrictedfiniteimages}
\lean{Oberwolfach.marked_restricted_finite_images}
\leanok
A remaining arc and spoke set supported on fixed embedded domains has an exact preimage in the restricted sample. The original occurrence IDs are preserved by the two sample embeddings.
\end{lemma}
\begin{proof}\leanok The stated transport or numerical bound is proved for the actual finite objects in the corresponding Lean module. It does not assert completion of the main theorem. \end{proof}

\begin{lemma}[Finite Divisibility Reflection]\label{bp:markedfinitedivisibilityreflection}
\lean{Oberwolfach.marked_finite_divisibility_reflects}
\leanok
Exact reflection of the four divisibility identities under injective row, hub and typed-occurrence maps.
\end{lemma}
\begin{proof}\leanok The stated transport or numerical bound is proved for the actual finite objects in the corresponding Lean module. It does not assert completion of the main theorem. \end{proof}

\begin{lemma}[Restricted Finite Remainder]\label{bp:markedrestrictedfiniteremainder}
\lean{Oberwolfach.marked_restricted_finite_remainder}
\leanok
Reindex the actual remainder, with its four exact integer balances, on the next vertex pair. All finite resources and opposite-endpoint balance values are transported by their specified embeddings.
\end{lemma}
\begin{proof}\leanok The stated transport or numerical bound is proved for the actual finite objects in the corresponding Lean module. It does not assert completion of the main theorem. \end{proof}

\begin{lemma}[Restricted Reference Resources]\label{bp:markedrestrictedreferenceresources}
\lean{Oberwolfach.marked_restricted_reference_resources}
\leanok
The reference sample on a nested invariant pair is exactly the induced resource set of the original sample. This concerns the fixed reference, before any adaptive deletion.
\end{lemma}
\begin{proof}\leanok The stated transport or numerical bound is proved for the actual finite objects in the corresponding Lean module. It does not assert completion of the main theorem. \end{proof}

\begin{lemma}[Local Estimate Monotonicity]\label{bp:markedlocalestimatemonotonicity}
\lean{Oberwolfach.markedLocalConfigurations_mono,Oberwolfach.markedLocalSpokeLoads_mono}
\leanok
Weakening numerical bounds retains the same domains and coefficients.
\end{lemma}
\begin{proof}\leanok The stated transport or numerical bound is proved for the actual finite objects in the corresponding Lean module. It does not assert completion of the main theorem. \end{proof}

\begin{lemma}[Finite Resource Preimages]\label{bp:markedfiniteresourcepreimages}
\lean{Oberwolfach.marked_finite_resource_preimages,Oberwolfach.marked_finite_resources_determine_sets}
\leanok
Recover the actual typed arc and spoke sets of any finite subset of a host's physical resources. The construction is used for the next error set after cover-down, rather than assuming a new typed representation.
\end{lemma}
\begin{proof}\leanok The stated transport or numerical bound is proved for the actual finite objects in the corresponding Lean module. It does not assert completion of the main theorem. \end{proof}

\begin{lemma}[Global Cover Input Margins]\label{bp:markedglobalcoverinputmargins}
\lean{Oberwolfach.marked_global_cover_input_margins}
\leanok
A single global deletion allowance supplies the potential-load and deep-stage configuration inputs at every cover-down level. The weights, root predicates and normalizing order are fixed before deletion.
\end{lemma}
\begin{proof}\leanok The stated transport or numerical bound is proved for the actual finite objects in the corresponding Lean module. It does not assert completion of the main theorem. \end{proof}

\begin{lemma}[Local Cover Input Margins]\label{bp:markedlocalcoverinputmargins}
\lean{Oberwolfach.marked_local_cover_input_margins}
\leanok
The actual protected-star incidence allowance supplies the local cover-down inputs with their fixed local order and fixed coefficient function. The hub domain may include hubs outside the protected pair.
\end{lemma}
\begin{proof}\leanok The stated transport or numerical bound is proved for the actual finite objects in the corresponding Lean module. It does not assert completion of the main theorem. \end{proof}

\begin{lemma}[Next Sample Reference Error]\label{bp:markednextsamplereferenceerror}
\lean{Oberwolfach.marked_next_sample_reference_error}
\leanok
The actual next reference sample inherits the cover-down deletion bounds under its specified row, hub and typed-arc maps. The error is the newly removed inner resource set, not a sum of errors from earlier levels.
\end{lemma}
\begin{proof}\leanok The stated transport or numerical bound is proved for the actual finite objects in the corresponding Lean module. It does not assert completion of the main theorem. \end{proof}

\begin{lemma}[Iteration Tolerance Choices]\label{bp:markediterationtolerancechoices}
\lean{Oberwolfach.marked_iteration_global_tolerance_choice,Oberwolfach.marked_iteration_local_tolerance_choice}
\leanok
Explicit choices of the two uniform deletion allowances. The global allowance precedes the later analytic tolerance x; the star allowance is chosen after x and before the inner balance and order thresholds.
\end{lemma}
\begin{proof}\leanok The stated transport or numerical bound is proved for the actual finite objects in the corresponding Lean module. It does not assert completion of the main theorem. \end{proof}

\begin{lemma}[Canonical Remainder State]\label{bp:markedcanonicalremainderstate}
\lean{Oberwolfach.marked_canonical_remainder_divisibility}
\leanok
Identifying the resource set of a remaining host also identifies its typed arc and spoke complements. Its exact divisibility then holds in the canonical deleted sample, with no added balance hypothesis.
\end{lemma}
\begin{proof}\leanok The stated transport or numerical bound is proved for the actual finite objects in the corresponding Lean module. It does not assert completion of the main theorem. \end{proof}

\begin{lemma}[Next Sample Canonical Remainder]\label{bp:markednextsamplecanonicalremainder}
\lean{Oberwolfach.marked_next_sample_canonical_remainder}
\leanok
The actual divisible remainder is the canonical deletion of the next fixed reference sample. This joins the exact occurrence reindexing with the two incidence bounds needed by the next cover-down call.
\end{lemma}
\begin{proof}\leanok The stated transport or numerical bound is proved for the actual finite objects in the corresponding Lean module. It does not assert completion of the main theorem. \end{proof}

\begin{lemma}[Resource Domain Monotonicity]\label{bp:markedresourcedomainmonotonicity}
\lean{Oberwolfach.markedStarResource_mono,Oberwolfach.markedInnerResource_mono}
\leanok
Nested row and hub domains induce nested physical stars and inner resource sets.
\end{lemma}
\begin{proof}\leanok The stated transport or numerical bound is proved for the actual finite objects in the corresponding Lean module. It does not assert completion of the main theorem. \end{proof}

\begin{lemma}[Nested Iteration Balance Inputs]\label{bp:markednestediterationbalanceinputs}
\lean{Oberwolfach.marked_nested_iteration_balance_inputs}
\leanok
The actual physically deleted sample has balance error at most the original absolute error plus (m+2) times the realised protected-star bound. At inner roots its balance values are unchanged. No normalization is altered.
\end{lemma}
\begin{proof}\leanok The stated transport or numerical bound is proved for the actual finite objects in the corresponding Lean module. It does not assert completion of the main theorem. \end{proof}

\section{The marked cover-down step from fixed reference estimates}

\begin{lemma}[Current Packing Remainder]\label{bp:markedcurrentpackingremainder}
\lean{Oberwolfach.marked_current_packing_original_remainder}
\leanok
Lift the actual packing in the current deleted sample to the original occurrence identifiers and compute its remaining typed host. Divisibility is taken from the current host, not from the fixed reference or its earlier deleted edge set.
\end{lemma}
\begin{proof}\leanok The packing and the transfer are derived from the displayed fixed reference estimates and actual finite deletions. Uniform allowances are chosen before graph orders and the separation parameter. \end{proof}

\begin{lemma}[Finite Resource Support]\label{bp:markedfiniteresourcesupport}
\lean{Oberwolfach.markedFiniteResources_support_iff}
\leanok
Equivalence between support of the actual typed sets and support of their physical resources on an invariant inner pair.
\end{lemma}
\begin{proof}\leanok The packing and the transfer are derived from the displayed fixed reference estimates and actual finite deletions. Uniform allowances are chosen before graph orders and the separation parameter. \end{proof}

\begin{lemma}[Cover Remainder As Fixed Complement]\label{bp:markedcoverremainderasfixedcomplement}
\lean{Oberwolfach.marked_cover_remainder_fixed_inner_complement}
\leanok
A cover-down packing supported outside the remaining inner host leaves exactly that fixed inner reference minus its newly consumed edges. The statement uses the actual original arc and spoke complements.
\end{lemma}
\begin{proof}\leanok The packing and the transfer are derived from the displayed fixed reference estimates and actual finite deletions. Uniform allowances are chosen before graph orders and the separation parameter. \end{proof}

\begin{lemma}[Cover Down Next State]\label{bp:markedcoverdownnextstate}
\lean{Oberwolfach.marked_cover_down_next_state}
\leanok
The actual divisible remainder is the canonical deletion of the next fixed reference sample. This joins the exact occurrence reindexing with the two incidence bounds needed by the next cover-down call.
\end{lemma}
\begin{proof}\leanok The packing and the transfer are derived from the displayed fixed reference estimates and actual finite deletions. Uniform allowances are chosen before graph orders and the separation parameter. \end{proof}

\begin{lemma}[Deleted Cover Fixed Remainder]\label{bp:markeddeletedcoverfixedremainder}
\lean{Oberwolfach.marked_deleted_cover_fixed_remainder}
\leanok
The output of cover-down in the actual deleted sample supplies a balanced remainder on the original occurrence identifiers. The removed inner resource set and both incidence bounds are measured in the fixed reference, with no earlier deletion added to the next error set.
\end{lemma}
\begin{proof}\leanok The packing and the transfer are derived from the displayed fixed reference estimates and actual finite deletions. Uniform allowances are chosen before graph orders and the separation parameter. \end{proof}

\begin{lemma}[Protected Cover Down From Fixed Reference]\label{bp:markedprotectedcoverdownfromfixedreference}
\lean{Oberwolfach.marked_protected_cover_down_from_fixed_reference}
\leanok
Apply protected cover-down to the actual finite deletion of a fixed reference. All analytic estimates of the current sample are derived from those of the reference and the displayed global and star allowances.
\end{lemma}
\begin{proof}\leanok The packing and the transfer are derived from the displayed fixed reference estimates and actual finite deletions. Uniform allowances are chosen before graph orders and the separation parameter. \end{proof}

\begin{lemma}[Fixed Reference Cover Down Step]\label{bp:markedfixedreferencecoverdownstep}
\lean{Oberwolfach.marked_fixed_reference_cover_down_step}
\leanok
Apply protected cover-down to the actual finite deletion of a fixed reference. All analytic estimates of the current sample are derived from those of the reference and the displayed global and star allowances.
\end{lemma}
\begin{proof}\leanok The packing and the transfer are derived from the displayed fixed reference estimates and actual finite deletions. Uniform allowances are chosen before graph orders and the separation parameter. \end{proof}

\begin{lemma}[Uniform Fixed Reference Cover Down Step]\label{bp:markeduniformfixedreferencecoverdownstep}
\lean{Oberwolfach.marked_uniform_fixed_reference_cover_down_step}
\leanok
Apply protected cover-down to the actual finite deletion of a fixed reference. All analytic estimates of the current sample are derived from those of the reference and the displayed global and star allowances.
\end{lemma}
\begin{proof}\leanok The packing and the transfer are derived from the displayed fixed reference estimates and actual finite deletions. Uniform allowances are chosen before graph orders and the separation parameter. \end{proof}

\begin{lemma}[Uniform Cover Down Next State]\label{bp:markeduniformcoverdownnextstate}
\lean{Oberwolfach.marked_uniform_cover_down_next_state}
\leanok
Apply protected cover-down to the actual finite deletion of a fixed reference. All analytic estimates of the current sample are derived from those of the reference and the displayed global and star allowances.
\end{lemma}
\begin{proof}\leanok The packing and the transfer are derived from the displayed fixed reference estimates and actual finite deletions. Uniform allowances are chosen before graph orders and the separation parameter. \end{proof}

\section{Exact typed host encoding and the first cover-down step}

\begin{lemma}[Host Rim Encoding]\label{bp:markedhostrimencoding}
\lean{Oberwolfach.markedHostPhysicalCoordinate_injective,Oberwolfach.markedHostRimEncoding_at,Oberwolfach.markedHostRimEncoding_exists_of_some,Oberwolfach.markedHostRimEncoding_arcs}
\leanok
A physically injective typed host can be represented by categorical rim coordinates. This is deterministic encoding of a given host, not a claim that it has the distribution of a fresh random sample.
\end{lemma}
\begin{proof}\leanok The actual typed occurrence maps preserve the original physical resources, directions and roots. The final cover-down packing is obtained from the proved deterministic constructor and transferred back to the given host. No distribution is assigned to the encoded outcome. \end{proof}

\begin{lemma}[Host Arc Equivalences]\label{bp:markedhostarcequivalences}
\lean{Oberwolfach.markedHostOrdinaryArcEquiv_endpoints,Oberwolfach.markedHostMarkedArcEquiv_endpoints}
\leanok
The categorical encoding preserves the actual ordinary and marked arc identifiers bijectively, including both directed endpoints.
\end{lemma}
\begin{proof}\leanok The actual typed occurrence maps preserve the original physical resources, directions and roots. The final cover-down packing is obtained from the proved deterministic constructor and transferred back to the given host. No distribution is assigned to the encoded outcome. \end{proof}

\begin{lemma}[Host Raw Encoding]\label{bp:markedhostrawencoding}
\lean{Oberwolfach.rawMarkedEncodeHost_spokes,Oberwolfach.rawMarkedEncodeHost_inner_coordinates,Oberwolfach.rawMarkedEncodeHost_restrict,Oberwolfach.rawMarkedEncodeHost_rims}
\leanok
Encoding a modified typed host keeps every original interval coordinate. On a protected inner row set the original physical coordinates can also be kept exactly. The operation makes no assertion about a new probability law.
\end{lemma}
\begin{proof}\leanok The actual typed occurrence maps preserve the original physical resources, directions and roots. The final cover-down packing is obtained from the proved deterministic constructor and transferred back to the given host. No distribution is assigned to the encoded outcome. \end{proof}

\begin{lemma}[Raw Host Arc Equivalences]\label{bp:markedrawhostarcequivalences}
\lean{Oberwolfach.marked_raw_host_encoding_exact}
\leanok
Exact occurrence maps for the encoded modified host, while its protected inner coordinates and all interval coordinates are retained.
\end{lemma}
\begin{proof}\leanok The actual typed occurrence maps preserve the original physical resources, directions and roots. The final cover-down packing is obtained from the proved deterministic constructor and transferred back to the given host. No distribution is assigned to the encoded outcome. \end{proof}

\begin{lemma}[Host Reindexing]\label{bp:markedhostreindexing}
\lean{Oberwolfach.markedResourceEmbedding_refl,Oberwolfach.markedHostReindex_resources,Oberwolfach.markedHostReindex_balance,Oberwolfach.markedHostReindex_domain_balance,Oberwolfach.markedHostReindex_divisible_iff,Oberwolfach.markedHostReindex_decomposition_iff}
\leanok
Reindex typed occurrences without changing vertices, hubs, or physical resources. The balance comparisons hold at every original root and on the same supplied domains.
\end{lemma}
\begin{proof}\leanok The actual typed occurrence maps preserve the original physical resources, directions and roots. The final cover-down packing is obtained from the proved deterministic constructor and transferred back to the given host. No distribution is assigned to the encoded outcome. \end{proof}

\begin{lemma}[Encoded Host Invariants]\label{bp:markedencodedhostinvariants}
\lean{Oberwolfach.marked_encoded_host_invariants}
\leanok
Exact finite resource, balance and decomposition statements for a modified host encoded in raw coordinates. No random law is assigned to the encoded outcome.
\end{lemma}
\begin{proof}\leanok The actual typed occurrence maps preserve the original physical resources, directions and roots. The final cover-down packing is obtained from the proved deterministic constructor and transferred back to the given host. No distribution is assigned to the encoded outcome. \end{proof}

\begin{lemma}[Host Packing Reindexing]\label{bp:markedhostpackingreindexing}
\lean{Oberwolfach.markedHostReindex_packing,Oberwolfach.markedHostReindex_packing_remainder}
\leanok
Reindex a packing between equivalent typed hosts and identify its actual remainder. Divisibility is derived by subtracting the packing in the destination host, not postulated for a resource set.
\end{lemma}
\begin{proof}\leanok The actual typed occurrence maps preserve the original physical resources, directions and roots. The final cover-down packing is obtained from the proved deterministic constructor and transferred back to the given host. No distribution is assigned to the encoded outcome. \end{proof}

\begin{lemma}[Host Reindex Incidences]\label{bp:markedhostreindexincidences}
\lean{Oberwolfach.markedHostReindex_incidence_card}
\leanok
Both physical and auxiliary incidence counts are invariant under exact typed reindexing. All roots are retained, and the removed resource set is compared without multiplying it by a sampling probability.
\end{lemma}
\begin{proof}\leanok The actual typed occurrence maps preserve the original physical resources, directions and roots. The final cover-down packing is obtained from the proved deterministic constructor and transferred back to the given host. No distribution is assigned to the encoded outcome. \end{proof}

\begin{lemma}[Host Cover Reindexing]\label{bp:markedhostcoverreindexing}
\lean{Oberwolfach.markedHostReindex_cover}
\leanok
Transfer an actual cover-down packing through a typed host equivalence, including its divisible supported remainder and both next-level bounds.
\end{lemma}
\begin{proof}\leanok The actual typed occurrence maps preserve the original physical resources, directions and roots. The final cover-down packing is obtained from the proved deterministic constructor and transferred back to the given host. No distribution is assigned to the encoded outcome. \end{proof}

\begin{lemma}[Encoded Host Matrices]\label{bp:markedencodedhostmatrices}
\lean{Oberwolfach.marked_encoded_host_matrices}
\leanok
The rim matrices of the encoded outcome are exactly the original ordinary and marked occurrence matrices, with directions unchanged.
\end{lemma}
\begin{proof}\leanok The actual typed occurrence maps preserve the original physical resources, directions and roots. The final cover-down packing is obtained from the proved deterministic constructor and transferred back to the given host. No distribution is assigned to the encoded outcome. \end{proof}

\begin{lemma}[Protected Cover Down Typed Host]\label{bp:markedprotectedcoverdowntypedhost}
\lean{Oberwolfach.marked_protected_cover_down_typed_host}
\leanok
The protected cover-down theorem in an order usable for iteration. The analytic input tolerance is chosen before the arbitrary requested star loss. The inner correction is then smaller than that loss; only inner-root imbalance needs the still smaller correction tolerance. All packings are constructed from the displayed graph inputs.
\end{lemma}
\begin{proof}\leanok The actual typed occurrence maps preserve the original physical resources, directions and roots. The final cover-down packing is obtained from the proved deterministic constructor and transferred back to the given host. No distribution is assigned to the encoded outcome. \end{proof}

\section{Reference states, reindexing and bounded cycle weights}

\begin{definition}[Reindexed Remainder]\label{bp:markedreindexedremainder-definitions}
\lean{Oberwolfach.MarkedReindexedRemainder}
\leanok
The canonical next-level remainder statement, as a predicate on its fixed raw outcome and spoke relation. This keeps equality transport of the whole dependent finite sample explicit.
\end{definition}

\begin{lemma}[Encoded Cover Next State]\label{bp:markedencodedcovernextstate}
\lean{Oberwolfach.marked_encoded_cover_next_state}
\leanok
The actual divisible remainder is the canonical deletion of the next fixed reference sample. This joins the exact occurrence reindexing with the two incidence bounds needed by the next cover-down call.
\end{lemma}
\begin{proof}\leanok The stated conclusion follows from the displayed hypotheses, via the checked Lean proof. This is not a proof of Theorem~1.1. \end{proof}

\begin{lemma}[Typed Host Cover Next State]\label{bp:markedtypedhostcovernextstate}
\lean{Oberwolfach.marked_typed_host_cover_next_state}
\leanok
The protected cover-down theorem in an order usable for iteration. The analytic input tolerance is chosen before the arbitrary requested star loss. The inner correction is then smaller than that loss; only inner-root imbalance needs the still smaller correction tolerance. All packings are constructed from the displayed graph inputs.
\end{lemma}
\begin{proof}\leanok The stated conclusion follows from the displayed hypotheses, via the checked Lean proof. This is not a proof of Theorem~1.1. \end{proof}

\begin{lemma}[Reindexed Resource Iteration]\label{bp:reindexedresourceiteration}
\lean{Oberwolfach.reindexed_resource_iteration,Oberwolfach.reindexed_resource_iteration_absorption}
\leanok
Finite iteration with a different state type at each level. All resource sets are compared in one ambient ground set; no identification of local vertex types or of local coefficient functions is made. The step constructor is an explicit premise of this algebraic lemma and must be supplied by the graph-specific cover-down theorem.
\end{lemma}
\begin{proof}\leanok The stated conclusion follows from the displayed hypotheses, via the checked Lean proof. This is not a proof of Theorem~1.1. \end{proof}

\begin{definition}[Reference Deletion State]\label{bp:markedreferencedeletionstate-definitions}
\lean{Oberwolfach.MarkedReferenceDeletion,Oberwolfach.MarkedReferenceDeletion.removed,Oberwolfach.MarkedReferenceDeletion.remaining,Oberwolfach.MarkedReferenceDeletion.Valid}
\leanok
A vortex state records the actual deleted typed arcs and spoke pairs of one fixed reference. Validity consists of exact divisibility and the two specified incidence bounds, with no deletion in the next inner pair.
\end{definition}

\begin{lemma}[Reference Deletion State]\label{bp:markedreferencedeletionstate}
\lean{Oberwolfach.marked_reindexed_remainder_state}
\leanok
A vortex state records the actual deleted typed arcs and spoke pairs of one fixed reference. Validity consists of exact divisibility and the two specified incidence bounds, with no deletion in the next inner pair.
\end{lemma}
\begin{proof}\leanok The stated conclusion follows from the displayed hypotheses, via the checked Lean proof. This is not a proof of Theorem~1.1. \end{proof}

\begin{lemma}[Reference State Remainder]\label{bp:markedreferencestateremainder}
\lean{Oberwolfach.MarkedReferenceDeletion.remaining_eq,Oberwolfach.MarkedReferenceDeletion.removed_subset,Oberwolfach.MarkedReferenceDeletion.original_remainder_divisible}
\leanok
Recover the divisible typed remainder of a valid reference-deletion state in the original occurrence identifiers.
\end{lemma}
\begin{proof}\leanok The stated conclusion follows from the displayed hypotheses, via the checked Lean proof. This is not a proof of Theorem~1.1. \end{proof}

\begin{lemma}[Uniform Reference State Step]\label{bp:markeduniformreferencestatestep}
\lean{Oberwolfach.marked_uniform_reference_state_step}
\leanok
Apply protected cover-down to the actual finite deletion of a fixed reference. All analytic estimates of the current sample are derived from those of the reference and the displayed global and star allowances.
\end{lemma}
\begin{proof}\leanok The stated conclusion follows from the displayed hypotheses, via the checked Lean proof. This is not a proof of Theorem~1.1. \end{proof}

\begin{lemma}[Reference State Step Decomposition]\label{bp:markedreferencestatestepdecomposition}
\lean{Oberwolfach.marked_reference_state_step_decomposition}
\leanok
The exact resources covered in one reference-state step. The packing must avoid the recorded deletion; the remaining resource identity is then proved by finite set subtraction.
\end{lemma}
\begin{proof}\leanok The stated conclusion follows from the displayed hypotheses, via the checked Lean proof. This is not a proof of Theorem~1.1. \end{proof}

\begin{lemma}[Resource Decomposition Embedding]\label{bp:markedresourcedecompositionembedding}
\lean{Oberwolfach.marked_resource_decomposition_embedding}
\leanok
Exact decompositions embed along permutation-compatible vertex and hub maps and type-preserving arc maps. The covered physical set is transported by the same injective resource map as the individual wheels.
\end{lemma}
\begin{proof}\leanok The stated conclusion follows from the displayed hypotheses, via the checked Lean proof. This is not a proof of Theorem~1.1. \end{proof}

\begin{lemma}[Positive Ordinary Cycle Weights And Pair Loads]\label{bp:positiveordinarycycleweightsandpairloads}
\lean{Oberwolfach.positive_ordinary_cycle_weights_with_pair_loads}
\leanok
The positive bounded cycle weighting and both physical pair-load bounds hold simultaneously for the same weighting. These are the conclusions of Theorem 8.2 needed by the ordinary wheel construction.
\end{lemma}
\begin{proof}\leanok The stated conclusion follows from the displayed hypotheses, via the checked Lean proof. This is not a proof of Theorem~1.1. \end{proof}

\section{Reference states, reindexing and bounded cycle weights}

\begin{lemma}[Constructed Reference State Step]\label{bp:markedconstructedreferencestatestep}
\lean{Oberwolfach.marked_constructed_reference_state_step}
\leanok
Apply protected cover-down to the actual finite deletion of a fixed reference. All analytic estimates of the current sample are derived from those of the reference and the displayed global and star allowances.
\end{lemma}
\begin{proof}\leanok The stated conclusion follows from the displayed hypotheses, via the checked Lean proof. This is not a proof of Theorem~1.1. \end{proof}

\section{Reference states, reindexing and bounded cycle weights}

\begin{definition}[Reference Level Data]\label{bp:markedreferenceleveldata-definitions}
\lean{Oberwolfach.MarkedReferenceLevelData}
\leanok
Fixed inputs for one step of marked cover-down. These are exactly the explicit deterministic graph, profile, geometric and weight hypotheses of that step, before any choice of the current bounded deletion.
\end{definition}

\begin{lemma}[Uniform Reference Level Constructor]\label{bp:markeduniformreferencelevelconstructor}
\lean{Oberwolfach.marked_uniform_reference_level_constructor}
\leanok
Apply protected cover-down to the actual finite deletion of a fixed reference. All analytic estimates of the current sample are derived from those of the reference and the displayed global and star allowances.
\end{lemma}
\begin{proof}\uses{bp:markedconstructedreferencestatestep}\leanok The stated conclusion follows from the displayed hypotheses, via the checked Lean proof. This is not a proof of Theorem~1.1. \end{proof}

\begin{definition}[Reference Level Step Property]\label{bp:markedreferencelevelstepproperty-definitions}
\lean{Oberwolfach.MarkedReferenceLevelStepProperty}
\leanok
The conclusion needed for iteration at one level: a packing constructed in the current reference host, exact decomposition of the discarded part, and a bounded divisible state on every compatible next reference. This is a predicate on the explicit graph/profile data, not an axiom or an assertion that these inputs hold for the sampled vortex.
\end{definition}

\begin{lemma}[Uniform Reference Level Step]\label{bp:markeduniformreferencelevelstep}
\lean{Oberwolfach.marked_uniform_reference_level_step}
\leanok
Apply protected cover-down to the actual finite deletion of a fixed reference. All analytic estimates of the current sample are derived from those of the reference and the displayed global and star allowances.
\end{lemma}
\begin{proof}\uses{bp:markeduniformreferencelevelconstructor}\leanok The stated conclusion follows from the displayed hypotheses, via the checked Lean proof. This is not a proof of Theorem~1.1. \end{proof}

\begin{lemma}[Reference State Transport]\label{bp:markedreferencestatetransport}
\lean{Oberwolfach.marked_reference_state_transport}
\leanok
A reference-deletion state is unchanged by equal descriptions of the underlying graph, raw outcome and spoke relations. In particular, the actual deleted and remaining physical sets and the full validity condition are transported together.
\end{lemma}
\begin{proof}\leanok The stated conclusion follows from the displayed hypotheses, via the checked Lean proof. This is not a proof of Theorem~1.1. \end{proof}

\begin{definition}[Reference Level Chain]\label{bp:markedreferencelevelchain-definitions}
\lean{Oberwolfach.MarkedReferenceLevelChain}
\leanok
A fixed sequence of nested reference graphs and their actual restrictions. The equalities express compatibility of the prescribed sample and spoke relations; no sampling law or cover-down packing is assumed here.
\end{definition}

\begin{definition}[Reference Level Embedding]\label{bp:markedreferencelevelembedding-definitions}
\lean{Oberwolfach.markedReferenceLevelArcEmbedding}
\leanok
Embed every typed occurrence and every wheel from a fixed vortex reference into the first reference, preserving the actual physical rim and spoke resources.
\end{definition}

\begin{lemma}[Reference Level Embedding]\label{bp:markedreferencelevelembedding}
\lean{Oberwolfach.markedReferenceLevelArcEmbedding_endpoints,Oberwolfach.marked_reference_level_decomposition_embedding}
\leanok
Embed every typed occurrence and every wheel from a fixed vortex reference into the first reference, preserving the actual physical rim and spoke resources.
\end{lemma}
\begin{proof}\uses{bp:markedresourcedecompositionembedding}\leanok The stated conclusion follows from the displayed hypotheses, via the checked Lean proof. This is not a proof of Theorem~1.1. \end{proof}

\begin{lemma}[Resource Embedding Composition]\label{bp:markedresourceembeddingcomposition}
\lean{Oberwolfach.markedResourceEmbedding_comp_apply}
\leanok
Composing nested vertex and hub embeddings composes their physical resource maps exactly, including the finite set representing a rim edge.
\end{lemma}
\begin{proof}\leanok The stated conclusion follows from the displayed hypotheses, via the checked Lean proof. This is not a proof of Theorem~1.1. \end{proof}

\begin{definition}[Reference Terminal Representation]\label{bp:markedreferenceterminalrepresentation-definitions}
\lean{Oberwolfach.MarkedReferenceTerminalRepresentation}
\leanok
The last physical remainder is represented by its actual typed arc and spoke sets on the final fixed reference. Their four divisibility equations and containment in the terminal pair are included.
\end{definition}

\begin{lemma}[Reference Finite Iteration]\label{bp:markedreferencefiniteiteration}
\lean{Oberwolfach.marked_reference_finite_iteration}
\leanok
Iteration of the constructed fixed-reference cover-down step. The individual level estimates are explicit inputs; no packing or next-state existence is assumed. The last remainder is retained on the last inner pair without requiring an additional analytic level at the terminal scale.
\end{lemma}
\begin{proof}\uses{bp:markedreferencelevelstepproperty-definitions,bp:markedreferencelevelembedding,bp:markedresourceembeddingcomposition,bp:reindexedresourceiteration}\leanok The stated conclusion follows from the displayed hypotheses, via the checked Lean proof. This is not a proof of Theorem~1.1. \end{proof}

\begin{lemma}[Uniform Reference Finite Iteration]\label{bp:markeduniformreferencefiniteiteration}
\lean{Oberwolfach.marked_uniform_reference_finite_iteration}
\leanok
Apply protected cover-down to the actual finite deletion of a fixed reference. All analytic estimates of the current sample are derived from those of the reference and the displayed global and star allowances.
\end{lemma}
\begin{proof}\uses{bp:markedreferencefiniteiteration,bp:markeduniformreferencelevelstep}\leanok The stated conclusion follows from the displayed hypotheses, via the checked Lean proof. This is not a proof of Theorem~1.1. \end{proof}

\section{Reference states, reindexing and bounded cycle weights}

\begin{definition}[Ambient Terminal Representation]\label{bp:markedambientterminalrepresentation-definitions}
\lean{Oberwolfach.MarkedAmbientTerminalRepresentation}
\leanok
A divisible terminal remainder represented in the original reference host, with its typed endpoints inside the prescribed terminal pair.
\end{definition}

\begin{lemma}[Reference Terminal Transfer]\label{bp:markedreferenceterminaltransfer}
\lean{Oberwolfach.marked_reference_terminal_transfer}
\leanok
The terminal remainder lifts from the last reference to the first reference as actual typed arcs and spokes. Its four balance equations and its physical resource set are preserved.
\end{lemma}
\begin{proof}\uses{bp:markedreferenceterminalrepresentation-definitions,bp:marked-finite-remainder-host-transfer}\leanok The stated conclusion follows from the displayed hypotheses, via the checked Lean proof. This is not a proof of Theorem~1.1. \end{proof}

\begin{lemma}[Reference Reserved Iteration]\label{bp:markedreferencereservediteration}
\lean{Oberwolfach.marked_reference_reserved_iteration}
\leanok
A reservoir is chosen before the initial deletion state. The actual constructed reference-level steps are iterated and their terminal typed remainder is absorbed in the first reference. This theorem still requires the displayed fixed level inputs and the reserved absorber property.
\end{lemma}
\begin{proof}\uses{bp:markedreferencefiniteiteration,bp:markedreferenceterminaltransfer,bp:marked-terminal-larger-exclusion}\leanok The stated conclusion follows from the displayed hypotheses, via the checked Lean proof. This is not a proof of Theorem~1.1. \end{proof}

\begin{lemma}[Uniform Reference Reserved Iteration]\label{bp:markeduniformreferencereservediteration}
\lean{Oberwolfach.marked_uniform_reference_reserved_iteration}
\leanok
Apply protected cover-down to the actual finite deletion of a fixed reference. All analytic estimates of the current sample are derived from those of the reference and the displayed global and star allowances.
\end{lemma}
\begin{proof}\uses{bp:markedreferencereservediteration,bp:markeduniformreferencelevelstep}\leanok The stated conclusion follows from the displayed hypotheses, via the checked Lean proof. This is not a proof of Theorem~1.1. \end{proof}

\begin{lemma}[Typed Host Cover Reference State]\label{bp:markedtypedhostcoverreferencestate}
\lean{Oberwolfach.marked_typed_host_cover_reference_state}
\leanok
The protected cover-down theorem in an order usable for iteration. The analytic input tolerance is chosen before the arbitrary requested star loss. The inner correction is then smaller than that loss; only inner-root imbalance needs the still smaller correction tolerance. All packings are constructed from the displayed graph inputs.
\end{lemma}
\begin{proof}\uses{bp:markedtypedhostcovernextstate,bp:markedreferencedeletionstate}\leanok The stated conclusion follows from the displayed hypotheses, via the checked Lean proof. This is not a proof of Theorem~1.1. \end{proof}

\begin{lemma}[Reference Level Geometry]\label{bp:markedreferencelevelgeometry}
\lean{Oberwolfach.marked_reference_level_in_first_pair,Oberwolfach.marked_reference_terminal_pair_in_first_pair}
\leanok
Nested embeddings place the terminal pair inside the first protected pair. No additional size, sampling or host assumption is used.
\end{lemma}
\begin{proof}\leanok The stated conclusion follows from the displayed hypotheses, via the checked Lean proof. This is not a proof of Theorem~1.1. \end{proof}

\begin{lemma}[Reference Absorber Deletion]\label{bp:markedreferenceabsorberdeletion}
\lean{Oberwolfach.marked_reference_absorber_deletion}
\leanok
Delete a previously constructed wheel-decomposable reservoir from a fixed reference. The state, its exact physical complement, the two incidence bounds and canonical divisibility are derived from the actual wheels.
\end{lemma}
\begin{proof}\uses{bp:marked-boundary-remainder,bp:marked-protected-active-estimates}\leanok The stated conclusion follows from the displayed hypotheses, via the checked Lean proof. This is not a proof of Theorem~1.1. \end{proof}

\begin{lemma}[Reference Reserved Completion]\label{bp:markedreferencereservedcompletion}
\lean{Oberwolfach.marked_reference_reserved_completion}
\leanok
Complete a fixed reference host using the reserved terminal absorber. The initial state is constructed by deleting its actual wheel decomposition; its divisibility and both loss bounds are not independent hypotheses.
\end{lemma}
\begin{proof}\uses{bp:markedreferencereservediteration,bp:markedreferenceabsorberdeletion,bp:markedreferencelevelgeometry,bp:marked-terminal-larger-exclusion}\leanok The stated conclusion follows from the displayed hypotheses, via the checked Lean proof. This is not a proof of Theorem~1.1. \end{proof}

\begin{lemma}[Uniform Reference Reserved Completion]\label{bp:markeduniformreferencereservedcompletion}
\lean{Oberwolfach.marked_uniform_reference_reserved_completion}
\leanok
Apply protected cover-down to the actual finite deletion of a fixed reference. All analytic estimates of the current sample are derived from those of the reference and the displayed global and star allowances.
\end{lemma}
\begin{proof}\uses{bp:markedreferencereservedcompletion,bp:markeduniformreferencelevelstep}\leanok The stated conclusion follows from the displayed hypotheses, via the checked Lean proof. This is not a proof of Theorem~1.1. \end{proof}

\begin{lemma}[Reference Absorber Common Event]\label{bp:markedreferenceabsorbercommonevent}
\lean{Oberwolfach.marked_reference_absorber_common_event,Oberwolfach.marked_reference_absorber_common_event_separate_thresholds}
\leanok
The original retained experiment simultaneously supplies the fixed reference estimates at all vortex levels and the reserved terminal absorber property. The reference conclusion remains uniform over later bounded deletions, which may depend on the entire retained outcome. A separate-threshold version chooses the level-size threshold before the absorber size and the host-order threshold afterwards.
\end{lemma}
\begin{proof}\uses{bp:markedreferenceadaptivevortexevent,bp:marked-terminal-larger-exclusion}\leanok The stated conclusion follows from the displayed hypotheses, via the checked Lean proof. This is not a proof of Theorem~1.1. \end{proof}

\begin{lemma}[Reference Absorber Configuration Event]\label{bp:markedreferenceabsorberconfigurationevent}
\lean{Oberwolfach.marked_reference_absorber_configuration_common_event,Oberwolfach.marked_reference_absorber_configuration_common_event_separate_thresholds}
\leanok
The original retained experiment simultaneously supplies the fixed reference estimates at all vortex levels and the reserved terminal absorber property. The reference conclusion remains uniform over later bounded deletions, which may depend on the entire retained outcome. A separate-threshold version chooses the level-size threshold before the absorber size and the host-order threshold afterwards.
\end{lemma}
\begin{proof}\uses{bp:markedreferenceabsorbercommonevent,bp:markedcommonadaptivevortexevent}\leanok The stated conclusion follows from the displayed hypotheses, via the checked Lean proof. This is not a proof of Theorem~1.1. \end{proof}

\begin{lemma}[Reference Absorber Configuration Balance Event]\label{bp:markedreferenceabsorberconfigurationbalanceevent}
\lean{Oberwolfach.marked_reference_absorber_configuration_balance_event,Oberwolfach.marked_reference_absorber_configuration_balance_event_separate_thresholds}
\leanok
The original retained experiment simultaneously supplies the fixed reference estimates at all vortex levels and the reserved terminal absorber property. The reference conclusion remains uniform over later bounded deletions, which may depend on the entire retained outcome. A separate-threshold version chooses the level-size threshold before the absorber size and the host-order threshold afterwards.
\end{lemma}
\begin{proof}\uses{bp:markedreferenceabsorberconfigurationevent,bp:markedfinitebalanceadaptivevortexevent}\leanok The stated conclusion follows from the displayed hypotheses, via the checked Lean proof. This is not a proof of Theorem~1.1. \end{proof}

\begin{lemma}[Reference Absorber Configuration Balance Neighbourhood Event]\label{bp:markedreferenceabsorberconfigurationbalanceneighbourhoodevent}
\lean{Oberwolfach.marked_reference_absorber_configuration_balance_neighbourhood_event,Oberwolfach.marked_reference_absorber_configuration_balance_neighbourhood_event_separate_thresholds}
\leanok
On one retained outcome, the reference and common-configuration estimates, fixed-domain balance and row and index neighbourhood estimates hold at all specified levels, together with the reserved terminal absorber property. The estimates remain valid for all later deletions satisfying the displayed bounds, including deletions chosen after observing the retained outcome.
\end{lemma}
\begin{proof}\uses{bp:markedreferenceabsorberconfigurationbalanceevent,bp:markedneighbourhoodadaptivevortexevent}\leanok The stated conclusion follows from the displayed hypotheses, via the checked Lean proof. This is not a proof of Theorem~1.1. \end{proof}

\begin{lemma}[Reference Absorber Configuration Balance Neighbourhood Load Event]\label{bp:markedreferenceabsorberconfigurationbalanceneighbourhoodloadevent}
\lean{Oberwolfach.marked_reference_absorber_configuration_balance_neighbourhood_load_event,Oberwolfach.marked_reference_absorber_configuration_balance_neighbourhood_load_event_separate_thresholds}
\leanok
On one retained outcome, the reference, configuration, fixed-domain balance, neighbourhood and weighted-load estimates hold at all specified levels, together with the reserved terminal absorber property. The weights and root domains are fixed before sampling, and the conclusions hold for every later deletion satisfying the displayed bounds.
\end{lemma}
\begin{proof}\uses{bp:markedreferenceabsorberconfigurationbalanceneighbourhoodevent,bp:markedpackingadaptivevortexevent}\leanok The stated conclusion follows from the displayed hypotheses, via the checked Lean proof. This is not a proof of Theorem~1.1. \end{proof}

\section{Local configurations and full-hub potential loads}

\begin{lemma}[Reference Absorber Configuration Balance Neighbourhood Load Local Event]\label{bp:markedreferenceabsorberconfigurationbalanceneighbourhoodloadlocalevent}
\lean{Oberwolfach.marked_reference_absorber_configuration_balance_neighbourhood_load_local_event,Oberwolfach.marked_reference_absorber_configuration_balance_neighbourhood_load_local_event_separate_thresholds}
\leanok
The common retained outcome also supplies local spoke-root potential loads on every fixed vortex level. The supplied weights and normalizations are fixed before sampling.
\end{lemma}
\begin{proof}\uses{bp:markedreferenceabsorberconfigurationbalanceneighbourhoodloadevent,bp:marked-retained-load-realization}\leanok The stated conclusion follows from the displayed hypotheses, via the checked Lean proof. This is not a proof of Theorem~1.1. \end{proof}

\begin{lemma}[Reference Absorber Configuration Balance Neighbourhood Load Local Config Event]\label{bp:markedreferenceabsorberconfigurationbalanceneighbourhoodloadlocalconfigevent}
\lean{Oberwolfach.marked_reference_absorber_configuration_balance_neighbourhood_load_local_config_event,Oberwolfach.marked_reference_absorber_configuration_balance_neighbourhood_load_local_config_event_separate_thresholds}
\leanok
The same retained outcome supplies local configuration counts and local spoke-root loads at every specified level, in addition to the reference estimates, balance, neighbourhood, packing loads, and reserved absorber property. The displayed parent and domain hypotheses remain explicit.
\end{lemma}
\begin{proof}\uses{bp:markedreferenceabsorberconfigurationbalanceneighbourhoodloadlocalevent,bp:marked-retained-configurations-vortex}\leanok The stated conclusion follows from the displayed hypotheses, via the checked Lean proof. This is not a proof of Theorem~1.1. \end{proof}

\begin{lemma}[Packing Full Hub Loads]\label{bp:markedpackingfullhubloads}
\lean{Oberwolfach.marked_packing_loads_on_all_hubs}
\leanok
For the identity hub embedding, the packing-domain load estimate contains the potential-load estimate for every root in the whole host.
\end{lemma}
\begin{proof}\uses{bp:marked-retained-load-realization}\leanok The stated conclusion follows from the displayed hypotheses, via the checked Lean proof. This is not a proof of Theorem~1.1. \end{proof}

\section{Protected absorber inputs after physical deletion}

\begin{lemma}[Common Potential Physical Deletion]\label{bp:commonpotentialphysicaldeletion}
\lean{Oberwolfach.common_potential_loads_after_physical_deletion}
\leanok
A fixed common-hub selector and fixed parent coefficients survive later bounded physical rim and spoke deletions. The load at every potential root is compared on the original vertex domain, including deleted roots.
\end{lemma}
\begin{proof}\uses{bp:marked-retained-balance-load-stability}\leanok The stated conclusion follows from the displayed hypotheses, via the checked Lean proof. This is not a proof of Theorem~1.1. \end{proof}

\begin{lemma}[Common Witness Physical Deletion]\label{bp:commonwitnessphysicaldeletion}
\lean{Oberwolfach.common_witness_after_physical_deletion}
\leanok
A fixed common-hub witness survives bounded physical deletions. The selector and parent coefficients are not resampled, the original domain and normalization are retained, and the loss is charged directly to the realized deleted degrees.
\end{lemma}
\begin{proof}\uses{bp:commonpotentialphysicaldeletion,bp:active-common-configurations,bp:common-joint-witness}\leanok The stated conclusion follows from the displayed hypotheses, via the checked Lean proof. This is not a proof of Theorem~1.1. \end{proof}

\begin{lemma}[Common Fixed Witness Profiles]\label{bp:markedcommonfixedwitnessprofiles}
\lean{Oberwolfach.marked_common_fixed_witness_profiles}
\leanok
Construct the fixed pair profiles from a single common-witness event. The selector for each pair of hubs is supplied and held fixed. The coefficients are the supplied parent coefficients, with their original normalization.
\end{lemma}
\begin{proof}\uses{bp:common-joint-witness}\leanok The stated conclusion follows from the displayed hypotheses, via the checked Lean proof. This is not a proof of Theorem~1.1. \end{proof}

\begin{lemma}[Common Profiles Physical Deletion]\label{bp:commonprofilesphysicaldeletion}
\lean{Oberwolfach.common_pair_profiles_after_physical_deletion}
\leanok
Construct common-pair absorber profiles after an arbitrary bounded physical deletion from one fixed family of original witnesses. The selector for each pair and every cycle coefficient are inherited from the original outcome; the only losses are the displayed deterministic degree costs.
\end{lemma}
\begin{proof}\uses{bp:commonwitnessphysicaldeletion,bp:markedcommonfixedwitnessprofiles}\leanok The stated conclusion follows from the displayed hypotheses, via the checked Lean proof. This is not a proof of Theorem~1.1. \end{proof}

\begin{lemma}[Common Profiles Adaptive Deletion]\label{bp:commonprofilesadaptivedeletion}
\lean{Oberwolfach.common_pair_profiles_for_all_physical_deletions}
\leanok
Construct common-pair absorber profiles after an arbitrary bounded physical deletion from one fixed family of original witnesses. The selector for each pair and every cycle coefficient are inherited from the original outcome; the only losses are the displayed deterministic degree costs.
\end{lemma}
\begin{proof}\uses{bp:commonprofilesphysicaldeletion}\leanok The stated conclusion follows from the displayed hypotheses, via the checked Lean proof. This is not a proof of Theorem~1.1. \end{proof}

\begin{definition}[Common Physical Profile Margins]\label{bp:commonphysicalprofilemargins-definitions}
\lean{Oberwolfach.MarkedCommonPairCycleProfiles.lowerCount}
\leanok
Choose one deletion fraction before the host and its later deletion. The loss leaves positive common-cycle/path coefficients and the requested profile error.
\end{definition}

\begin{lemma}[Common Physical Profile Margins]\label{bp:commonphysicalprofilemargins}
\lean{Oberwolfach.common_physical_profile_margins}
\leanok
Choose one deletion fraction before the host and its later deletion. The loss leaves positive common-cycle/path coefficients and the requested profile error.
\end{lemma}
\begin{proof}\leanok The stated conclusion follows from the displayed hypotheses, via the checked Lean proof. This is not a proof of Theorem~1.1. \end{proof}

\begin{lemma}[Common Profiles Uniform Deletion]\label{bp:commonprofilesuniformdeletion}
\lean{Oberwolfach.common_pair_profiles_with_preselected_deletion_margin}
\leanok
Construct common-pair absorber profiles after an arbitrary bounded physical deletion from one fixed family of original witnesses. The selector for each pair and every cycle coefficient are inherited from the original outcome; the only losses are the displayed deterministic degree costs.
\end{lemma}
\begin{proof}\uses{bp:commonprofilesadaptivedeletion,bp:commonphysicalprofilemargins}\leanok The stated conclusion follows from the displayed hypotheses, via the checked Lean proof. This is not a proof of Theorem~1.1. \end{proof}

\begin{lemma}[Host Counts Physical Deletion]\label{bp:markedhostcountsphysicaldeletion}
\lean{Oberwolfach.marked_host_extension_counts_arc_congr,Oberwolfach.marked_sample_host_extensions_after_physical_deletion}
\leanok
Extension counts in the actual physically deleted sample. The typed arc occurrences are identified by their exact deletion equivalences, and spoke losses are measured by the realized physical degrees.
\end{lemma}
\begin{proof}\uses{bp:marked-sampled-deletion-occurrences,bp:marked-local-extensions}\leanok The stated conclusion follows from the displayed hypotheses, via the checked Lean proof. This is not a proof of Theorem~1.1. \end{proof}

\begin{lemma}[Host Counts Physical Deletion Real]\label{bp:markedhostcountsphysicaldeletionreal}
\lean{Oberwolfach.marked_sample_host_extensions_after_physical_deletion_real}
\leanok
Degree bounds may be real. Taking the floor introduces no additive loss because the deleted degrees are integers.
\end{lemma}
\begin{proof}\uses{bp:markedhostcountsphysicaldeletion}\leanok The stated conclusion follows from the displayed hypotheses, via the checked Lean proof. This is not a proof of Theorem~1.1. \end{proof}

\begin{lemma}[Transformer Factor Threshold Gap]\label{bp:transformerfactorthresholdgap}
\lean{Oberwolfach.odd_transformer_factor_threshold_gap}
\leanok
The odd-cycle decomposition threshold can be strictly above the minimum-degree coefficient used by the transformer factor theorem.
\end{lemma}
\begin{proof}\uses{bp:threshold-odd}\leanok The stated conclusion follows from the displayed hypotheses, via the checked Lean proof. This is not a proof of Theorem~1.1. \end{proof}

\section{Adaptive absorber and protected completion}

\begin{lemma}[Finite Absorbers After Physical Deletion]\label{bp:markedfiniteabsorbersafterphysicaldeletion}
\lean{Oberwolfach.marked_finite_absorbers_after_physical_deletion}
\leanok
Construct finite absorbers in every later bounded physical deletion of one fixed sampled host. The coefficients, domains and common-pair selectors are inherited from the original host. The deletion tolerance is chosen before the host, its sample and every later deletion.
\end{lemma}
\begin{proof}\uses{bp:commonprofilesuniformdeletion,bp:markedhostcountsphysicaldeletionreal,bp:marked-finite-remainder-absorber}\leanok The stated conclusion follows from the displayed hypotheses, via the checked Lean proof. This is not a proof of Theorem~1.1. \end{proof}

\begin{lemma}[Reserved Absorbers After Physical Deletion]\label{bp:markedreservedabsorbersafterphysicaldeletion}
\lean{Oberwolfach.marked_reserved_absorbers_after_physical_deletion}
\leanok
Construct finite absorbers in every later bounded physical deletion of one fixed sampled host. The coefficients, domains and common-pair selectors are inherited from the original host. The deletion tolerance is chosen before the host, its sample and every later deletion.
\end{lemma}
\begin{proof}\uses{bp:markedfiniteabsorbersafterphysicaldeletion,bp:marked-terminal-larger-exclusion}\leanok The stated conclusion follows from the displayed hypotheses, via the checked Lean proof. This is not a proof of Theorem~1.1. \end{proof}

\begin{lemma}[Sampled Reserved Absorbers After Physical Deletion]\label{bp:markedsampledreservedabsorbersafterphysicaldeletion}
\lean{Oberwolfach.marked_sampled_reserved_absorbers_after_physical_deletion}
\leanok
The original marked retention event supplies a fixed terminal absorber for every later permitted terminal host and bounded forbidden set. The finite family of possible remainders is counted before any of its absorbers are chosen. No additional random sample or independence claim is used.
\end{lemma}
\begin{proof}\uses{bp:markedreservedabsorbersafterphysicaldeletion,bp:common-joint-witness,bp:marked-local-extensions}\leanok The stated conclusion follows from the displayed hypotheses, via the checked Lean proof. This is not a proof of Theorem~1.1. \end{proof}

\begin{lemma}[Reference Adaptive Absorber Joint Event]\label{bp:markedreferenceadaptiveabsorberjointevent}
\lean{Oberwolfach.marked_reference_adaptive_absorber_joint_event,Oberwolfach.marked_reference_adaptive_absorber_joint_event_separate_thresholds}
\leanok
The same retained outcome supplies local configuration counts and local spoke-root loads at every specified level, in addition to the reference estimates, balance, neighbourhood, packing loads, and reserved absorber property. The displayed parent and domain hypotheses remain explicit.
\end{lemma}
\begin{proof}\uses{bp:markedsampledreservedabsorbersafterphysicaldeletion,bp:markedreferenceabsorberconfigurationbalanceneighbourhoodloadlocalconfigevent}\leanok The stated conclusion follows from the displayed hypotheses, via the checked Lean proof. This is not a proof of Theorem~1.1. \end{proof}

\begin{lemma}[Reference Completion After Physical Deletion]\label{bp:markedreferencecompletionafterphysicaldeletion}
\lean{Oberwolfach.marked_reference_completion_after_physical_deletion}
\leanok
Complete a fixed reference host using the reserved terminal absorber. The initial state is constructed by deleting its actual wheel decomposition; its divisibility and both loss bounds are not independent hypotheses.
\end{lemma}
\begin{proof}\uses{bp:markedreferencereservedcompletion}\leanok The stated conclusion follows from the displayed hypotheses, via the checked Lean proof. This is not a proof of Theorem~1.1. \end{proof}

\begin{lemma}[Uniform Reference Completion After Physical Deletion]\label{bp:markeduniformreferencecompletionafterphysicaldeletion}
\lean{Oberwolfach.marked_uniform_reference_completion_after_physical_deletion}
\leanok
Apply protected cover-down to the actual finite deletion of a fixed reference. All analytic estimates of the current sample are derived from those of the reference and the displayed global and star allowances.
\end{lemma}
\begin{proof}\uses{bp:markedreferencecompletionafterphysicaldeletion,bp:markeduniformreferencelevelstep}\leanok The stated conclusion follows from the displayed hypotheses, via the checked Lean proof. This is not a proof of Theorem~1.1. \end{proof}

\section{Ordinary cycle profiles from neighbourhoods}

\begin{definition}[Ordinary Common Profiles From Geometry]\label{bp:ordinarycommonprofilesfromgeometry-definitions}\uses{bp:ordinary-local-absorbers}
\lean{Oberwolfach.ordinary_common_profiles_from_geometry}
\leanok
Construct the pair profiles with the explicit paired absorber template. The absorber counts follow uniformly from the actual common-neighbourhood bounds; the supplied weights and their original domains are unchanged.
\end{definition}

\section{Restricted hosts and ordinary pair profiles}

\begin{lemma}[Restricted Deletion Budget]\label{bp:markedrestricteddeletionbudget}
\lean{Oberwolfach.marked_restricted_deletion_budget}
\leanok
The same ambient physical deletion restricts to each fixed level, with its degree budget rescaled by the ratio of the ambient and level orders.
\end{lemma}
\begin{proof}\leanok The stated conclusion follows from the displayed hypotheses, via the checked Lean proof. This is not a proof of Theorem~1.1. \end{proof}

\begin{lemma}[Encoded Inner Deletion]\label{bp:markedencodedinnerdeletion}
\lean{Oberwolfach.marked_encoded_inner_deletion}
\leanok
The encoded host agrees with the original sampled coordinates on a fixed protected inner domain, even after a physical deletion. The induced deletion uses the actual pullback graph, not a new sampling law.
\end{lemma}
\begin{proof}\uses{bp:markedhostrawencoding}\leanok The stated conclusion follows from the displayed hypotheses, via the checked Lean proof. This is not a proof of Theorem~1.1. \end{proof}

\begin{lemma}[Retained Agreement Event]\label{bp:markedretainedagreementevent}
\lean{Oberwolfach.marked_retained_full_neighbourhood_with_agreement_event}
\leanok
The same retained outcome satisfies all row and index inequalities. Retention losses are bounded by realised changed spokes. Product targets and the original outer-domain error scales are retained.
\end{lemma}
\begin{proof}\uses{bp:markedreferenceinputstability,bp:marked-spoke-change-estimates}\leanok The stated conclusion follows from the displayed hypotheses, via the checked Lean proof. This is not a proof of Theorem~1.1. \end{proof}

\begin{definition}[Ordinary Pair Common Geometry]\label{bp:ordinarypaircommongeometry-definitions}\uses{bp:ordinarycommonprofilesfromgeometry-definitions,bp:ordinarypaircommongeometry}
\lean{Oberwolfach.ordinary_common_profiles_from_pair_geometry}
\leanok
The two-root common-neighbourhood input needed by the paired-absorber count follows from the concrete two-neighbourhood hypotheses of the ordinary relative-host construction, including sets of zero or one root.
\end{definition}

\begin{lemma}[Ordinary Pair Common Geometry]\label{bp:ordinarypaircommongeometry}
\lean{Oberwolfach.ordinary_pair_common_geometry}
\leanok
The two-root common-neighbourhood input needed by the paired-absorber count follows from the concrete two-neighbourhood hypotheses of the ordinary relative-host construction, including sets of zero or one root.
\end{lemma}
\begin{proof}\leanok The stated conclusion follows from the displayed hypotheses, via the checked Lean proof. This is not a proof of Theorem~1.1. \end{proof}

\begin{definition}[Ordinary Triangle Common Geometry]\label{bp:ordinarytrianglecommongeometry-definitions}\uses{bp:ordinarycommonprofilesfromgeometry-definitions,bp:ordinarytrianglecommongeometry}
\lean{Oberwolfach.ordinary_triangle_profiles_from_four_neighbourhoods}
\leanok
Four-root common-neighbourhood estimates also supply the zero- through three-root estimates required by the triangular paired-absorber count.
\end{definition}

\begin{lemma}[Ordinary Triangle Common Geometry]\label{bp:ordinarytrianglecommongeometry}
\lean{Oberwolfach.ordinary_triangle_common_geometry}
\leanok
Four-root common-neighbourhood estimates also supply the zero- through three-root estimates required by the triangular paired-absorber count.
\end{lemma}
\begin{proof}\leanok The stated conclusion follows from the displayed hypotheses, via the checked Lean proof. This is not a proof of Theorem~1.1. \end{proof}

\section{Ordinary absorber callers and marked retained agreement}

\begin{lemma}[Ordinary Profile Budgets]\label{bp:ordinaryprofilebudgets}
\lean{Oberwolfach.ordinary_profile_budgets}
\leanok
The order threshold for the fixed paired-absorber count is chosen after the cycle length, neighbourhood density and conflict-degree bound.
\end{lemma}
\begin{proof}\uses{bp:ordinary-local-absorbers}\leanok The stated conclusion follows from the displayed hypotheses, via the checked Lean proof. This is not a proof of Theorem~1.1. \end{proof}

\begin{lemma}[Ordinary Pair Profiles For Large Order]\label{bp:ordinarypairprofilesforlargeorder}
\lean{Oberwolfach.ordinary_pair_profiles_for_large_order}
\leanok
The fixed order threshold supplies the numerical prerequisites of the ordinary pair-profile construction. The graph and the fixed weights are introduced only after that threshold; no profile existence is assumed.
\end{lemma}
\begin{proof}\uses{bp:ordinarypaircommongeometry-definitions,bp:ordinaryprofilebudgets}\leanok The stated conclusion follows from the displayed hypotheses, via the checked Lean proof. This is not a proof of Theorem~1.1. \end{proof}

\begin{lemma}[Ordinary Triangle Profiles For Large Order]\label{bp:ordinarytriangleprofilesforlargeorder}
\lean{Oberwolfach.ordinary_triangle_profiles_for_large_order}
\leanok
In the triangular case the same fixed order threshold supplies the numerical paired-absorber prerequisites from four-neighbourhood bounds.
\end{lemma}
\begin{proof}\uses{bp:ordinarytrianglecommongeometry-definitions,bp:ordinaryprofilebudgets}\leanok The stated conclusion follows from the displayed hypotheses, via the checked Lean proof. This is not a proof of Theorem~1.1. \end{proof}

\begin{lemma}[Ordinary Triangle Ambient Absorber From Weights]\label{bp:ordinarytriangleambientabsorberfromweights}
\lean{Oberwolfach.ordinary_triangle_ambient_absorber_from_weights}
\leanok
The triangular relative ambient absorber from concrete four-neighbourhood and fixed cycle-weight hypotheses. The paired-absorber profile is derived.
\end{lemma}
\begin{proof}\uses{bp:ordinary-ambient-absorber,bp:ordinarytriangleprofilesforlargeorder}\leanok The stated conclusion follows from the displayed hypotheses, via the checked Lean proof. This is not a proof of Theorem~1.1. \end{proof}

\begin{lemma}[Ordinary Long Cycle Ambient Absorber From Weights]\label{bp:ordinarylongcycleambientabsorberfromweights}
\lean{Oberwolfach.ordinary_long_cycle_ambient_absorber_from_weights}
\leanok
The longer-cycle relative ambient absorber from concrete two-neighbourhood and fixed cycle-weight hypotheses. The paired-absorber profile is derived.
\end{lemma}
\begin{proof}\uses{bp:ordinary-ambient-absorber,bp:ordinarypairprofilesforlargeorder}\leanok The stated conclusion follows from the displayed hypotheses, via the checked Lean proof. This is not a proof of Theorem~1.1. \end{proof}

\begin{lemma}[Reference Adaptive Absorber Agreement Joint Event]\label{bp:markedreferenceadaptiveabsorberagreementjointevent}
\lean{Oberwolfach.marked_reference_adaptive_absorber_agreement_joint_event,Oberwolfach.marked_reference_adaptive_absorber_agreement_joint_event_separate_thresholds}
\leanok
The same retained outcome supplies local configuration counts and local spoke-root loads at every specified level, in addition to the reference estimates, balance, neighbourhood, packing loads, and reserved absorber property. The displayed parent and domain hypotheses remain explicit.
\end{lemma}
\begin{proof}\uses{bp:markedreferenceadaptiveabsorberjointevent,bp:markedretainedagreementevent}\leanok The stated conclusion follows from the displayed hypotheses, via the checked Lean proof. This is not a proof of Theorem~1.1. \end{proof}

\begin{lemma}[Retained Fixed Domain Agreement Event]\label{bp:markedretainedfixeddomainagreementevent}
\lean{Oberwolfach.marked_retained_fixed_domain_agreement_event}
\leanok
On one original retained sample, raw and retained spoke relations agree up to the prescribed row and column errors on every fixed domain, including at roots outside that domain.
\end{lemma}
\begin{proof}\uses{bp:marked-spoke-change-estimates}\leanok The stated conclusion follows from the displayed hypotheses, via the checked Lean proof. This is not a proof of Theorem~1.1. \end{proof}

\begin{lemma}[Retained Fixed Domain Agreement With Event]\label{bp:markedretainedfixeddomainagreementwithevent}
\lean{Oberwolfach.marked_retained_fixed_domain_agreement_with_event}
\leanok
On one original retained sample, raw and retained spoke relations agree up to the prescribed row and column errors on every fixed domain, including at roots outside that domain.
\end{lemma}
\begin{proof}\uses{bp:markedretainedfixeddomainagreementevent}\leanok The stated conclusion follows from the displayed hypotheses, via the checked Lean proof. This is not a proof of Theorem~1.1. \end{proof}

\begin{lemma}[Reference Domain Agreement Deletion]\label{bp:markedreferencedomainagreementdeletion}
\lean{Oberwolfach.marked_spoke_reference_change_after_deletion_on_fixed_domains}
\leanok
Agreement on each original fixed domain survives physical deletion; all roots remain ambient and the realised deletion loss is charged once.
\end{lemma}
\begin{proof}\uses{bp:markedreferenceinputstability,bp:marked-spoke-change-estimates}\leanok The stated conclusion follows from the displayed hypotheses, via the checked Lean proof. This is not a proof of Theorem~1.1. \end{proof}

\section{Physical-rim completion and exact encoding criteria}

\begin{lemma}[Retained Fixed Domain Agreement After Deletion]\label{bp:markedretainedfixeddomainagreementafterdeletion}
\lean{Oberwolfach.marked_retained_fixed_domain_agreement_after_deletion_with_event}
\leanok
On one original retained sample, raw and retained spoke relations agree up to the prescribed row and column errors on every fixed domain, including at roots outside that domain.
\end{lemma}
\begin{proof}\uses{bp:markedretainedfixeddomainagreementwithevent,bp:markedreferencedomainagreementdeletion}\leanok The stated conclusion follows from the displayed hypotheses, via the checked Lean proof. This is not a proof of Theorem~1.1. \end{proof}

\begin{lemma}[Encoded Host Deletion Obstruction]\label{bp:markedencodedhostdeletionobstruction}
\lean{Oberwolfach.encoded_host_not_rim_deletion_of_absent_edge}
\leanok
An encoded typed host can add a physical rim coordinate outside the protected set. A raw deletion can never create a previously absent rim.
\end{lemma}
\begin{proof}\uses{bp:markedhostrawencoding}\leanok The stated conclusion follows from the displayed hypotheses, via the checked Lean proof. This is not a proof of Theorem~1.1. \end{proof}

\begin{lemma}[Encoded Host Deletion Criterion]\label{bp:markedencodedhostdeletioncriterion}
\lean{Oberwolfach.encoded_host_rim_deletion_iff}
\leanok
An exact coordinate criterion for representing an encoded typed host as a raw rim deletion. This is a finite-set deletion criterion; degree and looplessness requirements for a physical deletion graph remain separate.
\end{lemma}
\begin{proof}\uses{bp:markedhostrawencoding}\leanok The stated conclusion follows from the displayed hypotheses, via the checked Lean proof. This is not a proof of Theorem~1.1. \end{proof}

\begin{lemma}[Reserved Property Physical Congruence]\label{bp:markedreservedpropertyphysicalcongruence}
\lean{Oberwolfach.marked_reserved_property_physical_congruence}
\leanok
The reserved-absorber property depends on the actual typed rim relations, not on unused coordinates of the raw outcome.
\end{lemma}
\begin{proof}\uses{bp:marked-terminal-larger-exclusion}\leanok The stated conclusion follows from the displayed hypotheses, via the checked Lean proof. This is not a proof of Theorem~1.1. \end{proof}

\begin{lemma}[Reference Completion Physical Rims]\label{bp:markedreferencecompletionphysicalrims}
\lean{Oberwolfach.marked_reference_completion_after_physical_deletion_physical_rims}
\leanok
Complete a fixed reference host using the reserved terminal absorber. The initial state is constructed by deleting its actual wheel decomposition; its divisibility and both loss bounds are not independent hypotheses.
\end{lemma}
\begin{proof}\uses{bp:markedreferencereservedcompletion,bp:markedreservedpropertyphysicalcongruence}\leanok The stated conclusion follows from the displayed hypotheses, via the checked Lean proof. This is not a proof of Theorem~1.1. \end{proof}

\begin{lemma}[Uniform Reference Completion Physical Rims]\label{bp:markeduniformreferencecompletionphysicalrims}
\lean{Oberwolfach.marked_uniform_reference_completion_after_physical_deletion_physical_rims}
\leanok
Apply protected cover-down to the actual finite deletion of a fixed reference. All analytic estimates of the current sample are derived from those of the reference and the displayed global and star allowances.
\end{lemma}
\begin{proof}\uses{bp:markedreferencecompletionphysicalrims,bp:markeduniformreferencelevelstep}\leanok The stated conclusion follows from the displayed hypotheses, via the checked Lean proof. This is not a proof of Theorem~1.1. \end{proof}

\section{Typed subhosts and protected completion}

\begin{lemma}[Spoke Deletion Representation]\label{bp:markedspokedeletionrepresentation}
\lean{Oberwolfach.marked_spoke_deletion_representation}
\leanok
A typed spoke subhost of a type-disjoint original host is obtained by one physical spoke deletion. Its row and column degrees count the deleted physical pairs exactly.
\end{lemma}
\begin{proof}\leanok The stated conclusion follows from the displayed hypotheses, via the checked Lean proof. This is not a proof of Theorem~1.1. \end{proof}

\begin{definition}[Typed Rim Subhost Deletion]\label{bp:markedtypedrimsubhostdeletion-definitions}
\lean{Oberwolfach.markedRemovedPhysicalRim}
\leanok
A typed physical rim subhost is obtained by deleting exactly its missing physical edges. No equality of unused raw coordinates is needed.
\end{definition}

\begin{lemma}[Typed Rim Subhost Deletion]\label{bp:markedtypedrimsubhostdeletion}
\lean{Oberwolfach.marked_typed_rim_subhost_deletion}
\leanok
A typed physical rim subhost is obtained by deleting exactly its missing physical edges. No equality of unused raw coordinates is needed.
\end{lemma}
\begin{proof}\uses{bp:marked-sample-physical-map}\leanok The stated conclusion follows from the displayed hypotheses, via the checked Lean proof. This is not a proof of Theorem~1.1. \end{proof}

\begin{lemma}[Uniform Reference Completion Typed Subhost]\label{bp:markeduniformreferencecompletiontypedsubhost}
\lean{Oberwolfach.marked_uniform_reference_completion_from_typed_subhost}
\leanok
Apply protected cover-down to the actual finite deletion of a fixed reference. All analytic estimates of the current sample are derived from those of the reference and the displayed global and star allowances.
\end{lemma}
\begin{proof}\uses{bp:markedspokedeletionrepresentation,bp:markedtypedrimsubhostdeletion,bp:markedreferencecompletionphysicalrims,bp:markeduniformreferencelevelstep}\leanok The stated conclusion follows from the displayed hypotheses, via the checked Lean proof. This is not a proof of Theorem~1.1. \end{proof}

\section{A fixed absorber in a changed host}

\begin{lemma}[Encoded-host typed containment]\label{bp:encoded-typed-containment}
\lean{Oberwolfach.marked_encoded_host_typed_containment_iff,Oberwolfach.marked_encoded_host_typed_containment_iff_outside}
\leanok
For an encoded typed host agreeing with the original sample on the protected inner pair, containment of all its typed rim arcs in the original sample is equivalent to checking the supplied arcs outside that pair. This is a criterion, not a proof that the actual preparation host is contained in the sample.
\end{lemma}
\begin{proof}\uses{bp:markedhostrawencoding}\leanok The encoding identifies precisely the supplied ordinary and marked arcs. On the inner pair the stated agreement supplies their membership in the sample.\end{proof}

\begin{lemma}[Choosing the terminal reservoir before the host]\label{bp:choose-fixed-reservoir}
\lean{Oberwolfach.marked_choose_reserved_terminal_reservoir}
\leanok
For a fixed terminal pair and a larger excluded pair, the reserved-absorber property supplies one bounded reservoir that works for every divisible typed remainder on the terminal pair. The reservoir is chosen before any later host.
\end{lemma}
\begin{proof}\uses{bp:marked-terminal-larger-exclusion}\leanok Apply the reserved-absorber property to the complete family of typed arcs and spokes on the fixed terminal pair.\end{proof}

\begin{lemma}[Reserved absorber in a changed host]\label{bp:reserved-changed-host}
\lean{Oberwolfach.marked_reserved_terminal_absorber_in_ambient_host,Oberwolfach.marked_reserved_terminal_completion_in_modified_subhost,Oberwolfach.marked_reserved_absorber_deletion_in_changed_host}
\leanok
A reservoir selected in a source host may be transferred to a later host that preserves its typed resources, even if the later host has additional edges. The reservoir avoids the larger protected pair. Deleting its wheel decomposition before cover-down preserves the inner induced resources, the displayed incidence budgets and divisibility.
\end{lemma}
\begin{proof}\uses{bp:marked-terminal-larger-exclusion,bp:marked-terminal-completion-between-hosts,bp:marked-local-host-transfer,bp:markedreferenceabsorberdeletion}\leanok Transfer the actual typed resources, then apply the physical deletion and remainder identities.\end{proof}

\begin{lemma}[Iteration with a fixed transferred absorber]\label{bp:fixed-transferred-iteration}
\lean{Oberwolfach.marked_reference_iteration_with_fixed_transferred_absorber}
\leanok
Given the stated reference-level data and step properties, a fixed reservoir chosen in another typed host absorbs the final remainder of the finite iteration whenever its typed resources are preserved and the final typed remainder agrees with the source host.
\end{lemma}
\begin{proof}\uses{bp:markedreferencefiniteiteration,bp:markedreferenceterminaltransfer,bp:marked-finite-remainder-host-transfer}\leanok Iterate the physical deletions, transfer the final typed remainder and combine the disjoint decompositions.\end{proof}

\begin{lemma}[Completion with a fixed transferred absorber]\label{bp:fixed-transferred-completion}
\lean{Oberwolfach.marked_reference_completion_with_fixed_transferred_absorber}
\leanok
A changed divisible host has a marked wheel decomposition under the displayed fixed-level, typed-agreement and deletion-budget hypotheses. The fixed source-host reservoir is removed before the first cover-down and is used after the finite iteration. This statement does not construct those level inputs or prove the main theorem.
\end{lemma}
\begin{proof}\uses{bp:marked-local-host-transfer,bp:marked-terminal-larger-exclusion,bp:markedreferenceabsorberdeletion,bp:fixed-transferred-iteration}\leanok Transfer and remove the reservoir, run the fixed-reservoir iteration and reunite its resources with the remainder.\end{proof}

\begin{lemma}[Terminal typed agreement on the protected pair]\label{bp:encoded-terminal-agreement}
\lean{Oberwolfach.marked_encoded_host_terminal_arc_agreement,Oberwolfach.marked_terminal_spoke_agreement}
\leanok
Every terminal ordinary or marked arc of the encoded host on the protected pair has an arc of the same endpoints and type in the original sample. Agreement of the two spoke types on the protected pair gives the corresponding terminal spoke agreement.
\end{lemma}
\begin{proof}\uses{bp:markedhostrawencoding}\leanok Use the prescribed inner agreement, and restrict it to the fixed terminal pair.\end{proof}

\section{Encoded completion from the original reservoir}

\begin{definition}[Encoded-host arc bijections]\label{bp:encoded-host-arc-bijections}
\lean{Oberwolfach.markedRawHostOrdinaryArcEquiv,Oberwolfach.markedRawHostMarkedArcEquiv}
\leanok
The ordinary and marked occurrences of an explicit typed host are in bijection with the corresponding occurrences of its encoded host, preserving both directed endpoints.
\end{definition}

\begin{lemma}[Reserved rim transfer to an encoded host]\label{bp:encoded-reserved-rim-transfer}
\lean{Oberwolfach.marked_encoded_host_preserves_reserved_rims}
\leanok
If a later explicit host contains every ordinary and marked rim occurrence of the reserved absorber with the same directed endpoints and type, these occurrences also appear in its encoded representation.
\end{lemma}
\begin{proof}\uses{bp:encoded-host-arc-bijections}\leanok Apply the endpoint-preserving occurrence bijections of the encoded host.\end{proof}

\begin{lemma}[Terminal agreement along the reference chain]\label{bp:reference-encoded-terminal-agreement}
\lean{Oberwolfach.marked_reference_encoded_terminal_arc_agreement,Oberwolfach.marked_reference_terminal_spoke_agreement}
\leanok
The terminal pair of the nested reference chain lies inside the first protected pair. Agreement of the later encoded host with the original sample there gives the terminal arc and spoke comparisons required for transfer back to the original sample.
\end{lemma}
\begin{proof}\uses{bp:encoded-terminal-agreement,bp:markedreferencelevelgeometry}\leanok Restrict the first-pair agreement to the terminal pair.\end{proof}

\begin{lemma}[Encoded completion with a fixed original reservoir]\label{bp:encoded-fixed-reservoir-completion}
\lean{Oberwolfach.marked_encoded_fixed_reservoir_completion}
\leanok
Given the stated level chain, step properties, first-pair agreement, typed retention of a fixed original-sample reservoir, divisibility and incidence budgets, the encoded later host has a marked wheel decomposition. Its changes outside the protected pair may include additions. The terminal arc and spoke comparison follows from the first-pair agreement; the reservoir's decomposition and resource support follow from its absorption property applied to the empty remainder.
\end{lemma}
\begin{proof}\uses{bp:encoded-reserved-rim-transfer,bp:reference-encoded-terminal-agreement,bp:fixed-transferred-completion}\leanok Transfer the reserved resources to the encoded host, use the nested-pair agreement and apply the fixed-reservoir completion.\end{proof}

\section{Two-hub weights on the original sample}

\begin{lemma}[Two-hub coordinate locality]\label{bp:ordinary-two-hub-coordinate-locality}
\lean{Oberwolfach.paired_coordinate_support_bounds,Oberwolfach.ordinary_pair_hub_coordinates_bounds,Oberwolfach.ordinary_pair_hub_resource_supported}
\leanok
A resource in the intersection of two fixed hubs' sampled spoke neighbourhoods reads the union of the two original coordinate sets. The read and reverse-incidence bounds are at most twice the original bounds.
\end{lemma}
\begin{proof}\leanok Take the union of the original supports and use the injections that insert each hub.\end{proof}

\begin{lemma}[Simultaneous two-hub load concentration]\label{bp:ordinary-two-hub-concentration}
\lean{Oberwolfach.ordinary_pair_hub_load_common_event}
\leanok
For fixed wheel coefficients, the potential loads on the common-spoke class of every pair of hubs concentrate simultaneously around their actual expectations under the original raw law. The event is common to all pairs; the statement does not assert uniformity over arbitrary outcome-dependent domains.
\end{lemma}
\begin{proof}\uses{bp:ordinary-two-hub-coordinate-locality}\leanok Apply the supported bounded-differences estimate and a union bound over hub pairs.\end{proof}

\begin{lemma}[Two-hub resource marginal]\label{bp:ordinary-two-hub-marginal}
\lean{Oberwolfach.ordinary_pair_hub_resource_probability}
\leanok
For two distinct hubs, an ordinary rim resource has its original marginal, while a spoke in their common neighbourhood has marginal equal to the product of the two individual spoke probabilities.
\end{lemma}
\begin{proof}\uses{bp:ordinary-two-hub-coordinate-locality}\leanok The two spoke events read disjoint coordinate families associated with different hubs.\end{proof}

\begin{lemma}[Two-hub candidate coordinate separation]\label{bp:ordinary-two-hub-independence}
\lean{Oberwolfach.ordinary_pair_hub_candidate_coordinates_disjoint}
\leanok
Distinct resources of one admissible cycle read disjoint primitive coordinate sets when its vertices avoid the permutation-conflict graph.
\end{lemma}
\begin{proof}\uses{bp:ordinary-two-hub-coordinate-locality}\leanok Compare the two original supports, using distinct hubs or the conflict-free cycle condition.\end{proof}


\begin{lemma}[Omitted-root two-hub pattern expectation]\label{bp:ordinary-two-hub-pattern-expectation}
\lean{Oberwolfach.ordinary_pair_hub_pattern_load_expectation,Oberwolfach.ordinary_pair_hub_pattern_marginal_expectation}
\leanok
For an admissible cycle, the expectation of the generic pattern load at a fixed root is the sum of its fixed coefficients multiplied by the product of the actual resource marginals, omitting the root. For distinct hubs, the factors are the original rim marginals and the products of the two spoke marginals. This is an exact identity under the original raw law; normalization against the parent cycle weights is a separate step.
\end{lemma}
\begin{proof}\uses{bp:ordinary-two-hub-marginal,bp:ordinary-two-hub-independence}\leanok The candidate resources read disjoint coordinates, so their state probabilities factor; substitute the checked marginal probabilities.\end{proof}

\begin{lemma}[Normalized parent cycle load]\label{bp:ordinary-normalized-parent-load}
\lean{Oberwolfach.ordinary_normalized_admissible_vertex_load}
\leanok
The original fractional cycle weights, restricted only by the conflict graph and divided by the parent degree parameter, have the stated vertex-load error controlled by the conflict loss and degree deviation. The weights themselves are not recomputed.
\end{lemma}
\begin{proof}\uses{bp:ordinary-seed-means}\leanok Combine the conflict-free parent vertex-load estimate with the parent degree deviation.\end{proof}

\section{Degree-normalized two-hub mean and concentration}

\begin{lemma}[Degree-normalized wheel seed]\label{bp:ordinary-degree-normalized-seed}
\lean{Oberwolfach.ordinary_normalized_spoke_seed,Oberwolfach.ordinary_degree_normalized_spoke_seed,Oberwolfach.ordinary_degree_normalized_seed_bounds}
\leanok
Dividing the original fixed cycle weights by half the parent degree parameter gives wheel coefficients at the concentration scale and an explicit parent spoke-load error controlled by the degree deviation and conflict loss.
\end{lemma}
\begin{proof}\uses{bp:ordinary-normalized-parent-load}\leanok Transfer the cycle vertex sum to the wheel seed and use the parent degree lower bound for the coefficient estimate.\end{proof}

\begin{lemma}[Spoke-root mean on a two-hub class]\label{bp:ordinary-two-hub-spoke-mean}
\lean{Oberwolfach.ordinary_pair_hub_marginal_on_candidate_ext,Oberwolfach.ordinary_pair_hub_pattern_spoke_mean_ext,Oberwolfach.ordinary_pair_hub_normalized_mean}
\leanok
For two distinct hubs with positive fixed spoke marginals, the omitted-spoke-root expectation factors into the rim and common-spoke sampling densities times the original seed's spoke load. After dividing by those fixed densities, its error is bounded by the parent degree and conflict error.
\end{lemma}
\begin{proof}\uses{bp:ordinary-two-hub-pattern-expectation,bp:ordinary-degree-normalized-seed}\leanok Evaluate the omitted-root product and apply the fixed-seed parent-load estimate.\end{proof}

\begin{lemma}[Common concentration for the fixed parent seed]\label{bp:ordinary-two-hub-seed-concentration}
\lean{Oberwolfach.ordinary_pair_hub_seed_concentration}
\leanok
The degree-normalized original cycle weights, fixed before sampling, satisfy a single concentration event for all pairs of hubs. The event is centered at the actual expectation. Combining it with the preceding mean in one representation and constructing the actual factor profiles remain separate obligations.
\end{lemma}
\begin{proof}\uses{bp:ordinary-two-hub-concentration,bp:ordinary-degree-normalized-seed}\leanok Apply the common concentration theorem using the checked coefficient bound.\end{proof}

\begin{lemma}[Pattern-load concentration on the original law]\label{bp:ordinary-pattern-concentration}
\lean{Oberwolfach.ordinaryPatternLoad_supported_tail,Oberwolfach.ordinaryPatternLoad_supported_common_event}
\leanok
For fixed wheel coefficients and locally supported resource states, the root-omitted pattern loads satisfy the displayed bounded-differences estimates simultaneously. The result is stated directly in the pattern-load representation.
\end{lemma}
\begin{proof}\leanok Apply the supported pattern-load tail estimate and a union bound over the fixed resource roots.\end{proof}

\begin{lemma}[Common normalized loads for two hubs]\label{bp:ordinary-two-hub-normalized-event}
\lean{Oberwolfach.ordinary_pair_hub_pattern_load_common_event_ext,Oberwolfach.ordinary_pair_hub_seed_pattern_concentration,Oberwolfach.ordinary_pair_hub_normalized_load_common_event}
\leanok
Under the displayed parent-weight, conflict, degree and sampling hypotheses, one event in the original sampling law controls the normalized omitted-spoke-root load at every vertex for every pair of distinct hubs. The coefficients and density normalization are fixed before sampling; the explicit error combines the concentration deviation with the parent degree and conflict losses. This result does not yet construct the factor profiles or prove the main theorem.
\end{lemma}
\begin{proof}\uses{bp:ordinary-pattern-concentration,bp:ordinary-two-hub-coordinate-locality,bp:ordinary-two-hub-spoke-mean,bp:ordinary-degree-normalized-seed}\leanok Specialize the common concentration event to the fixed degree-normalized seed and combine it with the checked mean, using the same original sample for all pairs and vertices.\end{proof}

\begin{lemma}[Omitted-spoke test and conflict deletion]\label{bp:ordinary-pattern-test-conflict}
\lean{Oberwolfach.resource_pattern_indicator_sum,Oberwolfach.ordinary_spoke_pattern_predicate,Oberwolfach.ordinary_pattern_spoke_sum,Oberwolfach.ordinary_pattern_load_delete_conflicts,Oberwolfach.ordinary_local_potential_div}
\leanok
The root-omitted pattern load is the sum of the fixed cycle weights over wheels satisfying the explicit rim and spoke tests. Deleting conflict edges leaves that load unchanged, and division of the fixed parent weights commutes with the local cycle potential. The identification with local weighted degree and the factor callers remain to be checked.
\end{lemma}
\begin{proof}\leanok Expand the pattern as an indicator, use the omitted-spoke resource characterization, and note that every candidate rim avoids the conflict graph.\end{proof}

\begin{lemma}[Parent packing load and identity restriction]\label{bp:marked-parent-load-identity}
\lean{Oberwolfach.marked_reference_parent_packing_load,Oberwolfach.raw_marked_restrict_pair_identity}
\leanok
The packing-load hypothesis on the full hub set gives the parent reference load at the stated density and degree parameters whenever its error is at most $\kappa/4$. Restricting the raw marked outcome by identity maps on the vertex and hub sets returns the original outcome. Establishing these hypotheses and the actual first-level embeddings in the simultaneous sampling construction remains open.
\end{lemma}
\begin{proof}\uses{bp:markedpackingfullhubloads}\leanok Apply the full-hub load estimate and use the identity maps on each type of raw coordinate.\end{proof}

\begin{lemma}[Zero deletion in the marked packing event]\label{bp:marked-packing-zero-deletion}
\lean{Oberwolfach.graph_bottom_edgeFinset_any_instance,Oberwolfach.marked_packing_zero_deletion_any_instance}
\leanok
The zero-deletion instance of the stated adaptive packing hypothesis is the undeleted packing-load estimate. The edge set of the empty graph is empty independently of the chosen finite edge-set enumeration. This implication alone does not establish the adaptive event.
\end{lemma}
\begin{proof}\leanok Set the deletion budget to zero and use empty rim and spoke deletion graphs.\end{proof}

\begin{lemma}[Parent reference load from adaptive packing]\label{bp:marked-parent-load-zero}
\lean{Oberwolfach.marked_reference_parent_load_from_zero_deletion_split}
\leanok
With the displayed density and degree parameters, the zero-deletion case of the stated adaptive packing hypothesis implies the parent reference load bound, provided its error is at most $\kappa/4$. The result retains the fixed parent sequence weights; the actual sampled hypothesis and numerical error bound still require verification.
\end{lemma}
\begin{proof}\uses{bp:marked-packing-zero-deletion,bp:marked-parent-load-identity}\leanok Apply the zero-deletion implication and the parent packing-load bound.\end{proof}

\begin{lemma}[Omitted-spoke pattern and local weighted degree]\label{bp:ordinary-pattern-local-degree}
\lean{Oberwolfach.mixed_ite_and,Oberwolfach.ordinary_pattern_spoke_sum_val,Oberwolfach.ordinary_spoke_indicator_sum_eq_local_mixed,Oberwolfach.ordinary_pattern_spoke_eq_local_bridge,Oberwolfach.ordinary_pattern_spoke_local_degree}
\leanok
The root-omitted pattern load equals the corresponding local cycle potential. At a vertex in the specified spoke domain, dividing the fixed parent weights by the prescribed constant gives a local weighted degree equal to the normalized pattern load. The statements apply to the displayed graphs and use the original parent weights.
\end{lemma}
\begin{proof}\uses{bp:ordinary-pattern-test-conflict}\leanok Express the load as a wheel sum, reindex the hub-selected wheels by their cycles, and use the local cycle restriction at a present root.\end{proof}

\begin{lemma}[Local degrees on two-hub sampled classes]\label{bp:ordinary-two-hub-local-degree}
\lean{Oberwolfach.ordinary_pattern_load_delete_conflicts_high,Oberwolfach.ordinary_pair_hub_local_degree_identity_high}
\leanok
For every realized sample and vertex in a two-hub common-spoke class, the local cycle-weighted degree in the sampled graph with conflict edges deleted equals the normalized root-omitted load in the original sampled graph. The coefficients are the fixed parent cycle weights.
\end{lemma}
\begin{proof}\uses{bp:ordinary-pattern-local-degree,bp:ordinary-pattern-test-conflict}\leanok Apply the local-degree identity and remove the conflict edges, which no candidate wheel uses.\end{proof}

\begin{lemma}[Two-sided bounds for normalized cycle weights]\label{bp:ordinary-normalized-cycle-bounds}
\lean{Oberwolfach.normalized_weight_bounds,Oberwolfach.normalized_cycle_weight_bounds}
\leanok
The displayed lower and upper bounds on fixed parent cycle weights imply two-sided bounds at the cycle-factor scale after division by half the parent degree and a positive density product bounded away from zero. These estimates do not select new weights.
\end{lemma}
\begin{proof}\leanok Use the degree bounds and monotonicity of division by the positive normalization factors.\end{proof}

\begin{lemma}[Fixed coefficient bounds for every two-hub class]\label{bp:ordinary-pair-hub-weight-bounds}
\lean{Oberwolfach.pair_hub_density_bounds,Oberwolfach.ordinary_local_cycle_weight_bounds,Oberwolfach.ordinary_pair_hub_cycle_weight_bounds,Oberwolfach.ordinary_pair_hub_named_local_degree_identity}
\leanok
Under the displayed fixed positive lower bounds on the rim and spoke densities, the original parent cycle weights induce pair-independent two-sided bounds for every realized two-hub local cycle family at the cycle-factor scale. The named local degree is exactly the normalized root-omitted pattern load at each present root. The sampling hypotheses and degree estimates still need to be checked at the factor callers.
\end{lemma}
\begin{proof}\uses{bp:ordinary-normalized-cycle-bounds,bp:ordinary-two-hub-local-degree}\leanok Bound the density product and apply the deterministic normalized-weight estimates to each restricted parent cycle.\end{proof}

\begin{lemma}[Simultaneous local weighted degrees]\label{bp:ordinary-two-hub-local-degree-event}
\lean{Oberwolfach.ordinary_pair_hub_local_degree_common_event}
\leanok
Under the displayed fixed-parent-weight and sampling hypotheses, one event in the original sampling law bounds the local cycle-weighted degrees at all present vertices in every two-hub common-spoke class, using the sampled graph with conflict edges deleted. The event retains the explicit normalization and error from the root-omitted load estimate. Its application to the factor profiles additionally requires the common-neighbourhood estimates and numerical parameter bounds.
\end{lemma}
\begin{proof}\uses{bp:ordinary-two-hub-normalized-event,bp:ordinary-two-hub-local-degree}\leanok Apply the deterministic local-degree identity simultaneously on the existing original-law event.\end{proof}

\begin{lemma}[Simultaneous sampled pair-profile inputs]\label{bp:ordinary-sampled-profile-inputs}
\lean{Oberwolfach.ordinary_two_hub_neighbourhood_geometry,Oberwolfach.ordinary_small_root_density,Oberwolfach.ordinary_dense_two_hub_geometry,Oberwolfach.ordinary_sampled_two_hub_profile_inputs,Oberwolfach.ordinary_sampled_profile_property_from_inputs,Oberwolfach.ordinary_sampled_profile_property_common_event}
\leanok
Under the displayed dense-parent, fixed-weight, sampling and numerical hypotheses, a single event in the original sampling law supplies, simultaneously for every pair of hubs, the common-neighbourhood inequalities and fixed local cycle weights with the required two-sided and weighted-degree bounds. Its probability has the explicit union-bound error. The neighbourhood estimate includes repeated vertices and hubs and subtracts the realized conflict-edge loss. These inputs do not by themselves prove the profile construction or the manuscript callers.
\end{lemma}
\begin{proof}\uses{bp:ordinary-two-hub-local-degree-event,bp:ordinary-pair-hub-weight-bounds}\leanok Intersect the raw row and degree events without an independence assumption, apply the deterministic neighbourhood and fixed-weight bounds, and use the displayed error margins.\end{proof}

\begin{lemma}[Choosing the sampled-profile margins]\label{bp:ordinary-sampled-profile-margins}
\lean{Oberwolfach.ordinary_sampled_profile_margins}
\leanok
For fixed positive parent-degree slack, density lower bounds and target weighted-degree tolerance, the neighbourhood and load error margins can be chosen before an order threshold so that the displayed numerical inequalities hold above that threshold. This does not establish the remaining caller or probability requirements.
\end{lemma}
\begin{proof}\leanok Choose the errors as fixed fractions of the available slack and take the order large enough to absorb the conflict-degree terms.\end{proof}

\begin{lemma}[Uniform sampled pair-profile inputs]\label{bp:ordinary-uniform-sampled-profile-inputs}
\lean{Oberwolfach.ordinary_profile_failure_margins,Oberwolfach.ordinary_profile_cardinality_bounds,Oberwolfach.ordinary_profile_error_bound,Oberwolfach.ordinary_profile_event_exists,Oberwolfach.ordinary_sampled_profile_property_exists,Oberwolfach.ordinary_uniform_sampled_profile_property_exists}
\leanok
Under the stated fixed positive slack and density hypotheses and the displayed parent graph and fractional-weight conditions, the error constants can be chosen before an order threshold so that an actual sample supplies all the two-hub profile inputs simultaneously. The theorem retains the original parent weights. The step from these inputs to the manuscript's factor and absorber callers remains to be checked.
\end{lemma}
\begin{proof}\uses{bp:ordinary-sampled-profile-inputs,bp:ordinary-sampled-profile-margins}\leanok Bound the two failure terms uniformly by inverse powers of the order and choose an outcome from the resulting positive-probability event.\end{proof}

\begin{lemma}[Transport between equal hub domains]\label{bp:ordinary-hub-domain-transport}
\lean{Oberwolfach.ordinary_hub_domains_equal,Oberwolfach.ordinary_local_weights_transport}
\leanok
Equivalent descriptions of the two-hub vertex sets define equal finite sets. A local cycle-weight family and its vertex degrees transport across equal domains without changing their values or bounds.\end{lemma}
\begin{proof}\leanok Compare memberships and transport along the resulting equality.\end{proof}

\begin{lemma}[Pair-profile transport and sampled longer-cycle profiles]\label{bp:ordinary-uniform-pair-profiles}
\lean{Oberwolfach.ordinary_local_weights_transport_explicit,Oberwolfach.ordinary_pair_profiles_from_high_domain_explicit,Oberwolfach.ordinary_sampled_property_to_pair_profiles,Oberwolfach.ordinary_uniform_sampled_pair_profiles_exists}
\leanok
Under the displayed parent-graph, fixed fractional-weight and sampling hypotheses, an outcome of the original sampling experiment admits the longer-cycle pair profiles. The profile constants and order threshold are chosen in the stated order. Equality of the finite hub domains and independence of their finite enumerations preserve the weights and weighted degrees. Applications to the manuscript's transformer and absorber constructions remain to be verified.
\end{lemma}
\begin{proof}\uses{bp:ordinary-uniform-sampled-profile-inputs,bp:ordinary-hub-domain-transport}\leanok Transport the sampled weights to the domain used by the existing profile theorem and apply that theorem.\end{proof}

\begin{lemma}[Four-root sampled triangle inputs]\label{bp:ordinary-triangle-sampled-inputs}
\lean{Oberwolfach.ordinary_many_root_two_hub_geometry,Oberwolfach.ordinary_root_density,Oberwolfach.dense_four_root_neighbourhood,Oberwolfach.ordinary_dense_four_root_two_hub_geometry,Oberwolfach.ordinary_sampled_triangle_profile_inputs,Oberwolfach.ordinary_sampled_triangle_property_from_inputs,Oberwolfach.ordinary_sampled_triangle_property_common_event}
\leanok
Under the displayed three-quarter parent-degree, fixed-weight and sampling hypotheses, one original-law event supplies the triangle profile inputs, including four-root common-neighbourhood bounds. Repeated roots are allowed and the conflict-edge loss is counted in the realized graph.\end{lemma}
\begin{proof}\uses{bp:ordinary-two-hub-local-degree-event}\leanok Combine the fixed-root row and local-degree events on the same law and apply the explicit density and deletion estimates.\end{proof}

\begin{lemma}[Uniform sampled triangle pair profiles]\label{bp:ordinary-uniform-triangle-pair-profiles}
\lean{Oberwolfach.ordinary_triangle_failure_margins,Oberwolfach.ordinary_triangle_profile_cardinality_bounds,Oberwolfach.ordinary_triangle_profile_error_bound,Oberwolfach.ordinary_sampled_triangle_margins,Oberwolfach.ordinary_sampled_triangle_property_exists,Oberwolfach.ordinary_uniform_sampled_triangle_property_exists,Oberwolfach.ordinary_triangle_profiles_from_high_domain_explicit,Oberwolfach.ordinary_triangle_property_to_pair_profiles,Oberwolfach.ordinary_uniform_sampled_triangle_pair_profiles_exists}
\leanok
Under the stated fixed positive slack and density hypotheses and the displayed parent graph and fractional-weight conditions, the constants can be chosen before an order threshold so that an actual sample admits the triangular pair profiles. This result is conditional on those inputs and does not prove the triangle transformer or the main theorem.\end{lemma}
\begin{proof}\uses{bp:ordinary-triangle-sampled-inputs,bp:ordinary-hub-domain-transport}\leanok Bound the failure terms, choose a sample on the common event, and apply the triangular profile theorem after transporting equal domains.\end{proof}

\begin{lemma}[Fixed slack and common sampled profiles]\label{bp:ordinary-fixed-slack-joint-profiles}
\lean{Oberwolfach.ordinary_sampled_profile_margins_fixed_beta,Oberwolfach.ordinary_sampled_triangle_margins_fixed_beta,Oberwolfach.ordinary_uniform_sampled_pair_profiles_and_inputs_joint_fixed_beta,Oberwolfach.ordinary_uniform_sampled_triangle_pair_profiles_and_inputs_joint_fixed_beta}
\leanok
Under the displayed parent, fractional-weight and sampling hypotheses, the positive profile slack is fixed before the weighted-degree tolerance. For an event of the original sampling experiment whose probability exceeds the explicit profile-failure bound, there is an outcome in that event satisfying both the sampled-input conditions and the corresponding pair-profile conclusion. The triangle and longer-cycle statements have separate hypotheses. The other event and its probability, and the remaining absorber hypotheses, must still be established for the manuscript's application.
\end{lemma}
\begin{proof}\uses{bp:ordinary-uniform-pair-profiles,bp:ordinary-uniform-triangle-pair-profiles}\leanok Choose the slack as a fixed positive function of the parent slack and density bounds, intersect the original-law events, and transport the fixed weights on the resulting outcome.\end{proof}

\begin{lemma}[Neighbourhood consequences for the absorber]\label{bp:ordinary-profile-absorber-geometry}
\lean{Oberwolfach.ordinary_long_profile_absorber_geometry,Oberwolfach.ordinary_triangle_profile_absorber_geometry}
\leanok
If the sampled two-hub profile inputs hold and the hub set is nonempty, they imply the displayed pair-neighbourhood conditions. The triangle inputs also imply the required four-neighbourhood conditions. These statements do not establish the other common-hub or embedding assumptions of the absorber.
\end{lemma}
\begin{proof}\uses{bp:ordinary-fixed-slack-joint-profiles}\leanok Choose one hub and compare the relevant finite sets; in the triangle case repeat roots in the four-root hypothesis.\end{proof}

\begin{lemma}[Sampled absorber inputs from the cycle threshold]\label{bp:ordinary-sampled-absorber-from-threshold}
\lean{Oberwolfach.ordinary_long_absorber_slack_choice,Oberwolfach.ordinary_triangle_absorber_slack_choice,Oberwolfach.ordinary_index_event_absorber_hub_bound_from_density,Oberwolfach.ordinary_preparation_probability_exceeds_profile_failure,Oberwolfach.ordinary_uniform_sampled_long_absorber_inputs_from_threshold,Oberwolfach.ordinary_uniform_sampled_triangle_absorber_inputs_from_threshold}
\leanok
With the displayed positive density bounds, homogeneous spoke sampling, nearly regular parent graph and cycle-threshold minimum degree, the parent weights are obtained from the positive fractional-cycle-weight theorem. One original-law outcome supplies the corresponding sampled profile inputs, pair profiles and bounded-order common-hub inequalities. The constants are chosen in the stated order, with the profile slack small enough for both the pair and hub estimates. These results do not establish the manuscript's full preparation hypotheses.
\end{lemma}
\begin{proof}\uses{bp:ordinary-fixed-slack-joint-profiles,bp:ordinary-profile-absorber-geometry}\leanok Choose the slack below both density margins, use the existing high-probability row and index event, and intersect it with the weighted profile event on the same sampling law.\end{proof}

\begin{lemma}[Relative absorber from sampled profile inputs]\label{bp:ordinary-absorber-from-sampled-inputs}
\lean{Oberwolfach.ordinary_long_cycle_ambient_absorber_from_sampled_inputs,Oberwolfach.ordinary_triangle_ambient_absorber_from_sampled_inputs}
\leanok
The sampled longer-cycle and triangle profile inputs give the respective relative absorber conclusions under the displayed bounded-remainder, divisibility, host-containment and common-hub conditions. For length four a separate three-neighbourhood hypothesis remains. The construction of these inputs for the actual manuscript caller is not yet verified.
\end{lemma}
\begin{proof}\uses{bp:ordinary-profile-absorber-geometry,bp:ordinary-sampled-absorber-from-threshold}\leanok Apply the profile and geometric implications and then the relative absorber theorem.\end{proof}

\begin{lemma}[Uniform ordinary absorbers on one sample]\label{bp:uniform-ordinary-sampled-absorbers}
\lean{Oberwolfach.ordinary_uniform_sampled_triangle_absorbers_from_threshold,Oberwolfach.ordinary_uniform_sampled_four_cycle_absorbers_from_threshold,Oberwolfach.ordinary_uniform_sampled_long_absorbers_from_threshold}
\leanok
For cycle lengths three, four, and at least five, the respective statements construct an outcome of the specified sampling experiment before the bounded remainder is selected. Under their displayed parent threshold, near-regularity, uniform spoke-density, and sampling hypotheses, every remainder satisfying the stated balance, divisibility, size and containment conditions has a relative absorber in that outcome. The four-cycle proof derives its three-neighbourhood estimate on the same event. These statements do not construct one fixed absorber for all remainders or discharge the manuscript's full preparation and factor-theorem callers.
\end{lemma}
\begin{proof}\uses{bp:ordinary-sampled-absorber-from-threshold,bp:ordinary-absorber-from-sampled-inputs}\leanok Choose the shared slack before the absorber tolerance, obtain all sampled inputs on one event, and apply the corresponding relative absorber theorem to each eligible remainder.\end{proof}

\begin{lemma}[Simultaneous marked load error margins]\label{bp:marked-reference-joint-load-margin}
\lean{Oberwolfach.marked_reference_parent_load_from_adaptive_error,Oberwolfach.marked_reference_joint_load_margin_choice}
\leanok
Positive numerical parameters can be chosen so that the explicit parent and local load errors are simultaneously below their respective quarter margins. The zero-deletion case of the adaptive packing hypothesis then yields the corresponding parent load bound. These conditional statements do not construct the marked reference-level data or its common sampling event.
\end{lemma}
\begin{proof}\leanok Choose the shared leading errors and the separate remaining errors as positive fractions of the available margins; apply the zero-deletion implication with the original weights.\end{proof}

\begin{lemma}[Ordinary absorbers and transformers on one sample]\label{bp:joint-ordinary-absorbers-transformers}
\lean{Oberwolfach.ordinary_uniform_sampled_triangle_absorbers_and_transformers_with_extra_event,Oberwolfach.ordinary_uniform_sampled_four_cycle_absorbers_and_transformers_with_extra_event,Oberwolfach.ordinary_uniform_sampled_long_absorbers_and_transformers_with_extra_event,Oberwolfach.ordinary_extra_event_probability_margin}
\leanok
Under the displayed parent threshold, near-regularity, homogeneous spoke-density and sampling hypotheses, one original-law outcome satisfies a specified additional event of sufficiently large probability and supports both the relative absorber and transformer conclusions for every later bounded input satisfying their respective hypotheses. The three statements cover cycle lengths three, four and at least five. The outcomes of the two constructions are not asserted to be mutually disjoint, and no single absorber is asserted to work for all remainders. Establishing the actual additional event and all manuscript-caller conditions remains open.
\end{lemma}
\begin{proof}\uses{bp:uniform-ordinary-sampled-absorbers,bp:ordinary-sampled-absorber-from-threshold}\leanok Choose the smaller of the two weighted-degree tolerances, preserve the original weights, and intersect the extra event with the row, index and profile events on the same sampling law.\end{proof}

\begin{lemma}[Marked parent loads on the retained sampling law]\label{bp:marked-parent-load-retained-law}
\lean{Oberwolfach.marked_parent_adaptive_load_common_event,Oberwolfach.marked_parent_and_local_reference_load_uniform_density,Oberwolfach.marked_reference_densities_load_with_extra_event}
\leanok
Under the stated weight, density, size and error hypotheses, a one-level specialization of the retained sampling theorem supplies the full-parent potential-load and local-spoke-load bounds. The reference-density form uses the original fixed weights and the normalizations \(s_0=(r+4)t\), \(p_{\rm ref}=\rho s_0/2\), \(k_{\rm ref}=p_{\rm ref}/(r+4)\) and \(D_{\rm ref}=2|U|/\rho\). It can be intersected with any further event on the same retained law with the displayed union-bound loss. This does not construct the nested reference levels or establish the remaining common-event and first-cover conditions.
\end{lemma}
\begin{proof}\uses{bp:marked-reference-joint-load-margin}\leanok Specialize the existing adaptive packing event to the full parent domains and identity embeddings at one level; use the zero-deletion case, the density identities and the original-law union bound.\end{proof}

\begin{lemma}[Parent weights and the initial vortex domains]\label{bp:marked-parent-vortex-weights}
\lean{Oberwolfach.marked_parent_reference_load_with_shared_vortex_weights,Oberwolfach.successor_weight_family_preserve_parent,Oberwolfach.vortex_full_successor_weight_families_with_fixed_parent,Oberwolfach.vortex_initial_parent_identities,Oberwolfach.vortexInitialVertexOrderIso,Oberwolfach.vortexInitialHubEquiv,Oberwolfach.vortexInitialPermutation,Oberwolfach.vortex_initial_permutation_commutes}
\leanok
Under the stated parent degree, regularity and density conditions, the reference parent-load event admits fixed successor-weight constants of the same form as the full-vortex weight theorem. A compatible supplied parent weight can be retained exactly while choosing the weights at the other vortex levels. The initial vortex sets are full, have canonical vertex and hub equivalences, and their reference degree equals the original degree; the induced initial permutation commutes with the original permutation. These facts do not themselves construct the compatible all-level event or protected-reservoir caller.
\end{lemma}
\begin{proof}\uses{bp:marked-parent-load-retained-law}\leanok Choose the shared successor constants and update the level-zero member of the weight family; use the vortex initial-set identities and the canonical subtype maps.\end{proof}


\begin{lemma}[Relabelling the retained sampling inputs]\label{bp:marked-retention-relabel}
\lean{Oberwolfach.vortex_initial_permutation_pow,Oberwolfach.vortex_initial_no_short_cycles,Oberwolfach.vortex_initial_permutation_close_iff,Oberwolfach.raw_marked_start_hub_mem_relabel,Oberwolfach.raw_marked_start_hub_relabel,Oberwolfach.raw_marked_retention_minimum_relabel,Oberwolfach.start_retention_law_relabel,Oberwolfach.marked_conditional_retention_law_relabel,Oberwolfach.markedRetainedPairSecond,Oberwolfach.markedRetainedPairRestriction,Oberwolfach.marked_retained_pair_conditional_map,Oberwolfach.marked_retained_pair_restriction_probability,Oberwolfach.raw_marked_restrict_pair_inverse,Oberwolfach.finiteLaw_sigma_dependent_transport,Oberwolfach.marked_vortex_local_parent_loads_from_adaptive}
\leanok
The canonical initial permutation preserves iterates, short-cycle conditions and closeness. Under bijective relabelling commuting with the permutations, the eligible hub sets and their common retention minimum agree. For each fixed raw outcome, the conditional law of uniformly retaining eligible hubs at all starts is transported exactly by the corresponding bijections. The raw outcomes are bijectively relabelled; the actual retained-outcome map preserves the conditional law, and a general dependent two-stage transport lemma is available. At an embedded level, the displayed adaptive and local hypotheses imply the stated load bounds. The complete two-stage retained law is transported exactly under these bijections, with its conditional choices tied to the same raw outcome. The all-level common event and protection of the absorber in its actual caller remain to be checked.
\end{lemma}
\begin{proof}\uses{bp:marked-parent-vortex-weights}\leanok Compare eligible sets via the commuting permutation maps; transport fixed-size subsets by their hub bijections, reindex the product over starts, and preserve the conditioning on the original raw outcome. The remaining load assertion is a specialization of the displayed adaptive inequalities.\end{proof}


\begin{lemma}[Initial-vortex typed sampling identity]\label{bp:marked-vortex-typed-relabel}
\lean{Oberwolfach.vortexInitialVertexEquiv,Oberwolfach.vortexInitialVertexOrderEmbedding,Oberwolfach.vortex_initial_embedding_equiv_agree,Oberwolfach.start_retained_hub_relabel,Oberwolfach.start_retained_hub_relabel_inverse,Oberwolfach.marked_retained_hub_restriction,Oberwolfach.marked_retained_start_restriction,Oberwolfach.retained_interval_state_relabel,Oberwolfach.marked_retained_state_restriction,Oberwolfach.marked_retained_spoke_restriction,Oberwolfach.marked_retained_rim_restriction,Oberwolfach.marked_retained_typed_restriction,Oberwolfach.marked_vortex_initial_typed_restriction,Oberwolfach.marked_vortex_initial_retention_law,Oberwolfach.markedRetainedSampleArcEquiv,Oberwolfach.marked_retained_arc_equiv_endpoints,Oberwolfach.marked_relabel_physical_twist}
\leanok
Under the canonical equivalences from the full initial vortex sets, the whole retained sampling law is unchanged, and the retained start and interval-state predicates, both spoke types and both rim types agree on the same outcome. The corresponding typed sample-arc sets are equivalent with endpoints transported by the vertex equivalence; in particular, physical edges using the permutation twist are transported exactly. This identifies the initial-domain representations. It does not identify a later modified host with the original sample or establish that a selected absorber survives its first cover-down.
\end{lemma}
\begin{proof}\uses{bp:marked-retention-relabel}\leanok Transport the eligible sets and retained choices, then use the commuting permutation to compare windows, states, typed rims and typed spokes. The bijection on ordered endpoint pairs induces the arc equivalence, and commutation also preserves the twisted physical edge.\end{proof}


\begin{lemma}[Weights and root-omitted loads under relabelling]\label{bp:marked-vortex-weight-load-relabel}
\lean{Oberwolfach.successor_sequence_valid_relabel,Oberwolfach.successor_sequence_weights_relabel,Oberwolfach.vortex_initial_successor_weights_relabel,Oberwolfach.marked_wheel_equiv_loads,Oberwolfach.raw_marked_marginal_relabel,Oberwolfach.resource_root_omitted_load_relabel,Oberwolfach.resource_omitted_pair_weight_relabel,Oberwolfach.marked_physical_root_omitted_load_relabel,Oberwolfach.marked_canonical_arc_membership_relabel,Oberwolfach.marked_typed_occurrence_mem_relabel,Oberwolfach.marked_typed_occurrence_image_relabel,Oberwolfach.marked_typed_root_omitted_load_relabel}
\leanok
A supplied successor-sequence weighting transports through the initial vortex equivalence with the same values, constants and edge and position sums. Typed-resource marginals are preserved under permutation-compatible embeddings. When the wheel, arc, row and hub maps are bijections, with the stated endpoint uniqueness and typed-relation agreement, the typed occurrence sets and their root-omitted loads transport exactly; the omitted root need not be retained. A separate identity treats physical resources. These results still require the actual level-data and fixed-weight substitutions and do not establish the protected-absorber caller.
\end{lemma}
\begin{proof}\uses{bp:marked-vortex-typed-relabel}\leanok Reindex the finite sums by the corresponding bijections. For typed occurrences, endpoint uniqueness identifies the directed arcs, while the two spoke labels are preserved separately.\end{proof}


\begin{lemma}[Initial parent-load transfer and one-sided rounding]\label{bp:marked-parent-load-transfer-rounding}
\lean{Oberwolfach.marked_retained_potential_load_relabel,Oberwolfach.marked_retained_parent_load_bound_relabel,Oberwolfach.marked_retained_parent_load_event_relabel,Oberwolfach.marked_vortex_initial_parent_load_bound,Oberwolfach.marked_potential_seed_relabel,Oberwolfach.bipartite_one_sided_rounding,Oberwolfach.bipartite_rounding_with_tail_ceilings}
\leanok
Under the stated bijections, a parent typed-potential-load bound for every nonzero marginal transfers to the initial vortex domains with the same sample, fixed weight and error. The transfer preserves an arbitrary additional event on the same probability space. Separately, a fractional bipartite selection with integral head degrees can be rounded so that those head degrees remain exact and every tail degree is at most the ceiling of its fractional degree. The parent event and its actual parameter substitutions, and the application of this rounding to the marked-indegree correction, remain to be established.
\end{lemma}
\begin{proof}\uses{bp:marked-vortex-weight-load-relabel}\leanok Apply the typed root-load and marginal identities to all roots simultaneously. For rounding, add a dummy edge at each tail to complete its degree to the chosen integer and apply bipartite integer-degree rounding.\end{proof}


\begin{lemma}[Marked-indegree correction by retyping arcs]\label{bp:marked-indegree-correction}
\lean{Oberwolfach.bipartite_one_sided_rounding_support,Oberwolfach.bipartite_rounding_with_tail_ceilings_support,Oberwolfach.uniform_incoming_weight_bounds,Oberwolfach.uniform_incoming_weight_columns,Oberwolfach.uniform_incoming_rounding,Oberwolfach.uniform_incoming_weight_upper,Oberwolfach.uniform_incoming_rounding_bounded,Oberwolfach.marked_indegree_flip_reduces,Oberwolfach.marked_indegree_flip_increases,Oberwolfach.marked_deviation_cast_abs,Oberwolfach.marked_flexible_retyping,Oberwolfach.retyped_physical_degree_le,Oberwolfach.marked_indegree_correction_physical,Oberwolfach.marked_indegree_correction_manuscript_bound,Oberwolfach.marked_indegree_correction_oriented_graph}
\leanok
Represent the directed arcs by a finite subtype of ordered vertex pairs. Under the deviation and flexible incoming-degree hypotheses of the manuscript's marked-indegree correction lemma, the arc types can be changed within the prescribed flexible set so that each vertex has the required marked indegree. Directions are unchanged, and the graph of physical pairs whose type changes satisfies the manuscript's displayed maximum-degree bound. The proof obtains the stronger bound \((1+1/\rho)\eta n+1\). This establishes the finite correction lemma; its use in the long preparation still requires the actual flexible-arc and degree estimates.
\end{lemma}
\begin{proof}\uses{bp:marked-parent-load-transfer-rounding}\leanok Give every eligible incoming arc at a head the same fractional weight. Complete the fractional tail degrees with dummy edges and round, retaining exact head degrees, the tail ceilings and zero support. Flip precisely the chosen types and bound the physical degree by the changed incoming and outgoing counts.\end{proof}


\begin{lemma}[Protected correction from candidate bounds]\label{bp:marked-protected-candidate-correction}
\lean{Oberwolfach.marked_indegree_correction_finsets,Oberwolfach.marked_indegree_correction_protect_reserved,Oberwolfach.incoming_candidate_loss,Oberwolfach.incoming_forbidden_union_loss,Oberwolfach.incoming_tail_domain_card_le,Oberwolfach.incoming_flexible_type_availability,Oberwolfach.incoming_flexible_availability_real,Oberwolfach.marked_sample_arc_endpoint_injective,Oberwolfach.marked_correction_from_candidate_bounds}
\leanok
For finite directed arc sets, the marked-indegree correction admits a set-theoretic formulation with the stated symmetric-difference and physical-degree conclusions. If the flexible set avoids reserved physical pairs, their arc types are unchanged. A deterministic candidate-loss lemma supplies the flexible incoming-degree hypotheses from the displayed baseline counts, losses, forbidden-set bounds and numerical margins. The proof still requires those hypotheses to be established for the particular sets in the long preparation; it does not prove that the later host or cover-down preserves the full absorber.
\end{lemma}
\begin{proof}\uses{bp:marked-indegree-correction}\leanok Compare the original and current incoming sets, subtract the realized forbidden arcs once, and apply the finite-set correction while excluding the reserved physical pairs.\end{proof}

\begin{lemma}[Fixed loads after changing the marked host]\label{bp:marked-modified-host-load}
\lean{Oberwolfach.marked_potential_load_after_host_modification,Oberwolfach.marked_potential_load_uniform_host_modification,Oberwolfach.marked_retained_load_after_host_modification,Oberwolfach.marked_modified_potential_load_relabel,Oberwolfach.marked_modified_parent_load_bound_relabel,Oberwolfach.marked_modified_local_load_bridge}
\leanok
For a fixed parent weight, a root-omitted typed potential-wheel load remains controlled after changing rim and spoke relations, provided the displayed outgoing, incoming, row and column incidence changes obey the prescribed bounds. The conclusion holds simultaneously for all later hosts satisfying these bounds and allows additions as well as deletions. The original weight, sample, root omission and marginal condition are preserved under full-domain relabelling. These results do not establish the incidence bounds for the particular host constructed in the preparation, or prove that the absorber survives the first cover-down.
\end{lemma}
\begin{proof}\uses{bp:marked-vortex-weight-load-relabel}\leanok
Bound the changes in the individual incidences and apply the fixed-weight stability estimate, then transport the resulting load under the compatible vertex and hub equivalences.
\end{proof}

\begin{lemma}[Neighbourhood estimates under physical changes]\label{bp:marked-modified-host-neighbourhoods}
\lean{Oberwolfach.marked_row_rim_physical_change,Oberwolfach.marked_spoke_physical_change_incidence,Oberwolfach.marked_neighbourhoods_after_physical_change}
\leanok
Suppose every change to a typed rim state is supported on an edge of a physical change graph, and every change to a typed spoke state is supported on an edge of a bipartite physical change graph. Their degrees bound the corresponding row and column changes, including the permuted head of a marked rim. The displayed neighbourhood bounds persist with the original domains and prescribed error scales, whether the changes add or delete states. The change-graph support and degree bounds for the preparation's actual host remain to be verified.
\end{lemma}
\begin{proof}\uses{bp:marked-physical-deletion-profiles}\leanok
Count the changed incidences using the physical change graphs and apply the established neighbourhood-stability inequalities.
\end{proof}

\begin{lemma}[Joint estimates under physical host changes]\label{bp:marked-joint-physical-changes}
\lean{Oberwolfach.marked_rim_change_of_raw_state_change,Oberwolfach.marked_rim_change_of_raw_state_on_host,Oberwolfach.marked_rim_physical_change_incidence,Oberwolfach.marked_modified_local_load_from_physical_changes,Oberwolfach.marked_modified_local_load_uniform_physical_changes,Oberwolfach.marked_modified_joint_estimates_on_one_sample}
\leanok
A change in a typed rim state is supported on the physical pair whose raw state changed; only pairs in the fixed host graph need be considered. If the resulting rim and spoke change graphs have the displayed degree bounds, then the original fixed-domain neighbourhood and fixed-weight load estimates hold together, for every later host satisfying these bounds. The exact support and degree estimates for the host of the preparation are still required.
\end{lemma}
\begin{proof}\uses{bp:marked-modified-host-load,bp:marked-modified-host-neighbourhoods}\leanok
Compare the raw physical states, count their changed typed incidences, and apply the two stability estimates to the same sample and fixed weight.
\end{proof}

\begin{lemma}[Fresh triangle factors at boundedly many transformer hubs]\label{bp:triangle-transformer-fresh-spoke-factors}
\lean{Oberwolfach.finite_fresh_sets,Oberwolfach.triangle_transformer_sample_factors_uniform_old_size,Oberwolfach.triangle_transformer_available_neighbours,Oberwolfach.triangle_transformer_available_sample_factors,Oberwolfach.triangle_transformer_simultaneous_fresh_factors,Oberwolfach.triangular_transformer_spoke_factors,Oberwolfach.triangular_transformer_fresh_spoke_factors}
\leanok
For a bounded family of hubs, choose the size and order thresholds before the host. Under the displayed minimum-degree, bounded spoke-path, divisibility, and near-full available-set hypotheses, the actual old spoke-path sets admit pairs of edge-disjoint triangle factors after adjoining fresh sets of the prescribed size. When each available set avoids all old sets, the fresh sets and the resulting factor vertex domains can be chosen pairwise disjoint. The construction of the available sets from the manuscript's hypotheses and assembly of the complete transformer remain to be checked.
\end{lemma}
\begin{proof}\uses{bp:triangle-transformer-sample-factors,bp:ordinary-transformer-spoke-sets}\leanok
Take maxima over the finitely many old-set sizes, derive the neighbourhood estimates from minimum degree, and choose the fresh sets successively while excluding the previously chosen vertices. Apply the factor result to the actual spoke-path sets.
\end{proof}

\begin{lemma}[Triangle hub trades and their assembly]\label{bp:triangle-transformer-nearcomplete-trades}
\lean{Oberwolfach.triangular_transformer_common_hub_available,Oberwolfach.triangular_transformer_factors_from_near_complete_spokes,Oberwolfach.ordinary_host_resources_monotone,Oberwolfach.triangle_factor_pair_hub_trade,Oberwolfach.triangular_transformer_wheel_trades_from_near_complete_spokes,Oberwolfach.triangle_hub_trade_rims_disjoint,Oberwolfach.triangle_hub_trade_spokes_disjoint,Oberwolfach.triangle_hub_trades_choose_simultaneously,Oberwolfach.triangular_transformer_simultaneous_wheel_trades,Oberwolfach.triangle_selected_hub_trades_finish}
\leanok
For boundedly many old spoke-path sets, near-completeness and the displayed forbidden-set budget give pairwise disjoint fresh domains. The two triangle factors at each hub yield wheel decompositions in the original host and the signed hub-switch identity. Their rim edges lie in the corresponding domain and have no edge between two old vertices; resources at different hubs are disjoint. Given the separately constructed vertex-wheel families, their support and root-spoke exclusions, and their exact signed identity, the selected trades assemble into a resource transformer of bounded additional order. The proof still has to establish those vertex-wheel assembly inputs and all parameter substitutions in the manuscript's original caller.
\end{lemma}
\begin{proof}\uses{bp:triangle-transformer-fresh-spoke-factors}\leanok
Construct common-hub available sets, choose the factors on disjoint domains, and turn each pair into a hub trade. The physical support and spoke exclusions permit the exact resource-identity assembly.
\end{proof}

\begin{lemma}[Uniform bounded-root triangle transformer in resource form]\label{bp:triangle-transformer-uniform-resource}
\lean{Oberwolfach.triangle_vertex_wheels_and_hub_trades_complete,Oberwolfach.triangle_four_graph_neighbours,Oberwolfach.triangle_four_neighbours_two_hubs,Oberwolfach.triangle_common_hub_neighbours,Oberwolfach.triangle_transformer_dense_edge_realization,Oberwolfach.triangle_dense_transformer_from_explicit_budgets,Oberwolfach.triangle_budgets_from_uniform_root_bounds,Oberwolfach.triangle_active_hubs_le_spokes,Oberwolfach.triangle_hub_spoke_count_le,Oberwolfach.triangle_transformer_uniform_bounded_inputs,Oberwolfach.triangle_root_spokes_bounded,Oberwolfach.triangle_root_complexity_bound,Oberwolfach.triangle_transformer_uniform_bounded_vertices,Oberwolfach.triangle_transformer_positive_beta,Oberwolfach.triangle_transformer_finite_resource_uniform}
\leanok
Under the displayed graph, root-support and edge-map hypotheses, the triangle resource transformer exists with its vertices uniformly bounded. The spoke error is chosen before the root-size bound; the graph and hub thresholds come afterwards. The conclusion includes containment in the original host, the root-edge exclusions and avoidance of prescribed old row and hub sets. The subsequent finite-pair result translates the embedded root graphs and their edge-bijective map to this construction; its use in the manuscript's higher-level caller remains to be checked.
\end{lemma}
\begin{proof}\uses{bp:triangle-transformer-nearcomplete-trades}\leanok
Derive the four-root and two-hub neighbourhood estimates from minimum degree and near-completeness, construct the edge and vertex wheels, and assemble the hub switches. Bound the finite root complexities and choose the constants in the stated order.
\end{proof}

\begin{lemma}[Finite-pair triangle transformer]\label{bp:triangle-finite-pair-transformer}
\lean{Oberwolfach.triangle_finite_pair_transformer,Oberwolfach.triangle_finite_pair_extended_transformer}
\leanok
Under the stated minimum-degree, near-completeness, bounded-root and forbidden-set hypotheses, a pair of disjoint embedded root graphs with a part-preserving edge-bijective map has a bounded transformer in the original extended host. Both wheel decompositions and the root and forbidden-vertex exclusions are included. Verification of the manuscript's actual higher-level applications remains separate.
\end{lemma}
\begin{proof}\uses{bp:triangle-transformer-uniform-resource}\leanok
Translate the finite pair and its resource map into the uniform construction and recover the extended subgraph and both decompositions.
\end{proof}

\begin{lemma}[Successive triangle transformers]\label{bp:triangle-sequential-transformers}
\lean{Oberwolfach.ordinary_sequential_variable_budget_selection,Oberwolfach.ordinary_sequential_recursive_budget_selection,Oberwolfach.triangle_sequential_resource_transformers}
\leanok
A fixed finite sequence of triangle transformers can be chosen to avoid an initial resource set and all previously selected transformers. The graph-order thresholds are chosen using recursively bounded forbidden sets after the common spoke error has been fixed. The construction of the particular roots and their embeddings in the absorber caller remains to be checked.
\end{lemma}
\begin{proof}\uses{bp:triangle-transformer-uniform-resource}\leanok
Include all endpoints of the previously used resources in the next forbidden sets, and apply the uniform transformer with the corresponding recursively chosen budget.
\end{proof}

\begin{lemma}[Conditional transfer of a reserved marked reservoir]\label{bp:marked-reserved-state-transfer}
\lean{Oberwolfach.marked_reserved_rim_state_agreement_transfers,Oberwolfach.marked_encoded_fixed_reservoir_completion_from_states}
\leanok
Agreement of the raw state on each reserved physical rim edge preserves its encoded type and direction. Together with the stated inclusion of reserved typed spokes, level data and other hypotheses, this permits the fixed-reservoir completion in the later host. The preparation and first cover must still establish these agreement and inclusion hypotheses for the actual absorber.
\end{lemma}
\begin{proof}\leanok
Apply the rim encoding to each reserved edge and invoke the fixed-reservoir completion with the resulting typed containment.
\end{proof}

\begin{lemma}[Direct dense-host triangle absorber]\label{bp:triangle-dense-host-absorber}
\lean{Oberwolfach.finite_four_absorber_bounds,Oberwolfach.triangle_four_resource_transformer_absorber,Oberwolfach.triangle_actual_diagram_absorber,Oberwolfach.triangle_bounded_diagram_absorber,Oberwolfach.triangle_diagram_dense_geometry,Oberwolfach.triangle_dense_host_absorber,Oberwolfach.triangle_dense_host_absorber_graph,Oberwolfach.triangle_manuscript_absorber}
\leanok
For a bounded balanced, triangle-divisible root graph embedded in a sufficiently large host with minimum degree at least $(3/4+\beta)|V|$ and sufficiently near-complete spokes, the stated bounded forbidden sets can be avoided by a bounded extended absorber. Both wheel decompositions, host containment and root-resource avoidance are included. The error is chosen before the root and forbidden-set bounds, and the two part-size thresholds afterwards. This does not establish that the sampled hosts in the general proof meet these dense-host hypotheses.
\end{lemma}
\begin{proof}\uses{bp:triangle-sequential-transformers}\leanok
Construct the five finite diagram embeddings using the common-neighbourhood bounds implied by minimum degree and near-complete spokes. Select four disjoint transformers with recursively bounded forbidden sets, assemble the absorber and recover its extended graph.
\end{proof}

\begin{lemma}[Common marked subhost and protected inner pair]\label{bp:marked-common-subhost}
\lean{Oberwolfach.marked_rim_common_subhost_deletion,Oberwolfach.marked_spoke_common_subhost_deletion,Oberwolfach.marked_common_subhost_physical_deletions,Oberwolfach.marked_common_subhost_agrees_on_inner}
\leanok
Given an original sample and a later typed host, delete from the original physical rim edges on which their raw states differ and delete original spokes whose type is absent from the later host. The resulting typed host belongs to both hosts. Its deletion degrees are the stated realized counts, and if the two hosts agree on the protected inner pair then the common host agrees there as well. This deterministic construction does not assert a new sampling law or supply bounds for the actual preparation.
\end{lemma}
\begin{proof}\leanok
Construct the two physical deletion graphs from the changed states and from the missing typed spokes. Check directions, types and the permutation of marked heads, and use agreement on the inner pair.
\end{proof}

\begin{lemma}[Reserved resources under changed rims and label exclusions]\label{bp:marked-reservoir-change-labels}
\lean{Oberwolfach.marked_reserved_rim_change_graph_transfers,Oberwolfach.marked_encoded_fixed_reservoir_completion_from_change_graph,Oberwolfach.marked_reserved_spokes_from_label_exclusions,Oberwolfach.marked_encoded_fixed_reservoir_completion_from_label_exclusions}
\leanok
If the actual rim-change graph avoids the reserved physical edges and the displayed retained and unused-label conditions hold for the reserved spokes, the fixed-reservoir completion applies under its stated level, divisibility and loss hypotheses. The preparation and first cover must still prove those hypotheses for the actual construction; the common-subhost route requires additional absorber and remainder transfers.
\end{lemma}
\begin{proof}\uses{bp:marked-reserved-state-transfer}\leanok
Transfer the unchanged rim states with their types and directions, derive spoke inclusion from the label exclusions, and apply the fixed-reservoir completion.
\end{proof}

\begin{lemma}[Conditional completion after typed host changes]\label{bp:marked-typed-change-completion}
\lean{Oberwolfach.marked_encode_after_protected_rim_deletion,Oberwolfach.marked_common_subhost_transfer_to_later,Oberwolfach.marked_common_terminal_arc_agreement,Oberwolfach.marked_common_reference_completion,Oberwolfach.marked_reference_completion_after_host_changes,Oberwolfach.marked_rim_host_change_count,Oberwolfach.marked_spoke_host_change_counts,Oberwolfach.marked_reference_completion_from_change_graphs,Oberwolfach.marked_rim_typed_intersection_deletion,Oberwolfach.marked_common_terminal_arc_agreement_typed,Oberwolfach.marked_common_reference_typed_completion,Oberwolfach.marked_reference_completion_after_typed_host_changes,Oberwolfach.marked_removed_typed_rim_count_of_change_graph,Oberwolfach.marked_reference_completion_from_typed_change_graphs}
\leanok
Suppose the actual reference-level data and step properties hold, the original and later typed hosts agree on the protected inner pair, and the realized numbers of original typed physical rims and spokes missing from the later host satisfy the stated degree bounds. Suppose also that one fixed original-sample event supplies a terminal absorber in every physical deletion within those bounds. Then the later divisible host has a marked-wheel decomposition, even when it has additional rims or spokes. The absorber is chosen in the common typed subhost after the later host is fixed. This is a conditional completion theorem: the manuscript's actual reference levels, common event, change estimates, divisibility and parameter substitutions have not yet been established here.
\end{lemma}
\begin{proof}\uses{bp:marked-common-subhost}\leanok
Delete exactly the original typed rims and spokes missing from the later host; their degree bounds follow from the displayed realized counts. The common subhost agrees with the later host on the protected terminal pair. Choose its fixed absorber using the one uniform deletion property, transfer its resources and every later terminal remainder with types and endpoints intact, and apply the fixed-reference iteration.
\end{proof}

\begin{lemma}[Reference estimates after physical host changes]\label{bp:marked-reference-fields-changes}
\lean{Oberwolfach.marked_reference_common_fields_after_changes,Oberwolfach.marked_reference_weighted_load_after_changes}
\leanok
Assume the displayed original reference estimates, the interval-coordinate agreement, and the stated degree and numerical error bounds for the physical rim and spoke change graphs. Then the seven displayed reference statistics, cut, common-configuration, neighbourhood and spoke-agreement fields hold for the modified host. The fixed root-omitted weighted loads also satisfy their stated error bound on the original root domain, with the original weight function. These deterministic conclusions allow both additions and deletions, but do not prove that the actual preparation supplies the hypotheses.
\end{lemma}
\begin{proof}\uses{bp:marked-joint-physical-changes,bp:marked-common-subhost}\leanok
Compare typed incidences through the physical change graphs. Transfer cycles and root-omitted paths through the common subhost, use the original error scales, and bound the weighted load change with the same fixed weights.
\end{proof}

\begin{lemma}[Conditional completion using shared physical change graphs]\label{bp:marked-shared-change-completion}
\lean{Oberwolfach.marked_reference_completion_uniform_size_order,Oberwolfach.marked_reference_completion_from_shared_change_graphs}
\leanok
Under the stated reference-level, step, protected-pair, divisibility and uniform absorber hypotheses, the same physical change graphs used in the reference estimates yield the missing-resource bounds for the typed completion. The order threshold can be fixed before the level data. Establishing these hypotheses and their simultaneous application to the actual preparation remains open.
\end{lemma}
\begin{proof}\uses{bp:marked-typed-change-completion}\leanok
Bound each missing original typed resource by the corresponding physical change graph and apply the typed completion after fixing the threshold from the displayed positive loss parameters.
\end{proof}

\begin{lemma}[Cycle-weighted vertex and pair counts]\label{bp:cycle-sampling-means}
\lean{Oberwolfach.ordinary_cycle_fractional_vertex_weight,Oberwolfach.ordinary_cycle_two_vertex_count,Oberwolfach.embedded_cycle_support_range,Oberwolfach.ordinary_cycle_pair_weight_bound,Oberwolfach.ordinary_cycle_sample_vertex_mean,Oberwolfach.ordinary_cycle_sample_pair_mean_bound}
\leanok
For the stated fractional weights on actual ordinary cycle subgraphs, the total weight through a vertex is half its graph degree. Under the stated individual upper bound, the total weight through two distinct vertices is at most the displayed constant. The corresponding means for independent sampling with the supplied fixed probabilities follow. These estimates do not themselves establish the sampling margins needed for the factor caller.
\end{lemma}
\begin{proof}\leanok
Count vertex-edge incidences in each cycle, encode cycle subgraphs by their chosen vertex embeddings, and use the fixed two-coordinate counting bound. Then sum the independent retention probabilities.
\end{proof}

\begin{lemma}[A common independent-cycle sampling event]\label{bp:cycle-sampling-joint}
\lean{Oberwolfach.independent_cycle_count_expectation,Oberwolfach.independent_cycle_count_chernoff,Oberwolfach.independent_cycle_count_large_tail,Oberwolfach.independent_cycle_count_zero_mean,Oberwolfach.independent_cycle_sampling_joint_event,Oberwolfach.independent_cycle_sampling_joint_exists}
\leanok
One independent sample with the supplied, possibly unequal, fixed cycle probabilities satisfies all tests in the displayed finite families when the stated sum of failure bounds is less than one. Tests of zero mean vanish on the selected positive-probability outcome. The required quantitative estimates for the manuscript's actual cycle family and all factor callers remain to be proved.
\end{lemma}
\begin{proof}\leanok
Apply the fixed-test Chernoff bounds and a union bound to one product law; use nonnegativity for the zero-mean tests.
\end{proof}

\begin{lemma}[Conditional near-spanning packing of retained cycles]\label{bp:retained-cycle-packing}
\lean{Oberwolfach.cycle_hypergraph_near_perfect,Oberwolfach.retained_ordinary_cycle_vertices_card,Oberwolfach.retained_ordinary_cycle_degree,Oberwolfach.retained_ordinary_cycle_codegree,Oberwolfach.ordinary_sampled_cycle_packing}
\leanok
If the retained ordinary cycles satisfy the displayed hypergraph degree, codegree, order and parameter bounds, they contain a pairwise vertex-disjoint family covering the stated fraction of vertices. The hypergraph vertices here are vertices of the original graph, and its edges are cycle identifiers. Establishing those bounds for the actual sample and completing the packing with an absorber remain open.
\end{lemma}
\begin{proof}\uses{bp:uniform-weighted-matching}\leanok
Apply the weighted matching theorem to the cycle vertex sets, and count incidences using their common size and the near-regularity hypothesis.
\end{proof}

\begin{lemma}[A sampled cycle packing from explicit margins]\label{bp:cycle-sampling-packing-bridge}
\lean{Oberwolfach.ordinary_sampling_packing_bridge,Oberwolfach.cycle_sampling_union_error_small}
\leanok
The displayed fractional cycle weights, probability bounds, degree margins and common-event failure bound yield a positive-probability sample containing the stated near-spanning vertex-disjoint cycle packing. The pair-mean estimate is derived from the supplied individual weight bound. Separately, a polynomial number of tests with the stated quadratic and linear exponential tails has total failure probability less than one for sufficiently large order. The proof has not yet derived all the bridge's numerical assumptions from the actual manuscript parameters or supplied the absorber and exact-factor completion.
\end{lemma}
\begin{proof}\uses{bp:cycle-sampling-means,bp:cycle-sampling-joint,bp:retained-cycle-packing}\leanok
Apply the fixed-family common event to vertex and positive-mean pair tests and handle zero-mean pairs on the same positive-probability outcome. The stated margins give the sampled degree and codegree hypotheses for the retained-cycle matching theorem. Exponential domination controls the separate abstract union bound.
\end{proof}

\begin{lemma}[Cycle sampling from quantitative mean bounds]\label{bp:cycle-packing-means}
\lean{Oberwolfach.cycle_sampling_finite_union_margin,Oberwolfach.cycle_sampling_relative_tail_bound,Oberwolfach.cycle_sampling_large_tail_half,Oberwolfach.cycle_sampling_failure_from_means,Oberwolfach.ordinary_sampling_packing_from_means,Oberwolfach.ordinary_cycle_sampling_scaled_probability,Oberwolfach.cycle_sampling_integer_degree}
\leanok
For the stated fractional weights, quadratic lower bounds on vertex means, linear upper bounds on pair means and displayed numerical margins give a positive-probability sample with a near-spanning vertex-disjoint cycle packing. The failure-sum premise is derived from these mean bounds. The separate scale and integer-degree estimates exhibit possible parameter choices. The actual factor callers and the absorber completion for this sampling route remain unproved; the separately checked almost-regular factor theorem still cannot be applied to an arbitrary original caller below its threshold.
\end{lemma}
\begin{proof}\uses{bp:cycle-sampling-packing-bridge}\leanok
Bound the finite union of fixed vertex and pair tests by the displayed exponential estimates, treat zero-mean pairs on the same sample and apply the sampling-to-packing bridge.
\end{proof}

\begin{lemma}[Deep marked reference fields under host changes]\label{bp:marked-deep-fields-changes}
\lean{Oberwolfach.marked_local_configurations_rim_inclusion,Oberwolfach.marked_local_configurations_spoke_inclusion,Oberwolfach.marked_local_configurations_after_host_change,Oberwolfach.marked_local_spoke_loads_after_physical_changes,Oberwolfach.marked_reference_deep_fields_after_changes}
\leanok
Under the stated original local configuration and fixed-weight spoke-load estimates, physical change-degree bounds, type conditions and numerical margins, the two displayed deep reference fields hold for the same modified host. The original root domain, weights, normalization, types and omitted root spokes are retained. The actual preparation, sampling and level iteration must still provide the hypotheses and verify the remaining fields.
\end{lemma}
\begin{proof}\uses{bp:marked-common-subhost,bp:marked-joint-physical-changes}\leanok
Transfer the cycle and rooted-path configurations through the typed common subhost and apply the fixed-weight spoke-load stability estimate to the same physical change graphs.
\end{proof}

\begin{lemma}[Parent potential from the adaptive event after host changes]\label{bp:marked-parent-adaptive-changes}
\lean{Oberwolfach.marked_reference_parent_potential_after_changes,Oberwolfach.marked_reference_parent_from_adaptive_event}
\leanok
The stated zero-deletion case of the adaptive packing event gives the original parent potential estimate. Under the displayed fixed-weight bounds and physical-change margins, the same normalized weight and marginal-root condition give the parent estimate in the later host. The actual event specialization and its constants remain to be checked.
\end{lemma}
\begin{proof}\uses{bp:marked-parent-load-zero,bp:marked-reference-fields-changes}\leanok
Use zero deletion on the original outcome and apply fixed-weight typed potential stability on the original root domain.
\end{proof}

\begin{lemma}[Uniform completion with a prescribed inner-size cap]\label{bp:marked-uniform-xicap-completion}
\lean{Oberwolfach.marked_reference_level_increase_xi,Oberwolfach.marked_reference_step_from_larger_xi,Oberwolfach.marked_uniform_reference_completion_with_xi_cap}
\leanok
For any positive cap fixed at the displayed stage, the stated uniform completion supplies a positive inner-size parameter below the cap. Its order threshold is fixed before the later colour types, level sequence and its length. The generic level steps are obtained from the proved uniform level-step theorem. Compatible level data, the single robust absorber event, the physical change bounds and protected-pair agreement are still required for the actual application.
\end{lemma}
\begin{proof}\uses{bp:markeduniformreferencelevelstep,bp:marked-typed-change-completion}\leanok
Increasing the permitted inner-set size changes no underlying level objects or step conclusion. Reduce the selected size parameter to the cap, apply the uniform level-step theorem to the enlarged bound and use typed common-subhost completion.
\end{proof}

\begin{lemma}[Protected-pair agreement from physical changes]\label{bp:marked-inner-protection-from-changes}
\lean{Oberwolfach.marked_protected_pair_from_change_graphs,Oberwolfach.marked_protected_pair_avoided_by_packing,Oberwolfach.marked_protected_pair_from_star_avoiding_packing,Oberwolfach.marked_uniform_reference_completion_from_inner_avoiding_changes}
\leanok
If the physical change graphs avoid the inner rim and spoke resources, the two typed hosts agree on that pair; the permutation in the marked physical head is included. A packing avoiding the inner star contributes no change there when the stated prescribed changes also avoid the pair. Under the remaining displayed hypotheses, the uniform completion applies. The manuscript's actual preparation must still establish these change-graph and star-avoidance hypotheses.
\end{lemma}
\begin{proof}\uses{bp:markeduniformreferencelevelstep,bp:marked-typed-change-completion}\leanok
Translate the resource avoidance into physical graph avoidance, use invariance of the inner vertex set under the permutation and apply the typed completion with the uniform level steps.
\end{proof}

\begin{lemma}[Protection through the two first cover-down deletions]\label{bp:marked-two-stage-active-protection}
\lean{Oberwolfach.marked_two_stage_cover_protects_active_pair,Oberwolfach.marked_two_stage_cover_preserves_active_raw_pair,Oberwolfach.marked_two_stage_cover_preserves_reference_inner_raw_pair}
\leanok
The star-filtered first correction and the inner-resource-free deep packing preserve the restricted raw sample and both spoke relations on fixed subdomains of the active pair. The reference-domain specialization uses the stated vertex and hub image equalities; it restricts hubs to the designated inner hub set. The packing properties, compatibility of the actual level data, and the absorber's protected domain remain to be supplied at the application.
\end{lemma}
\begin{proof}\uses{bp:marked-inner-protection-from-changes}\leanok
The star-filtered physical resources and the first packing avoid the active pair. The deep packing avoids the larger inner pair. Their physical deletion graphs therefore avoid the fixed active domains, and restriction commutes with the successive deletions.
\end{proof}

\begin{lemma}[Fixed active-domain neighbourhood event]\label{bp:marked-active-domains-joint-event}
\lean{Oberwolfach.vortex_sdiff_scales,Oberwolfach.vortex_sdiff_size_bounds,Oberwolfach.marked_neighbourhood_active_vortex_event,Oberwolfach.marked_active_vortex_neighbourhood_with_existing_event}
\leanok
For fixed nested domains satisfying the displayed size and invariance conditions, the active differences retain the stated proportion and obey geometric decay. Their row and index estimates hold on one retained outcome for all later admissible deletions. This event can be combined with an existing event on the same law, with the stated sum of failure probabilities. The actual domains, margins and pre-exposure choice remain to be verified at the application.
\end{lemma}
\begin{proof}\leanok
Bound the sizes after deleting the deeper sets, apply the retained neighbourhood theorem to the fixed differences, and intersect the events on the same probability space.
\end{proof}

\begin{lemma}[Reference-chain active geometry]\label{bp:marked-active-chain-geometry}
\lean{Oberwolfach.marked_reference_next_inner_images,Oberwolfach.marked_reference_active_pair_scales,Oberwolfach.marked_reference_active_ambient_bounds,Oberwolfach.marked_reference_active_minimum_uniform_all_types,Oberwolfach.marked_reference_bounded_ambient_comm,Oberwolfach.marked_reference_bounded_next_inner_scales}
\leanok
The next reference level's inner sets map onto the current level's deeper sets. Thus the existing cut-size bounds give geometric decay, while the existing relative-size bounds control ambient-to-inner ratios. A size threshold may be chosen before the later colour types and level objects. A finite nonempty chain can be extended to all indices by repeating level zero. Constructing the actual chain and showing that its domains were fixed before sampling remain separate obligations.
\end{lemma}
\begin{proof}\uses{bp:marked-active-domains-joint-event}\leanok
Use the chain's image and preimage identities, the cut and size bounds, and the permutation-commuting embeddings.
\end{proof}

\begin{lemma}[Odd factor-caller threshold comparison]\label{bp:odd-factor-caller-threshold-gap}
\lean{Oberwolfach.odd_factor_caller_below_cycle_threshold}
\leanok
For odd cycle length and a sufficiently small positive caller margin, the displayed factor-density lower bound is below the ordinary cycle-decomposition threshold. This shows that the proposed above-threshold replacement cannot be applied from that displayed bound alone.
\end{lemma}
\begin{proof}\leanok
Compare the odd cycle-decomposition lower bound with the factor density and the specified margin.
\end{proof}

\begin{lemma}[Positive probability for the combined reference events]\label{bp:marked-joint-failure-budget}
\lean{Oberwolfach.joint_event_positive_mass_witness,Oberwolfach.marked_reference_joint_failure_positive_witness}
\leanok
If the stated fixed failure budget is below one, positive inverse-power exponents yield an order threshold for which a common event with the displayed probability bound has a positive-mass outcome. The application must still choose its constants and establish the bound for its actual fixed domains.
\end{lemma}
\begin{proof}\uses{bp:marked-active-domains-joint-event}\leanok
Combine the fixed errors and the two inverse-power terms, then use their convergence to zero and positivity of the resulting probability.
\end{proof}

\begin{lemma}[Fixed-vortex event and degree inputs]\label{bp:fixed-vortex-joint-witness}
\lean{Oberwolfach.marked_fixed_vortex_active_joint_positive_witness,Oberwolfach.marked_vortex_reference_degree_inputs,Oberwolfach.marked_vortex_reference_scale_inputs}
\leanok
The previously selected paired vortex and its preserving permutation supply the active domains and the stated geometric, degree, density-centre and scale inputs. Given the displayed baseline-event probability bound and a sufficiently small fixed failure budget, one positive-mass outcome satisfies both the baseline event and the active-domain estimates for all admissible later deletions. Specializing the baseline event, constructing the reference-level data and protecting the absorber remain separate obligations.
\end{lemma}
\begin{proof}\uses{bp:marked-active-domains-joint-event,bp:marked-active-chain-geometry,bp:marked-joint-failure-budget}\leanok
Use the actual vortex inclusions and its independently rounded vertex and hub sizes, transfer neighbourhood counts for all parent roots, and apply the common-event and positive-probability bounds on the original retained law.
\end{proof}

\begin{lemma}[Edge-disjoint triangles can intersect]\label{bp:edge-matching-not-vertex-packing}
\lean{Oberwolfach.edge_disjoint_triangles_can_share_a_vertex}
\leanok
Two explicitly listed triangles may have disjoint edge sets while sharing a vertex. Thus a matching of cycle edge-sets alone does not ensure a vertex-disjoint cycle packing.
\end{lemma}
\begin{proof}\leanok
Check the two finite edge and vertex sets.
\end{proof}

\begin{lemma}[Fixed-vortex successor weights]\label{bp:fixed-vortex-weights}
\lean{Oberwolfach.marked_fixed_vortex_successor_weights_with_fixed_parent}
\leanok
On the fixed vortex subtypes, the successor-weight family can be chosen before the sampling densities while retaining the supplied parent weight exactly, under the displayed graph and parameter conditions.
\end{lemma}
\begin{proof}\uses{bp:fixed-vortex-joint-witness}\leanok
Apply the successor-weight family theorem to the vortex degree inputs and replace the level-zero weight by the supplied parent weight.
\end{proof}

\begin{lemma}[One explicit failure budget]\label{bp:explicit-reference-failure-budget}
\lean{Oberwolfach.marked_reference_fixed_failure_one_hundredth,Oberwolfach.marked_reference_explicit_error_positive}
\leanok
Taking the eight fixed error parameters to be one hundredth gives total fixed error one tenth. With both power exponents equal to one, the displayed lower probability bound is positive from order 16 onward. Other parameter and event hypotheses remain to be verified for the application.
\end{lemma}
\begin{proof}\uses{bp:marked-joint-failure-budget}\leanok
Evaluate the fixed error sum and bound the remaining inverse-order term.
\end{proof}

\begin{lemma}[Deterministic fixed-vortex reference inputs]\label{bp:fixed-vortex-reference-inputs}
\lean{Oberwolfach.marked_fixed_vortex_deterministic_reference_inputs}
\leanok
For the preselected paired vortex, the stated successor-weight family, size and scale bounds, parent and inner degree estimates, and density-centre identities hold simultaneously. The weights are chosen before the later density parameters; the order threshold follows the prescribed event minimum. The remaining assumptions and all reference-level constructions still require verification.
\end{lemma}
\begin{proof}\uses{bp:fixed-vortex-joint-witness,bp:fixed-vortex-weights}\leanok
Combine the pre-exposure vortex geometry, parent and inner degree estimates, the exact density identities and the successor-weight theorem, preserving their parameter order.
\end{proof}

\begin{lemma}[Vortex mass and minimum domains]\label{bp:fixed-vortex-mass}
\lean{Oberwolfach.marked_vortex_retention_mass,Oberwolfach.marked_vortex_event_minimum_domains}
\leanok
The graph degree and hub-density lower bounds imply the displayed retained-sampling mass inequality once the parameter bound is met. The active-set size margins give the required minimum cardinalities of the inner and parent vertex and hub domains.
\end{lemma}
\begin{proof}\uses{bp:fixed-vortex-reference-inputs}\leanok
Multiply the hub-density and degree inequalities, and compare each inner set with its parent domain.
\end{proof}

\begin{lemma}[Fixed-vortex baseline geometry]\label{bp:fixed-vortex-baseline-geometry}
\lean{Oberwolfach.marked_fixed_vortex_baseline_geometry}
\leanok
The preselected paired vortex supplies the stated embedding compatibility, minimum domain sizes, parent hub-density ratios, scale and inner-set shrinkage, nestedness and permutation invariance simultaneously. Instantiating the baseline sampling event requires the remaining hypotheses and parameter choices.
\end{lemma}
\begin{proof}\uses{bp:fixed-vortex-joint-witness,bp:fixed-vortex-mass}\leanok
Combine the fixed-vortex size, permutation and subtype estimates before sampling.
\end{proof}

\begin{lemma}[Reference event on the fixed vortex]\label{bp:fixed-vortex-reference-event}
\lean{Oberwolfach.marked_fixed_vortex_reference_event}
\leanok
The reference-statistics, cut and configuration event holds on the actual fixed vortex, under the displayed density and parameter conditions, uniformly over later admissible physical deletions. The other components of the baseline event, including its absorber and load conditions, remain to be specialized and combined.
\end{lemma}
\begin{proof}\uses{bp:fixed-vortex-baseline-geometry,bp:fixed-vortex-reference-inputs}\leanok
Apply the reference-vortex event using the fixed domain embeddings and the original graph's vortex degree estimates.
\end{proof}

\begin{lemma}[A concrete common vortex outcome]\label{bp:fixed-vortex-reference-active-witness}
\lean{Oberwolfach.marked_fixed_vortex_reference_and_active_event,Oberwolfach.marked_fixed_vortex_reference_and_active_witness}
\leanok
Under the stated density, separation and parameter conditions, a single positive-mass retained outcome on the preselected vortex satisfies the reference-configuration, statistics and cut estimates together with the active-domain neighbourhood estimates, uniformly over their admissible later deletions. The balance and weighted-load events remain to be joined to this outcome; a further result below also joins a deletion-robust absorber event.
\end{lemma}
\begin{proof}\uses{bp:fixed-vortex-reference-event,bp:fixed-vortex-joint-witness,bp:explicit-reference-failure-budget,bp:fixed-vortex-mass}\leanok
Apply the two concrete events on the same retained law, derive the hub-mass condition from the original graph and choose explicit positive failure budgets.
\end{proof}

\begin{lemma}[Robust absorber event on the regular host]\label{bp:regular-host-robust-absorber-event}
\lean{Oberwolfach.marked_regular_host_robust_absorber_with_existing_event}
\leanok
The deletion-robust sampled absorber event specializes to the original regular host and may be intersected with an existing event on the same retained law. It remains uniform over all later deletions satisfying the stated bounds. The actual deletion budgets and level compatibility remain to be checked.
\end{lemma}
\begin{proof}\uses{bp:fixed-vortex-mass}\leanok
Derive the degree and hub-size inputs from regularity and the density bounds, then intersect the two events without assuming independence.
\end{proof}

\begin{lemma}[A common outcome with a robust absorber]\label{bp:fixed-vortex-reference-active-absorber}
\lean{Oberwolfach.marked_fixed_vortex_reference_active_absorber_event,Oberwolfach.marked_fixed_vortex_reference_active_absorber_witness,Oberwolfach.marked_fixed_vortex_reference_active_absorber_event_separate_thresholds,Oberwolfach.eventually_marked_reference_absorber_failure_below,Oberwolfach.marked_fixed_vortex_reference_active_absorber_probability_uniform,Oberwolfach.marked_fixed_vortex_reference_active_absorber_probability_separate_thresholds,Oberwolfach.long_assigned_fixed_vortex_reference_probability_separate}
\leanok
Under the displayed regular-host, vortex and density hypotheses, one retained outcome of positive mass satisfies the reference and active-domain estimates and the deletion-robust absorber property simultaneously. The latter holds for every later pair of physical deletion graphs whose global degrees meet its bound. A separate-threshold version chooses the vortex terminal-size threshold before the absorber size and the host-order threshold after that size; its probability can be bounded below by any fixed budget exceeding the stated constant failure terms. This probability estimate has also been specialized to the assigned two-colour sampling law, under the stated regular-host and vortex hypotheses. This does not yet establish the balance and weighted-load events or the degree bounds for the actual cover-down.
\end{lemma}
\begin{proof}\uses{bp:fixed-vortex-reference-active-witness,bp:regular-host-robust-absorber-event,bp:explicit-reference-failure-budget}\leanok
Intersect the events on their common sample space and choose fixed error budgets so that their union bound is positive.
\end{proof}

\begin{lemma}[Two successive physical deletions]\label{bp:robust-absorber-sequential-deletion}
\lean{Oberwolfach.marked_robust_absorber_after_two_physical_deletions}
\leanok
The deletion-robust absorber property holds after two successive physical deletions if the sum of their global degree bounds is at most the robustness margin. Verifying these bounds for the actual first cover-down remains an obligation.
\end{lemma}
\begin{proof}\uses{bp:marked-protected-active-deletion}\leanok
Successive physical deletions equal deletion by the union of the two graphs, whose degree is bounded by the sum of their degrees.
\end{proof}

\begin{lemma}[Outside degrees do not bound all degrees]\label{bp:outside-degree-obstruction}
\lean{Oberwolfach.protected_pair_outside_degree_does_not_control_global}
\leanok
A four-vertex star has no edge inside its centre, degree one at every vertex outside the centre, and degree three at the centre. Thus bounds outside a protected set and the absence of internal edges do not by themselves give global degree bounds. This example makes no assertion about the manuscript's particular packing.
\end{lemma}
\begin{proof}\leanok Check the six ordered edges.\end{proof}

\begin{lemma}[A fixed absorber before the later host]\label{bp:preselected-absorber-transfer}
\lean{Oberwolfach.marked_reserved_absorber_universal_host_transfer}
\leanok
A reserved absorber can be selected once in the original typed host, before the later host and terminal remainder are chosen. It remains an absorber for every divisible terminal remainder in any later host that preserves its actual typed resources and agrees on the prescribed terminal pair, under the displayed injectivity and size conditions. The actual preparation must still be shown to satisfy these preservation conditions.
\end{lemma}
\begin{proof}\uses{bp:marked-finite-remainder-host-transfer}\leanok
Select the absorber using all original terminal arcs and spokes. Transfer each later remainder and its absorbing decomposition using the specified type and resource correspondences.
\end{proof}

\begin{lemma}[Iteration with a fixed absorber]\label{bp:preselected-absorber-iteration}
\lean{Oberwolfach.marked_reference_iteration_with_fixed_absorber,Oberwolfach.marked_ambient_terminal_representation_absorbed,Oberwolfach.marked_reference_iteration_with_typed_fixed_absorber}
\leanok
Given the displayed reference chain, its level-step properties and one fixed absorber for every admissible terminal remainder, the finite iteration and the absorber decompose the union of the reserved resources and the remaining host. This does not construct the reference chain or establish preservation of the reserved resources.
\end{lemma}
\begin{proof}\uses{bp:markedreferencefiniteiteration,bp:markedreferenceterminaltransfer}\leanok
Apply the finite iteration, transfer its terminal remainder to the ambient host and use the fixed absorber.
\end{proof}

\begin{lemma}[Typed reserved resources after deletion]\label{bp:reserved-resources-two-deletions}
\lean{Oberwolfach.marked_reserved_resources_survive_physical_deletion,Oberwolfach.marked_reserved_resources_survive_two_deletions}
\leanok
If one or two physical deletion graphs avoid every rim pair and spoke of the reserved resource set, the original ordinary and marked arc directions and the spoke types of those resources persist. No global deletion-degree bound is needed. Avoidance by the actual cover-down remains to be established.
\end{lemma}
\begin{proof}\uses{bp:marked-protected-active-deletion}\leanok
Each surviving physical coordinate retains its original state; successive deletions equal deletion by the union.
\end{proof}

\begin{lemma}[Correction avoiding reserved rim pairs]\label{bp:correction-reserved-pair-count}
\lean{Oberwolfach.marked_correction_forbidden_incoming_bound,Oberwolfach.marked_correction_with_reserved_pair_count}
\leanok
At a fixed head, a physical rim pair determines the directed arc. Excluding a set of reserved pairs therefore removes at most its cardinality from each incoming candidate set. Under the displayed baseline, loss and margin inequalities, the marked-indegree correction can be chosen to preserve all those pairs. The sample-specific candidate and loss estimates still require verification.
\end{lemma}
\begin{proof}\uses{bp:marked-protected-candidate-correction}\leanok
Inject the excluded incoming arcs into the set of reserved pairs, and apply the protected correction with this bound.
\end{proof}

\begin{lemma}[A preselected absorber under adaptive deletions]\label{bp:adaptive-preselected-absorber}
\lean{Oberwolfach.marked_preselected_absorber_after_adaptive_deletions,Oberwolfach.marked_robust_sample_selects_absorber_for_adaptive_deletions}
\leanok
A sampled host with the displayed reserved-absorber property admits one fixed absorber before any later physical deletion or terminal remainder is chosen. The absorber works for every later deletion that avoids its resources and the terminal spoke pair and for every divisible terminal remainder of the resulting host. The robust sampled event supplies the initial property by taking empty deletions. The actual later deletion must still be shown to meet these avoidance conditions.
\end{lemma}
\begin{proof}\uses{bp:preselected-absorber-transfer,bp:reserved-resources-two-deletions,bp:fixed-vortex-reference-active-absorber}\leanok
Preserve the absorber's typed arcs and spokes, embed each later terminal arc into the original sample, and transfer the absorption property.
\end{proof}

\begin{lemma}[Reserved-resource graphs and prescribed packing]\label{bp:reserved-resource-prescribed-packing}
\lean{Oberwolfach.marked_reserved_rim_pairs_membership,Oberwolfach.marked_reserved_rim_pairs_card,Oberwolfach.marked_reserved_correction_forbidden_bound,Oberwolfach.marked_correction_protects_reserved_resources,Oberwolfach.marked_reserved_spoke_pairs_card,Oberwolfach.marked_reserved_resource_graphs,Oberwolfach.marked_reserved_graphs_avoid_internal_resources,Oberwolfach.marked_reserved_graphs_preserve_prescribed_resources,Oberwolfach.marked_packing_avoiding_reserved_graphs,Oberwolfach.marked_sample_prescribed_packing_avoiding_reserved_resources}
\leanok
The physical rim and spoke graphs of a reserved set have maximum degree at most its number of resources. If the reserved set avoids the prescribed inner pair, then prescribed arcs and spokes on that pair remain available when these graphs are removed. Under the displayed sampling, size, and degree hypotheses, the prescribed wheel packing can be chosen disjoint from the reserved set. The marked-indegree correction also preserves its rim pairs under the displayed candidate margins. This does not yet establish avoidance for the later deep packing or construct the full preparation.
\end{lemma}
\begin{proof}\uses{bp:correction-reserved-pair-count,bp:marked-protected-candidate-correction}\leanok
Construct separate physical rim and spoke graphs, apply their degree bounds to the existing packing theorem, and translate graph disjointness back to physical resources.
\end{proof}

\begin{lemma}[Deep cover while reserving physical resources]\label{bp:reserved-resource-deep-cover}
\lean{Oberwolfach.marked_deleted_sample_packing_avoids_reserved_resources,Oberwolfach.marked_sample_deep_cover_preserving_reserved_resources,Oberwolfach.marked_disjoint_packing_transfers_to_reserved_deletion}
\leanok
A packing disjoint from the reserved set transfers to the sample with the reserved rim and spoke graphs deleted, with exactly the same physical resources. Under the displayed sample, degree and parameter conditions, the deep-cover theorem applies after deleting those graphs together with earlier bounded deletions. Its packing avoids the reserved set, and its coverage conclusion concerns the reduced host. The caller must still verify its balance and divisibility inputs and connect the two stages to the reference iteration.
\end{lemma}
\begin{proof}\uses{bp:reserved-resource-prescribed-packing,bp:reserved-resources-two-deletions}\leanok
Transfer every typed wheel to the reduced host. Divide the deletion budget between the earlier changes and the reserved set, then apply the deep-cover theorem to the union of the physical deletion graphs.
\end{proof}

\begin{lemma}[Divisibility after reserving the absorber]\label{bp:reserved-both-remainders}
\lean{Oberwolfach.marked_reserved_graphs_of_packing,Oberwolfach.marked_reserved_reduced_host_divisible,Oberwolfach.marked_reserved_deletion_protects_inner_pair,Oberwolfach.marked_reserved_deletion_preserves_inner_arcs,Oberwolfach.marked_reserved_inner_remainder_transfer,Oberwolfach.marked_reserved_canonical_inner_remainder,Oberwolfach.marked_reserved_canonical_inner_spokes,Oberwolfach.marked_reserved_canonical_full_inner_remainder,Oberwolfach.marked_reserved_outer_and_inner_remainders_divisible}
\leanok
Removing a wheel-decomposable reserved set from a divisible host preserves global divisibility. If the reserved set avoids the invariant inner pair, the canonical inner remainder after prescribed arc and spoke exclusions remains divisible and has the same physical resources. The arc transfer is the actual deletion embedding, preserving types and multiplicities. The remaining task is to instantiate these conclusions together with the sample estimates in the two-stage cover-down.
\end{lemma}
\begin{proof}\uses{bp:reserved-resource-prescribed-packing,bp:marked-finite-remainder-host-transfer}\leanok
Identify the physical graphs of a decomposition of the reserved set, subtract its wheels, and transfer the exact inner arc and spoke sets across the deletion.
\end{proof}

\begin{lemma}[Protection through both cover-down deletions]\label{bp:reserved-two-stage-protection}
\lean{Oberwolfach.marked_first_deletion_preserves_reserved_resources,Oberwolfach.marked_two_stage_deletions_preserve_reserved_resources,Oberwolfach.marked_reserved_deleted_packing_lift}
\leanok
Suppose the prescribed arcs and spokes lie in a protected inner pair avoided by a fixed reserved set, and both successive packings avoid that set. Then their physical deletion graphs leave every reserved rim direction and spoke type present. A packing in a host with the reserved graphs removed lifts to the preceding host with the same physical resources and still avoids the reserve. The remaining caller must connect the chosen packings, divisibility conditions and sample estimates with these hypotheses.
\end{lemma}
\begin{proof}\uses{bp:reserved-resource-deep-cover,bp:reserved-resource-prescribed-packing,bp:reserved-resources-two-deletions}\leanok
The prescribed inner resources and the two packings are disjoint from the reserve. Check physical rim pairs and spokes separately, and lift the second packing without changing its resources.
\end{proof}

\begin{lemma}[Combined deletions and protection of the active pair]\label{bp:reserved-union-cover-protection}
\lean{Oberwolfach.marked_reserved_two_stage_lift_and_protection,Oberwolfach.marked_two_stage_cover_protects_reserve_and_active_pair,Oberwolfach.marked_reserved_lift_protects_active_pair_and_absorber,Oberwolfach.marked_reserved_union_sequential_identity,Oberwolfach.marked_equal_sample_wheels_rebase,Oberwolfach.marked_reserved_union_packing_to_sequential,Oberwolfach.marked_reserved_union_cover_protects_active_pair_and_absorber}
\leanok
Deleting the union of the first deletion and the reserved graphs gives the same marked sample as deleting them successively. A deep packing selected in that host lifts with the same physical resources; under the stated disjointness conditions, both the active pair and all typed resources of the reserved absorber survive the two deletions.
\end{lemma}
\begin{proof}\uses{bp:reserved-two-stage-protection,bp:reserved-resource-deep-cover}\leanok
Identify the two samples, rebase the typed wheels, and use the separate resource and active-pair protection statements.
\end{proof}

\begin{lemma}[Lifting prescribed inner resources]\label{bp:reserved-inner-resource-lift}
\lean{Oberwolfach.marked_reserved_inner_arc_set_lift,Oberwolfach.marked_reserved_inner_resources_lift}
\leanok
An arbitrary set of typed arcs inside the invariant inner pair lifts through the actual deletion embedding with the same cardinality and physical rim resources. The same holds for the mixed set obtained by adjoining any specified set of spokes; no divisibility of the omitted arcs is required.
\end{lemma}
\begin{proof}\uses{bp:reserved-both-remainders}\leanok
The reserved set avoids the inner pair, so every prescribed inner arc survives with its endpoints and type. Taking preimages under the injective deletion embedding gives the claimed identities.
\end{proof}

\begin{lemma}[Exclusions and inner resources after transfer]\label{bp:reserved-transfer-exactness}
\lean{Oberwolfach.marked_packing_exclusions_transfer_on_resources,Oberwolfach.marked_reserved_packing_transfer_preserves_exclusions,Oberwolfach.marked_reserved_packing_transfer_preserves_inner_exactness}
\leanok
A packing disjoint from the reserved resources transfers to the reduced host with its exact physical resources and disjointness intact. Prescribed arc and spoke exclusions, and the equality specifying its inner resources, transfer with the prescribed arcs lifted by the deletion embedding.
\end{lemma}
\begin{proof}\uses{bp:reserved-resource-deep-cover,bp:reserved-inner-resource-lift}\leanok
Use injectivity of the original physical arc map to transfer arc exclusions, and use equality of physical resource sets for spokes and the inner-resource identity.
\end{proof}

\begin{lemma}[Canonical exclusions in the reduced sample]\label{bp:canonical-star-preimage}
\lean{Oberwolfach.marked_sample_canonical_star_arcs_preimage,Oberwolfach.marked_packing_spokes_subset_host,Oberwolfach.marked_packing_disjoint_exclusions_give_host_subsets}
\leanok
The canonical set of inner arcs meeting the smaller vertex set is the preimage of the corresponding set in the original sample, even after arbitrary prescribed inner arc omissions. Every spoke of a marked packing belongs to a typed spoke relation of its host, so disjointness from prescribed arc and spoke sets gives the corresponding host-subset conditions.
\end{lemma}
\begin{proof}\uses{bp:reserved-transfer-exactness}\leanok
The deletion embedding preserves both endpoints and the arc type; the spoke statement follows directly from the definition of a marked wheel.
\end{proof}

\begin{lemma}[First packing in the canonical reduced host]\label{bp:canonical-packing-transfer}
\lean{Oberwolfach.marked_reserved_canonical_packing_transfer}
\leanok
If a first packing satisfies the canonical exclusion and inner-resource conditions and is disjoint from the reserved set, it transfers to the absorber-deleted host with the same physical resources. The transferred packing satisfies those conditions for the canonical reduced-host sets, with the prescribed omitted arcs lifted by the deletion embedding.
\end{lemma}
\begin{proof}\uses{bp:reserved-transfer-exactness,bp:canonical-star-preimage,bp:reserved-both-remainders}\leanok
Transfer the packing and use the equality of the canonical arc preimages and inner spoke sets.
\end{proof}

\begin{lemma}[Validity of inner omissions after reservation]\label{bp:reserved-inner-omissions}
\lean{Oberwolfach.marked_reserved_inner_omissions_remain_valid}\leanok
Prescribed inner arcs lift with both endpoints still in the inner vertex set. Prescribed inner spokes that were present in the original host remain present in the reduced host when the reserved set avoids the inner pair.
\end{lemma}
\begin{proof}\uses{bp:reserved-both-remainders}\leanok
Use the deletion embedding for arcs and the fact that the reserved spoke graph has no edge inside the inner pair.
\end{proof}

\begin{lemma}[Prescribed packing and a balanced reduced host]\label{bp:reserved-prescribed-balanced-host}
\lean{Oberwolfach.marked_inner_exactness_implies_star_exclusions,Oberwolfach.marked_packing_resource_equality_degree_counts,Oberwolfach.marked_reserved_canonical_deep_inputs_transfer,Oberwolfach.marked_reserved_prescribed_deep_host_balanced,Oberwolfach.marked_sample_prescribed_star_packing_avoiding_reserved_resources,Oberwolfach.marked_sample_prescribed_star_to_balanced_deep_host}
\leanok
Under the displayed sampled-host, size and divisibility hypotheses, a first packing can be chosen with the prescribed inner-star resources while avoiding the reserved absorber. Its exact inner-resource identity implies the canonical star exclusions. The packing transfers to the absorber-deleted host with the same rim and spoke degree counts, and the resulting deep host has the stated balance and cleared-neighbourhood properties. The two representations of its first-deletion graphs agree.
\end{lemma}
\begin{proof}\uses{bp:canonical-packing-transfer,bp:reserved-inner-omissions,bp:reserved-resource-prescribed-packing,bp:canonical-deletion-graphs}\leanok
Apply the prescribed packing theorem on the given sample. Transfer its resources and omissions, subtract the decomposable reserved set from the divisible host, and apply the deterministic deep-host balance lemma.
\end{proof}

\begin{lemma}[Equality of canonical deletion graphs]\label{bp:canonical-deletion-graphs}
\lean{Oberwolfach.marked_finite_resources_union_packing,Oberwolfach.marked_finite_resources_determine_projections,Oberwolfach.marked_reserved_canonical_deletion_resources_equal,Oberwolfach.marked_reserved_canonical_deletion_graphs_equal}
\leanok
When the reserved set avoids the inner pair and the transferred packing has the same physical resources, the canonical first-deletion rim and spoke graphs are identical in the original and absorber-deleted representations.
\end{lemma}
\begin{proof}\uses{bp:canonical-star-preimage,bp:reserved-inner-resource-lift,bp:canonical-packing-transfer}\leanok
Lift the canonical inner star exactly and use physical-resource equality to recover the rim and spoke graphs separately.
\end{proof}

\begin{lemma}[Two sampled cover-down stages with a reserve]\label{bp:two-stage-sampled-reserve}
\lean{Oberwolfach.marked_reserved_union_reverse_sequential_identity,Oberwolfach.marked_reserved_canonical_combined_deletion_identity,Oberwolfach.marked_deep_host_balance_transport,Oberwolfach.marked_reserved_canonical_combined_deep_balance,Oberwolfach.marked_deep_host_enumerate_roots,Oberwolfach.marked_sample_deep_cover_from_balanced_reserved_host,Oberwolfach.marked_two_stage_sample_cover_with_reserve}
\leanok
For \(r\ge 2\), under the theorem's displayed sample, density, size, weight, divisibility and deletion-budget hypotheses, the prescribed first packing and a deep packing can be selected on the same marked sample. Both avoid the fixed reserved set and their physical resources are disjoint. The first packing has the exact prescribed inner-star resources and the stated degree bounds; the deep packing has the stated degree, root, inner-avoidance and coverage conclusions on the host after deleting the first packing and reserve. The constants and sample conditions are explicit in the Lean statement.
\end{lemma}
\begin{proof}\uses{bp:reserved-prescribed-balanced-host,bp:reserved-union-cover-protection,bp:canonical-deletion-graphs,bp:reserved-resource-deep-cover}\leanok
Use the first sampled packing and its deterministic transfer to obtain the balanced combined-deletion host. Enumerate its small root sets and apply the sampled deep-cover theorem on the same outcome. The physical deletion graphs show that the resulting packing avoids both the reserve and the first packing.
\end{proof}

\begin{lemma}[Common sample and chosen cover-down scales]\label{bp:two-stage-common-scales}
\lean{Oberwolfach.marked_rows_with_spokes_tolerance_mono,Oberwolfach.marked_indices_with_spokes_tolerance_mono,Oberwolfach.marked_two_stage_sample_cover_common_tolerance,Oberwolfach.marked_two_stage_density_lower_bounds,Oberwolfach.marked_two_stage_sample_cover_canonical_densities,Oberwolfach.marked_two_stage_choose_inner_and_packing_scales,Oberwolfach.marked_two_stage_sample_cover_selected_scales,Oberwolfach.marked_two_stage_sample_cover_compact_caller}
\leanok
The two-stage theorem can use one row and index estimate at the smaller of the required tolerances. Its first-stage density lower bound follows from the displayed sample densities. After both stages' constants are fixed, positive inner and packing scales can be chosen so that the three deletion-degree budgets hold. The compact caller retains the other sample, weighted-load and divisibility hypotheses.
\end{lemma}
\begin{proof}\uses{bp:two-stage-sampled-reserve}\leanok
Relax the common estimates to each stage's tolerance and choose the scales before the sample and subsequent packings.
\end{proof}

\begin{lemma}[Parent-level inputs for the two-stage theorem]\label{bp:reference-two-stage-inputs}
\lean{Oberwolfach.marked_local_configurations_with_spokes_mono,Oberwolfach.marked_local_spoke_loads_tolerance_mono,Oberwolfach.marked_reference_level_two_stage_sample_inputs,Oberwolfach.marked_reference_choose_two_stage_configuration_parameter,Oberwolfach.marked_reference_choose_two_stage_error}
\leanok
The existing parent-level data imply the two-stage theorem's row and index estimates, local configuration counts, weighted-spoke load estimate, and type conditions when the stated tolerance inequalities hold. The configuration parameter can be chosen before the cover-down constants and the displayed error parameters afterwards. This does not construct the parent-level data or establish the common event for the actual vortex.
\end{lemma}
\begin{proof}\uses{bp:two-stage-common-scales}\leanok
Use the recorded fixed-domain estimates and weight function, the exact density normalization, and monotonicity of the error and lower-bound inequalities.
\end{proof}

\begin{lemma}[Uniform two-stage cover after protected deletions]\label{bp:reference-two-stage-protected}
\lean{Oberwolfach.marked_reference_two_stage_inputs_after_changes,Oberwolfach.marked_reference_two_stage_inputs_after_deletions,Oberwolfach.marked_reference_choose_uniform_change_margins,Oberwolfach.marked_reference_scale_neighbourhood_margins,Oberwolfach.marked_reference_weighted_change_coefficient_nonnegative,Oberwolfach.marked_reference_two_stage_uniform_deletion_inputs,Oberwolfach.marked_two_stage_complete_density_bounds,Oberwolfach.marked_reference_two_stage_uniform_canonical_deletion_inputs,Oberwolfach.marked_reference_two_stage_sample_cover_after_deletions,Oberwolfach.marked_reserved_decomposition_after_avoiding_deletions,Oberwolfach.marked_reference_two_stage_cover_with_protected_reserve,Oberwolfach.marked_supported_reserve_avoids_inner_subset}
\leanok
For a parent-level record with its original sampled outcome and fixed weight, the sample and weighted-load hypotheses for both cover-down stages survive every later physical deletion satisfying the stated degree bound. The required tolerances can be chosen in the displayed order and below arbitrary positive caps. A wheel-decomposable reserve stays decomposable when the deletions avoid its physical resources, and the two-stage theorem applies with that original decomposition. The reserve's support condition implies avoidance of an invariant inner pair contained in the avoided pair.
\end{lemma}
\begin{proof}\uses{bp:reference-two-stage-inputs,bp:two-stage-common-scales,bp:reserved-two-stage-protection}\leanok
Bound the effect of all allowed physical changes on the same fixed counts and weights. Apply the two-stage theorem to the resulting host, and transfer the reserved wheels across the deletions without changing their physical resources.
\end{proof}

\begin{lemma}[A preselected absorber through two cover-down stages]\label{bp:reference-preselected-two-stage}
\lean{Oberwolfach.marked_reference_two_stage_cover_with_supported_reserve,Oberwolfach.marked_preselected_absorber_decomposes_and_avoids_inner_after_deletions,Oberwolfach.marked_two_stage_deletion_spoke_graphs_avoid_active_pair,Oberwolfach.marked_two_stage_terminal_deletion_graphs,Oberwolfach.marked_two_stage_preserves_terminal_absorption,Oberwolfach.marked_reference_preselected_absorber_and_two_stage_cover,Oberwolfach.marked_packing_union_with_deleted_host_decomposition}
\leanok
Given the stated parent-level data and terminal absorber property, one supported absorber can be selected before subsequent physical deletions and both sampled cover-down packings. Its original decomposition and absorption property transfer through deletions that avoid its resources. The two packings preserve the active pair and the absorber; when the terminal pair is contained in the active pair and the previous deletion graph avoids its spokes, the same absorber works for every divisible terminal remainder. A packing and a decomposition in a host with its physical resources deleted lift and combine into a decomposition of their union.
\end{lemma}
\begin{proof}\uses{bp:reference-two-stage-protected,bp:reserved-union-cover-protection}\leanok
Select the absorber once and apply the uniform two-stage result. Lift the deep packing without changing its physical resources, check the physical deletion graphs on the reserve and active terminal pair, and invoke the absorber's universal property. Join disjoint resource decompositions after lifting the wheels.
\end{proof}

\begin{lemma}[Terminal absorption after two cover-down stages]\label{bp:two-stage-exact-absorption}
\lean{Oberwolfach.marked_two_stage_exact_decomposition_direct,Oberwolfach.bounded_card_fits_two_linear_reserves,Oberwolfach.marked_reference_level_lower_threshold}
\leanok
Under the displayed containment, avoidance and divisibility hypotheses, the two packings, the absorber and a terminal remainder have a wheel decomposition of their union in the original host. A fixed absorber-size bound fits inside two positive linear reserves for sufficiently large order. Identifying this union with the full host remains a separate obligation.
\end{lemma}
\begin{proof}\uses{bp:reference-preselected-two-stage}\leanok
The absorber decomposes its union with the terminal remainder after the two deletion stages. Lift this decomposition through each deletion and join it with the corresponding packing. The size statement follows by choosing the order threshold after the two reserve densities.
\end{proof}

\begin{lemma}[Boundary completion for six-wheels]\label{bp:six-wheel-boundary-completion}
\lean{Oberwolfach.marked_actual_boundary_finish,Oberwolfach.marked_actual_boundary_finish_support,Oberwolfach.marked_sample_actual_boundary_completion,Oberwolfach.marked_sample_actual_boundary_completion_local}
\leanok
The stated deterministic and sampled boundary-completion results hold for $r\geq 1$, so their wheel length $r+5$ may be six. Their graph, sampling and other hypotheses remain in force. This does not by itself establish the full six-wheel cover-down construction.
\end{lemma}
\begin{proof}\leanok
Apply the mixed boundary-finishing result with parameter $r+1\geq 2$, and use its conclusion in the two sampled-completion arguments.
\end{proof}

\begin{lemma}[Six-wheel prescribed packing and sampled outside completion]\label{bp:six-wheel-outside-completion}
\lean{Oberwolfach.five_wheel_path_to_single_root,Oberwolfach.marked_rim_wheel_on_inner_domain_five,Oberwolfach.marked_inner_rim_wheel_after_deletion_five,Oberwolfach.marked_inner_prescribed_rim_wheel_packing_five,Oberwolfach.marked_inner_prescribed_edge_packing_five,Oberwolfach.marked_sparse_inner_domain_edge_packing_five,Oberwolfach.marked_sparse_outside_subgraph_packing_five,Oberwolfach.marked_sampled_outside_cover_five,Oberwolfach.marked_sampled_outside_completion_five}
\leanok
Under their stated neighbourhood, deletion, sparseness and sampling hypotheses, the prescribed-edge packing and the sampled outside-completion results hold for six-wheels with their original resource and degree bounds. The existing approximate-cover argument requires $c_0\leq6$ for this case, but its proof chooses $c_0\geq16$. The full six-wheel cover-down remains open, but the manuscript application selects a sufficiently large even wheel length.
\end{lemma}
\begin{proof}\uses{bp:six-wheel-boundary-completion}\leanok
Choose the marked step before the last ordinary join, while avoiding the forbidden vertices, and apply the original physical-deletion, greedy-packing and sampling arguments. The approximate-cover caller uses the existing long-cycle input only under its displayed cutoff.
\end{proof}

\begin{lemma}[Nested vortex restrictions and finite resource assembly]\label{bp:vortex-restriction-finite-assembly}
\lean{Oberwolfach.marked_vortex_nested_next_images,Oberwolfach.marked_vortex_nested_next_deeper,Oberwolfach.marked_vortex_nested_commutes,Oberwolfach.marked_vortex_nested_raw_restriction,Oberwolfach.marked_vortex_nested_rim_agreement,Oberwolfach.marked_vortex_nested_spoke_agreement,Oberwolfach.marked_vortex_nested_resource_composition,Oberwolfach.marked_vortex_nested_decomposition_lift,Oberwolfach.marked_vortex_ambient_decomposition_lift,Oberwolfach.finite_resource_layer_decompositions,Oberwolfach.marked_vortex_decompositions_from_exact_layers}
\leanok
For the actual fixed vortex, nested vertex and hub embeddings preserve the graph, permutation, one raw outcome, typed rims and spokes, and the physical resource maps. Decompositions lift through these embeddings. Exact deletion layers in a finite decreasing sequence, together with a terminal decomposition, give a decomposition of the initial resource set. Constructing the reference records and proving the layer identities for the full iteration remain open.
\end{lemma}
\begin{proof}\leanok
Use the inclusions of the nested vortex sets and compose the raw coordinate maps. The resource maps compose on rim-edge sets and spoke pairs. Successive set differences are disjoint and telescope to the initial set; lift each decomposition and combine the disjoint pieces.
\end{proof}

\begin{lemma}[Physical vortex iteration and a preselected terminal absorber]\label{bp:vortex-physical-absorbing-iteration}
\lean{Oberwolfach.marked_vortex_step_physical_conservation,Oberwolfach.marked_vortex_physical_iteration,Oberwolfach.marked_terminal_absorber_decomposition_lift,Oberwolfach.marked_resource_layers_cumulative_deletion_protects,Oberwolfach.marked_cover_down_avoids_deeper_spokes,Oberwolfach.marked_vortex_iteration_with_cumulative_deletion}
\leanok
For actual vortex packings satisfying the stated local containment and disjointness hypotheses, each lifted packing is exactly the physical resource layer it removes. If successive reference states represent the preceding remainders, these layers assemble with the terminal remainder. The displayed reserve-avoidance, terminal protection, divisibility and representation hypotheses allow the same preselected absorber to complete the decomposition. Establishing these hypotheses for the actual sampled construction remains open.
\end{lemma}
\begin{proof}\uses{bp:vortex-restriction-finite-assembly}\leanok
Lift each one-step resource identity, telescope the finite layers, and lift the terminal absorber decomposition from the host after the cumulative physical deletion. The cumulative deletion avoids the reserve; the cover-down star condition gives the stated protection of deeper spokes.
\end{proof}

\begin{lemma}[Actual next and terminal states]\label{bp:vortex-exact-states}
\lean{Oberwolfach.marked_restricted_next_state_exact,Oberwolfach.marked_vortex_next_state_exact,Oberwolfach.marked_vortex_next_state_after_packing,Oberwolfach.marked_terminal_remainder_survives_cumulative_deletion,Oberwolfach.marked_vortex_valid_terminal_state_representation,Oberwolfach.marked_vortex_iteration_from_valid_terminal_state}
\leanok
A remainder whose typed arcs and spokes lie on the next vortex pair has a deletion state on the next fixed sample with the same physical image. A valid final state gives an exactly represented divisible terminal remainder, which survives deletion of all previously covered layers. Under the stated stage, reserve and terminal-protection hypotheses, the preselected absorber completes the finite iteration. Proving that the states and packings satisfy these hypotheses simultaneously remains open.
\end{lemma}
\begin{proof}\uses{bp:vortex-physical-absorbing-iteration}\leanok
Restrict the original sampled outcome without resampling and take preimages of the surviving typed arcs and spokes. The physical resource embeddings compose exactly. Lift the final state to the original sample, transfer its divisibility through deletion, and apply the finite absorbing assembly.
\end{proof}

\begin{lemma}[Even wheel length and the successor margin]\label{bp:marked-even-wheel-parameters}
\lean{Oberwolfach.marked_even_wheel_reference_parameters,Oberwolfach.marked_r_one_rim_cardinality,Oberwolfach.marked_even_wheel_parameters_with_successor_margin,Oberwolfach.marked_vortex_next_state_from_physical_support}
\leanok
The marked-wheel index counts the ordinary rim arcs, so its total rim length is one greater. For every positive margin, every sufficiently large even wheel length can be written as $r+5=2(k_0+2)$ with $r$ odd and with the stated successor-density margin. A one-step remainder supported on the next vortex pair also determines an exact next-level state. The remaining event, quantitative state validity and caller substitutions are not supplied by these statements.
\end{lemma}
\begin{proof}\uses{bp:vortex-exact-states}\leanok
Take the half of the even wheel length to define $k_0$ and bound the reciprocal term in the successor density. For the state, use invariance of the next vortex pair and reindex the original typed arcs and spokes.
\end{proof}

\begin{lemma}[Valid initial and next vortex states]\label{bp:vortex-valid-initial-next}
\lean{Oberwolfach.marked_universal_terminal_absorber_supported_in_host,Oberwolfach.marked_full_domain_valid_state_restriction,Oberwolfach.marked_vortex_initial_valid_state,Oberwolfach.marked_vortex_initial_absorber_state,Oberwolfach.marked_cover_down_next_valid_state,Oberwolfach.marked_vortex_cover_down_next_valid_state,Oberwolfach.marked_vortex_iteration_original_host_from_valid_terminal}
\leanok
Under the displayed initial divisibility, reserve avoidance and size bounds, a preselected universal absorber gives a valid initial state with exact physical resources. A divisible cover-down remainder, together with the stated deeper-pair and deletion estimates, gives a valid state at the actual next vortex level with the same ambient remainder. The resulting finite assembly does not require a separate reserve-containment hypothesis. Constructing the common sample and proving all these inputs at every stage remain open.
\end{lemma}
\begin{proof}\uses{bp:marked-even-wheel-parameters,bp:vortex-exact-states}\leanok
Use the empty terminal remainder to obtain a decomposition of the reserve. Transfer the valid initial deletion state along the full-domain embeddings. At a later level apply the canonical remainder reindexing to the fixed sample and compose the physical resource maps.
\end{proof}

\begin{lemma}[An actual valid vortex cover-down step]\label{bp:vortex-constructed-valid-step}
\lean{Oberwolfach.marked_vortex_constructed_valid_step}
\leanok
Given an actual level packing satisfying its disjointness, remaining-resource, divisibility and deletion-estimate hypotheses, there is a valid state at the next vortex level. Its ambient remainder is contained in the current one, and their physical difference has a wheel decomposition in the original host. These hypotheses must still be established simultaneously for the constructed packings.
\end{lemma}
\begin{proof}\uses{bp:vortex-valid-initial-next}\leanok
Reindex the divisible remainder and its deletion estimates to the next fixed pair, and use the physical conservation identity for the packing.
\end{proof}

\begin{lemma}[Fixed-absorber completion from valid level steps]\label{bp:vortex-fixed-absorber-completion}
\lean{Oberwolfach.marked_universal_terminal_absorber_is_fixed,Oberwolfach.marked_fixed_terminal_absorber_supported_in_host,Oberwolfach.marked_vortex_initial_fixed_absorber_state,Oberwolfach.marked_vortex_fixed_absorber_finite_iteration,Oberwolfach.marked_vortex_fixed_absorber_complete_from_steps,Oberwolfach.marked_select_fixed_terminal_absorber,Oberwolfach.marked_vortex_completion_from_selected_absorber}
\leanok
For one fixed marked sample, assume the manuscript's fixed original-host absorber, the stated initial divisibility and reserve bounds, and a valid cover-down step for every valid state at each level. Then the original host has a wheel decomposition. The argument does not require all terminal-pair spokes to survive earlier deletions. The absorber can also be selected from the stated reserved-absorber property before iteration; the common sampled event and its quantitative step and absorber inputs remain to be constructed.
\end{lemma}
\begin{proof}\uses{bp:vortex-constructed-valid-step,bp:vortex-valid-initial-next,bp:reference-preselected-two-stage}\leanok
Take empty deletion graphs to obtain the fixed-absorber property, construct the valid initial state, and iterate the physical deletion layers. Transfer the divisible final remainder to the original host and apply the fixed absorber.
\end{proof}

\begin{lemma}[Active weights on the nested vortex types]\label{bp:nested-vortex-active-weights}
\lean{Oberwolfach.markedVortexActiveVertexEquiv,Oberwolfach.markedVortexActiveHubEquiv,Oberwolfach.marked_nested_vortex_active_successor_weights,Oberwolfach.markedVortexActivePermutation,Oberwolfach.markedVortexActiveVertexEmbedding,Oberwolfach.markedVortexActiveHubEmbedding,Oberwolfach.marked_vortex_active_permutation_embedding,Oberwolfach.marked_vortex_active_degree_fields,Oberwolfach.marked_vortex_actual_active_inputs,Oberwolfach.marked_vortex_active_geometry_bounds,Oberwolfach.marked_nested_vortex_active_geometry}
\leanok
For a fixed vortex satisfying the displayed degree and parameter conditions, the physical active vertex and hub sets are equivalent to the active sets inside the vortex subtypes. The existing difference-domain theorem supplies weights on every active vertex set before sampling; reindexing preserves the given permutation, the weights and the degree target determined by the active hub count. For the original permutation, its restriction to each active set commutes with the physical embedding, and the vortex degree estimates hold on the nested active graph. The physical active sets also satisfy explicit size, hub-density and scale-decrease bounds. Constructing the complete level records and their simultaneous sampled estimates remains an obligation.
\end{lemma}
\begin{proof}\uses{bp:vortexsuccessorweightfamilies,bp:fixed-vortex-weights}\leanok
Use the equivalences induced by the identity on physical vertices and hubs, transport the permutation and weights, and identify the active hub cardinalities.
\end{proof}

\begin{lemma}[Local loads on the actual active vortex sets]\label{bp:actual-active-local-load}
\lean{Oberwolfach.marked_actual_active_local_load_event,Oberwolfach.marked_active_size_threshold,Oberwolfach.marked_actual_active_local_load_event_eventually,Oberwolfach.marked_active_seed_mass_bound,Oberwolfach.marked_actual_active_local_load_event_seeded,Oberwolfach.marked_actual_active_local_load_event_manuscript_seeds}
\leanok
Under the displayed graph, vortex, scale, profile and sampling hypotheses, the small constant is chosen before the degree parameter, and weights can be fixed on the active vertex sets before the profile and sample are chosen. The local spoke-root load estimates then hold simultaneously at every selected level with the displayed probability bound. The same retained sampling law permits intersection with any other event having a stated failure bound, with the two errors added. The explicit size conditions follow from a terminal-scale threshold chosen after the vortex ratio; the seed-mass condition follows from the degree and density bounds, and the manuscript seed choices are supplied by the final specialization. The specific reference, absorber and balance events and their parameter substitutions have yet to be combined with this result.
\end{lemma}
\begin{proof}\uses{bp:nested-vortex-active-weights,bp:marked-local-load-vortex}\leanok
Use the actual active-set weights, permutation and geometry to instantiate the retained local-load theorem, then apply the finite-law intersection bound.
\end{proof}

\begin{lemma}[A joint reference, absorber and active-load event]\label{bp:actual-reference-active-load}
\lean{Oberwolfach.MarkedReferenceActiveAbsorberProperty,Oberwolfach.marked_fixed_vortex_reference_active_absorber_event_named,Oberwolfach.marked_joint_actual_reference_active_load,Oberwolfach.marked_joint_actual_reference_active_load_witness}
\leanok
Under the displayed graph, vortex, profile and hierarchy conditions, one retained sampling law supplies the reference configurations and estimates, the active row and index estimates, the robust terminal-absorber property and the actual active local-load estimates simultaneously, with the displayed combined failure bound; a fixed choice of errors yields a positive-mass common outcome. The weights are chosen before the profiles and outcome; the small constant is chosen before the degree parameter. The balance, parent weighted-load and remaining level-data fields still have to be joined.
\end{lemma}
\begin{proof}\uses{bp:actual-active-local-load,bp:fixed-vortex-reference-active-absorber}\leanok
Apply the reference and absorber event with half the degree margin, then intersect it with the active local-load event on the same retained law. Use the manuscript seed densities and add the failure probabilities.
\end{proof}

\begin{lemma}[A common outcome with adaptive parent loads]\label{bp:actual-reference-active-parent-load}
\lean{Oberwolfach.marked_actual_parent_packing_event,Oberwolfach.marked_joint_actual_reference_active_parent_load,Oberwolfach.marked_joint_actual_reference_active_parent_load_witness}
\leanok
Under the displayed graph, vortex, profile and hierarchy conditions, the reference and robust-absorber event, the active local-load estimates and the parent packing and potential-root load estimates hold together on one retained outcome of positive mass. The parent estimates are uniform over the displayed later physical deletions. Both weight families are fixed before the profiles and outcome, and the small constant is chosen before the degree parameter. Balance and the remaining level-data and parameter substitutions are still obligations.
\end{lemma}
\begin{proof}\uses{bp:actual-reference-active-load,bp:marked-packing-load-vortex}\leanok
Specialize the adaptive parent-load event to the actual vortex and the manuscript seed law. Intersect it with the reference, absorber and active-load event on the same law, and choose fixed error budgets that give positive mass.
\end{proof}

\begin{lemma}[Balance on the actual vortex pairs]\label{bp:actual-vortex-balance-domains}
\lean{Oberwolfach.marked_vortex_fixed_deeper_degree_fields,Oberwolfach.marked_actual_inner_balance_event,Oberwolfach.marked_actual_deeper_balance_event,Oberwolfach.marked_finite_domain_balance_empty,Oberwolfach.marked_vortex_deeper_terminal_empty,Oberwolfach.marked_vortex_terminal_balance_zero,Oberwolfach.marked_vortex_terminal_balance_invariant,Oberwolfach.marked_balance_record_margin_choice,Oberwolfach.marked_balance_record_error_bounds}
\leanok
Under their displayed graph, vortex and sampling hypotheses, the inner and deeper fixed vortex pairs each satisfy a retained balance event uniformly over the specified later deletions. In the deeper-pair event the roots range over the parent level, not only over the next level. Balance on the appended empty terminal pair is identically zero, before and after any resource deletion. A positive choice of tolerances fitting the inner and global balance margins and its associated size inequalities are also checked. These are separate events; a common outcome with all other inputs has not yet been constructed.
\end{lemma}
\begin{proof}\uses{bp:markedfinitebalanceadaptivevortexevent}\leanok
Apply the finite balance event to the fixed inner and deeper pairs, using the original parent-domain degree estimates and the respective vortex size bounds.
\end{proof}

\begin{lemma}[A joint balance event on the actual vortex]\label{bp:joint-actual-balance}
\lean{Oberwolfach.marked_actual_joint_balance_events,Oberwolfach.marked_actual_joint_balance_witness,Oberwolfach.marked_reference_event_three_quarters,Oberwolfach.marked_joint_actual_reference_parent_balance_witness}
\leanok
For the actual fixed vortex pairs, Lean checks one probability bound for the intersection of the inner-pair balance events through level $L+1$ and the deeper-pair balance events through level $L$ and an arbitrary further event on the same retained law. The bound is the sum of the displayed failure bounds, and the deletion choices remain quantified inside the events. A specialization gives positive mass when the further event has probability at least three quarters. With the stated inputs and sufficiently large order, the reference, absorber and parent-load event satisfies this bound, yielding one positive-mass outcome. The reference event and inner balance extend through the final selected level; deeper balance is checked through the preceding level, with the empty endpoint treated separately.
\end{lemma}
\begin{proof}\uses{bp:actual-vortex-balance-domains,bp:actual-reference-active-parent-load}\leanok
Choose the smaller admissible vortex tolerance, intersect the inner and deeper balance events with the reference event on the same retained law, and use the stated failure budgets to obtain positive mass.
\end{proof}

\begin{lemma}[Inner cuts on the common retained outcome]\label{bp:joint-actual-inner-cuts}
\lean{Oberwolfach.marked_vortex_fixed_parent_degree_fields_including_terminal,Oberwolfach.marked_actual_inner_cut_event,Oberwolfach.marked_joint_actual_reference_parent_balance_cut_witness}
\leanok
The retained cut estimates hold on the actual fixed inner vortex sets, including the final nonempty set. Under the displayed hypotheses they hold on one positive-mass retained outcome together with the reference, absorber, parent-load, active-load and balance estimates. The two weight families are fixed before the profiles and outcome. The cut inequalities are expressed on the actual inner subtype; their transport and numerical substitution into the cover-down records remain to be checked.
\end{lemma}
\begin{proof}\uses{bp:joint-actual-balance,bp:marked-retained-cut-vortex}\leanok
Apply the retained cut event to the shifted vortex levels, using the degree estimate through the terminal level. Intersect it with the reference and load event and then the two balance events on the same law, with fixed failure budgets.
\end{proof}

\begin{lemma}[Transport of the inner cut fields]\label{bp:actual-inner-cut-fields}
\lean{Oberwolfach.marked_vortex_inner_vertex_equiv_nested,Oberwolfach.marked_vortex_inner_hub_equiv_nested,Oberwolfach.marked_cut_families_spoke_large_cut,Oberwolfach.marked_cut_families_rim_large_cut,Oberwolfach.marked_vortex_inner_rim_cut_transport,Oberwolfach.marked_vortex_inner_spoke_cut_transport,Oberwolfach.marked_vortex_inner_cut_fields,Oberwolfach.marked_vortex_inner_cut_fields_uniform_densities,Oberwolfach.marked_uniform_inner_cut_parameter_choice,Oberwolfach.marked_uniform_inner_cut_record_choice}
\leanok
Under the displayed nonemptiness and numerical inequalities, the retained rim and spoke cut bounds on the next vortex level give the normalized cut inequalities on the inner subsets of its parent. The identification uses the same retained outcome and the exact vertex and hub embeddings. A further checked version uses a fixed lower bound on both sampling densities, and positive cut parameters satisfying its inequalities can be chosen below arbitrary fixed positive caps before the actual densities. The remaining construction must verify those conditions for the constants used in the cover-down records.
\end{lemma}
\begin{proof}\uses{bp:joint-actual-inner-cuts,bp:marked-retained-cut-vortex}\leanok
Convert the retained cut-family bounds to normalized large cuts and transport the two matrices through the exact inner-domain equivalences.
\end{proof}

\begin{lemma}[One common outcome at every vortex level]\label{bp:joint-actual-all-levels}
\lean{Oberwolfach.marked_actual_inner_balance_witness,Oberwolfach.marked_joint_actual_reference_all_levels_witness}
\leanok
Under the displayed hypotheses and fixed failure budgets, one retained outcome of positive mass satisfies the reference, absorber, parent-load, active-load, inner-cut and inner-balance estimates at every selected level. Deeper balance holds through the preceding level; when $L+1$ is the final vortex level, the deeper pair at level $L$ is empty. The constants are chosen before the vortex and profiles. Construction of compatible cover-down records and identification of their hypotheses remain open.
\end{lemma}
\begin{proof}\uses{bp:joint-actual-balance,bp:actual-reference-active-parent-load,bp:marked-retained-cut-vortex,bp:actual-vortex-balance-domains}\leanok
Choose common tolerances and thresholds before the vortex, then combine the retained events on the same probability law. Handle a single selected level by the inner-balance witness and longer sequences by the joint inner- and deeper-balance witness.\end{proof}

\begin{lemma}[Balance inputs for the vortex level records]\label{bp:actual-balance-record-inputs}
\lean{Oberwolfach.marked_vortex_actual_balance_fields,Oberwolfach.marked_vortex_actual_balance_record_inputs}
\leanok
For the actual terminal vortex sequence and the stated common balance event, the inner and deeper balance estimates hold at every selected parent level, including the empty deeper pair at the final level. Under the displayed numerical comparisons, their error bounds meet the inner and global margins used by the level records. The full level records and the substitutions for their remaining fields are not yet constructed.
\end{lemma}
\begin{proof}\uses{bp:joint-actual-all-levels,bp:actual-vortex-balance-domains}\leanok
Use the inner and deeper event estimates at their respective levels, the zero-balance identity for the terminal pair, and the explicit error-margin inequalities.
\end{proof}

\begin{lemma}[Parent load inputs for the level records]\label{bp:parent-load-record-inputs}
\lean{Oberwolfach.marked_parent_load_base_fields,Oberwolfach.marked_parent_load_record_inputs,Oberwolfach.marked_parent_load_parameter_choice,Oberwolfach.marked_vortex_parent_load_record_inputs}
\leanok
The adaptive parent-load estimate supplies the original-sample local and potential-root load estimates by taking empty deletion graphs. The potential-root estimate retains the whole parent root domain. Under the displayed error comparisons these estimates meet the two record tolerances, and the error parameters can be chosen below arbitrary fixed positive caps. The estimate is also specialized to the actual fixed vortex graphs, retained sample, profiles and weights. Insertion into the complete level records and the required numerical comparisons remain to be checked.
\end{lemma}
\begin{proof}\uses{bp:joint-actual-all-levels,bp:markedpackingfullhubloads}\leanok
Take zero deletion, use the identity hub map for the full root set, and choose the fixed error parameters using the displayed inequalities.
\end{proof}

\begin{lemma}[Parent configurations on the common retained outcome]\label{bp:joint-actual-parent-configurations}
\lean{Oberwolfach.marked_actual_parent_configuration_event,Oberwolfach.marked_parent_configuration_record_input,Oberwolfach.marked_parent_configuration_parameter_choice,Oberwolfach.marked_joint_actual_reference_all_levels_configurations_witness}
\leanok
The local configuration estimates on the actual fixed parent levels hold on the same positive-mass retained outcome as the reference, absorber, load, cut and balance estimates. A separate coefficient choice gives the displayed parent configuration threshold below arbitrary fixed positive caps. The identification of those constants with the full level-record hierarchy and the other configuration inputs remain open.
\end{lemma}
\begin{proof}\uses{bp:joint-actual-all-levels,bp:marked-retained-configurations-vortex}\leanok
Specialize the retained configuration event to the fixed parent levels and intersect it with the existing events on the same law. Choose the coefficient and retention-error parameters using the stated inequalities.
\end{proof}

\begin{lemma}[Ordered common configuration event]\label{bp:ordered-common-configurations}
\lean{Oberwolfach.marked_joint_actual_reference_all_levels_common_configurations_witness_ordered,Oberwolfach.marked_parent_configuration_hierarchy_choice,Oberwolfach.marked_vortex_parent_configuration_record_inputs,Oberwolfach.marked_actual_active_configuration_event,Oberwolfach.marked_local_configurations_hub_restriction,Oberwolfach.marked_external_hub_configuration_record_input,Oberwolfach.marked_actual_next_active_all_parent_hubs_configuration_event,Oberwolfach.marked_vortex_external_hub_configuration_record_inputs,Oberwolfach.marked_vortex_next_active_embedding_composition,Oberwolfach.marked_vortex_next_active_permutation_commutes,Oberwolfach.marked_vortex_next_active_parent_image,Oberwolfach.rawMarked_pair_nested_rims,Oberwolfach.marked_vortex_nested_external_hub_record_input,Oberwolfach.marked_shared_configuration_hierarchy_choice,Oberwolfach.marked_joint_actual_reference_all_levels_external_hub_configurations_witness,Oberwolfach.marked_parent_to_active_rows,Oberwolfach.marked_active_row_parameter_choice}
\uses{bp:joint-actual-parent-configurations,bp:marked-retained-common-configurations}
\leanok
The reference and parent common-configuration coefficients can be chosen before the later error parameters. The stated balance, cut, reference, absorber, parent-load and local-configuration estimates hold on one retained outcome of positive mass. The numerical choice and monotonicity lemmas supply the displayed parent configuration inputs at every level. A separate event controls configurations with vertices and hubs in each active set, including the last level. This last event does not by itself control hubs outside the active set. A separate checked event provides the estimate for every parent hub and the next active vertex set, and the restriction lemma gives the required coefficient bound on a fixed subfamily of hubs. The next active vertex embedding, its image in the parent, its commutation with the restricted permutation, and the conversion to the record’s nested rim restriction are also checked. A shared numerical choice supplies both parent and active coefficients before the smaller errors. The stronger all-parent-hub configuration event is now intersected with the other displayed events on one positive-mass outcome. A deterministic embedding lemma and an explicit error-parameter choice also prepare the transfer of the parent row estimates. The remaining record identifications must still be checked. The construction of the complete level records remains open.
\end{lemma}
\begin{proof}\uses{bp:joint-actual-parent-configurations,bp:marked-retained-common-configurations}\leanok
The retained configuration events are intersected on the same probability law. The ordered reference wrappers make the coefficient choices before the subsequent error choices. The record input estimates follow from the explicit coefficient inequalities and monotonicity.
\end{proof}

\begin{lemma}[Fixed weights and larger hub domains]\label{bp:retained-load-fixed-weights}
\lean{Oberwolfach.marked_joint_actual_reference_all_levels_external_hub_configurations_with_weights,Oberwolfach.marked_actual_next_active_nondeeper_hubs_load_event,Oberwolfach.marked_parent_hubs_next_active_size,Oberwolfach.marked_larger_hub_load_error_choice}
\uses{bp:ordered-common-configurations}
\leanok
The all-level witness retains the deterministic conditions for both fixed weight families before the profiles and the retained outcome are chosen. Separately, the displayed load event applies on the next active vertices and the non-deeper parent hubs for any supplied weight family satisfying those conditions. Its error depends on the ratio between the two domain sizes; the checked numerical lemma permits the additional retention error to be chosen after this ratio is fixed. Their common-outcome intersection is supplied by the later checked lemma. The remaining level-record inputs and the main theorem remain open.
\end{lemma}
\begin{proof}\uses{bp:ordered-common-configurations}\leanok
Carry the deterministic conditions through the weight-event constructions without changing the chosen weights. Bound the parent hub set relative to the next active vertex set and apply the retained-load event on the fixed non-deeper hub domain. Choose the additional error using the displayed numerical inequality.
\end{proof}

\begin{lemma}[Active row and load inputs]\label{bp:active-nondeeper-record-inputs}
\lean{Oberwolfach.marked_actual_next_active_nondeeper_hubs_rows_event,Oberwolfach.marked_rows_with_spokes_rim_equality,Oberwolfach.marked_vortex_nested_active_rows_input,Oberwolfach.marked_local_spoke_loads_rim_equality,Oberwolfach.marked_vortex_nested_active_load_input,Oberwolfach.marked_larger_hub_load_tolerance}
\uses{bp:retained-load-fixed-weights}
\leanok
On the next active vertex set and the non-deeper parent hubs, a separate retained-sample event supplies the row estimates. The checked rim identities convert the direct row and spoke-load estimates into the nested representation used by a level record, keeping the active error scale, fixed weights and hub domain. The displayed tolerance lemmas apply under their explicit numerical bounds. Their common-outcome intersection is supplied by the later checked lemma; the simultaneous numerical comparisons and complete level records remain open.
\end{lemma}
\begin{proof}\uses{bp:retained-load-fixed-weights}\leanok
Apply the retained-row event on the fixed vortex domains. The nested and direct restrictions read the same physical rim states, and their supplied spoke relations coincide. Apply the stated numerical tolerance comparisons.
\end{proof}

\begin{lemma}[One outcome for active row and load estimates]\label{bp:joint-active-record-events}
\lean{Oberwolfach.marked_joint_actual_reference_all_levels_nondeeper_hub_loads_with_weights,Oberwolfach.marked_joint_actual_reference_all_levels_nondeeper_rows_loads_with_weights,Oberwolfach.marked_active_rows_loads_parameter_choice}
\uses{bp:active-nondeeper-record-inputs,bp:ordered-common-configurations}
\leanok
For the displayed hypotheses and fixed vortex, the deterministic conditions on the parent and active weight families and the reference, absorber, balance, cut, parent configuration, active configuration, active row and active spoke-load estimates hold with the stated domains on one retained outcome of positive mass. The two additional retention errors are chosen after the vortex ratio and before the profiles and outcome. The row and load tolerances admit a joint numerical choice under arbitrary positive caps, including the later retention-error choice. Their compatibility with all other record fields, the complete parameter hierarchy and the main theorem remain open.
\end{lemma}
\begin{proof}\uses{bp:active-nondeeper-record-inputs,bp:ordered-common-configurations}\leanok
Intersect the stated estimates on the same finite probability law. The sum of their displayed failure bounds is below the required margin once the chosen order threshold is large enough. The fixed weight families are selected before the profiles and retained outcome.
\end{proof}

\begin{lemma}[Nested active record inputs]\label{bp:joint-active-record-inputs}
\lean{Oberwolfach.marked_vortex_nested_active_rows_and_loads,Oberwolfach.marked_vortex_inner_nondeeper_hubs_image,Oberwolfach.marked_vortex_nondeeper_permutation_embedding,Oberwolfach.marked_vortex_next_active_image_invariant,Oberwolfach.marked_local_configurations_restrict_hubs,Oberwolfach.marked_vortex_nested_active_configurations_input}
\uses{bp:joint-active-record-events,bp:active-nondeeper-record-inputs}
\leanok
For the displayed hypotheses, the row and load estimates on the same retained outcome yield the corresponding nested record inputs using the same profiles and weights. The configuration estimate on all parent hubs restricts to the non-deeper hubs, with the stated coefficient comparisons. The active vertex and hub images and the permutation actions satisfy the displayed set and commutation identities. Constructing the full level record and a compatible hierarchy remains open.
\end{lemma}
\begin{proof}\uses{bp:joint-active-record-events,bp:active-nondeeper-record-inputs}\leanok
Identify the nested rim states and retain the supplied spoke relations and fixed weights. For configurations, specialize the all-parent-hub estimate to each non-deeper hub. The geometric identities follow from the nested set inclusions and permutation commutation.
\end{proof}

\begin{lemma}[Vortex level degree and size inputs]\label{bp:vortex-level-geometry-inputs}
\lean{Oberwolfach.marked_vortex_parent_level_degree_inputs,Oberwolfach.marked_vortex_inner_level_degree_inputs,Oberwolfach.marked_vortex_nondeeper_parent_images,Oberwolfach.marked_vortex_deeper_parent_permutation_invariant,Oberwolfach.marked_level_small_remainder_bounds,Oberwolfach.marked_vortex_small_remainder_inputs,Oberwolfach.marked_vortex_inner_active_size_inputs}
\uses{bp:joint-active-record-inputs}
\leanok
Under the displayed vortex and parameter hypotheses, the parent and inner degree bounds hold on every level and at every parent vertex. The deeper sets have the displayed smallness bounds, the non-deeper types have the stated images, and the inner and next-active vertex sets satisfy the displayed lower bounds. These are individual inputs, not a construction of the full level records.
\end{lemma}
\begin{proof}\uses{bp:joint-active-record-inputs}\leanok
Use the vortex degree estimates and consecutive-size bounds, preserving the two independently rounded vertex and hub sets. Apply the displayed comparisons of the vortex error and size ratio with the record tolerances.
\end{proof}

\begin{lemma}[Full-parent neighbourhoods on the common outcome]\label{bp:joint-full-parent-neighbourhood}
\lean{Oberwolfach.marked_actual_parent_full_neighbourhood_event,Oberwolfach.marked_joint_actual_reference_all_levels_full_parent_neighbourhood_with_weights}
\uses{bp:joint-active-record-events,bp:vortex-level-geometry-inputs}
\leanok
Under the displayed hypotheses, the full-parent row and index estimates hold at every level on the same positive-mass outcome as the existing absorber, balance, cut, configuration and active load estimates, with both weight families fixed before the profiles and sample. The error parameters must still be matched to the complete level-record hierarchy; the retained-versus-original spoke-agreement fields remain to be supplied.
\end{lemma}
\begin{proof}\uses{bp:joint-active-record-events,bp:vortex-level-geometry-inputs}\leanok
Apply the all-level retained row and index events to the full fixed parent vertex and hub domains and intersect their events with the existing event on the same law. Bound the combined failure probability before choosing the common outcome.
\end{proof}

\begin{lemma}[Parent spoke agreement on the common outcome]\label{bp:joint-parent-spoke-agreement}
\lean{Oberwolfach.marked_actual_parent_spoke_agreement_event,Oberwolfach.marked_joint_actual_reference_all_levels_parent_agreement_with_weights}
\uses{bp:joint-full-parent-neighbourhood}
\leanok
Under the displayed hypotheses, the row and column bounds comparing the original and retained spokes hold at every parent level on the same positive-mass outcome as the other stated events. The parent and active weight families remain fixed before the profiles and the sample. The numerical substitutions into complete level records and their common parameter hierarchy remain open.
\end{lemma}
\begin{proof}\uses{bp:joint-full-parent-neighbourhood}\leanok
Apply the retention event on the fixed parent vertex and hub domains, transport its realized discrepancy bounds to the parent types, and intersect it with the previous event before choosing the outcome.
\end{proof}

\begin{lemma}[Parent neighbourhood record inputs]\label{bp:parent-neighbourhood-record-inputs}
\lean{Oberwolfach.marked_parent_neighbourhood_parameter_choice,Oberwolfach.marked_parent_neighbourhood_record_inputs}
\uses{bp:joint-parent-spoke-agreement}
\leanok
The displayed parent neighbourhood and agreement estimates yield the four corresponding record bounds under the stated comparisons, using the original row and column domain sizes. The tolerances admit a choice under arbitrary positive upper caps. Compatibility with every other level-record parameter remains to be proved.
\end{lemma}
\begin{proof}\uses{bp:joint-parent-spoke-agreement}\leanok
Enlarge the error bounds using the size comparison of the two parent types; choose the two positive tolerances with room below the displayed caps.
\end{proof}

\begin{lemma}[Two-step vortex bounds and endpoint inputs]\label{bp:two-step-vortex-and-endpoints}
\lean{Oberwolfach.marked_reference_active_endpoint_base,Oberwolfach.marked_active_endpoint_record_inputs,Oberwolfach.marked_vortex_inner_upper_size_inputs,Oberwolfach.marked_active_hub_density_arithmetic,Oberwolfach.marked_vortex_active_density_inputs,Oberwolfach.marked_vortex_active_inner_size_input,Oberwolfach.marked_vortex_two_step_remainder_inputs,Oberwolfach.marked_deep_ratio_parameter_choice}
\uses{bp:parent-neighbourhood-record-inputs,bp:vortex-level-geometry-inputs}
\leanok
The adaptive endpoint event gives the displayed zero-deletion estimates on the original fixed inner-minus-deeper domains. Under the stated inequalities, these provide the endpoint error bounds. The deeper sets have order at most $\varepsilon^2$ relative to their parent sets, while their bounds relative to the inner sets use $\varepsilon$. The active density and smallness estimates and a compatible numerical choice for these two-step bounds hold under the displayed hypotheses. The latter assumes the stated upper bound on the geometric constant. All remaining level-record fields and the full hierarchy are still to be assembled.
\end{lemma}
\begin{proof}\uses{bp:parent-neighbourhood-record-inputs,bp:vortex-level-geometry-inputs}\leanok
Use zero deletions for the endpoint estimates. Apply the two successive vortex size bounds with the separately rounded vertex and hub sets, and then the stated numerical comparisons.
\end{proof}

\begin{lemma}[Reference parent normalization]\label{bp:reference-parent-normalization}
\lean{Oberwolfach.marked_reference_parent_base,Oberwolfach.marked_reference_parameter_normalization,Oberwolfach.marked_reference_parent_minimum_degree_comparison}
\uses{bp:two-step-vortex-and-endpoints}
\leanok
Taking zero deletions in the reference event gives the stated parent statistics and cut families. With the joint event parameter $\alpha/2$, their degree-margin argument is $\alpha/8$; the record parameter $\gamma=\alpha/16$ therefore gives $2\gamma=\alpha/8$. The displayed profile normalizations and the corresponding parent minimum-degree comparison hold exactly. Insertion into complete level records remains open.
\end{lemma}
\begin{proof}\uses{bp:two-step-vortex-and-endpoints}\leanok
Specialize the adaptive reference event to the empty deletion and apply the displayed algebraic identities.
\end{proof}

\begin{lemma}[Vortex level input bridges]\label{bp:level-input-bridges}
\lean{Oberwolfach.marked_vortex_endpoint_record_inputs,Oberwolfach.marked_cut_threshold_before_ratio,Oberwolfach.marked_fixed_geometry_deep_ratio_choice,Oberwolfach.marked_vortex_cut_deep_hierarchy,Oberwolfach.marked_vortex_parent_neighbourhood_record_inputs,Oberwolfach.marked_vortex_level_minimum_sizes,Oberwolfach.marked_vortex_parent_hub_sizes,Oberwolfach.marked_vortex_inner_hub_sizes,Oberwolfach.marked_vortex_retained_spoke_geometry}
\uses{bp:two-step-vortex-and-endpoints,bp:reference-parent-normalization}
\leanok
Under their displayed hypotheses, the common-outcome parent and endpoint estimates give the stated record inputs on the actual nested vortex types. The parent and inner hub densities, minimum sizes, and retained-spoke disjointness and successor separation hold on every level. The cut threshold and two-step error admit the displayed joint choice with the cut threshold fixed before the vortex ratio and with a fixed geometric constant. The full record construction, compatibility with all remaining choices, its iteration, and the main theorem remain open.
\end{lemma}
\begin{proof}\uses{bp:two-step-vortex-and-endpoints,bp:reference-parent-normalization}\leanok
Use the stated fixed-domain events and vortex embeddings, the vertex and hub size recursions, and the disjointness of the retained spoke states. Choose the displayed positive parameters in the indicated order.
\end{proof}

\begin{lemma}[Reference inputs on the actual vortex]\label{bp:reference-vortex-record-inputs}
\lean{Oberwolfach.marked_vortex_reference_record_inputs}
\uses{bp:reference-parent-normalization,bp:level-input-bridges}
\leanok
The displayed reference event supplies the parent statistics with the exact record normalizations, including $\gamma=\alpha/16$, and cut estimates with the fixed lower density $t_*$. The stronger actual-density cut field is supplied separately by the later common-outcome result. Constructing the complete record remains open.
\end{lemma}
\begin{proof}\uses{bp:reference-parent-normalization,bp:level-input-bridges}\leanok
Use the zero-deletion reference estimates and substitute the checked profile and margin identities.
\end{proof}

\begin{lemma}[Active weight and coefficient inputs]\label{bp:nested-weight-target}
\lean{Oberwolfach.marked_active_configuration_coefficient_choice,Oberwolfach.marked_vortex_nested_active_weight_record_input,Oberwolfach.marked_vortex_active_hub_cardinality,Oberwolfach.marked_vortex_active_hub_degree_target,Oberwolfach.marked_vortex_nested_active_weight_target}
\uses{bp:level-input-bridges,bp:joint-active-record-events}
\leanok
The fixed active weight family satisfies the stated condition on the nested graph without changing its values. The active hubs at the next vortex level and their nested parent representation have the same cardinality, so the displayed degree target is unchanged. The active configuration coefficient can be chosen independently of the vortex regularity error under the stated positive caps. The construction of complete level records and the main theorem remain open.
\end{lemma}
\begin{proof}\uses{bp:level-input-bridges,bp:joint-active-record-events}\leanok
Identify the physical hub sets through the nested embeddings and transfer the unchanged weight condition; choose the coefficient below each displayed cap.
\end{proof}

\begin{lemma}[Actual-density parent cuts on the common outcome]\label{bp:actual-reference-cut-density}
\lean{Oberwolfach.rawSelectedSpokeMatrix_large_cuts_target_uniform,Oberwolfach.rawMarked_spoke_cuts_target_uniform,Oberwolfach.marked_reference_cut_event_target_uniform,Oberwolfach.marked_cut_vortex_target_uniform_event,Oberwolfach.marked_actual_parent_cuts_target_event,Oberwolfach.marked_joint_actual_reference_all_levels_actual_cut_density_with_weights,Oberwolfach.marked_vortex_reference_actual_cut_record_inputs}
\uses{bp:reference-vortex-record-inputs,bp:joint-parent-spoke-agreement}
\leanok
Under the displayed hypotheses, a fixed positive lower density controls the probability bound while the cut conclusion uses the actual supplied density $t$. This stronger cut event holds at every fixed parent level on the same positive-mass outcome as the other displayed marked estimates. Together with the reference statistics, it supplies the stated record inputs with their actual density and margin normalizations. The complete record, hierarchy, and main theorem remain open.
\end{lemma}
\begin{proof}\uses{bp:reference-vortex-record-inputs,bp:joint-parent-spoke-agreement}\leanok
Use degree and quadratic-cut concentration with errors fixed by the lower density, and apply the deterministic large-cut estimate at the actual target density. Transfer the original-law event through the same retained sampling law and intersect it with the existing common event before selecting the outcome.
\end{proof}

\begin{lemma}[Compatible record error margins]\label{bp:joint-record-error-margins}
\lean{Oberwolfach.marked_deep_gap_uniform_transfer,Oberwolfach.marked_vortex_active_degree_record_inputs,Oberwolfach.marked_joint_record_error_choice,Oberwolfach.marked_vortex_error_parameter_choice,Oberwolfach.marked_uniform_vortex_cut_deep_hierarchy}
\uses{bp:actual-reference-cut-density,bp:nested-weight-target}
\leanok
The cut and deep-error parameters can be chosen using fixed lower density bounds, before the graph and supplied spoke densities are known; the resulting deep-error inequality then holds for those later densities. The vortex regularity error can be chosen below the simultaneous sampling and degree caps. The active degree target satisfies the stated record bounds when the actual active set is nonempty. Under the displayed positive caps, the parent, active, neighbourhood, row and balance errors can be chosen together; the larger-hub retention error is chosen after the vortex ratio. This does not complete the full parameter hierarchy or construct the level records.
\end{lemma}
\begin{proof}\uses{bp:actual-reference-cut-density,bp:nested-weight-target}\leanok
Use the two density lower bounds, the active degree estimates, and the checked numerical error choices with the stated quantifier order.
\end{proof}

\begin{lemma}[Finite reference chains]\label{bp:finite-reference-chain}
\lean{Oberwolfach.markedReferenceClip_next,Oberwolfach.markedFiniteLevels_prefix,Oberwolfach.markedFiniteLevels_tail,Oberwolfach.markedFiniteReferenceChainToChain,Oberwolfach.marked_finite_reference_iteration_inputs,Oberwolfach.marked_uniform_finite_reference_iteration_inputs,Oberwolfach.marked_uniform_finite_family_reference_iteration,Oberwolfach.marked_uniform_finite_family_reserved_completion}
\uses{bp:restriction}
\leanok
A nonempty finite sequence of compatible marked reference level records gives a reference chain for the iteration. The sequence is extended by repeating its last record; only the original finite prefix is used in the iteration. The transition maps, graph restrictions, sample restrictions and spoke relations on that prefix are unchanged. With the checked uniform cover-down hierarchy, the finite family supplies the iteration steps and the corresponding reserved-absorber completion, subject to the stated absorber, divisibility and size hypotheses.
\end{lemma}
\begin{proof}\uses{bp:restriction}\leanok
Index both endpoints of each finite transition, and use the clipped index to extend the records and their ambient maps. For every transition used by the iteration the clipped indices are consecutive.
\end{proof}

\begin{lemma}[Completion with the common-sample absorber]\label{bp:sampled-fixed-absorber-completion}
\lean{Oberwolfach.marked_vortex_sampled_event_original_absorber,Oberwolfach.marked_vortex_completion_from_common_sample_absorber}
\uses{bp:vortex-fixed-absorber-completion,bp:actual-reference-cut-density}
\leanok
The absorber part of the common retained-sample event supplies the reserved-absorber property on the original sampled host. Given the displayed valid cover-down steps, divisibility and size bounds, one can select the absorber before the iteration, avoid the first inner pair and obtain a decomposition of the original host. The construction of those steps from the sampled level records remains open.
\end{lemma}
\begin{proof}\uses{bp:vortex-fixed-absorber-completion,bp:actual-reference-cut-density}\leanok
Take empty physical deletion graphs in the common event and apply the fixed-absorber completion on the same outcome.
\end{proof}

\begin{lemma}[Density-uniform inner margins]\label{bp:uniform-inner-margins}
\lean{Oberwolfach.marked_inner_factor_deletion_tolerance_uniform,Oberwolfach.marked_protected_inner_exception_order_uniform,Oberwolfach.marked_protected_inner_budget_scale_uniform}
\uses{bp:joint-record-error-margins}
\leanok
The inner-factor deletion tolerance, exceptional-set order threshold and protected inner error budget can be chosen using positive lower bounds on the relevant densities, before the actual densities are supplied. The coefficient in the error bound retains its dependence on the actual graph density; the proof controls it by the reciprocal of its fixed lower bound. The complete uniform cover-down hierarchy is not established by these inequalities alone.
\end{lemma}
\begin{proof}\uses{bp:joint-record-error-margins}\leanok
Apply the existing estimates at the fixed lower bounds and use monotonicity, including the reversed inequality for reciprocals.
\end{proof}

\begin{lemma}[Density-uniform global reference inputs]\label{bp:uniform-global-reference}
\lean{Oberwolfach.marked_active_reference_scale_uniform,Oberwolfach.marked_reference_cut_density_mono,Oberwolfach.marked_active_global_reference_conditions_uniform}
\uses{bp:uniform-inner-margins,bp:actual-reference-cut-density}
\leanok
Fix the displayed structural parameters and positive lower bounds for the graph, hub and spoke densities. One reference accuracy and, for each separation radius, one order threshold work for every later choice of actual densities above these bounds. The actual sampling profiles, degree target and weights remain unchanged; the four active cut bounds use the fixed lower target.
\end{lemma}
\begin{proof}\uses{bp:uniform-inner-margins,bp:actual-reference-cut-density}\leanok
Choose the numerical accuracy and order threshold from the lower bounds. The actual-density cut hypothesis implies the weaker fixed-target cut hypothesis, while the degree and fixed-weight estimates retain the actual parameters.
\end{proof}

\begin{lemma}[Density-uniform parent-reference active cover-down]\label{bp:uniform-parent-active-cover}
\lean{Oberwolfach.marked_sample_fresh_inner_hub_event_uniform,Oberwolfach.marked_sample_local_boundary_selection_uniform,Oberwolfach.marked_sample_profile_boundary_selection_uniform,Oberwolfach.marked_sample_quantitative_boundary_selection_uniform,Oberwolfach.marked_sample_robust_boundary_selection_uniform,Oberwolfach.marked_sample_actual_boundary_selection_local_uniform,Oberwolfach.marked_sample_actual_boundary_completion_local_uniform,Oberwolfach.marked_active_cover_down_local_uniform_spoke_density,Oberwolfach.marked_active_cover_down_global_conditions_uniform_spoke,Oberwolfach.marked_active_cover_down_canonical_roots_uniform_spoke,Oberwolfach.marked_active_cover_down_reference_uniform_densities}
\uses{bp:uniform-global-reference}
\leanok
Fix the displayed structural parameters and positive lower density bounds. The parent-reference and active cover-down constants, including their order threshold for each fixed separation radius, can be chosen before the actual graph, hub and spoke densities. The construction uses the original profiles, degree targets, fixed weights and one common physical sample. The endpoint class and gap inputs remain explicit.
\end{lemma}
\begin{proof}\uses{bp:uniform-global-reference}\leanok
Choose the fresh hub, boundary selection and completion parameters from the lower spoke-density bound, and apply the uniform parent-reference inputs. The available cut and hub-density bounds are weakened to the fixed lower targets without changing the underlying sampled host.
\end{proof}


\begin{lemma}[Uniform protected cover-down from packings]\label{bp:uniform-protected-cover}
\lean{Oberwolfach.marked_active_cover_down_endpoint_uniform_densities,Oberwolfach.marked_active_cover_down_after_packings_uniform_densities,Oberwolfach.marked_active_cover_down_from_deep_bounds_uniform_densities,Oberwolfach.marked_active_cover_down_separate_geometry_uniform_densities,Oberwolfach.marked_protected_cover_down_separate_geometry_from_packings_uniform_densities}
\uses{bp:uniform-parent-active-cover,bp:uniform-inner-margins}
\leanok
Given the stated lower bounds and previously constructed deep packings, the active and protected cover-down constants and the order threshold can be fixed before the actual graph, hub and spoke densities. The original profiles, physical deletion maps and weights are retained. The existence and uniform construction of the deep packings remain separate dependencies.
\end{lemma}
\begin{proof}\uses{bp:uniform-parent-active-cover,bp:uniform-inner-margins}\leanok
Choose fixed endpoint targets from the lower spoke-density bound and use the uniform deletion and error budgets. The resulting active packing is combined with the supplied packings using the established protected completion.
\end{proof}


\begin{lemma}[Density-uniform reference-level step]\label{bp:uniform-reference-level-density}
\lean{Oberwolfach.marked_deep_endpoint_profiles_uniform,Oberwolfach.marked_deep_buffered_endpoint_profiles_uniform,Oberwolfach.marked_sample_row_covers_after_hub_rims_uniform,Oberwolfach.marked_sample_hub_and_row_covers_uniform,Oberwolfach.marked_sample_deep_cover_after_physical_deletion_uniform,Oberwolfach.marked_reserved_sample_deep_cover_uniform,Oberwolfach.marked_corrected_sample_deep_cover_controlled_uniform,Oberwolfach.marked_constructed_deep_canonical_estimates_uniform,Oberwolfach.marked_protected_cover_down_separated_tolerances_uniform_densities,Oberwolfach.marked_protected_cover_down_iterable_uniform_densities,Oberwolfach.marked_fixed_reference_cover_down_step_uniform_densities,Oberwolfach.marked_uniform_reference_level_step_uniform_densities}
\uses{bp:uniform-protected-cover,bp:uniform-inner-margins}
\leanok
For fixed structural parameters and positive lower bounds on the densities, the constants and order threshold in the reference-level step can be chosen before the actual densities. For every level-data record satisfying its hypotheses, the step property follows, including protection and incidence estimates. The record itself remains an input and is not constructed by this lemma.
\end{lemma}
\begin{proof}\uses{bp:uniform-protected-cover,bp:uniform-inner-margins}\leanok
Choose the buffered endpoint, deep-packing and global and local deletion tolerances from the lower density bounds. Apply them to the actual profiles and fixed reference, and then transfer the protected packing to the next deletion state.
\end{proof}


\begin{lemma}[Density-uniform finite iteration and reserved completion]\label{bp:uniform-finite-completion-density}
\lean{Oberwolfach.marked_uniform_finite_reference_iteration_inputs_uniform_densities,Oberwolfach.marked_uniform_reference_finite_iteration_uniform_densities,Oberwolfach.marked_uniform_reference_reserved_completion_uniform_densities,Oberwolfach.marked_uniform_finite_family_reference_iteration_uniform_densities,Oberwolfach.marked_uniform_finite_family_reserved_completion_uniform_densities}
\uses{bp:uniform-reference-level-density,bp:finite-reference-chain,bp:sampled-fixed-absorber-completion}
\leanok
The finite reference iteration and its reserved terminal completion admit constants chosen before the actual densities, for each compatible finite family of valid records with common density parameters and the specified initial state and absorber. The common-event construction of those inputs is a separate obligation.
\end{lemma}
\begin{proof}\uses{bp:uniform-reference-level-density,bp:finite-reference-chain,bp:sampled-fixed-absorber-completion}\leanok
Apply the density-uniform reference-level step to each supplied record and use the existing finite iteration and terminal-reservoir completion.
\end{proof}


\begin{lemma}[Configuration choices in the reference-level hierarchy]\label{bp:ordered-level-configuration}
\lean{Oberwolfach.marked_shared_configuration_delayed_active_choice,Oberwolfach.marked_reference_level_step_ordered_configuration,Oberwolfach.marked_reference_level_step_chosen_configuration}
\uses{bp:uniform-reference-level-density}
\leanok
Choose the parent configuration coefficient and the local-configuration error before applying the density-uniform level theorem. Once the theorem supplies a positive active coefficient cap, choose a positive active coefficient below that cap and the parent coefficient, while preserving the estimates required for both the parent and active configurations. The resulting level step remains conditional on the actual level record.
\end{lemma}
\begin{proof}\uses{bp:uniform-reference-level-density}\leanok
Use the parent configuration hierarchy and choose the active coefficient as the minimum of the parent coefficient, a fixed density lower bound and the later cap. Apply the density-uniform reference-level step with the chosen parent coefficient.
\end{proof}


\begin{lemma}[Parent and active record fields on one sampled outcome]\label{bp:joint-parent-record-fields}
\lean{Oberwolfach.marked_vortex_joint_configuration_record_inputs,Oberwolfach.marked_joint_actual_reference_all_levels_parent_balance_fields_with_weights}
\uses{bp:ordered-level-configuration,bp:reference-vortex-record-inputs,bp:joint-record-error-margins}
\leanok
The existing common sampling construction yields one positive-mass outcome on which the displayed parent reference statistics, actual-density cuts, parent and active configurations, parent neighbourhood and spoke-agreement bounds, fixed-weight parent loads, and inner and deeper balance estimates hold at every selected level. The later deterministic choices must satisfy the displayed numerical inequalities, and the reference tolerances must be set to the indicated multiples of the level parameter before sampling. The remaining record fields and compatibility are separate obligations.
\end{lemma}
\begin{proof}\uses{bp:ordered-level-configuration,bp:reference-vortex-record-inputs,bp:joint-record-error-margins}\leanok
Use the same outcome and the parent and active weights supplied by the joint event. Apply the existing record-input transfers, retaining their domains and the event's simultaneous quantifiers. The terminal deeper-pair balance uses the stated terminal-index equality.
\end{proof}


\begin{lemma}[Analytic record fields on one retained outcome]\label{bp:joint-analytic-record-fields}
\lean{Oberwolfach.marked_vortex_inner_domains_nonempty,Oberwolfach.marked_joint_actual_reference_all_levels_endpoint_record_fields_with_weights}
\uses{bp:joint-parent-record-fields,bp:joint-actual-inner-cuts,bp:uniform-inner-margins}
\leanok
The common positive-mass event supplies, for all selected levels, the previously listed parent fields together with the active row and load bounds, inner cuts and endpoint neighbourhood fields. The displayed tolerances and density margins must hold, the reference tolerances are set before sampling, and the endpoint order must be specialized to the order required by the record. The terminal deeper pair is handled using the stated terminal-index equality. Constructing the full records and verifying the joint parameter hierarchy remain separate tasks.
\end{lemma}
\begin{proof}\uses{bp:joint-parent-record-fields,bp:joint-actual-inner-cuts,bp:uniform-inner-margins}\leanok
Apply the nested-domain row, load and endpoint transfers to the same fixed sample and weights. The vortex estimates imply that the inner vertex and hub sets are nonempty, including at the final nonempty level, so the inner-cut transfer applies.
\end{proof}


\begin{lemma}[Cut and error choices in the record hierarchy]\label{bp:ordered-record-error-choices}
\lean{Oberwolfach.marked_joint_record_ordered_error_choices,Oberwolfach.marked_joint_actual_record_ordered_error_choices}
\uses{bp:joint-record-error-margins,bp:uniform-inner-margins}
\leanok
The inner-cut error and cut threshold can be chosen before the balance tolerance is supplied. After any positive balance tolerance is given, the parent and active sampling errors, endpoint error and balance errors can be chosen to satisfy the displayed margins. The wider-hub sampling error is chosen later, after the vortex ratio is fixed.
\end{lemma}
\begin{proof}\uses{bp:joint-record-error-margins,bp:uniform-inner-margins}\leanok
Choose the cut error below a fixed positive density margin, apply the cut-parameter choice, and then apply the joint record error choice to the returned balance tolerance.
\end{proof}


\begin{lemma}[Compatible density bounds for sampling and the reference step]\label{bp:sample-reference-density-hierarchy}
\lean{Oberwolfach.marked_sample_reference_density_hierarchy}
\uses{bp:uniform-global-reference}
\leanok
Fixed lower bounds on the graph and total spoke densities give compatible positive lower bounds for the actual spoke density and the reference rim and marked-rim densities. They satisfy the required upper bounds and the sampling event's reference-probability inequality.
\end{lemma}
\begin{proof}\uses{bp:uniform-global-reference}\leanok
Use the same density products in the sample and reference targets, divide by the fixed cycle length and compare with the lower bounds.
\end{proof}


\begin{lemma}[Early vortex geometry and level-step constants]\label{bp:early-vortex-level-geometry}
\lean{Oberwolfach.marked_vortex_uniform_early_geometry,Oberwolfach.marked_reference_level_step_with_early_vortex_geometry}
\uses{bp:ordered-level-configuration,bp:ordered-record-error-choices,bp:uniform-inner-margins}
\leanok
The cut threshold and an upper bound on the vortex ratio can be fixed before choosing the density of the next active set. The resulting deep-set error satisfies the required inequality uniformly for all densities above the fixed lower bounds. The ratio is also small enough for the active-set size and inner-degree margins. The reference-level conclusion applies to every supplied valid record; this lemma does not construct such records.
\end{lemma}
\begin{proof}\uses{bp:ordered-level-configuration,bp:ordered-record-error-choices,bp:uniform-inner-margins}\leanok
Choose the cut and ratio bounds from the lower density parameters and the geometric constant. Monotonicity transfers the deep-set inequality to all admissible densities. Apply the level-step theorem with the resulting inner density, and choose the sampling errors after its balance tolerance has been supplied.
\end{proof}

\begin{lemma}[Vortex size and degree fields]\label{bp:assembled-vortex-size-degree}
\lean{Oberwolfach.marked_vortex_assembled_size_fields,Oberwolfach.marked_vortex_assembled_degree_fields}
\uses{bp:early-vortex-level-geometry}
\leanok
The vortex estimates give the displayed parent, inner and deeper size bounds and the parent and inner degree bounds at every selected level, with constants fixed from the original density and degree margins. In particular, the inner degree estimate holds for every parent vertex. The hypotheses and exact constants are those in the Lean statements; these estimates do not by themselves construct a sampled record.
\end{lemma}
\begin{proof}\uses{bp:early-vortex-level-geometry}\leanok
Combine the fixed-vortex size, degree and active-domain estimates, using the same vortex and the specified numerical margins at every level.
\end{proof}

\begin{lemma}[Exact sample and reference density identities]\label{bp:sample-reference-density-identities}
\lean{Oberwolfach.marked_sample_reference_density_identities}
\uses{bp:sample-reference-density-hierarchy}
\leanok
The density lower bounds in the common sample and the reference step agree exactly after dividing the total spoke density by the cycle length, including the reference-probability normalization.
\end{lemma}
\begin{proof}\uses{bp:sample-reference-density-hierarchy}\leanok
Use the definitions and the positivity of the cycle length.
\end{proof}


\begin{lemma}[Uniform actual vortex record geometry]\label{bp:uniform-vortex-record-geometry}
\lean{Oberwolfach.marked_vortex_uniform_geometric_record_fields,Oberwolfach.marked_vortex_record_map_fields}
\uses{bp:assembled-vortex-size-degree,bp:early-vortex-level-geometry}
\leanok
A single terminal-size threshold, chosen before the graph and vortex, supplies the displayed parent, inner, deeper and next-active size and degree fields at every selected level. The accompanying map identities give the nested vertex and hub embeddings, invariant sets, and their physical compositions. These are the geometric and map inputs, not a construction of the complete records from a sampled outcome.
\end{lemma}
\begin{proof}\uses{bp:assembled-vortex-size-degree,bp:early-vortex-level-geometry}\leanok
Combine the vortex size and degree bounds with the uniform minimum-size threshold. Use the actual nested embeddings and restricted permutations for the map identities.
\end{proof}

\begin{theorem}[Near-spanning cycles in an almost-regular graph]\label{bp:almost-regular-near-spanning-cycles}
\lean{Oberwolfach.cycle_sampling_almost_regular_parameters,Oberwolfach.ordinary_almost_regular_fractional_cycle_packing,Oberwolfach.ordinary_above_threshold_almost_regular_cycle_packing,Oberwolfach.ordinary_above_threshold_almost_regular_near_spanning_cycles,Oberwolfach.ordinary_above_threshold_near_spanning_cycles}
\uses{bp:cycle-packing-means,bp:positive-ordinary-cycle-weights}
\leanok
For each fixed $r\ge3$ and positive $\alpha,\varepsilon$, there are $\eta>0$ and $n_0$ such that every graph of order at least $n_0$, minimum degree at least $(\delta_{C_r}+\alpha)n$, and degrees within $\eta n$ of a common real number has vertex-disjoint $r$-cycles covering at least $(1-\varepsilon)n$ vertices. No divisibility condition on $n$ is needed. The separate exact-factor theorem applies above the same threshold when $r\mid n$; neither result alone covers the manuscript's lower-density odd-cycle callers.
\end{theorem}
\begin{proof}\uses{bp:cycle-packing-means,bp:positive-ordinary-cycle-weights}\leanok
Obtain bounded fractional cycle weights from the ordinary threshold. Choose a fixed sampling scale, an integer degree and a linear codegree budget from the almost-regularity bound, and apply the common-event cycle-packing theorem. A final deterministic choice of the tolerances gives the stated uncovered fraction.
\end{proof}


\begin{lemma}[Separate terminal and ambient scales in the joint event]\label{bp:joint-event-separated-scales}
\lean{Oberwolfach.marked_fixed_vortex_reference_active_absorber_event_ordered_separate_scales,Oberwolfach.marked_joint_actual_reference_active_load_ordered_with_weight_separate_scales,Oberwolfach.marked_joint_actual_reference_active_parent_load_ordered_with_weights_separate_scales,Oberwolfach.marked_joint_actual_reference_all_levels_actual_cut_density_with_weights_separate_scales,Oberwolfach.marked_joint_actual_reference_all_levels_endpoint_record_fields_with_weights_separate_scales}
\uses{bp:joint-analytic-record-fields,bp:uniform-vortex-record-geometry}
\leanok
The common sampling construction admits a terminal-scale threshold chosen before the absorber capacity. After the capacity and forbidden-set bound are fixed, an ambient graph-order threshold is chosen separately. On one positive-probability outcome the construction supplies the stated reference, load, balance, cut, configuration, agreement and endpoint fields at all selected levels, with the original fixed weights and numerical margins. Constructing the full compatible level records and applying the reserved completion remain open.
\end{lemma}
\begin{proof}\uses{bp:joint-analytic-record-fields,bp:uniform-vortex-record-geometry}\leanok
Keep the threshold required for the reference and level estimates on the terminal scale; impose the capacity-dependent absorber and finite-union thresholds only on the ambient graph order. Intersect the same events under the original sampling law and apply the existing endpoint transfers.
\end{proof}

\begin{lemma}[Terminal vortex pair fits a fixed capacity]\label{bp:vortex-terminal-capacity}
\lean{Oberwolfach.marked_vortex_terminal_pair_capacity,Oberwolfach.marked_vortex_terminal_pair_image_capacity}
\uses{bp:uniform-vortex-record-geometry}
\leanok
Under the displayed regularity and density hypotheses, the vertex and hub sets in the final nonempty vortex pair have total size at most twice the terminal scale. The same bound holds for their images under injective maps. Thus this numerical capacity can be chosen once the terminal scale is fixed.
\end{lemma}
\begin{proof}\uses{bp:uniform-vortex-record-geometry}\leanok
The terminal vertex set has fewer than the terminal-scale number of vertices and the terminal hub set is no larger. Injective images preserve both cardinalities.
\end{proof}

\begin{lemma}[Vortex and sampling with capacity chosen in order]\label{bp:chosen-vortex-joint-fields}
\lean{Oberwolfach.marked_vortex_reserved_capacity_fields,Oberwolfach.marked_vortex_choose_scale_for_absorber,Oberwolfach.marked_absorber_loss_budget_after_capacity,Oberwolfach.marked_chosen_vortex_joint_endpoint_record_fields}
\uses{bp:joint-event-separated-scales,bp:vortex-terminal-capacity,bp:joint-analytic-record-fields}
\leanok
After the displayed sampling parameters and a positive vortex error and ratio are fixed, the terminal scale can be chosen above any prescribed bound and the absorber capacity set to twice that scale, before the ambient graph order. After the bounded absorber size is known, the ambient threshold also ensures both prescribed positive linear loss bounds. Every eligible host then has a vortex and preserving permutation whose first and final pairs meet the stated size bounds. The same fixed weight families and one retained outcome give the displayed analytic fields, subject to the stated numerical margins. The construction of the complete level records and the iteration steps remains open.
\end{lemma}
\begin{proof}\uses{bp:joint-event-separated-scales,bp:vortex-terminal-capacity,bp:joint-analytic-record-fields}\leanok
Combine the vortex and permutation existence theorem with the separated-scale joint event and the first-pair and terminal-pair size estimates. The scale can be enlarged to meet the prescribed bound before choosing the capacity and ambient graph threshold.
\end{proof}

\begin{lemma}[Preselected absorber and conditional completion on the physical host]\label{bp:joint-sample-protected-completion}
\lean{Oberwolfach.marked_vortex_preselect_absorber_from_joint_sample,Oberwolfach.marked_vortex_preselect_absorber_and_apply,Oberwolfach.marked_vortex_completion_from_joint_sample_and_steps}
\uses{bp:chosen-vortex-joint-fields,bp:sampled-fixed-absorber-completion}
\leanok
The absorber part of the common sampled event yields one fixed absorber in the original host that avoids the first inner pair and any specified bounded forbidden set. If the stated valid cover-down steps exist, the original host is divisible, and the absorber size satisfies the displayed loss bounds, the physical host has the stated wheel decomposition. The lemma does not prove those remaining inputs.
\end{lemma}
\begin{proof}\uses{bp:chosen-vortex-joint-fields,bp:sampled-fixed-absorber-completion}\leanok
Take the empty physical deletion in the sampled absorber event, select the absorber using the first and final vortex size bounds, and apply the fixed-absorber physical iteration to the given steps.
\end{proof}

\begin{lemma}[Compatible densities and vortex parameters]\label{bp:coupled-vortex-parameters}
\lean{Oberwolfach.marked_coupled_reference_density_bounds,Oberwolfach.marked_vortex_record_parameter_choices,Oberwolfach.marked_odd_length_successor_margin,Oberwolfach.marked_reference_level_step_with_early_vortex_geometry_and_cap}
\uses{bp:early-vortex-level-geometry,bp:chosen-vortex-joint-fields}
\leanok
The sampled-event density lower bounds imply the corresponding cover-down density inequalities with the same normalization. A sufficiently large odd internal length supplies the successor-cycle margin. After this length is fixed, the inner-set bound can be capped, the vortex ratio chosen, and then a positive vortex quality parameter chosen below the displayed geometric and sampling bounds.
\end{lemma}
\begin{proof}\uses{bp:early-vortex-level-geometry,bp:chosen-vortex-joint-fields}\leanok
Use the exact density identities and the Archimedean property, then reduce the inner-set bound by monotonicity and choose positive minima for the remaining parameters.
\end{proof}

\begin{lemma}[Uniform slack for changed reference fields]\label{bp:uniform-perturbed-reference-fields}
\lean{Oberwolfach.marked_reference_perturbation_slack,Oberwolfach.marked_reference_statistics_delta_mono,Oberwolfach.marked_common_fields_uniform_perturbation,Oberwolfach.marked_deep_perturbation_slack,Oberwolfach.marked_deep_fields_uniform_perturbation}
\uses{bp:coupled-vortex-parameters,bp:joint-analytic-record-fields}
\leanok
For fixed positive density lower bounds, the indicated stricter reference tolerances leave a positive physical-change allowance chosen before the host. The seven common reference fields, and separately the parent configuration and fixed-weight local-load fields, then persist under the stated bounded typed changes. These results retain the original coefficients and root domains; they do not construct the prepared residual or the remaining record fields.
\end{lemma}
\begin{proof}\uses{bp:coupled-vortex-parameters,bp:joint-analytic-record-fields}\leanok
Choose the allowance below the error, density, configuration and load margins, and apply the corresponding stability estimates to the same change graphs.
\end{proof}

\begin{lemma}[Strict margins for the common sampled event]\label{bp:strict-joint-errors}
\lean{Oberwolfach.marked_joint_record_strict_error_choice}
\uses{bp:uniform-perturbed-reference-fields}
\leanok
The parent, active, row and balance sampling errors can be chosen simultaneously with the stated additional fractions of their target budgets. The error for the wider hub domain can be chosen after the vortex ratio is fixed. This is a numerical choice and does not itself construct a common sampled outcome.
\end{lemma}
\begin{proof}\uses{bp:uniform-perturbed-reference-fields}\leanok
Apply the joint error-choice lemma with halved target parameters and retain its common choices.
\end{proof}

\begin{lemma}[Ambient-root balance under realised changes]\label{bp:finite-balance-physical-changes}
\lean{Oberwolfach.marked_restricted_balance_after_arbitrary_changes,Oberwolfach.marked_restricted_balance_after_physical_changes,Oberwolfach.marked_restricted_balance_of_finite_balance,Oberwolfach.marked_finite_balance_transfer_after_physical_changes,Oberwolfach.marked_finite_balance_uniform_perturbation}
\uses{bp:uniform-perturbed-reference-fields,bp:strict-joint-errors}
\leanok
The finite-domain balance estimates transfer through bounded realised rim and spoke changes at every ambient root, using the same fixed summation domains. The loss is charged directly to the degree bounds of the physical change graphs. If the displayed stricter original errors and lower bound on the domain size hold, one uniform positive allowance preserves both the local tolerance and the original ambient-order error scale.
\end{lemma}
\begin{proof}\uses{bp:uniform-perturbed-reference-fields,bp:strict-joint-errors}\leanok
Identify finite-domain and restricted balances, bound each changed incidence by the realised physical degrees, and choose the change allowance from the two remaining error budgets.
\end{proof}

\begin{lemma}[Active fields and endpoints after physical changes]\label{bp:active-physical-changes}
\lean{Oberwolfach.marked_local_change_degree_scale,Oberwolfach.marked_physical_changes_restrict,Oberwolfach.marked_active_fields_uniform_physical_changes,Oberwolfach.marked_active_rows_uniform_physical_changes,Oberwolfach.marked_endpoint_rows_uniform_physical_changes,Oberwolfach.marked_parent_potential_uniform_perturbation}
\uses{bp:uniform-perturbed-reference-fields,bp:strict-joint-errors,bp:finite-balance-physical-changes}
\leanok
Under the stated stronger sampled estimates, physical change bounds transfer the active configuration counts, fixed-weight spoke loads and row counts to the smaller active vertex domain while retaining the non-deeper hub domain. Endpoint row and index errors retain the parent-order scale, and parent root-omitted potential loads retain their original weights and root conditions. These conditional estimates do not yet assemble the actual level records or establish the prepared residual.
\end{lemma}
\begin{proof}\uses{bp:uniform-perturbed-reference-fields,bp:strict-joint-errors,bp:finite-balance-physical-changes}\leanok
Restrict the change graphs through the commuting embeddings and bound their degrees on the active scale. Apply the corresponding stability estimates with a common sufficiently small allowance.
\end{proof}

\begin{lemma}[Protected later levels and terminal balance]\label{bp:protected-later-levels-terminal-balance}
\lean{Oberwolfach.marked_vortex_level_zero_physical_maps,Oberwolfach.marked_vortex_terminal_deeper_balance_zero,Oberwolfach.marked_protected_inner_cuts_after_changes,Oberwolfach.marked_vortex_protected_later_levels,Oberwolfach.marked_vortex_encoded_later_agreement,Oberwolfach.marked_spoke_geometry_from_marked_sets,Oberwolfach.marked_spoke_geometry_restrict}
\uses{bp:active-physical-changes,bp:finite-balance-physical-changes}
\leanok
The first vortex level is identified with the original vertex and hub sets, with its permutation preserved. Changes avoiding the first inner pair preserve its cut conditions and leave the complete encoded raw outcome and spoke relations unchanged on every later level. The deeper domains at the final level are empty, so their balances vanish. The stated exclusion of ordinary spokes from the marked vertices and their successors implies the required spoke disjointness and successor separation, also after restriction to a commuting subdomain.
\end{lemma}
\begin{proof}\uses{bp:active-physical-changes,bp:finite-balance-physical-changes}\leanok
Use the nested vortex inclusions, the protected raw-coordinate encoding, and the indicator definitions of typed rims and spokes. The final balance is zero on empty domains.
\end{proof}

\begin{lemma}[Sampled balance through physical changes]\label{bp:sample-balance-changes}
\lean{Oberwolfach.marked_sample_balance_error_budgets,Oberwolfach.marked_sample_balance_uniform_physical_changes,Oberwolfach.marked_sample_balance_ambient_physical_changes}
\uses{bp:finite-balance-physical-changes,bp:strict-joint-errors,bp:protected-later-levels-terminal-balance}
\leanok
The sampled balance error, scaled by the size of its original fixed domain, supplies the local and ambient budgets required for physical-change stability. The ambient-order version also applies to empty domains and makes no positive domain-size assumption. Both conclusions remain conditional on the displayed sampled balance and change-graph inputs.
\end{lemma}
\begin{proof}\uses{bp:finite-balance-physical-changes,bp:strict-joint-errors,bp:protected-later-levels-terminal-balance}\leanok
Bound the domain size by the ambient order and apply the corresponding uniform finite-balance transfer.
\end{proof}

\begin{lemma}[Refined physical exposure and retained common events]\label{bp:refined-physical-exposure}
\lean{Oberwolfach.marked_expanded_physical_collapse_law,Oberwolfach.raw_expanded_marked_wheel_law_projection,Oberwolfach.finite_refined_exposure_common_event,Oberwolfach.raw_expanded_marked_retained_common_event,Oberwolfach.marked_joint_actual_reference_all_levels_probability}
\uses{bp:strict-joint-errors,bp:protected-later-levels-terminal-balance}
\leanok
One physical-edge exposure distinguishes selected ordinary and marked states, approximate ordinary and marked states, an unassigned state, and the outside state. Forgetting the last four classes yields the selected physical law, while the interval coordinates retain their original joint law. The selected-layer completion event has the stated positive probability bound. A separately established high-probability event on the richer exposure can therefore be intersected with it, including after conditional retention. The full preparation event and its probability estimate are not established by this lemma.
\end{lemma}
\begin{proof}\uses{bp:strict-joint-errors,bp:protected-later-levels-terminal-balance}\leanok
Compute the coordinate pushforwards, apply the product-law and conditional-retention projection identities, and use the union bound.
\end{proof}

\begin{lemma}[Degrees in the enriched physical exposure]\label{bp:enriched-physical-degrees}
\lean{Oberwolfach.raw_expanded_marked_all_colour_degrees,Oberwolfach.raw_expanded_marked_colour_degrees_highProbability,Oberwolfach.marked_expanded_class_product_projection,Oberwolfach.raw_expanded_marked_selected_and_degrees}
\uses{bp:refined-physical-exposure}
\leanok
For any two fixed graphs, the directed degrees and undirected class degrees of all six physical classes satisfy simultaneous concentration bounds. The estimates remain valid after a conditional retention depending on the projected selected outcome. The union of the first five classes, together with their orientation, has the indicated two-class product law. The same outcome can be chosen to satisfy any selected retained event with probability at least $7/10$. This covers only part of the preparation event.
\end{lemma}
\begin{proof}\uses{bp:refined-physical-exposure}\leanok
Project the enriched product sample to its physical coordinates and use the categorical degree bounds. Sum the class probabilities to obtain the two-class projection.
\end{proof}

\begin{lemma}[A common selected, degree and two-path event]\label{bp:enriched-preparation-partial}
\lean{Oberwolfach.raw_expanded_marked_colour_degrees_999,Oberwolfach.raw_expanded_marked_class_two_paths,Oberwolfach.raw_expanded_marked_class_two_paths_highProbability,Oberwolfach.raw_expanded_marked_selected_degrees_two_paths}
\uses{bp:refined-physical-exposure,bp:enriched-physical-degrees}
\leanok
For sufficiently large vertex sets, a single outcome can satisfy the selected retained event of probability at least $7/10$, all six physical-class degree estimates, and the directed two-path estimates for the union of the five sampled classes. The latter two events each have probability at least $999/1000$ on the enriched exposure. This does not establish the other conditions of the manuscript's preparation event.
\end{lemma}
\begin{proof}\uses{bp:refined-physical-exposure,bp:enriched-physical-degrees}\leanok
Apply the degree and two-path concentration estimates on the same refined law and combine their failure probabilities before intersecting with the selected retained event.
\end{proof}

\begin{lemma}[Raw starts for every interval colour]\label{bp:enriched-raw-starts}
\lean{Oberwolfach.raw_interval_all_colours_starts_event,Oberwolfach.raw_expanded_marked_all_colours_starts_event}
\uses{bp:refined-physical-exposure}
\leanok
The number of hubs offering a raw interval of colour $g$ at a fixed start concentrates simultaneously for all starts and colours, with mean $\sum_u q(u,g)/d$. The probabilities $q(u,g)$ may vary with the hub. The same estimate holds on the refined exposure after conditional retention of the selected colour.
\end{lemma}
\begin{proof}\uses{bp:refined-physical-exposure}\leanok
Apply the independent-column tail bound to each start and colour, take a union bound, and project the refined exposure to its interval coordinates.
\end{proof}

\begin{lemma}[A common event including both interval colours]\label{bp:enriched-starts-joint-event}
\lean{Oberwolfach.raw_expanded_starts_highProbability,Oberwolfach.raw_expanded_marked_selected_degrees_two_paths_starts}
\uses{bp:enriched-preparation-partial,bp:enriched-raw-starts}
\leanok
When the hub set is nonempty, its size is at most the vertex-set size, and there are at most two interval colours, the all-colour raw start counts hold with probability at least $999/1000$. A single positive-mass outcome also satisfies these counts, the six-class degree estimates, the sampled-class two-path estimates, and any selected retained event of probability at least $7/10$. The hub-dependent marginal probabilities are fixed before sampling.
\end{lemma}
\begin{proof}\uses{bp:enriched-preparation-partial,bp:enriched-raw-starts}\leanok
Use the exponential bound and a sufficiently large vertex set; combine the three failure probabilities and apply the retained common-event projection.
\end{proof}

\begin{lemma}[Number of intervals at each hub]\label{bp:hub-interval-counts}
\lean{Oberwolfach.rawIntervalHubStartCount_expectation,Oberwolfach.rawIntervalHubFamily_card,Oberwolfach.rawIntervalHubStartCount_tail_singleton,Oberwolfach.raw_expanded_marked_hub_starts_all_event}
\uses{bp:enriched-raw-starts,bp:refined-physical-exposure}
\leanok
For each hub and interval colour, the raw number of starts has mean $|V|q(u,g)/d$. It is also the number of intervals in that raw family. A bounded-differences estimate holds simultaneously for all hubs and colours, including when the fixed probabilities $q(u,g)$ depend on the hub. The estimate remains valid on the enriched exposure with conditional retention.
\end{lemma}
\begin{proof}\uses{bp:enriched-raw-starts,bp:refined-physical-exposure}\leanok
Use the marginal probability of each start, injectivity of the interval map on the kept starts, and the bounded-differences inequality with the local coordinate sets. Then take a union bound and project the enriched exposure to its interval coordinates.
\end{proof}

\begin{lemma}[A common event including the hub interval counts]\label{bp:hub-counts-joint-event}
\lean{Oberwolfach.raw_expanded_hub_starts_highProbability,Oberwolfach.raw_expanded_marked_selected_degrees_paths_and_interval_counts}
\uses{bp:hub-interval-counts,bp:enriched-starts-joint-event}
\leanok
For each fixed interval length $d$, a sufficiently large vertex set gives the per-hub, per-colour raw interval-count estimates with error $|V|^{3/4}/d$ and probability at least $999/1000$. Under the displayed size and short-cycle conditions, these estimates, the all-colour start counts, the physical degree and two-path estimates, and a selected retained event of probability at least $7/10$ hold on a single positive-mass outcome. The completion event and the remaining preparation requirements still have to be connected to this result in the actual manuscript proof.
\end{lemma}
\begin{proof}\uses{bp:hub-interval-counts,bp:enriched-starts-joint-event}\leanok
Specialise the bounded-differences bound to the given error, then sum the failure probabilities on the same enriched retained law.
\end{proof}

\begin{lemma}[Unassigned states on fixed domains]\label{bp:unassigned-fixed-domains}
\lean{Oberwolfach.rawIntervalState_none_iff,Oberwolfach.rawIntervalFamily_unassigned_probability,Oberwolfach.rawIntervalUnassignedHubCountOn_tail,Oberwolfach.rawIntervalUnassignedColumnCount_tail,Oberwolfach.raw_expanded_unassigned_fixed_domains_highProbability,Oberwolfach.raw_expanded_unassigned_rim_degree_highProbability}
\uses{bp:refined-physical-exposure,bp:hub-interval-counts}
\leanok
The unassigned auxiliary state excludes covered labels and assigned ordinary labels, but includes uncovered successor endpoints. Its probability is $1-\sum_g q(u,g)-\sum_j s(u,j)$. For vertex and hub sets fixed before exposure, the corresponding restricted and unrestricted counts concentrate at the displayed $|V|^{3/4}$ scale. The physical unassigned degrees into a fixed vertex set concentrate simultaneously at all ambient roots, including roots inside that set. The high-probability auxiliary statement assumes the fixed hub set is nonempty.
\end{lemma}
\begin{proof}\uses{bp:refined-physical-exposure,bp:hub-interval-counts}\leanok
Use the exact categorical partition and the local interval dependence for counts within a hub; use independence of distinct hub columns for counts at a fixed vertex. For physical degrees, take the root-specific graph having exactly the relevant neighbours of that root. Project all estimates to the enriched retained exposure.
\end{proof}

\begin{lemma}[Basic preparation estimates on one outcome]\label{bp:basic-preparation-common}
\lean{Oberwolfach.raw_expanded_selected_preparation_basic_common_event}
\uses{bp:hub-counts-joint-event,bp:unassigned-fixed-domains}
\leanok
Under the stated fixed-domain, size, and short-cycle conditions, one positive-mass outcome satisfies the selected retained event of probability at least $7/10$, the physical degree and directed two-path estimates, both interval-colour start and hub counts, the restricted and global unassigned auxiliary counts, and the restricted unassigned physical degrees. The sets are chosen before exposure. Approximate-packing, list, and adaptive-choice properties are not included in this result.
\end{lemma}
\begin{proof}\uses{bp:hub-counts-joint-event,bp:unassigned-fixed-domains}\leanok
Establish the seven stated events on the same enriched law, sum their failure probabilities, and intersect with the selected retained event.
\end{proof}

\begin{lemma}[Approximate-layer projection and potential concentration]\label{bp:approximate-projection-potentials}
\lean{Oberwolfach.marked_expanded_approximate_physical_law,Oberwolfach.raw_expanded_approximate_wheel_law_projection,Oberwolfach.raw_expanded_approximate_probability_after_retention,Oberwolfach.raw_expanded_approximate_marked_potentials_event}
\uses{bp:refined-physical-exposure}
\leanok
Projecting the enriched exposure onto the approximate ordinary and marked physical states gives exactly the approximate categorical law and leaves every raw interval coordinate unchanged. After any conditional retention based on the selected projection, a raw approximate event has the same marginal probability. Consequently the stated root-omitted marked potential concentration estimate applies on this same enriched exposure with its fixed supplied weights. The identification of its expectations with the manuscript's specific targets and the other approximate-packing estimates remain to be checked.
\end{lemma}
\begin{proof}\uses{bp:refined-physical-exposure}\leanok
Compute the coordinate pushforwards and their product law, sum out the conditional retention, and apply the marked potential concentration estimate.
\end{proof}

\begin{lemma}[High-probability approximate marked loads and an additional event]\label{bp:approximate-high-extra}
\lean{Oberwolfach.raw_expanded_approximate_marked_potentials_highProbability,Oberwolfach.raw_expanded_selected_preparation_with_extra_event}
\uses{bp:approximate-projection-potentials,bp:basic-preparation-common}
\leanok
For fixed interval and wheel lengths, coefficient bound and error tolerance, the root-omitted approximate marked potential estimates hold with probability at least $999/1000$ for all sufficiently large orders, uniformly in the remaining displayed inputs and fixed weights. Independently, the basic preparation common-event argument permits one further event of probability at least $999/1000$ on the same enriched exposure. These results can be used together only after their input parameters and fixed weight functions have been matched; the manuscript-specific target normalization and the remaining packing and list estimates are still open.
\end{lemma}
\begin{proof}\uses{bp:approximate-projection-potentials,bp:basic-preparation-common}\leanok
Bound the number of possible roots by four times the square of the vertex-set size and use the exponential tail. Add the further event to the existing failure budget.
\end{proof}

\begin{lemma}[Basic preparation and approximate marked loads on one outcome]\label{bp:approximate-and-basic-common}
\lean{Oberwolfach.raw_expanded_selected_preparation_with_approx_potentials}
\uses{bp:basic-preparation-common,bp:approximate-high-extra}
\leanok
For fixed interval length, marked wheel length, coefficient bound and tolerance, one positive-mass enriched retained outcome satisfies the displayed basic preparation estimates and the root-omitted approximate marked potential estimates with the supplied fixed weight function. The weights and their bounds are inputs to this intermediate statement; the identification of the expected loads with the manuscript's targets, the ordinary approximate potentials and the remaining packing and list estimates remain open.
\end{lemma}
\begin{proof}\uses{bp:basic-preparation-common,bp:approximate-high-extra}\leanok
Choose the order threshold after the fixed parameters, apply the approximate potential high-probability estimate, and use it as the additional event in the basic preparation common-event lemma.
\end{proof}

\begin{lemma}[Variable-hub marked seed and marked-spoke target]\label{bp:variable-hub-marked-spoke}
\lean{Oberwolfach.approximate_preparation_target_arithmetic,Oberwolfach.rawMarkedSeed_candidate_variable_hub_mean,Oberwolfach.rawMarkedSeed_variable_hub_expectation,Oberwolfach.marked_variable_hub_seed_upper,Oberwolfach.rawMarked_variable_marked_spoke_target,Oberwolfach.rawMarked_variable_marked_spoke_concentration_target}
\uses{bp:approximate-projection-potentials}
\leanok
The root-omitted marked seed expectation admits an exact formula with the prescribed hub-dependent interval and ordinary-spoke probabilities. At a marked-spoke root of positive marginal, its mean is within the stated positional error of one; the local root probability cancels, and no positivity is imposed on the marked probabilities at other hubs. The supplied lower bounds on the ordinary-spoke and physical densities give a uniform seed coefficient bound. The concentration estimate then yields the marked-spoke target with the sum of the sampling and positional errors. The other marked-cycle root calculations and the ordinary-cycle contributions are treated separately.
\end{lemma}
\begin{proof}\uses{bp:approximate-projection-potentials}\leanok
On the resources of one wheel, all auxiliary marginals use its own hub; apply the fixed-hub product identity and sum over wheels. Identify the marked-spoke sum with the positional sum of the original successor weighting, then use its support and positional bounds. Apply the triangle inequality to the concentration event.
\end{proof}

\begin{lemma}[Marked-spoke target on the basic preparation outcome]\label{bp:marked-spoke-common-target}
\lean{Oberwolfach.raw_expanded_selected_preparation_with_marked_spoke_target}
\uses{bp:approximate-and-basic-common,bp:variable-hub-marked-spoke}
\leanok
With the displayed fixed density lower bounds and supplied successor weighting, the marked-spoke target at every positive-marginal root holds on the same positive-mass enriched outcome as the basic preparation estimates and the approximate marked concentration event. The weights and lower bounds remain inputs to this intermediate result; verifying their actual manuscript choices and the other preparation properties is still required.
\end{lemma}
\begin{proof}\uses{bp:approximate-and-basic-common,bp:variable-hub-marked-spoke}\leanok
Use the uniform seed bound to apply the common-outcome lemma to the original hub-normalized seed. Convert its concentration conclusion at each positive marked-spoke root using the deterministic target estimate.
\end{proof}

\begin{lemma}[Marked-cycle contributions to all four root types]\label{bp:marked-all-roots}
\lean{Oberwolfach.raw_expanded_selected_preparation_with_all_marked_contributions}
\uses{bp:marked-spoke-common-target}
\leanok
For a fixed marked type and ordinary colour, the same positive-mass preparation outcome satisfies the marked-cycle target estimates at all four root types. At an ordinary-spoke root the marked contribution has target $m q(u,g)/s(u,j)$; at an ordinary directed rim root it has target $mk/p$. At a marked rim root its target is one, using the supplied identity $\sum_u q(u,g)=Dk$. The probabilities, weights, and target identity are fixed inputs, whose manuscript instances still require verification.
\end{lemma}
\begin{proof}\uses{bp:marked-spoke-common-target}\leanok
Factor each root-omitted expectation using the hub-specific marginals. Express the incidence sums through the supplied positional bounds and apply the common concentration event. The marked closing arc has auxiliary head $z$ and physical head $\pi(z)$.
\end{proof}

\begin{lemma}[Ordered ordinary-cycle seed expectations]\label{bp:ordered-ordinary-seed}
\lean{Oberwolfach.orderedOrdinary_potential_expectation,Oberwolfach.orderedOrdinary_seed_spoke_expectation,Oberwolfach.orderedOrdinary_seed_rim_expectation,Oberwolfach.orderedOrdinary_spoke_weighted_sum,Oberwolfach.orderedOrdinary_rim_weighted_sum}
\uses{bp:approximate-projection-potentials}
\leanok
For the original hub-dependent coefficients, the root-omitted ordinary-cycle potential has the displayed exact spoke and directed-rim expectation formulas in the linked Lean declarations. Their incidence sums enumerate each ordered position once. These identities do not yet supply their manuscript positional bounds or a simultaneous concentration event.
\end{lemma}
\begin{proof}\uses{bp:approximate-projection-potentials}\leanok
The physical edges of a separated cycle and its spoke interval coordinates are pairwise disjoint. Factor the remaining resource probabilities after omitting the tested root, cancel the fixed seed denominator, and count positions.
\end{proof}

\begin{lemma}[Ordinary-cycle targets and simultaneous root concentration]\label{bp:ordered-ordinary-target-concentration}
\lean{Oberwolfach.orderedOrdinary_seed_spoke_target,Oberwolfach.orderedOrdinary_seed_rim_target,Oberwolfach.orderedOrdinary_omitted_pair_bound,Oberwolfach.raw_expanded_approximate_ordered_potentials_highProbability,Oberwolfach.raw_expanded_approximate_marked_potentials_highProbability_any,Oberwolfach.finiteLaw_high_events_common,Oberwolfach.preparation_balanced_means}
\uses{bp:ordered-ordinary-seed,bp:marked-all-roots}
\leanok
For a fixed cycle length and fixed pre-exposure coefficients, the directed ordinary-cycle potential concentrates simultaneously at every spoke and directed rim root, with any prescribed positive failure budget. The deterministic target bounds use the two-sided positional error. The same tunable budget is available for the marked potential; finite intersections need no independence. These are intermediate statements; the positional estimates and joint family event are addressed below, while their application at the full manuscript caller remains open.
\end{lemma}
\begin{proof}\uses{bp:ordered-ordinary-seed,bp:marked-all-roots}\leanok
Count candidates through their fixed physical vertices and hub. For two distinct rim requests, their physical edges are distinct and force at least three vertices. The resulting root and pair bounds imply the concentration estimate. Lift it to the enriched experiment by the exact projection, then apply the union bound and the fixed-coefficient balance identity.
\end{proof}

\begin{lemma}[Physical-cycle weights at ordered positions]\label{bp:ordinary-physical-position-transfer}
\lean{Oberwolfach.raw_expanded_selected_preparation_with_finite_extra_events,Oberwolfach.finite_uniform_threshold,Oberwolfach.ordered_cycle_representation_card,Oberwolfach.orderedCycle_position_fibre_unique,Oberwolfach.orderedCycle_position_vertex_card,Oberwolfach.separated_ordered_physical_edge_sum,Oberwolfach.separated_ordered_physical_vertex_sum,Oberwolfach.separated_physical_position_edge_error,Oberwolfach.separated_physical_position_vertex_error}
\uses{bp:ordered-ordinary-target-concentration}
\leanok
For a fixed embedded cycle of length $r$, the ordered representations are in bijection with the directed edges, giving $2r$ representations. At any fixed position, an admissible directed edge occurs in exactly one representation, and a vertex occurs with multiplicity equal to its degree in the cycle. Separation is a property of the vertex set. Thus the positional sums of the original physical-cycle weights equal the corresponding sums over conflict-free cycles, and the linked error bounds follow from the existing conflict-loss estimates. A finite family of additional events can be intersected with the preparation outcome without an independence hypothesis. The parameter budgets are addressed below; the full manuscript caller substitutions remain to be verified.
\end{lemma}
\begin{proof}\uses{bp:ordered-ordinary-target-concentration}\leanok
Two orderings with the same first directed edge agree by induction around the cycle. Rotations and reflections supply every directed edge as a first edge, and rotations transfer the count to every position. Partition separated host tuples by their underlying cycle, apply the exact multiplicities, and use the conflict-graph degree bound in the physical-cycle estimates.
\end{proof}

\begin{lemma}[Fixed seeds on one preparation sample]\label{bp:preparation-concrete-seed-families}
\lean{Oberwolfach.physicalOrdinaryTupleWeight_edge_positions,Oberwolfach.physicalOrdinaryTupleWeight_vertex_positions,Oberwolfach.ordinary_positional_error_budget,Oberwolfach.physicalOrdinary_seed_upper,Oberwolfach.raw_expanded_preparation_concrete_seed_families,Oberwolfach.physical_ordinary_sampled_spoke_target,Oberwolfach.physical_ordinary_sampled_rim_target,Oberwolfach.preparation_concentrated_balanced_sum,Oberwolfach.long_preparation_density_bounds,Oberwolfach.long_preparation_eta_bound,Oberwolfach.finite_length_selection_uniform_threshold}
\uses{bp:ordinary-physical-position-transfer,bp:marked-all-roots}
\leanok
For fixed finite families of ordinary and marked types, fixed positive lower bounds on the relevant densities, and the hypotheses displayed in the linked declarations, the original hub-dependent seed formulas satisfy the required coefficient bounds. One outcome has the stated root-omitted concentration estimates for all types together with the earlier preparation estimates and an additional selected event of probability at least $7/10$. The ordinary spoke and eligible rim targets follow from the original physical-cycle weights and the displayed error budgets. In the marked formula, the Lean parameter counts ordinary positions: the marked tuple has one additional vertex.

This result does not establish the selected event for the reserved absorber, the fixed-domain reference estimates, the cycle-factor callers or the theorem's conclusion. A finite-maximum argument permits a bounded set of short lengths to be chosen after a common order threshold, provided the required threshold has been established for each such set. The uniform choice of parameters at the full manuscript caller remains to be checked.
\end{lemma}
\begin{proof}\uses{bp:ordinary-physical-position-transfer,bp:marked-all-roots}\leanok
Transfer the physical cycle weights to ordered separated tuples, bound the lost cycles, and apply the concentration estimates with a finite union bound on the same enriched sample. The deterministic target and density inequalities use the fixed hub coefficients and their stated margins.
\end{proof}

\begin{lemma}[One selected outcome for all four root types]\label{bp:preparation-balanced-roots}
\lean{Oberwolfach.raw_expanded_ordinary_selected_lengths_highProbability,Oberwolfach.raw_expanded_preparation_selected_length_seed_families,Oberwolfach.rawMarked_variable_ordinary_spoke_concentration_target,Oberwolfach.rawMarked_ordinary_arc_concentration_target,Oberwolfach.rawMarked_closing_arc_concentration_target,Oberwolfach.combined_preparation_all_ordinary_roots,Oberwolfach.raw_expanded_preparation_balanced_all_roots,Oberwolfach.long_preparation_concrete_marginals,Oberwolfach.reference_event_seven_tenths}
\uses{bp:preparation-concrete-seed-families}
\leanok
With the hypotheses of the linked declarations, including the actual fixed cycle weights, numerical budgets and a selected event of probability at least $7/10$, one selected outcome has the earlier preparation estimates and all four balanced root estimates simultaneously. The set of short cycle lengths may be chosen after the order threshold within a fixed finite ambient set. The hub coefficient and density identities follow from the stated retained component fractions and marginals. The displayed reference-event failure bound implies the required probability for sufficiently large order if its fixed error terms sum to less than $3/10$.

The reference event itself, the needed fixed cycle weights, and the full caller substitutions have not been established here. In particular, this does not complete the protected first cover-down or prove Theorem~1.1.
\end{lemma}
\begin{proof}\uses{bp:preparation-concrete-seed-families}\leanok
Take a maximum of the finite collection of thresholds, intersect the root-concentration events on the same sample, and use the individual marked and ordinary estimates with the hub and rim balance identities. The stated polynomial error terms vanish as the order tends to infinity.
\end{proof}

\begin{lemma}[Preparation weights chosen before the allowed set]\label{bp:preparation-weights-allowed}
\lean{Oberwolfach.preparation_weights_from_allowed_threshold,Oberwolfach.selected_lengths_positive_fractional_weights,Oberwolfach.marked_preparation_length_choice,Oberwolfach.marked_preparation_weight_choice}
\uses{bp:successor-complete,bp:positive-ordinary-cycle-weights}
\leanok
Given a positive degree margin and a lower bound on the marked length, one can choose an even marked cycle length and its weight constants before choosing the finite family of short lengths. After choosing the positional tolerance, the degree tolerance and order threshold are independent of the allowed set. For any subsequently chosen nonempty allowed set of lengths at least three, the linked theorem constructs the bounded ordinary fractional cycle weights, their positional error budgets, and the marked successor weight, provided the graph has the displayed minimum degree and is sufficiently close to regular.

Applying this to the actual graph and the remaining hypotheses of the common-event theorem is still part of the unfinished preparation proof.
\end{lemma}
\begin{proof}\uses{bp:successor-complete,bp:positive-ordinary-cycle-weights}\leanok
Choose the marked length so that its successor-cycle threshold lies below the degree margin. Apply the successor-weight theorem and the ordinary fractional-cycle theorem, and take a maximum over the finitely many short-length thresholds. The supremum definition of the allowed-set threshold supplies the required degree bound for every selected length.
\end{proof}

\begin{lemma}[Reference witness with a positive-mass approximate extension]\label{bp:reference-positive-extension}
\lean{Oberwolfach.finiteLaw_sigma_positive_extension_probability,Oberwolfach.finiteLaw_sigma_reference_extension,Oberwolfach.raw_expanded_reference_positive_extension,Oberwolfach.raw_expanded_preparation_basic_plus_extra_highProbability,Oberwolfach.selected_retention_probability_bridge,Oberwolfach.long_preparation_interval_feasible,Oberwolfach.long_preparation_physical_feasible,Oberwolfach.long_preparation_unassigned_probability,Oberwolfach.long_preparation_marginal_bounds,Oberwolfach.long_preparation_residual_bounds}
\uses{bp:preparation-balanced-roots}
\leanok
For the actual expanded sampling law, a joint event of probability at least $99/100$ projects to a retained event of at least the same probability whose outcomes each have an extension in the joint event with positive joint mass. Hence a retained reference existence theorem that accepts arbitrary events of probability at least $99/100$ yields a joint outcome satisfying both the reference property and the joint event, without an independence assumption. The basic preparation estimates together with an additional event of probability at least $999/1000$ form such a joint event. The linked algebraic lemmas check the stated interval and physical-state probabilities under their displayed marginal inequalities.

This bridge does not verify the graph-specific hypotheses of the reference theorem or construct the full cover-down step.
\end{lemma}
\begin{proof}\uses{bp:preparation-balanced-roots}\leanok
Project only outcomes in the joint event having positive mass; removing zero-mass outcomes does not change its probability. The exact sigma-law projection identifies the retained law. Apply the retained existence theorem to the projected event and choose the positive-mass extension it supplies.
\end{proof}

\begin{lemma}[Chosen vortex with an additional retained event]\label{bp:chosen-vortex-extra}
\lean{Oberwolfach.marked_chosen_vortex_joint_endpoint_record_fields_with_extra,Oberwolfach.marked_preparation_reference_length_threshold}
\uses{bp:vortex-terminal-capacity}
\leanok
The chosen-vortex construction retains the all-level reference and endpoint conclusions when an arbitrary retained event of probability at least $99/100$ is specified after the vortex and its fixed weight families. A sufficiently large even marked length also admits the indices and density margin used by that reference construction.
\end{lemma}
\begin{proof}\uses{bp:vortex-terminal-capacity}\leanok
Choose the vortex and its scale as in the linked construction and use the all-level reference theorem with its additional-event input. The length comparison is an elementary threshold calculation.
\end{proof}

\begin{lemma}[Reference-compatible seeded root estimates]\label{bp:seeded-reference-roots}
\lean{Oberwolfach.raw_expanded_preparation_with_reference_extension,Oberwolfach.raw_expanded_families_with_reference,Oberwolfach.raw_expanded_selected_families_with_reference,Oberwolfach.raw_expanded_seeded_balanced_reference_all_roots}
\uses{bp:reference-positive-extension,bp:preparation-balanced-roots}
\leanok
Under the displayed input hypotheses, a retained reference principle that accepts every $99/100$-probability extra event can be combined with the fixed ordinary and marked seed formulas. One positive-mass enriched outcome then has the reference property, the basic preparation estimates and all four balanced root estimates. A bounded set of ordinary lengths may be selected after the order threshold.

The chosen-vortex theorem and this statement have not yet been assembled at the full manuscript caller. The hypotheses needed for the actual cover-down and factor steps remain open.
\end{lemma}
\begin{proof}\uses{bp:reference-positive-extension,bp:preparation-balanced-roots}\leanok
Project the joint good event to retained outcomes with a positive-mass extension. Apply the reference principle to this event, select an extension, and use the fixed-weight root and balance estimates on that same outcome.
\end{proof}

\begin{lemma}[Cover-down deletion witness and vortex step]\label{bp:vortex-cover-witness-adapter}
\lean{Oberwolfach.marked_vortex_valid_step_from_cover_witness}
\uses{bp:vortex-constructed-valid-step}
\leanok
A cover-down witness with the displayed exact remainder identity and incidence bounds supplies the next valid vortex state and the corresponding physical layer decomposition. The statement does not construct that witness from the reference data.
\end{lemma}
\begin{proof}\uses{bp:vortex-constructed-valid-step}\leanok
The exact remainder identity identifies the witness deletion set with the set used in the next-state theorem; the existing conservation result then gives the layer decomposition.
\end{proof}

\begin{lemma}[Factor components and long-preparation marginals]\label{bp:actual-factor-marginals}
\lean{Oberwolfach.long_preparation_factor_family_actual_marginals,Oberwolfach.long_preparation_factor_family_seed_coefficients,Oberwolfach.long_preparation_family_retained_lengths,Oberwolfach.long_preparation_seed_indices,Oberwolfach.long_preparation_retained_coefficient_sum,Oberwolfach.long_preparation_uniform_seed_density,Oberwolfach.long_preparation_residual_cutoff}
\leanok
For a family of two-factors in the long-component case, classify vertices by their actual component lengths and retain a short length at a hub when its vertex fraction exceeds the fixed cutoff. The linked results prove the discarded-fraction bound, the original two-colour interval marginals and physical-state probabilities, and the fixed hub coefficient balances. Every globally retained short length occurs in a factor and belongs to the allowed set. The uniform positive density bound is chosen from the initial density lower bound, independently of the factor family.

These are deterministic inputs to the preparation argument. The common-event construction and its use in the full cover-down and factor steps remain to be assembled.
\end{lemma}
\begin{proof}\leanok
Partition each factor's vertices by component length, bound the discarded short fractions by their number times the cutoff, and substitute these fractions in the displayed marginal and coefficient formulas. The finite coefficient sum may be restricted to the globally retained indices because all other terms vanish.
\end{proof}

\begin{lemma}[Long and short factor classes]\label{bp:factor-class-reduction}
\lean{Oberwolfach.factor_long_or_abundant_short,Oberwolfach.factor_preparation_class,Oberwolfach.factor_class_pruning_bound,Oberwolfach.factor_retained_class_density}
\leanok
Each prescribed two-factor either has at least half its vertices on cycles longer than the fixed cutoff or has a short cycle length containing at least a $1/(2K)$ fraction of its vertices. The assignment takes the least such length and only assigns allowed lengths. Removing classes smaller than their respective cutoffs removes at most the sum of those cutoffs, and each remaining class has the corresponding density lower bound.
\end{lemma}
\begin{proof}\leanok
Count the vertices in each component-length class, apply the pigeonhole principle, and sum the disjoint class sizes.
\end{proof}

\begin{lemma}[Sequential removal of small classes]\label{bp:small-class-removal}
\lean{Oberwolfach.factor_sequential_removal,Oberwolfach.factor_small_removal_margin,Oberwolfach.factor_small_removal_integer_margin,Oberwolfach.factor_small_class_removal,Oberwolfach.factor_sequential_even_removal,Oberwolfach.factor_small_class_even_removal,Oberwolfach.factor_decomposition_single_copy,Oberwolfach.factor_small_class_from_decompositions,Oberwolfach.factor_pruned_family_removal}
\leanok
Assuming the stated single-copy embedding at every intermediate regular host, the small factors can be removed successively. Their copies and the remainder partition the original edges, the remaining degree is exactly the original degree minus twice the number removed, and the stated degree margin persists.

The even-degree formulation matches the manuscript and derives the single-copy input from a decomposition into repeated copies. That density-one preparation decomposition remains an open dependency of the main theorem.
\end{lemma}
\begin{proof}\leanok
Induct over the prescribed list, deleting the edges of the copy supplied at each step. At each vertex, the degree falls by exactly two.
\end{proof}

\begin{lemma}[Chosen reference law and joint extension]\label{bp:chosen-reference-extension}
\lean{Oberwolfach.marked_retained_reference_principle,Oberwolfach.marked_reference_positive_extension,Oberwolfach.long_preparation_reference_densities,Oberwolfach.long_preparation_initial_sizes,Oberwolfach.long_preparation_uniform_family_seed_coefficients}
\uses{bp:reference-positive-extension,bp:actual-factor-marginals}
\leanok
The selected retained outcome has the same sigma-law as the reference input of the joint sampling theorem. A reference property supplied uniformly for high-probability retained events therefore has a positive-mass extension on the joint experiment whenever the prescribed joint event has the stated probability. The linked calculations give the density and coefficient inputs from the regular host and actual factor family.

Instantiating the reference property with all chosen-vortex outputs and discharging every cover-down input remain open.
\end{lemma}
\begin{proof}\uses{bp:reference-positive-extension,bp:actual-factor-marginals}\leanok
Identify the retention law, project the joint event to positive-mass retained outcomes, and apply the reference property on that projection.
\end{proof}

\begin{lemma}[Protected partition for the actual remaining classes]\label{bp:pruned-protected-partition}
\lean{Oberwolfach.factor_small_class_budget,Oberwolfach.factor_remaining_family_size,Oberwolfach.factor_remaining_class_density_sum,Oberwolfach.factor_retained_class_actual_density,Oberwolfach.factor_pruned_family_parameters,Oberwolfach.factor_nonempty_class_indices,Oberwolfach.factor_nonempty_class_count_bound,Oberwolfach.factor_pruned_class_indices,Oberwolfach.protected_partition_bounded_classes,Oberwolfach.factor_pruned_protected_partition}
\leanok
For the actual family left after removing small classes, the residual degree is even and positive; the number of remaining factors is half this degree. The nonempty classes can be indexed by $\operatorname{Fin}(k)$ with $k\le K+1$, and their densities sum to one and have the stated fixed lower bound. Consequently the protected-partition event applies with one order threshold chosen before the family and class count.

This event has not yet been intersected with the class-specific preparation events under a single common exposure.
\end{lemma}
\begin{proof}\leanok
Count each remaining class and use the small-class budget. Take the maximum of the finitely many order thresholds, one for each possible positive class count, then apply the protected-partition theorem to the resulting density profile.
\end{proof}

\begin{lemma}[Coupled categorical class refinements]\label{bp:class-refinement-law}
\lean{Oberwolfach.categorical_refinement_class_marginal,Oberwolfach.finiteLaw_sigma_map_first,Oberwolfach.categorical_refinement_product_views}
\leanok
Refining a single categorical class choice with a conditional finite law preserves the original class colouring and gives each fixed class its exact refined marginal. The joint construction retains the dependence between different classes.

The conditional distributions and projection to the manuscript's physical and interval outcomes still have to be instantiated.
\end{lemma}
\begin{proof}\leanok
Sum the joint masses over the conditional states and use the product law over physical coordinates.
\end{proof}

\begin{lemma}[One shared physical and raw experiment for marked classes]\label{bp:marked-shared-raw-law}
\lean{Oberwolfach.marked_physical_option_refinement,Oberwolfach.marked_class_physical_refinement,Oberwolfach.marked_class_physical_product,Oberwolfach.class_shared_raw_product,Oberwolfach.marked_class_shared_raw_law}
\uses{bp:class-refinement-law}
\leanok
A single independent categorical class colour and a single fair orientation are sampled for each physical edge. Conditional on the class colour, its five physical states are sampled with the manuscript's normalized probabilities. Independent class-local interval variables are then sampled. For each marked class, the resulting physical and interval variables have exactly that class's expanded raw wheel law. The class-colour projection retains the original categorical colouring.

The simultaneous high-probability reference and preparation events, and the analogous ordinary-class views, still have to be established on this common experiment.
\end{lemma}
\begin{proof}\uses{bp:class-refinement-law}\leanok
Compute the five conditional masses, include the outside state with its original orientation, and use the product law on the disjoint physical and interval coordinates.
\end{proof}

\begin{lemma}[Class-specific retention after a common raw exposure]\label{bp:class-common-retention}
\lean{Oberwolfach.class_shared_retention_map,Oberwolfach.class_projected_retention_map,Oberwolfach.class_common_retained_inputs}
\uses{bp:marked-shared-raw-law}
\leanok
If each class view has its prescribed raw law and its collapse has the selected reference law, then sampling each class's retention kernel after the common raw exposure gives each class its prescribed retained-reference marginal. The joint law keeps the common raw outcome and does not assert independence between classes.
\end{lemma}
\begin{proof}\uses{bp:marked-shared-raw-law}\leanok
Project a dependent product of retention kernels to one class and then apply the exact collapse law.
\end{proof}

\begin{lemma}[A common raw exposure for ordinary and marked classes]\label{bp:mixed-common-raw-exposure}
\lean{Oberwolfach.marked_global_ordinary_view,Oberwolfach.marked_class_shared_ordinary_law}
\uses{bp:marked-shared-raw-law}
\leanok
The common categorical colouring and class-local interval samples have the exact ordinary raw law for every class as well as the expanded marked raw law for every class. The projection to class colours remains the original categorical colouring. The class views are not asserted to be independent.
\end{lemma}
\begin{proof}\uses{bp:marked-shared-raw-law}\leanok
Project a categorical class indicator to its Bernoulli law, and combine this physical view with the same class-local interval coordinates used for the marked view.
\end{proof}

\begin{lemma}[Uniform failure tolerance for the marked reference event]\label{bp:marked-uniform-failure}
\lean{Oberwolfach.marked_joint_reference_all_levels_uniform_failure}
\leanok
The joint all-level reference event has the stated probability bound with an arbitrarily small positive failure tolerance. Its geometric upper bound on the hierarchy parameter is chosen before the tolerance, preserving the parameter order needed to take a union over a bounded number of classes. The event retains the original root weights, fixed domains and additional-event error.
\end{lemma}
\begin{proof}\leanok
Choose the geometric parameters from the uniform balance, cut, configuration and fixed-weight concentration estimates, and then intersect their events using the requested failure tolerance.
\end{proof}

\begin{lemma}[High-probability ordinary preparation inputs]\label{bp:ordinary-high-prob-inputs}
\lean{Oberwolfach.ordinary_uniform_sampled_long_inputs_probability}
\leanok
The ordinary sampled profile, associated cycle profiles, rows and absorber-hub lower bounds hold simultaneously with failure probability at most $1/n+2/n^2$ under the assumptions of the ordinary preparation input theorem.
\end{lemma}
\begin{proof}\leanok
Apply its extension property to the complement of this conjunction and use the complement identity for the finite probability law.
\end{proof}

\begin{lemma}[One event for a bounded number of classes]\label{bp:class-common-event-budget}
\lean{Oberwolfach.class_common_event_union,Oberwolfach.class_common_failure_budget,Oberwolfach.class_simultaneous_events}
\uses{bp:pruned-protected-partition,bp:mixed-common-raw-exposure,bp:marked-uniform-failure,bp:ordinary-high-prob-inputs}
\leanok
For a fixed bound on the number of classes, a failure tolerance and an order threshold can be chosen so that a partition event and all class events hold at one positive-mass outcome, provided their exact class projections and stated probability bounds have been established. No independence between these events is assumed.

The full manuscript-specific class events, the retained reference marginals, and all their numerical substitutions must still be assembled into this interface.
\end{lemma}
\begin{proof}\leanok
Choose the failure budget from the fixed class bound and apply the union bound to the complement events on the common law.
\end{proof}

\begin{lemma}[Vanishing failure for the expanded marked preparation event]\label{bp:expanded-preparation-invcard}
\lean{Oberwolfach.raw_expanded_preparation_basic_plus_extra_invCard,Oberwolfach.raw_expanded_marked_retained_projection_probability,Oberwolfach.raw_expanded_preparation_with_retained_event_invCard}
\leanok
For the stated fixed domains, the seven basic marked events and any selected retained event hold together with failure at most $7/n$ plus the retained event's failure probability. The conditional retention is taken after the expanded physical sample and projects to the original selected retained law. The estimates do not assert uniformity over arbitrary domains chosen after sampling.
\end{lemma}
\begin{proof}\leanok
Use the inverse-order concentration bounds for the seven events, project the retained law exactly, and apply the union bound without an independence assumption.
\end{proof}

\begin{lemma}[Ordinary triangle input event]\label{bp:triangle-high-prob-inputs}
\lean{Oberwolfach.ordinary_uniform_sampled_triangle_inputs_probability,Oberwolfach.ordinary_class_event_budget}
\leanok
Under the ordinary triangle preparation assumptions, the sampled triangle profiles, cycle profiles and hub lower bounds hold simultaneously with failure at most $1/n+2/n^2$. This failure fits the common class-event budget for sufficiently large order.
\end{lemma}
\begin{proof}\leanok
Apply the existing extension statement to the complement of the conjunction and use the complement identity.
\end{proof}

\begin{lemma}[Pruned classes on one retained experiment]\label{bp:pruned-common-retained-event}
\lean{Oberwolfach.mixed_class_concrete_raw_views,Oberwolfach.class_retained_simultaneous_events_200,Oberwolfach.class_reference_basic_budget,Oberwolfach.mixed_class_concrete_retained_common_event,Oberwolfach.factor_pruned_common_retained_event}
\uses{bp:pruned-protected-partition,bp:expanded-preparation-invcard,bp:triangle-high-prob-inputs,bp:marked-uniform-failure,bp:ordinary-high-prob-inputs}
\leanok
For the actual pruned factor classes, the protected-partition event and all supplied class preparation events occur at one positive-mass outcome of the common categorical physical exposure with class-specific interval sampling and retention, whenever the stated per-class probability bounds hold. The failure tolerance and order threshold are chosen before the factor family and class count. The partition conclusion remains uniform over the permitted later choices of protected edges.

The ordinary view in this lemma records whole-class membership, rather than the selected ordinary layer. The corrected selected-layer projection is given in Lemma~\ref{bp:selected-ordinary-class-exposure}. The class-specific preparation probabilities must still be derived for all the manuscript's actual inputs, and the effects of the partition's bounded changes must be carried through their consumers. This statement alone does not establish the common preparation event needed for the main theorem.
\end{lemma}
\begin{proof}\leanok
Use the exact mixed-class projections and the protected partition bound for the pruned classes. Project the dependent retention law to each class and apply the common-event union bound.
\end{proof}

\begin{lemma}[One retained marked reference and preparation event]\label{bp:marked-reference-expanded-joint}
\lean{Oberwolfach.mixed_class_marked_retention_projection,Oberwolfach.mixed_class_ordinary_retention_projection,Oberwolfach.mixed_class_marked_preparation_uniform_invCard,Oberwolfach.marked_initial_index_card_bounds,Oberwolfach.marked_reference_and_expanded_preparation_uniform,Oberwolfach.marked_reference_expanded_preparation_and_extra_uniform}
\uses{bp:marked-uniform-failure,bp:expanded-preparation-invcard}
\leanok
For a fixed vortex and the displayed density, geometry and interval hypotheses, the selected all-level reference event and the seven expanded fixed-domain preparation events hold together under the original start-wise conditional retention law with failure at most $100\chi+107/n$. The geometric hierarchy bound is selected before the arbitrary positive $\chi$; the order threshold for the expanded preparation estimates is independent of the additional-event error. A further supplied high-probability event on this same sample can be intersected with an additive error.

The construction of the vortex and verification of the input hypotheses at every actual class remain to be supplied.
\end{lemma}
\begin{proof}\leanok
Derive the index-set size bounds from the host degree and density, project the enriched retained experiment to the selected law, and intersect the stated events.
\end{proof}

\begin{lemma}[Original seed formulas on the common marked sample]\label{bp:marked-reference-seed-families}
\lean{Oberwolfach.mixed_approximate_families_common,Oberwolfach.mixed_class_complete_failure_budget,Oberwolfach.approximate_seed_family_bounds,Oberwolfach.interval_spoke_marginals_at_most_one,Oberwolfach.marked_reference_with_approximate_families_uniform,Oberwolfach.marked_reference_with_concrete_seed_families_uniform}
\uses{bp:marked-reference-expanded-joint}
\leanok
Under the displayed fixed host-weight and coefficient bounds, the selected reference event, the seven expanded preparation events, and both finite families of approximate root-omitted load estimates hold on the same marked retained outcome with failure at most $100\chi+200/n$. The marked and ordinary candidate weights are the original normalized seed formulas, including the supplied hub coefficients; the estimates do not recompute them after sampling.

The host weights and coefficients still have to be supplied for the actual factor family, and the resulting events must be connected to the protected partition and the cover-down inputs.
\end{lemma}
\begin{proof}\leanok
Bound the displayed seed coefficients from the original cycle weights, intersect the two potential-load families with the reference event, and allocate the failure budget before the order threshold.
\end{proof}

\begin{lemma}[Fixed indices for the actual long-factor family]\label{bp:actual-family-fixed-seeds}
\lean{Oberwolfach.pad_indexed_functions_on,Oberwolfach.pad_indexed_functions_off,Oberwolfach.pad_indexed_functions_sum,Oberwolfach.retained_coefficient_padding,Oberwolfach.preparation_padded_weights_from_allowed_threshold,Oberwolfach.preparation_actual_padded_family_with_marginals,Oberwolfach.marked_preparation_reference_length,Oberwolfach.preparation_actual_reference_ready}
\uses{bp:actual-factor-marginals,bp:marked-reference-seed-families}
\leanok
For a prescribed family satisfying the stated long-component, density and regular-host hypotheses, the selected ordinary cycle weights, marked successor weight, original hub coefficients, and two-colour physical and interval probabilities have the displayed bounds and balance identities. The short-length index set is fixed before the factor family is selected: unused lengths are padded by zero, while every retained weight and its exact edge sums are unchanged. The marked length can be chosen to satisfy the all-level reference theorem's successor-density margin.

These inputs have not yet been substituted into the joint marked reference and preparation theorem at every actual class.
\end{lemma}
\begin{proof}\leanok
Use the positive cycle weights for the retained allowed lengths and extend them by zero; use the actual cycle fractions for the hub coefficients. The retained-length sum agrees with the padded sum, and the established density estimates give the fixed seed floor.
\end{proof}

\begin{lemma}[Retained-length roots and the actual retention law]\label{bp:actual-retention-and-roots}
\lean{Oberwolfach.expanded_approximate_seed_families_imply_balanced_roots,Oberwolfach.padded_approximate_seeds_balanced_roots,Oberwolfach.factor_pruned_actual_retention_common_event,Oberwolfach.mixed_class_preparation_probability}
\uses{bp:pruned-common-retained-event,bp:actual-family-fixed-seeds}
\leanok
A fixed-index approximate-load event restricts to the retained lengths and gives the four stated balanced root estimates when the displayed edge identities and coefficient balances hold. On the common pruned-class experiment, marked classes use the original start-wise uniform retention, and ordinary classes use a singleton retention law. The protected partition is simultaneous with any class events whose individual probability bounds have been supplied.

The proof has not yet discharged those bounds for all actual classes or constructed the compatible cover-down records.
\end{lemma}
\begin{proof}\leanok
Remove the zero coefficients outside the retained set, apply the root-omitted load identities, and project the dependent retention laws to the corresponding class.
\end{proof}

\begin{lemma}[Marked reference event with a chosen vortex]\label{bp:chosen-vortex-concrete-seeds}
\lean{Oberwolfach.marked_reference_concrete_seed_event_with_vortex,Oberwolfach.positive_vortex_beta_choice,Oberwolfach.marked_reference_concrete_seed_event_choose_beta,Oberwolfach.long_reference_fixed_scales}
\uses{bp:marked-reference-seed-families,bp:actual-family-fixed-seeds}
\leanok
Under the displayed regular-host, fixed-scale and seed-weight hypotheses, a vortex and preserving permutation can be chosen before the marked sampling so that the fixed-domain reference and concrete root-omitted load events hold together with the stated failure bound. The positive vortex parameter is selected inside the proof, and the long-factor sampling floors are fixed before the actual graph density.

The actual padded factor-family inputs have not yet been substituted into this event at every class; the common protected-partition event and cover-down obligations remain open.
\end{lemma}
\begin{proof}\leanok
Choose the vortex and the permutation first, select the hierarchy parameter below both permitted bounds, and apply the retained-sample event with the original fixed weights and coefficients.
\end{proof}

\begin{lemma}[Two factors in a nearly regular rooted host]\label{bp:near-regular-rooted-two-factors}
\lean{Oberwolfach.augmentation_degree_upper,Oberwolfach.near_regular_augmentation_bounds,Oberwolfach.two_near_regular_cycle_factors_at_hub,Oberwolfach.rooted_two_dense_graphs_exist,Oberwolfach.subset_explicit_complement_degrees,Oberwolfach.rooted_near_regular_sample,Oberwolfach.near_regular_host_rooted_two_factors}
\leanok
In a sufficiently large nearly regular graph whose common degree centre is above the relevant cycle-decomposition threshold, one can choose a bounded set disjoint from a given bounded root set and obtain edge-disjoint cycle factors on the new set and on its union with the roots, with no edge of the latter factor joining two roots. The size and divisibility of the new set are fixed before the ambient graph. The degree estimates at the sampled vertices and roots hold on the same sample.

The common-centre hypotheses remain to be verified for actual callers. The conclusion does not imply that each cycle meets the root set at most once, and therefore does not by itself supply the stronger relative-absorber factor interface or the general minimum-degree factor theorem.
\end{lemma}
\begin{proof}\leanok
Use one uniform rooted sample and intersect the lower-degree events for the graph and its complement. Apply the almost-regular factor theorem first on the new set and then after deleting its factor and the edges inside the root set; the fixed root bound controls both degree losses.
\end{proof}

\begin{lemma}[Root-separated factors with bounded conflicts and deletions]\label{bp:root-separated-conflict-factors}
\lean{Oberwolfach.root_separated_cycle_packing,Oberwolfach.root_packing_remaining_divisible,Oberwolfach.induced_bounded_factor_inputs,Oberwolfach.near_regular_root_separated_factors,Oberwolfach.rooted_two_dense_conflict_free,Oberwolfach.rooted_near_regular_sample_conflict_free,Oberwolfach.near_regular_host_conflict_free_root_separated_factors,Oberwolfach.regular_subgraph_common_neighborhood_conflict_free_factors,Oberwolfach.root_separated_factors_physical,Oberwolfach.root_separated_factors_ambient,Oberwolfach.regular_subgraph_common_neighborhood_physical_factors}
\uses{bp:near-regular-rooted-two-factors}
\leanok
The near-regular two-factor construction can be strengthened so that every cycle of the factor covering the roots contains at most one root. A single fixed-size sample is also independent in a given bounded-degree conflict graph when the prescribed roots are independent. The sample size is chosen before the conflict-degree bound. In a regular host above the cycle-decomposition threshold, the construction applies inside two nearly complete neighbourhoods after excluding a small forbidden set, and remains valid after a bounded-degree deletion of rim edges. The stated physical lift preserves cycle supports and edge-disjointness.

These statements require the displayed near-regular or regular host, neighbourhood-size, divisibility, and root-independence assumptions. Their hypotheses have not yet been checked at all manuscript callers, and they do not establish the general minimum-degree-only factor theorem or the marked absorber.
\end{lemma}
\begin{proof}\leanok
First select disjoint cycles, each through one root and no other root. Their union has size divisible by the cycle length, so the unused fresh vertices have divisible size. Complete this packing by the almost-regular factor theorem, then delete its edges and apply the theorem again on the fresh set. To impose conflict-freeness, intersect the two degree events for the graph and its complement with the conflict-free event on one uniform rooted sample.
\end{proof}

\begin{lemma}[Uniform ordinary hub trades and their assembly]\label{bp:uniform-ordinary-hub-assembly}
\lean{Oberwolfach.regular_host_spoke_hub_factors_uniform,Oberwolfach.regular_host_spoke_hub_trade_uniform,Oberwolfach.regular_host_spoke_hub_trade_uniform_root_bound,Oberwolfach.regular_host_indexed_hub_trades_uniform,Oberwolfach.assemble_indexed_ordinary_hub_trades,Oberwolfach.regular_host_indexed_hub_transformer_uniform}
\uses{bp:root-separated-conflict-factors}
\leanok
For a regular host above the ordinary cycle-decomposition threshold, the permitted neighbourhood error is chosen before the number and size of the roots. Under the displayed common-spoke, root-independence, bounded-conflict and bounded-forbidden-edge hypotheses, the construction chooses simultaneous separated hub trades, preserving both wheel decompositions and the signed spoke identities. Given the previously constructed wheel families and their stated resource identity, these trades assemble to a physical resource transformer; every new rim edge meets a chosen fresh set unless it was already present in the supplied positive family.

The required previous wheel families and their identity are explicit inputs to this statement. This does not yet verify their construction for all manuscript callers, the more general pseudorandom-host hypotheses, the marked absorber, or Theorem~1.1.
\end{lemma}
\begin{proof}\leanok
Fix the near-regularity tolerance before the root bound. Choose conflict-free fresh sets successively, accounting for the neighbourhoods of previously chosen sets. Apply the exact resource identity to combine the two sides of the hub switches with the supplied wheel families.
\end{proof}

\begin{lemma}[Actual class densities and a common physical vortex]\label{bp:actual-class-shared-vortex}
\lean{Oberwolfach.preparation_actual_reference_ready_fixed_scales,Oberwolfach.factor_pruned_initial_index_densities,Oberwolfach.vortex_extend_fixed_vertices,Oberwolfach.shared_vortex_with_permutation_exists,Oberwolfach.factor_pruned_shared_vortex}
\uses{bp:pruned-protected-partition,bp:chosen-vortex-concrete-seeds}
\leanok
For the actual nonempty classes remaining after the prescribed small classes are removed, the initial index densities are at least $2\min\{\kappa_L,\kappa_S\}$ and at most one. A common physical vertex vortex and a single preserving permutation can be selected, with a separate nested index sequence for each actual class. The construction is valid for any common terminal scale above the selected lower bound, so it can be made compatible with other independently established size thresholds.

The marked length and its weight bounds may be fixed before the positive density floor is selected; the reference scales are then fixed before the host, allowed set and factor family. These results do not yet supply the class-specific preparation events on the joint exposure, the cover-down level records, or the fixed-set theorem.
\end{lemma}
\begin{proof}\leanok
The pruning estimates give the density bounds for the subtypes representing the actual classes. Choose one vertex vortex and extend its index sets separately for each class. Since every extended vortex has the same underlying vertex vortex, a permutation preserving the shells of one preserves those of all. The reference scale bounds follow from the fixed positive density floor and the already constructed actual-family weight bounds.
\end{proof}

\begin{lemma}[Ordinary transformer from prescribed paths in a regular host]\label{bp:ordinary-transformer-from-paths-regular}
\lean{Oberwolfach.factor_pruned_actual_density_identities,Oberwolfach.regular_host_indexed_hub_transformer_uniform_named_subgraph,Oberwolfach.ordinary_transformer_from_paths_regular_host_succ,Oberwolfach.ordinary_transformer_from_paths_regular_host}
\uses{bp:actual-class-shared-vortex,bp:uniform-ordinary-hub-assembly}
\leanok
The density assigned to a retained class is exactly its initial index density, and its product with the remaining degree is twice the number of indices in that class. For an ordinary wheel transformer in a regular host above the cycle-decomposition threshold, the prescribed paths and wheel data can be combined with the uniform hub trades when the displayed near-complete spoke-neighbourhood and conflict conditions hold. The resulting transformer has the asserted physical support and avoids the specified old rim and spoke resources.

The regular-host and near-complete-spoke hypotheses are explicit here. Their validity for the manuscript's general pseudorandom hosts has not been shown, and this lemma does not establish the minimum-degree factor theorem or the marked absorber.
\end{lemma}
\begin{proof}\leanok
Use the prescribed paths to construct the vertex-wheel stage and its signed resource identity. Apply the uniform indexed hub construction on fresh vertices outside the forbidden set, and combine its switches with the vertex-wheel stage. The identities for class densities follow from the cardinality of the actual retained-class subtype and positivity of the remaining degree.
\end{proof}

\begin{lemma}[First vortex sets fit the protected-partition budget]\label{bp:actual-vortex-partition-budget}
\lean{Oberwolfach.vortex_first_level_total_bound,Oberwolfach.vortex_first_level_partition_budget,Oberwolfach.factor_pruned_shared_vortex_partition_budget}
\uses{bp:actual-class-shared-vortex,bp:pruned-protected-partition}
\leanok
If the common vortex contraction parameter is chosen sufficiently small in terms of the fixed class bound, the sum of the sizes of the first physical vertex sets for all actual retained classes meets the deterministic inner-set budget of the protected partition. The result holds for every terminal scale above the chosen lower bound.

This supplies the set-size hypothesis only; the common sampled event and the per-class protected absorbers are still to be constructed.
\end{lemma}
\begin{proof}\leanok
The first vortex set has size at most $\varepsilon n$. There are at most $K+1$ retained classes, so the sum is at most $(K+1)\varepsilon n$. Choose $\varepsilon$ to meet the stated budget.
\end{proof}

\begin{lemma}[The selected ordinary layer in the common exposure]\label{bp:selected-ordinary-class-exposure}
\lean{Oberwolfach.expanded_ordinary_selected_physical_law,Oberwolfach.expanded_ordinary_selected_wheel_law,Oberwolfach.marked_class_shared_selected_ordinary_law,Oberwolfach.mixed_class_selected_raw_views,Oberwolfach.mixed_class_selected_retained_common_event,Oberwolfach.factor_pruned_selected_common_retained_event,Oberwolfach.factor_pruned_selected_actual_retention_common_event,Oberwolfach.factor_pruned_vortex_selected_retention_common_event,Oberwolfach.selected_class_state_implies_membership}
\uses{bp:actual-vortex-partition-budget,bp:pruned-protected-partition}
\leanok
The original categorical physical exposure has a separate selected ordinary state in each class. Its indicator has probability $2p$, remains inside that class's colour graph and has the original independent interval-coordinate law. The marked view retains the complete categorical state and direction. With these corrected views, the actual class vortices and their common permutation may be chosen before exposure, and a single retained outcome satisfies the protected partition event and all class events provided the displayed per-class probability estimates hold.

The older whole-class membership projection has probability $\varrho$ and cannot be substituted for the selected layer. The per-class probability estimates and the resulting decompositions still require proof.
\end{lemma}
\begin{proof}\leanok
Project each class's expanded categorical state to its selected ordinary state and sum its two directions. The global class colour is unchanged, and the class-local interval variables are unchanged. Apply the finite union bound and the protected-partition event on the same exposure, followed by the original conditional retention laws.
\end{proof}

\begin{lemma}[Ordinary sampling parameters from the prescribed factors]\label{bp:ordinary-factor-sampling-parameters}
\lean{Oberwolfach.factor_pruned_class_hub_sizes,Oberwolfach.factor_pruned_reference_seed_lower_bounds,Oberwolfach.ordinary_discarded_short_frequency,Oberwolfach.ordinary_frequency_marginal_bounds,Oberwolfach.ordinary_actual_frequency_marginals,Oberwolfach.ordinary_factor_frequency_partition,Oberwolfach.ordinary_factor_actual_marginals,Oberwolfach.ordinary_interval_marginals_feasible,Oberwolfach.ordinary_class_physical_parameters,Oberwolfach.ordinary_class_reservation,Oberwolfach.ordinary_factor_family_sampling_parameters,Oberwolfach.ordinary_class_assigned_sampling_parameters,Oberwolfach.ordinary_uniform_selected_density_floor}
\uses{bp:actual-class-shared-vortex,bp:selected-ordinary-class-exposure}
\leanok
The actual retained classes have linear-size index sets. For a class assigned an abundant short cycle length, the assignment supplies the reserved fraction required by the construction. Removing that fraction and discarding short lengths below the specified cutoff gives the stated bounds on the marked, ordinary and unassigned index marginals. These marginals satisfy the interval feasibility inequality, and their selected, approximate and unassigned physical probabilities have the original class density and balance identities.

This establishes the sampling parameters, not the subsequent approximate packing or full preparation event.
\end{lemma}
\begin{proof}\leanok
Partition the vertices of each prescribed factor by component length. Bound the discarded short frequencies by the number of possible lengths times the cutoff; the discarded long fraction contributes at most one further cutoff. The remaining inequalities follow from the frequency identities. Summing the per-index balance and using $\varrho D=2|U|$ gives the physical balance and approximate-density lower bound.
\end{proof}

\begin{lemma}[Ordinary host and absorber inputs on one sample]\label{bp:ordinary-host-and-absorber-common}
\lean{Oberwolfach.ordinary_full_host_and_absorber_inputs_common,Oberwolfach.ordinary_full_host_and_triangle_inputs_common,Oberwolfach.ordinary_combined_failure_below_class_budget}
\uses{bp:ordinary-factor-sampling-parameters}
\leanok
Under the stated dense almost-regular host and fixed-parameter hypotheses, including the separate triangle degree requirement, the ordinary full-host potential estimates and absorber-input estimates hold on the same raw physical and interval outcome with failure probability at most $1/n+3/n^2$. Both use the same cycle weights chosen before the interval sample and permutation are supplied. This error is bounded by the class-event error budget for $n\ge1$.

The approximate packing, adaptive list choices and final ordinary preparation event remain to be established on this common exposure.
\end{lemma}
\begin{proof}\leanok
Choose one family of cycle weights from the host, then apply the two high-probability results with that family and the same sampling law. Intersect their events by the union bound; no independence between the events is used.
\end{proof}

\begin{lemma}[Actual ordinary-class sampling inputs, uniformly in the density]\label{bp:ordinary-actual-assigned-class-event}
\uses{bp:ordinary-factor-sampling-parameters,bp:ordinary-host-and-absorber-common}
\lean{Oberwolfach.ordinary_assigned_class_combined_event_uniform}
\leanok
For a fixed cycle length and the displayed numerical parameters, the floor for the class density, the host weights and the order threshold are chosen before the actual assigned factor class is supplied. The original cycle frequencies determine the marked-start and spoke probabilities; for these probabilities, the ordinary full-host and absorber-input estimates hold on the same raw sample with failure probability at most $1/n+3/n^2$. The statement concerns one assigned ordinary class and a single fixed length. The simultaneous class events, approximate packing and complete preparation still require their own inputs.
\end{lemma}

\begin{lemma}[Uniform ordinary inputs for a bounded range of lengths]\label{bp:ordinary-bounded-length-uniform-events}
\uses{bp:ordinary-actual-assigned-class-event,bp:triangle-high-prob-inputs}
\lean{Oberwolfach.ordinary_uniform_events_all_bounded_lengths}
\leanok
For a fixed upper bound on the lengths and a fixed positive density floor, one can choose common positive input scales and one order threshold for the assigned-class host and absorber-input event at every triangle or longer cycle length in the range. The source weights and auxiliary constants may depend on the length, but are chosen before the actual class and factor family. When a requested host-degree bound exceeds the order, the corresponding statement is proved vacuously. This gives separate probability estimates for the classes; their simultaneous realization with the partition and the remaining preparation events is not asserted here.
\end{lemma}

\begin{lemma}[Ordinary position capacity from the displayed bounds]\label{bp:ordinary-position-capacity}
\uses{bp:actual-residual-prescribed-set}
\lean{Oberwolfach.ordinary_natural_position_capacity_for_large_orders,Oberwolfach.residual_subtype_glue_inputs}
\leanok
Under the stated coefficient, assigned-edge, prescribed-set and remainder-size estimates, a sufficiently large order guarantees the exact natural-number capacity required for prescribed-position selection, including the ceiling in the edge-assignment bound. The order threshold depends only on the displayed fixed constants. The subtype gluing hypotheses follow from the equality between the old domain's intersection with the remainder and the prescribed set. Establishing the displayed estimates and the old embedding itself at the preparation caller remains separate.
\end{lemma}

\begin{lemma}[Extending lists: path and load ingredients]\label{bp:extending-list-path-pieces}
\lean{Oberwolfach.extending_list_after_neighbourhood_changes,Oberwolfach.extending_list_two_arc_choice,Oberwolfach.extending_list_one_arc_choice,Oberwolfach.extending_list_free_choice,Oberwolfach.extending_list_embed_internal_path,Oberwolfach.extending_list_embed_one_sided_path,Oberwolfach.prescribed_embedding_saturated_vertices,Oberwolfach.prescribed_embedding_margins_from_hypotheses,Oberwolfach.prescribed_embedding_greedy_slack}
\leanok
For an extending list, the stated one- and two-arc neighbourhood bounds give choices avoiding a specified set of vertices and a bounded-degree physical forbidden graph. Under the explicit cardinality budgets in the Lean statements, a path can be embedded between two distinct specified endpoints, or forwards or backwards from one specified endpoint. The saturation estimate bounds the number of excluded vertices by the square-root scale, and the numerical margins follow from the hypotheses of the manuscript's prescribed-embedding lemma.

These results do not yet embed a family of paths and cycles with prescribed images. In particular, a cycle with exactly one prescribed vertex has coinciding boundary images, so the manuscript proof's assertion that these boundary images are always distinct requires a separate case.
\end{lemma}
\begin{proof}\leanok
Use the extending-neighbourhood inequalities and subtract the forbidden vertices and physical neighbours. Induct on the number of internal path vertices, and count saturated vertices using the total number of free images. The numerical inequalities follow from the displayed square-root bound.
\end{proof}

\begin{lemma}[Uniform ordinary-class inputs and partial event budget]\label{bp:ordinary-uniform-partial-inputs}
\lean{Oberwolfach.factor_pruned_uniform_sampling_inputs,Oberwolfach.ordinary_same_sample_retained_class_budget,Oberwolfach.ordinary_frequency_weight_allocation,Oberwolfach.ordinary_factor_frequency_weight_allocation}
\uses{bp:ordinary-host-and-absorber-common,bp:ordinary-factor-sampling-parameters}
\leanok
Positive class-density and index-size floors can be fixed from the pruning constants before the actual prescribed factors are chosen. The retained short-cycle frequencies supply the displayed nonnegative fixed weight coefficients and their exact mass identity. The already proved ordinary host and absorber-input events fit the retained-class failure budget on the same sample, under their explicit hypotheses.

This does not establish the approximate-packing event or the complete preparation event.
\end{lemma}
\begin{proof}\leanok
Use the class-size estimates and density identities, then substitute the actual retained frequencies into the fixed-coefficient identities. Apply the retention projection and the union bound to the two already established events.
\end{proof}

\begin{lemma}[A cycle with one prescribed vertex]\label{bp:pinned-cycle-correction}
\lean{Oberwolfach.extending_list_embed_cycle,Oberwolfach.extending_list_embed_pinned_cycle,Oberwolfach.prescribed_embedding_cycle_slack,Oberwolfach.prescribed_embedding_cycle_budget,Oberwolfach.prescribed_embedding_label_budget,Oberwolfach.prescribed_embedding_assigned_and_unembedded_budget,Oberwolfach.prescribed_embedding_pinned_cycle_of_source_budgets,Oberwolfach.prescribed_embedding_load_update,Oberwolfach.degree_two_image_union_bound,Oberwolfach.bounded_degree_two_union_from_occurrences,Oberwolfach.oriented_physical_edge_unique}
\uses{bp:extending-list-path-pieces}
\leanok
A cycle of length at least three with exactly one prescribed vertex can be completed from an extending list using the given forbidden-edge graph, provided the explicit source-size, saturation and forbidden-degree bounds in the Lean statement hold. The original square-root hypotheses provide the necessary numerical margin. The source-budget lemma counts the assigned vertices and the free vertices of the component as disjoint parts of the same source factor.

This addresses a gap in the manuscript proof: in this case the free path has the same prescribed predecessor and successor, so their images need not be distinct. The checked single-cycle construction does not establish the simultaneous family embedding or its dynamic forbidden-edge and load invariants.
\end{lemma}
\begin{proof}\leanok
Choose a free image adjacent to the prescribed vertex. It is distinct from that prescribed image; embed the remaining path between these distinct endpoints. The source and saturation bounds supply enough available labels, while the specified degree bound controls forbidden physical edges.
\end{proof}

\begin{lemma}[Invariants for prescribed-image embeddings]\label{bp:prescribed-image-invariants}
\lean{Oberwolfach.partial_injective_image_degree,Oberwolfach.prescribed_embedding_used_degree_bound,Oberwolfach.prescribed_embedding_used_degree_from_partial_images,Oberwolfach.prescribed_embedding_total_size,Oberwolfach.prescribed_embedding_saturated_actual_occurrences,Oberwolfach.prescribed_embedding_load_history_bound,Oberwolfach.occurrence_index_projection_card_le,Oberwolfach.prescribed_embedding_disjoint_update,Oberwolfach.prescribed_embedding_injective_update,Oberwolfach.prescribed_embedding_update_preserves}
\uses{bp:pinned-cycle-correction}
\leanok
An injective partial image of a graph of maximum degree two has degree at most two at an image vertex and zero outside its image. For a family of such images, the specified bounds on prescribed and new vertex occurrences yield the stated bound on the union of used physical edges, including the separately forbidden graph. The source-size estimate bounds the sum of the source-factor sizes. Separately, a bound on the actual number of free occurrences gives the stated bound on the number of saturated host vertices; their inclusion in the source family still has to be checked at the caller. A sequence of choices that avoids saturated vertices preserves the stated per-vertex load bound.

One update at a previously unassigned source vertex preserves the earlier images and injectivity if its new image is unused. Adding physical edges to one factor preserves pairwise disjointness if those edges avoid all other factor edge sets. These are invariant lemmas under their displayed hypotheses; constructing every update from the manuscript's paths and cycles remains open.
\end{lemma}
\begin{proof}\leanok
Bound the degree in each partial image by its source degree, and sum over indices in which the vertex occurs. Count saturated vertices from the total number of actual free occurrences. For one update, use the unused-image and edge-disjointness hypotheses directly; the load assertion follows by induction over the updates.
\end{proof}

\begin{lemma}[Saturation and physical edges of selected paths]\label{bp:source-family-path-updates}
\lean{Oberwolfach.prescribed_embedding_saturation_from_source_family,Oberwolfach.prescribed_path_physical_edges_avoid,Oberwolfach.prescribed_one_sided_physical_edges_avoid,Oberwolfach.prescribed_path_disjoint_update,Oberwolfach.prescribed_one_sided_path_disjoint_update}
\uses{bp:prescribed-image-invariants,bp:extending-list-path-pieces}
\leanok
If the set of free occurrences is contained in the disjoint union of the source vertex sets, the original bounds on their sizes and number give the required upper bound on the number of saturated host vertices. The physical edges of a selected two-ended path, or of a selected forward or backward one-ended path, avoid the forbidden graph. Thus adding these edges to one factor preserves physical edge-disjointness, provided the forbidden graph contains the current physical edges of every factor.

The construction of the partial maps is carried out below for a supplied finite list of vertex-disjoint source components. Its identification with the input graphs and the original callers remains to be checked.
\end{lemma}
\begin{proof}\leanok
Count the source vertices in the dependent disjoint union and apply the saturation estimate. For a path, each consecutive physical edge avoids the forbidden graph by the path construction; the asserted update then follows from pairwise disjointness before the step.
\end{proof}

\begin{lemma}[Five component embedding steps]\label{bp:five-component-steps}
\lean{Oberwolfach.prescribed_free_path_component_step,Oberwolfach.prescribed_free_cycle_component_step,Oberwolfach.prescribed_one_ended_component_step,Oberwolfach.prescribed_two_ended_component_step,Oberwolfach.prescribed_pinned_cycle_component_step}
\uses{bp:source-family-path-updates,bp:pinned-cycle-correction}
\leanok
Under the explicit extending-list, cardinality, forbidden-degree and resource-containment hypotheses in the Lean statements, each of the following can be embedded in one step: a path with no prescribed boundary, a cycle with no prescribed vertex, a path with one prescribed boundary, a path with two distinct prescribed boundaries, and a cycle with exactly one prescribed vertex. Each construction extends the partial injection, preserves old images, uses allowed unsaturated images, and adds exactly the mapped physical edges while retaining pairwise edge-disjointness and avoiding the reserved graph.

The family-state construction below establishes these hypotheses at every step of a supplied component ordering.
\end{lemma}
\begin{proof}\leanok
Choose the host path or cycle from the extending list and match its ordered vertices to the free source vertices. The new images are distinct and avoid the old images. Map the source physical edges to the selected physical edges and apply the disjointness update.
\end{proof}

\begin{lemma}[One family component update]\label{bp:family-component-invariant}
\lean{Oberwolfach.prescribed_family_update,Oberwolfach.prescribed_physical_image_preserved,Oberwolfach.prescribed_component_occurrence_count,Oberwolfach.prescribed_occurrence_union,Oberwolfach.prescribed_component_load_update,Oberwolfach.prescribed_actual_component_load_step,Oberwolfach.prescribed_family_component_invariant}
\uses{bp:prescribed-image-invariants}
\leanok
Suppose one factor's injective partial map is extended on a disjoint set of source vertices and all new images avoid the currently saturated labels. If the displayed containment and edge-avoidance hypotheses hold, the family update preserves every partial injection and all earlier images, the bound on actual free occurrences at each host vertex, pairwise disjointness of physical edge sets and avoidance of the reserved graph by the added edges. Previously embedded source edges retain their physical images.

The bound survives choosing a whole component against one saturated set because injectivity makes each host vertex occur at most once in that component. The saturated set must be recomputed before the next step. This is a single-step invariant, not the full prescribed-family embedding.
\end{lemma}
\begin{proof}\leanok
Count the new occurrences by the indicator of membership in the component's image. Combine this identity with the saturation rule, the preservation of earlier assignments, and the physical-edge update.
\end{proof}

\begin{lemma}[The forbidden graph from the assigned source vertices]\label{bp:actual-forbidden-source-graph}
\lean{Oberwolfach.prescribed_family_domains_update,Oberwolfach.prescribed_family_free_index_count,Oberwolfach.prescribed_induced_degree_le,Oberwolfach.prescribed_forbidden_graph_from_partial_images,Oberwolfach.prescribed_actual_forbidden_graph,Oberwolfach.prescribed_source_edges_image_contained,Oberwolfach.prescribed_source_edge_update}
\uses{bp:family-component-invariant}
\leanok
For each factor, restrict its degree-two source graph to the assigned vertices and map it by the current partial injection. The union of these physical image graphs with the reserved graph has the stated degree bound, derived from the original prescribed-occurrence bound and the maintained free-occurrence bound. It contains the reserved graph and every edge in those partial image graphs.

A source edge with both endpoints assigned maps into the corresponding partial image graph. Extending the map without changing old assigned images preserves the physical images of old source edges. The full construction must still maintain that all tracked edges come from assigned source edges and that the chosen components cover the source factors.
\end{lemma}
\begin{proof}\leanok
Restriction cannot increase source degrees. Count, separately, factors in which a host vertex is a prescribed image or a free image, and apply the degree bound for the union of the image graphs. The source-edge assertions follow from the definitions of an induced graph and a mapped physical edge.
\end{proof}


\begin{lemma}[Family state and quantitative resources]\label{bp:family-state-resources}
\lean{Oberwolfach.prescribed_induced_family_state_initial,Oberwolfach.prescribed_family_state_step,Oberwolfach.prescribed_family_state_resources}
\uses{bp:family-component-invariant,bp:actual-forbidden-source-graph}
\leanok
The prescribed maps and induced source edges initialise a family state with exact physical edge images. A component update preserves this state. Under the original numerical bounds and the maintained state invariants, the saturated labels, the union of the partial physical image graphs and the reserved graph satisfy the cardinality and forbidden-degree bounds required by all five component embeddings.
\end{lemma}
\begin{proof}\leanok
Use the prescribed maps for the initial state. Bound the actual free occurrences and the degree of the union of the current image graphs. The stated numerical inequalities then give the component budgets, and the family update preserves all the recorded invariants.
\end{proof}

\begin{lemma}[An ordered path split at prescribed vertices]\label{bp:path-segment-schedule}
\lean{Oberwolfach.prescribed_path_segments_free_vertices,Oberwolfach.prescribed_path_segment_boundaries,Oberwolfach.prescribed_path_segment_boundaries_distinct,Oberwolfach.prescribed_path_segment_infix,Oberwolfach.prescribed_path_segment_inherits_chain,Oberwolfach.prescribed_path_segments_edge_union,Oberwolfach.prescribed_ordered_right_sided_edges,Oberwolfach.prescribed_empty_segment_edges}
\leanok
Splitting an ordered path at its prescribed vertices lists exactly its free vertices. The boundary vertices of each segment are prescribed, and two boundaries of a segment of a path with distinct vertices are distinct. Each segment is contiguous and inherits the source path's directed arcs. The union of its segment edge sets is exactly the original path edge set, including edges between consecutive prescribed vertices. The reversed ordering used for a segment with only a right boundary gives the same undirected source edges. A segment without free vertices contributes only edges of the initially prescribed induced source graph.
\end{lemma}
\begin{proof}\leanok
Induct along the ordered path, splitting at each prescribed vertex. The edge identity follows by splitting the ordered edge set at the same vertex; reversing a segment reverses the order of each edge's endpoints without changing the undirected edge.
\end{proof}

\begin{lemma}[Actual family component advances]\label{bp:family-original-advances}
\lean{Oberwolfach.prescribed_family_free_path_advance,Oberwolfach.prescribed_family_free_cycle_advance,Oberwolfach.prescribed_family_one_ended_advance,Oberwolfach.prescribed_family_two_ended_advance,Oberwolfach.prescribed_family_pinned_cycle_advance,Oberwolfach.prescribed_family_segment_advance}
\uses{bp:five-component-steps,bp:family-state-resources,bp:path-segment-schedule}
\leanok
Given an actual ordered source component of any of the five types, the original numerical hypotheses, the extending-list property and a valid current family state supply an embedding and a new valid state. A nonempty scheduled path segment can therefore be processed with its exact source edge set and directed target arcs, including the right-boundary case read in reverse. The step preserves the images of all previously assigned vertices and places its newly assigned vertices in their specified list. A cycle with one prescribed vertex is handled by the separate pinned-cycle step.

The finite induction below combines these advances for a supplied list of source components. Identification of that list with the source graphs and verification of the original lemma's callers remain to be proved.
\end{lemma}
\begin{proof}\leanok
Construct the current forbidden graph and saturation set from the family state. Apply the appropriate path or cycle embedding, then update the source and physical edges and the free-occurrence counts. For an ordered path segment, use its boundaries and inherited source arcs to select the applicable component step.
\end{proof}


\begin{lemma}[Induction over a supplied family of source components]\label{bp:actual-family-induction}
\lean{Oberwolfach.prescribed_segment_schedule_induction,Oberwolfach.prescribed_ordered_path_family_advance,Oberwolfach.prescribed_cycle_family_advance,Oberwolfach.prescribed_source_piece_advance,Oberwolfach.prescribed_source_family_induction}
\uses{bp:family-original-advances}
\leanok
Suppose a finite list specifies, for each factor, ordered source paths and cycles satisfying the stated path or cycle conditions, with vertex-disjoint pieces within that factor. Starting from a valid family state whose unprocessed free vertices are unassigned, the induction processes every listed piece. It preserves all previously assigned images, the common free-image load bound, pairwise disjointness of physical edge sets and the reserved graph condition. Every new image lies in its factor's specified list. Every listed source arc either has two initially prescribed endpoints or maps to an arc of the target digraph, and every listed source vertex and physical edge is recorded in the final state.

This result is conditional on the displayed list and its hypotheses. The ordered-piece embedding theorem below supplies the final degree and load conclusions. Transport from an arbitrary presentation of the manuscript source graphs, the two preparation callers, and the principal theorem remain to be checked.
\end{lemma}
\begin{proof}\leanok
Split a path at its prescribed vertices and process its segments in order, treating an empty segment as an initially prescribed edge. For a cycle, either use the free-cycle or pinned-cycle step, or split it into two paths at prescribed vertices; rotation preserves the directed and undirected edges. Iterate over the supplied list, using vertex-disjointness to keep the unprocessed vertices free and preservation of previous images to retain the conclusions from earlier pieces.
\end{proof}


\begin{lemma}[Prescribed embeddings for ordered source pieces]\label{bp:ordered-prescribed-embedding}
\lean{Oberwolfach.prescribed_source_degree_from_components,Oberwolfach.prescribed_source_model_piece_valid,Oberwolfach.prescribed_source_model_arc_graph,Oberwolfach.prescribed_family_embedding_from_ordered_pieces}
\uses{bp:actual-family-induction}
\leanok
Let the source factors be specified by a finite list of ordered directed paths, isolated vertices and directed cycles of length at least three, whose pieces are vertex-disjoint within each factor. Construct each source vertex set, directed arc relation and underlying graph directly from its pieces. Under the numerical hypotheses, prescribed-image hypotheses, target extending-list conditions and reserved-graph bound displayed in the Lean statement, there are injective extensions with the required lists, directed arcs, pairwise physical edge-disjointness and reserved-edge avoidance. The graph formed by the union of the physical images has degree at most $5\sqrt\delta\,n$, and each physical vertex occurs as a free image at most $\sqrt\delta\,n+1$ times.

The checked theorem uses one common finite type for source vertices and an explicit list of pieces. The result below transports it to a separate finite vertex type for each factor. Identification with the two preparation applications remains to be verified. This result does not prove the main theorem.
\end{lemma}
\begin{proof}\leanok
The listed paths and cycles determine their source vertices and arcs and have no loops. Vertex-disjointness gives at most one predecessor and successor at each vertex, so the underlying graph has maximum degree at most two. Initialise the family state with all prescribed edges and apply the finite component induction. The exact source-edge and free-occurrence identities yield the claimed edge-disjointness and load estimates; the degree estimate follows from the prescribed and free occurrence bounds.
\end{proof}


\begin{lemma}[Prescribed embeddings with separate source types]\label{bp:dependent-prescribed-embedding}
\lean{Oberwolfach.prescribed_lifted_pieces_apart,Oberwolfach.prescribed_lifted_source_vertices,Oberwolfach.prescribed_lifted_source_arc_membership,Oberwolfach.prescribed_lifted_initial_edges,Oberwolfach.prescribed_dependent_free_load_transfer,Oberwolfach.prescribed_lifted_full_edge_images,Oberwolfach.prescribed_dependent_family_embedding}
\uses{bp:ordered-prescribed-embedding}
\leanok
The ordered-piece embedding result remains valid when each factor has its own finite source vertex type. Its prescribed edge sets, physical edge images, directed arcs and free-image occurrence counts are those of the original factors. In particular, the degree, load, list, injectivity and edge-avoidance conclusions hold with the original numerical hypotheses.

The Lean statement presents the source factors by ordered, vertex-disjoint path and cycle pieces. The substitutions and preceding constructions in the ordinary and long preparation proofs still require verification, as do the other dependencies of Theorem~1.1.
\end{lemma}
\begin{proof}\leanok
Embed the source types in their dependent sum and relabel each listed piece. This preserves disjointness within each factor, the source graph and its arcs, prescribed maps and prescribed physical edges. Apply the common-type theorem and restrict the resulting maps to the appropriate fibers. The physical edge sets and the count of free occurrences agree under this restriction.
\end{proof}


\begin{lemma}[Prescribing a matching and isolated images]\label{bp:prescribed-positions}
\lean{Oberwolfach.prescribed_arc_count_eq_edge_count,Oberwolfach.prescribed_positions_with_arc_count}
\leanok
For a finite oriented source graph with underlying maximum degree at most two, the prescribed-position conclusion holds under the original edge-count and size inequalities. The new positions induce exactly the chosen directed matching and map bijectively onto its endpoints and the prescribed isolated labels, while the original injection is preserved. The directed edge count equals the edge count of the underlying graph. The path-and-cycle source representation and the two actual preparation substitutions still need to be connected to this statement.
\end{lemma}
\begin{proof}\leanok
Greedily choose source arcs away from the closed neighbourhood of the prescribed vertices and previously chosen endpoints; then choose mutually nonadjacent single positions. The original inequalities imply the selection budgets. Assign the selected arcs and isolated vertices to the target labels, preserving their directions; the separation proves that no additional arc is induced.
\end{proof}


\begin{lemma}[Prescribed positions followed by the family embedding]\label{bp:positions-family-embedding}
\lean{Oberwolfach.prescribed_piece_no_reverse,Oberwolfach.prescribed_source_model_degree,Oberwolfach.prescribed_dependent_position_family,Oberwolfach.prescribed_positions_and_family_embedding,Oberwolfach.prescribed_family_avoids_old_packing,Oberwolfach.prescribed_augmented_occurrence_bound,Oberwolfach.prescribed_assigned_matchings_disjoint,Oberwolfach.prescribed_padding_vertices}
\uses{bp:prescribed-positions,bp:dependent-prescribed-embedding}
\leanok
For separately typed source factors represented by vertex-disjoint directed paths and cycles, the source orientation and degree conditions follow from the component structure. Under the displayed position, list, occurrence and edge-disjointness hypotheses, the selected positions extend the initial injections and induce exactly the prescribed directed matchings. The ensuing family embedding covers their physical edges and the prescribed isolated images, and retains the physical disjointness, forbidden-edge and free-occurrence bounds. Assignment of each physical bad edge to a single factor gives disjoint matching edge sets, and its endpoint occurrences are bounded by its degree in the bad-edge graph. Previously used edges and the reserved graph remain avoided under the stated disjointness conditions. Omitted isolated source vertices can be added as singleton paths without changing arcs or edges.

The numerical and random-sampling inputs of the actual ordinary and marked preparation proofs have not yet been supplied to this interface. In particular, this result does not establish the cover-down construction or Theorem~1.1.
\end{lemma}
\begin{proof}\leanok
Select the source positions and use the exact directed-to-physical edge image identity. Apply the family embedding theorem to the resulting prescribed maps. Bound matching-endpoint occurrences by choosing an incident bad edge for each index; edge assignment makes this choice injective. Add singleton paths for any otherwise unlisted isolated vertices.
\end{proof}


\begin{lemma}[Embedding after the bad-edge assignment]\label{bp:assigned-family-embedding}
\lean{Oberwolfach.prescribed_positions_and_family_embedding_all_vertices,Oberwolfach.prescribed_assigned_arcs_physical,Oberwolfach.prescribed_assigned_arcs_are_matchings,Oberwolfach.prescribed_assigned_arc_card,Oberwolfach.prescribed_assigned_family_embedding}
\uses{bp:positions-family-embedding}
\leanok
Given the displayed source-piece, position-capacity, list and degree bounds, assign each physical bad edge to one factor so that each assigned edge class is a matching, and orient the edges in the target digraph. The construction embeds the separately typed factors, covers every bad edge and all prescribed isolated labels, and keeps their physical edge sets pairwise disjoint and disjoint from the reserved graph and the previously used packing. Each assigned directed matching has exactly as many edges as its physical assigned class. No extra source-coverage premise is needed: any omitted isolated vertices are represented by singleton paths.

The actual ordinary and marked preparation constructions must still establish the displayed source, assignment, numerical and common-event inputs before this result can be applied. The main theorem remains unproved.
\end{lemma}
\begin{proof}\leanok
Orient the assigned physical edges and apply the prescribed-position selection. Its image and endpoint identities imply the required multiplicity and physical disjointness conditions. Apply the family embedding theorem and use the exact initial and new physical edge sets to verify coverage and avoidance.
\end{proof}


\begin{lemma}[Assignment of bad edges to factors]\label{bp:graph-edge-assignment}
\lean{Oberwolfach.graph_edge_assignment_original,Oberwolfach.prepare_assigned_bad_edges}
\leanok
The edge-assignment lemma holds with the manuscript's list sizes, endpoint-forbidden bounds, maximum-degree bound and exact ceiling capacity. Its assignment is a function on the graph's edge set, including when the graph has no edges. For the preparation application, the chosen physical edges can be oriented in the supplied target digraph, and their assigned classes are matchings of the stated sizes.
\end{lemma}
\begin{proof}\leanok
Process the edges one at a time. Exclude forbidden indices at the two endpoints, indices assigned an incident edge, and indices already at capacity. The first two losses are bounded by the two endpoint bounds; the last is counted from the actual edges assigned so far. The manuscript's numerical inequalities leave an allowed index. The induction preserves all earlier choices and capacities.
\end{proof}

\begin{lemma}[Ordered source pieces for a two-factor]\label{bp:two-factor-piece-model}
\lean{Oberwolfach.two_factor_piece_model_exact}
\leanok
Every finite two-regular graph can be presented as a finite list of pairwise vertex-disjoint ordered cycle pieces. The corresponding source-piece vertex set is the whole vertex set, and its physical source graph is exactly the original graph. This does not yet construct the paths and cycles obtained after deleting the prescribed subgraphs in the preparation arguments.
\end{lemma}

\begin{lemma}[Residual graph after disjoint path and cycle deletions]\label{bp:residual-deletion-count}
\lean{Oberwolfach.two_factor_residual_graph_after_disjoint_deletions,Oberwolfach.two_factor_residual_position_edge_budget,Oberwolfach.disjoint_deleted_path_endpoints_survive}
\leanok
For a two-factor, delete the edges of vertex-disjoint simple paths and cycles and delete the path interiors and cycle vertices. Every edge incident with a deleted vertex is deleted, and the number of remaining edges plus the number of deleted paths equals the number of remaining vertices. If a prescribed set contains both endpoints of every deleted path, it has at least twice as many vertices as there are deleted paths, and these endpoints survive. Identifying the paths, the full prescribed set, and the ordered components of the remainder with those in the manuscript remains to be done.
\end{lemma}

\begin{lemma}[Residual graph and prescribed positions from extended paths]\label{bp:actual-residual-prescribed-set}
\uses{bp:residual-deletion-count}
\lean{Oberwolfach.residual_actual_graph_and_prescriptions}
\leanok
Given the stated disjoint two-edge-extended paths, selected segments and disjoint cycle components of a two-factor, construct the final paths and maximal selected-run subpaths, delete their edges and interiors and delete the selected components. For the resulting graph $Q$, surviving vertex set $R$ and prescribed set $A_0$, the formal statement proves $A_0\subseteq R$, its independence in $Q$, the exact prescribed-marker and residual-edge counts, and $|R|\le e(Q)+|A_0|$. It preserves the actual marker survival and maximal-run constructions. The path selection and quantitative preparation callers, and the remaining proof stages, are not yet established by this lemma.
\end{lemma}

\begin{lemma}[Fixed-orientation source pieces for residual factors]\label{bp:fixed-orientation-residual-pieces}
\lean{Oberwolfach.two_factor_family_fixed_orientation_for_all_subgraphs}
\leanok
For each member of a finite family of two-factors, one can fix an orientation into directed cycles before choosing any residual graphs. For every later subgraph and every surviving vertex set outside which that subgraph is isolated, the surviving vertices and edges admit a representation by disjoint ordered paths and cycles. The directed arcs in these pieces are exactly the original arcs whose physical edges survive. The construction also works for a supplied ordered family of oriented source cycles. Applying it to the manuscript's earlier partial embeddings and proving the sampling and numerical hypotheses remain separate obligations.
\end{lemma}

\begin{lemma}[Residual source model and prescribed positions on surviving vertices]\label{bp:residual-subtype-position-inputs}
\uses{bp:actual-residual-prescribed-set,bp:fixed-orientation-residual-pieces}
\lean{Oberwolfach.two_factor_family_residual_ready_for_positions,Oberwolfach.residual_actual_vertex_count,Oberwolfach.residual_actual_prescribed_set_bound,Oberwolfach.residual_actual_embedding_overlap,Oberwolfach.residual_actual_full_selection_comparison}
\leanok
The fixed orientations and disjoint residual pieces can be restricted to the actual surviving vertex types, with exactly the original surviving arcs and the induced residual graphs. For independent prescribed sets $A_0\subseteq R$ with $|R|\le e(Q)+|A_0|$, the prescribed-position arc-count condition follows. For the specified disjoint extended paths, selected runs, and deleted cycle components, the formal results also prove the exact remaining-vertex count, the stated bound on $|A_0|$, the comparison with the fully selected remainder, and that the selected-run and already embedded component domain meets $R$ exactly in $A_0$. The existence and compatibility of the earlier partial embeddings, the quantitative sampling hypotheses and the remaining main-theorem stages are separate obligations.
\end{lemma}


\begin{lemma}[List-based remainder completion and gluing]\label{bp:remainder-completion-gluing}
\lean{Oberwolfach.complete_remainder_with_image_and_shield,Oberwolfach.old_image_shield_from_lists,Oberwolfach.old_image_shield_from_forbidden,Oberwolfach.prescribed_subtype_complete_factor_glue,Oberwolfach.prescribed_glued_physical_edge_images,Oberwolfach.complete_remainder_and_glue_from_lists}
\uses{bp:graph-edge-assignment,bp:assigned-family-embedding}
\leanok
Under the explicit numerical, list, source-piece and graph-disjointness hypotheses in the Lean statement, the bad edges can be assigned and all prescribed factor remainders embedded. The embeddings cover every bad edge and prescribed singleton label, avoid both the reserved graph and the earlier packing, and glue injectively to the already embedded factor parts. The exact old and remainder physical edge images are preserved. The old-image exclusion can follow from the edge lists or the endpoint-forbidden sets; in both manuscript preparation applications the latter sets are empty and the lists supply the exclusion.

This is a deterministic conditional construction. The actual sampled lists, earlier partial maps, numerical bounds and the common event still have to be constructed and substituted. The original factor-edge partition, full preparation theorem, cover-down iteration and Theorem~1.1 remain unproved.
\end{lemma}
\begin{proof}\leanok
Apply the graph edge-assignment lemma, the position-selection theorem and the family embedding theorem. Retain the exact selected-image set and each selected endpoint's allowed edge. These identities show that the new images avoid the old image set. Transport the remainder from its own vertex type into the original factor and glue the two injections; the physical edge-image identities follow from their agreement on the respective domains.
\end{proof}

\begin{lemma}[Actual residual factors with assigned bad edges]\label{bp:actual-residual-assigned-gluing}
\uses{bp:residual-subtype-position-inputs,bp:remainder-completion-gluing}
\lean{Oberwolfach.ActualResidualData.graph_facts,Oberwolfach.ActualResidualData.old_domain_union_remaining,Oberwolfach.actual_factor_residuals_embed_and_glue,Oberwolfach.actual_factor_remainder_assigned_and_glued}
\leanok
Fix the orientations of a family of two-factors before choosing disjoint extended paths, selected runs and deleted cycle components. The actual surviving source graphs then admit the path-and-cycle models required for prescribed-position embedding. The previously embedded source domain includes the selected path interiors and meets each surviving domain exactly in the prescribed set. If the deleted cycle components have already been embedded, these two domains cover the whole source factor.

Under the explicit list, capacity, occurrence, earlier-injection and physical avoidance hypotheses in the Lean statement, the deterministic construction assigns every bad edge, embeds the actual remainders, protects the reserved and old physical edges, and glues the new and old maps to injections of the full source factors. This does not establish those hypotheses for the sampled preparation, nor does it show that the old embedded edges and new edges together give the complete decomposition. Those caller obligations and Theorem~1.1 remain open.
\end{lemma}

\begin{lemma}[Actual glued factor edge images]\label{bp:actual-glued-factor-edges}
\uses{bp:actual-residual-assigned-gluing}
\lean{Oberwolfach.ActualResidualData.source_edge_partition,Oberwolfach.ActualResidualData.glued_factor_edge_images,Oberwolfach.actual_piece_family_glued_physical_edges,Oberwolfach.actual_glued_factor_avoids_reserved_edges,Oberwolfach.actual_factor_remainder_complete_physical_family}
\leanok
For the actual selected paths and cycle components, the deleted source edges and residual source edges partition each original factor. Once the old and residual vertex maps have been glued, the union of their physical edge images is the image of the original factor. If the earlier factor-edge images are pairwise disjoint and contained in the old packing, the new full-factor images are pairwise disjoint and cover the assigned bad edges; if the old packing also avoids the reserved graph, the full-factor images avoid it. These are conditional conclusions: the earlier packing, partial maps and preparation estimates must still be constructed and shown to meet the displayed hypotheses.
\end{lemma}

\begin{lemma}[Spanning two-factor embeddings exhaust a regular host]\label{bp:factor-embeddings-exhaust-host}
\uses{bp:graphs}
\lean{Oberwolfach.regular_host_two_factor_edge_sets_exhaust,Oberwolfach.spanning_two_factor_embeddings_give_decomposition,Oberwolfach.even_regular_two_factor_embeddings_give_decomposition}
\leanok
Let $G$ be $D$-regular on $n$ vertices, where $D$ is even. If the $D/2$ prescribed spanning two-factors have injective vertex maps into $G$ whose edge images are contained in $G$ and are pairwise disjoint, then these images exhaust $E(G)$ and give a decomposition of $G$ into the prescribed factors. This is a finite counting and graph-isomorphism conclusion; the existence of the embeddings is not assumed to have been proved.
\end{lemma}

\begin{lemma}[Conditional decomposition from the actual residual construction]\label{bp:actual-residual-to-decomposition}
\uses{bp:actual-glued-factor-edges,bp:factor-embeddings-exhaust-host}
\lean{Oberwolfach.prescribed_piece_images_in_host,Oberwolfach.actual_glued_factor_images_in_host,Oberwolfach.actual_glued_residuals_give_factor_decomposition}
\leanok
Suppose that the earlier and residual maps of the actual factor remainders have been glued injectively, the earlier factor-edge images are pairwise disjoint and lie in a physical packing inside the regular host, and the new edge images are pairwise disjoint, avoid that packing, and follow directions whose underlying edges lie in the host. If every deleted cycle component has been embedded, then these data give the prescribed factor decomposition of the host. The construction of the old maps and packing, the directed sample, and all their required preparation and absorption properties is still open; this implication is not a proof of the main theorem.
\end{lemma}

\begin{lemma}[Actual retained long-class preparation inputs]\label{bp:pruned-long-class-inputs}
\uses{bp:actual-factor-marginals,bp:actual-family-fixed-seeds,bp:actual-class-shared-vortex}
\lean{Oberwolfach.long_class_assigned_frequency,Oberwolfach.long_assigned_factor_concrete_sampling,Oberwolfach.pruned_long_class_concrete_sampling,Oberwolfach.pruned_long_class_reference_inputs,Oberwolfach.pruned_long_class_preparation_reference_ready}
\leanok
If the long class remains after the prescribed small classes are removed, its members satisfy the long-component frequency required in the sampling construction. Their actual two-colour interval and physical-state probabilities obey the displayed balance and feasibility equations; the selected ordinary and marked probabilities have positive lower bounds chosen from the fixed class-density floor, independently of the realised class density. The regular graph remaining after pruning retains the required degree margin, so the original fixed-index cycle and successor weights and their reference inputs are supplied for this actual class. The asserted order of choosing constants is the one in the Lean statements. This does not yet establish the simultaneous marked and ordinary sampling event, construct the approximate packing, or prove the preparation proposition.
\end{lemma}

\begin{lemma}[Actual mixed-class laws on one vortex]\label{bp:pruned-mixed-class-laws}
\uses{bp:pruned-long-class-inputs,bp:actual-class-shared-vortex}
\lean{Oberwolfach.mixed_assigned_classes_sampling_witness,Oberwolfach.pruned_mixed_classes_sampling_witness,Oberwolfach.pruned_actual_mixed_vortex_common_event}
\leanok
For all nonempty classes remaining after pruning, the interval and spoke probabilities are the original formulas computed from the factors in that class. Their selected and approximate physical-state probabilities satisfy the exact class balance, the interval-law feasibility bounds, and the fixed positive density bounds in the Lean statements. The approximate and unassigned probabilities retain their original sums, so their weighted coefficients are not recomputed after sampling. The actual class vortices use one physical vortex and permutation, and the laws are realised on the same categorical exposure with conditional retention. If the stated class-event probability bounds hold, one outcome also has the protected-partition property and all those class events. Establishing the class-event bounds and constructing the preparation remain open.
\end{lemma}

\begin{lemma}[Actual ordinary-class threshold inputs]\label{bp:pruned-ordinary-class-inputs}
\uses{bp:actual-class-shared-vortex,bp:pruned-long-class-inputs}
\lean{Oberwolfach.pruned_actual_class_size_and_density,Oberwolfach.pruned_ordinary_class_threshold,Oberwolfach.ordinary_uniform_events_fixed_cycle_thresholds}
\leanok
Each retained class has linearly many indices, with its lower fraction fixed from the pruning cutoffs. A retained ordinary class has an allowed cycle length, and the actual regular graph after pruning remains above that cycle's threshold by the stated positive margin. For any fixed bounded range of ordinary lengths, including triangles, the corresponding event parameters and order threshold can be chosen from the fixed cycle thresholds, independently of the allowed set. These inputs do not yet supply the class events in the common exposure.
\end{lemma}

\begin{lemma}[Actual ordinary-class event substitutions]\label{bp:pruned-ordinary-class-events}
\uses{bp:pruned-ordinary-class-inputs}
\lean{Oberwolfach.pruned_ordinary_class_reference_inputs,Oberwolfach.pruned_ordinary_assigned_event_from_bound,Oberwolfach.pruned_triangle_assigned_event_from_bound}
\leanok
For an actual retained ordinary class, the fixed-length class-event estimate applies to its factors and the actual pruned regular graph, with the original fixed weights and interval probabilities. The triangle and longer-cycle estimates are handled separately. These statements use the corresponding fixed-length estimates as inputs, with their parameters and order threshold fixed before choosing the allowed set. They do not yet establish the joint preparation event for all classes.
\end{lemma}

\begin{lemma}[Class events on the mixed exposure]\label{bp:mixed-class-event-transfer}
\uses{bp:pruned-ordinary-class-events,bp:pruned-mixed-class-laws}
\lean{Oberwolfach.pruned_ordinary_retained_event_from_bound,Oberwolfach.pruned_triangle_retained_event_from_bound,Oberwolfach.mixed_ordinary_named_event_transfer,Oberwolfach.mixed_triangle_named_event_transfer,Oberwolfach.mixed_marked_assigned_law_identity,Oberwolfach.mixed_marked_assigned_retained_transfer,Oberwolfach.mixed_marked_assigned_retained_extra_transfer}
\leanok
The ordinary and triangle class estimates remain valid after their singleton retention step. The original ordinary product law is exactly the selected marginal of the mixed exposure, and the original marked product law is exactly its marked marginal with the five physical-state probabilities unchanged. Thus each displayed ordinary or marked event bound transfers to the corresponding retained mixed-class law, including the full error term for an additional marked event. The actual marked event bound, the additional packing and list events, and their simultaneous use in the preparation are not supplied by this transfer.
\end{lemma}

\begin{lemma}[Long-class weights and target estimates]\label{bp:long-class-uniform-d-targets}
\uses{bp:mixed-class-event-transfer,bp:pruned-long-class-inputs}
\lean{Oberwolfach.pruned_long_class_preparation_reference_ready_uniform_d,Oberwolfach.long_assigned_reference_recipe_identity,Oberwolfach.long_assigned_concrete_recipe_identity,Oberwolfach.raw_marked_all_targets_from_concentration,Oberwolfach.physical_ordinary_family_targets_from_concentration,Oberwolfach.physical_ordinary_frequency_targets_from_concentration}
\leanok
The cycle-weight bounds for the actual long class can be chosen before the interval length. The retained-frequency and reference-weight formulas agree with the original marked sampling recipe, including its approximate and unassigned densities. On any one sample satisfying the stated all-root estimates, the corresponding marked spoke and rim targets and the ordinary targets for all active lengths follow with their supplied frequency coefficients. The construction of the complete actual preparation and compatible iteration records remains open.
\end{lemma}

\begin{lemma}[Reserved spokes in the actual residual embedding]\label{bp:actual-reserved-spoke-assembly}
\uses{bp:actual-glued-factor-edges}
\lean{Oberwolfach.actual_reserved_matching_assignment_inputs,Oberwolfach.actual_residual_reservation_from_exact_resources,Oberwolfach.actual_factor_remainder_complete_with_reserved_spokes,Oberwolfach.actual_factor_remainder_complete_with_exact_reservation,Oberwolfach.actual_glued_factor_images_preserve_typed_rims,Oberwolfach.actual_factor_remainder_with_outer_reserved_lists,Oberwolfach.actual_reservation_canonical_assignment_exclusions,Oberwolfach.actual_factor_remainder_from_outer_codegrees,Oberwolfach.actual_outer_codegree_common_slack}
\leanok
The actual reserved ordinary spoke resources determine the vertex labels excluded from matching lists. Removing those labels at both endpoints loses at most twice the total number of reserved resources; with the displayed initial slack, the lists retain the size and occupied-image conditions required by the assignment. For the actual residual source graphs and their glued embeddings, the stated old-image, matching-list, inner-label and free-list exclusions preserve all reserved spokes. When the full glued factor images avoid the reserved rim graph, deleting those images also preserves the type and orientation of every reserved rim. The remaining starting estimates and compatibility of the resulting absorber with the cover-down records remain to be proved at the preparation caller.
\end{lemma}


\begin{lemma}[Common ordinary-spoke codegrees, including close pairs]\label{bp:all-pairs-reserved-spokes}
\uses{bp:actual-reserved-spoke-assembly}
\lean{Oberwolfach.raw_interval_two_labels_probability_lower,Oberwolfach.raw_interval_all_pair_spoke_codegree_event,Oberwolfach.raw_expanded_reserved_matching_with_event,Oberwolfach.raw_expanded_reserved_matching_high_probability,Oberwolfach.marked_initial_index_card_lower,Oberwolfach.mixed_marked_reserved_matching_event_uniform_d,Oberwolfach.raw_ordinary_spoke_survives_start_retention,Oberwolfach.mixed_retained_reserved_pairs,Oberwolfach.mixed_retained_reserved_lists,Oberwolfach.mixed_long_reserved_spoke_recipe}
\leanok
In the original interval experiment, every distinct pair has ordinary-spoke probability at least $5s_0^2/8$ at each hub, including pairs whose permutation windows meet. The corresponding lower codegree bound holds simultaneously for every pair, with the displayed exponential failure estimate. The codegree constant may be fixed before the interval length; the ordinary colour of the actual long-class recipe has density $1/100$. The event is measured under the same expanded physical law and persists under its conditional start retention. On this one event, the actual retained-spoke matching lists satisfy the displayed conclusions for every later choice obeying the loss and old-image bounds. These conclusions do not themselves construct the preparation or a compatible first cover-down record.
\end{lemma}

\begin{lemma}[Joint reference, approximate seeds and reserved spokes]\label{bp:marked-reference-reserved-spoke-caller}
\uses{bp:marked-reference-expanded-joint,bp:all-pairs-reserved-spokes}
\lean{Oberwolfach.marked_reference_with_approximate_and_reservation_uniform,Oberwolfach.marked_reference_with_concrete_seeds_and_reservation_uniform,Oberwolfach.marked_reference_concrete_reserved_seed_event_with_vortex,Oberwolfach.marked_reference_concrete_reserved_seed_event_choose_beta,Oberwolfach.mixed_reference_and_retained_reservation}
\leanok
The marked reference estimates, both approximate-packing seed families, and the all-pairs reserved-spoke event hold with the displayed common failure budget under the same expanded physical law and retention. The concrete seed formulas and the chosen vortex retain that event. A separate deterministic implication transfers its codegrees and matching-list conclusions to the retained ordinary spokes.
\end{lemma}

\begin{lemma}[First matching-list and ineligible-edge bounds]\label{bp:marked-first-matching-inputs}
\uses{bp:marked-reference-reserved-spoke-caller,bp:actual-reserved-spoke-assembly}
\lean{Oberwolfach.marked_initial_reserved_loss,Oberwolfach.marked_occupied_with_raw_approximate_intervals_avoid_outer,Oberwolfach.marked_first_matching_inputs,Oberwolfach.long_preparation_ineligible_degree,Oberwolfach.long_preparation_ineligible_edge_count,Oberwolfach.long_preparation_bad_graph_bounds,Oberwolfach.long_preparation_bad_graph_margins,Oberwolfach.long_required_physical_containment,Oberwolfach.long_unassigned_restricted_degree_from_row,Oberwolfach.long_required_vortex_row_margins,Oberwolfach.long_preparation_physical_unassigned_upper,Oberwolfach.marked_first_required_graph_inputs}
\leanok
For a bounded reserved set and a sufficiently small first vortex ratio, the first inner index set and reserved labels fit the codegree loss budget. If the previous images have the displayed interval and colour classification, they avoid the retained ordinary-spoke lists, so the first matching-list inputs follow. The ineligible physical graph has maximum degree at most $6d+2$ and at most $(3d+1)n$ edges. The categorical physical states imply the stated containment of the required graph when the remaining and selected graphs satisfy the displayed sample-identification conditions. The fixed-inner-set unassigned row estimate, without changing its error scale, and the unused approximate-edge degree bound give its degree and edge bounds and the matching-assignment degree margin. The same retained outcome then supplies the first matching-list inputs. Deriving the sample-identification conditions, previous-image classification, unused-packing bound, and remaining preparation and compatible absorber records from the actual constructions remains open.
\end{lemma}

\begin{lemma}[Sampled mixed resource matching and physical edges]\label{bp:sampled-mixed-physical-matching}
\uses{bp:fractional-resource-matching,bp:marked-reference-reserved-spoke-caller}
\lean{Oberwolfach.long_assigned_approximate_coefficients_sum_one,Oberwolfach.sampled_resource_degree_eq_pattern,Oberwolfach.mixed_sampled_resource_degree_from_targets,Oberwolfach.mixed_sampled_resource_codegree_bound,Oberwolfach.mixed_sampled_resource_matching,Oberwolfach.ordinary_raw_load_in_marked_universe,Oberwolfach.marked_retained_resource_mem_sample,Oberwolfach.markedSampledPhysicalEmbedding,Oberwolfach.marked_retained_candidate_physical_disjoint}
\leanok
The fixed marked and ordinary coefficients sum to one at the ordinary spokes and rims, under the displayed positivity and balance hypotheses. For one sampled outcome, restriction preserves the candidate weights and identifies each active resource degree with its root-omitted load. The ordinary candidate resources embed in the marked universe, preserving their loads and codegrees. Under the stated load, codegree, size and test-function estimates, the weighted matching theorem applies to the combined candidates. The retained typed resources embed into the physical edges, so disjointness of those resources gives physical disjointness. The construction of the actual combined candidate estimates and test functions, and the embedding of the resulting packing into the preparation, remain open.
\end{lemma}

\begin{lemma}[Retained seed targets and their combined degree]\label{bp:retained-mixed-seed-targets}
\uses{bp:sampled-mixed-physical-matching,bp:marked-reference-reserved-spoke-caller}
\lean{Oberwolfach.expanded_retention_restriction_probability,Oberwolfach.expanded_retention_restrictions_witness,Oberwolfach.expanded_retained_marked_spoke_positive_density,Oberwolfach.retained_marked_seed_load_stability,Oberwolfach.ordered_ordinary_spoke_rim_three_vertices,Oberwolfach.retained_ordinary_seed_pattern_stability,Oberwolfach.expanded_retained_marked_active_seed_bounds,Oberwolfach.retained_ordinary_active_seed_bounds,Oberwolfach.long_assigned_marked_seed_target,Oberwolfach.long_assigned_marked_named_seed_active_targets,Oberwolfach.long_assigned_ordinary_seed_target,Oberwolfach.long_assigned_ordinary_named_seed_active_targets,Oberwolfach.long_assigned_actual_mixed_active_loads,Oberwolfach.long_assigned_fraction_below_three_zero,Oberwolfach.long_assigned_mixed_sampled_unit_degree,Oberwolfach.mixed_sampled_resource_matching_active_targets,Oberwolfach.all_resource_sets_from_large_sets}
\leanok
The restriction event for retained starts has the same probability under the original interval law and under the common expanded physical experiment. It can therefore be intersected with any earlier common event, and a positive combined probability yields a positive-mass outcome. A retained marked spoke on such an outcome has positive supplied marked density. The displayed retention-stability estimates are auxiliary: the manuscript's approximate packing uses the raw colour-two intervals and raw ordinary values, while the final colour-one sample is trimmed. The exact long-class targets sum to one, and the formal length-one and length-two classes have zero contribution. The original assigned ordinary seed concentration gives its raw active-root target bound under the stated fractional-cycle equations, host-degree estimate and numerical budgets. The original assigned marked seed concentration gives the marked active-root target bound on a positive-mass outcome of the actual retained experiment, under its stated density and weight conditions. On one positive-mass retained outcome carrying both concentration events, the actual assigned mixed seed loads satisfy their targets under the original fractional-cycle equations, host-degree estimate and numerical budgets for retained lengths; discarded lengths contribute zero. It remains to obtain that positive-mass common outcome from the actual reference event, supply the host and numerical inputs, verify the remaining matching hypotheses, and construct the approximate packing and protected absorber records.
\end{lemma}

\begin{lemma}[Raw approximate resources and hub tests]\label{bp:raw-approximate-resource-hub-tests}
\uses{bp:sampled-mixed-physical-matching,bp:retained-mixed-seed-targets}
\lean{Oberwolfach.markedRawSampledPhysicalEmbedding,Oberwolfach.expanded_raw_marked_spoke_positive_density,Oberwolfach.expanded_raw_marked_active_seed_bounds,Oberwolfach.raw_ordinary_active_seed_bounds,Oberwolfach.expanded_raw_spoke_hub_counts_common_event,Oberwolfach.long_assigned_mixed_loads_from_seed_targets,Oberwolfach.long_mixed_seed_component_codegrees,Oberwolfach.typed_marked_seed_resource_size,Oberwolfach.ordered_ordinary_seed_resource_size,Oberwolfach.long_assigned_positive_seed_valid_length,Oberwolfach.mixed_sampled_resource_matching_positive_active_targets,Oberwolfach.ordered_ordinary_common_hub_spoke_count,Oberwolfach.common_marked_wheel_hub_spoke_count,Oberwolfach.mixed_restricted_hub_test_mass_lower,Oberwolfach.hub_test_mass_margin}
\leanok
The approximate candidate resources use the original raw colour-two and ordinary labels and the approximate physical rim states. Their active typed resources embed injectively into the physical edges. On a positive-mass expanded outcome, an active marked spoke has positive marked density, and the original marked and ordinary seed concentration estimates give their displayed active-root targets. Both raw spoke types have simultaneous hub-specific count estimates on the same expanded outcome. The marked and ordinary candidate seed weights have the displayed component codegree bounds, and genuine candidates have the displayed resource-size bounds. Zero-weight candidates may be excluded from the matching without changing its degree, codegree or test masses. For an individual family and hub, the displayed weighted-degree identity and spoke counts give the test-mass lower bound whenever their numerical hypotheses hold. The full construction of the source-specific joint event, verification of all numerical margins and test families, and the actual packing and protected absorber compatibility remain open.
\end{lemma}

\begin{lemma}[Common seed event and selected family frequencies]\label{bp:raw-common-seed-frequency}
\uses{bp:raw-approximate-resource-hub-tests}
\lean{Oberwolfach.expanded_mixed_seed_common_event,Oberwolfach.expanded_approximate_common_events_witness,Oberwolfach.long_active_typed_resource_card_le_cube,Oberwolfach.mixed_matched_family_hub_frequency,Oberwolfach.matching_frequency_margin,Oberwolfach.eventually_any_polynomial_exponential_failure_below,Oberwolfach.eventually_mixed_seed_failure_below,Oberwolfach.long_fixed_seed_uniform_bounds}
\leanok
For fixed seed weights satisfying the displayed size bounds, the marked root-omitted concentration estimate and all the specified ordinary seed estimates hold together on the raw approximate layer, with the sum of their explicit failure probabilities. This event can be joined to an earlier event of the same expanded-and-retained experiment. The selected-family test sum is exactly the number of selected candidates with the specified hub, and its error relative to the prescribed frequency is controlled by the unchanged seed-load, hub-count, and matching errors. For fixed lengths and uniform positive weight bounds, the explicit sum of marked and ordinary seed failure bounds tends to zero with a threshold independent of the particular weights and hub set. The explicit marked and ordinary seed formulas satisfy the displayed upper bounds under the existing parent-weight and spoke-probability hypotheses. Those hypotheses, the other numerical inequalities and the earlier event still have to be assembled for the manuscript caller; this result does not assert the existence of the completed packing.
\end{lemma}

\begin{lemma}[Retained lengths and the actual seed weights]\label{bp:retained-actual-seed-weights}
\uses{bp:raw-common-seed-frequency}
\lean{Oberwolfach.long_retained_padded_length_allowed,Oberwolfach.long_nonretained_seed_exact_target,Oberwolfach.long_ordinary_active_inputs_from_retained,Oberwolfach.long_assigned_actual_raw_law_parameters,Oberwolfach.long_retained_seed_weights_uniform_d,Oberwolfach.eventually_named_mixed_seed_failure_below}
\leanok
The indices whose retained frequency is nonzero at some hub have actual lengths in the allowed set. A padded index outside this retained set has identically zero ordinary target and load, so no cycle decomposition at that padded length is needed. With the preparation hypotheses, fixed parent and actual seed weights exist for the retained lengths, and the raw sampling parameters obey the probability constraints. Their uniform weight constants are chosen before the allowed set and factor family. These results do not yet combine all concentration events with the matching, the resulting partial factor images and the reserved absorber.
\end{lemma}

\begin{lemma}[Raw retained hub tests and active marked roots]\label{bp:raw-retained-hub-tests}
\uses{bp:retained-actual-seed-weights,bp:raw-common-seed-frequency}
\lean{Oberwolfach.long_retained_hub_tests_card_bound,Oberwolfach.long_retained_hub_coefficient_lower,Oberwolfach.raw_ordinary_common_hub_sample_count,Oberwolfach.long_retained_hub_tests_mass_from_raw,Oberwolfach.long_mixed_active_inputs_from_retained,Oberwolfach.expanded_marked_seed_event_to_active_targets}
\leanok
The fixed tests indexed by retained length--hub pairs form a polynomial-size family. Under the stated root-load and raw hub-count estimates, each has the explicit positive mass required by the weighted matching theorem, using the actual mixed wheel families and raw approximate outcome. The ordinary hub count in the typed resource state is exactly the count in the original raw law. The marked seed concentration event supplies its active-root targets on a positive-mass outcome; ordinary estimates are needed only for retained lengths. It remains to obtain all these hypotheses on one actual outcome, apply the matching with the original parameter hierarchy, and construct its factor images and protected absorber.
\end{lemma}

\begin{lemma}[Unused approximate rim edges]\label{bp:raw-approximate-unused-rim-degree}
\uses{bp:raw-approximate-resource-hub-tests}
\lean{Oberwolfach.long_approximate_edge_has_raw_resource,Oberwolfach.physical_rim_covered_eq_candidate_union,Oberwolfach.long_approximate_unused_degree_from_matching,Oberwolfach.mixed_restricted_hub_test_mass_deviation,Oberwolfach.long_required_vortex_rows_from_matching,Oberwolfach.mixed_matching_large_incident_rim_coverage}
\leanok
The weighted matching theorem applies simultaneously to precisely the incident rim-resource sets above its size threshold; smaller sets need no matching test. For the raw approximate sample, each eligible physical rim edge has a unique active typed resource. The union of the rim resources of selected candidates is exactly their physical rim-edge union. If their matching covers all but a proportion $\gamma$ of every incident resource set with at least $a|V|$ members, then every vertex has at most $\max\{a,\gamma\}|V|+6d+2$ unused approximate rim edges. The last term accounts for the explicitly defined ineligible graph. The resulting bound supplies the unused-edge input of the first-vortex required-graph estimates when the displayed parameter margins hold. The displayed hub-test deviation identity retains the fixed coefficients and original error scales. These are conditional estimates; it remains to construct the common matching event, verify the source-specific margins and identify its selected candidates with the partial factor images.
\end{lemma}

\begin{lemma}[Protection of the fixed absorber in the first two stages]\label{bp:raw-absorber-protection}
\uses{bp:raw-approximate-resource-hub-tests,bp:marked-first-matching-inputs}
\lean{Oberwolfach.selected_and_approximate_rim_graphs_disjoint,Oberwolfach.raw_approximate_spoke_avoids_retained_selected,Oberwolfach.raw_approximate_and_retained_selected_physical_disjoint,Oberwolfach.raw_sampled_physical_embedding_in_approximate_layer,Oberwolfach.physical_resource_to_finite_injective,Oberwolfach.fixed_absorber_selected_layer_subset,Oberwolfach.raw_matching_avoids_fixed_terminal_absorber,Oberwolfach.approximate_and_first_cover_protect_fixed_absorber,Oberwolfach.long_class_matching_avoids_fixed_absorber}
\leanok
The raw colour-two approximate resources are disjoint from the final retained colour-one resources, including after the actual trimming of marked starts. The physical images of every approximate resource therefore avoid the fixed absorber supplied in the selected host. Any physically disjoint first-cover packing made in the absorber-deleted selected host also avoids that same absorber. The latter statement is conditional on the existence of the packing in that host; it does not construct the first cover-down or identify the approximate matching with its factor images.
\end{lemma}

\begin{lemma}[Realizing sampled wheel candidates]\label{bp:raw-wheel-realizations}
\uses{bp:raw-approximate-resource-hub-tests,bp:raw-absorber-protection}
\lean{Oberwolfach.raw_marked_seed_candidate_lift,Oberwolfach.raw_restricted_marked_candidate_physical_lift,Oberwolfach.raw_marked_realization_properties,Oberwolfach.ordinary_seed_physical_resource_image,Oberwolfach.raw_restricted_ordinary_candidate_physical_lift,Oberwolfach.mixed_restricted_component_resources,Oberwolfach.mixed_positive_matching_physical_realizations,Oberwolfach.raw_ordinary_named_seed_active_targets}
\leanok
A restricted marked candidate has a wheel realization in the original raw approximate sample with the same tuple, hub and physical resource set, including the twisted closing edge. An ordinary candidate has exactly the physical rim and spoke resources represented by its restricted typed resources. Componentwise realizations combine without changing those resources. If their common physical map is injective, a positive-weight resource matching with nonempty candidates gives pairwise physically disjoint, distinct realizations. The named ordinary concentration event also yields its active-root target estimates under the displayed hypotheses. The source-specific realization is given below. Matching existence on the common event, the assignment of these wheels to partial factor images and all first-cover hypotheses remain open.
\end{lemma}

\begin{lemma}[Uniform ordinary densities in the long-class recipe]\label{bp:long-ordinary-density-lower}
\uses{bp:retained-actual-seed-weights}
\lean{Oberwolfach.long_assigned_actual_ordinary_density_lower,Oberwolfach.long_assigned_actual_marked_density_lower,Oberwolfach.long_assigned_actual_common_density_lower,Oberwolfach.long_assigned_actual_density_bounds}
\leanok
The long-class preparation hypotheses imply that each actual ordinary spoke weight is at least $1/4$ and that the raw approximate ordinary rim density $p=(\sum_u b_u)/D$ is at least $\rho_0/8$. Thus these two positive lower bounds follow from the existing recipe uniformly over the prescribed factor family. The marked rim density is at least $\rho_0/(6(m+1))$. One positive constant, chosen from $m$ and $\rho_0$, is a lower bound for both rim densities and every ordinary spoke weight; the same quantities are at most one. The remaining seed concentration, event-intersection and matching hypotheses still require their own verification.
\end{lemma}

\begin{lemma}[The actual raw matching and its wheel realizations]\label{bp:actual-raw-wheel-realizations}
\uses{bp:raw-wheel-realizations,bp:raw-absorber-protection}
\lean{Oberwolfach.long_raw_realize_component_resources,Oberwolfach.long_raw_realize_component_hub,Oberwolfach.long_raw_matching_realizations,Oberwolfach.long_raw_matching_realization_preserves_hub,Oberwolfach.long_raw_matching_realized_family_hub_count,Oberwolfach.long_raw_realizations_avoid_fixed_absorber,Oberwolfach.long_raw_realized_resource_union,Oberwolfach.fiberwise_injective_assignment}
\leanok
For the mixed raw marked and ordinary candidate families, each component is realized with exactly its physical rim and spoke resources and its original hub. A positive-weight resource matching with nonempty candidates gives distinct, pairwise physically disjoint realizations. The number in each length-class and hub is unchanged, and the union of their physical resources is exactly the physical image of the matching resources. Under the manuscript's actual colour choices these realizations avoid the fixed absorber in the retained selected host. A separate finite assignment lemma gives an injective key-preserving assignment whenever each target key class has sufficient capacity. Establishing the required class capacities and constructing the partial factor images, as well as the common matching event and first cover-down, remain open.
\end{lemma}

\begin{lemma}[Assigning selected ordinary wheels to components]\label{bp:long-wheel-component-assignment}
\uses{bp:actual-raw-wheel-realizations,bp:raw-common-seed-frequency}
\lean{Oberwolfach.card_of_uniform_fibers,Oberwolfach.component_size_vertex_count,Oberwolfach.long_preparation_fraction_component_count,Oberwolfach.approximate_selected_real_target_le_capacity,Oberwolfach.long_kept_target_le_component_count,Oberwolfach.long_selected_ordinary_count_le_components,Oberwolfach.selected_wheels_assign_to_components,Oberwolfach.long_selected_wheels_assign_components,Oberwolfach.long_raw_ordinary_candidates_assign_components}
\leanok
For a prescribed factor, its vertex fraction at a given length, multiplied by the order and divided by that length, is exactly the number of its components of that length. The retained fraction is no larger. If the selected ordinary wheel counts satisfy the manuscript's displayed relative frequency estimate with error parameter at most $\theta$, the selected wheels at all hubs can be assigned injectively to actual components of the prescribed factor with the correct hub and length. This applies to the original mixed raw candidate indices and retains multiplicities. It does not establish the simultaneous frequency estimates or construct the vertex maps of the partial factor images.
\end{lemma}

\begin{lemma}[Retained frequency estimates and interval locations]\label{bp:long-frequency-assignment-interval}
\uses{bp:long-wheel-component-assignment,bp:actual-raw-wheel-realizations,bp:long-ordinary-density-lower}
\lean{Oberwolfach.long_ordinary_seed_zero_at_discarded_hub,Oberwolfach.long_discarded_hub_selected_zero,Oberwolfach.retained_target_lower_bound,Oberwolfach.absolute_to_relative_target_error,Oberwolfach.long_retained_hub_padded_length_exact,Oberwolfach.long_raw_absolute_to_relative_frequency,Oberwolfach.long_padded_ordinary_assignment_from_absolute_frequency,Oberwolfach.sampled_mixed_hub_frequency_from_patterns,Oberwolfach.sampled_mixed_hub_absolute_frequency,Oberwolfach.long_assigned_seed_coefficient_le_four,Oberwolfach.ordinary_common_hub_active_roots_card_le,Oberwolfach.raw_marked_wheel_unique_interval,Oberwolfach.raw_marked_wheel_interval_unique_position,Oberwolfach.disjoint_marked_wheels_distinct_interval_positions}
\leanok
At a discarded hub--length pair the original seed weight is zero, so a positive-weight matching selects no wheel there. At a retained pair, the padded length is the true component length and the retained fraction gives the displayed lower bound for the number of components. Consequently, if the matching supplies the stated uniform additive frequency estimate at every retained pair and its numerical margin holds, all selected ordinary candidates can be assigned injectively to components at the same hub and of the right length. The intermediate sampled-load estimate retains the fixed weights and the original root-omitted counts; it allows the coefficient bound four supplied by the preparation estimates. Every marked wheel has a unique raw interval start at its hub and, under the stated no-short-orbit condition, a unique position within that interval; physically disjoint marked wheels at the same hub have different interval positions. The simultaneous matching-frequency event and the construction of the full partial factor images remain open.
\end{lemma}

\begin{lemma}[Ordinary hub targets and frequency margins]\label{bp:ordinary-hub-target-frequency}
\uses{bp:long-frequency-assignment-interval,bp:raw-retained-hub-tests}
\lean{Oberwolfach.long_kept_fraction_bounds,Oberwolfach.long_kept_frequency_target_abs_le_order,Oberwolfach.long_ordinary_hub_root_target,Oberwolfach.long_ordinary_hub_root_load_bound,Oberwolfach.long_raw_ordinary_hub_spoke_count,Oberwolfach.long_raw_hub_context_spokes,Oberwolfach.long_raw_context_hub_frequencies,Oberwolfach.long_ordinary_hub_absolute_frequency,Oberwolfach.long_all_hub_absolute_frequencies,Oberwolfach.retained_hub_matching_tests_to_indices,Oberwolfach.long_hub_frequency_parameters,Oberwolfach.long_hub_count_error_budget,Oberwolfach.long_hub_matching_error_margins}
\leanok
For the actual long-class proportions, the retained fraction lies between zero and one. The mixed root target on an ordinary hub spoke is exactly the original kept contribution divided by the fixed hub weight, and a raw ordinary wheel uses precisely its length in such spokes at its own hub. The stated load, hub-count and matching-test estimates imply additive frequency estimates simultaneously at all retained hubs, without an extra target-size assumption. This estimate now applies to the actual raw mixed candidate family and the same raw state, for the supplied fixed weights. The retained tests in the weighted matching have exactly the indicators used in these estimates. The numerical lemmas choose frequency and retained-test margins from fixed positive parameters before the factor family; after the matching theorem gives its error threshold, they choose load tolerances meeting that threshold. The count budget fits the fixed density lower bound. These results apply to a supplied common sample and matching satisfying their hypotheses. Constructing the joint event and the matching, the partial factor maps, and the first cover-down remains open.
\end{lemma}

\begin{lemma}[Ordinary component maps for selected wheels]\label{bp:long-ordinary-component-maps}
\uses{bp:long-wheel-component-assignment,bp:actual-raw-wheel-realizations,bp:long-frequency-assignment-interval}
\lean{Oberwolfach.component_to_ordered_ordinary_wheel_embedding,Oberwolfach.ordinary_component_wheel_family_maps,Oberwolfach.restricted_resource_disjoint_in_original_domain,Oberwolfach.ordinary_restricted_wheels_original_resources_disjoint,Oberwolfach.assigned_restricted_ordinary_wheel_maps,Oberwolfach.option_some_sigma_payload_index,Oberwolfach.option_some_sigma_payload_injective,Oberwolfach.mixed_restricted_some_payload_index,Oberwolfach.mixed_restricted_some_payload_resources,Oberwolfach.mixed_restricted_some_payload_injective,Oberwolfach.long_raw_selected_ordinary_index,Oberwolfach.long_raw_selected_ordinary_resources,Oberwolfach.long_raw_selected_ordinary_hub,Oberwolfach.long_raw_selected_ordinary_restricted_matching,Oberwolfach.long_raw_selected_ordinary_component_maps,Oberwolfach.factor_component_assignment_transport,Oberwolfach.long_raw_assigned_ordinary_component_maps,Oberwolfach.long_padded_selected_ordinary_is_retained,Oberwolfach.long_padded_selected_assignment_reindex,Oberwolfach.long_padded_selected_assignment_from_frequencies}
\leanok
A component of a prescribed two-factor maps injectively to the rim vertices of an ordered wheel of the same length, sending each component edge to a rim edge and each vertex to a wheel spoke. For wheels in a common hub, disjoint sampled resource sets give disjoint vertex images. The ordinary part of the actual mixed matching can be extracted without changing candidate indices, hubs or restricted resource sets. The assignment supplied by the retained frequency estimates can be reindexed to these extracted wheels; positive selected weight forces each padded length to be its true component length. The displayed component-map theorem applies to an injective hub- and length-preserving assignment. Combining these ingredients with the common matching event and marked part of the partial factors remains open.
\end{lemma}

\begin{lemma}[Partial factor maps from selected ordinary wheels]\label{bp:ordinary-partial-factor-maps}
\uses{bp:long-ordinary-component-maps,bp:ordinary-hub-target-frequency,bp:actual-raw-wheel-realizations}
\lean{Oberwolfach.matching_frequency_assignment_form,Oberwolfach.long_padded_component_maps_from_frequencies,Oberwolfach.component_family_embeddings_glue,Oberwolfach.component_family_partial_embedding,Oberwolfach.factor_component_hub_equiv_vertex,Oberwolfach.factor_component_hub_equiv_adj,Oberwolfach.fixed_hub_assigned_components_injective,Oberwolfach.assigned_factor_partial_embeddings,Oberwolfach.assigned_factor_component_source_disjoint,Oberwolfach.component_family_partial_graph_embedding,Oberwolfach.assigned_factor_partial_graph_embeddings,Oberwolfach.long_raw_assigned_partial_graph_embeddings,Oberwolfach.long_raw_component_maps_to_partial_factors,Oberwolfach.long_padded_ordinary_partial_factors_from_frequencies,Oberwolfach.long_raw_selected_ordinary_resource_active,Oberwolfach.long_raw_selected_ordinary_rim_active,Oberwolfach.long_raw_partial_factor_edges_active}
\leanok
For one positive-support matching satisfying the displayed retained frequency estimates, the selected ordinary candidates can be assigned to components of the prescribed factors and mapped to their wheel rims. At each hub, these component maps glue to one injection on the union of the assigned source components. Every source edge in that union maps to a selected rim edge, and the selected rim resources are active in the original raw approximate sample. The assignment, domains and wheel resources are those of the supplied matching; the construction makes no independent resampling. Establishing the common matching event, integrating the marked candidates, and completing preparation and cover-down remain open.
\end{lemma}

\begin{lemma}[Ordinary partial factors avoid the fixed absorber]\label{bp:ordinary-partials-absorber-safe}
\uses{bp:ordinary-partial-factor-maps,bp:raw-absorber-protection,bp:ordinary-hub-target-frequency}
\lean{Oberwolfach.long_raw_selected_rim_physical_resource,Oberwolfach.long_raw_selected_rim_relation_edge_mem,Oberwolfach.long_raw_selected_rim_edge_is_physical_resource,Oberwolfach.long_raw_selected_ordinary_physical_disjoint,Oberwolfach.long_raw_selected_rim_edges_disjoint,Oberwolfach.partial_factors_physical_edges_disjoint,Oberwolfach.long_raw_partial_factor_edges_disjoint,Oberwolfach.long_raw_partial_factors_pairwise_disjoint,Oberwolfach.long_padded_ordinary_disjoint_partials_from_frequencies,Oberwolfach.long_raw_context_hub_frequencies_for_assignment,Oberwolfach.long_raw_selected_ordinary_realized_resources,Oberwolfach.long_raw_selected_ordinary_avoids_fixed_absorber,Oberwolfach.long_raw_partial_factors_disjoint_and_absorber_safe,Oberwolfach.long_padded_ordinary_absorber_safe_partials_from_frequencies}
\leanok
For a supplied matching of the actual raw mixed candidate family and its ordinary partial factor maps, a single choice of the maps sends every included factor edge to a raw approximate rim, makes the images of different factors physically edge-disjoint, and avoids the fixed terminal absorber in the original retained selected host. The avoidance statement uses the same sampling outcome and retention outcome as the absorber. The source-specific frequency and positive-support hypotheses give these same absorber-avoiding maps directly, but the joint sampling event and matching, the marked partial factors, the remaining vertices, and the first cover-down are not constructed here.
\end{lemma}

\begin{lemma}[Counting the ordinary partial-factor domains]\label{bp:ordinary-partial-domain-count}
\uses{bp:ordinary-partials-absorber-safe,bp:ordinary-hub-target-frequency}
\lean{Oberwolfach.component_family_domain_card,Oberwolfach.long_raw_component_sizes_from_rim_cover,Oberwolfach.long_raw_partial_domain_card,Oberwolfach.long_raw_partial_domain_card_by_wheel_length,Oberwolfach.long_raw_partial_domain_complement_card,Oberwolfach.long_raw_partial_domain_count_from_component_maps,Oberwolfach.finite_lengths_by_class,Oberwolfach.long_raw_selected_hub_class_card,Oberwolfach.long_raw_partial_domain_card_by_wheel_class,Oberwolfach.long_raw_partial_domain_card_by_matching_counts,Oberwolfach.long_raw_partial_factors_disjoint_absorber_safe_and_counted,Oberwolfach.long_padded_ordinary_counted_safe_partials_from_frequencies}
\leanok
For the same ordinary maps which use raw rims, have disjoint physical edge images and avoid the fixed absorber, the number of vertices in a prescribed factor's source domain is exactly the sum, over the ordinary length classes, of the selected matching count at that hub times the corresponding cycle length. Its complement has the corresponding size $n$ minus this sum. The counts are those of the original supplied mixed matching; the original sample and retention outcome are unchanged. The result uses a supplied matching, the fixed absorber and its frequency estimates. It does not yet derive the numerical remaining-vertex estimates, construct the marked partial maps or prove the first cover-down.
\end{lemma}

\begin{lemma}[Quantitative ordinary remainder bounds]\label{bp:ordinary-partial-remainder-bounds}
\uses{bp:ordinary-partial-domain-count,bp:ordinary-hub-target-frequency}
\lean{Oberwolfach.class_frequency_domain_error,Oberwolfach.long_raw_absolute_frequency_all_classes,Oberwolfach.long_raw_matching_size_from_frequencies,Oberwolfach.finite_complement_approximation,Oberwolfach.long_padded_ordinary_safe_partial_domain_error_from_frequencies,Oberwolfach.long_padded_length_sum_bound,Oberwolfach.long_padded_length_error_bound,Oberwolfach.long_padded_error_parameter,Oberwolfach.long_padded_error_parameter_with_ceiling,Oberwolfach.long_padded_error_parameter_for_assignment,Oberwolfach.retained_complement_target_bounds,Oberwolfach.long_ordinary_complement_target,Oberwolfach.long_complement_count_bounds,Oberwolfach.long_ordinary_complement_final_bound,Oberwolfach.long_padded_ordinary_complement_bounds_from_frequencies}
\leanok
Write $\nu_i+2$ for the fixed padded cycle length of class $i$. For the same ordinary partial maps which avoid the fixed absorber and have pairwise disjoint physical edge images, the size of the source domain is within $E$ of $(1-\theta)n$ times the sum of the retained ordinary class fractions, and the size of its complement is within $E$ of the complementary target, where $E=(\max\{\varepsilon,0\})n\sum_i(\nu_i+2)$. When at least half of each prescribed factor's vertices lie in components longer than the cutoff, this also bounds the complement between $n$ times the long-component fraction minus $E$ and $n$ times that fraction plus $101\theta n/100+E$. The length-sum bound chooses a positive error tolerance, below the existing assignment ceiling, before $n$ so that $E\leq \delta n$. The matching, its retained estimates and the fixed absorber are supplied inputs; the common event, marked partial maps, compatibility with the complete preparation and the first cover-down remain open.
\end{lemma}

\begin{lemma}[Long components remain outside the ordinary domains]\label{bp:ordinary-partial-preserve-long}
\uses{bp:ordinary-partial-remainder-bounds}
\lean{Oberwolfach.long_raw_partial_domain_excludes_long,Oberwolfach.long_raw_partial_domain_contains_long_count,Oberwolfach.long_raw_partial_domain_excludes_long_of_length_bound,Oberwolfach.long_raw_partial_factors_disjoint_absorber_safe_and_exclude_long,Oberwolfach.long_padded_ordinary_safe_partials_exclude_long_from_frequencies,Oberwolfach.long_padded_safe_partial_domain_error_and_long_preservation,Oberwolfach.long_padded_ordinary_complement_bounds_preserve_long}
\leanok
When the cutoff $K$ is at least three, the same selected ordinary partial maps leave every vertex in a factor component longer than $K$ outside their source domains. This conclusion holds simultaneously with their raw-rim containment, pairwise physical edge disjointness, avoidance of the fixed absorber and the stated quantitative complement bounds. The supplied matching, retained frequency estimates and absorber remain inputs; the later marked packing and the actual first cover-down are not constructed here.
\end{lemma}

\begin{lemma}[Surviving matching and deletion-safe ordinary factors]\label{bp:surviving-matching-partial-factors}
\uses{bp:ordinary-partial-preserve-long,bp:actual-raw-wheel-realizations}
\lean{Oberwolfach.restricted_resource_nested_lift_set,Oberwolfach.restricted_resource_nested_lift_matching,Oberwolfach.restricted_resource_nested_lift_count,Oberwolfach.surviving_candidate_positive_weight_lift,Oberwolfach.long_raw_surviving_resource_not_deleted,Oberwolfach.long_raw_surviving_sample_spoke,Oberwolfach.long_raw_surviving_matching_lift,Oberwolfach.long_raw_lifted_matching_avoids_deleted_rims,Oberwolfach.long_raw_surviving_ordinary_rim_avoids_deleted,Oberwolfach.long_raw_partial_factors_deleted_safe_and_exclude_long,Oberwolfach.long_raw_partial_factors_from_surviving_matching}
\leanok
Delete any specified set of physical rim edges from the fixed raw sample, retaining its sampled spokes. A matching of the surviving candidate resources lifts to the original sample without changing its candidate identities, fixed weights or hub and length counts. Given the original component-map property and the fixed absorber, this lift yields ordinary partial factor maps whose physical edges avoid both the deleted rim edges and the absorber; they are disjoint across hubs, use raw rims and leave components longer than the cutoff outside their domains. This is conditional on the supplied matching and component-map property. It does not construct the simultaneous event, matching or marked partial factors, and does not establish the full first cover-down.
\end{lemma}

\begin{lemma}[Retained-frequency bounds after rim deletion]\label{bp:surviving-matching-complement}
\uses{bp:surviving-matching-partial-factors,bp:ordinary-partial-remainder-bounds}
\lean{Oberwolfach.surviving_candidate_fixed_weighted_sum,Oberwolfach.surviving_candidate_fixed_weighted_test,Oberwolfach.surviving_candidate_lift_range_iff,Oberwolfach.long_raw_surviving_resource_iff,Oberwolfach.long_raw_surviving_candidate_iff,Oberwolfach.long_padded_ordinary_deleted_safe_partials_from_frequencies,Oberwolfach.long_padded_ordinary_partials_from_surviving_frequencies,Oberwolfach.long_padded_deleted_safe_domain_error_and_long_preservation,Oberwolfach.long_padded_deleted_safe_complement_bounds_preserve_long,Oberwolfach.long_padded_complement_bounds_from_surviving_frequencies}
\leanok
The surviving resources and candidates are exactly those whose physical rims avoid the specified deleted edges. Fixed weighted tests agree for a surviving matching and its lift. Given that matching with positive original weights and the stated retained hub-and-length frequencies, and given the fixed absorber, one choice of ordinary partial factor maps avoids the deleted edges and absorber, preserves all components longer than the cutoff, and satisfies the earlier quantitative complement bounds. The matching, its simultaneous frequency estimates and the absorber are inputs; this does not construct the common event, marked maps or first cover-down.
\end{lemma}

\begin{lemma}[Fixed-weight loads after physical rim deletion]\label{bp:long-adaptive-rim-loads}
\uses{bp:surviving-matching-complement,bp:raw-approximate-resource-hub-tests}
\lean{Oberwolfach.long_raw_adaptive_omitted_load_after_rim_deletion,Oberwolfach.long_raw_adaptive_pattern_load_after_rim_deletion,Oberwolfach.long_raw_adaptive_pattern_load_target_bound,Oberwolfach.adaptive_deleted_pair_sum_bound,Oberwolfach.long_raw_adaptive_pattern_load_degree_bound,Oberwolfach.long_raw_surviving_state_is_raw_rim_deletion,Oberwolfach.long_raw_deleted_resource_iff_rim_changed,Oberwolfach.long_raw_deleted_resource_incidence_bound,Oberwolfach.long_raw_adaptive_pattern_load_physical_deletion,Oberwolfach.long_raw_ordinary_hub_loads_after_physical_deletion}
\leanok
For every realised raw sample and every later physical rim deletion, the change in a fixed-weight root-omitted load is bounded by the original omitted pair weights of the realised deleted resources, including at roots outside the surviving state. The filtered sample is exactly the existing raw-outcome rim deletion. If the physical deletion graph has maximum degree at most $\delta$, the deleted typed resources have auxiliary in- and out-incidence at most $2\delta$. Consequently, under the displayed near/far pair-weight hypotheses, the original fixed target changes by at most $2\delta((|S|+|T|)M+nN)$; the same implication holds simultaneously for the actual ordinary hub seed families. No sampling density multiplies this realised deletion loss. Establishing those near/far estimates for all required mixed seed families and connecting them to the common event remain open.
\end{lemma}

\begin{lemma}[Ordinary seed loads under bounded physical deletion]\label{bp:long-ordinary-deleted-load}
\uses{bp:long-adaptive-rim-loads}
\lean{Oberwolfach.ordered_ordinary_disjoint_rims_four_vertices,Oberwolfach.ordered_ordinary_disjoint_rims_all_lengths,Oberwolfach.ordered_ordinary_rim_pair_near_far,Oberwolfach.ordered_ordinary_spoke_rim_pair_near_far,Oberwolfach.resource_omitted_pair_weight_map,Oberwolfach.long_embedded_ordinary_rim_pair_near_far,Oberwolfach.long_embedded_ordinary_spoke_rim_near_far,Oberwolfach.long_embedded_ordinary_pair_marked_rim_zero,Oberwolfach.long_embedded_ordinary_pair_nonordinary_root_zero,Oberwolfach.long_ordinary_deleted_rim_pair_near_far,Oberwolfach.long_ordinary_root_vertices_card_le_two,Oberwolfach.long_ordinary_fixed_load_after_physical_deletion_uniform,Oberwolfach.long_assigned_ordinary_seed_load_after_physical_deletion}
\leanok
For a fixed ordinary wheel family of length $m+1$, the omitted pair weights are $O(1/n)$ when a deleted rim meets the root's vertices and $O(1/n^2)$ otherwise. This includes triangles; a triangle has no pair of vertex-disjoint rims. For original nonnegative seed weights bounded by $B/n^m$ and a physical deletion graph of maximum degree at most $\delta$, every root-omitted load changes by at most $10\delta(m+1)^4B/n$. The statement applies to the original assigned ordinary seed weights once their previously proved coefficient bounds are supplied, for every later deletion and without resampling. The corresponding marked seed bounds, use of the actual common event and the construction of a matching remain open.
\end{lemma}

\begin{lemma}[Mixed seed loads after every later deletion]\label{bp:long-mixed-uniform-deleted-load}
\uses{bp:long-ordinary-deleted-load,bp:long-adaptive-rim-loads,bp:raw-common-seed-frequency}
\lean{Oberwolfach.long_marked_deleted_rim_pair_near_far,Oberwolfach.long_marked_root_vertices_card_le_two,Oberwolfach.long_marked_fixed_load_after_physical_deletion_uniform,Oberwolfach.long_assigned_marked_seed_load_after_physical_deletion,Oberwolfach.long_assigned_mixed_seed_load_after_physical_deletion,Oberwolfach.long_assigned_named_seed_physical_deletion_stability,Oberwolfach.long_mixed_deletion_uniform_constant,Oberwolfach.long_assigned_named_seed_uniform_physical_deletion_stability,Oberwolfach.long_assigned_common_event_loads_after_every_deletion}
\leanok
The fixed assigned marked and ordinary seed weights both satisfy the corresponding near/far pair bounds. The existing coefficient bounds give one explicit constant, fixed before the raw outcome and deletion graph, that controls every mixed seed load after any later physical-rim deletion by the original load error plus a constant multiple of the deleted maximum degree divided by $n$. A single hypothesis for all original roots and seed families is used before quantifying over the later deletion graph. The sampled outcome, fixed weights and targets are unchanged, and no density multiplier is inserted into the realised deletion loss. Establishing the required original common event, total test masses, matching, marked factor maps and full preparation remains open.
\end{lemma}

\begin{lemma}[Retained hub tests after every later rim deletion]\label{bp:long-deleted-hub-tests}
\uses{bp:long-mixed-uniform-deleted-load,bp:surviving-matching-complement,bp:ordinary-hub-target-frequency}
\lean{Oberwolfach.long_raw_surviving_bool_state,Oberwolfach.long_raw_ordinary_hub_spoke_count_after_rim_deletion,Oberwolfach.equal_state_full_weighted_test_mass,Oberwolfach.equal_state_candidate_matching,Oberwolfach.equal_state_candidate_filter_card,Oberwolfach.long_assigned_common_event_raw_deleted_loads,Oberwolfach.long_assigned_hub_frequencies_after_every_deletion,Oberwolfach.long_assigned_surviving_hub_frequencies_after_every_deletion,Oberwolfach.long_assigned_retained_test_mass_after_every_deletion,Oberwolfach.long_assigned_surviving_retained_test_mass_after_every_deletion,Oberwolfach.long_raw_pattern_load_nonincreasing_after_rim_deletion}
\leanok
Physical rim deletion leaves the sampled ordinary hub-spoke counts unchanged. On one raw outcome satisfying the stated simultaneous load and spoke-count hypotheses, the fixed retained-hub test masses retain their lower bounds after every later deletion satisfying the displayed numerical margin. Any supplied matching on the surviving resources whose weighted tests satisfy the stated inequalities also has the corresponding ordinary hub-and-length frequency bounds. The filtered-state and raw-deletion formulations are connected by an identity on candidate identifiers, preserving both full test masses and selected test sums. These results do not construct the original common event, establish the remaining matching tests or produce a matching.
\end{lemma}

\begin{lemma}[Simultaneous incidence tests for a fixed mixed sample]\label{bp:mixed-all-incidence-tests}
\uses{bp:fractional-resource-matching}
\lean{Oberwolfach.mixed_sampled_resource_matching_all_incidence_sets,Oberwolfach.long_raw_seed_resource_sizes,Oberwolfach.mixed_sampled_resource_matching_all_and_large_incidence_sets,Oberwolfach.mixed_matching_rim_and_other_incidence_tests,Oberwolfach.equal_state_candidate_covered_union,Oberwolfach.equal_state_candidate_uncovered_incidence_card,Oberwolfach.mixed_matching_rim_and_other_incidence_tests_on_named_state}
\leanok
For the original nonnegative candidate weights and a fixed sampled state, a single positive-support resource matching simultaneously satisfies all supplied weighted tests and covers every member of any polynomial-size family of resource sets, leaving at most $\max\{an,\gamma |S|\}$ resources of a set $S$ uncovered, provided the displayed load, pair-weight and test-mass hypotheses of the Lean statement hold. For the actual long raw seed family, the resources of every candidate have the stated bounded size. The source-specific conditional incidence application is given below. Its original common event and the discharge of all load, pair-weight and test-mass hypotheses remain open.
\end{lemma}


\begin{lemma}[Physical residual of the same surviving matching]\label{bp:surviving-physical-residual}
\uses{bp:surviving-matching-partial-factors,bp:raw-approximate-unused-rim-degree}
\lean{Oberwolfach.physical_rim_residual_after_deletion,Oberwolfach.physical_rim_residual_degree_after_deletion,Oberwolfach.surviving_candidate_covered_resources_lift,Oberwolfach.long_raw_physical_residual_after_rim_deletion,Oberwolfach.long_raw_matching_residual_degree_after_rim_deletion,Oberwolfach.long_approximate_unused_degree_after_rim_deletion,Oberwolfach.long_required_vortex_rows_after_rim_deletion,Oberwolfach.long_raw_surviving_matching_physical_rims_avoid_deleted_graph}
\leanok
The covered-resource union of a supplied surviving matching is preserved exactly by its lift to the original raw sample. Its physical rim graph avoids the deleted graph. The original physical residual degree is at most the surviving incidence-coverage error plus the realised deletion degree; the approximate-graph bound also retains the ineligible-edge term. Under the displayed row, containment and numerical margins, this gives the existing first-vortex required-graph bounds. No sampling density multiplies the deletion loss. The matching, common event, deletion-degree estimate and remaining preparation fields are not constructed by these transfer statements.
\end{lemma}

\begin{lemma}[Actual physical, auxiliary and spoke tests on one matching]\label{bp:long-actual-incidence-matching}
\uses{bp:mixed-all-incidence-tests,bp:surviving-physical-residual,bp:long-assigned-joint-seed-event}
\lean{Oberwolfach.long_raw_incidence_test_index_card_le_cube,Oberwolfach.long_raw_surviving_resource_card_le_cube,Oberwolfach.long_raw_incidence_set_card_le_full,Oberwolfach.long_raw_auxiliary_incidence_card_le,Oberwolfach.long_raw_vertex_spoke_incidence_card_le,Oberwolfach.long_raw_hub_spoke_incidence_card_le,Oberwolfach.long_raw_surviving_matching_all_incidences,Oberwolfach.long_assigned_named_seed_pair_bounds,Oberwolfach.long_assigned_reference_matching_inputs_witness,Oberwolfach.long_matching_tolerances_exist,Oberwolfach.long_matching_uniform_tolerances_exist,Oberwolfach.long_matching_manuscript_uniform_tolerances_exist,Oberwolfach.hub_test_mass_manuscript_margin,Oberwolfach.long_assigned_retained_hub_test_manuscript_mass,Oberwolfach.long_assignment_uniform_tolerances_exist,Oberwolfach.long_assignment_manuscript_tolerances_exist,Oberwolfach.long_manuscript_matching_parameters_exist,Oberwolfach.long_manuscript_matching_incidence_margins,Oberwolfach.long_assigned_hub_test_manuscript_mass_from_sampled_patterns,Oberwolfach.long_retained_hub_tests_manuscript_mass_from_patterns,Oberwolfach.long_retained_hub_tests_manuscript_mass_from_mixed_targets,Oberwolfach.long_assigned_manuscript_test_mass_after_every_deletion,Oberwolfach.long_assigned_surviving_manuscript_test_mass_after_every_deletion,Oberwolfach.long_assigned_actual_matching_coefficients,Oberwolfach.long_assigned_host_matching_weights,Oberwolfach.long_raw_matching_retained_hub_and_incidence_tests,Oberwolfach.long_assigned_reference_matching_after_deletion,Oberwolfach.long_assigned_reference_matching_manuscript_cutoff,Oberwolfach.long_assigned_reference_manuscript_matching_uniform_deletions,Oberwolfach.long_assigned_reference_manuscript_matching_assignment_ready,Oberwolfach.long_assigned_reference_matching_uniform_deletions,Oberwolfach.long_assigned_reference_matching_with_uniform_failure,Oberwolfach.long_assigned_reference_matching_assignment_ready,Oberwolfach.long_raw_matching_incidence_linear_bounds}
\leanok
The auxiliary endpoint and the two typed spoke incidence families are defined on the actual surviving raw resource set. Together with the physical rim tests their index set has size at most $n^3$ for $n\ge4$ and $|U|\le n$. Auxiliary tests have size at most $4n$; each spoke test has size at most $n$. The source-specific matching caller derives these polynomial-size conditions and the actual candidate-size condition, and returns one positive-support matching satisfying all incidence and supplied fixed-weight tests, together with the original physical residual bound for its exact lift. The original weights are unchanged. On the same conditional reference-law witness, the actual assigned inputs satisfy the mixed target sum, nonnegative weight, pair bounds with an explicit positive total constant, polynomial retained-test count, test range, deletion-filtered loads and retained-test mass hypotheses. The displayed positivity, padded-length, host, failure and deletion-margin inputs remain explicit. Positive load and deletion tolerances satisfying the matching and mass inequalities can be chosen after the matching tolerance and before the order and sample. The original host and reference-event hypotheses, the actual deletion bound and the final matching invocation still have to be supplied, so this is not yet the unconditional preparation construction.

With fixed positive tolerances chosen before the order, the prescribed sampling coefficients and uniform seed and spoke failure estimates supply the matching margins and the ordinary assignment error and hub-count margins for every deletion within a fixed linear degree budget. The resulting matching is on the same retained outcome and gives linear auxiliary and typed-spoke bounds. Under the displayed sharper relative errors, the retained weighted-test mass has the manuscript constant $	heta n/(200K^2)$, and the original minimum formulas for the matching parameters satisfy the ordinary-assignment hierarchy after choosing the spacing sufficiently large; the same formulas imply the stated absolute incidence margins when the deletion degree is small. The matching theorem has also been invoked with that sharper bound; its uniform-tolerance and ordinary partial-map callers are also checked. The matching-compatible fractional-cycle and successor-sequence weights also follow from the separately stated host-preparation hypotheses, with padded lengths and retained indices identified. The probability bound for the prior reference event and the displayed host and absorber assumptions remain explicit.
\end{lemma}

\begin{lemma}[Assigned seed concentration on the actual expanded law]\label{bp:long-assigned-joint-seed-event}
\uses{bp:raw-common-seed-frequency,bp:retained-actual-seed-weights}
\lean{Oberwolfach.long_assigned_named_seed_common_event,Oberwolfach.expanded_raw_spoke_hub_counts_any_retention,Oberwolfach.long_raw_ordinary_hub_roots_eq_spoke_column,Oberwolfach.long_assigned_spoke_column_to_hub_root_counts,Oberwolfach.long_assigned_seed_and_spoke_common_event,Oberwolfach.long_assigned_seed_and_hub_count_common_event,Oberwolfach.eventually_assigned_spoke_failure_below,Oberwolfach.eventually_assigned_seed_spoke_failure_below,Oberwolfach.long_assigned_reference_seed_hub_witness,Oberwolfach.long_assigned_reference_seed_hub_load_witness,Oberwolfach.long_assigned_reference_seed_hub_load_mass_witness}
\leanok
The actual long-class recipe supplies all parameters and validity proofs of the original expanded law, and the fixed assigned marked and ordinary seed functions satisfy the required coefficient bounds. On the same expanded-and-retained experiment, the joint concentration event intersects any supplied prior event with the original explicit total failure bound. This statement does not combine separately chosen witnesses. The assigned seed concentrations and the matching hub-root counts also hold in one event, with the original assigned law and explicit combined failure bound. The combined seed and spoke failure is uniformly small beyond a threshold chosen before the order and hub set, for fixed positive parameters and linear spoke-count tolerance. A probability estimate for any specified event under the original retained reference law transfers to the expanded law; if its explicit failure budget together with the seed and spoke failures is below one, one positive-mass outcome satisfies the reference event, seed concentrations and hub-count bounds simultaneously. Under the stated positive-density, fractional-cycle, host-degree and numerical hypotheses, this same positive-mass outcome also satisfies the original assigned mixed active-load bounds. Under the displayed padded-length, hub-count and deletion margins, every retained hub test has weighted mass at least $\theta n/(1600K^2)$ after every allowed physical deletion. The reference-event estimate and its host and vortex hypotheses, these original host and numerical inputs, and the remaining matching parameter margins still have to be established.
\end{lemma}

\begin{lemma}[Ordinary partial factors from the surviving weighted tests]\label{bp:matching-tests-to-safe-ordinary-partials}
\uses{bp:long-deleted-hub-tests,bp:surviving-matching-complement}
\lean{Oberwolfach.long_padded_safe_partials_from_surviving_matching_tests,Oberwolfach.long_padded_safe_partials_from_retained_hub_matching_tests,Oberwolfach.long_assigned_reference_matching_and_ordinary_partials,Oberwolfach.long_assigned_reference_manuscript_matching_and_ordinary_partials,Oberwolfach.long_assigned_vortex_reference_matching_and_partials,Oberwolfach.long_assigned_vortex_reference_matching_and_selected_absorber,Oberwolfach.long_assigned_vortex_reference_manuscript_matching_and_partials,Oberwolfach.long_assigned_vortex_reference_manuscript_matching_and_selected_absorber,Oberwolfach.long_assigned_vortex_manuscript_parameters_and_selected_absorber}
\leanok
For a supplied regular host and vortex satisfying the stated numerical and weight hypotheses, the assigned-law reference and robust-absorber event and one matching now hold on a common retained outcome; that matching supplies these ordinary partial factors for every stated fixed absorber. Under the stated terminal-vortex, absorber-size and forbidden-set bounds, a fixed absorber is selected on that same outcome before the later deletions. This same-outcome construction also works with the manuscript matching cutoff and its displayed minimum formulas, with the spacing and all matching tolerances chosen before the graph and allowed set. The full multi-field preparation event, construction of the vortex and the remaining numerical hypotheses remain explicit. For one supplied surviving matching, the fixed weighted test inequalities and the original simultaneous load and spoke-count hypotheses give the retained frequencies and one family of ordinary partial factor maps. The same maps satisfy the component-count and complement estimates, preserve every source component longer than the cutoff, and avoid both the fixed absorber and deleted rim graph. The event, matching, fixed absorber and numerical margins remain prerequisites; marked partial maps and complete preparation remain open.
\end{lemma}

\begin{lemma}[Compatible marked preparation and first cover-down]\label{bp:marked-reserved-first-cover-caller}
\uses{bp:surviving-physical-residual,bp:long-actual-incidence-matching,bp:long-assigned-joint-seed-event,bp:matching-tests-to-safe-ordinary-partials,bp:marked-first-matching-inputs,bp:sampled-mixed-physical-matching,bp:retained-mixed-seed-targets,bp:raw-approximate-resource-hub-tests,bp:raw-approximate-unused-rim-degree,bp:raw-common-seed-frequency,bp:retained-actual-seed-weights,bp:raw-retained-hub-tests,bp:raw-absorber-protection,bp:raw-wheel-realizations,bp:long-ordinary-density-lower,bp:actual-raw-wheel-realizations,bp:long-wheel-component-assignment,bp:long-frequency-assignment-interval,bp:ordinary-hub-target-frequency,bp:long-ordinary-component-maps,bp:ordinary-partial-factor-maps,bp:ordinary-partials-absorber-safe,bp:ordinary-partial-domain-count,bp:ordinary-partial-remainder-bounds,bp:ordinary-partial-preserve-long,bp:surviving-matching-partial-factors,bp:surviving-matching-complement,bp:long-adaptive-rim-loads,bp:long-ordinary-deleted-load,bp:long-mixed-uniform-deleted-load,bp:long-deleted-hub-tests,bp:mixed-all-incidence-tests}
The actual partial factor images, required physical graph, remaining numerical margins and common sample must be assembled into compatible preparation and first cover-down records for the preselected absorber. This construction remains open.
\end{lemma}

\begin{lemma}[Uniform failure budget for the adaptive agreement event]\label{bp:adaptive-agreement-uniform-budget}
\lean{Oberwolfach.eventually_marked_adaptive_agreement_failure_below,Oberwolfach.marked_adaptive_agreement_failure_budgets_exist,Oberwolfach.marked_adaptive_agreement_named_probability_assigned}
\leanok
For the displayed positive tolerances, the agreement-inclusive adaptive event has a uniform fixed failure budget in the assigned physical sampling law. The level-size threshold is selected before the absorber size and the host-order threshold afterwards.
\end{lemma}
\begin{proof}\leanok The checked probability estimate is combined with the uniform bound on its polynomial error terms. \end{proof}

\begin{lemma}[Adaptive event, matching and reservoir on one outcome]\label{bp:adaptive-agreement-matching-reservoir}
\lean{Oberwolfach.long_assigned_full_adaptive_manuscript_parameters_and_reservoir,Oberwolfach.long_assigned_full_adaptive_canonical_centres_and_reservoir}
\leanok
Under the displayed regularity, level, density and weight hypotheses, the manuscript's matching parameters can be selected before the graph and allowed set. One retained outcome supplies the agreement-inclusive adaptive event, a matching with the stated deletion-uniform test bounds, and ordinary partial maps. On this outcome the reservoir-selection property supplies a fixed reservoir for each terminal pair satisfying the displayed size conditions, uniformly over its eligible remainders.
\end{lemma}
\begin{proof}\uses{bp:adaptive-agreement-uniform-budget}\leanok The assigned-law event is used as the prior event for the matching construction, and the reservoir selector is applied to the same retained outcome. The host index and reference-mass bounds follow from the displayed density hypotheses; the exact level-density identities follow by choosing their canonical centres. \end{proof}

\begin{lemma}[Projecting a uniform prior witness]\label{bp:finite-law-universal-prior}
\lean{Oberwolfach.finiteLaw_probability_of_universal_prior_witness,Oberwolfach.finiteLaw_join_projected_witness_with_endpoint,Oberwolfach.long_assigned_manuscript_matching_projected_probability,Oberwolfach.long_assigned_manuscript_matching_with_endpoint}
\leanok
A witness construction that succeeds for every prior event of probability at least $1-\varepsilon$ has a projected good event of probability greater than $\varepsilon$. In particular, the checked long-assignment matching and ordinary partial maps define such a projected event. The concrete matching construction can be combined with an endpoint construction that accepts any sufficiently probable extra event.
\end{lemma}
\begin{proof}\leanok Apply the universal prior hypothesis to the complement of the projected good event, and then use the endpoint hypothesis. \end{proof}

\begin{lemma}[Constructed vortex endpoint with a projected witness]\label{bp:chosen-vortex-projected-witness}
\lean{Oberwolfach.marked_chosen_vortex_endpoint_with_projected_witness,Oberwolfach.marked_chosen_vortex_endpoint_callback}
\leanok
The constructed vortex and its endpoint-record fields can be combined with any expanded witness whose projection satisfies the displayed uniform prior-event hypothesis. The resulting endpoint fields and expanded witness refer to the same retained outcome.
\end{lemma}
\begin{proof}\uses{bp:finite-law-universal-prior}\leanok Apply the projected-witness intersection lemma to the chosen vortex's extra-event interface. \end{proof}

\begin{lemma}[Canonical densities for a fixed vortex]\label{bp:fixed-vortex-canonical-densities}
\lean{Oberwolfach.marked_fixed_vortex_canonical_densities}
\leanok
For a supplied vortex with positive initial index density, the canonical balance and load centres satisfy the four exact density identities used in the adaptive event.
\end{lemma}
\begin{proof}\leanok Substitute the canonical centres and use the checked level-balance identities. \end{proof}

\begin{lemma}[Adaptive level inputs from a fixed vortex]\label{bp:fixed-vortex-adaptive-level-inputs}
\lean{Oberwolfach.marked_fixed_vortex_adaptive_level_inputs,Oberwolfach.marked_fixed_vortex_weighted_adaptive_level_inputs}
\leanok
A neighbourhood constant can be chosen from the positive vortex shrinkage and index-density bounds before the level-size threshold. For every sufficiently large supplied vortex, its actual vertex and index domains simultaneously satisfy the displayed level-size, nesting, permutation, neighbourhood-size, degree and canonical density conditions. Under the displayed successor-cycle margin, the appropriate successor-sequence weight families are obtained on the same levels.
\end{lemma}
\begin{proof}\uses{bp:fixed-vortex-canonical-densities}\leanok Combine the fixed-vortex geometry with the neighbourhood-size and degree estimates and the canonical density identities. \end{proof}

\begin{lemma}[Ordered adaptive matching and fixed terminal selection]\label{bp:ordered-adaptive-fixed-terminal}
\lean{Oberwolfach.long_assigned_adaptive_manuscript_fixed_selection_ordered,Oberwolfach.long_assigned_adaptive_manuscript_fixed_canonical_ordered}
\leanok
The original matching spacing can be chosen before its dependent weight bounds, and the absorber tolerance before the subsequent event tolerances. Under the displayed host and level hypotheses, one retained outcome supplies the adaptive agreement event, the matching and its ordinary partial maps, and a fixed-terminal-absorber selection for each compatible vortex satisfying the stated size and tolerance bounds. The canonical-centre version also derives the reference mass, host index bound and exact level-density identities from the displayed density hypotheses.
\end{lemma}
\begin{proof}\uses{bp:adaptive-agreement-matching-reservoir}\leanok Combine the reordered probability and matching estimates with the fixed-terminal selector and the canonical centre identities. \end{proof}

\begin{lemma}[A compatible vortex at the regular-host threshold]\label{bp:regular-host-compatible-vortex}
\lean{Oberwolfach.long_regular_host_absorber_compatible_vortex,Oberwolfach.marked_adaptive_auxiliary_constants,Oberwolfach.marked_adaptive_failure_tolerances_exist}
\leanok
After the displayed positive absorber and level tolerances are fixed, compatible positive vortex tolerances and a vortex-size threshold can be chosen before the allowed set. For each regular host satisfying the original allowed-cycle degree threshold and the stated density bounds, they give a vortex and a preserving permutation without short cycles. The displayed auxiliary reference and failure scales can also be selected positively.
\end{lemma}
\begin{proof}\leanok Choose the positive scales, derive the regular-host numerical bounds and apply the vortex and permutation existence theorem. \end{proof}

\begin{lemma}[Regular-host matching and reference inputs]\label{bp:regular-host-matching-reference-inputs}
\lean{Oberwolfach.long_host_matching_vortex_complete_inputs,Oberwolfach.long_host_adaptive_reference_parameters}
\leanok
Under their displayed positivity and regular-host hypotheses, the host and reference parameter lemmas supply a compatible vortex and preserving short-cycle-free permutation together with the host weights and seed budgets for that same permutation, the exact initial index-density identity, the host-degree bounds and zero regularity error. The reference parameters and their density margins are chosen with the stated quantifier order. These lemmas do not themselves invoke the supplied-vortex common event.
\end{lemma}
\begin{proof}\uses{bp:regular-host-compatible-vortex}\leanok Combine the previously established host-weighted vortex construction with the exact regular-host numerical identities, and select the positive reference scales. \end{proof}

\begin{lemma}[Conditional regular-host common-event specialization]\label{bp:regular-host-common-event-specialization}
\uses{bp:regular-host-matching-reference-inputs}
\lean{Oberwolfach.long_host_supplied_event_specialization}
\leanok
Given the displayed regular-host and supplied-vortex event continuations, together with their stated parameter margins, the vortex size is chosen before the allowed set and host order. For each host satisfying the displayed regularity, cycle-length, density and degree hypotheses, the construction supplies a compatible vortex and preserving short-cycle-free permutation, and weights for that same permutation. The matching, degree and budget premises of the supplied-vortex theorem are discharged, yielding its full common event on that vortex. This statement is conditional on the two continuation inputs; their joint selection and the iteration remain separate obligations.
\end{lemma}
\begin{proof}\leanok Apply the host continuation, verify the supplied-vortex premises with the resulting matching weights and exact host identities, and invoke the supplied-vortex event. \end{proof}

\begin{lemma}[Buffered factor paths and partial embeddings]\label{bp:buffered-factor-path-partial-embedding}
\lean{Oberwolfach.regular_two_buffered_path_edge_cover,Oberwolfach.regular_two_buffered_path_target_embedding,Oberwolfach.regular_two_buffered_path_family_target_embedding}
\leanok
In a 2-regular source graph, disjoint extended walks with the displayed two-vertex buffers have no edges between distinct interiors. If each interior has an injective target path of the same indexed length and the target path ranges are disjoint, the index-matching maps glue into a graph-preserving partial embedding. The construction preserves each path's vertex order. Its stated source-walk and target-path hypotheses still need to be supplied in the marked-factor application.
\end{lemma}
\begin{proof}\leanok Use the two neighbours of each interior vertex to establish edge coverage, match equal-length source and target sequences, and glue the local maps. \end{proof}

\begin{lemma}[Indexed simple target walks]\label{bp:indexed-simple-target-walks}
\lean{Oberwolfach.simple_walk_vertex_embedding_adjacent,Oberwolfach.simple_walk_vertex_embedding_support,Oberwolfach.simple_walk_length_embedding_adjacent,Oberwolfach.simple_walk_length_embedding_support}
\leanok
A simple walk gives an injective sequence of its vertices, with exactly the walk's support and consecutive adjacency. The sequence may be indexed by any natural number equal to the walk's length. This supplies the target-path interface needed by the buffered partial-embedding lemma.
\end{lemma}
\begin{proof}\leanok Index the walk support using its path property, and reindex along the stated equality of lengths. \end{proof}

\begin{lemma}[Reserved factor paths and target-walk partial maps]\label{bp:reserved-factor-target-walk-partials}
\uses{bp:buffered-factor-path-partial-embedding,bp:indexed-simple-target-walks}
\lean{Oberwolfach.long_reserved_paths_with_partial_embedding_glue,Oberwolfach.long_reserved_paths_to_disjoint_target_paths,Oberwolfach.regular_two_buffered_paths_to_disjoint_simple_walks,Oberwolfach.regular_two_buffered_embedding_edge_coverage}
\leanok
Under the original reserved-path component, length and floor-capacity hypotheses, the chosen buffered factor paths admit a graph-preserving partial embedding into any displayed family of matching-length, vertex-disjoint target paths. The target-path interface can also be supplied by simple walks of the stated lengths and disjoint supports. Every edge of an index-matched partial map lies along a consecutive pair on one of its target paths, in one of the two orientations. The construction of the marked-wheel target family and the transfer of cross-factor edge disjointness remain separate inputs.
\end{lemma}
\begin{proof}\leanok Combine the reserved path construction, 2-regular buffer separation, indexed path embeddings and edge provenance. \end{proof}

\begin{lemma}[Reserved factor paths mapped to simple walks]\label{bp:reserved-factor-simple-walk-map}
\uses{bp:reserved-factor-target-walk-partials}
\lean{Oberwolfach.long_reserved_paths_to_disjoint_simple_walks}
\leanok
Under the displayed original component, path-length and floor-capacity hypotheses, the buffered source paths can be selected before the target family. For every family of simple target walks of matching lengths with pairwise disjoint supports, there is an injective graph-preserving partial map taking each reserved source path to its target walk index by index. The following lemma applies this construction to the marked-wheel walk family and establishes cross-factor edge disjointness under the displayed decomposition hypotheses.
\end{lemma}
\begin{proof}\leanok Apply the reserved-path partial embedding theorem to the indexed vertex sequences of the simple walks. \end{proof}

\begin{lemma}[Marked decompositions give edge-disjoint partial factor maps]\label{bp:marked-edge-disjoint-partial-factor-maps}
\lean{Oberwolfach.marked_decomposition_to_edge_disjoint_reserved_partial_maps_with_counts}
\leanok
Under the displayed marked physical decomposition, per-interval injectivity and separation hypotheses, two-regular source-factor hypotheses, and reserved-component floor-capacity bounds, the reconstructed target walks admit compatible index-matching partial factor maps with pairwise disjoint physical edge images. The same witnesses retain the decomposition's coverage and hub-count assertions and give exact cardinalities for the partial source domains and their complements. This result assumes the decomposition and capacity; it does not construct the cover-down chain or prove the main theorem.
\end{lemma}
\begin{proof}\leanok Combine marked directed-walk reconstruction with the disjoint buffered source-path embeddings, physical edge coverage, and the exact finite-domain counts. \end{proof}

\begin{lemma}[A sufficient bound for the floor capacity]\label{bp:reserved-floor-slack}
\lean{Oberwolfach.sum_floor_capacity_of_slack}
\leanok
If the total component size is at least the block length times the sum of the demand and the number of components, the sum of the componentwise floor capacities is at least the demand. This lemma does not assert the premise from the manuscript's parameter hierarchy.
\end{lemma}
\begin{proof}\leanok Bound the loss in each component by one block and sum. \end{proof}

\begin{lemma}[Reserved capacity from the long-cycle cutoff]\label{bp:long-reserved-capacity-rpow}
\lean{Oberwolfach.long_reserved_capacity_from_rpow,Oberwolfach.marked_decomposition_to_edge_disjoint_maps_with_rpow_capacity}
\leanok
Suppose the long components and raw demand satisfy the displayed total-size and demand bounds, the cutoff inequality, and the fourth-root size threshold. Then the original sum-of-floors capacity required for reserving buffered paths holds. Consequently, under these same numerical bounds and the displayed marked-decomposition and geometric hypotheses, the edge-disjoint partial-factor-map construction with its exact domain counts applies. The decomposition and the input numerical bounds are hypotheses, not consequences of this lemma.
\end{lemma}
\begin{proof}\leanok Bound the number of long components by their total size divided by the cutoff, prove the real surplus using the cutoff inequality and the fourth-root threshold, and transfer this surplus to the integer componentwise floor sum. Apply the marked partial-factor-map theorem. \end{proof}
\begin{lemma}[Complete compatible buffered factor maps]\label{bp:buffered-factor-completion}
\lean{Oberwolfach.buffered_factor_family_gives_decomposition,Oberwolfach.selected_buffered_path_family_disjoint,Oberwolfach.selected_uniform_buffered_strict_interior_card}
\leanok
For an even-regular host and the displayed two-regular factors, compatible reserved-path and complementary embeddings with the stated boundary agreement, image separation, and cross-factor edge-disjointness yield the exact factor decomposition. Restricting the already embedded paths to the selected final intervals preserves their edge disjointness, and the strict interiors of equal-length selected paths have the displayed exact cardinality. The complementary embeddings and their compatibility are hypotheses here.
\end{lemma}
\begin{proof}\leanok Join the two embeddings using the path-neighbour structure of two-regular factors, transfer edge disjointness, and apply the regular-degree exact decomposition theorem. \end{proof}

\begin{lemma}[Actual residual output to exact decomposition]\label{bp:actual-residual-directed-decomposition}
\lean{Oberwolfach.actual_directed_residual_family_gives_decomposition}
\leanok
For an even-regular host and two-regular factors, the displayed actual residual data and compatible old and new maps give an exact factor decomposition when the completed edge images are pairwise disjoint, the old packing lies in the host, and the directed residual images are host edges. The residual data, maps, and disjointness remain hypotheses.
\end{lemma}
\begin{proof}\leanok Partition each factor's edge image into its old and induced residual parts, verify host containment using the directed source orientation, and apply the regular-degree decomposition theorem. \end{proof}

\begin{lemma}[Residual data and matching conditions]\label{bp:long-short-residual-matching}
\lean{Oberwolfach.long_paths_disjoint_short_component_sets,Oberwolfach.actual_residual_physical_outputs_give_decomposition,Oberwolfach.marked_walk_complement_image_separation,Oberwolfach.residual_spoke_card_from_complement_embedding}
\leanok
The displayed long-component and short-component bounds imply the path--cycle separation needed to construct actual residual data. Given the displayed residual embedding outputs, old-edge packing, directed host containment, and old/new edge avoidance, the existing construction yields an exact factor decomposition. Separately, reconstructed marked-walk coverage, endpoint agreement, and disjoint residual-spoke and complement sets imply the required within-factor image separation; an injective complement embedding yields the stated residual-spoke cardinality. These are conditional interfaces: the common-event data and the residual embedding hypotheses still have to be assembled.
\end{lemma}
\begin{proof}\leanok Use component support separation, the source-edge partition, the exact regular-degree decomposition theorem, and finite injective-image counts. \end{proof}

\begin{lemma}[Exact residual wheel identities and reconstruction]\label{bp:exact-residual-wheel-identities}
\lean{Oberwolfach.buffered_residual_spoke_card,Oberwolfach.marked_interval_roots_in_residual,Oberwolfach.marked_residual_ordinary_avoids_successors,Oberwolfach.residual_directed_degree_balance,Oberwolfach.residual_missing_load_balance,Oberwolfach.marked_residual_identity_gives_factor_decomposition}
\leanok
For disjoint buffered source paths and an injective complement map, the displayed residual-set identity determines the exact spoke count before completing the paths. Selected interval roots avoid their final endpoints and lie in the residual set if they avoid the complement images; ordinary residual labels avoid successors of the marked roots. For disjoint factor arc images with one incidence at each used label, the residual in- and out-degrees have the required missing-factor count when the stated start and end loads balance. Given reconstructed marked walks, their source maps, the residual identity, the displayed host containment and the cross-factor edge-avoidance hypotheses, the completed factors form an exact decomposition. These statements do not construct the residual host or discharge the remaining probabilistic and embedding hypotheses.
\end{lemma}
\begin{proof}\leanok Count injective images and disjoint arc slices, use the interval successor geometry, and glue the compatible partial embeddings. \end{proof}

\begin{lemma}[Uniform root loads balance the interval endpoints]\label{bp:uniform-root-load-endpoints}
\lean{Oberwolfach.marked_uniform_roots_give_endpoint_balance,Oberwolfach.marked_residual_spoke_balance,Oberwolfach.oriented_partial_embedding_incidence,Oberwolfach.buffered_outside_image_avoids_both_endpoints}
\leanok
For each factor, let the displayed positive-length permutation intervals have injective extended windows which are pairwise disjoint. If the selected roots have uniform multiplicity across labels, the start and finish multiplicities agree at every label. For buffered paths of uniform length $(m+1)d$, the displayed residual-set identity and complement embedding yield $|S_w|=(m+1)|Y_w|$. Injective oriented tail and head maps give the single-incidence counts, and initial and final endpoints agreed with the same injective complement map avoid its outside-path image. The construction of the required complement maps and their directed tail and head image sets remains a hypothesis of subsequent applications.
\end{lemma}
\begin{proof}\leanok Shift each marked interval by the same permutation, count the disjoint boundary images, and use injectivity of the complement and oriented maps. \end{proof}

\begin{lemma}[Residual degrees from buffered interval data]\label{bp:buffered-interval-residual-degrees}
\lean{Oberwolfach.residual_degrees_from_buffered_intervals,Oberwolfach.buffered_oriented_path_boundary_identity,Oberwolfach.fiberwise_permutation_preserves}
\leanok
Given the displayed uniform marked intervals, disjoint buffered paths, injective complement maps and agreed endpoint images, and the displayed oriented tail/head image descriptions, host containment, arc disjointness and initial degree equalities, both residual directed degrees equal the count of factors whose residual spoke set contains the label. A buffered source path that follows a permutation has the corresponding orbit-boundary identity. Fiberwise permutations glue to a permutation preserving the fibers. The construction of the actual oriented tail/head images and the residual host remains to be supplied.
\end{lemma}
\begin{proof}\leanok Combine uniform endpoint-load balance, injective oriented incidence, complement-image avoidance and the finite residual arc count. \end{proof}

\begin{lemma}[Orientation of the reserved source blocks]\label{bp:reserved-factor-orientation-complement}
\lean{Oberwolfach.factor_cycle_orientation_on_component,Oberwolfach.extend_component_cycle_choices_selected,Oberwolfach.factor_cycle_orientation_reserved_block_step,Oberwolfach.reserved_factor_orientation_complement}
\leanok
For a two-regular factor and the displayed injective allocation of positive-length reserved blocks, the already chosen cyclic orders extend to all components. Their cyclic successors define a global vertex permutation which follows the original reserved blocks. It maps the complement of the deleted path-tail vertices, namely the outside-path vertices together with the final endpoints, to the outside-path vertices together with the initial endpoints. The target complement embedding and its arc-image descriptions remain to be constructed.
\end{lemma}
\begin{proof}\leanok Extend the component isomorphisms, glue their cyclic permutations, and apply the buffered path boundary and complement identities. \end{proof}

\begin{lemma}[Canonical oriented residual arcs]\label{bp:canonical-oriented-residual-arcs}
\lean{Oberwolfach.buffered_canonical_oriented_arc_images,Oberwolfach.residual_degrees_from_reserved_factor_arcs}
\leanok
The displayed complement embedding and boundary agreements determine canonical injective tail and head maps for each factor. For the original allocated reserved blocks, these maps have the required target image sets, and the residual directed degree identities follow from the displayed uniform-root and host degree hypotheses, provided the canonical directed arcs lie in the residual host and are pairwise disjoint. The construction of the complement embeddings and these last containment and disjointness properties remain hypotheses.
\end{lemma}
\begin{proof}\leanok Use the reserved-factor permutation and its complement identity to construct the two maps, then apply the buffered-interval degree theorem. \end{proof}

\begin{lemma}[Factor-edge orientation and residual deletion]\label{bp:factor-oriented-residual-deletion}
\lean{Oberwolfach.factor_cycle_orientation_edge_iff,Oberwolfach.factor_orientation_no_two_cycles,Oberwolfach.factor_oriented_physical_edge_injective,Oberwolfach.buffered_factor_remaining_edges,Oberwolfach.directed_arcs_disjoint_of_physical_packing}
\leanok
For a finite two-regular factor, its chosen component successors orient each edge exactly once. When the displayed disjoint buffered walks follow that orientation, the edges remaining after their deletion are exactly the images of the tails outside the walks together with their final endpoints. Pairwise disjoint physical edge sets imply disjointness of directed arc families whose physical images they contain. The existence and packing of the actual residual embeddings remain to be established.
\end{lemma}
\begin{proof}\leanok Use the cycle orders and injectivity of the successor-edge map, the exact walk-edge and tail-domain identities, and monotonicity of disjointness under the physical edge map. \end{proof}

\begin{lemma}[Remaining-edge packing gives canonical arc disjointness]\label{bp:canonical-arcs-remaining-packing}
\lean{Oberwolfach.buffered_canonical_arcs_in_remaining_edge_image,Oberwolfach.buffered_canonical_arcs_disjoint_of_remaining_edge_packing}
\leanok
For the displayed two-regular source factors, buffered walks following their chosen orientations, and complement embeddings, the physical image of each canonical directed arc belongs to the mapped edge set of the corresponding factor after deleting those walks and restricting to the complement of their strict interiors. If these mapped remaining edge sets are pairwise disjoint, then the canonical directed arc sets are pairwise disjoint. The pairwise physical packing and the directed host containment remain hypotheses for the residual degree application.
\end{lemma}
\begin{proof}\leanok Use the exact oriented-edge deletion identity and the complement tail domain, then transfer pairwise physical-edge disjointness to the directed arcs. \end{proof}

\begin{lemma}[Reserved-factor degree identities from physical packing]\label{bp:reserved-degrees-physical-packing}
\lean{Oberwolfach.residual_degrees_from_reserved_factor_packing}
\leanok
For the original allocated reserved blocks and the displayed uniform interval and degree hypotheses, pairwise physical packing of the mapped remaining factor edges supplies the canonical directed-arc disjointness required for the residual in- and out-degree identities. The complement maps, endpoint agreements, directed host containment, and the physical packing remain hypotheses.
\end{lemma}
\begin{proof}\leanok Apply the reserved-block orientation and the transfer from remaining-edge packing to canonical directed-arc disjointness, then use the residual degree theorem. \end{proof}

\begin{lemma}[The retained sample supplies the interval inputs]\label{bp:retained-actual-matching-inputs}
\lean{Oberwolfach.retained_marked_root_membership,Oberwolfach.retained_marked_roots_uniform,Oberwolfach.retained_marked_interval_geometry,Oberwolfach.retained_marked_matching_inputs}
\leanok
For one retained interval sample and its given short-cycle-free permutation, index the final intervals by their retained starts. Their root sets are exactly the marked states of that same retained sample. Every label has exactly $dt$ marked memberships; the orbit segments are injective, the extended intervals are disjoint, the root and endpoint sets are disjoint, and the permutation maps every root into the root or endpoint set. The displayed lower bound on $d$ and short-cycle exclusion are the hypotheses used here.
\end{lemma}
\begin{proof}\leanok Identify roots with the retained interval state, apply the exact retained-load identity, and use the original interval separation and permutation geometry. \end{proof}

\begin{lemma}[Retained-sample residual degrees from restricted packing]\label{bp:retained-residual-degree-packing}
\lean{Oberwolfach.actual_restricted_edge_image_cover,Oberwolfach.actual_restricted_edge_images_packed,Oberwolfach.buffered_canonical_arcs_disjoint_of_actual_old_new_packing,Oberwolfach.retained_sample_residual_degree_identities,Oberwolfach.marked_variable_residual_spoke_balance}
\leanok
For the displayed actual residual data and compatible old, new and complement maps, the mapped factor edges remaining after deletion of the final paths lie in the union of the old edge images with those final edges removed and the new residual-graph edge images. Pairwise packing of those old and new families, together with their stated separation, gives canonical arc disjointness. The retained sample then supplies the interval injectivity, separation and uniform-load conditions in the reserved-factor in- and out-degree identities. The exact spoke-count identity also holds for variable interval lengths. The complement maps and their agreements, directed host containment, the old/new packing, and the residual host remain hypotheses.
\end{lemma}
\begin{proof}\leanok Split source edges according to the actual deletion, transfer physical packing to directed arcs, use the retained-sample interval identities, and count the complementary source vertices. \end{proof}

\begin{lemma}[Gluing on the complement of the final interiors]\label{bp:actual-complement-only-gluing}
\lean{Oberwolfach.actual_residual_embeddings_glue_on_complement,Oberwolfach.actual_final_path_interiors_are_deleted,Oberwolfach.full_embedding_cannot_avoid_nonempty_reserved_set}
\leanok
If the final-path interiors are disjoint from the actual residual source graph, the displayed compatible old and residual maps glue to an injection on their complement. The old map need only be injective on the old domain with those interiors removed. The displayed identification of final paths with actually deleted paths gives this domain separation. Conversely, an injection between equally sized finite source and target sets cannot avoid a nonempty set of target labels on the entire source, so full-source avoidance cannot be used to reserve such labels.
\end{lemma}
\begin{proof}\leanok Restrict the old-domain gluing theorem and its overlap to the complement of the interiors. The final-path support identities place the interiors in the deleted vertices. For the obstruction, an injection between finite sets of the same size is surjective. \end{proof}

\begin{lemma}[Retained roots in the residual spoke set]\label{bp:actual-retained-residual-spokes}
\lean{Oberwolfach.actual_outside_buffered_path_image_avoids_reserved,Oberwolfach.retained_actual_residual_spoke_identities}
\leanok
Under the displayed compatibility and restricted image-avoidance hypotheses, the image of the source vertices outside the final paths avoids the retained marked roots. For the same retained sample and permutation, the marked roots lie in the residual spoke set $S$, $|S|=c|Y|$, and $S\setminus Y$ avoids the successors of $Y$. The construction and compatibility of the embeddings and the displayed old, anchor and free-list avoidance remain hypotheses.
\end{lemma}
\begin{proof}\leanok Use the old/residual source-domain cover only outside the final paths, then apply the retained interval geometry and exact complement cardinality. \end{proof}

\begin{lemma}[Actual residual assignment with complement embeddings]\label{bp:actual-assignment-complement-gluing}
\lean{Oberwolfach.actual_factor_remainder_assigned_and_complement_glued,Oberwolfach.actual_complement_embedding_preserves_reserved_labels,Oberwolfach.actual_final_path_endpoints_prescribed,Oberwolfach.actual_final_endpoints_agree_with_complement}
\leanok
Under the displayed bad-edge assignment, capacity, list and packing hypotheses, the actual residual-source construction yields compatible injections on the complements of the final interiors. It requires injectivity and old-image exclusions only on the old domain with those interiors removed, provided the final interiors are disjoint from the residual source vertices. Separately, the restricted gluing construction preserves reserved labels outside the final paths under the stated old, anchor and free-list avoidance assumptions. The endpoints of the actual final paths are prescribed and belong to both the old and residual domains; agreement of the old map with the prescribed endpoint labels then transfers to the complement map. The verification of these hypotheses from the common event and their compatibility with the chosen source orientation remain open.
\end{lemma}
\begin{proof}\leanok Apply the prescribed-position remainder embedding with the restricted old domain, use the actual old/residual overlap, and restrict the glued injection to the complement of the final interiors. The endpoint memberships follow from the definitions of the prescribed sets. \end{proof}

\begin{lemma}[Ordered regular-host common event with matching inputs]\label{bp:ordered-actual-common-event-matching}
\lean{Oberwolfach.long_host_matching_vortex_alias_decidable,Oberwolfach.long_host_adaptive_shared_choices_ordered_decidable,Oberwolfach.long_host_actual_common_event_with_matching_ordered}
\leanok
The displayed host constants are chosen before the class count and the later event parameters, in the original order. Under the displayed regular-host, factor, density and tolerance hypotheses, the constructed vortex and its preserving short-cycle-free permutation admit the supplied-vortex common outcome together with the retained interval matching identities on the very same outcome. The matching certificate includes the exact retained-root load, orbit injectivity, extended-interval separation, endpoint avoidance and successor containment. The later compatible embedding and packing hypotheses, reference chain, iteration and final decomposition are not proved by this lemma.
\end{lemma}
\begin{proof}\leanok Align the finite-equality instances in the checked host continuation, choose the parameters in order, apply the supplied-vortex specialization, and attach the deterministic retained matching certificate to its existing witness. \end{proof}

\begin{lemma}[Constructed complement and retained spoke identities]\label{bp:retained-complement-spoke-construction}
\lean{Oberwolfach.selected_anchors_avoid_reserved_outside,Oberwolfach.actual_complement_embedding_preserves_reserved_labels_from_anchor_images,Oberwolfach.retained_actual_complement_and_spoke_construction}
\leanok
For the retained interval family, the displayed selected final source paths, and compatible old and residual maps satisfying the stated injection, image, list and avoidance conditions, there is a complement injection agreeing with both maps and with the prescribed final-path endpoints. Its outside-path image avoids the retained marked roots; the resulting residual spoke set $S$ contains those roots $Y$, satisfies $|S|=c|Y|$, and has $(S\setminus Y)\cap\pi(Y)=\varnothing$. The final interiors need not be embedded in this construction. Constructing the compatible maps and verifying the stated image and avoidance hypotheses from the common event remain open.
\end{lemma}
\begin{proof}\leanok Use the actual deletion and endpoint identities, glue only on the complement, derive restricted anchor avoidance from the selected-image identity, and apply the same-sample interval geometry and cardinality calculation. \end{proof}

\begin{lemma}[Source orientations and reserved blocks agree]\label{bp:actual-source-orientation-alignment}
\lean{Oberwolfach.actual_source_orientation_covers_and_has_unique_incoming,Oberwolfach.component_cycle_iso_with_first_edge,Oberwolfach.factor_cycle_orders_align_with_source_pieces,Oberwolfach.actual_factor_remainder_assigned_aligned_complement,Oberwolfach.factor_cycle_orientation_reserved_block_step_from_global}
\leanok
For the displayed two-regular factors and their actual source-piece orientations, the cycle orders of all components can be chosen to agree with those orientations. These orders are chosen before the residual data, and their successors follow the originally allocated reserved blocks. The conclusion uses the edge coverage, pairwise source-piece disjointness and validity already provided by the actual source-piece construction; it does not require a separate orientation-agreement hypothesis.
\end{lemma}
\begin{proof}\leanok The source-piece orientation covers every factor edge and has at most one incoming arc at each vertex. Choose a cyclic ordering beginning with one of its directed edges in each component, then propagate that orientation around the component. Restrict the resulting global orders to the allocated blocks. \end{proof}

\begin{lemma}[Residual degrees with the aligned source orientation]\label{bp:actual-aligned-residual-degrees}
\lean{Oberwolfach.actual_final_path_matches_buffered_sequence,Oberwolfach.retained_actual_complement_from_buffered_parents,Oberwolfach.actual_canonical_arcs_in_host_from_residual_digraph,Oberwolfach.retained_sample_residual_degrees_from_aligned_global_orientation}
\leanok
The actual deleted final paths have the central vertex sequences of their buffered parent walks. Under the displayed old-map, residual-embedding, image-avoidance and packing conditions, those paths give the complement injection and retained spoke identities. If the displayed old directed arcs lie in the host and the new residual digraph embeds there, the canonical complement arcs lie in the host. With the globally aligned cycle orders and the same retained sample, the stated residual in- and out-degree identities follow. The old-arc containment, embedding and physical packing inputs remain to be obtained from the common event.
\end{lemma}
\begin{proof}\leanok Identify the central paths, glue on their complement, split the canonical arcs into old and new source edges, and apply the retained-sample load and packing identities. \end{proof}

\begin{lemma}[Actual component cycles, paths and their index types]\label{bp:actual-oriented-cycle-walks}
\lean{Oberwolfach.forward_cycle_graph_walk,Oberwolfach.factor_component_has_forward_cycle,Oberwolfach.factor_component_family_has_forward_cycles,Oberwolfach.factor_short_cycles_disjoint_from_long_reserved_paths,Oberwolfach.actual_selected_runs_are_disjoint_simple_paths,Oberwolfach.actualResidualDataFromFamilies,Oberwolfach.firstSumIndexEquiv,Oberwolfach.secondSumIndexEquiv}
\leanok
A two-regular factor's selected component cycles may be chosen to follow its globally chosen cycle orders, with the entire component as support and distinct components giving disjoint supports. The original short/long component size bounds make these cycles disjoint from the reserved parent paths. Any such disjoint path and cycle families assemble into actual residual data; the original final-path and selected-path index types are equivalent to the corresponding selected subtypes. The actual selected runs are simple and mutually disjoint.
\end{lemma}
\begin{proof}\leanok Map a forward cyclic-graph walk through the chosen component isomorphism, and use disjoint component supports. Assemble the data from the supplied families and use the established deleted-path properties. \end{proof}

\begin{lemma}[Matching fixed walks and old directed arcs]\label{bp:actual-old-walk-arc-matching}
\lean{Oberwolfach.disjoint_simple_walk_families_match,Oberwolfach.actual_final_path_edge_set_matches_buffered_sequence,Oberwolfach.actual_reindexed_final_edge_union,Oberwolfach.actual_old_arc_host_from_reindexed_buffered_paths}
\leanok
Two families of pairwise disjoint simple walks of equal corresponding lengths admit a partial injection matching every indexed vertex. For the actual residual data, the non-final deleted edges are covered by the selected runs and cycles; the reserved blocks determine the run directions. If the selected target run and cycle images have the displayed directed host arcs, the old directed arcs lie in the host. The equality of final edge unions is preserved by an explicit equivalence between the retained and selected index types.
\end{lemma}
\begin{proof}\leanok Reindex each pair of walks by its vertex sequence and glue the disjoint embeddings. Then split the old deleted edges into selected runs and cycles and use uniqueness of the source orientation. \end{proof}

\begin{lemma}[Selected runs and cycles matched to target walks]\label{bp:selected-runs-targets}
\lean{Oberwolfach.actual_selected_runs_match_target_walks,Oberwolfach.actual_selected_runs_match_directed_target_walks,Oberwolfach.cycle_tail_support_eq,Oberwolfach.disjoint_cycle_families_match}
\leanok
The actual selected runs can be matched, without reallocating their source vertices, to any corresponding family of disjoint simple target walks of the same lengths. If the target walks follow the displayed directed host arcs, the matched runs do too. Likewise, families of disjoint source and target cycles of the same lengths can be matched in their cyclic orders. These statements require the displayed target-walk existence and disjointness conditions.
\end{lemma}
\begin{proof}\leanok Use the actual run disjointness and the fixed-walk matching. A cycle with its first edge removed is a simple path with the same vertex support, so apply the same matching to cycle tails. \end{proof}

\begin{lemma}[The actual old embedding and its list]\label{bp:old-embedding-from-pieces}
\lean{Oberwolfach.disjoint_cycle_families_match_directed,Oberwolfach.actual_old_piece_walks_are_disjoint_simple_paths,Oberwolfach.actual_old_pieces_match_target_walks,Oberwolfach.actual_old_piece_paths_avoid_final_interiors,Oberwolfach.actual_old_domain_without_final_interiors,Oberwolfach.actual_final_interiors_eq_buffered_strict,Oberwolfach.actual_old_embedding_from_pieces_and_anchors_buffered,Oberwolfach.extending_list_avoids_old_image}
\leanok
The selected source runs and cycle tails form one disjoint family and can be matched by one partial injection to corresponding disjoint target walks; directed target cycles retain their closing arcs. The actual old domain outside the final interiors is precisely the union of these pieces and the prescribed anchors, and the actual final interiors agree with the buffered strict interiors after the stated index equivalence. Compatible injective maps on pieces and anchors therefore glue to an injection on this exact domain. An extending list can be made disjoint from its old image by deleting that image and charging its cardinality against the extending-list parameter. Target-walk existence, anchor compatibility and the stated cardinality budget remain to be established from the common event.
\end{lemma}
\begin{proof}\leanok Combine the disjoint source walks, identify the final-interior and old-domain sets, glue the maps on the overlap, and apply stability of extending lists under deletion. \end{proof}

\begin{lemma}[Matching both interval classes in one outcome]\label{bp:mixed-final-approximate-matching}
\lean{Oberwolfach.retained_and_raw_interval_geometry,Oberwolfach.buffered_marker_anchor_embedding,Oberwolfach.mixed_outcome_buffered_anchor_embedding,Oberwolfach.actual_prescribed_anchor_embedding,Oberwolfach.mixed_actual_prescribed_anchor_embedding,Oberwolfach.actual_anchor_values_on_run_endpoints,Oberwolfach.actual_matched_anchor_agreement_from_endpoints,Oberwolfach.mixed_actual_matched_old_embedding}
\leanok
In the same mixed outcome, the retained final intervals and raw approximate intervals form one family with injective, pairwise disjoint extended permutation windows. For disjoint reserved source blocks indexed by this family, their c-spaced markers have one injective map to the original interval labels. Its restriction is the prescribed anchor injection on the actual residual data. After matching the selected old pieces to target walks with the indicated run endpoints, the two maps agree on their overlap. Their remaining image-separation condition is still to be established from the packing.
\end{lemma}
\begin{proof}\leanok Distinct raw start colours cannot have the same start; the original isolation rule separates their extended windows. Match the indexed source markers to those labels, restrict to the prescribed domain, and use that prescribed old-piece vertices are selected-run endpoints. \end{proof}

\begin{lemma}[Directed selected-wheel runs with physical edges]\label{bp:packed-marked-directed-run}
\lean{Oberwolfach.marked_single_wheel_directed_walk_with_edges,Oberwolfach.simple_directed_walk_interval_concatenation_with_edges,Oberwolfach.disjoint_marked_wheels_distinct_spoke_labels,Oberwolfach.packed_marked_wheels_form_orbit_run_with_edges,Oberwolfach.matched_actual_run_preserves_parent_markers}
\leanok
A selected consecutive sequence of physically disjoint marked wheels at one hub forms a simple directed run of the required length. It visits every original interval marker at the corresponding segment boundary, and each physical edge of the run belongs to one of the selected wheels. The stated terminal-label avoidance follows from injectivity of the extended permutation interval and exclusion of ordinary spoke labels. Matching this run to its actual source run preserves the original parent markers. The following lemma derives the required wheel family and separation between different resulting runs from the packing.
\end{lemma}
\begin{proof}\leanok Read each wheel as its directed physical rim walk, use physical spoke disjointness for different rims, and concatenate only the selected walks while retaining their indexed boundaries and edge sets. \end{proof}

\begin{lemma}[Mixed packing and anchor embedding]\label{bp:mixed-packing-buffered-embedding}
\lean{Oberwolfach.actual_selected_segments_choose_wheels_from_packing_positions,Oberwolfach.actual_selected_run_family_from_packed_wheels,Oberwolfach.ordinary_packed_wheel_directed_path_with_edges,Oberwolfach.actual_mixed_packing_target_pieces,Oberwolfach.actual_matched_old_embedding_without_cross_hypothesis,Oberwolfach.actual_mixed_packing_buffered_embedding}
\leanok
For the stated actual residual data and physical mixed packing at one hub, assume the selected positions are those with marked roots covered by that packing, the ordinary spokes avoid the extended intervals, and each assigned ordinary wheel has the length of its assigned source cycle. The marked wheels construct all selected target runs together; the ordinary wheels construct the target cycle tails, including their closing arcs. These pieces have pairwise disjoint supports and retain their physical rim edges and directions. Together with the original prescribed marker map they give an injection on the buffered old domain, with the source-run marker positions preserved. The physical packing and assignment hypotheses, their connection to the common random outcome, and the subsequent residual completion are still to be supplied in the global theorem.
\end{lemma}
\begin{proof}\leanok Choose wheels by the distinct interval labels, join consecutive marked rims, use physical spoke disjointness to separate different pieces, and trace every possible anchor collision to its unique source marker.\end{proof}

\begin{lemma}[Matching-to-packing and physical source edges]\label{bp:matching-to-physical-packing}
\lean{Oberwolfach.retained_and_raw_intervals_avoid_ordinary_spokes,Oberwolfach.packed_physical_spoke_positions_eq_marked_roots,Oberwolfach.mixed_outcome_physical_packing_buffered_embedding,Oberwolfach.long_raw_matching_gives_hub_packing,Oberwolfach.long_raw_assigned_wheel_packing,Oberwolfach.actual_old_arc_host_from_matched_pieces,Oberwolfach.actual_matched_piece_physical_edges_from_wheels,Oberwolfach.physical_disjointness_gives_symmetric_packing}
\leanok
In the same raw interval outcome, ordinary spokes avoid the retained and raw extended intervals, and covered physical spoke positions agree with the marked-wheel roots. The mixed matching gives physically disjoint marked and ordinary wheel families at each hub, and reindexing the assigned ordinary wheels preserves their sampled directions and spokes. Matching the source pieces transfers both directions and physical edge provenance, including each ordinary cycle's closing edge. Pairwise disjoint physical rim families then give the symmetric-edge packing hypotheses. These statements use the displayed matching, source assignment and orientation hypotheses; the global residual embedding, numerical bounds and reconstruction remain open.
\end{lemma}
\begin{proof}\leanok Unpack the original sampled resource states and the physical matching, reindex the assigned wheels, and transfer the matched walk edges to the original source pieces.\end{proof}

\begin{lemma}[Assigned sources and the old embedding at one hub]\label{bp:assigned-source-old-embedding}
\lean{Oberwolfach.global_assignment_reserved_source_families,Oberwolfach.actual_all_approx_markers_preserved,Oberwolfach.actual_approx_intervals_occupied,Oberwolfach.actualResidualDataFromPacking,Oberwolfach.long_raw_assigned_source_embedding,Oberwolfach.long_raw_physical_hub_rims_packed,Oberwolfach.long_raw_selected_wheel_rims_covered}
\leanok
Fix one global orientation of the factor and an injective global assignment of the selected ordinary wheels to components of the required lengths. At a hub, assuming the displayed long-component capacity and the short/long length conditions, one can choose disjoint buffered source paths and the assigned short cycles, define the selected positions from the same mixed matching, and construct the buffered old-domain injection. It preserves all prescribed markers, occupies every label in each extended approximate interval, and sends the old directed arcs to the same sampled physical host. The physical rim sets associated with different hubs are disjoint, and the selected wheels' rims lie in their respective hub sets. Deriving the quantitative capacity and all hypotheses simultaneously from the common event, completing the residual embedding and global theorem remain open.
\end{lemma}
\begin{proof}\leanok Restrict the global component assignment to a hub, allocate buffered blocks in the disjoint oriented long cycles, form the actual residual record, and apply the matching construction and marker-survival lemmas. Physical rim disjointness follows from the original matching of resources.\end{proof}


\begin{lemma}[Physical old-edge packing and interval demand]\label{bp:physical-old-edge-capacity}
\lean{Oberwolfach.long_raw_assigned_source_embedding,Oberwolfach.long_raw_matching_whole_old_packing,Oberwolfach.retained_raw_interval_count_le_raw_counts,Oberwolfach.long_assigned_actual_interval_capacity,Oberwolfach.long_previous_spacing_exceeds_cutoff,Oberwolfach.long_compatible_parameter_choices,Oberwolfach.long_assigned_adaptive_fixed_spacing_canonical}
\leanok
For the same mixed matching and the displayed component-assignment and source-orientation hypotheses, the actual old source edges lie in the matching's physical hub rims. Hence the old symmetric edge sets are pairwise disjoint across hubs. The actual retained and approximate interval demand is bounded by the two sampled start counts, and the original long-class hypothesis supplies the component mass needed for the floor-capacity bound, provided the displayed cutoff and large-order inequalities hold. The earlier auxiliary matching tolerances force spacing to exceed the cutoff; a separate fixed-spacing parameter choice and fixed-spacing canonical continuation are checked. Carrying the compatible choice through the whole common event, proving its start-count hypotheses on the same outcome, and completing the residual embedding and main theorem remain open.
\end{lemma}
\begin{proof}\leanok Transfer source-piece edges to physical rims, use disjointness of matching resources, count actual interval indices by their raw starts, and allocate blocks using the long-component mass.\end{proof}

\begin{lemma}[Compatible common event and source-path capacity]\label{bp:compatible-event-source-paths}
\lean{Oberwolfach.long_host_compatible_common_event_with_starts_matching,Oberwolfach.long_padded_component_maps_from_retained_hub_matching_tests,Oberwolfach.long_raw_source_embeddings_from_matching_and_starts,Oberwolfach.long_raw_matching_actual_old_packing,Oberwolfach.long_capacity_cutoff_succ,Oberwolfach.long_capacity_large_order}
\leanok
The regular-host common event can be chosen with the interval spacing fixed before the long-cycle cutoff, and with the raw hub-start counts holding on the same retained outcome as the matching. Its selected matching tests give the ordinary component maps. For one matching, the original long-class hypothesis, those component maps and the same-outcome start counts yield a single global component assignment and buffered source embeddings at all hubs. Their old physical edges are packed across hubs. The cutoff and large-order inequalities are available from the compatible parameter choice and sufficiently large order. The remaining simultaneous residual embeddings, reference-chain compatibility, iteration, absorber and reconstruction are not supplied by this lemma.
\end{lemma}
\begin{proof}\leanok Intersect the original prior with the high-probability start-count event, select the matching on that witness, and use the retained matching tests and component counts to allocate the buffered paths. The physical resource matching gives the old-edge packing.\end{proof}

\begin{lemma}[Protected unassigned lists from physical-neighbourhood tests]\label{bp:protected-neighbour-list}
\lean{Oberwolfach.actual_matched_old_image_avoids_list,Oberwolfach.long_raw_source_embeddings_from_matching_and_starts_with_protected_list,Oberwolfach.raw_unassigned_physical_neighbour_tests_high_probability,Oberwolfach.raw_neighbour_test_concentration_protected_list}
\leanok
The old source embedding obtained from the selected matching avoids the concrete unassigned list after the extended interval endpoints are removed. Neighbourhood tests may depend on the independently sampled physical-edge coordinates, and the unassigned counts on these tests hold with arbitrarily high probability. If the displayed neighbourhood-density, probability, concentration-error and endpoint-loss inequalities hold, this protected list is an extending list. Establishing those inequalities for the chosen common outcome and constructing the residual embedding remain separate obligations.
\end{lemma}

\begin{lemma}[Selected free host and protected extending lists]\label{bp:selected-free-host-lists}
\lean{Oberwolfach.long_host_compatible_selected_tests_with_list_margins,Oberwolfach.long_raw_matching_actual_old_packing_in_approximate_layer,Oberwolfach.long_selected_extending_lists_from_event}
\leanok
The compatible retained common event includes the matching, the hub start counts, the selected-layer unassigned neighbourhood counts and the selected-layer directed two-path bounds. The stated positive list and deletion tolerances and a sufficiently large spacing can be chosen under the original parameter assumptions, after which the displayed numerical bounds hold for sufficiently large order. The old symmetric-edge packing obtained from the matching is contained in the approximate physical layer. For an outcome satisfying the stated events and numerical bounds, the concrete protected unassigned lists extend in the selected host after a sparse physical deletion, and this host avoids the approximate packing. The simultaneous residual embedding and the later iteration remain separate obligations.
\end{lemma}

\begin{lemma}[Long-class matching lift]\label{bp:long-class-matching-lift}
\lean{Oberwolfach.long_padded_maps_and_packing_from_long_class_tests,Oberwolfach.long_raw_matching_actual_old_packing_avoids_selected_host}
\leanok
Under the stated long-class and retained matching tests, the matching lifted from the surviving resources has the required component maps and pairwise disjoint resources, and its physical edges avoid the deleted graph. The corresponding actual old-edge packing lies in the approximate layer and is disjoint from the selected host. The subsequent construction of simultaneous buffered complement embeddings is not asserted here.
\end{lemma}

\begin{lemma}[Buffered complement families with aligned cycle orders]\label{bp:buffered-complement-families}
\lean{Oberwolfach.actual_factor_dependent_complement_with_aligned_orders}
\leanok
Under the displayed capacity, sparse-bad-edge, extending-list, protected-old-image and packing hypotheses, a common orientation and aligned component orders support the simultaneous assignment and embedding of the residual graphs, glued to the prescribed old images on the complements of the reserved interiors. The resulting complement edge images are pairwise disjoint and avoid the specified old packing. These hypotheses remain to be derived for the chosen common outcome.
\end{lemma}

\begin{lemma}[Pairwise protected availability]\label{bp:pairwise-protected-availability}
\lean{Oberwolfach.raw_interval_two_unassigned_probability_from_slack,Oberwolfach.raw_unassigned_pairs_common_event,Oberwolfach.raw_interval_start_columns_event,Oberwolfach.raw_protected_pair_hubs_lower_from_counts,Oberwolfach.protected_pair_edge_hubs_avoid_old_images,Oberwolfach.protected_pair_edge_hubs_card_lower}
\leanok
For each distinct pair, the raw interval sample has the stated lower bound on its joint unassigned probability, without requiring independence between the two vertices. Simultaneous pair counts and start-column bounds yield the displayed protected-list lower bound; assigning an edge only to outer hubs containing both endpoints in their protected lists avoids the old image. Joining these bounds to the main common event is a separate obligation.
\end{lemma}

\begin{lemma}[Protected pairs on the retained common outcome]\label{bp:protected-pairs-common-outcome}
\lean{Oberwolfach.long_host_compatible_protected_pairs_with_list_margins,Oberwolfach.long_protected_pair_matching_hypotheses,Oberwolfach.Vortex.inner_hubs_margin}
\leanok
The regular-host outcome can be chosen to satisfy the original selected-host and matching tests together with the protected-pair count. The spacing and order margins can be selected from the original positive parameters. If the displayed vortex contraction and old-image avoidance bounds hold, the protected lists supply linearly many outer labels for each bad edge and discharge the old-image shielding and forbidden-hub conditions. The residual degree, capacity, multiplicity, and absorber-list hypotheses are not asserted by this lemma.
\end{lemma}

\begin{lemma}[A single protected source family]\label{bp:protected-source-family}
\lean{Oberwolfach.long_raw_protected_source_family_and_old_packing,Oberwolfach.actual_sum_buffered_interiors_deleted_with_parameters,Oberwolfach.protected_list_shields_absorber}
\leanok
For a matching with the stated component-map and start-count hypotheses, one component assignment yields a simultaneous family of actual old embeddings. Their reserved strict interiors are disjoint from the residual domains, their restricted old images avoid the corresponding protected lists, and the old physical edge images pack pairwise. If an absorber target lies in a protected list, it avoids that old image. The further residual and final-path hypotheses remain separate.
\end{lemma}

\begin{lemma}[Cycle orders before matching-dependent indices]\label{bp:late-indices-continuation}
\lean{Oberwolfach.actual_factor_late_indices_complement_with_aligned_orders,Oberwolfach.long_raw_aligned_protected_markers_and_selected_host,Oberwolfach.actual_late_complement_with_protected_pair_matching}
\leanok
The common orientations and aligned cycle orders can be fixed before the matching-dependent residual index types are selected. The same orders support the actual protected old source family and the later residual continuation, whose matching-label and old-shield conditions follow from the protected pairs. The stated capacity, residual-size, load and absorber-target hypotheses remain to be verified.
\end{lemma}

\begin{lemma}[Interval occupancy and prescribed load]\label{bp:interval-occupancy-prescribed-load}
\lean{Oberwolfach.long_host_compatible_occupancy_pairs_with_list_margins,Oberwolfach.actual_prescribed_multiplicity_le_extended_hubs,Oberwolfach.actual_prescribed_load_from_sampled_intervals,Oberwolfach.long_prescribed_load_margins}
\leanok
The selected common outcome also bounds the number of hubs whose extended sampled intervals contain any given vertex. Every actual prescribed old image is such a vertex of its hub's interval, including the terminal-point case. With the stated path-parameter and order margins, and separate sparse bad-edge and absorber-load bounds, this gives the prescribed-image load inequality required by the residual continuation.
\end{lemma}

\end{document}
